% Bourbaki schemes — page-by-page transcription
% Model: gemini-3.1-flash-lite-preview
% Markers written by the pipeline from the PDF page index
% Pages transcribed: 437/437
% Rebuilt: 2026-07-05

%% ===== Page 1 =====

\marginpar{Archives}
\marginpar{Grothendieck sept. 59}
\hfill \text{u\textsuperscript{o} 326 bis}

\begin{center}
CHAPITRE 0
PRÉLIMINAIRES.
\end{center}

\section*{§1. Algèbre commutative.}

Nous rappelons dans ce paragraphe certaines définitions et certains résultats d'algèbre commutative, dont nous nous servirons souvent par la suite. Nous omettrons la plupart des démonstrations, renvoyant par exemple aux chap. I et II de l'Algèbre commutative de N. Bourbaki.

Sauf mention expresse du contraire, tous les anneaux considérés seront commutatifs et munis d'un élément unité, et tous les modules sur un tel anneau seront supposés unitaires. Les homomorphismes d'anneaux seront toujours supposés transformer l'élément unité en élément unité, et sauf mention expresse du contraire, un sous-anneau d'un anneau $A$ sera supposé contenir l'élément unité de $A$. Un anneau intègre est un anneau dans lequel le produit d'une famille finie d'éléments $\neq 0$ est $\neq 0$; il revient au même de dire que dans un tel anneau on a $0 \neq 1$ et le produit de deux éléments $\neq 0$ est non nul. Un idéal premier $\mathfrak{p}$ d'un anneau $A$ est un idéal tel que $A/\mathfrak{p}$ soit intègre ; cela entraîne donc $\mathfrak{p} \neq A$. Pour qu'un anneau ait au moins un idéal premier, il faut et il suffit que $A \neq \{0\}$.

\section*{1. Racine d'un idéal.}

Soit $\mathfrak{a}$ un idéal d'un anneau $A$; la racine de $\mathfrak{a}$, notée $r(\mathfrak{a})$, est l'ensemble des $x \in A$ tels que $x^n \in \mathfrak{a}$ pour un entier $n > 0$ au moins; c'est un idéal contenant $\mathfrak{a}$. On a $r(r(\mathfrak{a})) = r(\mathfrak{a})$; la relation $\mathfrak{a} \subset \mathfrak{b}$ entraîne $r(\mathfrak{a}) \subset r(\mathfrak{b})$; la racine d'une intersection d'idéaux est l'intersection de leurs racines. Si $\varphi$ est un homomorphisme d'un anneau $A'$ dans $A$, on a $r(\varphi^{-1}(\mathfrak{a})) = \varphi^{-1}(r(\mathfrak{a}))$ pour tout idéal $\mathfrak{a} \subset A$. Pour qu'un idéal soit racine d'un idéal, il faut et il suffit qu'il soit intersection d'idéaux premiers. La racine d'un idéal est l'intersection des idéaux premiers minimaux parmi ceux qui contiennent $\mathfrak{a}$; si $A$ est noethérien, ces idéaux premiers minimaux sont en nombre fini.

\newpage

%% ===== Page 2 =====

\marginpar{ii}

2. \textbf{Modules et anneaux de fractions.}

On dit qu'une partie $S$ d'un anneau $A$ est multiplicative si $1 \in S$ et si le produit de deux éléments de $S$ est dans $S$. Les exemples qui seront les plus importants pour la suite sont: $1^\circ$ l'ensemble des puissances $f^n$ ($n \geq 0$) d'un élément $f \in A$; $2^\circ$ le complémentaire $A - \mathfrak{p}$ d'un idéal premier $\mathfrak{p}$ de $A$.

Soient $S$ un ensemble multiplicatif dans $A$, $M$ un $A$-module; dans l'ensemble $M \times S$, la relation entre couples $(m_1, s_1), (m_2, s_2)$:

"il existe $s \in S$ tel que $s(s_1m_2 - s_2m_1) = 0$"

est une relation d'équivalence. On désigne par $S^{-1}M$ l'ensemble quotient de $M \times S$ par cette relation, par $m/s$ l'image canonique dans $S^{-1}M$ du couple $(m, s)$; on appelle application canonique de $M$ dans $S^{-1}M$ l'application $i_M^S: m \to m/1$. Cette application n'est en général ni injective ni surjective; son noyau est l'ensemble des $m \in M$ tels qu'il existe un $s \in S$ pour lequel $sm = 0$.

Dans $S^{-1}M$ on définit une loi de groupe additif en prenant $m_1/s_1 + m_2/s_2 = (s_2m_1 + s_1m_2)/s_1s_2$ (on vérifie que c'est bien indépendant des expressions des éléments de $S^{-1}M$ considérés). Sur $S^{-1}A$ on définit aussi une loi multiplicative $(a_1/s_1)(a_2/s_2) = (a_1a_2)/(s_1s_2)$, et enfin une loi externe sur $S^{-1}M$ ayant $S^{-1}A$ comme ensemble d'opérateurs en posant $(a/s)(m/s') = (am)/(ss')$. On vérifie ainsi que $S^{-1}A$ est muni d'une structure d'anneau (dit \emph{anneau des fractions de $A$ à dénominateurs dans $S$}) et $S^{-1}M$ d'une structure de $S^{-1}A$-module (dit \emph{module des fractions de $M$ à dénominateurs dans $S$}).

L'anneau de fractions $S^{-1}A$ et l'application canonique $i_A^S$ sont solution d'un problème d'application universelle: tout homomorphisme $u$ de $A$ dans un anneau $B$ tel que $u(S)$ se compose d'éléments inversibles dans $B$ se factorise d'une seule manière en

\begin{tikzcd}
A \arrow[r, "i_A^S"] \arrow[rr, bend right=30, "u"'] & S^{-1}A \arrow[r, "u^*"] & B
\end{tikzcd}

où $u^*$ est un homomorphisme de l'anneau $S^{-1}A$ dans l'anneau $B$.

On définit un isomorphisme canonique $S^{-1}A \otimes_A M \xrightarrow{\sim} S^{-1}M$ en faisant cor...

\newpage

%% ===== Page 3 =====

- iii -

repondre à l'élément $(a/s) \otimes m$ l'élément $(am)/s$, l'isomorphisme réciproque appliquant $m/s$ sur $(1/s) \otimes m$.

Notons enfin que l'application $\mathfrak{p} \to (i_A)^{-1}(\mathfrak{p}')$ est un isomorphisme pour la structure d'ordre de l'ensemble des idéaux premiers de $S^{-1}A$ sur l'ensemble des idéaux premiers de $A$ qui ne rencontrent pas $S$.

3. \emph{Propriétés fonctorielles}.

Soient $M, N$ deux $A$-modules, $u$ un $A$-homomorphisme de $M$ dans $N$. Si $S$ est une partie multiplicative de $A$, on définit un $S^{-1}A$-homomorphisme de $S^{-1}M$ dans $S^{-1}N$, noté $S^{-1}u$, en posant $(S^{-1}u)(m/s) = u(m)/s$; si $S^{-1}M$ et $S^{-1}N$ sont identifiés canoniquement à $S^{-1}A \otimes_A M$ et $S^{-1}A \otimes_A N$, $S^{-1}u$ est identifié à $1 \otimes u$. Si $P$ est un troisième $A$-module, $v$ un $A$-homomorphisme de $N$ dans $P$, on a $S^{-1}(v \circ u) = (S^{-1}v) \circ (S^{-1}u)$; autrement dit, $M \to S^{-1}M$ est un foncteur covariant de la catégorie des $A$-modules dans celle des $S^{-1}A$-modules ($A$ et $S$ étant fixés). En outre ce foncteur est \emph{exact}, autrement dit, si la suite
$$M \xrightarrow{u} N \xrightarrow{v} P$$
est exacte, il en est de même de la suite
$$S^{-1}M \xrightarrow{S^{-1}u} S^{-1}N \xrightarrow{S^{-1}v} S^{-1}P.$$

Soit $(M_\alpha, \phi_{\beta\alpha})$ un système inductif de $A$-modules; alors $(S^{-1}M_\alpha, S^{-1}\phi_{\beta\alpha})$ est un système inductif de $S^{-1}A$-modules. Exprimant les $S^{-1}M_\alpha$ et $S^{-1}\phi_{\beta\alpha}$ comme des produits tensoriels, il résulte de la permutabilité des opérations de produit tensoriel et de limite inductive, que l'on a un isomorphisme canonique
$$S^{-1} \varinjlim M_\alpha \xrightarrow{\sim} \varinjlim S^{-1}M_\alpha.$$

Soient $M, N$ deux $A$-modules; il existe un $S^{-1}A$ isomorphisme canonique fonctoriel
$$(S^{-1}M) \otimes_{S^{-1}A} (S^{-1}N) \xrightarrow{\sim} S^{-1}(M \otimes_A N)$$
qui transforme $(m/s) \otimes (n/t)$ en $(m \otimes n)/st$. Utilisant une résolution projective de l'un des modules $M, N$, et le fait que $M \to S^{-1}M$ est un foncteur exact, on en conclut des isomorphismes canoniques fonctoriels (en $M$ et $N$)

\newpage

%% ===== Page 4 =====

- iv -

$$\operatorname{Tor}_i^{S^{-1}A}(S^{-1}M, S^{-1}N) \cong S^{-1}\operatorname{Tor}_i^A(M, N) \quad (i \ge 0).$$

On a de même un homomorphisme fonctoriel (en $M$ et $N$)

$$S^{-1}\operatorname{Hom}_A(M, N) \to \operatorname{Hom}_{S^{-1}A}(S^{-1}M, S^{-1}N)$$

qui, à $u/s$ fait correspondre l'homomorphisme $m/t \mapsto u(m)/st$. Lorsque $M$ est le conoyau d'un homomorphisme $A^p \to A^q$, l'homomorphisme précédent est un isomorphisme; c'est immédiat lorsque $M$ est de la forme $A^r$, et on passe de là au cas général en partant de la suite exacte $A^p \to A^q \to M \to 0$, et utilisant l'exactitude à droite des foncteurs $M \to S^{-1}M$ et $X \to \operatorname{Hom}(X, N)$. On notera que ce cas se présente toujours lorsque $A$ est noethérien et le $A$-module $M$ de type fini; en outre, sous ces hypothèses, en utilisant une résolution projective de $M$ par des $A$-modules de type fini, on déduit de l'isomorphisme précédent des isomorphismes fonctoriels

$$S^{-1}\operatorname{Ext}_A^i(M, N) \xrightarrow{\sim} \operatorname{Ext}_{S^{-1}A}^i(S^{-1}M, S^{-1}N) \quad (i \ge 0).$$

4. \emph{Changement de partie multiplicative.}

Soient $S, T$ deux parties multiplicatives d'un anneau $A$ telles que $S \subset T$; il existe un homomorphisme canonique $\rho_A^{T, S}$ (ou simplement $\rho^{T, S}$) de $S^{-1}A$ dans $T^{-1}A$, faisant correspondre à l'élément $a/s$ de $S^{-1}A$ l'élément noté $a/s$ dans $T^{-1}A$; on a $i_A^T = \rho_A^{T, S} \circ i_A^S$. Pour tout $A$-module $M$, il existe de même une application $S^{-1}A$-linéaire de $S^{-1}M$ dans $T^{-1}M$ (ce dernier étant considéré comme $S^{-1}A$-module grâce à l'homomorphisme $\rho_A^{T, S}$), faisant correspondre à l'élément $m/s$ de $S^{-1}M$ l'élément $m/s$ de $T^{-1}M$; on note cette application $\rho_M^{T, S}$ ou simplement $\rho_M^{T, S}$, et on a encore $i_M^T = \rho_M^{T, S} \circ i_M^S$. En outre, $\rho^{T, S}$ est un morphisme fonctoriel (ou transformation naturelle) du foncteur $M \to S^{-1}M$ dans le foncteur $M \to T^{-1}M$, autrement dit le diagramme

\begin{tikzcd}
S^{-1}M \arrow[r, "S^{-1}u"] \arrow[d, "\rho_M^{T, S}"] & S^{-1}N \arrow[d, "\rho_N^{T, S}"] \\
T^{-1}M \arrow[r, "T^{-1}u"] & T^{-1}N
\end{tikzcd}

\newpage

%% ===== Page 5 =====

est commutatif, pour tout homomorphisme $u: M \to N$.

Si $S, T, U$ sont trois parties multiplicatives de $A$ telles que $S \subset T \subset U$ on a
$$(1) \quad \rho^{U,S} = \rho^{U,T} \rho^{T,S}.$$

Considérons maintenant une famille filtrante croissante $(S_\alpha)$ de parties multiplicatives de $A$ (on écrira $\alpha \leq \beta$ pour $S_\alpha \subset S_\beta$), et soit $S$ la partie multiplicative $\bigcup S_\alpha$; posons $\rho_{\beta\alpha} = \rho_A^{S_\beta, S_\alpha}$; en vertu de (1), les homomorphismes $\rho_{\beta\alpha}$ définissent un anneau $A'$ limite inductive du système inductif d'anneaux $(S_\alpha^{-1}A, \rho_{\beta\alpha})$. Soient $\rho_\alpha$ l'application canonique de $S_\alpha^{-1}A$ dans $A'$, et posons $\varphi_\alpha = \rho_A^{S, S_\alpha}$; comme $\varphi_\alpha = \varphi_\beta \circ \rho_{\beta\alpha}$ pour $\alpha \leq \beta$ d'après (1), on peut définir un homomorphisme $\varphi: A' \to S^{-1}A$ tel que le diagramme

$$\begin{tikzcd}
& S_\alpha^{-1}A \arrow[d, "\rho_\alpha"] \arrow[dr, "\varphi_\alpha"] & \\
A' \arrow[r, "\varphi"] & S_\beta^{-1}A \arrow[r, "\rho_{\beta\alpha}"'] & S^{-1}A
\end{tikzcd}$$
soit commutatif. En fait, $\varphi$ est un isomorphisme : il est en effet immédiat par construction que $\varphi$ est surjectif. D'autre part, si $\rho_\alpha(a/s_\alpha) \in A'$ est tel que $\varphi(\rho_\alpha(a/s_\alpha)) = 0$, cela signifie que $a/s_\alpha = 0$ dans $S^{-1}A$, c'est-à-dire qu'il existe $s \in S$ tel que $sa = 0$ ; mais il y a un $\beta \geq \alpha$ tel que $s \in S_\beta$, et par suite, comme $\rho_\alpha(a/s_\alpha) = \rho_\beta(sa/ss_\alpha) = 0$, on voit que $\varphi$ est injectif. On traite de même le cas d'un $A$-module $M$, et on a ainsi défini des isomorphismes canoniques
$$\varinjlim S_\alpha^{-1}A \xrightarrow{\sim} (\varinjlim S_\alpha)^{-1}A, \quad \varinjlim S_\alpha^{-1}M \xrightarrow{\sim} (\varinjlim S_\alpha)^{-1}M$$
le second étant fonctoriel.

Il y a un cas important dans lequel l'homomorphisme $\rho^{T,S}$ est bijectif, savoir lorsque tout élément de $T$ est diviseur d'un élément de $S$; on identifie alors par $\rho^{T,S}$ les modules $S^{-1}M$ et $T^{-1}M$. On dit que $S$ est saturé si tout diviseur d'un élément de $S$ est dans $S$; en remplaçant $S$ par l'ensemble $T$ de tous les diviseurs des éléments de $S$ (qui est multiplicatif et saturé), on voit qu'on peut toujours si l'on veut se limiter à la considération de modules $S^{-1}M$ où

\newpage

%% ===== Page 6 =====

$S$ est saturé.

Considérons enfin deux parties multiplicatives $S_1, S_2$ de $A$; alors $S_1S_2$ est aussi une partie multiplicative de $A$. Désignons par $S_2'$ l'image canonique de $S_2$ dans l'anneau $S_1^{-1}A$, qui est une partie multiplicative de cet anneau. Il existe alors pour tout $A$-module $M$ un isomorphisme fonctoriel de $S_2'^{-1}(S_1^{-1}M)$ sur $(S_1S_2)^{-1}M$, qui fait correspondre à $(m/s_1)/(s_2/1)$ l'élément $m/s_1s_2$.

5. Changement d'anneau.

Soient $A, A'$ deux anneaux, $\varphi$ un homomorphisme de $A'$ dans $A$, $S$ (resp. $S'$) une partie multiplicative de $A$ (resp. $A'$), telle que $\varphi(S') \subset S$; il est immédiat qu'on définit un homomorphisme $\varphi^S$ de $S'^{-1}A'$ dans $S^{-1}A$ en posant $\varphi^S(a'/s') = \varphi(a')/\varphi(s')$.

Soit maintenant $M$ un $A$-module; l'homomorphisme $\varphi$ définit sur $M$ une structure de $A'$-module, en posant $a' \cdot m = \varphi(a')m$; on désignera ce $A'$-module par $M_{[\varphi]}$. Cela étant, il existe un homomorphisme fonctoriel

$$\sigma^S: S'^{-1}(M_{[\varphi]}) \to (S^{-1}M)_{[\varphi^S]}$$

de $S'^{-1}A'$-modules, faisant correspondre à tout élément $m/s'$ de $S'^{-1}(M_{[\varphi]})$ l'élément $\varphi(s')^{-1}m$ de $(S^{-1}M)_{[\varphi^S]}$; on vérifie immédiatement en effet que cette définition ne dépend pas de l'expression $m/s'$ de l'élément considéré. Si $T$ (resp. $T'$) est une seconde partie multiplicative de $A$ (resp. $A'$) telle que $S \subset T$ (resp. $S' \subset T'$) et $\varphi(T') \subset T$, les diagrammes

\begin{tikzcd}
S'^{-1}A' \arrow[r, "\varphi^S"] \arrow[d, "\rho^{T', S'}"'] & S^{-1}A \arrow[d, "\rho^{T, S}"] \\
T'^{-1}A' \arrow[r, "\varphi^T"] & T^{-1}A
\end{tikzcd}
et
\begin{tikzcd}
S'^{-1}(M_{[\varphi]}) \arrow[r, "\sigma^S"] \arrow[d, "\rho^{T', S'}"'] & (S^{-1}M)_{[\varphi^S]} \arrow[d, "\rho^{T, S}_{[\varphi^T]}"] \\
T'^{-1}(M_{[\varphi]}) \arrow[r, "\sigma^T"] & (T^{-1}M)_{[\varphi^T]}
\end{tikzcd}

sont commutatifs. En outre, si $S = \varphi(S')$, l'homomorphisme $\sigma^S$ est bijectif.

Considérons maintenant un $A'$-module $N'$; formons le produit tensoriel $N' \otimes_{A'} A[_{\varphi}]$, qui peut être considéré comme un $A$-module en posant $a \cdot (n' \otimes b) = n' \otimes (ab)$. Il existe alors un isomorphisme fonctoriel de $S^{-1}A$-modules

$$\tau^S: (S'^{-1}N') \otimes_{S'^{-1}A'} (S^{-1}A)[\varphi^S] \xrightarrow{\sim} S^{-1}(N' \otimes_{A'} A[\varphi])$$

\newpage

%% ===== Page 7 =====

qui, à l'élément $(n'/s') \otimes (a/s)$ fait correspondre l'élément $(n' \otimes a)/\varphi(s')s$; on vérifie en effet séparément que lorsqu'on remplace $n'/s'$ (resp. $a/s$) par une autre expression du même élément, $(n' \otimes a)/\varphi(s')s$ ne change pas; d'autre part, on peut définir un homomorphisme réciproque de $\tau^S$ en faisant correspondre à $(n' \otimes a)/s$ l'élément $(n'/1) \otimes (a/s)$ (on utilise là le fait ($n^o 2$) que $S^{-1}(N' \otimes_{A'} A[\varphi])$ est canoniquement isomorphe à $(N' \otimes_{A'} A[\varphi]) \otimes_A S^{-1}A$ donc à $N' \otimes_{A'} (S^{-1}A)[\varphi]$, en désignant par $\psi$ l'homomorphisme $a' \to \varphi(a')/1$ de $A'$ dans $S^{-1}A$).

Si $T$ (resp. $T'$) est une seconde partie de $A$ (resp. $A'$) telle que $S \subset T$ (resp. $S' \subset T'$) et $\varphi(T') \subset T$, le diagramme

\begin{tikzcd}
(S^{-1}N') \otimes_{S^{-1}A'} (S^{-1}A)[\varphi^S] \arrow[r, "\tau^S"] \arrow[d] & S^{-1}(N' \otimes_{A'} A[\varphi]) \arrow[d] \\
(T^{-1}N') \otimes_{T^{-1}A'} (T^{-1}A)[\varphi^{T'}] \arrow[r, "\tau^T"] & T^{-1}(N' \otimes_{A'} A[\varphi])
\end{tikzcd}

est commutatif (la flèche verticale de gauche correspondant à l'homomorphisme obtenu en appliquant $\rho^{T',S'}$ à $S^{-1}N'$ et $S^{-1}A'$, $\rho^{T,S}$ à $S^{-1}A$).

6. \emph{Identification du module $M_f$ à une limite inductive.}

Soient $M$ un $A$-module, $f$ un élément de $A$, et considérons la partie multiplicative $S$ de $A$ formée des $f^n$ ($n \geq 0$); nous désignerons le module des fractions $S^{-1}M$ par la notation abrégée $M_f$. Considérons une suite $(M_n)$ de $A$-modules, tous identiques à $M$, et pour tout couple d'entiers $m \leq n$, soit $\varphi_{nm}$ l'homomorphisme $x \to f^{n-m}x$ de $M_m$ dans $M_n$; il est immédiat que $((M_n), (\varphi_{nm}))$ est un système inductif de $A$-modules; soit $N = \varinjlim M_n$ la limite inductive de ce système. Nous allons définir un $A$-isomorphisme fonctoriel de $N$ sur $M_f$. Pour cela, remarquons que, pour tout $n$, $\theta_n : x \to x/f^n$ est un $A$-homomorphisme de $M = M_n$ dans $M_f$; et il résulte des définitions que l'on a $\theta_m = \theta_n \circ \varphi_{nm}$ pour $m \leq n$. Il existe donc un $A$-homomorphisme $\theta : N \to M_f$ tel que, si $\psi_n$ désigne l'homomorphisme canonique $M_n \to N$, on ait $\theta_n = \theta \circ \psi_n$ pour tout $n$. Comme par

\newpage

%% ===== Page 8 =====

-- viii --

hypothèse tout élément de $M_{\mathfrak{p}}$ est de la forme $z/f^n$ pour un $n$ au moins, il est clair que $\theta$ est surjectif. D'autre part, si $\theta(\phi_n(z)) = 0$, autrement dit $z/f^n = 0$, il existe un entier $k \geq 0$ tel que $f^k z = 0$, donc $\phi_{n+k, n}(z) = 0$, ce qui entraîne $\phi_n(z) = 0$. On peut donc identifier $M_{\mathfrak{p}}$ et $\varinjlim M_n$ au moyen de $\theta$.

Ecrivons maintenant $M_{f,n}$, $\phi_{nm}^f$ et $\phi_n^f$ au lieu de $M_n$, $\phi_{nm}$ et $\phi_n$. Soit $g$ un second élément de $A$. Comme $f^n$ divise $f^n g^n$, on a un homomorphisme fonctoriel $\rho_{fg, f} : M_{\mathfrak{p}} \to M_{fg}$ (n°4) ; si on identifie $M_{\mathfrak{p}}$ et $M_{fg}$ à $\varinjlim M_{f,n}$ et $\varinjlim M_{fg,n}$, $\rho_{fg, f}$ s'identifie à la limite inductive des applications $\rho_{fg, f}^n : M_{f,n} \to M_{fg,n}$ définies par $\rho_{fg, f}^n(z) = g^n z$.

En effet, cela résulte immédiatement de la commutativité du diagramme

$$\begin{tikzcd}
M_{f,n} \arrow[r, "\rho_{fg, f}^n"] \arrow[d, "\phi_n^f"] & M_{fg,n} \arrow[d, "\phi_n^{fg}"] \\
M_{\mathfrak{p}} \arrow[r, "\rho_{fg, f}"] & M_{fg}
\end{tikzcd}$$

\section*{§2. Faisceaux.}

Pour les notions et résultats fondamentaux de la théorie des faisceaux, nous renverrons le plus souvent aux ouvrages suivants: R. GODEMENT, Topologie algébrique et Théorie des faisceaux, I, \emph{Actual. Scient. et Ind.}, n°1252, Paris (Hermann), 1958 (cité G). J. P. SERRE, Faisceaux algébriques cohérents, \emph{Annals of Math.}, 61 (1955), 197-278 (cité FAC). A. GROTHENDIECK, Sur quelques points d'algèbre homologique, \emph{Tôhoku Math. J.}, 9 (1957), 119-221 (cité T).

Le plus souvent, nous utiliserons les notations et la terminologie de (G); en particulier, pour un faisceau $F$ sur un espace $X$, nous noterons $F(x)$ ou $F_x$ la fibre de $F$ au point $x$. Le \emph{support} d'un faisceau de groupes abéliens $F$ sur $X$ est l'ensemble des $x \in X$ tels que $F(x) \neq \{0\}$; cet ensemble n'est pas nécessairement fermé dans $X$.

Pour une partie $M$ de $X$, nous écrirons $F(M)$ ou $\Gamma(M, F)$ le groupe de sections de $F$ au-dessus de $M$; on écrira $\Gamma(F) = \Gamma(X, F)$. Si $u: F \to G$ est un homomorphisme

\newpage

%% ===== Page 9 =====

- ix -

phisme de faisceaux sur $X$, nous désignerons par $u_V$ l'homomorphisme correspondant $s \mapsto u \cdot s$ de $F(V)$ dans $G(V)$, par $u_x$ l'homomorphisme $F(x) \to G(x)$ limite inductive des $u_V$ lorsque $V$ parcourt les voisinages ouverts de $x$ dans $X$.

Enfin, si $A$ est un faisceau d'anneaux sur $X$, au lieu de parler de $A$-Modules (G, II, 2, 2), nous parlerons aussi, par abus de langage de \emph{faisceau de $A$-modules} ($x \in X$).

1. \emph{Images directes de faisceaux de groupes abéliens.}

Soient $X, Y$ deux espaces topologiques, $\psi: X \to Y$ une application continue. 
Si $F$ est un faisceau de groupes abéliens sur $X$, on sait que l'\emph{image directe} de $F$ par $\psi$ est le faisceau $\psi_*(F)$ de groupes abéliens sur $Y$, défini de la façon suivante: pour toute partie ouverte $U$ de $Y$, $\psi_*(F)(U) = F(\psi^{-1}(U))$ (G, II, 1, 13). 
Si $S$ est le support de $F$ et si $y \notin \psi(S)$, il résulte aussitôt de la définition que $\psi_*(F)(y) = \{0\}$; mais il faut noter que si $y$ est adhérent à $\psi(S)$ mais n'appartient pas à $\psi(S)$, on peut avoir $\psi_*(F)(y) \neq \{0\}$ même si $X = S$.

Soit $x \in X$; pour tout voisinage ouvert $U$ de $\psi(x)$ dans $Y$, et toute section $s \in \psi_*(F)(U) = F(\psi^{-1}(U))$, posons $\psi_{x,U}(s) = s(x)$; on définit ainsi un homomorphisme $\psi_{x,U}: \psi_*(F)(U) \to F(x)$, et il est clair que ces homomorphismes forment un système inductif; la limite inductive $\psi_x = \lim_{\to} \psi_{x,U}$ est donc un homomorphisme $\psi_*(F)(\psi(x)) \to F(x)$, qui en général n'est ni injectif ni surjectif.

Toutefois, si $\psi$ est un homéomorphisme de $X$ sur une partie $\psi(X)$ de $Y$, $\psi_x$ est un isomorphisme pour tout $x \in X$. Ceci s'applique en particulier à l'injection canonique $j$ d'une partie $X$ de $Y$ dans $Y$; le faisceau $j_*(F)$ induit alors $F$ sur $X$. Si de plus $X$ est une partie fermée de $Y$, $j_*(F)$ est le faisceau sur $Y$ qui induit $F$ sur $X$ et $0$ sur $Y-X$.

Considérons maintenant deux faisceaux $F_1, F_2$ de groupes abéliens sur $X$, et soit $u: F_1 \to F_2$ un homomorphisme. Pour tout ensemble ouvert $U \subseteq Y$, $u$ définit un homomorphisme $F_1(\psi^{-1}(U)) \to F_2(\psi^{-1}(U))$, et ces homomorphismes sont compatibles avec les opérations de restriction, donc définissent un homomorphisme

\newpage

%% ===== Page 10 =====

$\psi_*(u) : \psi_*(F_1) \to \psi_*(F_2)$. Pour toute section $s \in F_1(\psi^{-1}(U)) = \psi_*(F_1)(U)$, on a par définition $\psi_{x,U}(\psi_*(u) \cdot s) = u_x(\psi_{x,U}(s))$; la définition de la limite inductive montre donc que le diagramme

\begin{tikzcd}
F_1(x) \arrow[r, "\psi_x"] & \psi_*(F_1)(\psi(x)) \\
F_2(x) \arrow[u, "u_x"] \arrow[r, "\psi_x"] & \psi_*(F_2)(\psi(x)) \arrow[u, "\psi_*(u)(\psi(x))"']
\end{tikzcd}

est commutatif.

Si $v : F_2 \to F_3$ est un homomorphisme dans un troisième faisceau $F_3$ sur $X$, on a $\psi_*(v \circ u) = \psi_*(v) \circ \psi_*(u)$; autrement dit $\psi_*$ est un foncteur covariant de la catégorie des faisceaux de groupes abéliens sur $X$ dans celle des faisceaux de groupes abéliens sur $Y$. En outre, comme le foncteur $F \to F(X)$ est exact à gauche, et que l'exactitude se conserve par passage à la limite inductive, le foncteur $\psi_*$ est \emph{exact à gauche}.

Soient $Z$ un troisième espace topologique, $\psi' : Y \to Z$ une application continue, et $\psi'' = \psi' \circ \psi$. Il est clair que l'on a $\psi''_*(F) = \psi'_*(\psi_*(F))$ et $\psi''_x = \psi'_x \psi_x$ pour tout $x \in X$; en outre, pour tout homomorphisme $u$ de faisceaux sur $X$, on a $\psi''_*(u) = \psi'_*(\psi_*(u))$. En d'autres termes, $\psi'_*$ est le composé des foncteurs $\psi'_*$ et $\psi_*$, ce qu'on peut écrire

$$(\psi' \circ \psi)_* = \psi'_* \circ \psi_* .$$

Notons enfin que pour toute partie ouverte $U$ de $Y$, l'image par du faisceau induit $F|_{\psi^{-1}(U)}$ n'est autre que le faisceau induit $\psi_*(F)|U$.

2. \emph{Images réciproques de faisceaux de groupes abéliens.}

Les hypothèses sur $X, Y$ et $\psi$ étant les mêmes qu'au n°1, soit $G$ un faisceau de groupes abéliens sur $Y$. L'image réciproque $\psi^*(G)$ est le faisceau de groupes abéliens sur $X$ défini comme suit : une section $s$ de $\psi^*(G)$ au-dessus d'un ouvert $U$ de $X$ est une application de $U$ dans l'espace étalé $G$ qui est telle que, pour tout $x \in U$, il existe un voisinage $V$ de $\psi(x)$ dans $Y$, une section $s'$ de $G$ au-dessus de $V$ et un voisinage $W \subseteq \psi^{-1}(V)$ de $x$ de sorte que $s(z) =$

\newpage

%% ===== Page 11 =====

$s'(\psi(z))$ pour tout $z \in W$ ($G$,II,1.12). En particulier, si $X$ est une partie
quelconque de $Y$ et $j$ l'injection canonique $X \to Y$, $j^*(G)$ n'est autre que le
faisceau induit $G|X$.

Soient $G_1, G_2$ deux faisceaux de groupes abéliens sur $Y$, et $v: G_1 \to G_2$ un
homomorphisme. On définit un homomorphisme $\psi^*(v)$ de $\psi^*(G_1)$ dans $\psi^*(G_2)$ de
la façon suivante: si $s$ est une section de $\psi^*(G_1)$ au-dessus d'un ouvert $U$,
$\psi^*(v) \cdot s$ est définie comme la section qui au voisinage de chaque point $x \in U$,
coïncide avec $v \cdot s' \circ \psi$, si $s'$ est une section de $G_1$ au voisinage de $\psi(x)$ telle
que $s$ coïncide avec $s' \circ \psi$ au voisinage de $x$. Si $w: G_2 \to G_3$ est un homomorphisme dans un troisième faisceau $G_3$ sur $Y$, on a $\psi^*(w \cdot v) = \psi^*(w) \psi^*(v)$, autrement dit $\psi^*$ est un foncteur covariant de la catégorie des faisceaux de groupes abéliens sur $Y$ dans celle des faisceaux de groupes abéliens.

Soient $Z$ un troisième espace topologique, $\psi': Y \to Z$ une application continue, et $\psi'' = \psi' \circ \psi$. Si $H$ est un faisceau sur $Z$, on a $\psi''^*(H) = \psi^*(\psi'^*(H))$, et pour tout homomorphisme $w$ de faisceaux sur $Z$, on a $\psi''^*(w) = \psi^*(\psi'^*(w))$. En d'autres termes, $\psi''^*$ est le composé des foncteurs $\psi'^*$ et $\psi^*$, ce qu'on peut écrire
$$\psi''^* = \psi^* \psi'^*.$$

\section*{3. Relations entre images directes et images réciproques.}

Pour tout ouvert $V \subset Y$, il y a un homomorphisme injectif de $G(V)$ dans
$\psi^*(G)(\psi^{-1}(V))$, faisant correspondre à toute section $s \in G(V)$ la section
$z \to s'(\psi(z))$ de $\psi^*(G)$ au-dessus de $\psi^{-1}(V)$; ces homomorphismes étant compatibles avec les opérations de restriction, définissent (en vertu de la définition de $\psi_*$) un homomorphisme injectif $a: G \to \psi_*(\psi^*(G))$, évidemment fonctoriel; cet homomorphisme n'est pas surjectif en général. Toutefois, la définition des sections de $\psi^*(G)$ montre que l'homomorphisme $\psi^*_x(G)(\psi(x))$ dans $\psi^*(G)(x)$ est un isomorphisme, qui permet d'identifier ces deux fibres; en outre, pour tout homomorphisme $v: G_1 \to G_2$, le diagramme

\newpage

%% ===== Page 12 =====

- xii -

$$\begin{tikzcd}
\psi^*(G_1)(x) \arrow[d, "\psi^*(v)_x"] & G_1(x) \arrow[l, no head] \arrow[d, "v_x"] \\
\psi^*(G_2)(x) & G_2(\psi(x)) \arrow[l, no head]
\end{tikzcd}$$
est commutatif. Une conséquence immédiate de ce fait est que le foncteur $\psi^*$ est exact.

Soit maintenant $F$ un faisceau de groupes abéliens sur $X$; pour toute section $s$ de $\psi_*(\psi^*(F))$ au-dessus d'un ouvert $U \subset X$, on définit une section de $F$ au-dessus de $U$ en faisant correspondre à $x \in U$ l'élément $\psi_x(s(x))$ de $F(x)$; on vérifie immédiatement que c'est bien une application continue de $U$ dans l'espace étalé $F$. On a ainsi défini un homomorphisme fonctoriel $\beta : \psi^*(\psi_*(F)) \to F$, qui n'est en général ni injectif ni surjectif.

4. \emph{Morphismes de faisceaux compatibles avec une application continue.}

Les hypothèses et notations étant comme ci-dessus, les définitions de $\psi_*$ et $\psi^*$ permettent de définir deux bifoncteurs
$$(F, G) \to \operatorname{Hom}_Y(G, \psi_*(F))$$
$$(F, G) \to \operatorname{Hom}_X(\psi^*(G), F)$$
à valeurs dans la catégorie des groupes abéliens, covariants en $F$ et contravariants en $G$. Nous nous proposons d'établir un isomorphisme canonique fonctoriel
$$\operatorname{Hom}_Y(G, \psi_*(F)) \iff \operatorname{Hom}_X(\psi^*(G), F)$$
entre ces deux bifoncteurs. Pour cela, pour tout $\theta \in \operatorname{Hom}_Y(G, \psi_*(F))$ nous désignerons par $\theta^b$ l'homomorphisme composé
$$\theta^b : \psi^*(G) \xrightarrow{\psi^*(\theta)} \psi^*(\psi_*(F)) \xrightarrow{\beta} F$$
et pour tout $\omega \in \operatorname{Hom}_X(\psi^*(G), F)$, nous désignerons par $\omega^\#$ l'homomorphisme composé
$$\omega^\# : G \xrightarrow{\alpha} \psi_*(\psi^*(G)) \xrightarrow{\psi_*(\omega)} \psi_*(F).$$

Comme $\alpha, \beta$ sont fonctoriels, et que $\psi_*$ et $\psi^*$ sont des foncteurs additifs, il est immédiat de vérifier que $\theta \to \theta^b$ et $\omega \to \omega^\#$ sont des morphismes de bifoncteurs. Reste à montrer que $(\theta^b)^\# = \theta$ et $(\omega^\#)^b = \omega$. Or, si $s \in \Gamma(U, \psi^*(G))$,

\newpage

%% ===== Page 13 =====

- xiii -

$e^b \cdot s$ coïncide, au voisinage de chaque $x \in U$, avec $(\theta \cdot s') \circ \psi$, si $s$ coïncide avec $s' \circ \psi$ dans un voisinage de $x$; et de même, si $s' \in \Gamma(V, G)$, $\omega^\sharp \cdot s'$ est la section de $\psi_*(F)$ au-dessus de $V$ identique à la section $\omega \cdot s'$ de $F$ au-dessus de $\psi^{-1}(V)$; ces deux relations entraînent nos assertions.

Il résulte de ce qui précède que, étant donné le faisceau $G$ sur $Y$, le faisceau $\psi^*(G)$ est solution du \emph{problème universel} suivant: trouver un faisceau $F_0$ sur $X$ tel que tout homomorphisme $G \to \psi_*(F)$ se factorise en deux homomorphismes $G \to \psi_*(F_0)$ et $\psi_*(F_0) \to \psi_*(F)$, le second étant de la forme $\psi_*(v)$, où $v$ est un homomorphisme de $F_0$ dans $F$.

Par définition, nous dirons qu'un homomorphisme $\theta: G \to \psi_*(F)$ est un \emph{morphisme du faisceau $G$ dans le faisceau $F$, compatible avec $\psi$}, ou encore un \emph{$\psi$-morphisme de $G$ dans $F$}. On dit par convention que $\theta$ (en tant que $\psi$-morphisme) est \emph{injectif} (resp. \emph{surjectif}) si l'homomorphisme correspondant $\theta^b: \psi^*(G) \to F$ est injectif (resp. surjectif); cela n'entraîne pas que $\theta$, considéré comme homomorphisme de $G$ dans $\psi_*(F)$, soit injectif (resp. surjectif); toutefois, si $\theta^b$ est injectif, il en est de même de $\theta_{(x)}$ pour tout $x \in X$.

Nous allons maintenant voir que les couples $(X, F)$ formés d'un espace topologique $X$ et d'un faisceau de groupes abéliens $F$ sur $X$, forment une \emph{catégorie}, lorsqu'on définit un \emph{morphisme de $(X, F)$ dans $(Y, G)$} comme un couple $(\psi, \theta)$, où $\psi$ est une application continue de $X$ dans $Y$ et $\theta$ un $\psi$-morphisme de $G$ dans $F$. Pour cela, nous devons définir le \emph{composé} de deux morphismes $(\psi, \theta): (X, F) \to (Y, G)$ et $(\psi', \theta') : (Y, G) \to (Z, H)$; par définition, ce sera le morphisme $(\psi'', \theta'')$, où $\psi'' = \psi' \circ \psi$, et où $\theta''$ est l'homomorphisme composé
$$\theta'': H \xrightarrow{\theta'} \psi'_*(G) \xrightarrow{\psi'_*(\theta)} \psi'_*(\psi_*(F));$$
la vérification de la condition d'associativité est immédiate. En outre, on voit aussitôt que $\theta''^b$ n'est autre que l'homomorphisme
$$\theta''^b: \psi^*(\psi'_*(H)) \xrightarrow{\psi^*(\theta'^b)} \psi^*(G) \xrightarrow{\theta^b} F.$$
Comme $\psi^*$ est un foncteur exact, on en déduit aussitôt que si les deux mor-

\newpage

%% ===== Page 14 =====

-- xiv --

phismes $\psi, \psi'$ sont \emph{injectifs} (resp. \emph{surjectifs}), $\psi$ est aussi injectif (resp. surjectif). On vérifie alors que si $\psi$ est injectif (resp. surjectif) et $\theta$ surjectif (resp. injectif), $(\psi,\theta)$ est un \emph{monomorphisme} (resp. \emph{épimorphisme}) ($T, I, 1.1$) pour la catégorie ainsi définie.

En particulier, si $M$ est une partie de $X$, $j$ l'injection canonique $M \to X$, le morphisme $(j,\rho): (M, F|M) \to (X,F)$, où $\rho$ est l'application identique de $F|M$, est un monomorphisme; le composé d'un morphisme $(\psi,\theta)$ et de $(j,\rho)$ est appelé la \emph{restriction} de $(\psi,\theta)$ à $M$. De même, si $N$ est une partie de $Y$ contenant $\psi(X)$, et $h$ l'injection canonique $N \to Y$, on peut écrire $\psi = h\psi_0$, où $\psi_0$ est l'application $\psi$ considérée comme prenant ses valeurs dans $N$; on a évidemment $\psi^*(G|N) = \psi^*(G)$, et si $\theta_0$ est l'homomorphisme $G|N \to (\psi_0)_*(F)$ tel que $\theta_0^b$ soit égal à $\theta^b$, on peut considérer $(\psi,\theta)$ comme composé de $(h,\rho)$ et de $(\psi_0,\theta_0)$, $\rho$ étant l'application identique de $G|N$.

\section*{5. Espaces annelés.}

Toutes les définitions et tous les résultats précédents s'appliquent lorsqu'on remplace les faisceaux de groupes abéliens par des faisceaux dont les fibres sont dans une catégorie quelconque admettant des limites inductives ($T, I, 1.8$) (en exceptant naturellement les assertions faisant intervenir la structure de groupe abélien, comme par exemple le fait que $\operatorname{Hom}_X(G, \psi_*(F))$ est un groupe abélien; en général ce sera simplement un ensemble).

Nous utiliserons particulièrement les faisceaux \emph{d'anneaux}; nous appellerons \emph{espace annelé} un couple $(X,A)$ formé d'un espace topologique $X$ et d'un faisceau \emph{d'anneaux} $A$ sur $X$; on dit encore que $A$ est le \emph{faisceau structural} de l'espace annelé, et on le note $\mathcal{O}_X$ (et les fibres $\mathcal{O}_x$ ou $\mathcal{O}(x)$), lorsqu'aucune confusion n'en résulte. Pour toute partie $M$ de $X$, le couple $(M, A|M)$ est évidemment un espace annelé, dit induit sur $M$ par $(X,A)$ (ou encore la \emph{restriction} de $(X,A)$ à $M$). Nous considérerons toujours les espaces annelés comme formant une catégorie suivant la définition des morphismes donnée au n°4. Soient $(X,A)$, $(Y,B)$ deux espaces annelés, $\psi = (\psi,\theta)$ un morphisme de $(X,A)$ dans $(Y,B)$. Soit

\newpage

%% ===== Page 15 =====

$F$ un $A$-Module, $\psi_*(F)$ son image directe par $\psi$, qui est un faisceau de groupes abéliens sur $Y$; par définition, pour tout ouvert $U \subset Y$, $\psi_*(F)(U) = F(\psi^{-1}(U))$ est muni d'une structure de module par rapport à l'anneau $\psi_*(A)(U) = A(\psi^{-1}(U))$; ces structures étant compatibles par rapport aux opérations de restriction, définissent sur $\psi_*(F)$ une structure de $\psi_*(A)$-Module. Comme $\theta$ est un homomorphisme de $B$ dans $\psi_*(A)$, il définit sur $\psi_*(F)$ une structure de $B$-Module; nous dirons que ce $B$-Module est l'image directe de $F$ par le morphisme $\tilde{\psi}$, et nous le noterons $\tilde{\psi}_*(F)$. Si $F_1, F_2$ sont deux $A$-Modules sur $X$, $u$ un $A$-homomorphisme $F_1 \to F_2$, $\psi_*(u)$ est un $\psi_*(A)$-homomorphisme de $\psi_*(F_1)$ dans $\psi_*(F_2)$, et a fortiori un $B$-homomorphisme; en tant que $B$-homomorphisme, nous le noterons $\tilde{\psi}_*(u)$. On voit donc que $\tilde{\psi}_*$ est un \emph{foncteur covariant} (exact à gauche) de la catégorie des $A$-Modules dans celle des $B$-Modules.

Soient maintenant $G$ un $B$-Module, $\psi^*(G)$ son image réciproque, qui est un faisceau de groupes abéliens sur $X$. Il résulte immédiatement de la définition des sections d'une image réciproque de faisceaux que $\psi^*(G)$ est naturellement muni d'une structure de $\psi^*(B)$-Module. D'autre part, l'homomorphisme $\theta^b$ de $\psi^*(B)$ dans $A$ muni $A$ d'une structure de $\psi^*(B)$-Module, que nous noterons $A[\theta]$; et le produit tensoriel $\psi^*(G) \otimes_{\psi^*(B)} A[\theta]$ est muni naturellement d'une structure de $A$-Module. Nous dirons que ce $A$-Module est l'image réciproque de $G$ par le morphisme $\tilde{\psi}$, et nous le noterons $\tilde{\psi}^*(G)$. Si $G_1, G_2$ sont deux $B$-Modules sur $Y$, $v$ un $B$-homomorphisme $G_1 \to G_2$, $\psi^*(v)$ est, comme on le vérifie aussitôt, un $\psi^*(B)$-homomorphisme de $\psi^*(G_1)$ dans $\psi^*(G_2)$; par suite $\psi^*(v) \otimes 1$ est un $A$-homomorphisme de $\tilde{\psi}^*(G_1)$ dans $\tilde{\psi}^*(G_2)$; si on le note $\tilde{\psi}^*(v)$, on voit qu'on a défini $\tilde{\psi}^*$ comme un \emph{foncteur covariant} de la catégorie des $B$-Modules dans celle des $A$-Modules. Ici, ce foncteur (contrairement à $\psi^*$) n'est plus exact mais seulement \emph{exact à droite}.

Nous avons donc de nouveau défini deux bifoncteurs

$$(F, G) \to \operatorname{Hom}_B(G, \tilde{\psi}_*(F)) \quad\quad (F, G) \to \operatorname{Hom}_A(\tilde{\psi}^*(G), F)$$

\newpage

%% ===== Page 16 =====

à valeurs dans la catégorie des groupes abéliens, covariants en $F$ et contravariants en $G$, $F$ étant dans la catégorie des $A$-Modules et $G$ dans celle des $B$-Modules. On va voir encore qu'il y a un isomorphisme canonique fonctoriel
$$(1) \qquad \operatorname{Hom}_B(G, \psi_*(F)) \cong \operatorname{Hom}_A(\psi^*(G), F)$$
entre ces deux bifoncteurs. Pour cela nous définirons d'abord un homomorphisme fonctoriel $\lambda: G \to \psi_*(\psi^*(G))$ de $B$-Modules et un homomorphisme fonctoriel $\mu: \psi^*(\psi_*(F)) \to F$ de $A$-Modules. On a un homomorphisme naturel $j: G \to \psi_*(\psi^*(G))$, qui, sur chaque fibre, se réduit à $z \to z \otimes 1$, et qui est évidemment un $\psi^*(B)$-homomorphisme; on en déduit un $\psi_*(B)$-homomorphisme $\psi_*(j): \psi_*(G) \to \psi_*(\psi^*(G))$, qu'on peut aussi considérer comme un $B$-homomorphisme puisque $B$ s'injecte dans $\psi_*(\psi^*(B))$ par $a$. On prend alors pour $\lambda$ l'homomorphisme composé
$$\lambda: G \to \psi_*(\psi^*(G)) \xrightarrow{\psi_*(j)} \psi_*(\psi^*(G)).$$
D'autre part, $F$ peut être considéré comme un $\psi^*(B)$-Module au moyen de l'homomorphisme $e^b$ de $\psi^*(B)$ dans $A$; on vérifie alors aussitôt que l'homomorphisme $\mu: \psi^*(\psi_*(F)) \to F$ est un homomorphisme de $\psi^*(B)$-Modules. On a d'autre part un homomorphisme naturel $h: F \otimes_{\psi^*(B)} \mathcal{O}_e \to F$ de $A$-Modules, qui sur chaque fibre se réduit à $z \otimes a \to az$. On prend pour $\mu$ l'homomorphisme composé
$$\mu: \psi^*(\psi_*(F)) \xrightarrow{\theta \otimes 1} F \otimes_{\psi^*(B)} \mathcal{O}_e \xrightarrow{h} F.$$
Cela étant, pour $v \in \operatorname{Hom}_B(G, \psi_*(F))$ et $u \in \operatorname{Hom}_A(\psi^*(G), F)$, on pose
$$v^b: \psi^*(G) \xrightarrow{\psi^*(v)} \psi^*(\psi_*(F)) \xrightarrow{\mu} F$$
et
$$u^\sharp: G \xrightarrow{\lambda} \psi_*(\psi^*(G)) \xrightarrow{\psi_*(u)} \psi_*(F).$$
On vérifie comme au $n^\circ 4$ que $v \to v^b$ et $u \to u^\sharp$ sont fonctoriels et que l'on a $(v^b)^\sharp = v$ et $(u^\sharp)^b = u$.

Considérons les anneaux $A = \Gamma(A)$, $B = \Gamma(B)$ des sections des faisceaux d'anneaux $A, B$; le groupe abélien $\operatorname{Hom}_B(G, \psi_*(F))$ est naturellement muni d'une

\newpage

%% ===== Page 17 =====

- xvii -

structure de bimodule sur les anneaux $B = \Gamma(B)$ et $\Gamma(\psi_*(A))$ et comme ce dernier n'est autre que $\Gamma(A) = A$, $\operatorname{Hom}_B(G, \psi_*(F))$ est un $(A,B)$-bimodule. De même, $\operatorname{Hom}_A(\psi^*(G), F)$ est naturellement muni d'une structure de bimodule sur $A = \Gamma(A)$ et $\Gamma(\psi^*(B))$, et comme $\Gamma(B)$ s'identifie à un sous-anneau de $\Gamma(\psi^*(B))$ (n°2) $\operatorname{Hom}_A(\psi^*(G), F)$ est un $(A,B)$-bimodule. Cela étant, il est facile de vérifier que l'isomorphisme (1) est aussi un isomorphisme de $(A,B)$-bimodules.

6. \emph{Faisceaux cohérents et faisceaux quasi-cohérents.}

Rappelons que si $(X, \mathcal{O}_X)$ est un espace annelé, on dit qu'un $\mathcal{O}_X$-Module $F$ est \underline{de type fini} si, pour tout $x \in X$, il y a un voisinage ouvert $V$ de $x$, un entier $n > 0$ et un homomorphisme surjectif $\mathcal{O}_X^n|V \to F|V \to 0$; on dit que $F$ est \underline{cohérent} si $F$ est de type fini et si, pour toute partie ouverte $V$ de $X$ et tout homomorphisme $\mathcal{O}_X^n|V \to F|V$, le noyau de cet homomorphisme est un faisceau (sur $V$) de type fini. Pour les propriétés des faisceaux cohérents, nous renvoyons à (FAC), chap. I, §2. On sait en particulier que si $F$ est cohérent, pour tout $x \in X$ il y a un voisinage ouvert $V$ de $x$ tel que $F|V$ soit isomorphe au \underline{conoyau} d'un homomorphisme $\mathcal{O}_X^n|V \to \mathcal{O}_X^m|V$. Plus généralement, nous dirons qu'un $\mathcal{O}_X$-Module $F$ est \underline{quasi-cohérent} si, pour tout $x \in X$, il y a un voisinage ouvert $V$ de $x$ tel que $F|V$ soit isomorphe au \underline{conoyau} d'un homomorphisme de la forme $\mathcal{O}_X^{(I)}|V \to \mathcal{O}_X^{(J)}|V$, où $I$ et $J$ sont des ensembles d'indices arbitraires.

\newpage

%% ===== Page 18 =====

\marginpar{Archives Grothendieck-kpt59} \hfill 326. Bis

\section*{CHAPITRE I}
\section*{PRESCHEMAS.}

\subsection*{§1. Schémas affines.}

\noindent 1. Le spectre premier d'un anneau.

\noindent \underline{Notations}: Soit $A$ un anneau commutatif $M$ un $A$-module. Dans ce chapitre et les suivants, nous utiliserons constamment les notations suivantes:

$\operatorname{Spec}(A) = \text{ensemble des idéaux premiers de } A, \text{ appelé aussi \underline{spectre premier}}$
de $A$; pour un $x \in X = \operatorname{Spec} A$, il sera parfois commode d'écrire
$j_x$ au lieu de $x$.

$A_x = A_{j_x} = \text{\underline{anneau (local) des fractions }} S^{-1}A, \text{ où } S = A - j_x$.

$\mathfrak{m}_x = j_x A_x = \text{idéal maximal de } A_x$.

$k_x = A_x / \mathfrak{m}_x = \text{corps des restes de } A_x, \text{ isomorphe canoniquement au corps des}$
\phantom{$k_x =$} fractions de $A/j_x$, auquel on l'identifie.

$f(x) = \text{classe de } f \pmod{j_x} \text{ dans } A/j_x \subset k(x), \text{ pour } f \in A \text{ et } x \in X$.

\noindent \underline{Les relations}
$$f(x) = 0 \quad \text{et} \quad f \in j_x$$
sont donc équivalentes.

$M_x = M \otimes_A A_x = \text{module des fractions à dénominateurs dans } A - j_x$.

$r(E) = \text{racine de l'idéal de } A \text{ engendré par } E, \text{ pour tout } E \subset A$.

$V(E) = \text{ensemble des } x \in X \text{ tels que } E \subset j_x \text{ (ou encore ensemble des } x \in X$
\phantom{$V(E) =$} tels que $f(x) = 0 \text{ pour tout } f \in E \text{). On a donc}$
$$(1) \quad r(E) = \bigcap_{x \in V(E)} j_x.$$

$V(f) = V(\{f\}) \text{ pour } f \in A$.

$D(f) = X - V(f) = \text{ensemble des } x \in X \text{ tels que } f(x) \neq 0$.

\noindent \underline{Proposition 1}. On a les propriétés suivantes:

(i) $V(0) = X, \quad V(1) = \emptyset$.

(ii) La relation $E \subset E'$ entraîne $V(E) \supset V(E')$.

(iii) Pour toute famille $(E_\lambda)$ de parties de $A$, $V(\bigcup_\lambda E_\lambda) = \bigcap_\lambda V(E_\lambda)$.

(iv) $V(EE') = V(E) \cup V(E')$.

(v) $V(E) = V(r(E))$.

\newpage

%% ===== Page 19 =====

- 2 -

Les propriétés (i), (ii), (iii) sont triviales, et (v) résulte de (ii) et de la formule (1). Il est évident que $V(EE') \subset V(E) \cup V(E')$; inversement, si $x \notin V(E)$ et $x \notin V(E')$, il existe $f, f'$ dans $A$ tels que $f(x) \neq 0$ et $f'(x) \neq 0$ dans $k(x)$, d'où $f(x)f'(x) \neq 0$, $x \notin V(EE')$, ce qui prouve (iv).

La prop. 1 montre entre autres que les ensembles de la forme $V(E)$ sont les ensembles fermés d'une topologie sur $X$, que nous appellerons la \emph{topologie spectrale}; sauf mention expresse du contraire, on supposera toujours $\operatorname{Spec}(A)$ muni de la topologie spectrale.

Pour toute partie $Y$ de $X$, nous désignerons par $j(Y)$ l'ensemble de $f \in A$ tels que $f(y) = 0$ pour tout $y \in Y$; il revient au même de dire que $j(Y)$ est l'intersection des idéaux premiers $j_y$ pour $y \in Y$. Il est clair que la relation $Y \subset Y'$ entraîne $j(Y) \supset j(Y')$ et que l'on a
$$(2) \quad j(\bigcup_\lambda Y_\lambda) = \bigcap_\lambda j(Y_\lambda)$$
pour toute famille $(Y_\lambda)$ de parties de $X$.

\emph{Proposition 2.} a) Pour toute partie $E$ de $A$, on a $j(V(E)) = r(E)$.
b) Pour toute partie $Y$ de $X$, $V(j(Y)) = \bar{Y}$, adhérence de $Y$ dans $X$.

a) est conséquence immédiate des définitions et de (1); d'autre part, $V(j(Y))$ est fermé et contient $Y$; inversement, si $Y \subset V(E)$ on a $f(y) = 0$ pour tout $f \in E$ et tout $y \in Y$, donc $E \subset j(Y)$, $V(E) \supset V(j(Y))$ ce qui prouve b).

\emph{Corollaire 1.} Les parties fermées de $X$ et les idéaux de $A$ égaux à leurs racines (autrement dit les intersections d'idéaux premiers) se correspondent biunivoquement par les applications $Y \to j(Y)$, $a \to V(a)$; à la réunion $Y_1 \cup Y_2$ de deux parties fermées correspond $j(Y_1) \cap j(Y_2)$, et à l'intersection d'une famille quelconque $(Y_\lambda)$ de parties fermées correspond la racine de la somme des $j(Y_\lambda)$.

\emph{Corollaire 2.} Si $A$ est un anneau noethérien, $X = \operatorname{Spec}(A)$ est un espace noethérien.

On notera que la réciproque de ce corollaire est inexacte, comme le montre

\newpage

%% ===== Page 20 =====

- 3 -

l'exemple d'un anneau non noethérien ayant un seul idéal premier $\neq \{0\}$, par exemple un anneau de valuation non discrète de rang 1.

\emph{Corollaire 3.} Pour tout $x \in X$, l'adhérence de $\{x\}$ est l'ensemble des $y \in X$ tels que $j_y \supset j_x$. Pour que $\{x\}$ soit fermé, il faut et il suffit que $j_x$ soit maximal.

\emph{Corollaire 4.} Si $x$ et $y$ sont deux points distincts de $X$, il existe un voisinage ouvert de l'un des points $x, y$ qui ne contient pas l'autre (axiome $T_0$).

En effet, on a soit $j_x \not\supset j_y$, soit $j_y \not\supset j_x$, donc il existe un ensemble $V(E)$ contenant l'un des points $x, y$ et non l'autre.

D'après la prop. 1, (iv), pour deux éléments $f, g$ de $A$, on a
$$(3) \quad D(fg) = D(f) \cap D(g).$$

Notons aussi que la relation $D(f) = D(g)$ signifie, d'après la prop. 2 a) et la prop. 1, (v), que $r(f) = r(g)$, ou encore que les idéaux premiers minimaux contenant $(f)$ et $(g)$ sont les mêmes (en particulier, il en est ainsi lorsque $f = ug$, où $u$ est inversible).

\emph{Proposition 3.} a) Lorsque $f$ parcourt $A$, les ensembles $D(f)$ forment une base de la topologie de $X$.

b) Pour tout $f \in A$, $D(f)$ est quasi-compact.

a) Soit $V$ un ensemble ouvert dans $X$; par définition, on a $U = X - V(E)$ où $E$ est une partie de $A$, et $V(E) = \bigcap_{f \in E} V(f)$, d'où $U = \bigcup_{f \in E} D(f)$.

b) D'après a), il suffit de prouver que si $(f_\lambda)_{\lambda \in L}$ est une famille d'éléments de $A$ telle que $D(f) \subset \bigcup_{\lambda \in L} D(f_\lambda)$, il existe une partie finie $J$ de $L$ telle que $D(f) \subset \bigcup_{\lambda \in J} D(f_\lambda)$. Soit $a$ l'idéal engendré par les $f_\lambda$; on a par hypothèse $V(f) \supset V(a)$, donc $r(f) \subset r(a)$; comme $f \in r(f)$ il existe un entier $n > 0$ tel que $f^n \in a$. Mais alors $f^n$ appartient à l'idéal engendré par une sous-famille finie $(f_\lambda)_{\lambda \in J}$, et on a $V(f^n) = V(f) \supset V(a) = \bigcap_{\lambda \in J} V(f_\lambda)$, c'est-à-dire $D(f) \subset \bigcup_{\lambda \in J} D(f_\lambda)$.

\emph{Proposition 4.} Soit $N = r(0)$ l'idéal des éléments nilpotents de $A$. Pour

\newpage

%% ===== Page 21 =====

que l'espace $X = \operatorname{Spec}(A)$ soit irréductible, il faut et il suffit que $A/N$ soit
un anneau intègre (autrement dit que $N$ soit premier).

Notons d'abord que les idéaux premiers de $A/N$ étant en correspondance bi-
univoque avec ceux de $A$, les spectres premiers de $A$ et de $A/N$ s'identifient
canoniquement avec leurs topologies. On peut donc supposer $N = (0)$. Si $X$
est réductible, il existe deux parties fermées $X_1, X_2$ distinctes de $X$ et telles
que $X = X_1 \cup X_2$, d'où $j(X) = j(X_1) \cap j(X_2) = (0)$, les idéaux $j(X_1)$ et $j(X_2)$
étant distincts de $(0)$; donc $A$ n'est pas intègre. Inversement, si dans $A$ on
a $fg = 0$, $f \neq 0$, $g \neq 0$, on a $V(f) \neq X$, $V(g) \neq X$ puisque l'intersection des
idéaux premiers est $(0)$, et $X = V(f) \cup V(g)$.

Si $a$ est un idéal de $A$, il y a correspondance biunivoque entre les idéaux
premiers de $A/a$ et les idéaux premiers de $A$ contenant $a$; par suite $\operatorname{Spec}(A/a)$
s'identifie canoniquement au sous-espace fermé $V(a)$ de $\operatorname{Spec}(A)$. La prop.4
entraîne alors le

Corollaire. a) Dans la correspondance biunivoque entre parties fermées de
$X = \operatorname{Spec}(A)$ et idéaux de $A$ égaux à leurs racines, les parties fermées ir-
réductibles correspondent aux idéaux premiers de $A$.

b) L'application $x \to \overline{\{x\}}$ établit une correspondance biunivoque entre $X$
et l'ensemble des parties fermées irréductibles de $X$.

2. Propriétés fonctorielles des spectres premiers d'anneaux.

Soient $A, A'$ deux anneaux, $\varphi$ un homomorphisme de $A'$ dans $A$. Pour tout
idéal premier $x = j_x \in \operatorname{Spec}(A)$, l'anneau $A'/\varphi^{-1}(j_x)$ est canoniquement iso-
morphe à un sous-anneau de $A/j_x$, donc est intègre, et par suite $\varphi^{-1}(j_x)$ est
un idéal premier de $A'$; nous le noterons $^a\varphi(x)$, et nous avons ainsi défini
une application $^a\varphi$ de $X = \operatorname{Spec}(A)$ dans $X' = \operatorname{Spec}(A')$, que nous appellerons
l'application associée à l'homomorphisme $\varphi$. Nous désignerons par $\varphi^\#$ l'homo-
morphisme injectif de $A'/j_{\varphi(x)}$ dans $A/j_x$, déduit de $\varphi$ par passage aux
quotients, ainsi que son prolongement canonique en un monomorphisme $k(^a\varphi(x))$

\newpage

%% ===== Page 22 =====

$$\varphi^\ast(f'(^\alpha(x))) = (\varphi(f'))(x) \quad (x \in X).$$

Proposition 5. a) Pour toute partie $E'$ de $A'$, on a
$$(5) \qquad ^a\varphi^{-1}(V(E')) = V(\varphi(E')).$$
b) Pour tout idéal $a$ de $A$, on a
$$(6) \qquad \overline{^a\varphi(V(a))} = V(\varphi^{-1}(a)).$$

En effet, la relation $^a\varphi(x) \in V(E')$ est par définition équivalente à $E' \subset j_{^a\varphi(x)} = \varphi^{-1}(j_x)$, donc à $\varphi(E') \subset j_x$, et finalement à $x \in V(\varphi(E'))$ d'où a).

Pour démontrer b), on peut supposer $a$ égal à sa racine, puisque $V(r(a)) = V(a)$ et $^a\varphi^{-1}(r(a)) = r(^a\varphi^{-1}(a))$; si on pose $Y = V(a)$, $a' = j(^a\varphi(Y))$, on a $\overline{^a\varphi(Y)} = V(a')$ ($n^\circ 1, \text{prop.} 2 \text{ b}$); la relation $f' \in a'$ est par définition équivalente à $f'(x') = 0$ pour tout $x' \in ^a\varphi(Y)$, donc, en vertu de la formule (4), elle est aussi équivalente à $\varphi(f')(x) = 0$ pour tout $x \in Y$, ou encore à $\varphi(f') \in j(Y) = a$, puisque $a$ est égal à sa racine; d'où b).

Corollaire 1. L'application $^a\varphi$ est continue.

Remarquons que si $A''$ est un troisième anneau, $\varphi'$ un homomorphisme de $A''$ dans $A'$, on a $^a(\varphi' \circ \varphi) = ^a\varphi \circ ^a\varphi'$; ce résultat et le cor. 1 signifient que $A \to \operatorname{Spec}(A)$ est un foncteur contravariant de la catégorie des anneaux dans celle des espaces topologiques.

Corollaire 2. Supposons que $\varphi$ soit tel que tout $f \in A$ s'écrive $f = h\varphi(f')$, où $h$ est inversible dans $A$ (ce qui est en particulier le cas lorsque $\varphi$ est surjectif). Alors $^a\varphi$ est un homéomorphisme de $X$ sur $^a\varphi(X)$.

Montrons que pour toute partie $E \subset A$, il existe une partie $E'$ de $A'$ telle que $V(E) = V(\varphi(E'))$; en vertu de l'axiome $T_0$ (cor. 4 de la prop. 2) et de la formule (5), cela entraînera d'abord que $^a\varphi$ est injective, puis, toujours en vertu de la formule (5), que $^a\varphi$ est un homéomorphisme. Or, il suffit pour chaque $f \in E$ de choisir un $f' \in A'$ tel que $h\varphi(f') = f$ avec $h$ inversible dans

\newpage

%% ===== Page 23 =====

- 6 -

En particulier, lorsque $\varphi$ est l'homomorphisme canonique de $A$ sur un anneau quotient $A/\mathfrak{a}$, $^a\varphi$ est l'injection canonique de $V(\mathfrak{a})$ (identifié à $\operatorname{Spec}(A/\mathfrak{a})$) dans $X = \operatorname{Spec}(A)$.

Un autre cas particulier du cor. 2 donne:

\emph{Corollaire 3.} Si $S$ est une partie multiplicative de $A$, le spectre $\operatorname{Spec}(S^{-1}A)$ s'identifie (avec sa topologie) au sous-espace de $X = \operatorname{Spec}(A)$ formé des $\mathfrak{x}$ tels que $\mathfrak{j}_{\mathfrak{x}} \cap S = \emptyset$.

On sait en effet que les idéaux premiers de $S^{-1}A$ sont les idéaux $S^{-1}\mathfrak{j}_{\mathfrak{x}}$, où $\mathfrak{j}_{\mathfrak{x}} \cap S = \emptyset$ et que l'on a $\mathfrak{j}_{\mathfrak{x}} = i_S^{-1}(S^{-1}\mathfrak{j}_{\mathfrak{x}})$, en désignant par $i_S$ l'application canonique de $A$ dans $S^{-1}A$. Il suffit donc d'appliquer à $i_S$ le cor. 2.

\emph{Corollaire 4.} Pour que $^a\varphi(X')$ soit partout dense dans $X$, il faut et il suffit que tout élément du noyau $\operatorname{Ker} \varphi$ soit nilpotent.

En effet, appliquant la formule (6) à $\mathfrak{a} = (0)$, il vient $^a\varphi(X) = V(\operatorname{Ker} \varphi)$, et pour que $V(\operatorname{Ker} \varphi) = X$, il faut et il suffit que $\operatorname{Ker} \varphi$ soit contenu dans tous les idéaux premiers de $A$, donc dans l'idéal des éléments nilpotents de $A$.

\medskip

3. \underline{Faisceau associé à un module.}

Notations: Soient $A$ un anneau commutatif, $M$ un $A$-module $f$ un élément de $A$.
On utilisera les notations suivantes:
$S_f = \text{ensemble multiplicatif des } f^n \text{ pour } n \ge 0$.
$A_f = \text{anneau des fractions } S_f^{-1}A$.
$M_f = \text{module des fractions } S_f^{-1}M$.

Soit $S_f'$ la partie multiplicative (saturée) de $A$ formée des $g \in A$ qui divisent un des éléments $f^n$ de $S_f$. Le lemme suivant résulte des définitions et du cor. 1 de la prop. 2 du n° 1:

\emph{Lemme 1.} Les conditions suivantes sont équivalentes:
(i) $g \in S_f'$; (ii) $S_g \subset S_f'$; (iii) $f \in r(g)$; (iv) $r(f) \subset r(g)$; (v) $V(g) \subset V(f)$; (vi) $D(f) \subset D(g)$.

\newpage

%% ===== Page 24 =====

Le $A_f$-module $M_f$ s'identifie canoniquement à $S_f^{-1}M$ (chap. 0, §1, n°4) si $D(f) = D(g)$, on a donc (lemme 1) $M_f = M_g$. Plus généralement, si $D(f) \supset D(g)$, donc $S_f' \subset S_g'$, on sait (chap. 0, §1, n°4) qu'il existe un homomorphisme fonctoriel $\rho_{g,f} : M_f \to M_g$, et si $D(f) \supset D(g) \supset D(h)$, on a
$$(7) \qquad \rho_{h,g} \circ \rho_{g,f} = \rho_{h,f}$$

Lorsque $f$ parcourt $A - \mathfrak{j}_x$ (pour un $x$ donné dans $X = \operatorname{Spec}(A)$), les ensembles $S_f'$ forment un ensemble filtrant croissant, car pour deux éléments $f,g \in A - \mathfrak{j}_x$, $S_f'$ et $S_g'$ sont contenus dans $S_{fg}'$; comme la réunion des $S_f'$ pour $f \in A - \mathfrak{j}_x$ est $A - \mathfrak{j}_x$, on en conclut (chap. 0, §1, n°4) que le $A$-module $M_x$ s'identifie canoniquement à la limite inductive $\varinjlim M_f$, relativement à la famille d'homomorphismes $(\rho_{g,f})$. Nous désignerons par $\rho_f^x$ l'homomorphisme canonique $M_f \to M_x$ pour $f \in A - \mathfrak{j}_x$ (ou, ce qui revient au même, $x \in D(f)$).

Pour toute partie ouverte $V$ de $X = \operatorname{Spec}(A)$, nous désignerons par $S(V,M)$ l'ensemble des fonctions $F$ définies dans $V$ et telles que $F(x) \in M_x$ pour tout $x \in V$. Si $V,W$ sont deux parties ouvertes de $X$ telles que $W \subset V$, nous désignerons par $\rho_{W,V}^M$ l'application de restriction $S(V,M) \to S(W,M)$.

Pour tout $f \in A$, on définit un homomorphisme de groupes abéliens $\sigma_f : M_f \to S(D(f),M)$, que nous désignerons aussi par $\xi \to \tilde{\xi}$, en posant, pour tout $\xi \in M_f$ et tout $x \in D(f)$
$$(8) \qquad \tilde{\xi}(x) = \rho_f^x(\xi).$$

Cette définition, et la formule $\rho_{g,f}^x = \rho_{g,f} \circ \rho_f^x$ pour $x \in D(g) \subset D(f)$, montrent que le diagramme
$$(9) \qquad \begin{tikzcd} M_f \arrow[r, "\sigma_f"] \arrow[d, "\rho_{g,f}"] & S(D(f),M) \arrow[d, "\rho_{D(g),D(f)}^M"] \\ M_g \arrow[r, "\sigma_g"] & S(D(g),M) \end{tikzcd}$$
est commutatif (pour $D(f) \supset D(g)$).

Les groupes abéliens $S(V,M)$ et les homomorphismes $\rho_{W,V}^M$ définissent un

\newpage

%% ===== Page 25 =====

- 8 -

\emph{préfaisceau} $S(M)$ sur $X$ (G,1.1.9); en outre $S(M)$ est en fait un \emph{faisceau} de groupes abéliens, les axiomes des faisceaux (G,II,1.1) étant trivialement vérifiés; $S(V,M)$ s'identifie alors canoniquement au groupe des sections de $S(M)$ au-dessus de l'ensemble ouvert $V$.

\emph{Définition 1}. On appelle \emph{faisceau associé} au $A$-module $M$ et on note $\tilde{M}$ le sous-faisceau de $S(M)$ dont les sections au-dessus d'une partie ouverte $V$ de $X$ sont les sections $s \in S(V,M)$ telles que, pour tout $x \in V$, il existe un voisinage ouvert de $x$ contenu dans $V$, de la forme $D(f)$, et pour lequel la restriction de $s$ à $D(f)$ soit de la forme $\tilde{\xi}$, où $\xi \in M_f$.

Qu'il s'agisse bien d'un sous-faisceau de $S(M)$ résulte de la commutativité du diagramme (9): si $D(g) \subset D(f)$, la restriction à $D(g)$ de tout élément de la forme $\theta_f(\xi)$, où $\xi \in M_f$, est de la forme $\theta_g(\eta)$ où $\eta \in M_g$, donc l'image de $D(f)$ par $\tilde{\eta} = \theta_f(\eta)$, qui est une partie ouverte de $S(M)$, est contenue dans $M$.

Si on applique ce qui précède au cas $M = A$, tous les homomorphismes qui interviennent sont des homomorphismes d'anneaux, donc le faisceau $\tilde{A}$ est muni canoniquement d'une structure de faisceau d'anneaux. On dira que ce faisceau d'anneaux $\tilde{A}$ est le \emph{faisceau structural} sur le spectre premier $X = \operatorname{Spec}(A)$, et on le notera aussi $\mathcal{O}_X$. Si on revient au cas général d'un $A$-module $M$, il résulte alors de la déf. 1 que, pour tout ensemble ouvert $V \subset X$, $\tilde{M}(V)$ est muni canoniquement d'une structure de $\tilde{A}(V)$-module, et que si $V \supset W$ sont deux parties ouvertes de $X$, l'application de restriction $\tilde{M}(V) \to \tilde{M}(W)$ est compatible avec les homomorphismes d'anneaux $\tilde{A}(V) \to \tilde{A}(W)$; par suite (G,II,2.2), $\tilde{M}$ est un $\tilde{A}$-Module.

Soient $M, N$ deux $A$-modules, $u$ un homomorphisme $M \to N$; pour tout $f \in A$ il correspond à $u$ un homomorphisme $u_f$ du $A_f$-module $M_f$ dans le $A_f$-module $N_f$; de même, pour tout $x \in X$, il correspond à $u$ un $A_x$-homomorphisme $u_x: M_x \to N_x$ (chap. 0, §1, n°3). On vérifie aussitôt que si $x \in D(f)$, $\xi \in M_f$, $\eta = u_f(\xi)$, on a $\tilde{\eta}(x) = u_x(\tilde{\xi}(x))$; si on désigne par $\tilde{u}_f$ l'homomorphisme de $S(D(f), M)$

\newpage

%% ===== Page 26 =====

dans $S(D(f),N)$ tel que $(\tilde{u}_f(F))(x) = u_x(F(x))$, il est clair que les $\tilde{u}_f$ définissent un homomorphisme du faisceau $S(M) \to S(N)$, et que ce qui précède montre que cet homomorphisme applique $\tilde{M}$ dans $\tilde{N}$; nous désignerons par $\tilde{u}$ sa restriction à $\tilde{M}$. Si $P$ est un troisième $A$-module, $v$ un homomorphisme $N \to P$, et $w = v \circ u$, il est immédiat que $\tilde{w} = \tilde{v} \circ \tilde{u}$. On a donc défini ainsi un \emph{foncteur covariant} $M \to \tilde{M}$ de la catégorie des $A$-modules dans celle des $\tilde{A}$-Modules.

\emph{Proposition 6.} Pour tout $f \in A$, l'ensemble ouvert $D(f) \subset X$ s'identifie au spectre premier $\operatorname{Spec}(A_f)$, et le faisceau $\tilde{M}_f$ associé à $M_f$ s'identifie canoniquement à la restriction $\tilde{M}|D(f)$.

La première assertion est un cas particulier du cor.3 de la prop. du n°2. Pour tout $x \in D(f)$, $M_x$ s'identifie canoniquement au module des fractions de $M_f$ dont les dénominateurs sont dans l'idéal de $A_f$ qui correspond canoniquement à $j_x$; et de même si $D(f) \supset D(g)$, $M_g$ s'identifie au module des fractions de $M_g$ dont les dénominateurs sont les puissances de l'image de $g$ dans $A_f$ (chap.0, §1, n°4). L'identification de $\tilde{M}_f$ à $\tilde{M}|D(f)$ résulte alors aussitôt des définitions.

\emph{Théorème 1.} Pour tout $A$-module $M$, et tout $f \in A$, l'homomorphisme $\theta_f$ de $M_f$ dans $\Gamma(D(f), \tilde{M})$ est bijectif.

On notera que, si $M = A$, $\theta_f$ est un homomorphisme de structure d'anneau; le th.1 entraînera donc que, si on identifie les anneaux $A_f$ et $\Gamma(D(f), \tilde{A})$ au moyen de $\theta_f$, l'homomorphisme $\theta_f : M \to \Gamma(D(f), \tilde{M})$ sera un isomorphisme de modules.

Notons d'abord que $\theta_f$ est injectif : en effet, si $\xi \in M_f$ est tel que $\xi(x) = 0$ pour tout $x \in D(f)$, cela signifie que pour tout idéal premier $p$ de $A_f$, il existe $h \notin p$ tel que $h\xi = 0$; l'annulateur de $\xi$ n'est donc contenu dans aucun idéal premier de $A_f$, et par suite est $A_f$ lui-même, autrement dit $\xi = 0$.

Notons d'autre part que, pour démontrer le th.1, il suffit de le faire lorsque $f = 1$; on passe au cas général en "localisant" à l'aide de la prop.6.

\newpage

%% ===== Page 27 =====

Reste donc à prouver que pour $f = 1$, $\theta_f$ est surjectif. Soit donc $s$ une section de $\tilde{M}$ au-dessus de $X$ tout entier; en vertu de la déf. 1 et de la prop. 3 du n° 1, il existe un recouvrement fini $(D(f_i))_{i \in I}$ de $X$ tel que, pour tout $i \in I$, la restriction $s_i$ de $s$ à $D(f_i)$ soit de la forme $\theta_{f_i}(\xi_i)$, où $\xi_i \in M_{f_i}$. Si $i, j$ sont deux indices de $I$, en écrivant que les restrictions de $s_i$ et $s_j$ à $D(f_i) \cap D(f_j) = D(f_i f_j)$ sont égales, et tenant compte de ce que $\theta_f$ est injectif et du diagramme (9), on obtient
$$(10) \qquad \rho_{f_i f_j, f_i}(\xi_i) = \rho_{f_i f_j, f_j}(\xi_j).$$
Par définition, on peut écrire pour chaque $i \in I$, $\xi_i = z_i / f_i^{n_i}$ où $z_i \in M$, et en multipliant chaque $z_i$ par une puissance de $f_i$, on peut supposer tous les $n_i$ égaux à un même $n$. Par définition, (10) signifie qu'il existe $m_{ij}$ tel que $(f_i f_j)^{m_{ij}} (f_j^{n_i} z_i - f_i^{n_j} z_j) = 0$, et on peut encore supposer tous les $m_{ij}$ égaux à un même $m$; remplaçant alors $z_i$ par $f_i^m z_i$, on se ramène au cas où $n = 0$, autrement dit au cas où l'on a
$$(11) \qquad f_j^n z_i = f_i^n z_j$$
quels que soient $i$ et $j$. Or on a $D(f_i^n) = D(f_i)$, et comme les $D(f_i)$ forment un recouvrement de $X$, l'idéal engendré par les $f_i^n$ est $A$; autrement dit, il existe des $g_i \in A$ tels que $\sum g_i f_i^n = 1$. Considérons alors l'élément $z = \sum g_i z_i$ de $M$; d'après (11), on a $f_j^n z = \sum g_i f_j^n z_i = \sum g_i f_i^n z_j = z_j$; d'où par définition $\xi_j = z/1$. Le diagramme (9) montre alors que, pour tout $i$, $s_i$ est la restriction à $D(f_i)$ de $\theta_1(z)$, ce qui prouve que $s = \theta_1(z)$ et achève la démonstration.

\emph{Corollaire 1.} Pour tout $x \in X$, la fibre $\tilde{M}(x)$ du faisceau $\tilde{M}$ s'identifie canoniquement au module $M_x$.

En effet, $\tilde{M}(x)$ est la limite inductive des $\Gamma(D(f), \tilde{M})$ pour $x \in D(f)$ et $M_x$ est limite inductive des $M_f$ (chap. 0, § 1, n° 4).

\emph{Corollaire 2.} Si $M$ et $N$ sont deux $A$-modules, $\operatorname{Hom}_{\tilde{A}}(\tilde{M}, \tilde{N})$ est canoniquement isomorphe à $\operatorname{Hom}_A(M, N)$.

\newpage

%% ===== Page 28 =====

- 11 -

En effet, on a défini au début de ce $n^{\circ}$ un homomorphisme canonique $u \to \tilde{u}$ de $\operatorname{Hom}_A(M, N)$ dans $\operatorname{Hom}(\tilde{M}, \tilde{N})$. On a d'autre part l'homomorphisme canonique $v \to \Gamma(v)$ de $\operatorname{Hom}(\tilde{M}, \tilde{N})$ dans $\operatorname{Hom}_A(\Gamma(\tilde{A}), \Gamma(\tilde{N}))$, de la théorie des faisceaux ; mais comme $\Gamma(\tilde{A})$ (resp. $\Gamma(\tilde{N})$) s'identifie canoniquement à $A$ (resp. $M, N$), cela donne un homomorphisme canonique $\operatorname{Hom}(\tilde{M}, \tilde{N}) \to \operatorname{Hom}_A(M, N)$. Il reste à vérifier que les deux homomorphismes ainsi définis sont réciproques l'un de l'autre, ce qui découle aisément des définitions.

On notera en outre qu'il résulte du cor. 1 que si $u \in \operatorname{Hom}_A(M, N)$, la restriction de l'homomorphisme correspondant $\tilde{u}$ à la fibre $\tilde{M}_x$ s'identifie à l'homomorphisme $u_x: M_x \to N_x$.

\emph{Corollaire 3.} Le foncteur $M \to \tilde{M}$ est exact.

En vertu de la remarque précédente, il suffit de vérifier que $M \to M_x$ est un foncteur exact, résultat connu (chap. 0, §1, $n^{\circ}3$).

\emph{Corollaire 4.} a) Soit $u$ un homomorphisme d'un $A$-module $M$ dans un $A$-module $N$; alors les faisceaux associés à $\operatorname{Ker} u$, $\operatorname{Im} u$, $\operatorname{Coker} u$, sont respectivement $\operatorname{Ker} \tilde{u}$, $\operatorname{Im} \tilde{u}$, $\operatorname{Coker} \tilde{u}$.

b) Si $M$ est une limite inductive (resp. somme directe) d'une famille de $A$-modules $M_{\lambda}$, $\tilde{M}$ est limite inductive (resp. somme directe) des $\tilde{M}_{\lambda}$.

a) Il suffit d'appliquer le fait que $M \to \tilde{M}$ est un foncteur exact, aux suites exactes de $A$-modules
$$0 \to \operatorname{Ker} u \to M \to \operatorname{Im} u \to 0$$
$$0 \to \operatorname{Im} u \to N \to \operatorname{Coker} u \to 0.$$

b) Soit $(M_{\lambda}, \epsilon_{\mu\lambda})$ un système inductif de $A$-modules, de limite inductive $M$, et soit $\epsilon_{\lambda}$ l'homomorphisme canonique $M_{\lambda} \to M$. Comme on a $\tilde{\epsilon}_{\mu\lambda} \circ \tilde{\epsilon}_{\lambda} = \tilde{\epsilon}_{\mu\lambda}$ et $\tilde{\epsilon}_{\lambda} = \tilde{\epsilon}_{\mu\lambda} \circ \tilde{\epsilon}_{\mu}$ pour $\lambda \le \mu \le \nu$, $(\tilde{M}_{\lambda}, \tilde{\epsilon}_{\mu\lambda})$ est un système inductif de faisceaux sur $X$, et si on désigne par $h_{\lambda}$ l'homomorphisme canonique $\tilde{M}_{\lambda} \to \lim_{\to} \tilde{M}_{\lambda}$, il y a un homomorphisme unique $v: \lim_{\to} \tilde{M}_{\lambda} \to \tilde{M}$ tel que $v \circ h_{\lambda} = \tilde{\epsilon}_{\lambda}$. Pour voir que $v$ est bijectif, il suffit de vérifier que, pour tout $x \in X$, $v_x$ est une bijection de $(\lim_{\to} \tilde{M}_{\lambda})_x$ sur $\tilde{M}_x$; mais on a $(\tilde{M})_x = M_x$ (cor. 1) et $(\lim_{\to} \tilde{M}_{\lambda})_x = $

\newpage

%% ===== Page 29 =====

- 12 -

$\varinjlim (\tilde{M}_\lambda)_x = \varinjlim (M_\lambda)_x = M_x$ (chap.0, §1, n°4); d'autre part, il résulte aussi des définitions que $(\tilde{\varepsilon}_\lambda)_x$ et $(h_\lambda)_x$ sont tous deux égaux à l'application canonique de $(M_\lambda)_x$ dans $M_x$ ; comme $(\tilde{\varepsilon}_\lambda)_x = v_x \circ (h_\lambda)_x$, $v_x$ est l'identité.

Enfin, si $M$ est somme directe de deux $A$-modules $N, P$, il est immédiat que $\tilde{M} = \tilde{N} \oplus \tilde{P}$ ; toute somme directe étant limite inductive de sommes directes finies, les assertions de b) sont démontrées.

On notera que le cor. 4 prouve que les faisceaux associés aux $A$-modules forment une catégorie abélienne (T, I, 1.4).

Corollaire 5. \emph{Dans la catégorie des faisceaux associés aux $A$-modules, le foncteur $\Gamma$ est exact.}

En effet, soit $\tilde{M} \xrightarrow{u} \tilde{N} \xrightarrow{v} \tilde{P}$ une suite exacte de faisceaux associés à des $A$-modules $M, N, P$. Si $Q = \operatorname{Im} u$ et $R = \operatorname{Ker} v$, on a $\tilde{Q} = \operatorname{Im} \tilde{u} = \operatorname{Ker} \tilde{v} = \tilde{R}$ (cor. 4), d'où le corollaire.

Corollaire 6. \emph{Si $M$ et $N$ sont deux $A$-modules, le faisceau associé à $M \otimes_A N$ s'identifie canoniquement à $\tilde{M} \otimes \tilde{N}$. Si de plus $M$ est un conoyau d'un homomorphisme $A^p \to A^q$, le faisceau associé à $\operatorname{Hom}_A(M, N)$ s'identifie canoniquement à $\operatorname{Hom}(\tilde{M}, \tilde{N})$.}

Le faisceau $F = \tilde{M} \otimes \tilde{N}$ est associé au préfaisceau $F(U) = \Gamma(U, \tilde{M}) \otimes_{\Gamma(U, \tilde{A})} \Gamma(U, \tilde{N})$ ($U$ parcourant une base de la topologie de $X$). Or, si $U = D(f)$ avec $f \in A$, $F(D(f))$ s'identifie canoniquement à $M_f \otimes_{A_f} N_f$ en vertu du th. 1 et de la prop. 6. Mais on sait que ce $A_f$-module est canoniquement isomorphe à $(M \otimes_A N)_f$, qui lui-même est canoniquement isomorphe à $\Gamma(D(f), (M \otimes_A N)^{\sim})$ (th. 1 et prop. 6). En outre, les isomorphismes canoniques
$$ F(D(f)) \to \Gamma(D(f), (M \otimes_A N)^{\sim}) $$
ainsi obtenus vérifient les conditions de compatibilité avec les opérateurs de restriction (chap. 0, §1, n°4), donc définissent un isomorphisme canonique fonctoriel $\tilde{M} \otimes \tilde{N} \to (M \otimes_A N)^{\sim}$.

De même, le faisceau $G = \operatorname{Hom}(\tilde{M}, \tilde{N})$ est associé au préfaisceau $G(U) = $

\newpage

%% ===== Page 30 =====

- 13 -

\begin{itemize}
    \item[$\operatorname{Hom}_{\mathcal{O}_U}(\tilde{M}|_U, \tilde{N}|_U)$], $U$ parcourant une base de la topologie de $X$. Or si on prend $U = D(f)$, $G(D(f))$ s'identifie canoniquement à $\operatorname{Hom}_{A_f}(M_f, N_f)$ (cor. 2 et prop. 6), lequel s'identifie lui-même à $(\operatorname{Hom}_A(M, N))_f$ en vertu de l'hypothèse sur $M$ (chap. 0, §1, n° 3). Finalement, $(\operatorname{Hom}_A(M, N))_f$ s'identifie à $\Gamma(D(f), (\operatorname{Hom}_A(M, N))^{\sim})$ (th. 1 et prop. 6), et les isomorphismes canoniques ainsi définis sont compatibles avec les opérateurs de restriction (chap. 0, §1, n° 4); il définissent donc un isomorphisme canonique $\operatorname{Hom}(\tilde{M}, \tilde{N}) \xrightarrow{\sim} (\operatorname{Hom}_A(M, N))^{\sim}$.
\end{itemize}

\subsection*{4. Faisceaux quasi-cohérents sur un spectre premier.}

\textbf{Théorème 2.} Soient $X$ le spectre premier d'un anneau $A$, $V$ une partie ouverte quasi-compacte de $X$, $F$ un faisceau de $\mathcal{O}_X$-modules défini sur $V$. Les quatre conditions suivantes sont équivalentes :

a) Il existe un $A$-module $M$ tel que $F$ soit isomorphe à $\tilde{M}|_V$.

b) Il existe un recouvrement ouvert fini $(V_i)$ de $V$ par des ensembles de la forme $V_i = D(f_i)$ contenus dans $V$ tels que, pour tout $i$, $F|_{V_i}$ soit isomorphe à un faisceau de la forme $\tilde{M}_i$, où $M_i$ est un $A_{f_i}$-module.

c) Le faisceau $F$ est quasi-cohérent (chap. 0, §2, n° 6).

d) Les deux propriétés suivantes ont lieu :

1) Pour tout $f \in A$ tel que $D(f) \subset V$ et toute section $s \in \Gamma(D(f), F)$ il existe un entier $n \ge 0$ tel que $f^n s$ se prolonge en une section de $F$ sur $V$.

2) Pour tout $f \in A$ tel que $D(f) \subset V$ et toute section $t \in \Gamma(V, F)$, telle que la restriction de $t$ à $D(f)$ soit 0, il existe un entier $n \ge 0$ tel que $f^n t = 0$.

(Dans l'écriture des conditions d 1) et d 2), on a tacitement identifié $A$ et $\Gamma(\tilde{A})$ en vertu du th. 1 du n° 3).

Le fait que a) entraîne b) est conséquence immédiate de la prop. 6 du n° 3 et du fait que les $D(f)$ forment une base de la topologie de $X$. Comme tout $A$-module est isomorphe au conoyau d'un homomorphisme $A^{(I)} \to A^{(J)}$, le cor. 4 du th. 1 (n° 3) montre que tout faisceau associé à un $A$-module est

\newpage

%% ===== Page 31 =====

- 14 -

quasi-cohérent; donc b) entraîne c). Réciproquement, si $\mathcal{F}$ est quasi-cohérent, tout $x \in V$ possède un voisinage de la forme $D(f) \subset V$ tel que $\mathcal{F}|D(f)$ soit isomorphe à un faisceau conoyau d'un homomorphisme $\mathcal{A}_f^{(I)} \to \mathcal{A}_f^{(J)}$, donc au faisceau associé au module $N$, conoyau de l'homomorphisme $A_f^{(I)} \to A_f^{(J)}$ correspondant (cor. 2 et 4 du th. 1, n°3); comme $V$ est quasi-compact, il est clair que c) entraîne b).

Pour prouver que b) entraîne d) 1) et d) 2), supposons d'abord que l'on ait $V = D(g)$ pour un $g \in A$ et que $\mathcal{F}$ soit isomorphe à un faisceau $\tilde{N}$ associé à un $A$-module $N$; remplaçant $X$ par $V$ et $A$ par $A_g$ (prop. 6), on peut se ramener au cas où $g = 1$. Alors $\Gamma(D(f), \tilde{N})$ et $N_f$ s'identifient canoniquement par le th. 1 et la prop. 6 du n°3, donc une section $s \in \Gamma(D(f), \tilde{N})$ s'identifie à un élément de la forme $z/f^n$, où $z \in N$; la section $f^n s$ s'identifie à l'élément $z/1$ de $N_f$ et est par suite restriction à $D(f)$ de la section de $\tilde{N}$ sur $X$ identifiée à l'élément $z \in N$; d'où d) 1) dans ce cas. De même, $t \in \Gamma(X, \tilde{N})$ est identifié à un élément $z' \in N$, la restriction de $t$ à $D(f)$ est identifiée à l'image $z'/1$ de $z'$ dans $N_f$, et dire que cette image est nulle signifie qu'il existe $n \ge 0$ tel que $f^n z' = 0$ dans $N$, ou, ce qui revient au même, $f^n t = 0$.

Pour achever de prouver que b) entraîne d) 1) et d) 2), il suffira d'établir le lemme suivant:

\emph{Lemme 2.} Supposons que $V$ soit réunion finie d'ensembles de la forme $D(g_i)$, et que chacun des faisceaux $\mathcal{F}|D(g_i)$ et $\mathcal{F}|D(g_i) \cap D(g_j) = \mathcal{F}|D(g_i g_j)$ vérifie d) 1) et d) 2); alors $\mathcal{F}$ possède les deux propriétés suivantes:

d') 1) Pour tout $f \in A$ et toute section $s \in \Gamma(D(f) \cap V, \mathcal{F})$, il existe un entier $n \ge 0$ tel que $f^n s$ se prolonge en une section de $\mathcal{F}$ sur $V$.

d') 2) Pour tout $f \in A$ et toute section $t \in \Gamma(V, \mathcal{F})$ telle que la restriction de $t$ à $D(f) \cap V$ soit 0, il existe un entier $n \ge 0$ tel que $f^n t = 0$.

Prouvons d'abord d') 2); comme $D(f) \cap D(g_i) = D(fg_i)$, il existe pour chaque $i$ un entier $n_i$ tel que la restriction de $(fg_i)^{n_i} t$ soit nulle;

\newpage

%% ===== Page 32 =====

\marginpar{15}
comme l'image de $g_i$ dans $A_{g_i}$ est inversible, la restriction de $f^{n_i}t$ à $D(g_i)$ est aussi nulle; prenant pour $n$ le plus grand des $n_i$, on a d' 2).

Pour démontrer d' 1), appliquons d' 1) au faisceau $F|D(g_i)$: il existe un entier $n_i$ et une section $s'_i$ de $F|D(g_i)$ prolongeant la restriction de $(fg_i)^{n_i}s$ à $D(fg_i)$; comme l'image de $g_i$ dans $A_{g_i}$ est inversible, il y a une section $s_i$ de $F|D(g_i)$ telle que $s'_i = g_i^{n_i}s_i$, et $s_i$ prolonge la restriction de $f^{n_i}s$ à $D(fg_i)$; on peut en outre supposer tous les $n_i$ égaux à un même entier $n$. Par construction, la restriction de $s_i - s_j$ à $D(f) \cap D(g_i) \cap D(g_j) = D(fg_ig_j)$ est nulle; d'après d' 2) appliqué au faisceau $F|D(g_ig_j)$, il existe un entier $m_{ij}$ tel que la restriction à $D(g_ig_j)$ de $(fg_ig_j)^{m_{ij}}(s_i - s_j)$ soit nulle; comme l'image de $g_ig_j$ dans $A_{g_ig_j}$ est inversible, la restriction de $f^{m_{ij}}(s_i - s_j)$ à $D(g_ig_j)$ est nulle. On peut alors supposer tous les $m_{ij}$ égaux à un même entier $m$, et il existe donc une section $s \in \Gamma(V,F)$ prolongeant les $f^m s_i$; cette section prolonge par suite $f^{n+m}s$, d'où d' 1).

Reste à montrer que d' 1) et d' 2) entraînent a). Montrons d'abord que d' 1) et d' 2) entraînent que ces mêmes conditions sont vérifiées pour tout faisceau $F|D(g)$, où $g \in A$ est tel que $D(g) \subset V$: c'est évident pour d' 1); d'autre part, si $t \in \Gamma(D(g),F)$ est telle que sa restriction à $D(f) \subset D(g)$ soit nulle, il existe par d' 1) un entier $m \ge 0$ tel que $g^m t$ se prolonge en une section $s$ de $F$ sur $V$; appliquant d' 2), on voit qu'il existe un entier $n \ge 0$ tel que $f^n g^m t = 0$, et comme l'image de $g$ dans $A_g$ est inversible, $f^n t = 0$.

Cela étant, comme $V$ est quasi-compact, le lemme 2 prouve que les conditions d' 1) et d' 2) sont vérifiées. Considérons alors le $A$-module $M = \Gamma(V,F)$, et définissons un homomorphisme de faisceaux $u: \tilde{M} \to j_*(F)$, où $j$ est l'injection canonique $V \to X$. Comme les $D(f)$ forment une base de la topologie de $X$, il suffit pour chaque $f \in A$ de définir un homomorphisme $u_f: M_f \to \Gamma(D(f), j_*(F)) = \Gamma(D(f) \cap V, F)$, avec les conditions de compatibilité

\newpage

%% ===== Page 33 =====

\begin{itemize}
\item[-- 16 --]
\end{itemize}

usuel. Comme l'image canonique de $f$ dans $A_g$ est inversible, l'homomorphisme de restriction $M = \Gamma(V, \mathcal{F}) \to \Gamma(D(f) \cap V, \mathcal{F})$ se factorise en $M \to M_g \to \Gamma(D(f) \cap V, \mathcal{F})$, et la vérification des conditions de compatibilité pour $D(g) \subset D(f)$ est immédiate. Cela étant, montrons que la condition $d'1)$ (resp. $d'2)$) entraîne que chacun des $u_g$ est surjectif (resp. injectif), ce qui prouvera que $u$ est bijectif, et par suite que $\mathcal{F}$ est la restriction à $V$ d'un faisceau quasi-cohérent sur $X$. Or, si $s \in \Gamma(D(f) \cap V, \mathcal{F})$, il existe d'après $d'1)$ un entier $n \ge 0$ tel que $f^n s$ se prolonge en une section $z \in M$; on a alors $u_g(z/f^n) = s$, donc $u_g$ est surjectif. De même, si $z \in M$ est tel que $u_g(z/1) = 0$, cela signifie que la restriction à $D(f) \cap V$ de la section $z$ est nulle; d'après $d'2)$, il existe un entier tel que $f^n z = 0$, d'où $z/1 = 0$ dans $M_g$, et $u_g$ est donc injectif. C.Q.F.D.

\textit{Corollaire.} Tout faisceau quasi-cohérent sur un ouvert quasi-compact de $X$ est induit par un faisceau quasi-cohérent sur $X$.

\section*{5. Faisceaux cohérents sur un spectre premier.}

\textit{Théorème 3.} Soient $A$ un anneau noethérien, $X = \operatorname{Spec}(A)$ son spectre premier, $V$ une partie ouverte de $X$, $\mathcal{F}$ un faisceau de $\mathcal{O}_X$-modules défini sur $V$. Les conditions suivantes sont équivalentes:

a) $\mathcal{F}$ est cohérent.

b) $\mathcal{F}$ est de type fini et quasi-cohérent.

c) Il existe un $A$-module $M$ de type fini tel que $\mathcal{F}$ soit isomorphe au faisceau $\widetilde{M}|_V$.

a) implique trivialement b). Pour voir que b) entraîne c), remarquons déjà, puisque $V$ est quasi-compact, que $\mathcal{F}$ est isomorphe à un faisceau $\widetilde{N}|_V$, où $N$ est un $A$-module (th.2). Pour tout $x \in V$, il existe par hypothèse un voisinage $D(f) \subset V$ et un homomorphisme surjectif $\widetilde{A}^n_x \to \mathcal{F}_x \to 0$; cet homomorphisme est de la forme $\widetilde{u}$, où $u$ est un homomorphisme surjectif $A^n_f \to N_f \to 0$ ($n^\circ 3$, cor.2 et 5 du th.1); donc $N_x$ est un $A_x$-module de type fini. Comme $V$

\newpage

%% ===== Page 34 =====

est quasi-compact, il y a un recouvrement ouvert fini $(D(f_i))$ de $V$ tel que chacun des $N_{f_i}$ soit un $A_{f_i}$-module de type fini ; par suite il y a une partie finie $S$ de $N$ dont l'image canonique dans chacun des $N_{f_i}$ soit un système de générateurs de ce module. Soit $M$ le sous-module de $N$ engendré par $S$ ; la suite $0 \to M \to N$ étant exacte, il en est de même de $0 \to \widetilde{M} \to \widetilde{N}$ ($n^\circ 3, \text{cor.} 3$ du th. 1), donc aussi de $0 \to \widetilde{M}|_V \to \widetilde{N}|_V$ ; mais comme par construction la restriction de $\widetilde{M}|_V \to \widetilde{N}|_V$ à chacun des $D(f_i)$ est surjective, $\widetilde{M}|_V \to \widetilde{N}|_V$ est surjective, autrement dit $\mathcal{F}$ est isomorphe à $\widetilde{M}|_V$.

Montrons enfin que c) entraîne a). Il est clair que $\mathcal{F}$ est alors de type fini (th. 1) ; en outre, la question étant locale, on peut se borner au cas où $V = D(f)$. Comme $A_f$ est noethérien, on voit finalement que tout revient à prouver que le noyau d'un homomorphisme $\widetilde{A^n} \to \widetilde{M}$, où $M$ est un $A$-module, est de type fini. Or, un tel homomorphisme est de la forme $\widetilde{u}$, où $u$ est un homomorphisme $A^n \to M$ ($n^\circ 3, \text{cor.} 2$ du th. 1), et si $P = \ker u$, on a $\widetilde{P} = \ker \widetilde{u}$ (cor. 4 du th. 1). Comme $A$ est noethérien, $P$ est de type fini, ce qui achève la démonstration.

\emph{Corollaire.} \underline{Sous les hypothèses du th. 3, tout faisceau cohérent de $\mathcal{O}_X$-modules défini sur une partie ouverte $V$ de $X$ est la restriction à $V$ d'un faisceau cohérent de $\mathcal{O}_X$-modules sur $X$.}

6. \underline{Propriétés fonctorielles des faisceaux quasi-cohérents sur un spectre premier.}

Soient $A, A'$ deux anneaux, $\varphi$ un homomorphisme de $A'$ dans $A$, $\varphi^*$ l'application continue associée de $X = \operatorname{Spec}(A)$ dans $X' = \operatorname{Spec}(A')$ ($n^\circ 2$). Nous allons définir un homomorphisme canonique $\mathcal{O}_{X'} \to \varphi_*(\mathcal{O}_X)$ (autrement dit, un $\varphi^*$-morphisme de $\mathcal{O}_{X'}$ dans $\mathcal{O}_X$ ; chap. 0, \S 2, $n^\circ 4$). Pour tout $f' \in A'$, posons $f = \varphi(f')$ ; il résulte de (5) que l'on a $\varphi^{-1}(D(f')) = D(f)$. Les anneaux $\Gamma(D(f'), \widetilde{A'})$ et $\Gamma(D(f), \widetilde{A})$ s'identifient aux anneaux de fractions $A'_{f'}$ et $A_f$ ($n^\circ 3, \text{th.} 1$ et \text{prop.} 6). Or, l'homomorphisme $\varphi$ définit canoniquement un homomorphisme $\varphi_f : A'_{f'} \to A_f$ (chap. 0, \S 1, $n^\circ 5$), autrement dit, on a un homomorphisme

\newpage

%% ===== Page 35 =====

phisme d'anneaux $\Gamma(D(f'), \tilde{A}') \to \Gamma(\varphi^{-1}(D(f')), \tilde{A})$. En outre, ces homomorphismes satisfont aux conditions de compatibilité: pour $D(g') \subset D(f')$ le diagramme
$$
\begin{tikzcd}
\Gamma(D(f'), \tilde{A}') \arrow[r] \arrow[d] & \Gamma(\varphi^{-1}(D(f')), \tilde{A}) \arrow[d] \\
\Gamma(D(g'), \tilde{A}') \arrow[r] & \Gamma(\varphi^{-1}(D(g')), \tilde{A})
\end{tikzcd}
$$
est commutatif (chap.0, §1, n°5); d'où l'homomorphisme de faisceaux annoncé, que nous désignerons par $\tilde{\varphi}$. Le couple $\tilde{\varphi} = (\varphi, \tilde{\varphi})$ est donc un morphisme de l'espace annelé $(X, \mathcal{O}_X)$ dans l'espace annelé $(X', \mathcal{O}_{X'})$ (chap.0, §2, n°5).

Notons en outre que si on pose $x' = \varphi(x)$, l'homomorphisme $\tilde{\varphi}_{x'}$ n'est autre que l'homomorphisme $A'_{x'} \to A_x$ déduit canoniquement de $\varphi$ : $A' \to A$ (chap.0, §1, n°5). En effet, tout $z' \in A'_{x'}$ s'écrit $g'/f'$, où $f' \notin \mathfrak{p}_{x'} \subset A'$; $D(f')$ est donc un voisinage de $x'$, et l'homomorphisme $\tilde{\varphi}_{D(f')} : \Gamma(D(f'), \tilde{A}') \to \Gamma(\varphi^{-1}(D(f')), \tilde{A})$ déduit de $\tilde{\varphi}$ n'est autre que $\varphi_{f'}$; en considérant la section $s \in \Gamma(D(f'), \tilde{A}')$ correspondant à $g'/f' \in A'_{f'}$ on obtient bien $\tilde{\varphi}(z') = \varphi(g')/\varphi(f')$ dans $A_x$.

Soit $M$ un $A$-module; rappelons que l'on désigne par $M_{[\varphi]}$ le $A'$-module défini sur $M$ par l'homomorphisme $\varphi : A' \to A$.

\emph{Proposition 7.} Il existe un isomorphisme canonique fonctoriel du $\mathcal{O}_{X'}$-module $\tilde{M}_{[\varphi]}$ sur l'image directe $\varphi_*(\tilde{M})$.

Posons pour abréger $M' = M_{[\varphi]}$, et pour tout $f' \in A'$, posons encore $f = \varphi(f')$. Les modules de sections $\Gamma(D(f'), \tilde{M}')$ et $\Gamma(D(f), \tilde{M})$ s'identifient respectivement aux modules $M'_{f'}$ et $M_f$ (sur $A'_{f'}$ et $A_f$ respectivement); en outre, le $A'_{f'}$-module $(M_f)_{[\varphi_{f'}]}$ est canoniquement isomorphe à $M'_{f'}$ (chap.0, §1, n°5). 
On a donc un isomorphisme fonctoriel de $\Gamma(D(f'), \tilde{A}')$-modules
$$\Gamma(D(f'), \tilde{M}') \xrightarrow{\sim} \Gamma(\varphi^{-1}(D(f')), \tilde{M})_{[\varphi]}$$
et ces isomorphismes satisfont aux conditions de compatibilité usuelles (chap.0, §1, n°5), donc définissent l'isomorphisme fonctoriel annoncé.

\emph{Corollaire 1.} Le foncteur image directe $\varphi_*$ est exact sur la catégorie des faisceaux quasi-cohérents sur $X$.

\newpage

%% ===== Page 36 =====

\noindent En effet, il est clair que le foncteur $M \to M_{[\varphi]}$ est exact, et on sait qu'il en est de même du foncteur $M' \to \tilde{M}'$ ($n^{\circ}3$, cor.3 du th.1).

\noindent \emph{Corollaire} 2. Soit $U'$ une partie ouverte quasi-compacte de $X'$ telle que $U = \varphi^{-1}(U')$ soit quasi-compact; l'image directe $\varphi_*$ d'un faisceau quasi-cohérent de $\mathcal{O}_{X'}$-modules sur $U$ est un faisceau quasi-cohérent de $\mathcal{O}_X$-modules sur $U'$.

\noindent En effet, un faisceau quasi-cohérent $F$ sur $U$ est de la forme $\tilde{M}|U$ où $M$ est un $A$-module ($n^{\circ}4$, th.2); comme $\varphi_*(F) = \tilde{M}_{[ \varphi ]}|U'$, le corollaire résulte de la prop.7 et du th.2.

\noindent Soit maintenant $N'$ un $A'$-module, et associons-lui le $A$-module $N = N' \otimes_{A'} A_{[\varphi]}$.

\noindent \emph{Proposition} 8. Il existe un isomorphisme canonique fonctoriel en $\mathcal{O}_{X'}$-Module $\varphi^*(\tilde{N}') \text{ sur } \tilde{N}$.

\noindent Remarquons d'abord que $j: z' \to z' \otimes 1$ est un $A'$-homomorphisme de $N'$ dans $N_{[\varphi]}$; en effet, par définition, pour $f' \in A'$, on a $(f'z') \otimes 1 = z' \otimes \varphi(f') = \varphi(f')(z' \otimes 1)$. On en déduit ($n^{\circ}3$, cor.2 du th.1) un homomorphisme $\tilde{j}: \tilde{N}' \to \tilde{N}_{[\varphi]}$ de $\mathcal{O}_{X'}$-Modules, et en vertu de la prop.7, on peut considérer que $\tilde{j}$ applique $\tilde{N}'$ dans $\varphi_*(\tilde{N})$. Il correspond canoniquement à cet homomorphisme $\tilde{j}$ un homomorphisme $h = \tilde{j}^{\flat}$ de $\varphi^*(\tilde{N}')$ dans $\tilde{N}$ (chap.0, §2, $n^{\circ}5$); on va voir que pour chaque fibre $h_x$ est bijectif. Posons $x' = \varphi(x)$, et soit $f' \in A'$ tel que $x' \in D(f')$; soit $f = \varphi(f')$. L'anneau $\Gamma(D(f), A)$ s'identifie à $A_f$, les modules $\Gamma(D(f), \tilde{N})$ et $\Gamma(D(f'), \tilde{N}')$ à $N_f$ et $N'_{f'}$, respectivement; soient $s' \in \Gamma(D(f'), \tilde{N}')$, identifié à $n'/f'^p$ ($n' \in N'$), son image par $\tilde{j}_D(f')$ dans $\Gamma(D(f), \tilde{N})$; $s$ est identifiée à $(n' \otimes 1)/f^p$. Soit d'autre part $t \in \Gamma(D(f), \tilde{A})$ identifié à $g/f^q$ ($g \in A$); alors, par définition, on a $h_{s'}(x') \otimes t(x)) = t(x)s(x)$ (chap.0, §2, $n^{\circ}5$). Mais on peut identifier canoniquement $N_f$ à $N'_{f'} \otimes_{A'_{f'}} (A_f)_{[\varphi]}$ (chap.0, §1, $n^{\circ}5$); $s$ correspond alors à l'élément $(n'/f'^p) \otimes 1$, et la section $s \to t(y)s(y)$ à $(n'/f'^p) \otimes (g/f^q)$.

\newpage

%% ===== Page 37 =====

\begin{itemize}
\item[] Les diagrammes de compatibilité du chap.0, §1, n°5 montrent alors que $h_x$ n'est autre que l'isomorphisme canonique $\tilde{N}'_x \otimes_{A'_x} (A_x)[\phi_x] \to N_x = (\tilde{N}' \otimes_{A'} A[\phi])_x$.
\end{itemize}

\noindent \underline{Corollaire 1.} Les sections de $\tilde{\phi}^*(\tilde{N}')$ de la forme $s' \otimes 1$, où $s'$ parcourt le $A'$-module $\Gamma(\tilde{N}')$, engendrent le $A$-module $\Gamma(\tilde{\phi}^*(\tilde{N}'))$.

\noindent En effet, l'application $s' \to s' \otimes 1$ s'identifie à l'application $z' \to z' \otimes 1$ de $N'$ dans $N$, lorsqu'on identifie $N'$ et $N$ à $\Gamma(\tilde{N}')$ et $\Gamma(\tilde{N})$ respectivement ($n^\circ 3$, th. 1).

\noindent \underline{Corollaire 2.} Soit $U'$ une partie ouverte quasi-compacte de $X'$ telle que $U = \phi^{-1}(U')$ soit quasi-compact; l'image réciproque par $\tilde{\phi}$ d'un faisceau quasi-cohérent de $\mathcal{O}_{X'}$-modules sur $U'$ est un faisceau quasi-cohérent de $\mathcal{O}_X$-modules sur $U$.

\noindent En effet, un faisceau quasi-cohérent $G$ sur $U'$ est de la forme $\tilde{N}'|_U$ où $N'$ est un $A'$-module ($n^\circ 4$, th. 2); comme $\tilde{\phi}^*(G) = \tilde{\phi}^*(\tilde{N}')|_U$, le corollaire résulte de la prop. 8 et du th. 2.

\noindent Soient $A''$ un troisième anneau, $\phi'$ un homomorphisme $A'' \to A'$, et posons $\phi'' = \phi \circ \phi'$. Il résulte aussitôt des définitions que ${}^a\phi'' = ({}^a\phi') \circ ({}^a\phi)$, et $\tilde{\phi}'' = \tilde{\phi}' \circ \tilde{\phi}$ (chap. 0, §1, n°5). On en conclut que l'on a $\tilde{\phi}'' = \tilde{\phi}' \circ \tilde{\phi}$.

\section*{7. Caractérisation des morphismes des schémas affines.}

\noindent On dit qu'un espace annelé $(X, \mathcal{O}_X)$ est un \emph{schéma affine} s'il est isomorphe à un espace annelé de la forme $(\operatorname{Spec}(A), \tilde{A})$, où $A$ est un anneau; on dit alors que $\Gamma(X, \mathcal{O}_X)$, qui s'identifie à l'anneau $A$ ($n^\circ 3$, th. 1) est l'anneau du schéma affine $(X, \mathcal{O}_X)$ et on le note $A(X)$ quand aucune confusion n'en résulte.

\noindent Soient $A, B$ deux anneaux, $(X, \mathcal{O}_X)$ et $(Y, \mathcal{O}_Y)$ les schémas affines correspondants sur les spectres premiers $X = \operatorname{Spec}(A)$, $Y = \operatorname{Spec}(B)$. On a vu ($n^\circ 6$) qu'à tout homomorphisme d'anneaux $\phi: B \to A$ correspond un morphisme $\Phi = ({}^a\phi, \tilde{\phi}): (X, \mathcal{O}_X) \to (Y, \mathcal{O}_Y)$. On notera que $\phi$ est entièrement déterminé par $\tilde{\Phi}$, car on a par définition $\phi = \Gamma(\tilde{\phi}): \Gamma(\tilde{B}) \to \Gamma({}^a\phi_*(\tilde{A})) = \Gamma(\tilde{A})$.

\noindent \underline{Théorème 4.} Pour qu'un morphisme $(\psi, \theta)$ de $(X, \mathcal{O}_X)$ dans $(Y, \mathcal{O}_Y)$ soit de la forme $({}^a\phi, \tilde{\phi})$, où $\phi$ est un homomorphisme de $B$ dans $A$, il faut et il suffit que, pour

\newpage

%% ===== Page 38 =====

tout $x \in X$, l'image réciproque par $\theta_x^b$ de l'idéal maximal de $B_{\psi(x)}$ soit
l'idéal maximal de $A_x$.

La condition est nécessaire, car on a vu au $n^\circ 6$ que $\theta_x^b$ est l'homomorphisme de $B_{\psi(x)}$ dans $A_x$ déduit canoniquement de $\psi$; par définition de $^a\psi(x) = \psi^{-1}(\mathfrak{j}_x)$, cet homomorphisme a la propriété requise.

Prouvons que la condition est suffisante. Par définition, $\theta$ est un homomorphisme de $\mathcal{O}_Y$ dans $\psi_*(\mathcal{O}_X)$, et on en déduit canoniquement un homomorphisme d'anneaux

$$\varphi = \Gamma(\theta) : B = \Gamma(Y, \mathcal{O}_Y) \to \Gamma(Y, \psi_*(\mathcal{O}_X)) = \Gamma(X, \mathcal{O}_X) = A.$$

L'hypothèse sur $\theta^b$ permet de déduire de cet homomorphisme, par passage aux quotients, un homomorphisme $\theta^x$ du corps des restes $k(\psi(x))$ dans le corps des restes $k(x)$, tel que, pour toute section $f \in \Gamma(Y, \mathcal{O}_Y) = B$, on ait $\theta^x(f(\psi(x))) = \varphi(f)(x)$. La relation $f(\psi(x)) = 0$ est donc équivalente à $\varphi(f)(x) = 0$, ce qui signifie que $\psi(x) = ^a\varphi(x)$ et s'écrit encore $\psi = ^a\varphi$. On sait aussi que le diagramme

$$\begin{tikzcd}
B = \Gamma(Y, \mathcal{O}_Y) \arrow[r, "\varphi"] \arrow[d] & \Gamma(X, \mathcal{O}_X) = A \arrow[d] \\
B_{\psi(x)} \arrow[r, "\theta_x^b"] & A_x
\end{tikzcd}$$

est commutatif, ce qui (vu la caractérisation d'un anneau de fractions comme solution d'un problème d'application universelle) signifie que $\theta_x^b$ est égal à l'homomorphisme $\varphi_x : B_{\psi(x)} \to A_x$ déduit de $\varphi$; comme la donnée des $\theta_x^b$ caractérise entièrement $\theta$, on en conclut que $\theta = \tilde{\varphi}$, par définition de $\tilde{\varphi}$ ($n^\circ 6$).

Nous dirons qu'un morphisme $(\psi, \theta)$ d'espaces annelés satisfaisant à la condition du th. 4 est un morphisme de schémas affines.

\emph{Corollaire.} Si $(X, \mathcal{O}_X), (Y, \mathcal{O}_Y)$ sont deux schémas affines, il existe un isomorphisme canonique de l'ensemble de morphismes de schémas affines $\operatorname{Hom}((X, \mathcal{O}_X), (Y, \mathcal{O}_Y))$ sur l'ensemble d'homomorphismes $\operatorname{Hom}(B, A)$, où $A = \Gamma(\mathcal{O}_X)$ et $B = \Gamma(\mathcal{O}_Y)$.

On peut encore dire que les foncteurs $A \to (\operatorname{Spec}(A), \tilde{A})$ et $(X, \mathcal{O}_X) \to \Gamma(\mathcal{O}_X)$ définissent une équivalence de la catégorie des anneaux commutatifs et de la catégorie duale de la catégorie des schémas affines ($T, I, 1.2$).

\newpage

%% ===== Page 39 =====

\marginpar{Archive Grothendieck sept 59} \hfill 326 bis

§2. Préschémas; morphismes; préschémas sur un préschéma.

1. Définition des préschémas.

Etant donné un espace annelé $(X, \mathcal{O}_X)$, on dit qu'une partie ouverte $V$ de $X$ est un \emph{ouvert affine} si l'espace annelé $(V, \mathcal{O}_X|V)$ est un schéma affine (§1, n°7).

\emph{Définition 1}. On appelle \emph{préschéma} un espace annelé $(X, \mathcal{O}_X)$ tel que tout point de $X$ admette un voisinage ouvert affine. On dit que $X$ est l'espace de base du préschéma $(X, \mathcal{O}_X)$.

\emph{Lemme 1}. Si $(X, \mathcal{O}_X)$ est un préschéma, les ensembles ouverts affines forment une base de la topologie de $X$.

En effet, si $V$ est un voisinage ouvert quelconque de $x \in X$, il existe par hypothèse un voisinage ouvert $W$ de $x$ tel que $(W, \mathcal{O}_X|W)$ soit un schéma affine; désignons par $A$ son anneau. Dans l'espace $W$, $V \cap W$ est un voisinage ouvert de $x$; donc il existe $f \in A$ tel que $D(f)$ soit un voisinage ouvert de $x$ contenu dans $V \cap W$ (§1, n°1, prop.3a)). L'espace annelé $(D(f), \mathcal{O}_X|D(f))$ est alors un schéma affine d'anneau isomorphe à $A_f$ (§1, n°3, prop.6), d'où le lemme.

\emph{Lemme 2}. Si $x$ et $y$ sont deux points distincts de $X$, il existe un voisinage ouvert de l'un des points $x, y$ qui ne contient pas l'autre (axiome $(T_0)$).

C'est évident si les deux points ne sont pas dans un même ouvert affine (déf.1); et s'ils sont dans un même ouvert affine, cela résulte du §1, n°1, cor.3 de la prop.2.

\emph{Proposition 1}. L'application $x \to \overline{\{x\}}$ établit une correspondance biunivoque entre $X$ et l'ensemble des parties fermées irréductibles de $X$. Pour qu'un ensemble ouvert $U \subset X$ rencontre $\overline{\{x\}}$, il faut et il suffit que $x \in U$.

Si $V$ est un voisinage affine de $x$, $\{x\} \cap V$ est l'adhérence de $\{x\}$ dans $V$, donc est irréductible (§1, n°1, cor. de la prop.4), et comme cet ensemble est dense dans $\{x\}$, $\{x\}$ est irréductible. Inversement, si $Y$ est une partie fermée irréductible de $X$, $y \in Y$, et $U$ un voisinage affine de $y$, $U \cap Y$ est partout dense dans $Y$ et est irréductible, donc (§1, n°1, cor. de la prop.4),

\newpage

%% ===== Page 40 =====

$U \cap Y$ est l'adhérence dans $U$ d'un point unique $x$, et a fortiori $Y \subset \overline{\{x\}}$ est l'adhérence de $x$ dans $X$; ce point est le seul point $x'$ tel que $Y = \overline{\{x'\}}$, car un tel point doit être dans $U \cap Y$, et par suite $U \cap Y$ est son adhérence dans $U$, ce qui entraîne $x' = x$. La seconde assertion de la prop. 1 est évidente et valable dans tout espace topologique.

\textbf{Définition 2.} Soient $(X, \mathcal{O}_X)$ un préschéma, et $Y$ une partie fermée irréductible de $X$. L'unique point $y \in X$ tel que $Y = \overline{\{y\}}$ est appelé le \emph{point générique} de $Y$.

Si $Y, Y'$ sont deux parties fermées irréductibles de $X$, la condition $Y \subset Y'$ équivaut donc à $y \in Y'$, si $y$ est le point générique de $Y$.

Si $Y$ est une partie fermée irréductible de $X$, son anneau local $\mathcal{O}_y$ se note aussi $\mathcal{O}_{X,Y}$ et s'appelle l'\emph{anneau local de $X$ le long de $Y$}, ou \emph{l'anneau local de $Y$ dans $X$}.

\textbf{Proposition 2.} Si $(X, \mathcal{O}_X)$ est un préschéma, pour toute partie ouverte $V$ de $X$, l'espace annelé $(V, \mathcal{O}_X|V)$ est un préschéma.

Cela résulte aussitôt du lemme 1 et de la déf. 1. On dira que $(V, \mathcal{O}_X|V)$ est le préschéma \emph{induit} sur $V$ par $(X, \mathcal{O}_X)$, ou la \emph{restriction} de $(X, \mathcal{O}_X)$ à $V$.

\subsection*{2. Morphismes de préschémas.}

\textbf{Définition 3.} Étant donnés deux préschémas $(X, \mathcal{O}_X), (Y, \mathcal{O}_Y)$, on appelle \emph{morphisme (de préschéma) de $(X, \mathcal{O}_X)$ dans $(Y, \mathcal{O}_Y)$} tout morphisme $(\psi, \theta)$ d'espaces annelés tel que, pour tout $x \in X$, l'image réciproque par $\theta_x^b$ de l'idéal maximal de $\mathcal{O}_{Y, \psi(x)}$ soit l'idéal maximal de $\mathcal{O}_{X,x}$.

Par passage au quotient, l'application $\theta_x^b: \mathcal{O}_{Y, \psi(x)} \to \mathcal{O}_{X,x}$ donne donc un monomorphisme $\theta_x^\#: k(\psi(x)) \to k(x)$, ce qui permet donc de considérer $k(x)$ comme une extension de $k(\psi(x))$.

Le composé $(\psi', \theta')$ de deux morphismes de préschémas $(\psi, \theta)$ et $(\psi', \theta')$ est encore un morphisme de préschéma, comme il résulte de l'expression de $\theta^{b}$ à l'aide de $\theta^b$ et $\theta'^b$, et du fait que $\psi_x^\# \circ \theta_{\psi(x)}^\#$ est un isomorphisme (chap. 0, §2, n$^\circ$s 3 et 4). On en conclut que les préschémas forment une

\newpage

%% ===== Page 41 =====

\addtocounter{page}{23}
\section*{catégorie.}

\emph{Convention de notations}: Dans toute la suite de ce chapitre, et lorsque cela ne risquera pas d'entraîner confusion, on \underline{supprimera} dans la notation d'un préschéma (resp. d'un morphisme) le faisceau structural (resp. le morphisme de faisceaux structuraux). Si $U$ est une partie ouverte de l'espace de base d'un préschéma $(X, \mathcal{O}_X)$, lorsqu'on parlera de $U$ comme d'un préschéma, il s'agira du préschéma induit sur $U$. Par contre, lorsqu'on parlera d'un point $x \in X$, où $X$ est un préschéma, il s'agira toujours d'un point de l'espace de base $X$ (et non de l'espace étalé $X$). De même, si $\psi$ est un morphisme de préschémas $X \to Y$, $M$ (resp. $N$) une partie de l'espace de base $X$ (resp. $Y$), les notations $\psi(M)$ et $\psi^{-1}(N)$ désigneront, sauf mention expresse du contraire, les sous-espaces correspondants des espaces de base $Y, X$ respectivement, et non des préschémas ayant ces sous-espaces pour espaces de base.

\emph{Proposition 2}. Soient $(X, \mathcal{O}_X)$ un préschéma, $(S, \mathcal{O}_S)$ un schéma affine d'anneau $A$. Il existe une correspondance biunivoque canonique entre les \underline{morphismes du préschéma $(X, \mathcal{O}_X)$ dans le préschéma $(S, \mathcal{O}_S)$} et les \underline{homomorphismes de $A$ dans l'anneau $\Gamma(X, \mathcal{O}_X)$}.

Notons d'abord que, si $(X, \mathcal{O}_X)$ et $(Y, \mathcal{O}_Y)$ sont deux espaces annelés quelconques, un morphisme $(\psi, \theta)$ dans $(Y, \mathcal{O}_Y)$ définit canoniquement un homomorphisme d'anneaux $\Gamma(\theta): \Gamma(Y, \mathcal{O}_Y) \to \Gamma(X, \psi_*(\mathcal{O}_X)) = \Gamma(X, \mathcal{O}_X)$. Dans le cas considéré, tout revient à voir qu'un homomorphisme quelconque $\varphi$ de $A$ dans $\Gamma(X, \mathcal{O}_X)$ correspond de cette façon à un morphisme de préschéma et un seul. Or, il y a par hypothèse un recouvrement ouvert $(V_\alpha)$ de $X$ par des ouverts affines; par composition de $\varphi$ avec l'homomorphisme de restriction $\Gamma(X, \mathcal{O}_X) \to \Gamma(V_\alpha, \mathcal{O}_X|_{V_\alpha})$ on obtient un homomorphisme $\varphi_\alpha: A \to \Gamma(V_\alpha, \mathcal{O}_X|_{V_\alpha})$ qui correspond à un morphisme unique $(\psi_\alpha, \theta_\alpha)$ du préschéma $(V_\alpha, \mathcal{O}_X|_{V_\alpha})$ dans $(S, \mathcal{O}_S)$, en vertu du th. 4 du \S 1, n$^\circ$ 7. En outre, pour tout couple $(\alpha, \beta)$, tout point de $V_\alpha \cap V_\beta$ admet un voisinage ouvert affine $W$ contenu dans $V_\alpha \cap V_\beta$ (n$^\circ$ 1, lemme 1); il

\newpage

%% ===== Page 42 =====

est clair qu'en composant $\varphi_\alpha$ et $\varphi_\beta$ avec les homomorphismes de restriction à $W$, on obtient le même homomorphisme $\Gamma(S, \mathcal{O}_S) \to \Gamma(W, \mathcal{O}_X)$, donc, en vertu des relations $(\theta_\alpha)_x = (\varphi_\alpha)_x$ pour tout $x \in V_\alpha$, les restrictions à $W$ des morphismes $(\psi_\alpha, \theta_\alpha)$ et $(\psi_\beta, \theta_\beta)$ coïncident. On en conclut qu'il y a un morphisme d'espaces annelés $(\psi, \theta) : (X, \mathcal{O}_X) \to (S, \mathcal{O}_S)$ et un seul dont la restriction à chaque $V_\alpha$ est $(\psi_\alpha, \theta_\alpha)$, et il est clair que ce morphisme est un morphisme de préschémas et est tel que $\Gamma(\theta) = \varphi$.

On appelle \emph{schéma local} un schéma affine dont l'anneau $A$ est local; il existe alors dans $X = \operatorname{Spec}(A)$ un seul point $a$ tel que $\{a\}$ soit fermé, et pour tout autre point $b \in X$, on a $a \in \overline{\{b\}}$ ($\S 1, n^\circ 1, \text{cor.} 3 \text{ de la prop.} 2$). Pour tout préschéma $(Y, \mathcal{O}_Y)$ et tout point $y \in Y$, le schéma local $\operatorname{Spec}(\mathcal{O}_y)$ est appelé le \emph{préschéma local} de $Y$ au point $y$. Soit $V$ un ouvert affine de $Y$ contenant $y$, $B$ l'anneau de $(V, \mathcal{O}_Y|V)$; $\mathcal{O}_y$ s'identifie canoniquement à $B_y$, et l'homomorphisme canonique $i_V : B \to B_y$ correspond par suite à un morphisme de préschémas $\operatorname{Spec}(\mathcal{O}_y) \to V$. Si on compose ce morphisme avec le monomorphisme canonique $V \to Y$, on a donc un morphisme $\operatorname{Spec}(\mathcal{O}_y) \to Y$, qui est indépendant de l'ouvert affine $V$ (contenant $y$) choisi : on le voit en remarquant que si $V'$ est un second ouvert affine contenant $y$, il existe un troisième ouvert affine $W$ contenant $y$ tel que $W \subset V \cap V'$ (n°1, lemme 1); on peut donc se limiter au cas où $V \subset V'$, et notre assertion est alors évidente, compte tenu du th. 4 du $\S 1, n^\circ 7$. Le morphisme $\operatorname{Spec}(\mathcal{O}_y) \to Y$ ainsi défini est dit \emph{canonique}.

\emph{Proposition 4.} Soit $(Y, \mathcal{O}_Y)$ un préschéma; pour tout $y \in Y$, soit $(\psi, \theta)$ le morphisme canonique $(\operatorname{Spec}(\mathcal{O}_y), \mathcal{O}_y) \to (Y, \mathcal{O}_Y)$. Alors $\psi$ est un homéomorphisme de $\operatorname{Spec}(\mathcal{O}_y)$ sur le sous-espace de $Y$ formé des $z$ tels que $y \in \overline{\{z\}}$; en outre, si $p = \psi^{-1}(z)$, $\theta_z : \mathcal{O}_{Y,z} \to (\mathcal{O}_y)_p$ est un isomorphisme.

En effet, tout $z$ tel que $y \in \overline{\{z\}}$ appartient à tout ouvert affine contenant $y$, et on peut donc se ramener au cas où $(Y, \mathcal{O}_Y)$ est un schéma affine; dans ce cas, la première assertion est un cas particulier du cor. 3 de la

\newpage

%% ===== Page 43 =====

\noindent --- 26 ---

\noindent et de la correspondance biunivoque entre idéaux premiers de $\mathcal{O}_y$ et idéaux premiers de $B = \mathcal{O}_y$ contenus dans $j_y$ (chap.0, §1, n°2).

\noindent Il y a donc correspondance biunivoque canonique entre $\operatorname{Spec}(\mathcal{O}_y)$ et l'ensemble des parties fermées irréductibles contenant $y$.

\noindent \emph{Corollaire}. Pour que $y \in Y$ soit le point générique d'une composante irréductible de $Y$, il faut et il suffit que le seul idéal premier de l'anneau local $\mathcal{O}_y$ soit son idéal maximal (autrement dit, que $\mathcal{O}_y$ soit de dimension zéro).

\noindent \emph{Proposition 5}. Soient $(X, \mathcal{O}_X)$ un schéma local d'anneau $A$, a l'unique point fermé de $X$, $(Y, \mathcal{O}_Y)$ un préschéma. Tout morphisme $u = (\psi, \varphi) : (X, \mathcal{O}_X) \to (Y, \mathcal{O}_Y)$ se factorise de façon unique en $X \to \operatorname{Spec}(\mathcal{O}_{\psi(a)}) \to (Y, \mathcal{O}_Y)$, où la seconde flèche désigne le morphisme canonique, et la première correspond à un homomorphisme de $\mathcal{O}_{\psi(a)}$ dans $A$. Cela établit une correspondance biunivoque canonique entre l'ensemble des morphismes $(X, \mathcal{O}_X) \to (Y, \mathcal{O}_Y)$ et l'ensemble des homomorphismes $\mathcal{O}_y \to A$ ($y \in Y$).

\noindent En effet, pour tout $x \in X$, on a $a \in \overline{\{x\}}$, donc $\psi(a) \in \overline{\{\psi(x)\}}$, ce qui prouve que $\psi(X)$ est contenu dans tout ouvert affine contenant $\psi(a)$. On peut donc se ramener au cas où $(Y, \mathcal{O}_Y)$ est un schéma affine d'anneau $B$, et on a alors $u = (\psi, \varphi)$, où $\varphi \in \operatorname{Hom}(B, A)$ (§1, n°7, th.4); en outre on a $\varphi^{-1}(\mathfrak{m}_a) = j_{\psi(a)}$ et par suite l'image par $\varphi$ de tout élément de $B \setminus j_{\psi(a)}$ est inversible dans l'anneau local $A$; la factorisation de l'énoncé résulte donc de la propriété universelle des anneaux de fractions (chap.0, §1, n°2). Inversement, à tout homomorphisme $\mathcal{O}_y \to A$ correspond un morphisme unique $X \to \operatorname{Spec}(\mathcal{O}_y)$ (§1, n°7, th.4), et en le composant avec le morphisme canonique $\operatorname{Spec}(\mathcal{O}_y) \to Y$, on obtient un morphisme $X \to Y$, ce qui achève de démontrer la proposition.

\noindent Les schémas affines dont l'anneau est un corps $K$ ont un espace de base réduit à un seul point; si on tient compte de tout homomorphisme $\mathcal{O}_y \to K$ correspond biunivoquement à un monomorphisme $k(y) \to K$ du corps des

\newpage

%% ===== Page 44 =====

restes de $\mathcal{O}_y$ dans $K$, on voit que:

Corollaire. Si $(X, \mathcal{O}_X)$ est un schéma local dont l'anneau est un corps $K$, il y a correspondance biunivoque entre l'ensemble des morphismes $(X, \mathcal{O}_X) \to (Y, \mathcal{O}_y)$ et l'ensemble des monomorphismes $k(y) \to K$ ($y \in Y$).

En particulier, si $K, K'$ sont deux corps, il y a correspondance biunivoque entre les monomorphismes $K \to K'$ et les morphismes $\operatorname{Spec}(K') \to \operatorname{Spec}(K)$.

3. \emph{Préschémas au-dessus d'un préschéma.}

Définition 4. Etant donné un préschéma $(S, \mathcal{O}_S)$, on dit que la donnée d'un préschéma $(X, \mathcal{O}_X)$ et d'un morphisme de préschémas $(\psi, \theta): (X, \mathcal{O}_X) \to (S, \mathcal{O}_S)$ définit un préschéma $(X, \mathcal{O}_X)$ au-dessus du préschéma $(S, \mathcal{O}_S)$, ou un $S$-préschéma. Le morphisme $(\psi, \theta)$ est appelé le morphisme structural du $S$-préschéma $(X, \mathcal{O}_X)$. Lorsque $(S, \mathcal{O}_S)$ est un schéma affine d'anneau $A$, on dit aussi que $(X, \mathcal{O}_X)$ muni de $(\psi, \theta)$ est un préschéma au-dessus de l'anneau $A$ (ou un $A$-préschéma).

La prop. 3 prouve que la donnée d'un préschéma au-dessus d'un anneau $A$ équivaut à la donnée d'un espace topologique $X$, muni d'un faisceau $\mathcal{B}$ de $A$-algèbres, tel que l'espace annelé $(X, \mathcal{B})$ soit un préschéma. Un préschéma quelconque peut donc toujours être considéré comme un préschéma au-dessus de l'anneau $A$.

Si $\psi: X \to S$ est le morphisme structural d'un $S$-préschéma $X$, on dit qu'un point $x \in X$ est au-dessus d'un point $s \in S$ si $\psi(x) = s$.

Soient $X, Y$ deux préschémas au-dessus d'un préschéma $S$; un morphisme $X \to Y$ de préschémas au-dessus de $S$ (ou $S$-morphisme) est par définition un morphisme de préschémas $\psi: X \to Y$ tel que le diagramme
$$\begin{tikzcd}
X \arrow[rr, "\psi"] \arrow[dr] & & Y \arrow[dl] \\
& S &
\end{tikzcd}$$
(où les flèches obliques sont les morphismes structuraux) soit commutatif : cela entraîne que pour tout $s \in S$ et tout $x \in X$ au-dessus de $s$, $\psi(x)$ doit

\newpage

%% ===== Page 45 =====

\noindent -- 28 --

\noindent aussi être au-dessus de $s$. Il résulte aussitôt de cette définition que le composé de deux $S$-morphismes est un $S$-morphisme; les $S$-préschémas forment donc une catégorie. Lorsque $S$ est un schéma affine d'anneau $A$, on dira aussi $A$-morphisme au lieu de $S$-morphisme.

\noindent Si $\psi: X \to Y$ est un $S$-morphisme, la restriction de $\psi$ à toute partie ouverte $U$ de $X$ est un $S$-morphisme $U \to Y$. Soient $(U_\alpha)$ un recouvrement ouvert de l'espace de base $X$, et pour chaque $\alpha$, soit $\psi_\alpha: U_\alpha \to Y$ un $S$-morphisme; si, pour tout couple d'indices $(\alpha, \beta)$, les restrictions de $\psi_\alpha$ et $\psi_\beta$ à $U_\alpha \cap U_\beta$ coïncident, il existe un $S$-morphisme $X \to Y$ et un seul dont la restriction à chaque $U_\alpha$ soit $\psi_\alpha$.

\noindent Si $\psi: X \to Y$ est un $S$-morphisme tel que $\psi(X) \subset V$, où $V$ est une partie ouverte de $Y$, alors $\psi$, considéré comme morphisme de $X$ dans $V$ (chap. 0, §2, n$^\circ$4) est encore un $S$-morphisme.

\noindent Soit $S' \to S$ un morphisme de préschémas; pour tout $S'$-préschéma $X$, le morphisme composé $X \to S' \to S$ définit $X$ comme $S$-préschéma; réciproquement, si $S'$ est le préschéma induit par $S$ sur un ouvert de son espace de base, et si $X$ est un $S$-préschéma dont l'image par le morphisme structural est contenue dans $S'$, alors $X$ est un $S'$-préschéma. Dans ce dernier cas, si $Y$ est un second $S$-préschéma dont le morphisme structural applique $Y$ dans $S'$, tout $S$-morphisme de $X$ dans $Y$ est un $S'$-morphisme, et réciproquement.

\newpage

%% ===== Page 46 =====

\marginpar{Archive, Grothendieck - Sept 59}

Additions to chapter I, \S 2 :

p.31 , before prop.7 , add :

Corollaire .— Soient $T$ un schéma affine d'anneau $D$, $\alpha = (\rho, \tilde{\rho})$ (resp. $\beta = (\sigma, \tilde{\sigma})$) un morphisme $T \to X$ (resp. $T \to Y$), où $\rho$ (resp. $\sigma$) est un $A$-homomorphisme de $B$ (resp. $C$) dans $D$ ; alors $\alpha \times_S \beta = (\tau, \tilde{\tau})$, où $\tau$ est l'homomorphisme de $B \otimes_A C$ dans $D$ tel que $\tau(b \otimes c) = \rho(b)\sigma(c)$.

p.35 , before prop.8 , add :

En premier lieu, $(X, Y) \to X \times_S Y$ est un bifoncteur covariant dans la catégorie des $S$-préschémas : il suffit en effet de remarquer que le diagramme

\begin{tikzcd}
X \times Y \arrow[r, "f \times 1"] \arrow[d] & X' \times Y \arrow[r, "f' \times 1"] \arrow[d] & X'' \times Y \arrow[d] \\
X \arrow[r, "f"] & X' \arrow[r, "f'"] & X''
\end{tikzcd}

est commutatif.

p.35 , after prop.8 , add :

Corollaire .— Soient $X, Y$ deux $S$-préschémas, $\varphi : X \to S$, $\psi : Y \to S$ leurs morphismes structuraux. Si on identifie canoniquement $X$ à $X \times_S S$ et $Y$ à $S \times_S Y$, les projections $X \times_S Y \to X$ et $X \times_S Y \to Y$ s'identifient respectivement à $1 \times \psi$ et $\varphi \times 1$.

La vérification est immédiate.

p.36 , at the middle of the page, add :

Il est immédiat que $X \to X^{S'}$ est un foncteur covariant de la catégorie des $S$-préschémas dans celle des $S'$-préschémas, lorsqu'on définit $f^{S'}$ comme égal à $f \times_S 1$ pour tout $S$-morphisme $f : X \to Y$.

\newpage

%% ===== Page 47 =====

- ii -

P.37 , after cor.1 , add :

Corollaire 2 .— Soient $Y$ un $S$-préschéma , $f : X \to Y$ un morphisme qui fait de $X$ un $Y$-préschéma (et par suite aussi un $S$-préschéma) . Le préschéma $X^{S'}$ s'identifie alors au produit $X \times_Y S'$, la projection $X \times_Y S' \to Y^{S'}$ s'identifiant à $f^{S'}$.

Soit $\psi : Y \to S$ le morphisme structural ; on a le diagramme commutatif
$$
\begin{tikzcd}
S' \arrow[d] \arrow[r] & Y^{S'} \arrow[d] \arrow[r, "f^{S'}"] & X^{S'} \arrow[d] \\
S \arrow[r, "\psi"] & Y \arrow[r, "f"] & X
\end{tikzcd}
$$
Or $Y^{S'}$ s'identifie à $S'$, et $X^{S'}$ à $S' \times_Y X$ ; tenant compte de la prop. 9 et du cor. de la prop.8 , on en déduit le corollaire .

P.41 , after prop.15 , add :

Corollaire .— Si le $S$-morphisme $f : X \to Y$ est géométriquement injectif il en est de même de $f^{S'} : X^{S'} \to Y^{S'}$.

\newpage

%% ===== Page 48 =====

- 11 -

P.37, after cor.1, add :

\emph{Corollaire 2.} — Soient $Y$ un $S$-préschéma, $f : X \to Y$ un morphisme qui fait de $X$ un $Y$-préschéma (et par suite aussi un $S$-préschéma). Le préschéma $X^{S'}$ s'identifie alors au produit $X \times_Y S'$, la projection $X \times_Y S' \to Y^{S'}$ s'identifiant à $f^{S'}$.

Soit $\psi : Y \to S$ le morphisme structural ; on a le diagramme commutatif
\[
\begin{tikzcd}
X^{S'} \arrow[r, "f^{S'}"] \arrow[d] & Y^{S'} \arrow[d] \\
X \arrow[r, "f"] & Y \arrow[d, "\psi"] \\
& S
\end{tikzcd}
\]
Or $Y^{S'}$ s'identifie à $S'$, et $X^{S'}$ à $S' \times_S X$ ; tenant compte de la prop. 9 et du cor. de la prop. 8, on en déduit le corollaire.

P.41, after prop. 15, add :

\emph{Corollaire.} — Si le $S$-morphisme $f : X \to Y$ est géométriquement injectif, il en est de même de $f^{S'} : X^{S'} \to Y^{S'}$.

\newpage

%% ===== Page 49 =====

\noindent P. 57, before section 5, add :

Remarque. — Supposons que $f : X \to Y$ soit un monomorphisme. Alors $\Delta_X$ est un isomorphisme de $X$ sur $X \times_Y X$. En effet, soient $p_1, p_2$ les projections de $X \times_Y X$ ; pour tout couple de $Y$-morphismes $u, v$ de $T$ dans $X \times_Y X$, on a $f \circ p_1 \circ u = f \circ p_1 \circ v$, $f \circ p_2 \circ u = f \circ p_2 \circ v$ par définition d'un $Y$-morphisme, d'où $p_1 \circ u = p_1 \circ v$, $p_2 \circ u = p_2 \circ v$ et par suite $u = v$. Appliquons en particulier ceci au $Y$-morphisme $\Delta_X \circ p_1$ de $X \times_Y X$ dans lui-même ; il est nécessairement égal à l'identité, donc en vertu de (1), $p_1$ et $\Delta_X$ sont des isomorphismes réciproques.

\newpage

%% ===== Page 50 =====

- iv -

p.57, before section 5, add :

Remarque . - Supposons que $f : X \to Y$ soit un monomorphisme . Alors $\Delta_X$ est un isomorphisme de $X$ sur $X \times_Y X$ . En effet, soient $p_1, p_2$ les projections de $X \times_Y X$ ; pour tout couple de $Y$-morphismes $u, v$ de $T$ dans $X \times_Y X$, on a $f \circ p_1 \circ u = f \circ p_1 \circ v$ , $f \circ p_2 \circ u = f \circ p_2 \circ v$ par définition d'un $Y$-morphisme, d'où $p_1 \circ u = p_1 \circ v$ , $p_2 \circ u = p_2 \circ v$ et par suite $u = v$ . Appliquons en particulier ceci au $Y$-morphisme $\Delta_X p_1$ de $X \times_Y X$ dans lui-même ; il est nécessairement égal à l'identité, donc en vertu de (1), $p_1$ et $\Delta_X$ sont des isomorphismes réciproques.

\newpage

%% ===== Page 51 =====

\hfill 31 \hfill

tout $A$-homomorphisme $f$ de $B \otimes_A C$ dans une $A$-algèbre $L$, on associe les homomorphismes $f \circ u$ et $f \circ v$, on définit une bijection de l'ensemble $\operatorname{Hom}_A(B \otimes_A C, L)$ sur le produit $\operatorname{Hom}_A(B, L) \times \operatorname{Hom}_A(C, L)$, ce qui résulte immédiatement des définitions et de la relation $b \otimes c = (b \otimes 1)(1 \otimes c)$.

Proposition 7. — Soient $X, Y$ deux $S$-préschémas, $\varphi : X \to S, \psi : Y \to S$ leurs morphismes structuraux. Si $S'$ une partie ouverte de $S$ telle que $\varphi(X) \subset S', \psi(Y) \subset S'$. Tout produit des $S$-préschémas $X, Y$ est aussi un produit des $S'$-préschémas $X, Y$, et réciproquement.

Soit $(Z', p_1, p_2)$ un produit de $X, Y$ considérés comme $S'$-préschémas ; $p_1$ et $p_2$ sont aussi des $S$-morphismes. Soit alors $T$ un $S$-préschéma ; s'il existe un $S$-morphisme de $T$ dans $X, Y$ ou $Z'$, il résulte des définitions que le morphisme structural de $T$ dans $S$ applique $T$ dans $S'$, donc $T$ peut être considéré comme $S'$-préschéma, et tout $S$-morphisme de $T$ dans $X, Y$ ou $Z'$ est un $S'$-morphisme ; on en conclut que $(Z', p_1, p_2)$ est un produit de $X, Y$ considérés comme $S$-préschémas. Inversement, si $(Z, p_1, p_2)$ est un produit de $X, Y$ considérés comme $S$-préschémas, et si $\vartheta$ est le morphisme structural $Z \to S$ : on a nécessairement $\vartheta(Z) = \varphi(p_1(Z)) \subset S'$, donc $Z$ peut être considéré comme $S'$-préschéma, $p_1, p_2$ comme des $S'$-morphismes ; tout $S$-morphisme d'un $S'$-préschéma $T$ dans $X, Y$ ou $Z$ étant aussi un $S$-morphisme, il est clair que $(Z, p_1, p_2)$ est un produit de $X, Y$ considérés comme $S'$-préschémas.

Théorème 1. — Etant donnés deux $S$-préschémas $X, Y$, il existe un produit $X \times_S Y$.

Nous procéderons en plusieurs étapes.

Lemme 3. — Soient $(Z, p, q)$ un produit de $X$ et $Y$, $U, V$ des parties ouvertes de $X, Y$ respectivement. Si on pose $W = p^{-1}(U) \cap q^{-1}(V)$, le triplet formé de $W$ et des restrictions de $p$ et $q$ à $W$ (considérées comme des morphismes $W \to U, W \to V$ respectivement) est un produit de $U$ et $V$.

En effet, si $T$ est un $S$-préschéma, on peut identifier les $S$-morphismes $T \to W$ et les $S$-morphismes $T \to Z$ appliquant $T$ dans $W$. Si alors

\newpage

%% ===== Page 52 =====

- 32 -

$$ g : T \to X , \quad h : T \to Y $$
sont deux $S$-morphismes quelconques, on peut les considérer comme des $S$-morphismes de $T$ dans $X$ et $Y$ respectivement, et par hypothèse il y a donc un $S$-morphisme et un seul $f : T \to Z$ tel que $p \circ f = g$, $h = q \circ f$. Comme $p(f(T)) \subset U$, $q(f(T)) \subset V$, on a $f(T) \subset p^{-1}(U) \cap q^{-1}(V) = W$, d'où notre assertion.

Leme 4. — Soient $Z$ un $S$-préschéma, $p : Z \to X$, $q : Z \to Y$ deux $S$-morphismes, $(U_\alpha)$ un recouvrement ouvert de $X$, $(V_\lambda)$ un recouvrement ouvert de $Y$. On suppose que pour tout couple $(\alpha, \lambda)$, le $S$-préschéma $W_{\alpha\lambda} = p^{-1}(U_\alpha) \cap q^{-1}(V_\lambda)$ et les restrictions de $p$ et $q$ à $W_{\alpha\lambda}$ constituent un produit de $U_\alpha$ et $V_\lambda$. Alors $(Z, p, q)$ est un produit de $X$ et $Y$.

Montrons d'abord que si $f_1, f_2$ sont deux $S$-morphismes $T \to Z$, les relations $p \circ f_1 = p \circ f_2$ et $q \circ f_1 = q \circ f_2$ entraînent $f_1 = f_2$. En effet, $Z$ est réunion des $W_{\alpha\lambda}$, donc les $f_1^{-1}(W_{\alpha\lambda})$ forment un recouvrement ouvert de $T$, et il en est de même des $f_2^{-1}(W_{\alpha\lambda})$. En outre, on a $f_1^{-1}(W_{\alpha\lambda}) = f_1^{-1}(p^{-1}(U_\alpha)) \cap f_1^{-1}(q^{-1}(V_\lambda)) = f_2^{-1}(p^{-1}(U_\alpha)) \cap f_2^{-1}(q^{-1}(V_\lambda)) = f_2^{-1}(W_{\alpha\lambda})$ par hypothèse, et tout revient à voir que les restrictions de $f_1$ et $f_2$ à $f_1^{-1}(W_{\alpha\lambda}) = f_2^{-1}(W_{\alpha\lambda})$ sont identiques pour tout couple d'indices. Mais comme ces restrictions peuvent être considérées comme des $S$-morphismes de $f_1^{-1}(W_{\alpha\lambda})$ dans $W_{\alpha\lambda}$, notre assertion résulte de l'hypothèse et de la déf. 5.

Supposons maintenant donnés deux $S$-morphismes $g : T \to X$, $h : T \to Y$. Posons $T_{\alpha\lambda} = g^{-1}(U_\alpha) \cap h^{-1}(V_\lambda)$ ; les $T_{\alpha\lambda}$ forment un recouvrement ouvert de $T$. Par hypothèse, il existe un $S$-morphisme $f_{\alpha\lambda}$ tel que $p \circ f_{\alpha\lambda}$ et $q \circ f_{\alpha\lambda}$ soient les restrictions respectives de $g$ et $h$ à $T_{\alpha\lambda}$. En outre, montrons que les restrictions de $f_{\alpha\lambda}$ et de $f_{\beta\mu}$ à $T_{\alpha\lambda} \cap T_{\beta\mu}$ coïncident, ce qui achèvera de prouver le lemme. Or, les images de $T_{\alpha\lambda} \cap T_{\beta\mu}$ par $f_{\alpha\lambda}$ et par $f_{\beta\mu}$ sont contenues dans $W_{\alpha\lambda} \cap W_{\beta\mu}$ par définition. Comme $W_{\alpha\lambda} \cap W_{\beta\mu} = p^{-1}(U_\alpha \cap U_\beta) \cap q^{-1}(V_\lambda \cap V_\mu)$, il résulte du lemme 3 que $W_{\alpha\lambda} \cap W_{\beta\mu}$ et les restrictions à ce préschéma de $p$ et $q$ constituent un produit de $U_\alpha \cap U_\beta$ et de $V_\lambda \cap V_\mu$. Comme $p \circ f_{\alpha\lambda}$ et $p \circ f_{\beta\mu}$ co-

\newpage

%% ===== Page 53 =====

\text{Incident dans } T_{\alpha\lambda} \cap T_{\beta\mu} \text{ et qu'il en est de même de } q \circ f_{\alpha\lambda} \text{ et } q \circ f_{\beta\mu} \text{, on} \\
\text{voit que } f_{\alpha\lambda} \text{ et } f_{\beta\mu} \text{ coïncident dans } T_{\alpha\lambda} \cap T_{\beta\mu} \text{; q.f.d.} \\

\text{a) Soient } (U_{\alpha}) \text{ un recouvrement ouvert de } X \text{, } (V_{\lambda}) \text{ un recouvrement} \\
\text{ouvert de } Y \text{, et supposons que pour tout couple d'indices } (\alpha, \lambda) \text{, il} \\
\text{existe un produit de } U_{\alpha} \text{ et } V_{\lambda} \text{; alors il existe un produit de } X \text{ et } Y. \\
\text{Appliquant le lemme 3 aux ouverts } U_{\alpha} \cap U_{\beta} \text{, } V_{\lambda} \cap V_{\mu} \text{, on voit qu'il e-} \\
\text{xiste un produit des S-préschémas induits par } X \text{ et } Y \text{ respectivement} \\
\text{sur ces ouverts ; en outre, si on pose } i=(\alpha, \lambda) \text{, } j=(\beta, \mu) \text{, il y a} \\
\text{un isomorphisme canonique } h_{ij} \text{ (resp. } h_{ji} \text{) de ce produit sur un S-pré-} \\
\text{schéma } W_{ij} \text{ (resp. } W_{ji} \text{) induit par } U_{\alpha} \times V_{\lambda} \text{ (resp. } U_{\beta} \times V_{\mu} \text{) sur un ou-} \\
\text{vert ; } f_{ij}=h_{ij} \circ h_{ji}^{-1} \text{ est donc un isomorphisme de } W_{ji} \text{ sur } W_{ij} \text{ . En ou-} \\
\text{tre, pour un troisième couple } k=(\gamma, \nu) \text{, on a } f_{ik}=f_{ij} \circ f_{jk} \text{ dans} \\
\text{ } W_{ik} \cap W_{kj} \text{, comme il résulte de l'application du lemme 3 aux ouverts} \\
\text{ } U_{\alpha} \cap U_{\beta} \cap U_{\gamma} \text{, et } V_{\lambda} \cap V_{\mu} \cap V_{\nu} \text{, dans } U_{\rho} \text{ et } V_{\sigma} \text{ respectivement . Il y a par} \\
\text{suite un espace annelé } Z \text{, un recouvrement ouvert } \{Z_i\} \text{ de l'espace} \\
\text{de base de cet espace annelé (noté encore } Z \text{ par nos conventions) , et} \\
\text{pour chaque } i \text{ un isomorphisme } \phi_i \text{ de l'espace annelé induit } Z_i \text{ sur} \\
\text{ } U_{\alpha} \times V_{\lambda} \text{, de sorte que, pour tout couple } (i, j) \text{, on ait } f_{ij}=\phi_i \circ \phi_j^{-1} \\
\text{ (FAC, I, § 1, n°4) ; de plus, on a } \phi_i(Z_i \cap Z_j) = W_{ij} \text{ . Si } p_i, q_i, \theta_i \text{ sont} \\
\text{les projections et le morphisme structural de } U_{\alpha} \times V_{\lambda} \text{, on constate} \\
\text{aussitôt que } p_i \circ \phi_i = p_j \circ \phi_j \text{ dans } Z_i \cap Z_j \text{, et de même pour les deux au-} \\
\text{tres morphismes . On peut donc définir des morphismes d'espaces anne-} \\
\text{lés } p : Z \to X \text{ (resp. } q : Z \to Y \text{, } \theta : Z \to S) \text{ par le condition que } p \text{ (resp.} \\
\text{ } q, \theta) \text{ coïncide avec } p_i \circ \phi_i \text{ (resp. } q_i \circ \phi_i, \theta_i \circ \phi_i \text{) dans chacun des } Z_i \text{. Il} \\
\text{est clair que } Z \text{ est un préschéma (par définition des préschémas) , et} \\
\text{les morphismes de préschémas étant caractérisés par des conditions} \\
\text{ponctuelles (n°2) , } p, q \text{ et } \theta \text{ sont de tels morphismes ; } Z \text{, muni de} \\
\text{ } \theta \text{, est alors un S-préschéma . Montrons maintenant que } Z_i = p^{-1}(U_{\alpha}) \cap q^{-1}(V_{\lambda}) \\
\text{est égal à } Z_i \text{ . Pour tout indice } j=(\beta, \mu) \text{, on a } Z_j \cap Z_i = \phi_j^{-1} (\phi_j(Z_j \cap Z_i) \cap \phi_j(Z_i))

\newpage

%% ===== Page 54 =====

- 34 -

\hrulefill

Or, $p_j^{-1}(U_\alpha \cap U_\beta) \cap q_j^{-1}(V_\lambda \cap V_\mu) = p_j^{-1}(U_\alpha \cap U_\beta) \cap q_j^{-1}(V_\lambda \cap V_\mu)$; en vertu du lemme 3, les restrictions de $p_j$ et $q_j$ à $p_j^{-1}(U_\alpha) \cap q_j^{-1}(V_\lambda)$ définissent sur ce $S$-préschéma une structure de produit de $U_\alpha \cap U_\beta$ et $V_\lambda \cap V_\mu$ ; mais l'unicité du produit entraîne alors que $p_j^{-1}(U_\alpha) \cap q_j^{-1}(V_\lambda) = W_{ji}$. On a par suite $Z_j \cap Z_i = Z_j \cap Z_i$ pour tout $j$, d'où $Z_i = Z_j$. On déduit alors du lemme 4 que $(Z, p, q)$ est un produit de $X$ et $Y$.

b) Soient $\varphi : X \to S$, $\psi : Y \to S$ les morphismes structuraux de $X$ et $Y$, $(S_i)$ un recouvrement ouvert de $S$, et posons $X_i = \varphi^{-1}(S_i)$, $Y_i = \psi^{-1}(S_i)$. Si chacun des produits $X_i \times_S Y_i$ existe, alors $X \times_S Y$ existe.

D'après a), tout revient à prouver que les produits $X_i \times_S Y_i$ existent quels que soient $i$ et $j$. Posons $X_{ij} = X_i \cap X_j = \varphi^{-1}(S_i \cap S_j)$, $Y_{ij} = Y_i \cap Y_j = \psi^{-1}(S_i \cap S_j)$; en vertu du lemme 3, le produit $X_{ij} \times_S Y_{ij}$ existe. Notons maintenant que si $T$ est un $S$-préschéma et s'il existe des $S$-morphismes $g : T \to X_i$, $h : T \to Y_i$, on a nécessairement $\varphi(g(T)) = \psi(h(T)) \subset S_i \cap S_j$; entrainant dit $g(T) \subset X_{ij}$ et $h(T) \subset Y_{ij}$; il est alors immédiat que $Z_{ij}$ est un produit de $X_i$ et $Y_i$.

c) Nous pouvons maintenant achever de démontrer le th. 1. Si $S$ est un schéma affine, il y a des recouvrements $(U_\alpha), (V_\lambda)$ respectivement, formés d'ouverts affines ; comme $U_\alpha \times_S V_\lambda$ existe en vertu de la prop. 6, il en est de même de $X \times_S Y$ par a). Si $S$ est un préschéma quelconque, il y a un recouvrement $(S_i)$ de $S$ formé d'ouverts affines, et si $\varphi : X \to S$, $\psi : Y \to S$ sont les morphismes structuraux, et si on pose $X_i = \varphi^{-1}(S_i)$, $Y_i = \psi^{-1}(S_i)$, les produits $X_i \times_S Y_i$ existent d'après ce qui précède ; mais alors les produits $X_i \times_S Y_i$ existent aussi (prop. 7), donc il en est de même de $X \times_S Y$ d'après b).

\underline{Corollaire}

\hrulefill \ \ \ \ \ \ \ \ \ \ \ \ \ \ \ \ \ \ \ \ \ \ \ \ \ \ \ \ \ \ \ \ \ \ \ \ \ \ \ \ \ \ \ \ \ \ \ \ \ \ \ \ \ \ \ \ \ \ \ \ \ \ \ \ \ \ \ \ \ \ \ \ \ \ \ \ \ \ \ \ \ \ \ \ \ \ \ \ \ \ \ \ \ \ \ \ \ \ \ \ \ \ \ \ \ \ \ \ \ \ \ \ \ \ \ \ \ \ \ \ \ \ \ \ \ \ \ \ \ \ \ \ \ \ \ \ \ \ \ \ \ \ \ \ \ \ \ \ \ \ \ \ \ \ \ \ \ \ \ \ \ \ \ \ \ \ \ \ \ \ \ \ \ \ \ \ \ \ \ \ \ \ \ \ \ \ \ \ \ \ \ \ \ \ \ \ \ \ \ \ \ \ \ \ \ \ \ \ \ \ \ \ \ \ \ \ \ \ \ \ \ \ \ \ \ \ \ \ \ \ \ \ \ \ \ \ \ \ \ \ \ \ \ \ \ \ \ \ \ \ \ \ \ \ \ \ \ \ \ \ \ \ \ \ \ \ \ \ \ \ \ \ \ \ \ \ \ \ \ \ \ \ \ \ \ \ \ \ \ \ \ \ \ \ \ \ \ \ \ \ \ \ \ \ \ \ \ \ \ \ \ \ \ \ \ \ \ \ \ \ \ \ \ \ \ \ \ \ \ \ \ \ \ \ \ \ \ \ \ \ \ \ \ \ \ \ \ \ \ \ \ \ \ \ \ \ \ \ \ \ \ \ \ \ \ \ \ \ \ \ \ \ \ \ \ \ \ \ \ \ \ \ \ \ \ \ \ \ \ \ \ \ \ \ \ \ \ \ \ \ \ \ \ \ \ \ \ \ \ \ \ \ \ \ \ \ \ \ \ \ \ \ \ \ \ \ \ \ \ \ \ \ \ \ \ \ \ \ \ \ \ \ \ \ \ \ \ \ \ \ \ \ \ \ \ \ \ \ \ \ \ \ \ \ \ \ \ \ \ \ \ \ \ \ \ \ \ \ \ \ \ \ \ \ \ \ \ \ \ \ \ \ \ \ \ \ \ \ \ \ \ \ \ \ \ \ \ \ \ \ \ \ \ \ \ \ \ \ \ \ \ \ \ \ \ \ \ \ \ \ \ \ \ \ \ \ \ \ \ \ \ \ \ \ \ \ \ \ \ \ \ \ \ \ \ \ \ \ \ \ \ \ \ \ \ \ \ \ \ \ \ \ \ \ \ \ \ \ \ \ \ \ \ \ \ \ \ \ \ \ \ \ \ \ \ \ \ \ \ \ \ \ \ \ \ \ \ \ \ \ \ \ \ \ \ \ \ \ \ \ \ \ \ \ \ \ \ \ \ \ \ \ \ \ \ \ \ \ \ \ \ \ \ \ \ \ \ \ \ \ \ \ \ \ \ \ \ \ \ \ \ \ \ \ \ \ \ \ \ \ \ \ \ \ \ \ \ \ \ \ \ \ \ \ \ \ \ \ \ \ \ \ \ \ \ \ \ \ \ \ \ \ \ \ \ \ \ \ \ \ \ \ \ \ \ \ \ \ \ \ \ \ \ \ \ \ \ \ \ \ \ \ \ \ \ \ \ \ \ \ \ \ \ \ \ \ \ \ \ \ \ \ \ \ \ \ \ \ \ \ \ \ \ \ \ \ \ \ \ \ \ \ \ \ \ \ \ \ \ \ \ \ \ \ \ \ \ \ \ \ \ \ \ \ \ \ \ \ \ \ \ \ \ \ \ \ \ \ \ \ \ \ \ \ \ \ \ \ \ \ \ \ \ \ \ \ \ \ \ \ \ \ \ \ \ \ \ \ \ \ \ \ \ \ \ \ \ \ \ \ \ \ \ \ \ \ \ \ \ \ \ \ \ \ \ \ \ \ \ \ \ \ \ \ \ \ \ \ \ \ \ \ \ \ \ \ \ \ \ \ \ \ \ \ \ \ \ \ \ \ \ \ \ \ \ \ \ \ \ \ \ \ \ \ \ \ \ \ \ \ \ \ \ \ \ \ \ \ \ \ \ \ \ \ \ \ \ \ \ \ \ \ \ \ \ \ \ \ \ \ \ \ \ \ \ \ \ \ \ \ \ \ \ \ \ \ \ \ \ \ \ \ \ \ \ \ \ \ \ \ \ \ \ \ \ \ \ \ \ \ \ \ \ \ \ \ \ \ \ \ \ \ \ \ \ \ \ \ \ \ \ \ \ \ \ \ \ \ \ \ \ \ \ \ \ \ \ \ \ \ \ \ \ \ \ \ \ \ \ \ \ \ \ \ \ \ \ \ \ \ \ \ \ \ \ \ \ \ \ \ \ \ \ \ \ \ \ \ \ \ \ \ \ \ \ \ \ \ \ \ \ \ \ \ \ \ \ \ \ \ \ \ \ \ \ \ \ \ \ \ \ \ \ \ \ \ \ \ \ \ \ \ \ \ \ \ \ \ \ \ \ \ \ \ \ \ \ \ \ \ \ \ \ \ \ \ \ \ \ \ \ \ \ \ \ \ \ \ \ \ \ \ \ \ \ \ \ \ \ \ \ \ \ \ \ \ \ \ \ \ \ \ \ \ \ \ \ \ \ \ \ \ \ \ \ \ \ \ \ \ \ \ \ \ \ \ \ \ \ \ \ \ \ \ \ \ \ \ \ \ \ \ \ \ \ \ \ \ \ \ \ \ \ \ \ \ \ \ \ \ \ \ \ \ \ \ \ \ \ \ \ \ \ \ \ \ \ \ \ \ \ \ \ \ \ \ \ \ \ \ \ \ \ \ \ \ \ \ \ \ \ \ \ \ \ \ \ \ \ \ \ \ \ \ \ \ \ \ \ \ \ \ \ \ \ \ \ \ \ \ \ \ \ \ \ \ \ \ \ \ \ \ \ \ \ \ \ \ \ \ \ \ \ \ \ \ \ \ \ \ \ \ \ \ \ \ \ \ \ \ \ \ \ \ \ \ \ \ \ \ \ \ \ \ \ \ \ \ \ \ \ \ \ \ \ \ \ \ \ \ \ \ \ \ \ \ \ \ \ \ \ \ \ \ \ \ \ \ \ \ \ \ \ \ \ \ \ \ \ \ \ \ \ \ \ \ \ \ \ \ \ \ \ \ \ \ \ \ \ \ \ \ \ \ \ \ \ \ \ \ \ \ \ \ \ \ \ \ \ \ \ \ \ \ \ \ \ \ \ \ \ \ \ \ \ \ \ \ \ \ \ \ \ \ \ \ \ \ \ \ \ \ \ \ \ \ \ \ \ \ \ \ \ \ \ \ \ \ \ \ \ \ \ \ \ \ \ \ \ \ \ \ \ \ \ \ \ \ \ \ \ \ \ \ \ \ \ \ \ \ \ \ \ \ \ \ \ \ \ \ \ \ \ \ \ \ \ \ \ \ \ \ \ \ \ \ \ \ \ \ \ \ \ \ \ \ \ \ \ \ \ \ \ \ \ \ \ \ \ \ \ \ \ \ \ \ \ \ \ \ \ \ \ \ \ \ \ \ \ \ \ \ \ \ \ \ \ \ \ \ \ \ \ \ \ \ \ \ \ \ \ \ \ \ \ \ \ \ \ \ \ \ \ \ \ \ \ \ \ \ \ \ \ \ \ \ \ \ \ \ \ \ \ \ \ \ \ \ \ \ \ \ \ \ \ \ \ \ \ \ \ \ \ \ \ \ \ \ \ \ \ \ \ \ \ \ \ \ \ \ \ \ \ \ \ \ \ \ \ \ \ \ \ \ \ \ \ \ \ \ \ \ \ \ \ \ \ \ \ \ \ \ \ \ \ \ \ \ \ \ \ \ \ \ \ \ \ \ \ \ \ \ \ \ \ \ \ \ \ \ \ \ \ \ \ \ \ \ \ \ \ \ \ \ \ \ \ \ \ \ \ \ \ \ \ \ \ \ \ \ \ \ \ \ \ \ \ \ \ \ \ \ \ \ \ \ \ \ \ \ \ \ \ \ \ \ \ \ \ \ \ \ \ \ \ \ \ \ \ \ \ \ \ \ \ \ \ \ \ \ \ \ \ \ \ \ \ \ \ \ \ \ \ \ \ \ \ \ \ \ \ \ \ \ \ \ \ \ \ \ \ \ \ \ \ \ \ \ \ \ \ \ \ \ \ \ \ \ \ \ \ \ \ \ \ \ \ \ \ \ \ \ \ \ \ \ \ \ \ \ \ \ \ \ \ \ \ \ \ \ \ \ \ \ \ \ \ \ \ \ \ \ \ \ \ \ \ \ \ \ \ \ \ \ \ \ \ \ \ \ \ \ \ \ \ \ \ \ \ \ \ \ \ \ \ \ \ \ \ \ \ \ \ \ \ \ \ \ \ \ \ \ \ \ \ \ \ \ \ \ \ \ \ \ \ \ \ \ \ \ \ \ \ \ \ \ \ \ \ \ \ \ \ \ \ \ \ \ \ \ \ \ \ \ \ \ \ \ \ \ \ \ \ \ \ \ \ \ \ \ \ \ \ \ \ \ \ \ \ \ \ \ \ \ \ \ \ \ \ \ \ \ \ \ \ \ \ \ \ \ \ \ \ \ \ \ \ \ \ \ \ \ \ \ \ \ \ \ \ \ \ \ \ \ \ \ \ \ \ \ \ \ \ \ \ \ \ \ \ \ \ \ \ \ \ \ \ \ \ \ \ \ \ \ \ \ \ \ \ \ \ \ \ \ \ \ \ \ \ \ \ \ \ \ \ \ \ \ \ \ \ \ \ \ \ \ \ \ \ \ \ \ \ \ \ \ \ \ \ \ \ \ \ \ \ \ \ \ \ \ \ \ \ \ \ \ \ \ \ \ \ \ \ \ \ \ \ \ \ \ \ \ \ \ \ \ \ \ \ \ \ \ \ \ \ \ \ \ \ \ \ \ \ \ \ \ \ \ \ \ \ \ \ \ \ \ \ \ \ \ \ \ \ \ \ \ \ \ \ \ \ \ \ \ \ \ \ \ \ \ \ \ \ \ \ \ \ \ \ \ \ \ \ \ \ \ \ \ \ \ \ \ \ \ \ \ \ \ \ \ \ \ \ \ \ \ \ \ \ \ \ \ \ \ \ \ \ \ \ \ \ \ \ \ \ \ \ \ \ \ \ \ \ \ \ \ \ \ \ \ \ \ \ \ \ \ \ \ \ \ \ \ \ \ \ \ \ \ \ \ \ \ \ \ \ \ \ \ \ \ \ \ \ \ \ \ \ \ \ \ \ \ \ \ \ \ \ \ \ \ \ \ \ \ \ \ \ \ \ \ \ \ \ \ \ \ \ \ \ \ \ \ \ \ \ \ \ \ \ \ \ \ \ \ \ \ \ \ \ \ \ \ \ \ \ \ \ \ \ \ \ \ \ \ \ \ \ \ \ \ \ \ \ \ \ \ \ \ \ \ \ \ \ \ \ \ \ \ \ \ \ \ \ \ \ \ \ \ \ \ \ \ \ \ \ \ \ \ \ \ \ \ \ \ \ \ \ \ \ \ \ \ \ \ \ \ \ \ \ \ \ \ \ \ \ \ \ \ \ \ \ \ \ \ \ \ \ \ \ \ \ \ \ \ \ \ \ \ \ \ \ \ \ \ \ \ \ \ \ \ \ \ \ \ \ \ \ \ \ \ \ \ \ \ \ \ \ \ \ \ \ \ \ \ \ \ \ \ \ \ \ \ \ \ \ \ \ \ \ \ \ \ \ \ \ \ \ \ \ \ \ \ \ \ \ \ \ \ \ \ \ \ \ \ \ \ \ \ \ \ \ \ \ \ \ \ \ \ \ \ \ \ \ \ \ \ \ \ \ \ \ \ \ \ \ \ \ \ \ \ \ \ \ \ \ \ \ \ \ \ \ \ \ \ \ \ \ \ \ \ \ \ \ \ \ \ \ \ \ \ \ \ \ \ \ \ \ \ \ \ \ \ \ \ \ \ \ \ \ \ \ \ \ \ \ \ \ \ \ \ \ \ \ \ \ \ \ \ \ \ \ \ \ \ \ \ \ \ \ \ \ \ \ \ \ \ \ \ \ \ \ \ \ \ \ \ \ \ \ \ \ \ \ \ \ \ \ \ \ \ \ \ \ \ \ \ \ \ \ \ \ \ \ \ \ \ \ \ \ \ \ \ \ \ \ \ \ \ \ \ \ \ \ \ \ \ \ \ \ \ \ \ \ \ \ \ \ \ \ \ \ \ \ \ \ \ \ \ \ \ \ \ \ \ \ \ \ \ \ \ \ \ \ \ \ \ \ \ \ \ \ \ \ \ \ \ \ \ \ \ \ \ \ \ \ \ \ \ \ \ \ \ \ \ \ \ \ \ \ \ \ \ \ \ \ \ \ \ \ \ \ \ \ \ \ \ \ \ \ \ \ \ \ \ \ \ \ \ \ \ \ \ \ \ \ \ \ \ \ \ \ \ \ \ \ \ \ \ \ \ \ \ \ \ \ \ \ \ \ \ \ \ \ \ \ \ \ \ \ \ \ \ \ \ \ \ \ \ \ \ \ \ \ \ \ \ \ \ \ \ \ \ \ \ \ \ \ \ \ \ \ \ \ \ \ \ \ \ \ \ \ \ \ \ \ \ \ \ \ \ \ \ \ \ \ \ \ \ \ \ \ \ \ \ \ \ \ \ \ \ \ \ \ \ \ \ \ \ \ \ \ \ \ \ \ \ \ \ \ \ \ \ \ \ \ \ \ \ \ \ \ \ \ \ \ \ \ \ \ \ \ \ \ \ \ \ \ \ \ \ \ \ \ \ \ \ \ \ \ \ \ \ \ \ \ \ \ \ \ \ \ \ \ \ \ \ \ \ \ \ \ \ \ \ \ \ \ \ \ \ \ \ \ \ \ \ \ \ \ \ \ \ \ \ \ \ \ \ \ \ \ \ \ \ \ \ \ \ \ \ \ \ \ \ \ \ \ \ \ \ \ \ \ \ \ \ \ \ \ \ \ \ \ \ \ \ \ \ \ \ \ \ \ \ \ \ \ \ \ \ \ \ \ \ \ \ \ \ \ \ \ \ \ \ \ \ \ \ \ \ \ \ \ \ \ \ \ \ \ \ \ \ \ \ \ \ \ \ \ \ \ \ \ \ \ \ \ \ \ \ \ \ \ \ \ \ \ \ \ \ \ \ \ \ \ \ \ \ \ \ \ \ \ \ \ \ \ \ \ \ \ \ \ \ \ \ \ \ \ \ \ \ \ \ \ \ \ \ \ \ \ \ \ \ \ \ \ \ \ \ \ \ \ \ \ \ \ \ \ \ \ \ \ \ \ \ \ \ \ \ \ \ \ \ \ \ \ \ \ \ \ \ \ \ \ \ \ \ \ \ \ \ \ \ \ \ \ \ \ \ \ \ \ \ \ \ \ \ \ \ \ \ \ \ \ \ \ \ \ \ \ \ \ \ \ \ \ \ \ \ \ \ \ \ \ \ \ \ \ \ \ \ \ \ \ \ \ \ \ \ \ \ \ \ \ \ \ \ \ \ \ \ \ \ \ \ \ \ \ \ \ \ \ \ \ \ \ \ \ \ \ \ \ \ \ \ \ \ \ \ \ \ \ \ \ \ \ \ \ \ \ \ \ \ \ \ \ \ \ \ \ \ \ \ \ \ \ \ \ \ \ \ \ \ \ \ \ \ \ \ \ \ \ \ \ \ \ \ \ \ \ \ \ \ \ \ \ \ \ \ \ \ \ \ \ \ \ \ \ \ \ \ \ \ \ \ \ \ \ \ \ \ \ \ \ \ \ \ \ \ \ \ \ \ \ \ \ \ \ \ \ \ \ \ \ \ \ \ \ \ \ \ \ \ \ \ \ \ \ \ \ \ \ \ \ \ \ \ \ \ \ \ \ \ \ \ \ \ \ \ \ \ \ \ \ \ \ \ \ \ \ \ \ \ \ \ \ \ \ \ \ \ \ \ \ \ \ \ \ \ \ \ \ \ \ \ \ \ \ \ \ \ \ \ \ \ \ \ \ \ \ \ \ \ \ \ \ \ \ \ \ \ \ \ \ \ \ \ \ \ \ \ \ \ \ \ \ \ \ \ \ \ \ \ \ \ \ \ \ \ \ \ \ \ \ \ \ \ \ \ \ \ \ \ \ \ \ \ \ \ \ \ \ \ \ \ \ \ \ \ \ \ \ \ \ \ \ \ \ \ \ \ \ \ \ \ \ \ \ \ \ \ \ \ \ \ \ \ \ \ \ \ \ \ \ \ \ \ \ \ \ \ \ \ \ \ \ \ \ \ \ \ \ \ \ \ \ \ \ \ \ \ \ \ \ \ \ \ \ \ \ \ \ \ \ \ \ \ \ \ \ \ \ \ \ \ \ \ \ \ \ \ \ \ \ \ \ \ \ \ \ \ \ \ \ \ \ \ \ \ \ \ \ \ \ \ \ \ \ \ \ \ \ \ \ \ \ \ \ \ \ \ \ \ \ \ \ \ \ \ \ \ \ \ \ \ \ \ \ \ \ \ \ \ \ \ \ \ \ \ \ \ \ \ \ \ \ \ \ \ \ \ \ \ \ \ \ \ \ \ \ \ \ \ \ \ \ \ \ \ \ \ \ \ \ \ \ \ \ \ \ \ \ \ \ \ \ \ \ \ \ \ \ \ \ \ \ \ \ \ \ \ \ \ \ \ \ \ \ \ \ \ \ \ \ \ \ \ \ \ \ \ \ \ \ \ \ \ \ \ \ \ \ \ \ \ \ \ \ \ \ \ \ \ \ \ \ \ \ \ \ \ \ \ \ \ \ \ \ \ \ \ \ \ \ \ \ \ \ \ \ \ \ \ \ \ \ \ \ \ \ \ \ \ \ \ \ \ \ \ \ \ \ \ \ \ \ \ \ \ \ \ \ \ \ \ \ \ \ \ \ \ \ \ \ \ \ \ \ \ \ \ \ \ \ \ \ \ \ \ \ \ \ \ \ \ \ \ \ \ \ \ \ \ \ \ \ \ \ \ \ \ \ \ \ \ \ \ \ \ \ \ \ \ \ \ \ \ \ \ \ \ \ \ \ \ \ \ \ \ \ \ \ \ \ \ \ \ \ \ \ \ \ \ \ \ \ \ \ \ \ \ \ \ \ \ \ \ \ \ \ \ \ \ \ \ \ \ \ \ \ \ \ \ \ \ \ \ \ \ \ \ \ \ \ \ \ \ \ \ \ \ \ \ \ \ \ \ \ \ \ \ \ \ \ \ \ \ \ \ \ \ \ \ \ \ \ \ \ \ \ \ \ \ \ \ \ \ \ \ \ \ \ \ \ \ \ \ \ \ \ \ \ \ \ \ \ \ \ \ \ \ \ \ \ \ \ \ \ \ \ \ \ \ \ \ \ \ \ \ \ \ \ \ \ \ \ \ \ \ \ \ \ \ \ \ \ \ \ \ \ \ \ \ \ \ \ \ \ \ \ \ \ \ \ \ \ \ \ \ \ \ \ \ \ \ \ \ \ \ \ \ \ \ \ \ \ \ \ \ \ \ \ \ \ \ \ \ \ \ \ \ \ \ \ \ \ \ \ \ \ \ \ \ \ \ \ \ \ \ \ \ \ \ \ \ \ \ \ \ \ \ \ \ \ \ \ \ \ \ \ \ \ \ \ \ \ \ \ \ \ \ \ \ \ \ \ \ \ \ \ \ \ \ \ \ \ \ \ \ \ \ \ \ \ \ \ \ \ \ \ \ \ \ \ \ \ \ \ \ \ \ \ \ \ \ \ \ \ \ \ \ \ \ \ \ \ \ \ \ \ \ \ \ \ \ \ \ \ \ \ \ \ \ \ \ \ \ \ \ \ \ \ \ \ \ \ \ \ \ \ \ \ \ \ \ \ \ \ \ \ \ \ \ \ \ \ \ \ \ \ \ \ \ \ \ \ \ \ \ \ \ \ \ \ \ \ \ \ \ \ \ \ \ \ \ \ \ \ \ \ \ \ \ \ \ \ \ \ \ \ \ \ \ \ \ \ \ \ \ \ \ \ \ \ \ \ \ \ \ \ \ \ \ \ \ \ \ \ \ \ \ \ \ \ \ \ \ \ \ \ \ \ \ \ \ \ \ \ \ \ \ \ \ \ \ \ \ \ \ \ \ \ \ \ \ \ \ \ \ \ \ \ \ \ \ \ \ \ \ \ \ \ \ \ \ \ \ \ \ \ \ \ \ \ \ \ \ \ \ \ \ \ \ \ \ \ \ \ \ \ \ \ \ \ \ \ \ \ \ \ \ \ \ \ \ \ \ \ \ \ \ \ \ \ \ \ \ \ \ \ \ \ \ \ \ \ \ \ \ \ \ \ \ \ \ \ \ \ \ \ \ \ \ \ \ \ \ \ \ \ \ \ \ \ \ \ \ \ \ \ \ \ \ \ \ \ \ \ \ \ \ \ \ \ \ \ \ \ \ \ \ \ \ \ \ \ \ \ \ \ \ \ \ \ \ \ \ \ \ \ \ \ \ \ \ \ \ \ \ \ \ \ \ \ \ \ \ \ \ \ \ \ \ \ \ \ \ \ \ \ \ \ \ \ \ \ \ \ \ \ \ \ \ \ \ \ \ \ \ \ \ \ \ \ \ \ \ \ \ \ \ \ \ \ \ \ \ \ \ \ \ \ \ \ \ \ \ \ \ \ \ \ \ \ \ \ \ \ \ \ \ \ \ \ \ \ \ \ \ \ \ \ \ \ \ \ \ \ \ \ \ \ \ \ \ \ \ \ \ \ \ \ \ \ \ \ \ \ \ \ \ \ \ \ \ \ \ \ \ \ \ \ \ \ \ \ \ \ \ \ \ \ \ \ \ \ \ \ \ \ \ \ \ \ \ \ \ \ \ \ \ \ \ \ \ \ \ \ \ \ \ \ \ \ \ \ \ \ \ \ \ \ \ \ \ \ \ \ \ \ \ \ \ \ \ \ \ \ \ \ \ \ \ \ \ \ \ \ \ \ \ \ \ \ \ \ \ \ \ \ \ \ \ \ \ \ \ \ \ \ \ \ \ \ \ \ \ \ \ \ \ \ \ \ \ \ \ \ \ \ \ \ \ \ \ \ \ \ \ \ \ \ \ \ \ \ \ \ \ \ \ \ \ \ \ \ \ \ \ \ \ \ \ \ \ \ \ \ \ \ \ \ \ \ \ \ \ \ \ \ \ \ \ \ \ \ \ \ \ \ \ \ \ \ \ \ \ \ \ \ \ \ \ \ \ \ \ \ \ \ \ \ \ \ \ \ \ \ \ \ \ \ \ \ \ \ \ \ \ \ \ \ \ \ \ \ \ \ \ \ \ \ \ \ \ \ \ \ \ \ \ \ \ \ \ \ \ \ \ \ \ \ \ \ \ \ \ \ \ \ \ \ \ \ \ \ \ \ \ \ \ \ \ \ \ \ \ \ \ \ \ \ \ \ \ \ \ \ \ \ \ \ \ \ \ \ \ \ \ \ \ \ \ \ \ \ \ \ \ \ \ \ \ \ \ \ \ \ \ \ \ \ \ \ \ \ \ \ \ \ \ \ \ \ \ \ \ \ \ \ \ \ \ \ \ \ \ \ \ \ \ \ \ \ \ \ \ \ \ \ \ \ \ \ \ \ \ \ \ \ \ \ \ \ \ \ \ \ \ \ \ \ \ \ \ \ \ \ \ \ \ \ \ \ \ \ \ \ \ \ \ \ \ \ \ \ \ \ \ \ \ \ \ \ \ \ \ \ \ \ \ \ \ \ \ \ \ \ \ \ \ \ \ \ \ \ \ \ \ \ \ \ \ \ \ \ \ \ \ \ \ \ \ \ \ \ \ \ \ \ \ \ \ \ \ \ \ \ \ \ \ \ \ \ \ \ \ \ \ \ \ \ \ \ \ \ \ \ \ \ \ \ \ \ \ \ \ \ \ \ \ \ \ \ \ \ \ \ \ \ \ \ \ \ \ \ \ \ \ \ \ \ \ \ \ \ \ \ \ \ \ \ \ \ \ \ \ \ \ \ \ \ \ \ \ \ \ \ \ \ \ \ \ \ \ \ \ \ \ \ \ \ \ \ \ \ \ \ \ \ \ \ \ \ \ \ \ \ \ \ \ \ \ \ \ \ \ \ \ \ \ \ \ \ \ \ \ \ \ \ \ \ \ \ \ \ \ \ \ \ \ \ \ \ \ \ \ \ \ \ \ \ \ \ \ \ \ \ \ \ \ \ \ \ \ \ \ \ \ \ \ \ \ \ \ \ \ \ \ \ \ \ \ \ \ \ \ \ \ \ \ \ \ \ \ \ \ \ \ \ \ \ \ \ \ \ \ \ \ \ \ \ \ \ \ \ \ \ \ \ \ \ \ \ \ \ \ \ \ \ \ \ \ \ \ \ \ \ \ \ \ \ \ \ \ \ \ \ \ \ \ \ \ \ \ \ \ \ \ \ \ \ \ \ \ \ \ \ \ \ \ \ \ \ \ \ \ \ \ \ \ \ \ \ \ \ \ \ \ \ \ \ \ \ \ \ \ \ \ \ \ \ \ \ \ \ \ \ \ \ \ \ \ \ \ \ \ \ \ \ \ \ \ \ \ \ \ \ \ \ \ \ \ \ \ \ \ \ \ \ \ \ \ \ \ \ \ \ \ \ \ \ \ \ \ \ \ \ \ \ \ \ \ \ \ \ \ \ \ \ \ \ \ \ \ \ \ \ \ \ \ \ \ \ \ \ \ \ \ \ \ \ \ \ \ \ \ \ \ \ \ \ \ \ \ \ \ \ \ \ \ \ \ \ \ \ \ \ \ \ \ \ \ \ \ \ \ \ \ \ \ \ \ \ \ \ \ \ \ \ \ \ \ \ \ \ \ \ \ \ \ \ \ \ \ \ \ \ \ \ \ \ \ \ \ \ \ \ \ \ \ \ \ \ \ \ \ \ \ \ \ \ \ \ \ \ \ \ \ \ \ \ \ \ \ \ \ \ \ \ \ \ \ \ \ \ \ \ \ \ \ \ \ \ \ \ \ \ \ \ \ \ \ \ \ \ \ \ \ \ \ \ \ \ \ \ \ \ \ \ \ \ \ \ \ \ \ \ \ \ \ \ \ \ \ \ \ \ \ \ \ \ \ \ \ \ \ \ \ \ \ \ \ \ \ \ \ \ \ \ \ \ \ \ \ \ \ \ \ \ \ \ \ \ \ \ \ \ \ \ \ \ \ \ \ \ \ \ \ \ \ \ \ \ \ \ \ \ \ \ \ \ \ \ \ \ \ \ \ \ \ \ \ \ \ \ \ \ \ \ \ \ \ \ \ \ \ \ \ \ \ \ \ \ \ \ \ \ \ \ \ \ \ \ \ \ \ \ \ \ \ \ \ \ \ \ \ \ \ \ \ \ \ \ \ \ \ \ \ \ \ \ \ \ \ \ \ \ \ \ \ \ \ \ \ \ \ \ \ \ \ \ \ \ \ \ \ \ \ \ \ \ \ \ \ \ \ \ \ \ \ \ \ \ \ \ \ \ \ \ \ \ \ \ \ \ \ \ \ \ \ \ \ \ \ \ \ \ \ \ \ \ \ \ \ \ \ \ \ \ \ \ \ \ \ \ \ \ \ \ \ \ \ \ \ \ \ \ \ \ \ \ \ \ \ \ \ \ \ \ \ \ \ \ \ \ \ \ \ \ \ \ \ \ \ \ \ \ \ \ \ \ \ \ \ \ \ \ \ \ \ \ \ \ \ \ \ \ \ \ \ \ \ \ \ \ \ \ \ \ \ \ \ \ \ \ \ \ \ \ \ \ \ \ \ \ \ \ \ \ \ \ \ \ \ \ \ \ \ \ \ \ \ \ \ \ \ \ \ \ \ \ \ \ \ \ \ \ \ \ \ \ \ \ \ \ \ \ \ \ \ \ \ \ \ \ \ \ \ \ \ \ \ \ \ \ \ \ \ \ \ \ \ \ \ \ \ \ \ \ \ \ \ \ \ \ \ \ \ \ \ \ \ \ \ \ \ \ \ \ \ \ \ \ \ \ \ \ \ \ \ \ \ \ \ \ \ \ \ \ \ \ \ \ \ \ \ \ \ \ \ \ \ \ \ \ \ \ \ \ \ \ \ \ \ \ \ \ \ \ \ \ \ \ \ \ \ \ \ \ \ \ \ \ \ \ \ \ \ \ \ \ \ \ \ \ \ \ \ \ \ \ \ \ \ \ \ \ \ \ \ \ \ \ \ \ \ \ \ \ \ \ \ \ \ \ \ \ \ \ \ \ \ \ \ \ \ \ \ \ \ \ \ \ \ \ \ \ \ \ \ \ \ \ \ \ \ \ \ \ \ \ \ \ \ \ \ \ \ \ \ \ \ \ \ \ \ \ \ \ \ \ \ \ \ \ \ \ \ \ \ \ \ \ \ \ \ \ \ \ \ \ \ \ \ \ \ \ \ \ \ \ \ \ \ \ \ \ \ \ \ \ \ \ \ \ \ \ \ \ \ \ \ \ \ \ \ \ \ \ \ \ \ \ \ \ \ \ \ \ \ \ \ \ \ \ \ \ \ \ \ \ \ \ \ \ \ \ \ \ \ \ \ \ \ \ \ \ \ \ \ \ \ \ \ \ \ \ \ \ \ \ \ \ \ \ \ \ \ \ \ \ \ \ \ \ \ \ \ \ \ \ \ \ \ \ \ \ \ \ \ \ \ \ \ \ \ \ \ \ \ \ \ \ \ \ \ \ \ \ \ \ \ \ \ \ \ \ \ \ \ \ \ \ \ \ \ \ \ \ \ \ \ \ \ \ \ \ \ \ \ \ \ \ \ \ \ \ \ \ \ \ \ \ \ \ \ \ \ \ \ \ \ \ \ \ \ \ \ \ \ \ \ \ \ \ \ \ \ \ \ \ \ \ \ \ \ \ \ \ \ \ \ \ \ \ \ \ \ \ \ \ \ \ \ \ \ \ \ \ \ \ \ \ \ \ \ \ \ \ \ \ \ \ \ \ \ \ \ \ \ \ \ \ \ \ \ \ \ \ \ \ \ \ \ \ \ \ \ \ \ \ \ \ \ \ \ \ \ \ \ \ \ \ \ \ \ \ \ \ \ \ \ \ \ \ \ \ \ \ \ \ \ \ \ \ \ \ \ \ \ \ \ \ \ \ \ \ \ \ \ \ \ \ \ \ \ \ \ \ \ \ \ \ \ \ \ \ \ \ \ \ \ \ \ \ \ \ \ \ \ \ \ \ \ \ \ \ \ \ \ \ \ \ \ \ \ \ \ \ \ \ \ \ \ \ \ \ \ \ \ \ \ \ \ \ \ \ \ \ \ \ \ \ \ \ \ \ \ \ \ \ \ \ \ \ \ \ \ \ \ \ \ \ \ \ \ \ \ \ \ \ \ \ \ \ \ \ \ \ \ \ \ \ \ \ \ \ \ \ \ \ \ \ \ \ \ \ \ \ \ \ \ \ \ \ \ \ \ \ \ \ \ \ \ \ \ \ \ \ \ \ \ \ \ \ \ \ \ \ \ \ \ \ \ \ \ \ \ \ \ \ \ \ \ \ \ \ \ \ \ \ \ \ \ \ \ \ \ \ \ \ \ \ \ \ \ \ \ \ \ \ \ \ \ \ \ \ \ \ \ \ \ \ \ \ \ \ \ \ \ \ \ \ \ \ \ \ \ \ \ \ \ \ \ \ \ \ \ \ \ \ \ \ \ \ \ \ \ \ \ \ \ \ \ \ \ \ \ \ \ \ \ \ \ \ \ \ \ \ \ \ \ \ \ \ \ \ \ \ \ \ \ \ \ \ \ \ \ \ \ \ \ \ \ \ \ \ \ \ \ \ \ \ \ \ \ \ \ \ \ \ \ \ \ \ \ \ \ \ \ \ \ \ \ \ \ \ \ \ \ \ \ \ \ \ \ \ \ \ \ \ \ \ \ \ \ \ \ \ \ \ \ \ \ \ \ \ \ \ \ \ \ \ \ \ \ \ \ \ \ \ \ \ \ \ \ \ \ \ \ \ \ \ \ \ \ \ \ \ \ \ \ \ \ \ \ \ \ \ \ \ \ \ \ \ \ \ \ \ \ \ \ \ \ \ \ \ \ \ \ \ \ \ \ \ \ \ \ \ \ \ \ \ \ \ \ \ \ \ \ \ \ \ \ \ \ \ \ \ \ \ \ \ \ \ \ \ \ \ \ \ \ \ \ \ \ \ \ \ \ \ \ \ \ \ \ \ \ \ \ \ \ \ \ \ \ \ \ \ \ \ \ \ \ \ \ \ \ \ \ \ \ \ \ \ \ \ \ \ \ \ \ \ \ \ \ \ \ \ \ \ \ \ \ \ \ \ \ \ \ \ \ \ \ \ \ \ \ \ \ \ \ \ \ \ \ \ \ \ \ \ \ \ \ \ \ \ \ \ \ \ \ \ \ \ \ \ \ \ \ \ \ \ \ \ \ \ \ \ \ \ \ \ \ \ \ \ \ \ \ \ \ \ \ \ \ \ \ \ \ \ \ \ \ \ \ \ \ \ \ \ \ \ \ \ \ \ \ \ \ \ \ \ \ \ \ \ \ \ \ \ \ \ \ \ \ \ \ \ \ \ \ \ \ \ \ \ \ \ \ \ \ \ \ \ \ \ \ \ \ \ \ \ \ \ \ \ \ \ \ \ \ \ \ \ \ \ \ \ \ \ \ \ \ \ \ \ \ \ \ \ \ \ \ \ \ \ \ \ \ \ \ \ \ \ \ \ \ \ \ \ \ \ \ \ \ \ \ \ \ \ \ \ \ \ \ \ \ \ \ \ \ \ \ \ \ \ \ \ \ \ \ \ \ \ \ \ \ \ \ \ \ \ \ \ \ \ \ \ \ \ \ \ \ \ \ \ \ \ \ \ \ \ \ \ \ \ \ \ \ \ \ \ \ \ \ \ \ \ \ \ \ \ \ \ \ \ \ \ \ \ \ \ \ \ \ \ \ \ \ \ \ \ \ \ \ \ \ \ \ \ \ \ \ \ \ \ \ \ \ \ \ \ \ \ \ \ \ \ \ \ \ \ \ \ \ \ \ \ \ \ \ \ \ \ \ \ \ \ \ \ \ \ \ \ \ \ \ \ \ \ \ \ \ \ \ \ \ \ \ \ \ \ \ \ \ \ \ \ \ \ \ \ \ \ \ \ \ \ \ \ \ \ \ \ \ \ \ \ \ \ \ \ \ \ \ \ \ \ \ \ \ \ \ \ \ \ \ \ \ \ \ \ \ \ \ \ \ \ \ \ \ \ \ \ \ \ \ \ \ \ \ \ \ \ \ \ \ \ \ \ \ \ \ \ \ \ \ \ \ \ \ \ \ \ \ \ \ \ \ \ \ \ \ \ \ \ \ \ \ \ \ \ \ \ \ \ \ \ \ \ \ \ \ \ \ \ \ \ \ \ \ \ \ \ \ \ \ \ \ \ \ \ \ \ \ \ \ \ \ \ \ \ \ \ \ \ \ \ \ \ \ \ \ \ \ \ \ \ \ \ \ \ \ \ \ \ \ \ \ \ \ \ \ \ \ \ \ \ \ \ \ \ \ \ \ \ \ \ \ \ \ \ \ \ \ \ \ \ \ \ \ \ \ \ \ \ \ \ \ \ \ \ \ \ \ \ \ \ \ \ \ \ \ \ \ \ \ \ \ \ \ \ \ \ \ \ \ \ \ \ \ \ \ \ \ \ \ \ \ \ \ \ \ \ \ \ \ \ \ \ \ \ \ \ \ \ \ \ \ \ \ \ \ \ \ \ \ \ \ \ \ \ \ \ \ \ \ \ \ \ \ \ \ \ \ \ \ \ \ \ \ \ \ \ \ \ \ \ \ \ \ \ \ \ \ \ \ \ \ \ \ \ \ \ \ \ \ \ \ \ \ \ \ \ \ \ \ \ \ \ \ \ \ \ \ \ \ \ \ \ \ \ \ \ \ \ \ \ \ \ \ \ \ \ \ \ \ \ \ \ \ \ \ \ \ \ \ \ \ \ \ \ \ \ \ \ \ \ \ \ \ \ \ \ \ \ \ \ \ \ \ \ \ \ \ \ \ \ \ \ \ \ \ \ \ \ \ \ \ \ \ \ \ \ \ \ \ \ \ \ \ \ \ \ \ \ \ \ \ \ \ \ \ \ \ \ \ \ \ \ \ \ \ \ \ \ \ \ \ \ \ \ \ \ \ \ \ \ \ \ \ \ \ \ \ \ \ \ \ \ \ \ \ \ \ \ \ \ \ \ \ \ \ \ \ \ \ \ \ \ \ \ \ \ \ \ \ \ \ \ \ \ \ \ \ \ \ \ \ \ \ \ \ \ \ \ \ \ \ \ \ \ \ \ \ \ \ \ \ \ \ \ \ \ \ \ \ \ \ \ \ \ \ \ \ \ \ \ \ \ \ \ \ \ \ \ \ \ \ \ \ \ \ \ \ \ \ \ \ \ \ \ \ \ \ \ \ \ \ \ \ \ \ \ \ \ \ \ \ \ \ \ \ \ \ \ \ \ \ \ \ \ \ \ \ \ \ \ \ \ \ \ \ \ \ \ \ \ \ \ \ \ \ \ \ \ \ \ \ \ \ \ \ \ \ \ \ \ \ \ \ \ \ \ \ \ \ \ \ \ \ \ \ \ \ \ \ \ \ \ \ \ \ \ \ \ \ \ \ \ \ \ \ \ \ \ \ \ \ \ \ \ \ \ \ \ \ \ \ \ \ \ \ \ \ \ \ \ \ \ \ \ \ \ \ \ \ \ \ \ \ \ \ \ \ \ \ \ \ \ \ \ \ \ \ \ \ \ \ \ \ \ \ \ \ \ \ \ \ \ \ \ \ \ \ \ \ \ \ \ \ \ \ \ \ \ \ \ \ \ \ \ \ \ \ \ \ \ \ \ \ \ \ \ \ \ \ \ \ \ \ \ \ \ \ \ \ \ \ \ \ \ \ \ \ \ \ \ \ \ \ \ \ \ \ \ \ \ \ \ \ \ \ \ \ \ \ \ \ \ \ \ \ \ \ \ \ \ \ \ \ \ \ \ \ \ \ \ \ \ \ \ \ \ \ \ \ \ \ \ \ \ \ \ \ \ \ \ \ \ \ \ \ \ \ \ \ \ \ \ \ \ \ \ \ \ \ \ \ \ \ \ \ \ \ \ \ \ \ \ \ \ \ \ \ \ \ \ \ \ \ \ \ \ \ \ \ \ \ \ \ \ \ \ \ \ \ \ \ \ \ \ \ \ \ \ \ \ \ \ \ \ \ \ \ \ \ \ \ \ \ \ \ \ \ \ \ \ \ \ \ \ \ \ \ \ \ \ \ \ \ \ \ \ \ \ \ \ \ \ \ \ \ \ \ \ \ \ \ \ \ \ \ \ \ \ \ \ \ \ \ \ \ \ \ \ \ \ \ \ \ \ \ \ \ \ \ \ \ \ \ \ \ \ \ \ \ \ \ \ \ \ \ \ \ \ \ \ \ \ \ \ \ \ \ \ \ \ \ \ \ \ \ \ \ \ \ \ \ \ \ \ \ \ \ \ \ \ \ \ \ \ \ \ \ \ \ \ \ \ \ \ \ \ \ \ \ \ \ \ \ \ \ \ \ \ \ \ \ \ \ \ \ \ \ \ \ \ \ \ \ \ \ \ \ \ \ \ \ \ \ \ \ \ \ \ \ \ \ \ \ \ \ \ \ \ \ \ \ \ \ \ \ \ \ \ \ \ \ \ \ \ \ \ \ \ \ \ \ \ \ \ \ \ \ \ \ \ \ \ \ \ \ \ \ \ \ \ \ \ \ \ \ \ \ \ \ \ \ \ \ \ \ \ \ \ \ \ \ \ \ \ \ \ \ \ \ \ \ \ \ \ \ \ \ \ \ \ \ \ \ \ \ \ \ \ \ \ \ \ \ \ \ \ \ \ \ \ \ \ \ \ \ \ \ \ \ \ \ \ \ \ \ \ \ \ \ \ \ \ \ \ \ \ \ \ \ \ \ \ \ \ \ \ \ \ \ \ \ \ \ \ \ \ \ \ \ \ \ \ \ \ \ \ \ \ \ \ \ \ \ \ \ \ \ \ \ \ \ \ \ \ \ \ \ \ \ \ \ \ \ \ \ \ \ \ \ \ \ \ \ \ \ \ \ \ \ \ \ \ \ \ \ \ \ \ \ \ \ \ \ \ \ \ \ \ \ \ \ \ \ \ \ \ \ \ \ \ \ \ \ \ \ \ \ \ \ \ \ \ \ \ \ \ \ \ \ \ \ \ \ \ \ \ \ \ \ \ \ \ \ \ \ \ \ \ \ \ \ \ \ \ \ \ \ \ \ \ \ \ \ \ \ \ \ \ \ \ \ \ \ \ \ \ \ \ \ \ \ \ \ \ \ \ \ \ \ \ \ \ \ \ \ \ \ \ \ \ \ \ \ \ \ \ \ \ \ \ \ \ \ \ \ \ \ \ \ \ \ \ \ \ \ \ \ \ \ \ \ \ \ \ \ \ \ \ \ \ \ \ \ \ \ \ \ \ \ \ \ \ \ \ \ \ \ \ \ \ \ \ \ \ \ \ \ \ \ \ \ \ \ \ \ \ \ \ \ \ \ \ \ \ \ \ \ \ \ \ \ \ \ \ \ \ \ \ \ \ \ \ \ \ \ \ \ \ \ \ \ \ \ \ \ \ \ \ \ \ \ \ \ \ \ \ \ \ \ \ \ \ \ \ \ \ \ \ \ \ \ \ \ \ \ \ \ \ \ \ \ \ \ \ \ \ \ \ \ \ \ \ \ \ \ \ \ \ \ \ \ \ \ \ \ \ \ \ \ \ \ \ \ \ \ \ \ \ \ \ \ \ \ \ \ \ \ \ \ \ \ \ \ \ \ \ \ \ \ \ \ \ \ \ \ \ \ \ \ \ \ \ \ \ \ \ \ \ \ \ \ \ \ \ \ \ \ \ \ \ \ \ \ \ \ \ \ \ \ \ \ \ \ \ \ \ \ \ \ \ \ \ \ \ \ \ \ \ \ \ \ \ \ \ \ \ \ \ \ \ \ \ \ \ \ \ \ \ \ \ \ \ \ \ \ \ \ \ \ \ \ \ \ \ \ \ \ \ \ \ \ \ \ \ \ \ \ \ \ \ \ \ \ \ \ \ \ \ \ \ \ \ \ \ \ \ \ \ \ \ \ \ \ \ \ \ \ \ \ \ \ \ \ \ \ \ \ \ \ \ \ \ \ \ \ \ \ \ \ \ \ \ \ \ \ \ \ \ \ \ \ \ \ \ \ \ \ \ \ \ \ \ \ \ \ \ \ \ \ \ \ \ \ \ \ \ \ \ \ \ \ \ \ \ \ \ \ \ \ \ \ \ \ \ \ \ \ \ \ \ \ \ \ \ \ \ \ \ \ \ \ \ \ \ \ \ \ \ \ \ \ \ \ \ \ \ \ \ \ \ \ \ \ \ \ \ \ \ \ \ \ \ \ \ \ \ \ \ \ \ \ \ \ \ \ \ \ \ \ \ \ \ \ \ \ \ \ \ \ \ \ \ \ \ \ \ \ \ \ \ \ \ \ \ \ \ \ \ \ \ \ \ \ \ \ \ \ \ \ \ \ \ \ \ \ \ \ \ \ \ \ \ \ \ \ \ \ \ \ \ \ \ \ \ \ \ \ \ \ \ \ \ \ \ \ \ \ \ \ \ \ \ \ \ \ \ \ \ \ \ \ \ \ \ \ \ \ \ \ \ \ \ \ \ \ \ \ \ \ \ \ \ \ \ \ \ \ \ \ \ \ \ \ \ \ \ \ \ \ \ \ \ \ \ \ \ \ \ \ \ \ \ \ \ \ \ \ \ \ \ \ \ \ \ \ \ \ \ \ \ \ \ \ \ \ \ \ \ \ \ \ \ \ \ \ \ \ \ \ \ \ \ \ \ \ \ \ \ \ \ \ \ \ \ \ \ \ \ \ \ \ \ \ \ \ \ \ \ \ \ \ \ \ \ \ \ \ \ \ \ \ \ \ \ \ \ \ \ \ \ \ \ \ \ \ \ \ \ \ \ \ \ \ \ \ \ \ \ \ \ \ \ \ \ \ \ \ \ \ \ \ \ \ \ \ \ \ \ \ \ \ \ \ \ \ \ \ \ \ \ \ \ \ \ \ \ \ \ \ \ \ \ \ \ \ \ \ \ \ \ \ \ \ \ \ \ \ \ \ \ \ \ \ \ \ \ \ \ \ \ \ \ \ \ \ \ \ \ \ \ \ \ \ \ \ \ \ \ \ \ \ \ \ \ \ \ \ \ \ \ \ \ \ \ \ \ \ \ \ \ \ \ \ \ \ \ \ \ \ \ \ \ \ \ \ \ \ \ \ \ \ \ \ \ \ \ \ \ \ \ \ \ \ \ \ \ \ \ \ \ \ \ \ \ \ \ \ \ \ \ \ \ \ \ \ \ \ \ \ \ \ \ \ \ \ \ \ \ \ \ \ \ \ \ \ \ \ \ \ \ \ \ \ \ \ \ \ \ \ \ \ \ \ \ \ \ \ \ \ \ \ \ \ \ \ \ \ \ \ \ \ \ \ \ \ \ \ \ \ \ \ \ \ \ \ \ \ \ \ \ \ \ \ \ \ \ \ \ \ \ \ \ \ \ \ \ \ \ \ \ \ \ \ \ \ \ \ \ \ \ \ \ \ \ \ \ \ \ \ \ \ \ \ \ \ \ \ \ \ \ \ \ \ \ \ \ \ \ \ \ \ \ \ \ \ \ \ \ \ \ \ \ \ \ \ \ \ \ \ \ \ \ \ \ \ \ \ \ \ \ \ \ \ \ \ \ \ \ \ \ \ \ \ \ \ \ \ \ \ \ \ \ \ \ \ \ \ \ \ \ \ \ \ \ \ \ \ \ \ \ \ \ \ \ \ \ \ \ \ \ \ \ \ \ \ \ \ \ \ \ \ \ \ \ \ \ \ \ \ \ \ \ \ \ \ \ \ \ \ \ \ \ \ \ \ \ \ \ \ \ \ \ \ \ \ \ \ \ \ \ \ \ \ \ \ \ \ \ \ \ \ \ \ \ \ \ \ \ \ \ \ \ \ \ \ \ \ \ \ \ \ \ \ \ \ \ \ \ \ \ \ \ \ \ \ \ \ \ \ \ \ \ \ \ \ \ \ \ \ \ \ \ \ \ \ \ \ \ \ \ \ \ \ \ \ \ \ \ \ \ \ \ \ \ \ \ \ \ \ \ \ \ \ \ \ \ \ \ \ \ \ \ \ \ \ \ \ \ \ \ \ \ \ \ \ \ \ \ \ \ \ \ \ \ \ \ \ \ \ \ \ \ \ \ \ \ \ \ \ \ \ \ \ \ \ \ \ \ \ \ \ \ \ \ \ \ \ \ \ \ \ \ \ \ \ \ \ \ \ \ \ \ \ \ \ \ \ \ \ \ \ \ \ \ \ \ \ \ \ \ \ \ \ \ \ \ \ \ \ \ \ \ \ \ \ \ \ \ \ \ \ \ \ \ \ \ \ \ \ \ \ \ \ \ \ \ \ \ \ \ \ \ \ \ \ \ \ \ \ \ \ \ \ \ \ \ \ \ \ \ \ \ \ \ \ \ \ \ \ \ \ \ \ \ \ \ \ \ \ \ \ \ \ \ \ \ \ \ \ \ \ \ \ \ \ \ \ \ \ \ \ \ \ \ \ \ \ \ \ \ \ \ \ \ \ \ \ \ \ \ \ \ \ \ \ \ \ \ \ \ \ \ \ \ \ \ \ \ \ \ \ \ \ \ \ \ \ \ \ \ \ \ \ \ \ \ \ \ \ \ \ \ \ \ \ \ \ \ \ \ \ \ \ \ \ \ \ \ \ \ \ \ \ \ \ \ \ \ \ \ \ \ \ \ \ \ \ \ \ \ \ \ \ \ \ \ \ \ \ \ \ \ \ \ \ \ \ \ \ \ \ \ \ \ \ \ \ \ \ \ \ \ \ \ \ \ \ \ \ \ \ \ \ \ \ \ \ \ \ \ \ \ \ \ \ \ \ \ \ \ \ \ \ \ \ \ \ \ \ \ \ \ \ \ \ \ \ \ \ \ \ \ \ \ \ \ \ \ \ \ \ \ \ \ \ \ \ \ \ \ \ \ \ \ \ \ \ \ \ \ \ \ \ \ \ \ \ \ \ \ \ \ \ \ \ \ \ \ \ \ \ \ \ \ \ \ \ \ \ \ \ \ \ \ \ \ \ \ \ \ \ \ \ \ \ \ \ \ \ \ \ \ \ \ \ \ \ \ \ \ \ \ \ \ \ \ \ \ \ \ \ \ \ \ \ \ \ \ \ \ \ \ \ \ \ \ \ \ \ \ \ \ \ \ \ \ \ \ \ \ \ \ \ \ \ \ \ \ \ \ \ \ \ \ \ \ \ \ \ \ \ \ \ \ \ \ \ \ \ \ \ \ \ \ \ \ \ \ \ \ \ \ \ \ \ \ \ \ \ \ \ \ \ \ \ \ \ \ \ \ \ \ \ \ \ \ \ \ \ \ \ \ \ \ \ \ \ \ \ \ \ \ \ \ \ \ \ \ \ \ \ \ \ \ \ \ \ \ \ \ \ \ \ \ \ \ \ \ \ \ \ \ \ \ \ \ \ \ \ \ \ \ \ \ \ \ \ \ \ \ \ \ \ \ \ \ \ \ \ \ \ \ \ \ \ \ \ \ \ \ \ \ \ \ \ \ \ \ \ \ \ \ \ \ \ \ \ \ \ \ \ \ \ \ \ \ \ \ \ \ \ \ \ \ \ \ \ \ \ \ \ \ \ \ \ \ \ \ \ \ \ \ \ \ \ \ \ \ \ \ \ \ \ \ \ \ \ \ \ \ \ \ \ \ \ \ \ \ \ \ \ \ \ \ \ \ \ \ \ \ \ \ \ \ \ \ \ \ \ \ \ \ \ \ \ \ \ \ \ \ \ \ \ \ \ \ \ \ \ \ \ \ \ \ \ \ \ \ \ \ \ \ \ \ \ \ \ \ \ \ \ \ \ \ \ \ \ \ \ \ \ \ \ \ \ \ \ \ \ \ \ \ \ \ \ \ \ \ \ \ \ \ \ \ \ \ \ \ \ \ \ \ \ \ \ \ \ \ \ \ \ \ \ \ \ \ \ \ \ \ \ \ \ \ \ \ \ \ \ \ \ \ \ \ \ \ \ \ \ \ \ \ \ \ \ \ \ \ \ \ \ \ \ \ \ \ \ \ \ \ \ \ \ \ \ \ \ \ \ \ \ \ \ \ \ \ \ \ \ \ \ \ \ \ \ \ \ \ \ \ \ \ \ \ \ \ \ \ \ \ \ \ \ \ \ \ \ \ \ \ \ \ \ \ \ \ \ \ \ \ \ \ \ \ \ \ \ \ \ \ \ \ \ \ \ \ \ \ \ \ \ \ \ \ \ \ \ \ \ \ \ \ \ \ \ \ \ \ \ \ \ \ \ \ \ \ \ \ \ \ \ \ \ \ \ \ \ \ \ \ \ \ \ \ \ \ \ \ \ \ \ \ \ \ \ \ \ \ \ \ \ \ \ \ \ \ \ \ \ \ \ \ \ \ \ \ \ \ \ \ \ \ \ \ \ \ \ \ \ \ \ \ \ \ \ \ \ \ \ \ \ \ \ \ \ \ \ \ \ \ \ \ \ \ \ \ \ \ \ \ \ \ \ \ \ \ \ \ \ \ \ \ \ \ \ \ \ \ \ \ \ \ \ \ \ \ \ \ \ \ \ \ \ \ \ \ \ \ \ \ \ \ \ \ \ \ \ \ \ \ \ \ \ \ \ \ \ \ \ \ \ \ \ \ \ \ \ \ \ \ \ \ \ \ \ \ \ \ \ \ \ \ \ \ \ \ \ \ \ \ \ \ \ \ \ \ \ \ \ \ \ \ \ \ \ \ \ \ \ \ \ \ \ \ \ \ \ \ \ \ \ \ \ \ \ \ \ \ \ \ \ \ \ \ \ \ \ \ \ \ \ \ \ \ \ \ \ \ \ \ \ \ \ \ \ \ \ \ \ \ \ \ \ \ \ \ \ \ \ \ \ \ \ \ \ \ \ \ \ \ \ \ \ \ \ \ \ \ \ \ \ \ \ \ \ \ \ \ \ \ \ \ \ \ \ \ \ \ \ \ \ \ \ \ \ \ \ \ \ \ \ \ \ \ \ \ \ \ \ \ \ \ \ \ \ \ \ \ \ \ \ \ \ \ \ \ \ \ \ \ \ \ \ \ \ \ \ \ \ \ \ \ \ \ \ \ \ \ \ \ \ \ \ \ \ \ \ \ \ \ \ \ \ \ \ \ \ \ \ \ \ \ \ \ \ \ \ \ \ \ \ \ \ \ \ \ \ \ \ \ \ \ \ \ \ \ \ \ \ \ \ \ \ \ \ \ \ \ \ \ \ \ \ \ \ \ \ \ \ \ \ \ \ \ \ \ \ \ \ \ \ \ \ \ \ \ \ \ \ \ \ \ \ \ \ \ \ \ \ \ \ \ \ \ \ \ \ \ \ \ \ \ \ \ \ \ \ \ \ \ \ \ \ \ \ \ \ \ \ \ \ \ \ \ \ \ \ \ \ \ \ \ \ \ \ \ \ \ \ \ \ \ \ \ \ \ \ \ \ \ \ \ \ \ \ \ \ \ \ \ \ \ \ \ \ \ \ \ \ \ \ \ \ \ \ \ \ \ \ \ \ \ \ \ \ \ \ \ \ \ \ \ \ \ \ \ \ \ \ \ \ \ \ \ \ \ \ \ \ \ \ \ \ \ \ \ \ \ \ \ \ \ \ \ \ \ \ \ \ \ \ \ \ \ \ \ \ \ \ \ \ \ \ \ \ \ \ \ \ \ \ \ \ \ \ \ \ \ \ \ \ \ \ \ \ \ \ \ \ \ \ \ \ \ \ \ \ \ \ \ \ \ \ \ \ \ \ \ \ \ \ \ \ \ \ \ \ \ \ \ \ \ \ \ \ \ \ \ \ \ \ \ \ \ \ \ \ \ \ \ \ \ \ \ \ \ \ \ \ \ \ \ \ \ \ \ \ \ \ \ \ \ \ \ \ \ \ \ \ \ \ \ \ \ \ \ \ \ \ \ \ \ \ \ \ \ \ \ \ \ \ \ \ \ \ \ \ \ \ \ \ \ \ \ \ \ \ \ \ \ \ \ \ \ \ \ \ \ \ \ \ \ \ \ \ \ \ \ \ \ \ \ \ \ \ \ \ \ \ \ \ \ \ \ \ \ \ \ \ \ \ \ \ \ \ \ \ \ \ \ \ \ \ \ \ \ \ \ \ \ \ \ \ \ \ \ \ \ \ \ \ \ \ \ \ \ \ \ \ \ \ \ \ \ \ \ \ \ \ \ \ \ \ \ \ \ \ \ \ \ \ \ \ \ \ \ \ \ \ \ \ \ \ \ \ \ \ \ \ \ \ \ \ \ \ \ \ \ \ \ \ \ \ \ \ \ \ \ \ \ \ \ \ \ \ \ \ \ \ \ \ \ \ \ \ \ \ \ \ \ \ \ \ \ \ \ \ \ \ \ \ \ \ \ \ \ \ \ \ \ \ \ \ \ \ \ \ \ \ \ \ \ \ \ \ \ \ \ \ \ \ \ \ \ \ \ \ \ \ \ \ \ \ \ \ \ \ \ \ \ \ \ \ \ \ \ \ \ \ \ \ \ \ \ \ \ \ \ \ \ \ \ \ \ \ \ \ \ \ \ \ \ \ \ \ \ \ \ \ \ \ \ \ \ \ \ \ \ \ \ \ \ \ \ \ \ \ \ \ \ \ \ \ \ \ \ \ \ \ \ \ \ \ \ \ \ \ \ \ \ \ \ \ \ \ \ \ \ \ \ \ \ \ \ \ \ \ \ \ \ \ \ \ \ \ \ \ \ \ \ \ \ \ \ \ \ \ \ \ \ \ \ \ \ \ \ \ \ \ \ \ \ \ \ \ \ \ \ \ \ \ \ \ \ \ \ \ \ \ \ \ \ \ \ \ \ \ \ \ \ \ \ \ \ \ \ \ \ \ \ \ \ \ \ \ \ \ \ \ \ \ \ \ \ \ \ \ \ \ \ \ \ \ \ \ \ \ \ \ \ \ \ \ \ \ \ \ \ \ \ \ \ \ \ \ \ \ \ \ \ \ \ \ \ \ \ \ \ \ \ \ \ \ \ \ \ \ \ \ \ \ \ \ \ \ \ \ \ \ \ \ \ \ \ \ \ \ \ \ \ \ \ \ \ \ \ \ \ \ \ \ \ \ \ \ \ \ \ \ \ \ \ \ \ \ \ \ \ \ \ \ \ \ \ \ \ \ \ \ \ \ \ \ \ \ \ \ \ \ \ \ \ \ \ \ \ \ \ \ \ \ \ \ \ \ \ \ \ \ \ \ \ \ \ \ \ \ \ \ \ \ \ \ \ \ \ \ \ \ \ \ \ \ \ \ \ \ \ \ \ \ \ \ \ \ \ \ \ \ \ \ \ \ \ \ \ \ \ \ \ \ \ \ \ \ \ \ \ \ \ \ \ \ \ \ \ \ \ \ \ \ \ \ \ \ \ \ \ \ \ \ \ \ \ \ \ \ \ \ \ \ \ \ \ \ \ \ \ \ \ \ \ \ \ \ \ \ \ \ \ \ \ \ \ \ \ \ \ \ \ \ \ \ \ \ \ \ \ \ \ \ \ \ \ \ \ \ \ \ \ \ \ \ \ \ \ \ \ \ \ \ \ \ \ \ \ \ \ \ \ \ \ \ \ \ \ \ \ \ \ \ \ \ \ \ \ \ \ \ \ \ \ \ \ \ \ \ \ \ \ \ \ \ \ \ \ \ \ \ \ \ \ \ \ \ \ \ \ \ \ \ \ \ \ \ \ \ \ \ \ \ \ \ \ \ \ \ \ \ \ \ \ \ \ \ \ \ \ \ \ \ \ \ \ \ \ \ \ \ \ \ \ \ \ \ \ \ \ \ \ \ \ \ \ \ \ \ \ \ \ \ \ \ \ \ \ \ \ \ \ \ \ \ \ \ \ \ \ \ \ \ \ \ \ \ \ \ \ \ \ \ \ \ \ \ \ \ \ \ \ \ \ \ \ \ \ \ \ \ \ \ \ \ \ \ \ \ \ \ \ \ \ \ \ \ \ \ \ \ \ \ \ \ \ \ \ \ \ \ \ \ \ \ \ \ \ \ \ \ \ \ \ \ \ \ \ \ \ \ \ \ \ \ \ \ \ \ \ \ \ \ \ \ \ \ \ \ \ \ \ \ \ \ \ \ \ \ \ \ \ \ \ \ \ \ \ \ \ \ \ \ \ \ \ \ \ \ \ \ \ \ \ \ \ \ \ \ \ \ \ \ \ \ \ \ \ \ \ \ \ \ \ \ \ \ \ \ \ \ \ \ \ \ \ \ \ \ \ \ \ \ \ \ \ \ \ \ \ \ \ \ \ \ \ \ \ \ \ \ \ \ \ \ \ \ \ \ \ \ \ \ \ \ \ \ \ \ \ \ \ \ \ \ \ \ \ \ \ \ \ \ \ \ \ \ \ \ \ \ \ \ \ \ \ \ \ \ \ \ \ \ \ \ \ \ \ \ \ \ \ \ \ \ \ \ \ \ \ \ \ \ \ \ \ \ \ \ \ \ \ \ \ \ \ \ \ \ \ \ \ \ \ \ \ \ \ \ \ \ \ \ \ \ \ \ \ \ \ \ \ \ \ \ \ \ \ \ \ \ \ \ \ \ \ \ \ \ \ \ \ \ \ \ \ \ \ \ \ \ \ \ \ \ \ \ \ \ \ \ \ \ \ \ \ \ \ \ \ \ \ \ \ \ \ \ \ \ \ \ \ \ \ \ \ \ \ \ \ \ \ \ \ \ \ \ \ \ \ \ \ \ \ \ \ \ \ \ \ \ \ \ \ \ \ \ \ \ \ \ \ \ \ \ \ \ \ \ \ \ \ \ \ \ \ \ \ \ \ \ \ \ \ \ \ \ \ \ \ \ \ \ \ \ \ \ \ \ \ \ \ \ \ \ \ \ \ \ \ \ \ \ \ \ \ \ \ \ \ \ \ \ \ \ \ \ \ \ \ \ \ \ \ \ \ \ \ \ \ \ \ \ \ \ \ \ \ \ \ \ \ \ \ \ \ \ \ \ \ \ \ \ \ \ \ \ \ \ \ \ \ \ \ \ \ \ \ \ \ \ \ \ \ \ \ \ \ \ \ \ \ \ \ \ \ \ \ \ \ \ \ \ \ \ \ \ \ \ \ \ \ \ \ \ \ \ \ \ \ \ \ \ \ \ \ \ \ \ \ \ \ \ \ \ \ \ \ \ \ \ \ \ \ \ \ \ \ \ \ \ \ \ \ \ \ \ \ \ \ \ \ \ \ \ \ \ \ \ \ \ \ \ \ \ \ \ \ \ \ \ \ \ \ \ \ \ \ \ \ \ \ \ \ \ \ \ \ \ \ \ \ \ \ \ \ \ \ \ \ \ \ \ \ \ \ \ \ \ \ \ \ \ \ \ \ \ \ \ \ \ \ \ \ \ \ \ \ \ \ \ \ \ \ \ \ \ \ \ \ \ \ \ \ \ \ \ \ \ \ \ \ \ \ \ \ \ \ \ \ \ \ \ \ \ \ \ \ \ \ \ \ \ \ \ \ \ \ \ \ \ \ \ \ \ \ \ \ \ \ \ \ \ \ \ \ \ \ \ \ \ \ \ \ \ \ \ \ \ \ \ \ \ \ \ \ \ \ \ \ \ \ \ \ \ \ \ \ \ \ \ \ \ \ \ \ \ \ \ \ \ \ \ \ \ \ \ \ \ \ \ \ \ \ \ \ \ \ \ \ \ \ \ \ \ \ \ \ \ \ \ \ \ \ \ \ \ \ \ \ \ \ \ \ \ \ \ \ \ \ \ \ \ \ \ \ \ \ \ \ \ \ \ \ \ \ \ \ \ \ \ \ \ \ \ \ \ \ \ \ \ \ \ \ \ \ \ \ \ \ \ \ \ \ \ \ \ \ \ \ \ \ \ \ \ \ \ \ \ \ \ \ \ \ \ \ \ \ \ \ \ \ \ \ \ \ \ \ \ \ \ \ \ \ \ \ \ \ \ \ \ \ \ \ \ \ \ \ \ \ \ \ \ \ \ \ \ \ \ \ \ \ \ \ \ \ \ \ \ \ \ \ \ \ \ \ \ \ \ \ \ \ \ \ \ \ \ \ \ \ \ \ \ \ \ \ \ \ \ \ \ \ \ \ \ \ \ \ \ \ \ \ \ \ \ \ \ \ \ \ \ \ \ \ \ \ \ \ \ \ \ \ \ \ \ \ \ \ \ \ \ \ \ \ \ \ \ \ \ \ \ \ \ \ \ \ \ \ \ \ \ \ \ \ \ \ \ \ \ \ \ \ \ \ \ \ \ \ \ \ \ \ \ \ \ \ \ \ \ \ \ \ \ \ \ \ \ \ \ \ \ \ \ \ \ \ \ \ \ \ \ \ \ \ \ \ \ \ \ \ \ \ \ \ \ \ \ \ \ \ \ \ \ \ \ \ \ \ \ \ \ \ \ \ \ \ \ \ \ \ \ \ \ \ \ \ \ \ \ \ \ \ \ \ \ \ \ \ \ \ \ \ \ \ \ \ \ \ \ \ \ \ \ \ \ \ \ \ \ \ \ \ \ \ \ \ \ \ \ \ \ \ \ \ \ \ \ \ \ \ \ \ \ \ \ \ \ \ \ \ \ \ \ \ \ \ \ \ \ \ \ \ \ \ \ \ \ \ \ \ \ \ \ \ \ \ \ \ \ \ \ \ \ \ \ \ \ \ \ \ \ \ \ \ \ \ \ \ \ \ \ \ \ \ \ \ \ \ \ \ \ \ \ \ \ \ \ \ \ \ \ \ \ \ \ \ \ \ \ \ \ \ \ \ \ \ \ \ \ \ \ \ \ \ \ \ \ \ \ \ \ \ \ \ \ \ \ \ \ \ \ \ \ \ \ \ \ \ \ \ \ \ \ \ \ \ \ \ \ \ \ \ \ \ \ \ \ \ \ \ \ \ \ \ \ \ \ \ \ \ \ \ \ \ \ \ \ \ \ \ \ \ \ \ \ \ \ \ \ \ \ \ \ \ \ \ \ \ \ \ \ \ \ \ \ \ \ \ \ \ \ \ \ \ \ \ \ \ \ \ \ \ \ \ \ \ \ \ \ \ \ \ \ \ \ \ \ \ \ \ \ \ \ \ \ \ \ \ \ \ \ \ \ \ \ \ \ \ \ \ \ \ \ \ \ \ \ \ \ \ \ \ \ \ \ \ \ \ \ \ \ \ \ \ \ \ \ \ \ \ \ \ \ \ \ \ \ \ \ \ \ \ \ \ \ \ \ \ \ \ \ \ \ \ \ \ \ \ \ \ \ \ \ \ \ \ \ \ \ \ \ \ \ \ \ \ \ \ \ \ \ \ \ \ \ \ \ \ \ \ \ \ \ \ \ \ \ \ \ \ \ \ \ \ \ \ \ \ \ \ \ \ \ \ \ \ \ \ \ \ \ \ \ \ \ \ \ \ \ \ \ \ \ \ \ \ \ \ \ \ \ \ \ \ \ \ \ \ \ \ \ \ \ \ \ \ \ \ \ \ \ \ \ \ \ \ \ \ \ \ \ \ \ \ \ \ \ \ \ \ \ \ \ \ \ \ \ \ \ \ \ \ \ \ \ \ \ \ \ \ \ \ \ \ \ \ \ \ \ \ \ \ \ \ \ \ \ \ \ \ \ \ \ \ \ \ \ \ \ \ \ \ \ \ \ \ \ \ \ \ \ \ \ \ \ \ \ \ \ \ \ \ \ \ \ \ \ \ \ \ \ \ \ \ \ \ \ \ \ \ \ \ \ \ \ \ \ \ \ \ \ \ \ \ \ \ \ \ \ \ \ \ \ \ \ \ \ \ \ \ \ \ \ \ \ \ \ \ \ \ \ \ \ \ \ \ \ \ \ \ \ \ \ \ \ \ \ \ \ \ \ \ \ \ \ \ \ \ \ \ \ \ \ \ \ \ \ \ \ \ \ \ \ \ \ \ \ \ \ \ \ \ \ \ \ \ \ \ \ \ \ \ \ \ \ \ \ \ \ \ \ \ \ \ \ \ \ \ \ \ \ \ \ \ \ \ \ \ \ \ \ \ \ \ \ \ \ \ \ \ \ \ \ \ \ \ \ \ \ \ \ \ \ \ \ \ \ \ \ \ \ \ \ \ \ \ \ \ \ \ \ \ \ \ \ \ \ \ \ \ \ \ \ \ \ \ \ \ \ \ \ \ \ \ \ \ \ \ \ \ \ \ \ \ \ \ \ \ \ \ \ \ \ \ \ \ \ \ \ \ \ \ \ \ \ \ \ \ \ \ \ \ \ \ \ \ \ \ \ \ \ \ \ \ \ \ \ \ \ \ \ \ \ \ \ \ \ \ \ \ \ \ \ \ \ \ \ \ \ \ \ \ \ \ \ \ \ \ \ \ \ \ \ \ \ \ \ \ \ \ \ \ \ \ \ \ \ \ \ \ \ \ \ \ \ \ \ \ \ \ \ \ \ \ \ \ \ \ \ \ \ \ \ \ \ \ \ \ \ \ \ \ \ \ \ \ \ \ \ \ \ \ \ \ \ \ \ \ \ \ \ \ \ \ \ \ \ \ \ \ \ \ \ \ \ \ \ \ \ \ \ \ \ \ \ \ \ \ \ \ \ \ \ \ \ \ \ \ \ \ \ \ \ \ \ \ \ \ \ \ \ \ \ \ \

\newpage

%% ===== Page 55 =====

- 35 -

dans $\varphi^{-1}(S')$ (resp. $\psi^{-1}(S')$). Alors le produit $U \times_S V$ s'identifie canoniquement au préschéma induit par $Z$ sur $p^{-1}(U) \cap q^{-1}(V)$ (considéré comme $S$-préschéma). En outre, si $f : T \to X$, $g : T \to Y$ sont des $S$-morphismes tels que $f(T) \subset U$, $g(T) \subset V$, le morphisme $(f,g)_S$ s'identifie à la restriction de $(f,g)_S$ à $p^{-1}(U) \cap q^{-1}(V)$.

Cela résulte de la prop. 7 et du lemme 3.

5. \emph{Propriétés formelles du produit : changement de préschéma de base.}
Le lecteur remarquera que toutes les propriétés de ce n° sont valables sans modification dans toute catégorie où la notion de produit de deux objets quelconques existe (car il est clair que la notion de $S$-objet et de $S$-morphisme peut se définir exactement comme au n° 3 pour tout objet $S$ de la catégorie).

\emph{Proposition 8.} - Pour tout $S$-préschéma $X$, la première (resp. seconde) projection de $X \times_S S$ (resp. $S \times_S X$) est un isomorphisme fonctoriel de $X \times_S S$ (resp. $S \times_S X$) sur $X$, dont l'isomorphisme réciproque est $(1_X, \varphi)$ (resp. $(\varphi, 1_X)$), en désignant par $\varphi$ le morphisme structural $X \to S$ : on peut donc écrire
$$ X \times_S S = S \times_S X = X $$

Il suffit de prouver que $(X, 1_X, \varphi)$ est un produit de $X$ et $S$. Or, si $T$ est un $S$-préschéma, le seul $S$-morphisme de $T$ dans $S$ est nécessairement le morphisme structural $\psi : T \to S$. Si alors $f$ est un $S$-morphisme de $T$ dans $X$, on a nécessairement $f = \varphi \circ f$, d'où notre assertion.

On peut définir de la même manière qu'au n° 4 le produit d'un nombre fini quelconque $n$ de $S$-préschémas ; l'existence de ces produits résulte du th. 1 du n° 4 par récurrence sur $n$, en remarquant que
$(X_1 \times_S \dots \times_S X_{n-1}) \times_S X_n$ satisfait à la définition du produit. L'unicité du produit entraîne, comme dans toute catégorie, ses propriétés de commutativité et d'associativité.

Soient $S', S$ deux préschémas, $\varphi : S' \to S$ un morphisme, qui fait de $S'$ un $S$-préschéma. Pour tout $S$-préschéma $X$, considérons le produit

\newpage

%% ===== Page 56 =====

- 36 -
$X \times_{S} S'$, et soient $p$ et $\pi'$ ses projections dans $X$ et $S'$ respectivement. Muni de $\pi'$, ce produit est un $S'$-préschéma ; quand on le considère comme tel, on le désigne par $X^{S'}$ ou $X_{\varphi}$, et on dit que c'est le préschéma obtenu par extension du préschéma de base de $S$ à $S'$, au moyen de $\varphi$. On notera que si $\pi$ est le morphisme structural de $X$, le diagramme
$$
\begin{tikzcd}
X \arrow[d, "\pi"] & X^{S'} \arrow[l, "p"] \arrow[d, "\pi'"] \\
S \arrow[r, "\varphi"] & S'
\end{tikzcd}
$$
est commutatif ($\theta$ étant le morphisme structural de $X \times_{S} S'$, considéré comme $S$-préschéma). Lorsque $S$ et $S'$ sont des schémas affines d'anneaux $A, A'$ respectivement, on écrit aussi $X^{A'}$ ou $X \otimes_{A} A'$ au lieu de $X^{S'}$; si $X$ est lui-même affine d'anneau $B$, $X^{A'}$ est affine d'anneau $B \otimes_{A} A'$, obtenu par extension à $A'$ de l'anneau des scalaires de l'algèbre $B$ sur $A$.

Le préschéma $X^{S'}$ peut encore être considéré comme solution d'un problème d'application universelle : tout $S'$-préschéma $T$ est aussi un $S$-préschéma ; tout morphisme $g : T \to X$ s'écrit alors d'une seule manière $g = p \circ f$, où $f$ est un $S'$-morphisme $T \to X^{S'}$, comme il résulte de la définition du produit appliquée aux morphismes $g$ et $\varphi : T \to S'$.

Proposition 9 ("transitivité de l'extension des préschémas de base"). Soient $S''$ un préschéma, $\varphi' : S'' \to S'$ un morphisme de $S''$ dans $S'$. Pour tout $S$-préschéma $X$, il existe un isomorphisme canonique fonctoriel du $S''$-préschéma $(X^{S'})^{S''}$ sur le $S''$-préschéma $X^{S''}$.

En effet, soient $T$ un $S''$-préschéma, $\psi$ son morphisme structural, $g$ un $S$-morphisme de $T$ dans $X$ ($T$ étant considéré comme $S$-préschéma de morphisme structural $\varphi \circ \varphi'$). Comme $T$ est aussi un $S'$-préschéma de morphisme structural $\varphi' \circ \psi$, on peut écrire $g = p \circ g'$, où $g'$ est un $S'$-morphisme $T \to X^{S'}$, puis $g' = p' \circ g''$, où $g''$ est un $S''$-morphisme $T \to (X^{S'})^{S''}$.

\newpage

%% ===== Page 57 =====

- 37 -
\begin{tikzcd}
X & X \times_S S' \arrow[l] \arrow[d, "\pi"] & (X \times_S S') \times_{S'} S'' \arrow[l] \arrow[d, "\pi''"] \\
S & S' \arrow[l, "\varphi"] & S'' \arrow[l, "\varphi'"]
\end{tikzcd}
D'où la proposition en raison de l'unicité de la solution d'un problème universel.

Ce résultat s'exprime encore en écrivant l'égalité (à un isomorphisme près) $(X^{S'})^{S''} = X^{S''}$, si aucune confusion n'est à craindre, ou aussi (à un isomorphisme près)
$$(1) \quad (X \times_S S') \times_{S'} S'' = X \times_S S''$$
\emph{Corollaire 1}. - Si $X$ et $Y$ sont deux $S$-préschémas, il existe un isomorphisme canonique fonctoriel du $S'$-préschéma $X^{S'} \times_{S'} Y^{S'}$ sur le $S'$-préschéma $(X \times_S Y)^{S'}$.

En effet, on a (à des isomorphismes canoniques près)
$$(X \times_S S') \times_{S'} (Y \times_S S') = X \times_S (Y \times_S S') = (X \times_S Y) \times_S S'$$
en vertu de (1) et de l'associativité des produits de $S$-préschémas.

6. \emph{Morphismes surjectifs}.

\emph{Définition 6}. - On dit qu'un morphisme $(\varphi, \theta) : (X, \mathcal{O}_X) \to (Y, \mathcal{O}_Y)$ de préschémas est surjectif si l'application continue $\varphi$ est surjective.

On notera que cela n'entraîne nullement que $(\varphi, \theta)$ soit un épimorphisme d'espaces annelés ($\theta$ n'est pas nécessairement injectif).

Il est clair que le composé de deux morphismes surjectifs est surjectif.

\emph{Proposition 10}. - Soient $X, Y$ deux $S$-préschémas, $x, y$ des points appartenant respectivement aux espaces de base $X$ et $Y$. Pour qu'il existe un point de l'espace de base du produit $X \times_S Y$ dont les projections soient égales à $x$ et $y$, il faut et il suffit que $x$ et $y$ soient au-dessus d'un même point de l'espace de base de $S$.

La condition est évidemment nécessaire. Si on la suppose remplie, il existe des monomorphismes du corps $k(s)$ dans chacun des deux corps $k(x), k(y)$. Il existe d'autre part des monomorphismes de $k(x)$ et $k(y)$

\newpage

%% ===== Page 58 =====

dans un même corps $K$ tels que les monomorphismes composés
$k(s) \to k(x) \to K$ et $k(s) \to k(y) \to K$ soient identiques (Bourbaki, Alg., chap. V, §4, prop. 2). On en déduit que si $\{s\} = \operatorname{Spec}(K)$, les morphismes correspondants $\{s\} \to X, \{s\} \to Y$ (n°2, cor. de la prop. 5) sont des $S$-morphismes, et par suite définissent un $S$-morphisme $f: \{s\} \to X \times_S Y$.
Si on pose $z=f(s)$, il est clair par définition que $p_1(z)=x$ et $p_2(z)=y$.

Il y a donc une application surjective canonique de l'espace de base $X \times_S Y$ sur la partie de l'ensemble produit $X \times Y$ des espaces de base $X, Y$, formée des couples $(x, y)$ qui sont au-dessus d'un même point de $S$. Mais il faut noter que cette application n'est pas injective en général (autrement dit, il peut exister plusieurs $z$ distincts tels que $p_1(z)$ et $p_2(z)$ soient donnés).

Corollaire. — Soient $f: X \to Y$ un $S$-morphisme, $f^{S'}: X^{S'} \to Y^{S'}$ le $S'$-morphisme déduit de $f$ par une extension $S' \to S$ du pré-schéma de base. Soient $p$ (resp. $q$) la projection $X^{S'} \to X$ (resp. $Y^{S'} \to Y$) ; pour toute partie $M$ de l'espace de base $X$, on a $q^{-1}(f(M)) = f^{S'}(p^{-1}(M))$.

En effet (cor. 2 de la prop. 9), $X^{S'}$ s'identifie au produit $X \times_S Y^{S'}$ par le diagramme
\begin{tikzcd}
X \arrow[d] & X^{S'} \arrow[l, "p"] \arrow[d, "f^{S'}"] \\
Y & Y^{S'} \arrow[l, "q"]
\end{tikzcd}

En vertu de la prop. 10, la relation $q(y')=f(x)$ pour $x \in X, y' \in Y^{S'}$, équivaut donc à l'existence d'un $x' \in X^{S'}$ tel que $p(x') = x$ et $f^{S'}(x')=y'$, d'où le corollaire.

Proposition 11. — a) Si un $S$-morphisme $f: X \to Y$ est surjectif, tout $S'$-morphisme $f^{S'}: X^{S'} \to Y^{S'}$ déduit de $f$ par extension du pré-schéma de base est surjectif.

b) Si $f: X' \to X, g: Y' \to Y$ sont des $S$-morphismes surjectifs, le $S$-morphisme $f \times_S g: X' \times_S Y' \to X \times_S Y$ est surjectif.

La première assertion est une conséquence immédiate du cor. de la

\newpage

%% ===== Page 59 =====

\noindent -- 39 --

\noindent prop. 10. La seconde s'en déduit en remarquant que $f \times_S g$ est le morphisme composé
$$X' \times_{S'} Y' \xrightarrow{f \times 1} X \times_S Y' \xrightarrow{1 \times g} X \times_S Y$$

\noindent Proposition 12. -- Pour qu'un morphisme $f : X \to Y$ soit surjectif, il faut et il suffit que pour tout corps $K$ et tout morphisme $\operatorname{Spec}(K) \to Y$, il existe une extension $K'$ de $k$ et un morphisme $\operatorname{Spec}(K') \to X$ rendant commutatif le diagramme
$$\begin{tikzcd}
X \arrow[d, "f"] & \operatorname{Spec}(K') \arrow[l] \\
Y & \operatorname{Spec}(K) \arrow[l] \arrow[u]
\end{tikzcd}$$

\noindent La condition est suffisante, car pour tout $y \in Y$, il suffit de l'appliquer à un morphisme $\operatorname{Spec}(k(y)) \to Y$ correspondant à un monomorphisme $k(y) \to K$, $K$ étant une extension de $k(y)$ ($n^\circ 2$, cor. de la prop. 5). Inversement, supposons $f$ surjectif, et soit $y$ l'image de l'unique élément de $\operatorname{Spec}(k(y))$; il existe $x \in X$ tel que $f(x) = y$, donc $f$ définit un monomorphisme $k(y) \to k(x)$; il suffit alors de prendre $K'$ tel qu'il existe des monomorphismes $k(x) \to K'$ et $k(y) \to K'$ pour lesquelles $k(y) \to k(x) \to K'$ et $k(y) \to K'$ coïncident; le morphisme $\operatorname{Spec}(K') \to X$ correspondant à $k(x) \to K'$ répond à la question.

\noindent 7. Morphismes géométriquement injectifs.

\noindent Définition 7. -- On dit qu'un morphisme $f : X \to Y$ de préschémas est géométriquement injectif si, pour tout corps $K$, l'application $u \mapsto f \circ u$ de $\operatorname{Hom}(\operatorname{Spec}(K), X)$ dans $\operatorname{Hom}(\operatorname{Spec}(K), Y)$ est injective.

\noindent Pour que $f$ soit géométriquement injectif, il suffit que la condition de la déf. 7 soit vérifiée pour tout corps algébriquement clos. En effet, si $K$ est un corps quelconque, $K'$ une extension algébriquement close de $K$, le diagramme
$$\begin{tikzcd}
\operatorname{Hom}(\operatorname{Spec}(K), X) \arrow[r, "\alpha"] \arrow[d, "\psi"] & \operatorname{Hom}(\operatorname{Spec}(K), Y) \arrow[d, "\psi'"] \\
\operatorname{Hom}(\operatorname{Spec}(K'), X) \arrow[r, "\alpha'"] & \operatorname{Hom}(\operatorname{Spec}(K'), Y)
\end{tikzcd}$$
est commutatif, $\psi$ et $\psi'$ provenant du morphisme $\operatorname{Spec}(K') \to \operatorname{Spec}(K)$

\newpage

%% ===== Page 60 =====

et $\alpha, \alpha'$ sont les homomorphismes $u \mapsto f \circ u$. Or $\varphi$ est injectif, et il en est de même de $\alpha'$ par hypothèse, donc $\alpha$ est nécessairement injectif.

Il est clair que le composé de deux morphismes géométriquement injectifs est géométriquement injectif.

Proposition 13. — Pour qu'un morphisme $f = (\varphi, \theta) : X \to Y$ soit géométriquement injectif, il faut et il suffit que $\varphi$ soit injectif et que, pour tout $x \in X$, le monomorphisme $\theta^x : k(\varphi(x)) \to k(x)$ fasse de $k(x)$ une extension radicielle de $k(\varphi(x))$.

Supposons $f$ géométriquement injectif ; montrons d'abord que la relation $\varphi(x_1) = \varphi(x_2) = y$ entraîne nécessairement $x_1 = x_2$. En effet, il existe un corps $K$ et des monomorphismes $k(x_1) \to K$, $k(x_2) \to K$ tels que les monomorphismes $k(y) \to k(x_1) \to K$ et $k(y) \to k(x_2) \to K$ soient identiques ; les morphismes correspondants $u_1 : \operatorname{Spec}(K) \to X$, $u_2 : \operatorname{Spec}(K) \to X$ sont donc tels que $f \circ u_1 = f \circ u_2$, donc $u_1 = u_2$ par hypothèse, et cela entraîne $x_1 = x_2$. Considérons maintenant $k(x)$ comme extension de $k(\varphi(x))$ au moyen de $\theta^x$ ; si $k(x)$ n'est pas extension radicielle de $k(\varphi(x))$, il existe deux $k(\varphi(x))$-isomorphismes distincts de $k(x)$ dans une extension algébriquement close $K$ de $k(\varphi(x))$ et les deux morphismes correspondants de $\operatorname{Spec}(K)$ dans $X$ violeraient l'hypothèse. Inversement, il est clair que les conditions de l'énoncé sont suffisantes pour que $f$ soit géométriquement injectif.

Corollaire. — Si un morphisme $f = (\varphi, \theta) : X \to Y$ est tel que $\varphi$ soit injectif et que, pour tout $x \in X$, l'homomorphisme $\theta^x : k(\varphi(x)) \to k(x)$ soit bijectif, $f$ est géométriquement injectif.

Proposition 14. — Si $f : X \to Y$, $g : Y \to Z$ sont deux morphismes tels que $g \circ f$ soit géométriquement injectif, alors $f$ est géométriquement injectif.

En effet, l'application composée
$$\operatorname{Hom}(\operatorname{Spec}(K), X) \to \operatorname{Hom}(\operatorname{Spec}(K), Y) \to \operatorname{Hom}(\operatorname{Spec}(K), Z)$$

\newpage

%% ===== Page 61 =====

- 41 -

étant injective, il en est de même de $\operatorname{Hom}(\operatorname{Spec}(K), X) \to \operatorname{Hom}(\operatorname{Spec}(K), Y)$.

Proposition 15. - Si les $S$-morphismes $f : X \to X'$, $g : Y \to Y'$ sont géométriquement injectifs, il en est de même de $f \times g$.

En effet, $\operatorname{Hom}(\operatorname{Spec}(K), X \times_S Y)$ s'identifie à l'ensemble des couples $(u, v) \in \operatorname{Hom}(\operatorname{Spec}(K), X) \times \operatorname{Hom}(\operatorname{Spec}(K), Y)$ tels que $f \circ u = g \circ v$, en désignant par $f$ et $g$ les morphismes structuraux de $X$ et $Y$. On a une identification analogue pour $\operatorname{Hom}(\operatorname{Spec}(K), X' \times_S Y')$, et l'application $\operatorname{Hom}(\operatorname{Spec}(K), X \times_S Y) \to \operatorname{Hom}(\operatorname{Spec}(K), X' \times_S Y')$ correspond alors à $(u, v) \to (f \circ u, g \circ v)$; la proposition résulte trivialement de là.

Proposition 16. - Soient $f : X \to Y$, $j : Y' \to Y$ deux morphismes, tels que $j$ soit une injection géométrique. Alors la projection de $X \times_Y Y'$ dans $X$ est une injection géométrique, et l'image par cette projection est l'image réciproque $f^{-1}(Y')$.

En effet, (cor. de la prop. 8), si on identifie canoniquement $X$ et $X \times_Y Y$, la projection $X \times_Y Y' \to X$ est identifiée à $1 \times j$, donc est une injection géométrique par la prop. 15. En outre, en vertu de la prop. 10 du n° 6, les points $x \in X$ qui sont images de points de $X \times_Y Y'$ sont ceux pour lesquels il existe $y' \in Y'$ tels que $f(x) = j(y')$, d'où la proposition.

Considérons en particulier un point quelconque $y \in Y$, et prenons $Y' = \operatorname{Spec}(k(y))$; l'application identique $k(y) \to k(y)$ définit canoniquement un morphisme $Y' \to Y$, qui applique l'unique point de $Y'$ sur $y$ et est géométriquement injective (cor. de la prop. 13). La prop. 16 montre que l'ensemble de base du préschéma $X \times_Y Y'$ s'identifie canoniquement avec $f^{-1}(y)$. La projection $X \times_Y Y' \to Y'$ définit d'ailleurs $X \times_Y Y'$ comme un préschéma sur $k(y)$; c'est toujours de ce préschéma qu'il sera question lorsque nous parlerons de $f^{-1}(y)$ comme d'un préschéma.

\newpage

%% ===== Page 62 =====

\marginpar{Archives Grothendieck sept. 89}

\section*{§ 3. Sous-préschémas ; immersions. Morphismes et préschémas séparés.}

\subsection*{1. Sous-préschémas.}

Comme la notion de faisceau quasi-cohérent (chap. 0, § 2, n°6) est locale, un faisceau quasi-cohérent sur un préschéma $X$ peut être défini par la condition que, pour tout ouvert affine $V$ de $X$, $F|_V$ est isomorphe à un faisceau associé à un $\Gamma(V, \mathcal{O}_X)$-module (§ 1, n°4, th. 2). Il est clair que le faisceau structural $\mathcal{O}_X$ est quasi-cohérent et que noyaux, conoyaux, images, limites inductives et sommes directes de faisceaux quasi-cohérents sur $X$ sont quasi-cohérents (§ 1, n°3, cor. 4 du th. 1).

Proposition 1. — Soit $I$ un faisceau quasi-cohérent d'idéaux sur un préschéma $X$. Le support $Y$ du faisceau $\mathcal{O}_X/I$ est alors fermé, et si on désigne par $\mathcal{O}_Y$ la restriction de $\mathcal{O}_X/I$ à $Y$, $(Y, \mathcal{O}_Y)$ est un préschéma.

Il suffit évidemment (§ 2, n°1, lemme 1) de considérer le cas où $(X, \mathcal{O}_X)$ est un schéma affine, et de montrer que dans ce cas $Y$ est fermé dans $X$, et $(Y, \mathcal{O}_Y)$ un schéma affine. En effet, si $X = \operatorname{Spec}(A)$, on a $\mathcal{O}_X = \tilde{A}$ et $I = \tilde{J}$, où $J$ est un idéal de $A$, et $Y$ est égal à la partie fermée $V(J)$ de $X$ et s'identifie au spectre premier de l'anneau $B = A/J$ (§ 1, n°2, cor. 2 de la prop. 5). De plus, si $\varphi$ est l'homomorphisme canonique $A \to A/J$, l'image directe $\varphi_*(\tilde{B})$ s'identifie canoniquement au faisceau $\tilde{J} = \mathcal{O}_X/I$ (§ 1, n°6, prop. 7 et n°3, cor. 4 au th. 1).

Nous dirons que $(Y, \mathcal{O}_Y)$ est le sous-préschéma de $(X, \mathcal{O}_X)$ défini par le faisceau d'idéaux $I$ ; c'est un cas particulier de la notion générale de sous-préschéma :

Définition 1. — On dit qu'un espace annelé $(Y, \mathcal{O}_Y)$ est un sous-préschéma d'un préschéma $(X, \mathcal{O}_X)$ si $Y$ est un sous-espace (resp. un sous-espace fermé) de $X$ et si, pour tout $y \in Y$, il existe un voisinage ouvert affine $V$ de $y$ dans $X$ tel que $(V \cap Y, \mathcal{O}_Y|_{V \cap Y})$ soit un sous-préschéma ...

\newpage

%% ===== Page 63 =====

ne du schéma affine $(V, \mathcal{O}_V)$ défini par un faisceau quasi-cohérent d'idéaux.

Il est clair d'après la prop. 1 qu'un sous-préschéma $(Y, \mathcal{O}_Y)$ est un préschéma ; on notera en outre que $Y \cap V$ est fermé dans $V$, donc $Y$ est nécessairement localement fermé dans $X$. Si $V$ est ouvert dans $X$, le préschéma induit par $(X, \mathcal{O}_X)$ sur $V$ est un sous-préschéma en vertu du lemme 1 du $\S 2, n^\circ 1$.

\emph{Proposition 2.} -- Pour qu'un sous-préschéma $(Y, \mathcal{O}_Y)$ du préschéma $(X, \mathcal{O}_X)$ soit fermé, il faut et il suffit qu'il soit défini par un faisceau de $\mathcal{O}_X$-idéaux ; il y a correspondance bijective entre ces faisceaux et les sous-préschémas fermés de $(X, \mathcal{O}_X)$.

Il suffit de prouver que tout sous-préschéma fermé $(Y, \mathcal{O}_Y)$ est défini par un faisceau quasi-cohérent unique. Or, par définition, pour tout $y \in Y$, on a $\mathcal{O}_Y(y) = \mathcal{O}_X(y) / I(y)$, où $I(y)$ est un idéal bien déterminé de $\mathcal{O}_X(y)$ ; posons $I(x) = \mathcal{O}_X(x)$ pour $x \notin Y$. Reste à voir que les $I(x)$ sont les fibres d'un faisceau quasi-cohérent d'idéaux $I$ tel que $\mathcal{O}_Y = (\mathcal{O}_X/I)|_Y$ ; comme la question est locale, cela résulte de la définition et de l'hypothèse que $Y$ est fermé.

\emph{Proposition 3.} -- Un sous-préschéma (resp. un sous-préschéma fermé) d'un sous-préschéma (resp. sous-préschéma fermé) de $X$ s'identifie canoniquement à un sous-préschéma (resp. sous-préschéma fermé) de $X$.

Une partie fermée d'un sous-espace fermé de l'espace de base $X$ étant un sous-espace fermé de $X$, il est clair que la question est locale et qu'on peut donc supposer $X$ affine ; la proposition résulte donc de l'identification canonique de $A/J$ et de $(A/I)/(J/I)$ lorsque $I, J$ sont deux idéaux d'un anneau $A$ tels que $I \subset J$.

Nous ferons toujours dans la suite l'identification précédente. En particulier, si $(Y, \mathcal{O}_Y)$ est un sous-préschéma de $(X, \mathcal{O}_X)$, l'ensemble $Y$, étant localement fermé, est l'intersection d'un ensemble ouvert $U$ et d'un ensemble fermé ; on en conclut que $(Y, \mathcal{O}_Y)$ est un sous-préschéma.

\newpage

%% ===== Page 64 =====

- 44 -

schéma fermé du préschéma $(V, \mathcal{O}_X|V)$ induit par $(X, \mathcal{O}_X)$ sur $V$.

Soit $(Y, \mathcal{O}_Y)$ un sous-préschéma d'un préschéma $(X, \mathcal{O}_X)$ ; désignons par $j$ l'injection canonique $Y \to X$ ; on sait que l'image réciproque $j^*(\mathcal{O}_X)$ est la restriction $\mathcal{O}_X|Y$. Si, pour tout $y \in Y$, on désigne par $h_y$ l'homomorphisme canonique de la fibre $\mathcal{O}_X(y)$ sur la fibre $\mathcal{O}_Y(y)$, ces homomorphismes définissent un homomorphisme $h$ du faisceau $\mathcal{O}_Y = j^*(\mathcal{O}_X)$ sur $\mathcal{O}_Y$ : il suffit en effet de le vérifier localement, c'est-à-dire que l'on peut supposer $(X, \mathcal{O}_X)$ affine et le sous-préschéma $(Y, \mathcal{O}_Y)$ fermé ; et dans ce cas les $h_y$ ne sont autres que les restrictions aux fibres $\mathcal{O}_X(y)$ de l'homomorphisme de faisceaux $\mathcal{O}_X|Y \to (\mathcal{O}_X/I)|Y$ si $I$ est le faisceau d'idéaux qui définit $(Y, \mathcal{O}_Y)$. Nous avons donc défini un homomorphisme (chap. 0, § 2, n°4) du préschéma $(Y, \mathcal{O}_Y)$ dans le préschéma $(X, \mathcal{O}_X)$, qui sera appelé le morphisme d'injection $j$.

Conformément aux définitions générales (T, I, 1.1) nous dirons qu'un morphisme $f : Z \to X$ est majoré par le morphisme d'injection $j : Y \to X$ d'un sous-préschéma $Y$ de $X$, si $f$ se factorise en
$$Z \xrightarrow{g} Y \xrightarrow{j} X$$
(nécessairement unique puisque $j$ est un monomorphisme) où $g$ est un morphisme ; cela entraîne que $f(Z) \subset Y$ et que, pour tout $y=f(z) \in Y$, l'homomorphisme $\mathcal{O}_X(y) \to \mathcal{O}_Z(z)$ correspondant à $f$ se factorise en $\mathcal{O}_X(y) \to \mathcal{O}_Y(y) \to \mathcal{O}_Z(z)$ ; le noyau de $\mathcal{O}_X(y) \to \mathcal{O}_Z(z)$ contient donc le noyau $I_Y(y)$ de $\mathcal{O}_X(y) \to \mathcal{O}_Y(y)$. En particulier, on dit qu'un sous-préschéma $Z$ de $X$ est majoré par $Y$ si le morphisme d'injection $Z \to X$ est majoré par $j$. Pour qu'il en soit ainsi, on peut se ramener au cas où $X$ est affine et $Z$ fermé (la question étant locale, en raison de l'unicité de la solution $g$ de $j'=j \circ g$). On voit alors que la relation $Z \le Y$ équivaut à dire que l'espace de base $Z$ est un sous-espace de $Y$, et que pour tout $z \in Z$, si on pose $\mathcal{O}_Y(z)=\mathcal{O}_Y(z)/I_Y(z)$, $\mathcal{O}_Z(z)=\mathcal{O}_Z(z)/I_Z(z)$, on doit avoir $I_Z(z) \supset I_Y(z)$. La relation $Z \le Y$ est évidemment une relation d'ordre dans l'ensemble des sous-préschémas de $X$.

\newpage

%% ===== Page 65 =====

- 45 -

2. Morphismes d'immersion.

Définition 2. - On dit qu'un morphisme $f : Y \to X$ est une immersion (resp. une immersion fermée) s'il se factorise en un isomorphisme $g : Y \to Z$ sur un sous-préschéma (resp. un sous-préschéma fermé) $Z$ de $X$ et en le morphisme d'injection $j : Z \to X$.

Le sous-préschéma $Z$ et l'isomorphisme $g$ sont alors déterminés de façon unique, car si $Z'$ est un second sous-préschéma, $j'$ l'injection $Z' \to X$ et $g'$ un isomorphisme $Y \to Z'$ tel que $j \circ g = j' \circ g'$, on en déduit $j' = j \circ g \circ g'^{-1}$, d'où $Z' \leq Z$ et on montre de même que $Z \leq Z'$, donc $Z' = Z$, et comme $j$ est un monomorphisme de préschémas, $g' = g$.

Il est clair qu'une immersion est un monomorphisme de préschémas et une injection géométrique, puisqu'il en est ainsi de $j$ (cf. §2, n°7, cor. de la prop. 13).

Proposition 4. - Pour qu'un morphisme $f = (\varphi, \theta) : Y \to X$ soit une immersion (resp. une immersion fermée), il faut et il suffit que $\varphi$ soit un homéomorphisme de $Y$ sur une partie localement fermée (resp. fermée) de $X$, et que pour tout $y \in Y$, l'homomorphisme $\theta_y : \mathcal{O}_{X, \varphi(y)} \to \mathcal{O}_{Y, y}$ soit surjectif.

Les conditions étant évidemment nécessaires, montrons qu'elles sont suffisantes. Considérons d'abord le cas particulier où l'on suppose que $(X, \mathcal{O}_X)$ est un schéma affine, et que $\varphi(Y)$ est fermé dans $X$. Posons $Z = \varphi(Y)$; on sait (chap. 0, §2, n°1) que le faisceau $\varphi_*(\mathcal{O}_Y)$ a alors pour support $Z$ et que, si on désigne par $\mathcal{O}_Z$ sa restriction à $Z$, l'espace annelé $(Z, \mathcal{O}_Z)$ se déduit de $(Y, \mathcal{O}_Y)$ par transport de structure au moyen de l'homéomorphisme $\varphi$ (considéré comme application de $Y$ sur $Z$). Montrons que le faisceau $\varphi_*(\mathcal{O}_Y)$ sur $X$ est quasi-cohérent. En effet, pour tout $x \notin Z$, $\varphi_*(\mathcal{O}_Y)$, restreint à un voisinage de $x$, est nul. Si au contraire $x \in Z$, on a $x = \varphi(y)$ pour un $y \in Y$ bien déterminé; soit $V$ un voisinage ouvert affine de $y$ dans $Y$; $\varphi(V)$ est alors ouvert dans $Z$, donc trace sur $Z$ d'un ouvert $U$ de $X$, et la restric-

\newpage

%% ===== Page 66 =====

- 46 -

tion à $U$ de $\psi_*(\mathcal{O}_Y)$ est identique à la restriction à $U$ de l'image directe $(\psi_V)_*(\mathcal{O}_Y|V)$, où $\psi_V$ est la restriction à $V$ de $\psi$. Or, la restriction à $(V, \mathcal{O}_Y|V)$ du morphisme $(\psi, \theta)$ est un morphisme de ce préschéma dans $(X, \mathcal{O}_X)$, et par suite est de la forme $(\psi, \tilde{\varphi})$, où $\varphi$ est un homomorphisme de l'anneau $A = \Gamma(X, \mathcal{O}_X)$ dans l'anneau $\Gamma(V, \mathcal{O}_Y)$ (§ 1, n°7, th. 4) ; on en conclut que l'image directe $(\psi_V)_*(\mathcal{O}_Y|V)$ est un faisceau quasi-cohérent (§ 1, n°6, cor. 2 de la prop. 7), ce qui prouve notre assertion en raison du caractère local des faisceaux quasi-cohérents. En outre, l'hypothèse que $\psi$ est un homéomorphisme entraîne que pour tout $y \in Y$, $\psi_y$ est un isomorphisme d'anneau de $\psi_*(\mathcal{O}_Y)(\psi(y))$ sur $\mathcal{O}_Y(y)$ ; comme le diagramme

$$\begin{tikzcd}
\mathcal{O}_X(\psi(y)) \arrow[r] \arrow[d, "\psi_y"] & \psi_*(\mathcal{O}_Y)(\psi(y)) \arrow[d, "\theta_y"] \\
\mathcal{O}_Y(y) \arrow[r, "\theta_y^b"] & \mathcal{O}_Y(y)
\end{tikzcd}$$

est commutatif et que les deux flèches verticales sont des isomorphismes, l'hypothèse que $\theta_y^b$ est surjectif entraîne qu'il en est de même de $\theta_y$. Comme le support de $\psi_*(\mathcal{O}_Y)$ est $Z = \psi(Y)$, $\theta$ est un homomorphisme surjectif de $\mathcal{O}_X$ dans le faisceau quasi-cohérent $\psi_*(\mathcal{O}_Y)$. Par suite, il existe un isomorphisme unique $\omega$ d'un faisceau quotient $\tilde{\mathcal{J}}$ ($\mathcal{J}$ idéal de $A$) sur $\psi_*(\mathcal{O}_Y)$ qui, composé avec l'homomorphisme canonique de $\tilde{A}$ sur $\tilde{A}/\mathcal{J}$, donne $\theta$ (§ 1, n°3, cor. 4 du th. 1) ; si $\mathcal{O}_Z$ désigne la restriction de $\tilde{A}/\mathcal{J}$ à $Z = \psi(Y)$, $(Z, \mathcal{O}_Z)$ est un sous-préschéma de $(X, \mathcal{O}_X)$, et $f$ se factorise en l'injection de ce sous-préschéma dans $(X, \mathcal{O}_X)$, et l'isomorphisme $(\psi_0, \omega)$, où $\psi_0$ est $f$ considérée comme une application de $Y$ sur $Z$, et $\omega$ la restriction de $\omega$ à $\mathcal{O}_Z$.

Passons au cas général. Soit $U$ un ensemble ouvert affine dans $X$ tel que $U \cap \psi(Y)$ soit fermé dans $U$ et non vide. En restreignant $f$ au préschéma induit par $(Y, \mathcal{O}_Y)$ sur l'ouvert $\psi^{-1}(U)$, et en le considérant comme un morphisme de préschéma dans le préschéma induit par $(X, \mathcal{O}_X)$ sur $V$, on est ramené au premier cas ; la restriction de $f$

\newpage

%% ===== Page 67 =====

- 47 -

à $\psi(U)$ est donc une immersion fermée $\psi(U) \to U$, se factorisant canoniquement en $j_U \circ \varepsilon_U$, où $\varepsilon_U$ est un isomorphisme du préschéma $\psi'(U)$ sur un sous-préschéma $Z_U$ de $U$, et $j_U$ l'injection de $Z_U$ dans $U$. Soit $V$ un second ouvert affine de $X$ tel que $V \subset U$. Comme la restriction $Z_V$ de $Z_U$ à $V$ est un sous-préschéma du préschéma $V$, la restriction de $f$ à $\psi^{-1}(V)$ se factorise en $j_V \circ \varepsilon_V$, où $j_V$ est l'injection de $Z_V$ dans $V$ et $\varepsilon_V$ un isomorphisme de $\psi^{-1}(V)$ sur $Z_V$. En raison de l'unicité de la factorisation canonique d'une immersion, on voit que l'on a nécessairement $Z_U = Z_V$, $\varepsilon_U = \varepsilon_V$. On en conclut (déf.1, n°1) qu'il y a un sous-préschéma $Z$ de $X$, dont l'espace de base est $\psi(Y)$ et dont la restriction à chaque $U \cap \psi(Y)$ est $Z_U$ ; les $\varepsilon_U$ sont alors les restrictions aux $\psi^{-1}(U)$ d'un isomorphisme $g$ de $Y$ sur $Z$, tel que $f = j \circ g$, où $j$ est l'injection $Z \to X$. C.Q.F.D.

Corollaire 1. - Soit $(\psi, \theta)$ un morphisme $Y \to X$, $(V_\lambda)$ un recouvrement de $\psi(Y)$ par des ouverts de $X$ (resp. un recouvrement ouvert de $X$). Pour que $(\psi, \theta)$ soit une immersion (resp. une immersion fermée) il faut et il suffit que sa restriction à chacun des préschémas induits par $Y$ sur les $\psi^{-1}(V_\lambda)$ soit une immersion (resp. une immersion fermée).

Comme $\theta_y^\sharp$ est alors surjectif pour tout $y \in Y$, tout revient à vérifier que $\psi$ est un homéomorphisme de $Y$ sur une partie localement fermée (resp. fermée) $\psi(Y)$ de $X$. Or $\psi$ est évidemment injectif et transforme tout voisinage de $y$ dans $Y$ en un voisinage de $\psi(y)$ dans $\psi(Y)$ en vertu de l'hypothèse ; d'autre part, $\psi(Y) \cap V_\lambda$ est localement fermé (resp. fermé) dans $V_\lambda$, donc $\psi(Y)$ est localement fermé (resp. fermé) dans $X$.

Corollaire 2. - Soit $(X, \mathcal{O}_X)$ un schéma affine. Pour que $(\psi, \theta)$ soit

\newpage

%% ===== Page 68 =====

- 48 -

une immersion fermée de $(Y, \mathcal{O}_Y)$ dans $(X, \mathcal{O}_X)$, il faut et il suffit que $(Y, \mathcal{O}_Y)$ soit un schéma affine, et que l'homomorphisme $\Gamma(\psi) : \Gamma(\mathcal{O}_X) \to \Gamma(\mathcal{O}_Y)$ soit surjectif.

Corollaire 3. — Le composé de deux immersions (resp. de deux immersions fermées) est une immersion (resp. une immersion fermée).

Cela résulte aussitôt de la prop. 4, compte tenu du fait qu'une partie localement fermée d'une partie localement fermée est localement fermée.

Proposition 5. — Soient $\alpha : X' \to X$, $\beta : Y' \to Y$ deux $S$-morphismes : si $\alpha$ et $\beta$ sont des immersions (resp. des immersions fermées), $\alpha \times_S \beta$ est une immersion (resp. une immersion fermée). En outre, si $\alpha$ (resp. $\beta$) identifie $X'$ (resp. $Y'$) à un sous-schéma $X''$ (resp. $Y''$) de $X$ (resp. $Y$), $\alpha \times_S \beta$ identifie $X' \times_S Y'$ au sous-préschéma $p^{-1}(X'') \cap q^{-1}(Y'')$ de $X \times_S Y$, en désignant par $p$ et $q$ les projections de $X \times_S Y$ dans $X$ et $Y$ respectivement.

[text redacted/crossed out: compte tenu de la déf. 2,] Compte tenu de la déf. 2, on peut se restreindre au cas où $X'$ et $Y'$ sont des sous-préschémas, et $\alpha$ et $\beta$ les morphismes d'injection. La proposition a déjà été établie pour les sous-préschémas induits sur les ouverts (§ 2, n°4, cor. du th. 1) ; comme tout sous-préschéma est un sous-préschéma fermé d'un préschéma induit sur un ouvert, on est ramené au cas où $X'$ et $Y'$ sont des sous-préschémas fermés.

Montrons d'abord qu'on peut supposer que $S$ est affine. En effet, soit $(S_\lambda)$ un recouvrement de $S$ formé d'ouverts affines ; si $\psi$ et $\theta$ sont les morphismes structuraux de $X$ et $Y$, soit $X_\lambda = \psi^{-1}(S_\lambda)$, $Y_\lambda = \theta^{-1}(S_\lambda)$. La restriction $X'_\lambda$ (resp. $Y'_\lambda$) de $X'$ (resp. $Y'$) à $X_\lambda$ (resp. $Y_\lambda$) est un sous-préschéma fermé de $X_\lambda$ (resp. $Y_\lambda$), les préschémas $X_\lambda, Y_\lambda, X'_\lambda, Y'_\lambda$ peuvent être considérés comme des $S_\lambda$-préschémas, et les produits $X_\lambda \times_{S_\lambda} Y_\lambda$ et $X'_\lambda \times_{S_\lambda} Y'_\lambda$ (resp. $X'_\lambda \times_{S_\lambda} Y_\lambda$ et $X'_\lambda \times_S Y'_\lambda$) sont identiques (§ 2, n°4, prop. 7). Si la proposition est vraie lorsque $S$ est affine,

\newpage

%% ===== Page 69 =====

- 49 -
La restriction de $\alpha \times_S \beta$ à chacun des $X'_\lambda \times_S Y'_\mu$ sera donc une immersion ($\S 2, n^\circ 4, \text{cor. du th. 1}$). Comme le produit $X'_\lambda \times_S Y'_\mu$ s'identifie à $(X'_\lambda \cap X'_\lambda) \times_S (Y'_\mu \cap Y'_\mu)$ ($\S 2, n^\circ 4, \text{partie b) du th. 1}$), la restriction de $\alpha \times_S \beta$ à chacun des $X'_\lambda \times_S Y'_\mu$ est encore une immersion ; il en est par suite de même de $\alpha \times_S \beta$ en vertu du cor. 1 de la prop. 4.

En second lieu, montrons qu'on peut aussi supposer que $X$ et $Y$ sont affines. En effet, soit $(U_i)$ (resp. $(V_j)$) un recouvrement de $X$ (resp. $Y$) par des ouverts affines, et soit $X'_i$ (resp. $Y'_j$) la restriction de $X'$ (resp. $Y'$) à $U_i$ (resp. $V_j$), qui est un sous-préschéma fermé de $U_i$ (resp. $V_j$). En outre, $U_i \times_S V_j$ s'identifie à la restriction de $X \times_S Y$ à $p^{-1}(U_i) \cap q^{-1}(V_j)$ ($\S 2, \text{cor. du th. 1}$) ; et de même, si $p', q'$ sont les projections de $X' \times_S Y'$, $X'_i \times_S Y'_j$ s'identifie à la restriction de $X' \times_S Y'$ à $p'^{-1}(U_i) \cap q'^{-1}(V_j)$. Posons $\gamma = \alpha \times_S \beta$ ; on a par définition $p \circ \gamma = \alpha \circ p'$; comme $X'_i = \alpha^{-1}(U_i)$, $Y'_j = \beta^{-1}(V_j)$, on a aussi $p'^{-1}(X'_i) = \gamma^{-1}(p^{-1}(U_i))$, $q'^{-1}(Y'_j) = \gamma^{-1}(q^{-1}(V_j))$, d'où $p'^{-1}(X'_i) \cap q'^{-1}(Y'_j) = \gamma^{-1}(p^{-1}(U_i) \cap q^{-1}(V_j)) = \gamma^{-1}(U_i \times_S V_j)$ ; on conclut comme dans la première partie du raisonnement.

Supposons donc $X, Y, S$ affines, et soient $B, C, A$ leurs anneaux respectifs. Alors $B$ et $C$ sont des $A$-algèbres, $X'$ et $Y'$ des schémas dont les anneaux sont des algèbres quotients $B', C'$ de $B$ et $C$ respectivement. En outre, on a $\alpha = (\rho, \tilde{\rho})$, $\beta = (\sigma, \tilde{\sigma})$, où $\rho$ et $\sigma$ sont respectivement les homomorphismes canoniques $B \to B'$, $C \to C'$ ($\S 1, n^\circ 7, \text{th. 4}$). Cela étant, on sait que $X \times_S Y$ est un schéma affine, dont l'anneau est $B \otimes_A C$ ; de même, $X' \times_S Y'$ est un schéma affine dont l'anneau est $B' \otimes_A C'$, et $\alpha \times_S \beta = (\tau, \tilde{\tau})$, où $\tau$ est l'homomorphisme de $B \otimes_A C$ dans $B' \otimes_A C'$ ; comme cet homomorphisme est surjectif, $\alpha \times_S \beta$ est bien une immersion. En outre, si $I$ (resp. $J$) est le noyau de $\rho$ (resp. $\sigma$), le noyau de $\tau$ est $I \otimes_A C + B \otimes_A J$, et en remontant aux définitions ($\S 1$), on voit aussitôt que dans le spectre premier de

\newpage

%% ===== Page 70 =====

$B \otimes_A C$, cet idéal correspond à l'ensemble fermé $\bar{\rho}(X') \cap \bar{\sigma}(Y')$ (en se rappelant que $\rho = (\alpha, \tilde{\alpha})$ et $\sigma = (\beta, \tilde{\beta})$, où $\alpha$ (resp. $\beta$) est l'homomorphisme $b \to b \otimes 1$ (resp. $c \to 1 \otimes c$)).

\emph{Corollaire 1.} — Si $f : X \to Y$ est une immersion (resp. une immersion fermée) et un $S$-morphisme, $f_{S'}$ est une immersion (resp. une immersion fermée) pour toute extension $S' \to S$ du préschéma de base.

\emph{Corollaire 2.} — Deux sous-préschémas quelconques $Y', Y''$ de $X$ admettent une borne inférieure $Y' \cap Y''$ pour la relation d'ordre entre préschémas ; $Y$ est canoniquement isomorphe à $Y' \times_X Y''$ et est fermé si $Y'$ et $Y''$ le sont.

Soient $j' : Y' \to X$, $j'' : Y'' \to X$ les morphismes d'injection, qui font de $Y'$ et $Y''$ des $X$-préschémas ; en vertu de la prop. 5, $j' \times_X j''$ est une immersion (resp. une immersion fermée si $Y'$ et $Y''$ sont fermés) de $Y' \times_X Y''$ dans $X \times_X X = X$. En outre, si $Z$ est un sous-préschéma de $X$ tel que $Z \subseteq Y'$ et $Z \subseteq Y''$, l'injection $Z \to X$ se factorise en $Z \xrightarrow{h'} Y'$ et $Z \xrightarrow{h''} Y''$, donc $h'$ et $h''$ sont des $X$-morphismes, et il y a par suite un $X$-morphisme $g : Z \to Y' \times_X Y''$ tel que $h' = p' \circ g$, $h'' = p'' \circ g$, en désignant par $p'$ et $p''$ les projections. On en conclut que si $Y$ est le sous-préschéma de $X$ image canonique de $Y' \times_X Y''$, on a $Z \subseteq Y$, ce qui démontre le corollaire.

\emph{Corollaire 3.} — Soit $f : X \to Y$ un morphisme, $Y'$ un sous-préschéma (resp. un sous-préschéma fermé) de $Y$. Il existe alors un sous-préschéma (resp. un sous-préschéma fermé) de $X$ dont l'espace de base est l'image réciproque $f^{-1}(Y')$ et qui est canoniquement isomorphe à $X \times_Y Y'$.

En effet, si $j : Y' \to Y$ est le morphisme d'injection, $1_X \times_Y j$ est une immersion (resp. une immersion fermée) en vertu de la prop. 5 ; et on a déjà vu (§ 2, n$^\circ$7, prop. 16) que l'image de $X \times_Y Y'$ par cette immersion a un espace de base qui s'identifie canoniquement à $f^{-1}(Y')$. Quand on parlera de $f^{-1}(Y')$ comme d'un sous-préschéma de $X$, il s'agira toujours...

\newpage

%% ===== Page 71 =====

- 51 -

3. Pr\text{é}sch\text{é}mas r\text{é}duits.

Proposition 6. -- Soit $(X, \mathcal{O}_X)$ un pr\text{é}sch\text{é}ma, et pour tout $x \in X$, soit $N(x)$ l'id\text{é}al de $\mathcal{O}_X(x)$ form\text{é} des \text{é}l\text{é}ments nilpotents. Les $N(x)$ sont les fibres d'un faisceau quasi-coh\text{é}rent d'id\text{é}aux $N$. Lorsque $X$ est un sch\text{é}ma affine, on a $N = \tilde{N}$, o\text{ù} $N$ est l'id\text{é}al des \text{é}l\text{é}ments nilpotents de l'anneau $A = \mathcal{O}_X(X)$.

La question \text{é}tant locale, on est ramen\text{é} \`a d\text{é}montrer la derni\text{e}re assertion. On sait que $\tilde{N}$ est un faisceau quasi-coh\text{é}rent, et que sa fibre au point $x$ est l'id\text{é}al $N_x$ de l'anneau de fractions $\mathcal{O}_{X,x}$ ; tout revient \`a prouver que $N(x) \subset N_x$, l'inclusion oppos\text{é}e \text{é}tant \text{é}vidente. Or, soit $z/s$ un \text{é}l\text{é}ment de $N_x$, o\text{ù} $z \in A, s \notin j_x$ ; par hypoth\text{e}se il existe $k$ tel que $(z/s)^k = 0$, ce qui signifie qu'il existe $t \notin j_x$ tel que $tz^k = 0$. On en conclut $(tz)^k = 0$, et par suite $z/s = (tz)/(ts)$ appartient \`a $N_x$.

D\text{é}finition 3. -- On dit qu'un pr\text{é}sch\text{é}ma $(X, \mathcal{O}_X)$ est r\text{é}duit si aucun des anneaux locaux $\mathcal{O}_{X,x}$ ne contient d'\text{é}l\text{é}ments nilpotents $\neq 0$. Le sous-pr\text{é}sch\text{é}ma quasi-coh\text{é}rent associ\text{é} au faisceau d'id\text{é}aux $N$ est appel\text{é} le pr\text{é}sch\text{é}ma r\text{é}duit associ\text{é} \`a $X$ et not\text{é} aussi $X_{\text{red}}$.

Dire que $X$ est r\text{é}duit signifie donc que $X_{\text{red}} = X$.

Corollaire 1. -- Les espaces de base de $X$ et $X_{\text{red}}$ sont identiques.
Cela r\text{é}sulte de ce que pour tout $x \in X$, l'\text{é}l\text{é}ment unit\text{é} de $\mathcal{O}_x$ est $\neq 0$ donc $N(x) \neq \mathcal{O}_x(x)$.

Proposition 7. -- Pour tout sous-espace localement ferm\text{é} $Y$ de $X$, il existe un sous-pr\text{é}sch\text{é}ma r\text{é}duit et un seul de $X$ ayant $Y$ pour espace de base ; c'est le plus petit des sous-pr\text{é}sch\text{é}mas de $X$ ayant $Y$ pour espace de base.

Si $X$ est affine d'anneau $A$, et $Y$ ferm\text{é} dans $X$, la proposition est imm\text{é}diate : $j(Y)$ est le plus grand id\text{é}al $a$ tel que $V(a) = Y$, et il est \text{é}gal \`a sa racine ($\S 1, n^\circ 1, \text{prop.} 2$), donc $A/j(Y)$ n'a pas d'\text{é}l\text{é}ments...

\newpage

%% ===== Page 72 =====

- 52 -

ment nilpotent $\neq 0$ et pour tout idéal $\mathfrak{a} \subsetneq \mathfrak{j}(Y)$ et tel que $V(\mathfrak{a}) = Y$, $V/\mathfrak{a}$ possède de tels éléments.

Dans le cas général, pour tout ouvert affine $U \subset X$ tel que $U \cap Y$ soit fermé dans $U$, considérons le sous-préschéma fermé $Y_U$ de $U$ défini par le faisceau d'idéaux associé à l'idéal $\mathfrak{j}(U \cap Y)$ de $A(U)$ ; c'est le plus petit sous-préschéma de $U$ ayant $U \cap Y$ comme espace de base. En outre, si $V$ est un ouvert affine contenu dans $U$, $Y_V$ est induit par $Y_U$ sur $V$, car il suffit de le vérifier pour $V=D(f)$, où $f \in A(U)$ ; posons $A=A(U)$, et soit l'homomorphisme canonique $A \to A/\mathfrak{j}(Y \cap U) = B$ ; il lui correspond une surjection $\varphi : A_f \to B_{\varphi(f)}$, et tout revient à voir que le noyau de $\varphi$ est l'idéal $\mathfrak{j}(Y \cap V)$. Or, $\varphi(z/f^n) = 0$, cela signifie que $\varphi(z)/(\varphi(f))^n = 0$ dans $B_{\varphi(f)}$ ; et par suite $\varphi(f^k z) = 0$, ou encore $f^k z \in \mathfrak{j}(Y \cap U)$ pour un entier $k$ ; mais cette relation signifie que $f \in \mathfrak{j}_y$ ou $z \in \mathfrak{j}_y$ pour tout $y \in U$, et comme les $y$ tels que $f \notin \mathfrak{j}_y$ sont précisément les points $y \in V \cap D(f)$, $\varphi(f^k z) = 0$ équivaut finalement à $z/f^n \in \mathfrak{j}(Y \cap V)$. Les $Y_U$ sont donc les restrictions à $U$ d'un préschéma $(X, \mathcal{O}_X)$ ayant $Y$ pour espace de base, puisque les $U \cap Y$, où $U$ parcourt les ouverts affines de $X$ tels que $Y \cap U$ soit fermé dans $U$, forment une base de la topologie de $Y$. Si $Z$ est un autre sous-préschéma ayant $Y$ comme espace de base, on a nécessairement $Z_U \supset Y_U$ pour un $U$ moins, et par suite $Z_U$ n'est pas réduit, donc il en est de même de $Z$.

Soit $f=(f, \theta) : X \to Y$ un morphisme de préschémas ; l'homomorphisme $\theta_x : \mathcal{O}_Y(f(x)) \to \mathcal{O}_X(x)$ applique tout élément nilpotent de $\mathcal{O}_Y(f(x))$ sur un élément nilpotent de $\mathcal{O}_X(x)$ ; par passage aux quotients, on déduit donc $\theta_x$ un homomorphisme $\omega : f^*(\mathcal{N}_Y) \to \mathcal{N}_X$ ; et par suite $(f, \omega^{\#})$ est un morphisme $X_{red} \to Y_{red}$ que nous noterons $f_{red}$ et appellerons le morphisme réduit de $f$. Il est immédiat que pour deux morphismes $f, g$, on a $(f \circ g)_{red} = f_{red} \circ g_{red}$ ; on voit ainsi que $X \to X_{red}$ est un foncteur covariant.

\newpage

%% ===== Page 73 =====

riant . En outre, le diagramme
$$\begin{tikzcd}
X_{\text{red}} \arrow[r, "f_{\text{red}}"] \arrow[d] & Y_{\text{red}} \arrow[d] \\
X \arrow[r, "f"] & Y
\end{tikzcd}$$
est commutatif, les flèches verticales étant les injections . On notera en particulier que si $X$ est réduit, tout morphisme $f : X \to Y$ se factorise en $X \xrightarrow{f_{\text{red}}} Y_{\text{red}} \to Y$ ; en d'autres termes, ce morphisme est majoré par le morphisme d'injection $Y_{\text{red}} \to Y$.

Proposition 8. — Soit $f$ un morphisme de préschémas ; si $f$ est surjectif (resp. géométriquement injectif, une immersion, une immersion fermée), il en est de même de $f_{\text{red}}$.

La proposition est triviale si $f$ est surjectif ; si $f$ est géométriquement injectif, la proposition résulte de ce que pour tout $x \in X$, le corps $k(x)$ est le même pour le préschéma $X$ et le préschéma $X_{\text{red}}$ (cf. § 2, n°7, prop. 13) ; enfin, si $f=(\varphi, \theta)$ est une immersion (resp. une immersion fermée), la proposition résulte de ce que si $\theta_x^{\#}$ est surjectif, il en est de même de l'homomorphisme obtenu en passant aux quotients par les racines de $0$ dans $\mathcal{O}_{\varphi(x)}$ et $\mathcal{O}_x$ (cf. n°2, prop. 4).

Proposition 9. — Si $X$ et $Y$ sont des $S$-préschémas, les préschémas $X_{\text{red}} \times_{S_{\text{red}}} Y_{\text{red}}$ et $X_{\text{red}} \times_S Y_{\text{red}}$ sont identiques, et s'identifient canoniquement à un sous-préschéma de $X \times_S Y$ ayant même espace de base que ce produit.

L'identification canonique de $X_{\text{red}} \times_{S_{\text{red}}} Y_{\text{red}}$ à un sous-préschéma de $X \times_S Y$ ayant même espace de base résulte de la prop. 5 du n°2, compte tenu du cor. de la prop. 6. D'autre part, si $\varphi$ et $\psi$ sont les morphismes structuraux $X_{\text{red}} \to S$, $Y_{\text{red}} \to S$, ils se factorisent par $S_{\text{red}}$, et si $f : T \to X_{\text{red}}$ ; $g : T \to Y_{\text{red}}$ sont des $S$-morphismes, on a $\varphi \circ f = \psi \circ g$ donc (puisque $S_{\text{red}} \to S$ est un monomorphisme) $\varphi_{\text{red}} \circ f = \psi_{\text{red}} \circ g$ ; $T$ peut donc être considéré comme un $S_{\text{red}}$-préschéma et $f$ et $g$ des $S_{\text{red}}$-morphismes. La réciproque étant évidente, on en déduit l'identité.

\newpage

%% ===== Page 74 =====

- 54 -

\begin{itemize}
\item des produits $X_{\mathrm{red}} \times_{S_{\mathrm{red}}} Y_{\mathrm{red}}$ et $X_{\mathrm{red}} \times_S Y_{\mathrm{red}}$ en vertu de la définition de ces derniers.
\end{itemize}

On notera que si $X$ et $Y$ sont réduits, il n'en est pas nécessairement de même de $X \times_S Y$, car le produit tensoriel de deux algèbres sans élément nilpotent peut avoir de tels éléments.

\section*{4. Diagonale et graphe d'un morphisme}

Soit $X$ un $S$-préschéma ; on appelle \emph{morphisme diagonal} de $X$ dans $X \times_S X$ et on note $\Delta_X$ ou $\Delta$ le $S$-morphisme $(1_X, 1_X)_S$, autrement dit l'unique $S$-morphisme tel que
$$(1) \quad p_1 \circ \Delta_X = p_2 \circ \Delta_X = 1_X$$
en désignant par $p_1, p_2$ les projections de $X \times_S X$.

Proposition 10. --- Le \emph{morphisme diagonal} $\Delta_X$ est une immersion de $X$ dans $X \times_S X$.

En effet, les applications continues $p_1$ et $\Delta_X$ des espaces de base étant telles que $p_1 \circ \Delta_X$ soit l'identité, $\Delta_X$ est un homéomorphisme de $X$ sur $\Delta_X(X)$. De même, l'homomorphisme composé $\mathcal{O}_{X \times_S X} \to \Delta_X(\mathcal{O}_{X \times_S X}) \to \mathcal{O}_X$ des homomorphismes correspondant à $p_1$ et $\Delta_X$ respectivement étant l'identité, l'homomorphisme correspondant à $\Delta_X$ est surjectif ; la proposition résulte donc de la prop. 4 du n$^\circ$ 2.

Le sous-préschéma de $X \times_S X$ associé à l'immersion $\Delta_X$ est encore appelé la \emph{diagonale} de $X \times_S X$.

Proposition 11. --- Soient $X, Y$ deux $S$-préschémas ; si on identifie canoniquement les produits $(X \times X) \times (X \times Y)$ à $(X \times X) \times (X \times Y)$, le morphisme $\Delta_{X \times Y}$ s'identifie à $\Delta_X \times \Delta_Y$.

En effet, si $p_1, q_1$ sont les premières projections $X \times X \to X, Y \times Y \to Y$, la première projection $(X \times X) \times (X \times Y) \to X \times Y$ s'identifie à $p_1 \times q_1$, et on a $(p_1 \times q_1) \circ (\Delta_X \times \Delta_Y) = (p_1 \circ \Delta_X) \times (q_1 \circ \Delta_Y) = 1_{X \times Y}$.

Corollaire. --- Pour toute extension $S' \to S$ du préschéma de base, $\Delta_{X^{S'}}$ s'identifie canoniquement à $(\Delta_X)^{S'}$.

Il suffit de remarquer que $(X \times_S X)^{S'}$ s'identifie canoniquement à $X^{S'} \times_{S'} X^{S'}$.

\newpage

%% ===== Page 75 =====

$X^{S'} \times_{S'} X^{S'}$ (§ 2, n°5, cor. 1 de la prop. 9).

Proposition 12. — Si $f: X \to Y$ est un $S$-morphisme, le diagramme

$$
\begin{tikzcd}
& X \arrow[d, "(1_X, f)"] \arrow[dr, "\Delta_X"] & \\
Y \arrow[r, "\Delta_Y"] & Y \times_S Y \arrow[r, "f \times_S 1_Y"] & X \times_S Y \arrow[r, "1_X \times_S f"] & X \times_S X \arrow[d, "f \times_S f"] \\
& & & Y \times_S Y
\end{tikzcd}
$$

est commutatif. En outre $X$ s'identifie au produit des $(Y \times_S Y)$-préschémas $Y$ et $X \times_S Y$, les projections s'identifiant à $f$ et $(1_X, f)_S$.

La vérification de la commutativité du diagramme est triviale, compte tenu de la définition du produit et de (1) ; nous la laissons au lecteur. Désignons par $p_1, p_2$ les projections de $Y \times_S Y$, par $p, q$ celles de $X \times_S Y$. Soient $T$ un $S$-préschéma, $g: T \to Y$, $h: T \to X \times_S Y$ deux $S$-morphismes tels que $\Delta_Y \circ g = (f \times_S 1_Y) \circ h$. On déduit de là $g = p_1 \circ \Delta_Y \circ g = p_1 \circ (f \times_S 1_Y) \circ h = f \circ h$, et $g = p_2 \circ \Delta_Y \circ g = p_2 \circ (f \times_S 1_Y) \circ h = q \circ h$; en outre $p \circ (1_X, f)_S \circ h = p \circ h = p \circ (f \times_S 1_Y) \circ h = g = q \circ h$. En outre, si $u_1, u_2$ sont deux $S$-morphismes de $T$ dans $X$ tels que $(1_X; f)_S \circ u_1 = (1_X, f)_S \circ u_2$, on déduit de $p \circ (1_X, f)_S = 1_X$ que $u_1 = u_2$, ce qui achève de prouver la proposition.

Corollaire. — Si $X$ est un sous-préschéma de $Y$, la diagonale $\Delta_X(X)$ s'identifie à un sous-préschéma de $\Delta_Y(Y)$, dont l'espace de base s'identifie à $\Delta_Y(Y) \cap p_1^{-1}(X) = \Delta_Y(Y) \cap p_2^{-1}(X)$.

Il suffit d'appliquer la prop. 12 au cas où $f$ est l'injection $X \to Y$ ; on sait alors que $f \times_S f$ est une immersion, identifiant l'espace de base de $X \times_S X$ à $p_1^{-1}(X) \cap p_2^{-1}(X)$ (n°2, prop. 5) ; en outre, si $z \in \Delta_Y(Y) \cap p_1^{-1}(X)$, on a $z = \Delta_Y(y)$ et $y = p(z) \in X$, donc $y = f(y)$, et $z = \Delta_Y(f(y))$ appartient à $\Delta_X(X)$ en vertu de la commutativité du diagramme (2).

Soient $X, Y$ deux $S$-préschémas, $\varphi: S \to T$ un morphisme, qui fait de tout $S$-préschéma un $T$-préschéma. Soient $f: X \to S$, $g: Y \to S$ les

\newpage

%% ===== Page 76 =====

- 56 -

\begin{itemize}
    \item morphismes structuraux, et $p, q$ les projections de $X \times_S Y$, $\pi = f \circ p = g \circ q$ le morphisme structural $X \times_S Y \to S$.
\end{itemize}

\textbf{Proposition 13.} — Le diagramme
$$
\begin{tikzcd}
X \times_S Y \arrow[r, "{(p, q)_T}"] \arrow[d, "\pi"'] & X \times_T Y \arrow[d, "f \times_T g"] \\
S \arrow[r, "\Delta_S"'] & S \times_S S
\end{tikzcd}
$$
est commutatif. En outre, $X \times_S Y$ s'identifie au produit des $T$-préschémas $S$ et $X \times_T Y$, les projections s'identifiant à $\pi$ et $(p, q)_T$.

Nous laissons de nouveau au lecteur la vérification de la commutativité de (3). Soient $Z$ un préschéma, $u: Z \to S$ et $v: Z \to X \times_T Y$ deux morphismes tels que $\Delta_S \circ u = (f \times_T g) \circ v$. Soient $p_1, p_2$ les projections de $S \times_S S$, $p', q'$ celles de $X \times_T Y$. On a
$$ u = p_1 \circ \Delta_S \circ u = p_1 \circ (f \times_T g) \circ v = f \circ p' \circ v $$
$$ u = p_2 \circ \Delta_S \circ u = p_2 \circ (f \times_T g) \circ v = g \circ q' \circ v $$

Vu la définition du produit $X \times_S Y$, ces deux égalités entraînent qu'il existe un morphisme $w: Z \to X \times_S Y$ tel que $p' \circ v = p \circ w$ et $q' \circ v = q \circ w$, d'où $u = (f \circ p) \circ w = \pi \circ w$. D'autre part, on a $p = p' \circ (p, q)_T$ et $q = q' \circ (p, q)_T$, donc $p' \circ v = p' \circ (p, q)_T \circ w$ et $q' \circ v = q' \circ (p, q)_T \circ w$, donc $v = (p, q)_T \circ w$. Enfin, si $w_1, w_2$ sont deux morphismes de $Z$ dans $X \times_S Y$ tels que $(p, q)_T \circ w_1 = (p, q)_T \circ w_2$, les relations $p' \circ (p, q)_T = p, q' \circ (p, q)_T = q$ entraînent que $p \circ w_1 = p \circ w_2$ et $q \circ w_1 = q \circ w_2$, donc $w_1 = w_2$, ce qui achève la démonstration.

\textbf{Corollaire 1.} — Le morphisme $(p, q)_T$ est une immersion, s'identifiant à $1_{X \times_S Y} \times \Delta_S$ (où on pose $P = S \times_S S$).
Cela résulte de la prop. 13 et du \marginpar{de la prop 5 du n°2, 2, n°5, cor. de la prop. 8}

\textbf{Corollaire 2.} — Si $X, Y$ sont deux $S$-préschémas, $f: X \to Y$ un $S$-morphisme; alors $(1_X, f)_S$ est une immersion de $X$ dans $X \times_S Y$, appelée le morphisme graphe de $f$, tel que, si $p, q$ sont les projections de $X \times_S Y$, on ait $f \circ q \circ \Gamma_f = p \circ \Gamma_f$ et $1_X = p \circ \Gamma_f$.

En effet, si on considère $X$ comme un $Y$-préschéma au moyen de $f$,

\newpage

%% ===== Page 77 =====

- 57 -

on sait que $(1_X, f)_Y$ est un isomorphisme de $X$ sur $X \times_Y Y$ ; d'autre part, il est immédiat que, si $p, q$ désignent les projections de $X \times_Y Y$, $(1_X, f)_S$ se factorise en
$$ X \xrightarrow{(1_X, f)_Y} X \times_Y Y \xrightarrow{(p, q)_S} X \times_S Y $$ ;
c'est donc une immersion en vertu du cor. 1. La vérification de la seconde assertion est immédiate.

Le sous-préschéma de $X \times_S Y$ associé à l'immersion $\Gamma_f$ est aussi appelé le \emph{graphe} du morphisme $f$.

5. \emph{Morphismes et préschémas séparés.}

Définition 4. -- On dit qu'un morphisme de préschémas $f : X \to Y$ est \emph{séparé} si le morphisme diagonal $X \to X \times_Y X$ est une immersion fermée ; on dit aussi alors que $X$ est un préschéma séparé au-dessus de $Y$. On dit qu'un préschéma $X$ est séparé s'il est séparé au-dessus de $Z$ ; on dit aussi alors que $X$ est un schéma.

Proposition 14. -- Tout \emph{morphisme} de préschémas $f : X \to Y$ (et en particulier tout \emph{morphisme} immersion) est séparé.

En effet (n° 4, Remarque), le préschéma $X \times_Y X$ s'identifie alors à $X$ et à $\Delta_X(X)$.

Proposition 15. -- Si $f : X \to Y$ et $g : Y \to Z$ sont séparés, $g \circ f$ est séparé.

Désignons en effet par $\Delta_{X|Y}$ et $\Delta_{X|Z}$ les applications diagonales de $X$ dans $X \times_Y X$ et $X \times_Z X$ respectivement ($X$ étant considéré comme $Y$-préschéma par $f$ et comme $Z$-préschéma par $g \circ f$). On vérifie immédiatement que le diagramme
$$
\begin{tikzcd}
X \arrow[d, "\Delta_{X|Z}"'] \arrow[r, "\Delta_{X|Y}"] & X \times_Y X \arrow[d, "j"] \\
X \times_Z X \arrow[r, phantom, ""{coordinate, name=Y}] & X \times_Z X
\end{tikzcd}
$$
est commutatif, $j$ désignant l'immersion canonique (n° 4). Comme le composé de deux immersions fermées est une immersion fermée, la proposition résulte donc du lemme suivant :

\newpage

%% ===== Page 78 =====

\noindent Lemme 1. --- Soit $S \to T$ un morphisme séparé. Si $X$ et $Y$ sont deux $S$-préschémas, l'immersion canonique $X \times_S Y \to X \times_T Y$ est fermée.

En effet, il résulte du n°4, cor. 1 de la prop. 13 que cette immersion s'identifie à un produit de deux immersions fermées, et le lemme est donc conséquence de la prop. 5 du n°2.

Corollaire. --- Si $f : X \to Y$ est séparé, la restriction de $f$ à tout sous-préschéma de $X$ est séparée.

C'est une conséquence immédiate des prop. 14 et 15.

Proposition 16. --- Si $f : X \to Y$ est un $S$-morphisme séparé, le $S'$-morphisme $f^{S'} : X^{S'} \to Y^{S'}$ est séparé pour toute extension $S' \to S$ du préschéma de base.

En effet ($\S$ 2, n°5, cor. 2 de la prop. 9), $X^{S'} \times_{Y^{S'}} X^{S'}$ s'identifie canoniquement à $(X \times_Y X) \times_S S'$, et on vérifie immédiatement que le morphisme diagonal $\Delta_{S'}$ s'identifie alors à $\Delta \times 1_{S'}$ ; la proposition résulte donc de la prop. 5 du n°2.

Corollaire. --- Si $f : X \to X'$ et $g : Y \to Y'$ sont deux $S$-morphismes séparés, $f \times_S g$ est séparé.

En effet, $f \times_S g$ se factorise en
$$ X \times_S Y \xrightarrow{f \times 1_Y} X' \times_S Y \xrightarrow{1_{X'} \times g} X' \times_S Y' $$
et chacun des deux facteurs est séparé par la prop. 16 ; on conclut à l'aide de la prop. 15.

Proposition 17. --- Soit $(Y_\lambda)$ un recouvrement ouvert d'un préschéma $Y$ ; pour qu'un morphisme $f : X \to Y$ soit séparé, il faut et il suffit que chacune de ses restrictions $f^{-1}(Y_\lambda) \to Y_\lambda$ soit séparée.

Si on pose $X_\lambda = f^{-1}(Y_\lambda)$, la nécessité de la condition résulte du cor. de la prop. 15 et du fait que les produits $X_\lambda \times_{Y_\lambda} X_\lambda$ et $X_\lambda \times_Y X_\lambda$ sont identiques ($\S$ 2, n°4, prop. 7). Inversement, supposons que chaque $X_\lambda$ soit séparé au-dessus de $Y_\lambda$ ; comme les $X_\lambda \times_Y X_\lambda$ forment un recouvrement ouvert de $X \times_Y X$, il s'agit de montrer que l'intersection de la diagonale $\Delta(X)$ et de $X_\lambda \times_Y X_\lambda$ est fermée dans ce dernier espace.

\newpage

%% ===== Page 79 =====

Lenne 1. — Soit $S \to T$ un morphisme séparé. Si $X$ et $Y$ sont deux $S$-préschémas, l'immersion canonique $X \times_S Y \to X \times_T Y$ est fermée.

En effet, il résulte du n°4, cor.1 de la prop.13 que cette immersion s'identifie à un produit de deux immersions fermées, et le lemme est donc conséquence de la prop.5 du n°2.

Corollaire. — Si $f: X \to Y$ est séparé, la restriction de $f$ à tout sous-préschéma de $X$ est séparée.

C'est une conséquence immédiate des prop.14 et 15.

Proposition 16. — Si $f: X \to Y$ est un $S$-morphisme séparé, le $S'$-morphisme $f_{S'}: X_{S'} \to Y_{S'}$ est séparé pour toute extension $S' \to S$ du préschéma de base.

En effet (§ 2, n°5, cor.2 de la prop.9), $X_{S'} \times_{Y_{S'}} X_{S'}$ s'identifie canoniquement à $(X \times_Y X)_{S'} = (X \times_Y X) \times_S S'$, et on vérifie immédiatement que le morphisme diagonal $\Delta_{S'}$ s'identifie alors à $\Delta_{S} \times_S \text{id}_{S'}$; la proposition résulte donc de la prop.5 du n°2.

Corollaire. — Si $f: X \to X'$ et $g: Y \to Y'$ sont deux $S$-morphismes séparés, $f \times_S g$ est séparé.

En effet, $f \times_S g$ se factorise en
$$ X \times_S Y \xrightarrow{f \times_S \text{id}_Y} X' \times_S Y \xrightarrow{\text{id}_{X'} \times_S g} X' \times_S Y' $$
et chacun des deux facteurs est séparé par la prop.16 ; on conclut à l'aide de la prop.15.

Proposition 17. — Soit $(Y_\lambda)$ un recouvrement ouvert d'un préschéma $Y$ ; pour qu'un morphisme $f: X \to Y$ soit séparé, il faut et il suffit que chacune de ses restrictions $f^{-1}(Y_\lambda) \to Y_\lambda$ soit séparée.

Si on pose $X_\lambda = f^{-1}(Y_\lambda)$, la nécessité de la condition résulte du cor. de la prop.15 et du fait que les produits $X_\lambda \times_Y X_\lambda$ et $X_\lambda \times_{Y_\lambda} X_\lambda$ sont identiques (§ 2, n°4, prop.7). Inversement, supposons que chaque $X_\lambda$ soit séparé au-dessus de $Y_\lambda$ ; comme les $X_\lambda \times_Y X_\lambda$ forment un recouvrement ouvert de $X \times_Y X$, il s'agit de montrer que l'intersection de la diagonale $\Delta_X(X)$ et de $X_\lambda \times_Y X_\lambda$ est fermée dans ce dernier espace.

\newpage

%% ===== Page 80 =====

séparé, il faut et il suffit que $X$ soit séparé (autrement dit : que $X$ soit un schéma).

On observera en effet que le critère de la prop. 18 ne dépend pas du schéma affine $Y$ considéré.

Corollaire 3. — Pour qu'un morphisme $f : X \to Y$ soit séparé, il faut que pour tout ensemble ouvert séparé $U \subset Y$ (i.e. tel que le préschéma induit soit séparé), $f^{-1}(U)$ soit séparé, et il suffit qu'il en soit ainsi pour tout ouvert affine $U \subset Y$.

La nécessité de la condition résulte des prop. 17 et 15; la suffisance résulte de la prop. 17 et du cor. 2 de la prop. 18, compte tenu de l'existence de recouvrements ouverts affines de $Y$.

Corollaire 4. — Soient $f : X \to Y$, $g : Y \to Z$ deux morphismes. Si $g \circ f$ est séparé, $f$ est séparé.

Soient $(Z_\lambda)$ un recouvrement de $Z$ par des ouverts affines, $Y_\lambda = g^{-1}(Z_\lambda)$, $X_\lambda = f^{-1}(Y_\lambda)$; en vertu de la prop. 17, $X_\lambda$ est séparé au-dessus de $Z_\lambda$, et il suffit de prouver que $X_\lambda$ est séparé au-dessus de $Y_\lambda$. On peut donc supposer $Z$ affine. Soit alors $(U_\alpha)$ un recouvrement de $Y$ par des ouverts affines, et soit $V_\alpha = f^{-1}(U_\alpha)$; en vertu de la prop. 17, il suffit de montrer que $V_\alpha$ est séparé au-dessus de $U_\alpha$; mais $V_\alpha$ étant séparé au-dessus de $X$ (prop. 14), l'est aussi au-dessus de $Z$ (prop. 15), donc est un schéma (cor. 2), et par suite est séparé au-dessus de $U_\alpha$ (cor. 2).

Proposition 19. — Pour qu'un morphisme $f : X \to Y$ soit séparé, il faut et il suffit que $f_{\operatorname{red}} : X_{\operatorname{red}} \to Y_{\operatorname{red}}$ le soit.

En effet (nº 3, prop. 9), les préschémas $X_{\operatorname{red}} \times_{Y_{\operatorname{red}}} X_{\operatorname{red}}$ et $X_{\operatorname{red}} \times_Y X_{\operatorname{red}}$ s'identifient canoniquement, et si on désigne par $j$ l'injection $X_{\operatorname{red}} \to X$, on vérifie aussitôt que le diagramme
\begin{tikzcd}
X_{\operatorname{red}} \times_{Y_{\operatorname{red}}} X_{\operatorname{red}} \arrow[r, "\Delta_{\operatorname{red}}"] \arrow[d, "j \times j"'] & X_{\operatorname{red}} \times_{Y_{\operatorname{red}}} X_{\operatorname{red}} \arrow[d, "j \times j"] \\
X \times_Y X \arrow[r, "\Delta_X"] & X \times_Y X
\end{tikzcd}

\newpage

%% ===== Page 81 =====

et la proposition résulte de ce que les flèches verticales sont des homéomorphismes des espaces de base.

Corollaire. — Soit $X$ un préschéma, et supposons que l'espace de base $X$ soit réunion d'une famille finie de parties fermées $X_i$ ($1 \le i \le n$). On considère pour chaque $i$ le sous-préschéma réduit ayant $X_i$ pour espace de base. Pour qu'un morphisme $f : X \to Y$ soit séparé, il faut et il suffit que ses restrictions $X_i \to Y$ le soient.

La nécessité résulte des prop. 14 et 15. Inversement, si la condition est satisfaite, et si $p_1, p_2$ sont les projections de $X \times_Y X$, le sous-espace $\Delta_{X_i}(X_i)$ s'identifie au sous-espace $\Delta(X) \cap p_1^{-1}(X_i)$ de l'espace de base de $X \times_Y X$ ; ces sous-espaces étant fermés, il en est de même de leur réunion $\Delta(X)$.

Si en particulier l'espace de base de $X$ est réunion finie de ses composantes irréductibles, le corollaire précédent ramène la notion de séparation au cas des schémas irréductibles.

6. \emph{Premières conséquences de la séparation.}

Proposition 20. — Soient $Y$ un $S$-préschéma séparé au-dessus de $S$, $f : X \to Y$ un $S$-morphisme. Alors le morphisme graphe $\Gamma_f : X \to X \times_S Y$ est une immersion fermée.

La démonstration est la même que celle du cor. 2 de la prop. 13, compte tenu de ce que l'immersion canonique $X \times_Y Y \to X \times_S Y$ est fermée (n°6, lemme 1).

Corollaire 1. — Soient $f : X \to Y$, $g : Y \to Z$ deux morphismes, $g$ étant séparé. Si $g \circ f$ est une immersion (resp. une immersion fermée), il en est de même de $f$.

En effet, $f$ se factorise en : $X \xrightarrow{\Gamma_f} X \times_Z Y \xrightarrow{p_2} Y$, $p_2$ étant la seconde projection. D'autre part, $p_2$ s'identifie à $(g \circ f) \times_Z Y$ (§ 2, n°5, cor. de la prop. 8) ; si $g \circ f$ est une immersion (resp. une immersion fermée) il en est de même de $p_2$ (n°2, prop. 5) ; d'autre part

\newpage

%% ===== Page 82 =====

$\Gamma_f$ est une immersion fermée puisque $g$ est séparé (prop. 20), donc $p_2 \circ \Gamma_f$ est une immersion (resp. une immersion fermée).

Remarque. — Si on ne suppose pas $g$ séparé, le même raisonnement montre que $f$ est une immersion, compte tenu du cor. 2 de la prop. 13 du n°4.

Corollaire 2. — Soit $f: X \to Y$ un morphisme séparé. Tout morphisme $s: Y \to X$, section de $f$ (i.e. tel que $f \circ s = 1_Y$) est une immersion fermée.

C'est un cas particulier du cor. 1.

Corollaire 3. — Soit $Y$ un schéma, $f: X \to Y$ un morphisme. Pour tout ouvert affine $U$ de $X$ et tout ouvert affine $V$ de $Y$, $U \cap f^{-1}(V)$ est affine.

Soient $p_1, p_2$ les projections de $X \times_Y Y$; le sous-espace $U \cap f^{-1}(V)$ est l'image par $p_1$ de $\Gamma_f(X) \cap p_1^{-1}(U) \cap p_2^{-1}(V)$. Or, $p_1^{-1}(U) \cap p_2^{-1}(V)$ s'identifie (en tant que préschéma) à $U \times_V V$ (n°4, cor. du th. 1), et est par suite un schéma affine (§ 2, n°4, prop. 6) ; comme $\Gamma_f(X)$ est fermé dans $X \times_Y Y$ (prop. 20), $\Gamma_f(X) \cap p_1^{-1}(U) \cap p_2^{-1}(V)$ est fermé dans $U \times_V V$, et par suite le préschéma induit par le sous-préschéma de $X \times_Y Y$ associé à $\Gamma_f$, sur la partie ouverte $\Gamma_f(X) \cap p_1^{-1}(U) \cap p_2^{-1}(V)$ de son espace de base $\Gamma_f(X)$, est un sous-préschéma fermé d'un schéma affine, et par suite un schéma affine (n°2, cor. 2 de la prop. 4). Le corollaire résulte alors du fait que $\Gamma_f$ est une immersion.

Remarques. — Si on suppose réciproquement que la conclusion de la prop. 20 est vérifiée lorsque $f$ est l'identité $1_Y$, on en conclut que $Y$ est séparé au-dessus de $S$ ; de même si on suppose le cor. 1 valide pour les morphismes $Y \xrightarrow{\Delta_Y} Y \times_S Y \xrightarrow{p_1} Y$, on conclut que $\Delta_Y$ est une immersion fermée, donc que $Y$ est séparé au-dessus de $S$ ; enfin, la validité de la conclusion du cor. 2 pour le morphisme $\Delta_Y$, qui est une section de la projection $Y \times_S Y \to Y$, implique que $Y$ est fermé au-dessus de $S$.

\newpage

%% ===== Page 83 =====

7. Exemples. —

Soit $(X, \mathcal{O}_X)$ un préschéma, et soit $(U_\alpha)$ un recouvrement de $X$ par des ouverts affines. Désignons par $X_\alpha$ le schéma affine $(U_\alpha, \mathcal{O}_X / U_\alpha)$, par $j_{\beta\alpha} : X_{\alpha\beta} \to X_\alpha$ le morphisme identique ; les restrictions $j'_{\beta\alpha}, j'_{\beta\gamma}, j'_{\alpha\gamma}$ de $j_{\beta\alpha}, j_{\gamma\beta}, j_{\alpha\gamma}$ respectivement à l'espace annelé induit sur $U_\alpha \cap U_\beta \cap U_\gamma$, $X_{\alpha\beta\gamma}$ respectivement sont telles que $j'_{\alpha\gamma} \circ j'_{\gamma\beta} = j'_{\alpha\beta}$.

Inversement, donnons-nous une famille $(X_\alpha)$ de schémas affines, et pour chaque couple d'indices, soit $U_{\alpha\beta}$ une partie ouverte de $X_\alpha$ et soit $u_{\beta\alpha}$ un isomorphisme d'espaces annelés de $U_{\alpha\beta}$ sur $U_{\beta\alpha}$ ; supposons en outre que $u_{\alpha\gamma} \circ u_{\gamma\beta} \circ u_{\beta\alpha} = 1$ dans l'espace annelé induit sur $U_{\alpha\beta} \cap U_{\beta\gamma} \cap U_{\gamma\alpha}$, quels que soient les indices $\alpha, \beta, \gamma$. Alors, il existe un espace annelé $X$, une famille $(U_\alpha)$ d'ouverts affines de $X$ formant un recouvrement de $X$, et pour chaque $\alpha$ un isomorphisme $v_\alpha$ de $X_\alpha$ sur l'espace annelé induit sur $U_\alpha$, de sorte que $v_\alpha$ applique $U_{\alpha\beta}$ sur $U_\alpha \cap U_\beta$ et l'on ait $u_{\beta\alpha} = v_\beta^{-1} \circ v_\alpha$ pour tout couple d'indices. Il est clair par définition que $X$ est un préschéma, et on a ainsi un procédé général de définition des préschémas. Lorsque les $U_{\alpha\beta}$ sont des ouverts affines, les $u_{\beta\alpha}$ sont des isomorphismes de schémas affines, et sont donc définis par les isomorphismes correspondants des anneaux $A(U_{\alpha\beta})$ (§ 1, n°7, th.4) ; et (les $U_{\alpha\beta} \cap U_{\beta\gamma}$ sont aussi des ouverts affines), la vérification des conditions de compatibilité se ramènera aussi à une vérification analogue pour les isomorphismes d'anneaux correspondants (prop. 12).

Par exemple, prenons deux schémas affines isomorphes $X_1, X_2$ spectres des anneaux de polynômes $k[s], k[t]$, où $k$ est un corps, et deux indéterminées ; prenons $U_{12} = D(s)$, $U_{21} = D(t)$, et pour unique isomorphisme correspondant à l'isomorphisme d'anneaux $\varphi : k[s]/s^n \to k[t]/(1/t^n)$. Il n'y a pas de condition de compatibilité, et le préschéma $X$ ainsi défini est séparé.

\newpage

%% ===== Page 84 =====

- 64 -

les critères de la prop. 18 du n°5 étant trivialement vérifiés. Toutefois, le schéma ainsi obtenu n'est pas affine, car on constate aisément que les seules sections au-dessus de $X_1$ qui se prolongent en des sections au-dessus de $X$ tout entier sont les constantes, et le spectre de $\Gamma(X, \mathcal{O}_X)$ est donc réduit à un seul point.

Avec le même choix de $X_1, X_2, U_1, U_2$ mais en prenant pour $u_{12}$ l'isomorphisme qui correspond à l'isomorphisme d'anneaux $f(s) \to f(t)$, on obtient cette fois un préschéma non séparé $X$, car la première condition de la prop. 18 est satisfaite, mais non la seconde.

On peut aussi donner des exemples où aucune des deux conditions de la prop. 18 n'est vérifiée. Remarquons d'abord que dans le spectre premier $Y$ de l'anneau de polynômes $A=K[s, t]$ ($s, t$ indéterminées, $K$ corps), l'ouvert / réunion de $D(s)$ et de $D(t)$ n'est pas un ouvert affine : en effet, on constate aussitôt que les seules sections de $\mathcal{O}_Y$ au-dessus de $U$ sont celles qui se prolongent en sections au-dessus de $X$ tout entier ; si $U$ était un ouvert affine, le morphisme d'injection $j: U \to Y$ serait donc un isomorphisme (§ 1, n°7, th. 4), ce qui est absurde puisque $U \neq Y$.

Cela étant, prenons deux schémas affines $Y_1, Y_2$, spectres premiers des anneaux $A_1 = K[s_1, t_1]$, $A_2 = K[s_2, t_2]$ ; prenons $U_{12} = D(s_1) \cup D(t_1)$, $U_{21} = D(s_2) \cup D(t_2)$, et pour $u_{12}$ la restriction à $U_{21}$ de l'isomorphisme $Y_2 \to Y_1$ correspondant à l'isomorphisme d'anneaux $f(s_1, t_1) \to f(s_2, t_2)$ ; on a ainsi un exemple où aucune des conditions de la prop. 18 n'est satisfaite.

Enfin, pour avoir un exemple de spectre premier dans lequel il y a des ouverts non quasi-compacts, il suffit de prendre l'anneau $C(K)$ des fonctions continues numériques sur un espace compact $K$ dans lequel il y a un point $a$ n'ayant pas de système fondamental dénombrable de voisinages ; si $\mathfrak{m}$ est l'idéal premier maximal de $C(K)$ correspondant à $a$, $\{\mathfrak{m}\}$ n'est pas intersection finie de $V(f)$ dans $X$.

\newpage

%% ===== Page 85 =====

\marginpar{Archives \\ Grothendieck - Sept 59}

- 65 -

§ 4. Conditions de finitude.

1. Préschémas noethériens, localement noethériens, artiniens.

Définition 1. - On dit qu'un préschéma $X$ est noethérien (resp. localement noethérien) s'il est réunion finie (resp. réunion) d'ouverts affines $V_\alpha$ tels que l'anneau de chacun des schémas induits sur les $V_\alpha$ soit noethérien.

Si un préschéma $X$ est réunion finie (resp. réunion) d'ouverts $W_\lambda$ tels que les préschémas induits sur les $W_\lambda$ soient noethériens (resp. localement noethériens), il est clair que $X$ est noethérien (resp. localement noethérien).

Proposition 1. - Pour qu'un préschéma $X$ soit noethérien, il faut et suffit qu'il soit localement noethérien et que son espace de base soit quasi-compact. L'espace de base de $X$ est alors noethérien.

La première assertion est immédiate. La seconde résulte de ce que toute réunion finie de sous-espaces noethériens est noethérien (cf. § 1, n°1, cor. 2 de la prop. 2).

Rappelons que tout espace noethérien est quasi-compact et que tout sous-espace d'un espace noethérien est noethérien. Un espace noethérien n'a qu'un nombre fini de composantes irréductibles. En outre :

Lemme 1. - Dans un espace noethérien vérifiant l'axiome $(T_0)$, toute partie fermée contient une partie fermée réduite à un point. Toute partie ouverte de $X$ contenant tous les points fermés est identique à $X$.

La première assertion résulte de ce que toute partie fermée contient un ensemble fermé minimal, et en vertu de $(T_0)$, un tel ensemble est réduit à un point. L'hypothèse sur $U$ entraîne alors que son complémentaire $U$, qui est fermé, est nécessairement vide.

Proposition 2. - Soit $X$ un schéma affine d'anneau $A$. Les conditions suivantes sont équivalentes : a) $X$ est noethérien ; b) $X$ est localement noethérien ; c) $A$ est noethérien.

\newpage

%% ===== Page 86 =====

- 66 -
L'équivalence de a) et b) résulte de la prop. 1 et de ce que tout schéma affine a un espace de base quasi-compact ($\S$ 1, n$^\circ$3, prop. 3) ; il est clair en outre que c) entraîne a). Pour voir que a) entraîne c), remarquons qu'il y a un recouvrement fini $(V_i)$ de $X$ par des ouverts affines tels que l'anneau $A_i$ du préschéma induit sur $V_i$ soit noethérien. Soit alors $(a_n)$ une suite croissante d'idéaux de $A$ ; il lui correspond canoniquement une suite croissante $(\tilde{a}_n)$ de faisceaux d'idéaux dans $\mathcal{O}_X$ ; pour voir que la suite $(a_n)$ est stationnaire, il suffit de montrer que la suite $(\tilde{a}_n)$ l'est. Or, la restriction $\tilde{a}_i | V_i$ est un faisceau quasi cohérent d'idéaux sur $V_i$, étant l'image réciproque de $\tilde{a}_n$ par l'injection canonique $V_i \to X$ ($\S$ 1, n$^\circ$6, cor. 2 de la prop. 8) ; $\tilde{a}_n | V_i$ est donc de la forme $\tilde{a}_{ni}$, où $a_{ni}$ est un idéal de $A_i$ ($\S$ 1, n$^\circ$3, th. 1). Comme $A_i$ est noethérien, il existe un indice $n_0$ tel que $\tilde{a}_{n+1, i} = \tilde{a}_{ni}$ pour $n \geqslant n_0$ et pour tout $i$ ; on en conclut que $\tilde{a}_{n+1} = \tilde{a}_n$ pour $n \geqslant n_0$, ce qui achève la démonstration.

Proposition 3. - Tout sous-préschéma d'un préschéma noethérien (resp. localement noethérien) est noethérien (resp. localement noethérien).

Il suffit de faire la démonstration pour un préschéma noethérien ; en outre, on est aussitôt ramené au cas où $X$ est un schéma affine. Comme tout sous-préschéma de $X$ est un sous-préschéma fermé d'un préschéma induit sur un ouvert, on peut se borner à considérer le cas d'un sous-préschéma $Y$ fermé ou induit sur un ouvert de $X$. Le cas où $Y$ est fermé est immédiat, car si $A$ est l'anneau de $X$, on sait que $Y$ est un schéma affine d'anneau $A/I$, où $I$ est un idéal de $A$ ($\S$ 3, n$^\circ$2, cor. 2 de la prop. 4) ; comme $A$ est noethérien (prop. 2), il en est de même de $A/I$.

Supposons maintenant $Y$ ouvert dans $X$ ; l'espace de base $Y$ est noethérien, donc quasi-compact, et par suite réunion finie d'ouverts $D(f_i)$ ($f_i \in A$) ; tout revient à démontrer la proposition lorsque $Y = D(f)$ avec $f \in A$. Mais alors $Y$ est un schéma affine dont l'anneau est iso-

\newpage

%% ===== Page 87 =====

-- 67 --

phe à $A_f$ (§ 1, n°3, prop. 6) ; comme $A$ est noethérien (prop. 2), il en est de même de $A_f$.

On notera que le produit de deux préschémas noethériens n'est pas nécessairement noethérien, même si ces préschémas sont affines, car le produit tensoriel de deux algèbres noethériennes n'est pas nécessairement un anneau noethérien.

Définition 2. — On dit qu'un schéma affine est artinien si son anneau est artinien.

Proposition 4. — Étant donné un préschéma $X$, les conditions suivantes sont équivalentes :

a) $X$ est un schéma artinien ;

b) $X$ est noethérien et son espace de base est discret ;

c) $X$ est noethérien et les points de son espace de base sont fermés (condition $T_1$).

Lorsqu'il en est ainsi, l'espace de base de $X$ est fini, et l'anneau de $X$ est composé direct des anneaux locaux (artiniens) des points de $X$.

On sait que a) entraîne la dernière assertion (Bourbaki, Alg. comm., chap. § ) ; tout idéal premier de $A$ est alors maximal et est l'image réciproque de l'idéal maximal de l'un des composants locaux de $A$, donc l'espace $X$ est fini et discret ; a) entraîne b) et b) entraîne évidemment c). Pour voir que c) entraîne a), montrons d'abord que $X$ est alors fini : on peut en effet se ramener alors au cas où $X$ est affine, et on sait qu'un anneau noethérien dont tous les idéaux premiers sont maximaux est artinien (Bourbaki, loc. cit.), d'où notre assertion. L'espace de base $X$ est alors discret, somme topologique d'un nombre fini de points $x_i$, et les anneaux locaux $\mathcal{O}_{X,x_i} = A_i$ sont artiniens ; il est clair alors que $X$ est isomorphe au schéma affine spectre premier de l'anneau $A$ composé direct des $A_i$.

\newpage

%% ===== Page 88 =====

2. Morphismes de type fini.

Définition 2. — On dit qu'un morphisme $f : X \to Y$ est de type fini si $Y$ est réunion d'une famille $(V_\alpha)$ d'ouverts affines ayant les propriétés suivantes : $f^{-1}(V_\alpha)$ est réunion finie d'ouverts affines $U_{\alpha i}$ tels que chacun des anneaux $A(U_{\alpha i})$ soit une algèbre de type fini sur $A(V_\alpha)$. On dit alors aussi que $X$ est un préschéma de type fini sur $Y$, ou un $Y$-préschéma de type fini.

Proposition 5. — Si $f : X \to Y$ est de type fini, alors, pour tout ouvert affine $W$ de $Y$, $f^{-1}(W)$ est réunion finie d'ouverts affines dont les anneaux sont des algèbres de type fini sur $A(W)$.

Démontrons d'abord le lemme suivant :

Lemme 2. — Soit $T \subset Y$ un ouvert affine tel que $f^{-1}(T)$ soit réunion finie d'ouverts affines $Z_j$ tels que $A(Z_j)$ soit une algèbre de type fini sur $A(T)$, alors la même propriété est valable pour tout ouvert affine de la forme $D(g)$ ($g \in A(T)$).

En effet, si $u_j$ est l'homomorphisme $A(T) \to A(Z_j)$ correspondant à la restriction de $f$ à $Z_j$, et si on pose $g_j = u_j(g)$, on a $f^{-1}(D(g)) \cap Z_j = D(g_j)$ (§ 1, n°2, formule (5)) ; or $A(D(g_j)) = (A(Z_j))_{g_j} = A(Z_j)[1/g_j]$ est de type fini sur $A(Z_j)$, et a fortiori sur $A(T)$ et aussi sur $A(D(g)) = A(T)[1/g]$, d'où le lemme.

Cela étant, comme $W$ est quasi-compact, il existe un recouvrement ouvert fini de $W$ par des ensembles de la forme $D(g_i)$ ($g_i \in A(W)$) dont chacun est contenu dans un des $V_\alpha$ ; le lemme montre que chacun des $f^{-1}(D(g_i))$ admet un recouvrement fini par des ouverts affines $U_{ij}$ tels que $A(U_{ij})$ soit une algèbre de type fini sur $A(D(g_i))$ ; comme $A(D(g_i)) = A(W)[1/g_i]$ est de type fini sur $A(W)$, la proposition est démontrée.

On voit donc que la notion de préschéma de type fini sur $Y$ est locale sur $Y$.

Proposition 6. — Soient $X, Y$ deux schémas affines ; pour que $X$ soit

\newpage

%% ===== Page 89 =====

de type fini sur $Y$, il faut et il suffit que $A(X)$ soit une algèbre de type fini sur $A(Y)$.

La condition étant évidemment suffisante, prouvons qu'elle est nécessaire. Posons $A=A(Y)$, $B=A(X)$; en vertu de la prop. 5, il existe un recouvrement ouvert affine fini $(V_i)$ de $X$ tel que chacun des anneaux $A(V_i)$ soit de type fini sur $A$. En outre, les $V_i$ étant quasi-compacts, on peut recouvrir chacun d'eux par un nombre fini d'ouverts de la forme $D(g_{ij}) \subset V_i$ où $g_{ij} \in B$; si $\varphi_i$ est l'homomorphisme canonique $B \to A(V_i)$, $A(V_i)_{g_{ij}} = (A(V_i))_{\varphi_i(g_{ij})}$, donc $B_{g_{ij}}$ est de type fini sur $A$. On peut donc désormais supposer que $V_i = D(g_i)$, où $g_i \in B$; par hypothèse, il existe une partie finie $T_i$ de $B$ et un entier $n_i$ tels que $B_{g_i}$ soit l'algèbre engendrée sur $A$ par les éléments $b_i/g_i^{n_i}$, où $b_i$ parcourt $T_i$. Comme les $g_i$ sont en nombre fini, on peut d'ailleurs supposer tous les $n_i$ égaux à un même entier $n$. En outre, les $D(g_i)$ forment un recouvrement de $X$, l'idéal engendré dans $B$ par les $g_i$ est par suite égal à $B$, autrement dit il y a des $h_i \in B$ tels que $\sum h_i g_i = 1$. Soit alors $T$ la partie finie de $B$ réunion des $T_i$, de l'ensemble des $g_i$ et de l'ensemble des $h_i$; montrons que le sous-anneau $B' = A[T]$ de $B$ est égal à $B$. Par hypothèse, pour tout $b \in B$ et tout $i$, l'image canonique de $b$ dans $B_{g_i}$ est de la forme $b_i/g_i^{n_i}$, où $b_i \in B'$; en multipliant les $b_i$ par des puissances convenables de $g_i$, on peut supposer tous les $n_i$ égaux à un même entier $n$. Cela signifie encore que, dans $B$, il existe un entier $N \geqslant n$ (dépendant de $b$) tel que $g_i^N b \in B'$ pour tout $i$ : or, dans l'anneau $B'$, les $g_i$ engendrent l'idéal $B'$, puisqu'il en est ainsi des $g_i$; il y a donc des $c_i \in B'$ tels que $\sum c_i g_i^N = 1$, d'où $b \in B'$, c.q.f.d.

\emph{Proposition 7.} - Une immersion $f : X \to Y$ est de type fini dans les cas suivants :

a) $f$ est une immersion fermée.

\newpage

%% ===== Page 90 =====

- 70 -

b) $f$ est une immersion ouverte, $Y$ est un schéma et l'espace de base de $X$ est quasi-compact ;

c) l'espace de base de $Y$ est localement noethérien ;

d) l'espace de base de $X$ est noethérien.

On peut toujours supposer que $X$ est un sous-préschéma de $Y$.

Dans le cas a), on peut supposer $Y$ affine ; $X$ étant un sous-préschéma fermé de $Y$, est un schéma affine (cf. 3, n°2, cor. 2 de la prop. 4), dont l'anneau est isomorphe à un anneau quotient $A(Y)/I$, qui est de type fini sur $A(Y)$. Dans le cas b), il y a un recouvrement fini $(V_i)$ de $X$ par des ouverts affines ; pour tout ouvert affine $W$ de $Y$, $W \cap V_i$ est par hypothèse un ouvert affine (cf. 3, n°6, prop. 20), donc réunion d'un nombre fini d'ouverts de la forme $D(g_{ij})$ où $g_{ij} \in A(Y)$, et $A(D(g_{ij})) = A(Y)[1/g_{ij}]$ est une algèbre de type fini sur $A(Y)$. Dans le cas c), on peut supposer $Y$ affine, et alors l'espace de base de $X$ est noethérien ; on est donc ramené au cas d). Dans ce dernier cas, il existe encore un recouvrement de $X$ par un nombre fini d'ouverts affines $D(g_i) \subset Y$ ($g_i \in A(Y)$) tels que $X \cap D(g_i)$ soit fermé dans $D(g_i)$ ; alors $A(X \cap D(g_i))$ est de type fini sur $A(D(g_i))$, d'après le cas a), et $A(D(g_i)) = A(Y)_{g_i}$ est de type fini sur $A(Y)$.

Proposition 8. — Le composé de deux morphismes de type fini est de type fini.

Soient $f : X \to Y$, $g : Y \to Z$ deux morphismes de type fini. Soit $U$ un ouvert affine dans $Z$ tel que $g^{-1}(U)$ admette un recouvrement fini par des ouverts affines $V_i$ tels que $A(V_i)$ soit une algèbre de type fini sur $A(U)$ ; d'après la prop. 5, chacun des $f^{-1}(V_i)$ admet un recouvrement fini par des ouverts affines $W_{ij}$ tels que $A(W_{ij})$ soit de type fini sur $A(V_i)$ et par suite de type fini sur $A(U)$, d'où la proposition.

Corollaire. — Soit $f : X \to Y$ un morphisme de type fini ; si $T$ est un sous-préschéma de $X$ qui est fermé, ou dont l'espace de base est noethérien, ou si $X$ est localement noethérien, la restriction...

\newpage

%% ===== Page 91 =====

- 71 -

Cela résulte des prop. 7 et 8.

Proposition 9. — Soit $X$ un préschéma de type fini au-dessus de $Y$ ; pour tout morphisme $Y' \to Y$, $X'$ est de type fini au-dessus de $Y'$.

Soit $f$ le morphisme $X \to Y$, et soient $p, q$ les projections $X' \to X$ et $X' \to Y'$. Soit $V$ un ouvert affine dans $Y$ tel que $f^{-1}(V)$ soit réunion finie d'ouverts affines $W_i$ dont chacun est tel que $A(W_i)$ soit de type fini sur $A(V)$. Soit $V'$ un ouvert affine de $Y'$ contenu dans $g^{-1}(V)$ ; comme $f \circ p = g \circ q$, $q^{-1}(V')$ est contenu dans la réunion des $p^{-1}(W_i)$ ; d'autre part, l'intersection $p^{-1}(W_i) \cap q^{-1}(V')$ s'identifie au produit $W_i \times_Y Y'$, qui est un schéma affine d'anneau $A(W_i) \otimes_{A(Y)} A(Y')$ ; ce dernier étant par hypothèse une algèbre de type fini sur $A(Y')$, la proposition est démontrée.

Corollaire 1. — Si $f : X \to X'$, $g : Y \to Y'$ sont deux $S$-morphismes de type fini, il en est de même de $f \times_S g$.

Tout d'abord si $g = 1_Y$, on déduit le résultat de la prop. 9 ; on remarque que $X \times_S Y$ s'identifie à $X \times_S (X' \times_S Y)$ et $X \times_S Y \to X' \times_S Y$. On passe de là au cas général en écrivant $f \times_S g = (f \times_S 1_{Y'}) \circ (1_X \times_S g)$ et appliquant la prop. 8.

Corollaire 2. — Soient $f : X \to Y$, $g : Y \to Z$ deux morphismes. Si $g \circ f$ est de type fini, et si en outre $g$ est séparé, ou $X$ noethérien, ou $X \times_Z Y$ localement noethérien, $f$ est de type fini.

On factorise $f : X \to X \times_Z Y \to Y$, $p_2$ s'identifiant à $(g \circ f) \times_Z 1_Y$. En vertu du cor. 1, $p_2$ est de type fini, et il suffit en vertu de la prop. 8 de montrer que $f_p$ est de type fini ; comme c'est une immersion (resp. une immersion fermée si $g$ est séparé) (§ 3, n°7, prop. 22) il suffit d'appliquer la prop. 7.

Proposition 10. — Soit $f : X \to Y$ un morphisme de type fini ; si $Y$ est noethérien (resp. localement noethérien), $X$ est noethérien (resp. localement noethérien).

\newpage

%% ===== Page 92 =====

- 72 -

On peut se borner à faire la démonstration lorsque $Y$ est noethérien. Alors $Y$ est réunion finie d'ouverts affines $V_i$ tels que les $A(V_i)$ soient noethériens. D'après la prop. 5, chacun des $f^{-1}(V_i)$ est réunion d'un nombre fini d'ouverts affines $W_{ij}$ tels que $A(W_{ij})$ soit une algèbre de type fini sur $A(V_i)$, donc un anneau noethérien ; $X$ est par suite noethérien.

\noindent\textbf{Corollaire 1.} --- Soit $X$ un préschéma de type fini sur $S$. Pour toute extension $S' \to S$ telle que $S'$ soit noethérien (resp. localement noethérien), $X^{S'}$ est noethérien (resp. localement noethérien).

En effet, $X^{S'}$ est de type fini sur $S'$ (prop. 9), et le corollaire résulte de la prop. 10.

\noindent\textbf{Corollaire 2.} --- Soit $X$ un préschéma de type fini sur un préschéma noethérien $S$. Alors tout $S$-morphisme $f : X \to Y$ est de type fini.

En effet, si $\varphi : X \to S$, $\psi : Y \to S$ sont les morphismes structuraux, on a $\varphi = \psi \circ f$, et $X$ est noethérien d'après la prop. 10 ; donc $f$ est de type fini en vertu du cor. 2 de la prop. 9.

\noindent\textbf{Proposition 11.} --- Pour qu'un morphisme de type fini $f : X \to Y$ soit surjectif, il faut et il suffit que pour tout corps algébriquement clos $\Omega$, l'application $X(\Omega) \to Y(\Omega)$ correspondant à $f$ (§ 2, n°6) soit surjective.

La condition est suffisante en vertu de la prop. 12 du § 2, n°7. Inversement, supposons $f$ surjectif et soit $g : \{\xi\} = \operatorname{Spec}(\Omega) \to Y$ un morphisme, $\Omega$ étant algébriquement clos. Posons $y = g(\xi)$ ; par hypothèse $f^{-1}(y)$ n'est pas vide, et le morphisme $f' : f^{-1}(y) = X \times_Y \{\xi\} \to \{\xi\}$ est de type fini (prop. 9) ; donc $f^{-1}(y)$ est réunion finie d'ouverts affines $Z_i$ tels que $A(Z_i)$ soit une algèbre de type fini sur $\Omega$. Considérons un de ces schémas $Z_i$ ; en vertu du th. des zéros de Hilbert il existe un homomorphisme $A(Z_i) \to \Omega$, et le morphisme correspondant $h : \{\xi\} \to Z_i$ est tel que $f \circ h$ soit l'identité. La commutativité du diagramme

\newpage

%% ===== Page 93 =====

\begin{center}
\begin{tikzcd}
X \arrow[r, "f"] & Y \\
Z_i \arrow[u, hook] \arrow[r] & \{\xi\} \arrow[u, "g"']
\end{tikzcd}
\end{center}
achève alors de prouver la proposition.

\noindent \textsc{Corollaire}. --- Soient $K$ un corps, $X$ et $Y$ deux préschémas de type fini sur $K$, $f: X \to Y$ un $K$-morphisme. Pour que $f$ soit surjectif, il faut et il suffit que pour toute extension algébriquement close $\Omega$ de $K$, l'application $X(\Omega)_K \to Y(\Omega)_K$ correspondant à $f$ (§ 2, n°6) soit surjective.

Pour voir que la condition est suffisante, on peut encore appliquer la prop. 12 du § 2, n°7, en remarquant que pour tout $y \in Y$, $\kappa(y)$ est une extension de $K$; la nécessité résulte de la prop. 11, appliquée à un $K$-morphisme $\operatorname{Spec}(\Omega) \to Y$, en remarquant que $f$ est nécessairement de type fini (cor. 2 de la prop. 10).

Etant donné un corps $K$, on appelle $K$-préschéma algébrique un préschéma de type fini sur $K$ ; $K$ est appelé le corps de définition de $X$ ; si $X$ est un schéma, on dit aussi que c'est un $K$-schéma algébrique.

\noindent \textsc{Proposition 12}. --- Soit $X$ un préschéma algébrique sur un corps $K$. Les propriétés suivantes sont équivalentes :

a) $X$ est artinien ;

b) l'espace de base $X$ est discret ;

c) l'espace de base $X$ est fini ;

d) les points de l'espace de base $X$ sont fermés ;

e) on a $X = \operatorname{Spec}(A)$, où $A$ est une $K$-algèbre de rang fini.

Il résulte de l'hypothèse que $X$ est noethérien (prop. 10) ; l'équivalence de a), b), d) et e) (et ces conditions entraînent e)) résulte alors de la prop. 4 du n°1. Reste à voir que c) entraîne b) ; on peut se borner au cas où $X$ est affine. Alors $A(X)$ est une $K$-algèbre de type fini dans laquelle il n'y a qu'un nombre fini d'idéaux premiers ; on sait (Bourbaki, Alg. comm.,

\newpage

%% ===== Page 94 =====

- 74 -

chap., $\S$ 4) qu'une telle algèbre est de rang fini sur $K$, donc un anneau artinien, et par suite $X$ est un espace discret.

Lorsque les conditions de la prop. 12 sont satisfaites, on dit que $X$ est un schéma fini sur $K$, de rang $[A:K]$.

\emph{Corollaire 1}. — Soit $X$ un schéma fini sur un corps $K$. Pour toute extension $K'$ de $K$, $X \otimes_K K'$ est un schéma fini sur $K'$, et son rang sur $K'$ est égal au rang de $X$ sur $K$.

En effet : $[A \otimes_K K'] = [A:K]$.

\emph{Corollaire 2}. — Soit $X$ un schéma fini sur un corps $K$; on pose $n = \sum_{x \in X} [k(x):K]_s$ ; alors, pour toute extension algébriquement close $\Omega$ de $K$, l'espace de base de $X \otimes_K \Omega$ a exactement $n$ points.

On peut évidemment se borner au cas où l'anneau $A=A(X)$ est local ; soient $\mathfrak{m}$ son idéal maximal, $L = A/\mathfrak{m}$ son corps résiduel, $L_0$ la plus grande extension séparable de $K$ contenue dans $L$. Le radical de $A \otimes_K \Omega$ est l'image réciproque du radical de $L \otimes_K \Omega = L_0 \otimes_K \Omega = (L_0 \otimes_K \Omega) \otimes_{L_0} L$ ; comme $L_0 \otimes_K \Omega$ est composé direct de $[L_0:K]$ corps isomorphes à $\Omega$, et que $\Omega \otimes_{L_0} L$ est un anneau local de corps résiduel $\Omega$, on voit que $A \otimes_K \Omega$ est composé direct de $n$ anneaux locaux.

Le nombre $n$ est appelé le rang séparable de $A$ sur $K$, ou le nombre géométrique de points de $X$. Il résulte aussitôt de cette définition que pour toute extension $K'$ de $K$, $X \otimes_K K'$ a même nombre géométrique de points que $X$. Si on désigne ce nombre par $n(X)$, il est clair que si $X$ est somme de deux schémas finis $Y, Z$ sur $K$, on a $n(X) = n(Y) + n(Z)$. En outre, si $X, Y$ sont deux schémas finis sur $K$, on a $n(X \times Y) = n(X)n(Y)$ : en effet, on peut encore se ramener au cas où $X$ et $Y$ sont des spectres premiers d'anneaux locaux $A, B$ ; si $E, F$ sont les corps résiduels de $A$ et $B$, le nombre de composants locaux de $A \otimes_K B$ est égal au nombre de composants locaux de $E \otimes_K F$, et on conclut comme dans le cor. 2 ci-dessus.

\newpage

%% ===== Page 95 =====

- 75 -

tout $y \in Y$, la fibre $f^{-1}(y)$ est un préschéma algébrique sur le corps résiduel $\kappa(y)$.

Comme $f^{-1}(y) = X \times_Y \operatorname{Spec}(\kappa(y))$, la proposition résulte de la prop. 9.

\emph{Proposition 14}. — Soient $f : X \to Y$, $g : Y' \to Y$ deux morphismes ; posons $X' = X \times_Y Y'$ et $f' = f \times_Y \operatorname{id}_{Y'} : X' \to Y'$. Si $y' \in Y'$ est tel que la fibre $f^{-1}(y)$, où $y = g(y')$, est un schéma algébrique fini sur $\kappa(y')$, alors $f'^{-1}(y')$ est un schéma algébrique fini sur $\kappa(y')$, ayant même rang et même nombre géométrique de points que $f^{-1}(y)$.

Cela résulte aussitôt des cor. 1 et 2 de la prop. 12 ci-dessus et du § 2, n°7, prop. 18.

3. \emph{Détermination locale d'un morphisme}.

\emph{Proposition 15}. — Soient $X$ et $Y$ deux $S$-préschémas, $Y$ étant de type fini sur $S$ ; soient $x \in X$, $y \in Y$ au-dessus d'un même point $s \in S$.

(i) Si deux $S$-morphismes $f=(\varphi, \theta)$, $f'=(\varphi', \theta')$ sont tels que $\varphi(x) = \varphi'(x) = y$ et que les $\mathcal{O}_s$-homomorphismes $\theta_y$ et $\theta'_y$ de $\mathcal{O}_{Y,y}$ dans $\mathcal{O}_{X,x}$ soient identiques, $f$ et $f'$ coïncident dans un voisinage de $x$.

(ii) Supposons en outre $S$ localement noethérien. Pour tout $\mathcal{O}_s$-homomorphisme $\varphi : \mathcal{O}_{Y,y} \to \mathcal{O}_{X,x}$, il existe un $S$-morphisme $f=(\varphi, \theta)$ d'un voisinage ouvert $U$ de $x$ dans $Y$, tel que $\varphi(x) = y$ et $\theta_y = \varphi$.

(1) On se ramène aussitôt au cas où $S, X, Y$ sont affines d'anneaux respectifs $A, B, C$, et $f$ et $f'$ étant de la forme $(\varphi, \varphi^\sharp)$ et $(\varphi', \varphi'^\sharp)$, où $\varphi$ et $\varphi'$ sont deux $A$-homomorphismes de $C$ dans $B$ tels que $\varphi^{-1}(j_x) = \varphi'^{-1}(j_x) = j_y$, et les homomorphismes $\varphi_x$ et $\varphi'_x$ de $C_y$ dans $B_x$, déduits de $\varphi$ et $\varphi'$, sont identiques ; on peut en outre supposer que $C$ est une $A$-algèbre de type fini. Soient $c_i$ ($1 \le i \le n$) des générateurs de $C$ sur $A$, et posons $b_i = \varphi(c_i)$, $b'_i = \varphi'(c_i)$ ; par hypothèse, on a $b_i \equiv b'_i \pmod{j_x}$ dans l'anneau de fractions $B_x$. Autrement dit, il existe des éléments $s_i \in B \setminus j_x$ tels que $s_i(b_i - b'_i) = 0$, et on peut évidemment supposer tous les $s_i$ égaux à un même élément $s$.

\newpage

%% ===== Page 96 =====

- 76 -

$g \in B_{j_x}$. Par suite, on a $b_i/1 = b'_i/1$ dans l'anneau de fractions $B_g$, et si $i_g$ est l'homomorphisme canonique $B \to B_g$, on voit donc que l'on a $i_g \varphi = i_g \varphi'$; on en conclut que les restrictions de $f$ et $f'$ à $D(g)$ sont identiques.

(ii) On peut se ramener à la même situation que dans (i) et supposer en outre que $A$ est noethérien. Soient $c_i$ ($1 \le i \le n$) des générateurs de l'algèbre $C$ sur $A$, et soit $\alpha : A[X_1, \dots, X_n] \to C$ l'homomorphisme de l'algèbre de polynômes $A[X_1, \dots, X_n]$ transformant $X_i$ en $c_i$ pour $1 \le i \le n$ ; soit d'autre part $i_y$ l'homomorphisme canonique $C \to \mathcal{O}_y$, et considérons l'homomorphisme composé
$$\beta : A[X_1, \dots, X_n] \to C \xrightarrow{i_y} \mathcal{O}_y \to B_x$$
Désignons par $a$ le noyau de $\beta$ ; comme $A$ est noethérien, il en est de même de $A[X_1, \dots, X_n]$, et par suite $a$ admet un système fini de générateurs $Q_j(X_1, \dots, X_n)$ ($1 \le j \le m$). D'autre part, chacun des éléments $\varphi(i_y(c_i))$ peut s'écrire $b_i/s_i$ où $b_i \in B$ et $s_i \in j_x$ ; on peut en outre supposer tous les $s_i$ égaux à un même élément $g \in B_{j_x}$. Cela étant, on a par hypothèse $Q_j(b_1/g, \dots, b_n/g) = 0$ dans $B_x$ ; posons $Q_j(X_1, \dots, X_n) = P_j(X_1, \dots, X_n) / T^{k_j}$, où $P_j$ est homogène de degré $k_j$ ; on peut d'ailleurs supposer tous les $k_j$ égaux à un même nombre $k$. Soit alors $d_j = P_j(b_1, \dots, b_n, g) \in B$. Par hypothèse, on a $t_j d_j = 0$ pour un $t_j \in B \setminus j_x$ ; on peut évidemment supposer tous les $t_j$ égaux à un même élément $h \in B \setminus j_x$ ; on en conclut que $P_j(hb_1, \dots, hb_n, hg) = 0$. Cela étant, considérons l'homomorphisme $\rho$ de $A[X_1, \dots, X_n]$ dans l'anneau de fractions $B_{hg}$ qui applique $X_i$ sur $hb_i/hg$ ; l'image de $a$ par cet homomorphisme est 0, et il en est a fortiori de même de l'image par $\rho$ du noyau $\alpha^{-1}(0)$ ; donc $\rho$ se factorise en $A[X_1, \dots, X_n] / a \to C \to B_{hg}$, avec $\bar{\rho}(c_i) = hb_i/hg$, et il est clair que si $i_x$ est l'homomorphisme canonique $B_{hg} \to B_x$, le diagramme

\newpage

%% ===== Page 97 =====

\begin{tikzcd}
& C \arrow[d, "i_y"] \arrow[r, "\gamma"] & B_{hg} \arrow[d, "i_x"] \\
& C_y \arrow[r, "\varphi"] & B_x
\end{tikzcd}
est commutatif ; $f = (\varphi, \tilde{\varphi})$ est donc un $S$-morphisme du voisinage $D(hg)$ de $x$ dans $Y$ répondant à la question.

Proposition 16. — Soient $f: X \to Y$ un morphisme de type fini, $x$ un point de $X$, $y = f(x)$.

(i) Pour que $f$ soit une immersion locale au point $x$ ($\S$ 3, $n^\circ 3$, déf. 3), il faut et il suffit que $\mathcal{O}_{Y,y} \to \mathcal{O}_{X,x}$ soit surjectif.

(ii) On suppose en outre $Y$ localement noethérien. Pour que $f$ soit un isomorphisme local au point $x$ ($\S$ 3, $n^\circ 3$, déf. 4), il faut et il suffit que $\mathcal{O}_{Y,y} \to \mathcal{O}_{X,x}$ soit un isomorphisme.

(ii) D'après la prop. 15, il existe un voisinage ouvert $V$ de $y$ et un morphisme $g: V \to X$ tels que $g \circ f$ (resp. $f \circ g$) coïncide avec l'identité dans un voisinage de $x$ (resp. $y$), d'où aussitôt la conclusion.

(i) On se ramène au cas où $X$ et $Y$ sont affines, d'anneaux respectifs $A, B$; on a $f = (\varphi, \tilde{\varphi})$, où $\varphi$ est un homomorphisme $B \to A$, qui fait de $A$ une $B$-algèbre de type fini ; on a $\varphi^{-1}(j_x) = j_y$, et l'homomorphisme $\varphi_x : B_y \to A_x$ déduit de $\varphi$ est surjectif. Soit $(a_i)$ un système fini de générateurs de $A$ sur $B$ ; par hypothèse, il existe des $b_i \in B$ et $c \in B - j_y$ tels que $a_i/1 = \varphi(b_i)/\varphi(c)$ dans l'anneau $A_x$ ; cela implique par définition qu'il existe $a \in A - j_x$ tel que, si on pose $g = \varphi(c)$, on ait $a(a_i/1 - \varphi(b_i)/\varphi(c)) = 0$ dans l'anneau $A_x$. Cela étant, il existe par hypothèse un polynôme $Q(X_1, \dots, X_n)$ à coefficients dans l'anneau $\varphi(B)$ tel que $a = Q(a_1, \dots, a_n)$ ; posons $Q(X_1/T, \dots, X_n/T) = P(X_1, \dots, X_n)/T^m$, où $P$ est homogène de degré $m$ dans l'anneau $A_g$ ; on a $a/1 = a^m \varphi(b_1) \dots \varphi(b_n) \cdot \varphi(c) / g^m = a^m \tilde{a}(a) / g^m$. Mais comme dans $A_g$, $g/1 = (c/1)(\varphi(c)/1)$ est inversible, il en est de même de $a/1$ et $\varphi(c)/1$, et on peut donc écrire $a/1 = (\varphi(a)/1)(\varphi(c)/1)^{-m}$ ; on en conclut que $\varphi(a)/1$ est aussi inversible.

\newpage

%% ===== Page 98 =====

- 78 -

inversible dans $A_{\mathfrak{g}}$. Posons alors $h=cd$ ; comme $\varphi(h)/1$ est inversible dans $A_{\mathfrak{g}}$, l'homomorphisme composé $B \to B_h \to A_{\mathfrak{g}}$ se factorise en $B \to B_h \to A_{\mathfrak{g}}$. Montrons que $\gamma$ est surjectif : il suffit de vérifier que l'image de $\varphi$ dans $A_{\mathfrak{g}}$ contient les $t_i/1$ et $(g/1)^{-1}$ ; or, on a $(g/1)^{-1} = (\varphi(c)/1)^{-1} (\varphi(d)/1)^{-1} = \gamma(d^m/h^m)$, et $a/1 = \gamma(b_i d^{m+1}/h^m)$ donc $(a\varphi(b_i))/1 = \gamma(b_i d^{m+1}/h^m)$, et comme $t_i/1 = (a\varphi(b_i))/1 (g/1)^{-1}$, notre assertion est démontrée. On en conclut aussitôt (§ 3, n°2, cor. 2 de la prop. 4) que la restriction de $f$ à $D(g)$, égale à $(\gamma, \tilde{\gamma})$, est une immersion fermée de $D(g)$ dans $D(h)$.

Corollaire. — Soit $f : X \to Y$ un morphisme de type fini. On suppose $X$ irréductible, et on désigne par $x$ son point générique et on pose $y=f(x)$. Pour qu'il existe un ouvert non vide $U$ de $X$ tel que la restriction de $f$ au sous-préschéma induit sur $U$ soit une immersion fermée dans un sous-préschéma induit sur un ouvert de $Y$, il faut et il suffit que $\theta_x^b : \mathcal{O}_y \to \mathcal{O}_x$ soit surjectif. Si de plus $Y$ est irréductible et localement noethérien, pour qu'il existe un ouvert non vide $U$ de $X$ tel que la restriction de $f$ au sous-préschéma induit sur $U$ soit un isomorphisme sur un sous-préschéma induit sur un ouvert de $Y$, il faut et il suffit que $y$ soit le point générique de $Y$ (autrement dit, que $f(x)$ soit dense dans $Y$) et que $\theta_x^b$ soit un isomorphisme.

La première assertion résulte de la prop. 16, compte tenu de ce que tout ouvert non vide de $X$ contient $x$. Comme $\{x\}$ est dense dans $X$, $\{y\}$ est dense dans $f(x)$ et pour que $f(x)$ contienne un ouvert non vide de $Y$, il est donc nécessaire que $\{y\}$ soit dense dans $Y$, autrement dit que $y$ soit le point générique de $Y$ (lorsque $Y$ est supposé irréductible) ; la seconde assertion résulte alors de la prop. 16.

Lorsque les conditions de la seconde partie du cor. de la prop. 16 sont vérifiées, on voit donc que $\theta_x^b$ est un isomorphisme de l'anneau des fonctions rationnelles sur $Y$, sur l'anneau des fonctions rationnelles sur $X$. On dit alors que $f$ est un morphisme birationnel.

\newpage

%% ===== Page 99 =====

inversible dans $A_g$. Posons alors $h=cd$ ; comme $\varphi(h)/1$ est inversible dans $A_g$, l'homomorphisme composé $B \to B_g \to A_g$ se factorise en $B \to B_h \to A_g$. Montrons que $\gamma$ est surjectif : il suffit de vérifier que l'image de $\gamma$ dans $A_g$ contient les $t_i/1$ et $(g/1)^{-1}$ ; or, $(g/1)^{-1} = (\varphi(c)/1)^{-1}(\varphi(d)/1)^{-1} = \varphi(c^n/h)$, et $a/1 = \gamma(a^{n+1}/h^n)$ donc $(a(b_1))/1 = \gamma(b_1 a^{n+1}/h^n)$, et comme $t_1/1 = (\varphi(b_1)/1)(\varphi(g)/1)^{-1}$, notre assertion est démontrée. On en conclut aussitôt (§ 3, n$^\circ$ 2, cor. 2 de la prop. 4) que la restriction de $f$ à $D(g)$, égale à $(\gamma, \varphi)$, est une immersion fermée de $D(g)$ dans $D(h)$.

Corollaire. — Soit $f: X \to Y$ un morphisme de type fini. On suppose $X$ irréductible, on désigne par $x$ son point générique et on pose $y = f(x)$. Pour qu'il existe un ouvert non vide $U$ de $X$ tel que la restriction de $f$ au sous-préschéma induit sur $U$ soit une immersion fermée dans un sous-préschéma induit sur un ouvert de $Y$, il faut et il suffit que $\theta^b_x : \mathcal{O}_y \to \mathcal{O}_x$ soit surjectif. Si de plus $Y$ est irréductible et localement noethérien, pour qu'il existe un ouvert non vide $U$ de $X$ tel que la restriction de $f$ au sous-préschéma induit sur $U$ soit un isomorphisme sur un sous-préschéma induit sur un ouvert de $Y$, il faut et il suffit que $y$ soit le point générique de $Y$ (autrement dit, que $\varphi(x)$ soit dense dans $Y$) et que $\theta^b_x$ soit un isomorphisme.

La première assertion résulte de la prop. 16, compte tenu de ce que tout ouvert non vide de $X$ contient $x$. Comme $\{x\}$ est dense dans $X$, $\{y\}$ est dense dans $\varphi(x)$ et pour que $\varphi(x)$ contienne un ouvert non vide de $Y$, il est donc nécessaire que $\{y\}$ soit dense dans $Y$, autrement dit que $y$ soit le point générique de $Y$ (lorsque $X$ est supposé irréductible) ; la seconde assertion résulte alors de la prop. 16.

Lorsque les conditions de la seconde partie du cor. de la prop. 16 sont vérifiées, on voit donc que $\theta^b_x$ est un isomorphisme de l'anneau des fonctions rationnelles sur $Y$, sur l'anneau des fonctions rationnelles sur $X$. On dit alors que $f$ est un morphisme birationnel.

\newpage

%% ===== Page 100 =====

\marginpar{- vi -}

The end of § 2, from n°6 on, should be replaced by the following text:

6. Points d'un préschéma à valeurs dans un autre ; points géométriques.

Soit $X$ un préschéma ; pour tout préschéma $T$, on désignera l'ensemble des morphismes de $T$ dans $X$ par $\operatorname{Hom}(T,X)$ ou encore par $X(T)$, et les éléments de cet ensemble seront encore appelés les points de $X$ à valeurs dans $T$. Si on associe à tout morphisme $f : T \to T'$ l'application $u \mapsto u \circ f$ de $X(T')$ dans $X(T)$, on voit que, pour $X$ fixé, $T \mapsto X(T)$ est un foncteur contravariant de la catégorie des préschémas dans celle des ensembles. En outre, tout morphisme de préschémas $g : X \to Y$ définit un homomorphisme fonctoriel $X(T) \to Y(T)$, faisant correspondre $g \circ v$ à $v \in X(T)$.

Etant donnés trois ensembles $P, Q, R$ et deux applications $\varphi : P \to R$, $\psi : Q \to R$, on appelle produit fibré de $P$ et $Q$ au-dessus de $R$ la partie de l'ensemble produit $P \times Q$ formée des couples $(p,q)$ tels que $\varphi(p) = \psi(q)$; on le note $P \times_R Q$. La définition du produit de préschémas (n°4, déf.1) peut encore s'interpréter, avec les notations précédentes, par la formule

$$(1) \quad (X \times_S Y)(T) = X(T) \times_{S(T)} Y(T)$$

et les applications $X(T) \to S(T)$ et $Y(T) \to S(T)$ correspondant aux morphismes structuraux $X \to S, Y \to S$.

Si on se donne un préschéma $S$ et qu'on ne considère que les $S$-préschémas et les $S$-morphismes (n°3), on désignera de même l'ensemble des $S$-morphismes $T \to X$ par $\operatorname{Hom}_S(T,X)$ ou $X(T)_S$ (en se permettant de supprimer l'indice $S$ quand aucune confusion n'est possible). La formule (1) s'écrit alors aussi

$$(2) \quad (X \times_S Y)(T)_S = X(T)_S \times Y(T)_S$$

et plus généralement, si $Z$ est un $S$-préschéma, $X, Y$ des $Z$-préschémas (donc ipso facto des $S$-préschémas), on a

\newpage

%% ===== Page 101 =====

(3)
$$(X \times_S Y)(T) = X(T) \times_{S(T)} Y(T)$$
Lorsque dans ce qui précède $T$ (resp. $S$) est de la forme $\operatorname{Spec}(B)$ (resp. $\operatorname{Spec}(A)$), on se permet de remplacer $T$ (resp. $S$) par $B$ (resp. $A$) dans les notations précédentes, et on parle alors de points de $X$ à valeurs dans l'anneau $B$, ou de points du $A$-préschéma $X$ à valeurs dans la $A$-algèbre $B$ pour les éléments de $X(B)$ et de $X(B)_A$ respectivement. On notera que $X(B)$ et $X(B)_A$ sont maintenant des foncteurs covariants de $B$. On écrira de même $X(T)_A$ pour l'ensemble des points du $A$-préschéma $X$ à valeurs dans le $A$-préschéma $T$.

Considérons en particulier le cas où $T$ est de la forme $\operatorname{Spec}(A)$, où $A$ est un anneau local ; les éléments de $X(A)$ correspondent alors biunivoquement aux homomorphismes $\mathcal{O}_{X,x} \to A$ ($x \in X$) ($n^\circ 2$, prop. 5) ; on dit que le point $x$ de l'espace de base de $X$ est la \emph{localité} du point de $X$ à valeurs dans $A$ auquel il correspond.

Plus particulièrement, on appellera points géométriques d'un préschéma $X$ les points de $X$ à valeurs dans un corps $K$ : la donnée d'un tel point revient donc à la donnée de sa localité $x$ dans l'espace de base de $X$, et d'une extension $K$ de $\kappa(x)$ ; $K$ sera appelé le \emph{corps des valeurs} du point géométrique correspondant. On définit ainsi une application $X(K) \to X$, appliquant un point géométrique sur sa localité.

Si $X$ et $\operatorname{Spec}(K)$ sont des $S$-préschémas (autrement dit, si $K$ est considéré comme extension d'un corps résiduel $\kappa(s)$, où $s \in S$), un élément de $X(K)_S$, ou, comme on dit encore, un point de $X$ à valeurs dans l'extension $K$ de $\kappa(s)$, consiste en la donnée d'un $\kappa(s)$-isomorphisme d'un corps résiduel $\kappa(x)$ dans $K$, où $x$ est un point de $X$ au-dessus de $s$ (donc $\kappa(x)$ une extension de $\kappa(s)$).

Lemme 5. — Soient $X_i$ ($1 \le i \le n$) des $S$-préschémas, $s$ un point de l'espace de base $S$, $x_i$ ($1 \le i \le n$) un point de l'espace de base $X_i$ au-dessus de $s$. Il existe alors une extension $K$ de $\kappa(s)$ et un point géo-

\newpage

%% ===== Page 102 =====

- viii -

\noindent $(Y = X_1 \times_S X_2 \times \dots \times X_n)$, à valeurs dans $K$, dont les projections soient localisées aux $x_i$.

En effet, il existe des $\kappa(s)$-monomorphismes $\kappa(x_i) \to K$ dans une même extension $K$ de $\kappa(s)$ ; les composés $\kappa(s) \to \kappa(x_i) \to K$ sont donc tous identiques (Bourbaki, Alg., chap. V, § 4, prop. 2). On en déduit que les morphismes correspondants $\operatorname{Spec}(K) \to X_i$ sont des $S$-morphismes, et par suite définissent un $S$-morphisme unique $\operatorname{Spec}(K) \to Y$. Si $y$ est le point correspondant de $Y$, il est clair que les projections sur chacun des $X_i$ est $x_i$.

\noindent \textbf{Proposition 10.} — Soient $X_i$ ($1 \leq i \leq n$) des $S$-préschémas et pour chaque indice $i$ soit $x_i$ un point de l'espace de base $X_i$. Pour qu'il existe un point $y$ de l'espace de base de $Y = X_1 \times_S X_2 \times \dots \times X_n$ dont $x_i$ soit la projection d'indice $i$ ($1 \leq i \leq n$), il faut et il suffit que les $x_i$ soient au-dessus d'un même point $s$ de l'espace de base $S$.

La condition est évidemment nécessaire ; le lemme 5 montre qu'elle est suffisante.

En d'autres termes, si l'on désigne par $|X|$ l'ensemble sous-jacent à l'espace de base de $X$, on voit qu'on a une application surjective canonique $|X \times_S Y| \to |X| \times_{|S|} |Y|$ ; il faut noter que cette application n'est pas injective en général (autrement dit, il peut exister plusieurs points distincts ayant mêmes projections).

\noindent \textbf{Corollaire.} — Soient $f: X \to Y$ un $S$-morphisme, $f^{S'} : X^{S'} \to Y^{S'}$ le $S'$-morphisme déduit de $f$ par une extension $S' \to S$ du préschéma de base. Soit $p$ (resp. $q$) la projection $X^{S'} \to X$ (resp. $Y^{S'} \to Y$) ; pour toute partie $A$ de l'espace de base $X$, on a $q^{-1}(f(A)) = f^{S'}(p^{-1}(A))$.

En effet (n° 5, cor. 2 de la prop. 9), $X^{S'}$ s'identifie au produit $X \times_Y Y^{S'}$ par le diagramme
$$\begin{tikzcd}
X^{S'} \arrow[r, "p"] \arrow[d, "f^{S'}"] & X \arrow[d, "f"] \\
Y^{S'} \arrow[r, "q"] & Y
\end{tikzcd}$$

\newpage

%% ===== Page 103 =====

- ix -

En vertu de la prop. 10 , la relation $q(y')=p(x)$ pour $x \in X$ , $y' \in Y^{S'}$ , équivaut à l'existence d'un $x' \in X^{S'}$ tel que $p(x')=x$ et $f^{S'}(x')=y'$ , d'où le corollaire .

7 . \emph{Surjections et injections . Fibres .}

Définition 6 . - On dit qu'un morphisme $(\psi, \theta) : (X, \mathcal{O}_X) \to (Y, \mathcal{O}_Y)$ de préschémas est surjectif , ou est une surjection , si l'application continue $\psi$ est surjective .

On notera que cela n'entraîne nullement que $(\psi, \theta)$ soit un épimorphisme d'espaces annelés ($\theta$ n'est pas nécessairement injectif) .

Il est clair que le composé $g \circ f$ de deux surjections est une surjection , et que si inversement $g \circ f$ est une surjection , alors $g$ est une surjection .

Proposition 11 . - a) Si un $S$-morphisme $f : X \to Y$ est surjectif , tout $S'$-morphisme $f^{S'} : X^{S'} \to Y^{S'}$ déduit de $f$ par extension du préschéma de base est surjectif .

b) Si $f : X \to X'$ , $g : Y \to Y'$ sont des $S$-morphismes surjectifs , le $S$-morphisme $f \times_S g : X \times_S Y \to X' \times_S Y'$ est surjectif .

La première assertion est une conséquence immédiate du cor. de la prop. 10 du n° 6 . La seconde s'en déduit en remarquant que $f \times_S g$ est le morphisme composé
$$ X \times_S Y \xrightarrow{f \times 1_Y} X' \times_S Y \xrightarrow{1_{X'} \times g} X' \times_S Y' \quad . $$

Proposition 12 . - Pour qu'un morphisme $f : X \to Y$ soit surjectif , il faut et il suffit que pour tout corps $K$ et tout morphisme $\operatorname{Spec}(K) \to Y$ , il existe une extension $K'$ de $K$ et un morphisme $\operatorname{Spec}(K') \to X$ rendant commutatif le diagramme
$$
\begin{tikzcd}
X \arrow[d, "f"] & \operatorname{Spec}(K') \arrow[l] \arrow[d] \\
Y & \operatorname{Spec}(K) \arrow[l]
\end{tikzcd}
$$

La condition est suffisante , car pour tout $y \in Y$ , il suffit de l'appliquer à un morphisme $\operatorname{Spec}(K) \to Y$ correspondant à un monomorphisme $\kappa(y) \to K$ , $K$ étant une extension de $\kappa(y)$ (n° 2 , cor. de la prop. 5) .

\newpage

%% ===== Page 104 =====

Inversement, supposons $f$ surjectif, et soit $y \in Y$ l'image de l'unique élément de $\operatorname{Spec}(K)$ ; il existe $x \in X$ tel que $f(x)=y$, donc $f$ définit un monomorphisme $\kappa(y) \to \kappa(x)$ ; il suffit alors de prendre $K'$ extension de $\kappa(y)$, tel qu'il existe des $\kappa(y)$-monomorphismes de $\kappa(x)$ et $K$ dans $K'$ ; le morphisme $\operatorname{Spec}(K') \to X$ correspondant à $\kappa(x) \to K'$ répond à la question.

Avec le langage introduit au n°6, on peut dire encore que tout point géométrique de $Y$ à valeurs dans $K$ provient d'un point géométrique de $X$ à valeurs dans une extension de $K$.

Définition 7. — On dit qu'un morphisme de préschémas $f : X \to Y$ est \emph{géométriquement injectif}, ou est une \emph{injection géométrique}, si pour tout corps $K$, l'application $X(K) \to Y(K)$ est injective.

Pour que $f$ soit géométriquement injectif, il suffit que la condition de la déf. 7 soit vérifiée pour tout corps algébriquement clos. En effet, si $K$ est un corps quelconque, $K'$ une extension algébriquement close de $K$, le diagramme
$$\begin{tikzcd}
X(K) \arrow[r, "\alpha"] \arrow[d, "\varphi"] & Y(K) \arrow[d, "\varphi'"] \\
X(K') \arrow[r, "\alpha'"] & Y(K')
\end{tikzcd}$$
est commutatif, $\varphi$ et $\varphi'$ provenant du morphisme $\operatorname{Spec}(K') \to \operatorname{Spec}(K)$, et $\alpha, \alpha'$ correspondant à $f$. Or $\varphi$ est injectif et il en est de même de $\varphi'$ par hypothèse, donc $\alpha$ est nécessairement injectif.

Proposition 13. — Soient $f : X \to Y$, $g : Y \to Z$ deux morphismes : si $f$ et $g$ sont géométriquement injectifs, il en est de même de $g \circ f$ ; inversement, si $g \circ f$ est géométriquement injectif, il en est de même de $f$.

Compte tenu de la définition, la proposition revient aux assertions correspondantes pour les applications $X(K) \to Y(K) \to Z(K)$, qui sont évidentes.

Proposition 14. — Si les S-morphismes $f : X \to X'$, $g : Y \to Y'$ sont géométriquement injectifs, il en est de même de $f \times_S g$.

\newpage

%% ===== Page 105 =====

Inversement , supposons $f$ surjectif , et soit $y \in Y$ l'image de l'unique élément de $\operatorname{Spec}(K)$ ; il existe $x \in X$ tel que $f(x)=y$, donc $f$ définit un monomorphisme $\kappa(y) \to \kappa(x)$ ; il suffit alors de prendre $K'$ extension de $\kappa(y)$, tel qu'il existe des $\kappa(y)$-monomorphismes de $\kappa(x)$ et $K$ dans $K'$ ; le morphisme $\operatorname{Spec}(K') \to X$ correspondant à $\kappa(x) \to K'$ répond à la question .

Avec le langage introduit au $n^\circ 6$ , on peut dire encore que tout point géométrique de $Y$ provient d'un point géométrique de $X$ à valeurs dans une extension de $K$.

Définition 7 .- On dit qu'un morphisme de préschémas $f : X \to Y$ est géométriquement injectif , ou est une injection géométrique , si pour tout corps $K$ , l'application $X(K) \to Y(K)$ est injective .

Pour que $f$ soit géométriquement injectif , il suffit que la condition de la déf. 7 soit vérifiée pour tout corps $K$ algébriquement clos . En effet , si $k$ est un corps quelconque , $K'$ une extension algébriquement close de $k$ , le diagramme
$$\begin{tikzcd}
X(K) \arrow[r, "\alpha"] \arrow[d, "\phi"] & Y(K) \arrow[d, "\phi'"] \\
X(K') \arrow[r, "\alpha'"] & Y(K')
\end{tikzcd}$$
est commutatif , $\phi$ et $\phi'$ provenant du morphisme $\operatorname{Spec}(K') \to \operatorname{Spec}(K)$ , et $\alpha, \alpha'$ correspondant à $f$ . Or $y$ est injectif et il en est de même de $\alpha'$ par hypothèse , donc $\alpha$ est nécessairement injectif .

Proposition 13 .- Soient $f : X \to Y$ , $g : Y \to Z$ deux morphismes ; si $f$ et $g$ sont géométriquement injectifs , il en est de même de $gf$ ; inversement , si $gf$ est géométriquement injectif , il en est de même de $f$.

Compte tenu de la définition , la proposition revient aux assertions correspondantes pour les applications $X(K) \to Y(K) \to Z(K)$, qui sont évidentes .

Proposition 14 .- Si les $S$-morphismes $f : X \to X'$ , $g : Y \to Y'$ sont géométriquement injectifs , il en est de même de $f \times_S g$.

\newpage

%% ===== Page 106 =====

- xii -

phisme canonique $A \to S^{-1}A$ ; on sait que $\varphi$ est un homéomorphisme de $\operatorname{Spec}(S^{-1}A)$ sur un sous-espace de $\operatorname{Spec}(A)$ ($\S$ 1, n$^\circ$ 2, cor. 3 de la prop. 5), et d'autre part $\varphi_x$ donne par passage aux quotients un isomorphisme des corps résiduels pour tout $x \in \varphi(y)$.

\textbf{Proposition 16.} --- Soient $f : X \to Y$ une injection géométrique, et soit $g : Y' \to Y$ un morphisme ; posons $X' = X \times_Y Y'$. Le morphisme $f' : X' \to Y'$ est géométriquement injectif, et l'image par $g'$ de l'espace de base de $X'$ est l'image réciproque $g^{-1}(f(X))$.

La première assertion a déjà été démontrée (cor. de la prop. 14) ; en outre (n$^\circ$ 6, cor. de la prop. 10), les points $y' \in Y'$ qui sont images de points de $X'$ sont ceux pour lesquels il existe $x \in X$ tels que $f(x) = g(y')$.

\textbf{Corollaire.} --- Pour tout corps $K$, l'ensemble $X'(K)$ s'identifie au sous-ensemble de $Y'(K)$ image réciproque du sous-ensemble $X(K)$ de $Y(K)$ par l'application $Y'(K) \to Y(K)$ correspondant à $g$.

Cela résulte de la commutativité du diagramme
\begin{center}
\begin{tikzcd}
X'(K) \arrow[r] \arrow[d] & Y'(K) \arrow[d] \\
X(K) \arrow[r] & Y(K)
\end{tikzcd}
\end{center}

\textbf{Proposition 17.} --- Soient $f : X \to Y$ un morphisme, $y$ un point de $Y$. La projection $p : X \times_Y \operatorname{Spec}(\kappa(y)) \to X$ est un homéomorphisme de l'espace de base du $\kappa(y)$-préschéma $X \times_Y \operatorname{Spec}(\kappa(y))$ sur la fibre $f^{-1}(y)$, munie de la topologie induite par celle de l'espace de base de $X$.

On sait déjà, puisque $\operatorname{Spec}(\kappa(y)) \to Y$ est géométriquement injective (cor. 1 de la prop. 15) que $p$ identifie l'espace de base de $X \times_Y \operatorname{Spec}(\kappa(y))$ à $f^{-1}(y)$ en tant qu'ensemble (prop. 16) ; tout revient à démontrer que $p$ est un homéomorphisme. Comme on peut remplacer $Y$ par tout ouvert contenant $y$ (n$^\circ$ 4, prop. 7), la question est locale sur $X$ et $Y$, et on peut donc supposer $X = \operatorname{Spec}(A)$ et $Y = \operatorname{Spec}(B)$ affines, $A$ étant donc une $B$-algèbre. La proposition résulte alors du cor. 2 de la prop. 5, $\S$ 1, n$^\circ$ 2, appliqué à l'homomorphisme $A \to A \otimes_B \kappa(y)$.

Par la suite, on considérera toujours les fibres $f^{-1}(y)$ d'un morphis-

\newpage

%% ===== Page 107 =====

- xiii -

$f: X \to Y$ comme muni de la structure de $\kappa(y)$-préschéma transportée par la projection $p$. Le cor. de la prop. 16 montre en outre que les points de $X$ à valeurs dans une extension $K$ de $\kappa(y)$ sont identifiés aux points de $f^{-1}(y)$ à valeurs dans $K$.

("transitivité des fibres")

Proposition 18. — Soient $f : X \to Y, g : Y' \to Y$ deux morphismes ; posons $X' = X \times_Y Y'$, et $f' = f^{Y'} : X' \to Y'$. Pour tout $y' \in Y'$, le préschéma $f'^{-1}(y') = f^{-1}(y) \otimes_{\kappa(y)} \kappa(y')$, en posant $y = g(y')$.

Si on revient aux définitions, cela revient à montrer que les deux préschémas $(X \times_Y \operatorname{Spec}(\kappa(y))) \times_{\operatorname{Spec}(\kappa(y))} \operatorname{Spec}(\kappa(y'))$ et $(X \times_Y Y') \times_{Y'} \operatorname{Spec}(\kappa(y'))$ sont canoniquement isomorphes, ce qui résulte de la transitivité de l'extension du préschéma de base ($n^\circ 5$, prop. 9) et de la commutativité du diagramme
$$
\begin{tikzcd}
\operatorname{Spec}(\kappa(y)) & \operatorname{Spec}(\kappa(y')) \\
Y \arrow[u] & Y' \arrow[u]
\end{tikzcd}
$$

Proposition 19. — Soient $f : X \to Y$ un morphisme, $y$ un point de $Y$. La projection $p : X \times_Y \operatorname{Spec}(\mathcal{O}_y) \to X$ est un homéomorphisme de l'espace de base de $X \times_Y \operatorname{Spec}(\mathcal{O}_y)$ sur l'image réciproque par $f$ de l'espace de base de $\operatorname{Spec}(\mathcal{O}_y)$, munie de la topologie induite par celle de l'espace de base de $X$.

Comme dans la prop. 17, tout revient à prouver que $p$ est un homéomorphisme (cor. 1 de la prop. 15 et prop. 16). Comme l'image de $\operatorname{Spec}(\mathcal{O}_y)$ dans $Y$ est contenue dans tout ouvert affine contenant $y$ ($n^\circ 2$), on se ramène encore comme dans la prop. 17 au cas où $X = \operatorname{Spec}(A)$ et $Y = \operatorname{Spec}(B)$, $A$ étant une algèbre sur $B$. Alors $X \times_Y \operatorname{Spec}(\mathcal{O}_y)$ est le spectre de $A \otimes_B \mathcal{O}_y$, qui s'identifie à $S^{-1}A$, où $S$ est l'image dans $A$ de $B - \mathfrak{p}_y$ (chap. 0, $\S 1, n^\circ 2$) ; et on sait que $\operatorname{Spec}(S^{-1}A) \to \operatorname{Spec}(A)$ est un homéomorphisme sur un sous-espace de $X$ ($\S 1, n^\circ 2$, cor. 3 de la prop. 5).

\newpage

%% ===== Page 108 =====

- xiv -

P. 45, replace définition 2 by the following :

Définition 2. — On dit qu'un morphisme $f : Y \to X$ est une immersion (resp. une immersion fermée, une immersion ouverte) s'il se factorise en un isomorphisme $g : Y \to Z$ sur un sous-préschéma $Z$ de $X$ (resp. un sous-préschéma fermé $Z$ de $X$, un sous-préschéma induit sur un ouvert $Z$ de $X$) et en le morphisme d'injection $j : Z \to X$.

P. 45, replace prop. 4 by the following :

Proposition 4. — a) Pour qu'un morphisme $f = (\varphi, \theta) : Y \to X$ soit une immersion ouverte, il faut et il suffit que $\varphi$ soit un homéomorphisme de $Y$ sur une partie ouverte de $X$, et que pour tout $y \in Y$, l'homomorphisme $\theta_y : \mathcal{O}_{X, \varphi(y)} \to \mathcal{O}_{Y, y}$ soit bijectif.

b) Pour qu'un morphisme $f = (\varphi, \theta) : Y \to X$ soit une immersion (resp. une immersion fermée), il faut et il suffit que $\varphi$ soit un homéomorphisme de $Y$ sur une partie localement fermée (resp. fermée) de $X$, et que pour tout $y \in Y$, l'homomorphisme $\theta_y : \mathcal{O}_{X, \varphi(y)} \to \mathcal{O}_{Y, y}$ soit surjectif.

a) Les conditions sont évidemment nécessaires. Inversement, si elles sont remplies, il est clair que $\theta^b$ est un isomorphisme de $\mathcal{O}_Y$ sur $\varphi^*(\mathcal{O}_X)$, et que $\varphi^*(\mathcal{O}_X)$ est le faisceau déduit par transport de structure au moyen de $\varphi^{-1} \theta$, à partir de $\mathcal{O}_X / \varphi(Y)$ ; d'où la conclusion.

b)...

XX P. 47 in cor. 1, p. 48 in cor. 3 and prop. 5, p. 50 in cor. 1, introduce in the "resp." the mention of "immersions ouvertes" ("préschéma induit sur un ouvert" in cor. 3, p. 50). P. 47, in the proof of cor. 1, line 1 of the proof, after "bijectif", insert : (resp. bijectif s'il s'agit d'immersions ouvertes) ; in the other "resp." of the proof, introduce "ouvert". In the proof of cor. 3, p. 50, in the "resp." introduce also "une immersion ouverte".

P. 51, insert a new n°3 :

\newpage

%% ===== Page 109 =====

3. Immersions locales et isomorphismes locaux.

Définition 3. — Soit $f: X \to Y$ un morphisme de préschémas. On dit que $f$ est une immersion locale en un point $x$ de l'espace de base $X$ s'il existe un voisinage ouvert $U$ de $x$ dans $X$ et un voisinage ouvert $V$ de $f(x)$ dans $Y$ tels que la restriction de $f$ au sous-préschéma induit $U$ soit une immersion fermée de $U$ dans le sous-préschéma induit $V$. On dit que $f$ est une immersion locale si $f$ est une immersion locale en tout point $x$ de $X$.

Définition 4. — On dit qu'un morphisme $f: X \to Y$ est un isomorphisme local en un point $x \in X$ s'il existe un voisinage ouvert $U$ de $x$ dans $X$ tel que la restriction de $f$ au sous-préschéma induit $U$ soit une immersion ouverte de $U$ dans $Y$. On dit que $f$ est un isomorphisme local si $f$ est un isomorphisme local en tout point de $X$.

Une immersion (resp. une immersion fermée) $f: X \to Y$ peut donc se caractériser comme une immersion locale telle que $f$ soit un homéomorphisme de l'espace de base $X$ sur une partie localement fermée (resp. fermée) de $Y$. Une immersion ouverte peut se caractériser comme un isomorphisme local tel que $f$ soit un homéomorphisme de l'espace de base $X$ sur une partie ouverte de $Y$.

Proposition 6. — Le composé de deux immersions locales (resp. de deux isomorphismes locaux) est une immersion locale (resp. un isomorphisme local).

Cela résulte aussitôt de la transitivité des immersions (resp. fermées, ouvertes) (n°2, cor. 3 de la prop. 4) et du fait que si $f$ est un homéomorphisme de $X$ sur une partie fermée de $Y$, pour tout ouvert $U \subset X$, $f(U)$ est ouvert dans $f(X)$, donc il existe une partie ouverte $V$ de $Y$ telle que $f(U) = V \cap f(X)$, et $f(U)$ est donc fermé dans $V$.

Proposition 7. — Soient $f: X \to X'$, $g: Y \to Y'$ deux $S$-morphismes. Si $f$ et $g$ sont des immersions locales (resp. des isomorphismes locaux), il en est de même de $f \times g$.

\newpage

%% ===== Page 110 =====

- xvi -

En effet, soient $z \in X \times_S Y$, $z' \in X' \times_{S'} Y'$ par $f \times g$, $p, q$ les projections de $X \times_S Y$, $p', q'$ celles de $X' \times_{S'} Y'$. Il existe par hypothèse des voisinages ouverts $U, U', V, V'$ de $z=p(z), x'=p'(z')$, $y=q(z), y'=q'(z')$ respectivement, tels que les restrictions de $f$ et $g$ à $U$ et $V$ soient des immersions fermées (resp. ouvertes) dans $U'$ et $V'$ respectivement. Comme l'espace de base de $U \times_S V$ et celui de $U' \times_{S'} V'$ s'identifient aux voisinages ouverts $p^{-1}(U) \cap q^{-1}(V)$ et $p'^{-1}(U') \cap q'^{-1}(V')$ de $z$ et $z'$ respectivement, la proposition résulte de la prop. 5 du n° 2.

\emph{Corollaire} .- Si un $S$-morphisme $f$ est une immersion locale (resp. un isomorphisme local), il en est de même de $f_{S'}$ pour toute extension $S'$ du préschéma de base.

From p. 51 on, section $n$ is replaced by section $n+1$, definition $n$ (resp. prop. $n$) by definition $n+2$ (resp. prop. $n+2$).

p. 53, prop. 10, in the statement of the proposition and in the proof introduce in the "resp." the notion of "immersion ouverte", and in the last line but one of the proof, after "surjectif" insert: (resp. bijectif s'il s'agit d'immersions ouvertes).

p. 57, before the final "Remarque" of section 5, insert :

\emph{Corollaire} .- Soient $f : X \to Y$, $g : Y \to Z$ deux morphismes ; si $g \circ f$ est une immersion (resp. une immersion locale), il en est de même de $f$.

En effet, $f$ se factorise en $X \xrightarrow{\Gamma_f} X \times_Z Y \xrightarrow{p_2} Y$. D'autre part, $p_2$ s'identifie à $(g \circ f) \times_Z 1_Y$ ($\S$ 2, n° 5, cor. de la prop. 8) ; si $g \circ f$ est une immersion (resp. une immersion locale) il en est de même de $p_2$ (n° 2, prop. 5 et n° 3, prop. 7) ; en vertu du cor. 2, $p_2 \circ \Gamma_f$ est une immersion (resp. une immersion locale) (n° 2, cor. 3 de la prop. 4 et n° 3, prop. 6).

Le même raisonnement vaudrait pour les immersions fermées si on savait que $\Gamma_f$ est une immersion fermée ; on va introduire au n° 6 une condition assurant qu'il en est bien ainsi.

\newpage

%% ===== Page 111 =====

- xvii -

P. 61, replace cor. 1 by :

\emph{Corollaire 1.} — Soient $f : X \to Y$, $g : Y \to Z$ deux morphismes, $g$ étant séparé. Si $g \circ f$ est une immersion fermée, il en est de même de $f$.

Le raisonnement est le même que dans le cor. 3 de la prop. 15 du n$^\circ 5$, tenant compte de ce que $\Gamma_f$ est une immersion fermée, et du cor. 3 de la prop. 4 du n$^\circ 2$.

Suppress the "Remarque" on top of p. 62.

P. 51, before n$^\circ 3$, add :

La commutativité du diagramme
$$\begin{tikzcd}
X \times_{Y'} Y' \arrow[r, "f^{Y'}"] \arrow[d, "p"] & Y' \arrow[d, "j"] \\
X \arrow[r, "f"] & Y
\end{tikzcd}$$
montre que la restriction de $f$ au sous-préschéma $f^{-1}(Y')$ se factorise en $f^{-1}(Y') \xrightarrow{g} X \times_{Y'} Y' \xrightarrow{f^{Y'}} Y' \xrightarrow{j} Y$, en désignant par $g$ l'isomorphisme réciproque de l'isomorphisme associé à $p = 1_X \times_Y j$. En outre :

\emph{Corollaire 4.} — Soient $f : X \to Y$, $g : Y \to Z$ deux morphismes, $h = g \circ f$ leur composé. Pour tout sous-préschéma $Z'$ de $Z$, il existe un isomorphisme canonique de $h^{-1}(Z')$ sur $g^{-1}(f^{-1}(Z'))$.

Cela résulte aussitôt de l'isomorphisme $X \times_Y (Y \times_Z Z') = X \times_Z Z'$ (§ 2, n$^\circ 5$, prop. 9).

\newpage

%% ===== Page 112 =====

- xvii -

P.61, replace cor.1 by :

Corollaire 1. - Soient $f: X \to Y$, $g: Y \to Z$ deux morphismes, $g$ étant séparé. Si $g \circ f$ est une immersion fermée, il en est de même de $f$.

Le raisonnement est le même que dans le cor.3 de la prop.15 du n°5, tenant compte de ce que $\Gamma_f$ est une immersion fermée, et du cor.3 de la prop.4 du n°2.

Suppress the "Remarque" on top of p.62.

P.51, before n°3, add :

La commutativité du diagramme
$$\begin{tikzcd}
X \times_Y Y' \arrow[r, "f^{Y'}"] \arrow[d, "p"] & Y' \arrow[d, "j"] \\
X \arrow[r, "f"] & Y
\end{tikzcd}$$
montre que la restriction de $f$ au sous-préschéma $f^{-1}(Y')$ se factorise en $f^{-1}(Y') \to X \times_Y Y' \xrightarrow{f^{Y'}} Y' \xrightarrow{j} Y$, en désignant par $g$ l'isomorphisme réciproque de l'isomorphisme associé à $p = 1_X \times_Y j$. En outre :

Corollaire 4. - Soient $f: X \to Y$, $g: Y \to Z$ deux morphismes, $h = g \circ f$ leur composé. Pour tout sous-préschéma $Z'$ de $Z$, il existe un isomorphisme canonique de $h^{-1}(Z')$ sur $g^{-1}(f^{-1}(Z'))$.

Cela résulte aussitôt de l'isomorphisme $X \times_Y (Y \times_Z Z') = X \times_Z Z'$ (§ 2, n°5, prop.9).

\newpage

%% ===== Page 113 =====

- xix -

P. 54, after prop. 11, add :

\emph{Corollaire}. - Les schémas $(X \times_S Y)_{\text{red}}$ et $(X_{\text{red}} \times_{S_{\text{red}}} Y_{\text{red}})$ s'identifient canoniquement.

Cela résulte aussitôt des prop. 9 et 11.

P. 54, before prop. 12, add :

Si $f : T \to X$, $g : T \to Y$ sont deux $S$-morphismes, on vérifie aussitôt que l'on a $(f,g)_S = (f \times_S g) \circ \Delta_T$.

\newpage

%% ===== Page 114 =====

\marginpar{Archive\\ Grothendieck sept}

\section*{§ 5. Applications rationnelles}

\subsection*{1. Pseudo-morphismes et applications rationnelles.}

Soient $X, Y$ deux préschémas, $U, V$ deux ouverts denses dans $X$, $f$ (resp. $g$) un morphisme de $U$ (resp. $V$) dans $Y$ ; nous dirons que $f$ et $g$ sont équivalents s'ils coïncident dans un ouvert dense dans $U \cap V$. Comme une intersection finie d'ouverts denses dans $X$ est dense, il est clair que cette relation est bien une relation d'équivalence.

Définition 1. --- Etant donnés deux préschémas $X, Y$, on appelle pseudo-morphisme de $X$ dans $Y$ une classe d'équivalence de morphismes de parties ouvertes denses de $X$ dans $Y$. Si $X$ et $Y$ sont des $S$-préschémas, on dit qu'un pseudo-morphisme de $X$ dans $Y$ est un pseudo-$S$-morphisme s'il existe un représentant dans cette classe qui est un $S$-morphisme. On appelle pseudo-$S$-section d'un $S$-préschéma $X$ un pseudo-$S$-morphisme de $S$ dans $X$. On appelle pseudo-fonction sur $X$ toute pseudo-$X$-section du préschéma $X \times_{\mathbb{Z}} \mathbb{Z}[T]$ sur $X$ ($T$ indéterminée).

Les pseudo-$X$-sections sur $X$ s'identifient aux classes d'équivalence des sections du faisceau structural $\mathcal{O}_X$ au-dessus d'ouverts partout denses de $X$ (deux telles sections étant équivalentes si elles coïncident dans un ouvert partout dense contenu dans l'intersection de leurs ensembles de définition). En effet, soit $(V_\alpha)$ un recouvrement d'un ouvert partout dense $U$ de $X$ par des ouverts affines. Soit $u : U \to \mathbb{Z}[T]$ un $U$-morphisme, et soit $u_\alpha$ sa restriction à $V_\alpha = \operatorname{Spec}(A_\alpha)$ ; le préschéma $V_\alpha \times_{\mathbb{Z}} \mathbb{Z}[T]$ est affine d'anneau $A_\alpha[T]$, et $u_\alpha$ correspond à un $A_\alpha$-homomorphisme $A_\alpha[T] \to A_\alpha$ (II, n$^\circ$2, prop. 3). Or un tel homomorphisme est complètement déterminé par la donnée de l'image de $T$, soit $f_\alpha \in \Gamma(V_\alpha, \mathcal{O}_X)$, et si on écrit que les restrictions de $u_\alpha$ et $u_\beta$ à un ouvert affine $W \subset V_\alpha \cap V_\beta$ coïncident, on voit aussitôt que $f_\alpha$ et $f_\beta$ coïncident dans $W$, donc la famille

\newpage

%% ===== Page 115 =====

(80) -

$(f_\alpha)$ est formée des restrictions aux $V_\alpha$ d'une section $f$ de $\mathcal{O}_X$ au-dessus de $U$, et réciproquement, il est clair qu'une telle section détermine un $U$-morphisme $U \to \mathbb{Z}[T]$, d'où notre assertion (ces considérations seront généralisées au chap. II, § 2, n°8). On voit donc que les pseudo-fonctions sur $X$ forment un anneau.

Lorsque $X$ est irréductible, tout ouvert non vide est dense dans $X$; on peut encore dire que les ouverts non vides de $X$ sont les voisinages ouverts du point générique $x$ de $X$. Dire que deux morphismes de parties ouvertes non vides de $X$ dans $Y$ sont équivalents signifie donc dans ce cas qu'ils ont même germe au point $x$. Autrement dit les pseudo-morphismes $X \to Y$ sont les germes de morphismes de $X$ dans $Y$ au point générique $x$ de $Y$. En particulier :

Proposition 1. - Si $X$ est un préschéma irréductible, l'anneau des pseudo-fonctions sur $X$ s'identifie à l'anneau local $\mathcal{O}_{X,x}$ du point générique $x$ de $X$. C'est un anneau local de dimension 0, et par suite un anneau local artinien lorsque $X$ est noethérien ; c'est un corps lorsque $X$ est intègre.

Nous avons en effet démontré la première assertion ; les autres se déduisent du § 2, n°2, cor. de la prop. 4 et des propriétés des anneaux locaux de dimension 0 (Bourbaki, Alg. comm.).

Corollaire. - Si $X = \operatorname{Spec}(A)$, où $A$ est intègre, le corps des pseudo-fonctions sur $X$ s'identifie au corps des fractions de $A$.

Supposons maintenant que $X$ ait un nombre fini de composantes irréductibles $X_i$ ($1 \le i \le n$) ; soit $X'_i$ l'ouvert dense de $X$, complémentaire de la réunion des $X_i \cap X_j$ ($j \ne i$). Il est clair que pour tout ouvert partout dense $U$ de $X$, $U_i = U \cap X'_i$ est un ouvert non vide de $X_i$ ($1 \le i \le n$) et les $U_i$ sont deux à deux sans point commun. Or, se donner un morphisme de $U = \bigcup U_i$ dans $Y$ revient à se donner (arbitrairement) un morphisme de chacun des $U_i$ dans $Y$. Autrement dit :

\newpage

%% ===== Page 116 =====

- 81 -

\noindent Proposition 2. — Soient $X, Y$ deux préschémas (resp. S-préschémas), tels que $X$ ait un nombre fini de composantes irréductibles $X_i$. L'ensemble des pseudo-morphismes (resp. S-pseudo-morphismes) de $X$ dans $Y$ s'identifie au produit des ensembles de pseudo-morphismes (resp. pseudo-S-morphismes) des $X_i$ dans $Y$.

\noindent Corollaire 1. — Soit $X$ un préschéma noethérien. Alors l'anneau des pseudo-fonctions sur $X$ est un anneau artinien, dont les composantes locaux sont les anneaux de pseudo-fonctions sur les composantes irréductibles de $X$.

\noindent Corollaire 2. — Soient $A$ un anneau noethérien, $X = \operatorname{Spec}(A)$. L'anneau des pseudo-fonctions sur $X$ s'identifie canoniquement à $S^{-1}A$, où $S$ est le complémentaire de la réunion des idéaux premiers minimaux de $A$.

\noindent Cela résultera du lemme suivant :

\noindent Lemme 1. — Pour qu'un élément $f \in A$ soit tel que l'ouvert $D(f)$ soit partout dense, il faut et il suffit que $f \in S$, et tout ouvert dense dans $X$ contient un ouvert de la forme $D(f)$, où $f \in S$.

\noindent En effet, si ce lemme est démontré, comme l'anneau des sections $\Gamma(D(f), \mathcal{O}_X)$ s'identifie à $A_f$ (§ 1, n$^\circ$3, prop. 6 et th. 1), il résulte du fait que les $D(f)$ avec $f \in S$ forment un ensemble cofinal dans l'ensemble ordonné (par $\supset$) des ouverts denses dans $X$, et de la définition 1, que l'anneau des pseudo-fonctions sur $X$ s'identifie à la limite inductive des $A_f$ pour $f \in S$ (pour la relation d'ordre "$g$ est multiple de $f$"), c'est-à-dire à $S^{-1}A$ (chap. 0, § 1, n$^\circ$4).

\noindent Pour démontrer le lemme 1, remarquons que si $D(f)$ est dense dans $X$, $D(f) \cap X_i \neq \emptyset$ pour $1 \leq i \leq n$ (et réciproquement, mais cela signifie que $f \notin \mathfrak{p}_i$, si $\mathfrak{p}_i = \mathfrak{j}(X_i)$), donc $f \in S$, d'où la première assertion du lemme. D'autre part si $U$ est un ouvert dense dans $X$, $U \cap X_i \neq \emptyset$ pour tout $i$, donc il existe $f_i \in A$ tel que $D(f_i) \subset U \cap X_i$ ($1 \leq i \leq n$). Par définition, cela signifie...

\newpage

%% ===== Page 117 =====

- 82 -

le complémentaire de $U$ est de la forme $V(\mathfrak{a})$, où $\mathfrak{a}$ est un idéal qui n'est contenu dans aucun des $\mathfrak{p}_i$ ; il n'est donc pas contenu dans leur réunion, et il y a par suite $f \in \mathfrak{a}$ appartenant à $S$ ; d'où $D(f) \subset U$, ce qui achève la démonstration.

Supposons de nouveau $X$ irréductible, de point générique $x$. Comme tout ouvert non vide $U$ de $X$ contient $x$, et par suite contient aussi tout $z \in \overline{\{x\}}$, tout morphisme $U \to Y$ peut se composer avec le morphisme canonique $\operatorname{Spec}(\mathcal{O}_x) \to X$ ($\S$ 2, n°2) et deux morphismes de parties ouvertes non vides de $X$ qui coïncident dans une partie ouverte non vide de $X$ donnent le même morphisme $\operatorname{Spec}(\mathcal{O}_x) \to Y$ par composition. Autrement dit, à tout pseudo-morphisme $X \to Y$ correspond ainsi un morphisme $\operatorname{Spec}(\mathcal{O}_x) \to Y$ bien déterminé.

Proposition 3. - Soient $X$ et $Y$ deux $S$-préschémas ; on suppose $X$ irréductible, de point générique $x$, et $Y$ de type fini sur $S$. Soit $\xi = \operatorname{Spec}(\mathcal{O}_x)$. Deux pseudo-$S$-morphismes de $X$ dans $Y$ auxquels correspond le même $S$-morphisme $\xi \to Y$ sont identiques. Inversement, si $S$ est localement noethérien, tout $S$-morphisme de $\xi$ dans $Y$ correspond à un pseudo-$S$-morphisme de $X$ dans $Y$.

Compte tenu de ce que tout ouvert non vide dans $X$ est partout dense, cela résulte aussitôt de la prop. 15 du $\S$ 4, n°3.

Corollaire 1. - On suppose $S$ localement noethérien, et les autres hypothèses de la prop. 3 satisfaites. Les pseudo-morphismes de $X$ dans $Y$ s'identifient alors aux points du $S$-préschéma $Y$, à valeurs dans le $S$-préschéma $\xi = \operatorname{Spec}(\mathcal{O}_x)$.

Cela n'est autre que la prop. 3, avec la terminologie introduite au $\S$ 2, n°6.

Corollaire 2. - On suppose remplies les conditions du cor. 1. Soit $s$ l'image de $x$ dans $S$. La donnée d'un pseudo-$S$-morphisme $X \to Y$ équivaut à la donnée d'un point $y$ de $Y$ au-dessus de $s$, et d'un $\mathcal{O}_s$-homo-

\newpage

%% ===== Page 118 =====

- 83 -

morphisme de $\mathcal{O}_y$ dans $R = \mathcal{O}_x$ tel que l'image réciproque de l'idéal maximal de $\mathcal{O}_x$ soit l'idéal maximal de $\mathcal{O}_y$.

Cela résulte du cor. 1, ainsi que du § 2, n° 2, prop. 5.

En particulier :

Corollaire 3. — Sous les mêmes conditions, les pseudo-$S$-morphismes de $X$ dans $Y$ ne dépendent que du $S$-préschéma $\xi = \operatorname{Spec}(\mathcal{O}_x)$ et restent les mêmes en particulier quand on remplace $X$ par $\operatorname{Spec}(\mathcal{O}_z)$, où $z \in X$.

En effet, comme $z \in \{x\}$, $x$ appartient à $\operatorname{Spec}(\mathcal{O}_z)$ et en est par suite le point générique.

Lorsque $X$ est intègre, $R = \mathcal{O}_x$ est un corps ; les corollaires précédents se spécialisent en :

Corollaire 4. — On suppose vérifiées les conditions du cor. 1 et en outre que $X$ est intègre. Alors les pseudo-$S$-morphismes de $X$ dans $Y$ s'identifient aux points géométriques de $Y \otimes_S \kappa(s)$ à valeurs dans l'extension $R$ de $\kappa(s)$, autrement dit chacun d'eux équivaut à la donnée d'un point $y \in Y$ au-dessus de $s$ et d'un $\kappa(s)$-homomorphisme de $\kappa(y)$ dans $R = \kappa(x)$.

Les points de $Y$ au-dessus de $s$ s'identifient en effet à ceux de $Y \otimes_S \kappa(s)$ (§ 2, n° 7, prop. 17) et les $\mathcal{O}_y$-homomorphismes $\mathcal{O}_y \to R$ vérifiant la condition du cor. 2 aux $\kappa(s)$-homomorphismes $\kappa(y) \to R$.

Plus particulièrement :

Corollaire 5. — Soient $k$ un corps, $X, Y$ deux préschémas algébriques sur $k$ ; on suppose $X$ intègre, et on désigne par $R$ le corps des pseudo-fonctions sur $X$. Alors les pseudo-$k$-morphismes de $X$ dans $Y$ s'identifient aux points géométriques de $Y$ à valeurs dans l'extension $R$ de $k$.

\newpage

%% ===== Page 119 =====

- 84 -

2. Applications rationnelles .

Définition 2. — Soient $X$ un $S$-préschéma réduit, $Y$ un $S$-préschéma. On appelle $S$-application rationnelle de $X$ dans $Y$ un $S$-morphisme $f$ d'un ouvert partout dense $U \subset X$ dans $Y$ qui ne peut être prolongé en un morphisme (dans $Y$) d'un ouvert contenant strictement $U$. On dit que $U$ est le domaine de définition de $f$, et que $f$ est définie en $x$ si $x \in U$. On appelle $S$-section rationnelle de $X$ une $S$-application rationnelle de $S$ dans $X$. On appelle fonction rationnelle sur $X$ toute $X$-section rationnelle du préschéma $X \otimes_{\mathbb{Z}} T$ ($T$ indéterminée).

Comme au n°1, on voit aussitôt qu'une fonction rationnelle sur $X$ s'identifie à une section de $\mathcal{O}_X$ au-dessus d'un ouvert dense $U$, qui ne peut être prolongée à un ouvert strictement plus grand.

Il est clair que tout $S$-morphisme $X \to Y$ de $S$-préschémas est une $S$-application rationnelle.

Lorsque $X$ et $Y$ sont des préschémas sans préschéma de base explicité, une application rationnelle de $X$ dans $Y$ sera une $\mathbb{Z}$-application rationnelle par définition.

Proposition 4. — Soient $X, Y$ deux $S$-préschémas ; on suppose $X$ réduit et $Y$ séparé sur $S$. Soient $U_1, U_2$ deux ouverts partout denses de $X$, $f_1, f_2 : U_i \to Y$ deux $S$-morphismes ($i=1,2$) tels qu'il existe un ouvert $V \subset U_1 \cap U_2$, dense dans $X$ et dans lequel $f_1$ et $f_2$ coïncident. Alors $f_1$ et $f_2$ coïncident dans $U_1 \cap U_2$.

On peut évidemment se borner au cas où $U_1 = U_2 = U$. Comme $X$ est réduit, c'est le plus petit sous-préschéma fermé de $X$ majorant $V$ (§ 3, n°4, prop. 9). Soit alors $g = (f_1, f_2)_S : X \to Y \times_S Y$ ; comme par hypothèse la diagonale $\Delta = \Delta_Y(Y)$ est un sous-préschéma fermé de $Y \times_S Y$, $g^{-1}(\Delta)$ est un sous-préschéma fermé de $X$ (§ 3, n°2, cor. 3 de la prop. 5).

Si $h : V \to Y$ est la restriction commune de $f_1$ et $f_2$ à $V$, la restriction de $g$ à $V$ est $g' = (h, h)_S$ qui se factorise en $\Delta_Y \circ h$ ; comme $\Delta_Y(\Delta) = Y$,

\newpage

%% ===== Page 120 =====

- 85 -

on a $g^{-1}(T) = V$, et par suite $Z$ est un sous-préschéma fermé de $X$ induisant $V$, donc majorant $V$, et cela entraîne $Z = X$. De la relation $g^{-1}(T) = X$, on déduit que $g$ se factorise en $\Delta_Y \circ f$, où $f$ est un morphisme $X \to Y$, ce qui entraîne $f_1 = f_2 = f$.

Proposition 5. — Les hypothèses sur $X$ et $Y$ étant celles de la prop. 4, soit $U$ un ouvert partout dense de $X$ et soit $f : U \to Y$ un $S$-morphisme. Il existe alors une $S$-application rationnelle et une seule $F$ de $X$ dans $Y$ prolongeant $f$.

En effet, comme deux prolongements $f_1, f_2$ de $f$ à des ouverts $U_1, U_2$ contenant $U$ coïncident dans $U_1 \cap U_2$, la réunion des ouverts $V \supset U$ auxquels $f$ est prolongeable est le plus grand ouvert auquel $f$ est prolongeable, et ce prolongement est unique.

Corollaire 1. — Sous les conditions de la prop. 5, il y a correspondance biunivoque entre les $S$-morphismes $U \to Y$ et les $S$-applications rationnelles de $X$ dans $Y$, définies en tout point de $U$.

Corollaire 2. — Soient $S$ un schéma, $X$ un $S$-préschéma réduit, $Y$ un $S$-préschéma séparé sur $S$, $f : U \to Y$ un $S$-morphisme d'une partie ouverte dense de $X$ dans $Y$. Si $F$ est la $S$-application rationnelle de $X$ dans $Y$ prolongeant $f$, $F$ est un $S$-morphisme (et est donc la $S$-application rationnelle de $X$ dans $Y$ prolongeant $f$).

En effet, si $p : X \to S$, $q : Y \to S$ sont les morphismes structuraux, il suffit de prouver que $q \circ F = p$, ce qui résulte aussitôt de la prop. 4 appliquée à ces deux morphismes du domaine de définition de $F$ dans $S$.

Corollaire 3. — Soient $X, Y$ deux $S$-préschémas ; on suppose $X$ réduit, $X$ et $Y$ séparés sur $S$. Soient $p : Y \to X$ un $S$-morphisme (faisant de $Y$ un $X$-préschéma), $U$ un ouvert partout dense de $X$, $f$ une $U$-section de $Y$ ; alors l'application rationnelle $F$ de $X$ dans $Y$ prolongeant $f$ est une $X$-section rationnelle.

\newpage

%% ===== Page 121 =====

- 86 -

tionnelle de $Y$.

Il faut prouver que $p \circ f$ est l'identité dans le domaine de définition de $f$ ; puisque $X$ est séparé sur $S$, cela résulte de la prop. 4.

Pour toute $S$-application rationnelle de $X$ dans $Y$, désignons par $cl(f)$ le pseudo-$S$-morphisme de $X$ dans $Y$ auquel appartient $f$.

Proposition 6. - Sous les hypothèses de la prop. 4, l'application $f \to cl(f)$ est une bijection de l'ensemble des $S$-applications rationnelles de $X$ dans $Y$ sur l'ensemble des pseudo-$S$-morphismes de $X$ dans $Y$.

Cela résulte aussitôt de la prop. 5.

Sous les conditions envisagées, on identifiera le plus souvent l'ensemble des $S$-applications rationnelles de $X$ dans $Y$ et l'ensemble des pseudo-$S$-morphismes de $X$ dans $Y$, et en particulier les pseudo-fonctions sur $X$ et les fonctions rationnelles sur $X$.

Soient $X, Y$ deux $S$-préschémas, $X$ étant réduit et $Y$ séparé sur $S$. Soient $f$ une $S$-application rationnelle de $X$ dans $Y$, et soit $x$ un point de $X$ ; on peut composer $f$ avec le $S$-morphisme canonique $\operatorname{Spec}(\mathcal{O}_x) \to X$ pourvu que le domaine de définition de $f$ rencontre l'espace de base de $\operatorname{Spec}(\mathcal{O}_x)$, identifié à l'ensemble des $z \in X$ tels que $x \in \{z\}$. Ceci aura lieu dans les cas suivants :

$1^\circ$ $X$ est irréductible (donc intègre) ; car alors $x$ est adhérent au point générique de $X$.

$2^\circ$ $X$ est localement noethérien ; notre assertion résulte en effet alors du

Lemme 2. - Soient $X$ un préschéma localement noethérien, $x$ un point de $X$. Les composantes irréductibles de $\operatorname{Spec}(\mathcal{O}_x)$ sont les traces sur $\operatorname{Spec}(\mathcal{O}_x)$ des composantes irréductibles de $X$ contenant $x$. Pour qu'un ouvert $U \subset X$ soit tel que $U \cap \operatorname{Spec}(\mathcal{O}_x)$ soit dense dans $\operatorname{Spec}(\mathcal{O}_x)$, il faut et il suffit qu'il rencontre les composantes irréductibles de

\newpage

%% ===== Page 122 =====

- 87 -

$X$ contenant $x$ (ce qui a lieu en particulier si $U$ est dense dans $X$).

La seconde assertion découle évidemment de la première, et il suffit donc de démontrer celle-ci. Comme $\operatorname{Spec}(\mathcal{O}_x)$ est contenu dans tout ouvert affine contenant $x$, on peut supposer $X$ affine et noethérien, d'anneau $A$. Comme les idéaux premiers de $A_x$ correspondent biunivoquement aux idéaux premiers de $A$ contenus dans $\mathfrak{j}_x$, les idéaux premiers minimaux de $A_x$ correspondent aux idéaux premiers minimaux de $A$ contenus dans $\mathfrak{j}_x$, d'où le lemme.

Cela étant, supposons que l'on soit dans l'un des deux cas précités si $U$ est le domaine de définition de $f$, désignons par $f'$ l'application rationnelle de $\operatorname{Spec}(\mathcal{O}_x)$ dans $Y$ qui coïncide avec $f$ dans $U \cap \operatorname{Spec}(\mathcal{O}_x)$; nous dirons que cette application rationnelle correspond à $f$.

Proposition 7. - Soient $S$ un préschéma localement noethérien, $X$ un $S$-préschéma réduit, et $Y$ un $S$-préschéma séparé et de type fini sur $S$. On suppose en outre $X$ irréductible ou noethérien. Soient alors $f$ une $S$-application rationnelle de $X$ dans $Y$, $x$ un point de $X$. Pour que $f$ soit définie au point $x$, il faut et il suffit que l'application rationnelle $f'$ de $\operatorname{Spec}(\mathcal{O}_x)$ dans $Y$, correspondant à $f$, soit un morphisme.
($\operatorname{Spec}(\mathcal{O}_x)$ étant contenu dans tout ouvert contenant $x$)

La condition étant évidemment nécessaire, prouvons qu'elle est suffisante. D'après le § 4, n°3, prop. 15, il existe un voisinage ouvert $U$ de $x$ dans $X$ et un $S$-morphisme $g$ de $U$ dans $Y$, induisant $f'$ sur $\operatorname{Spec}(\mathcal{O}_x)$. Si $X$ est irréductible, on peut, en vertu de la prop. 5, supposer que $g$ est une $S$-application rationnelle, puisque $U$ est dense dans $X$. En outre, le point générique de $X$ appartient à $\operatorname{Spec}(\mathcal{O}_x)$ et au domaine de définition de $f$, donc $f$ et $g$ coïncident en ce point, et par suite dans un ensemble ouvert non vide de $X$ (§ 4, n°3, prop. 15).

\newpage

%% ===== Page 123 =====

Comme $f$ et $g$ sont des $S$-applications rationnelles, elles coïncident (prop.5) et par suite $f$ est définie en $x$. Si $X$ est localement noethérien, on peut supposer $U$ noethérien ; il n'y a alors qu'un nombre fini de composantes irréductibles $X_i$ de $X$ contenant $x$ (lame 2), et on voit comme ci-dessus que $f$ et $g$ coïncident dans un ensemble ouvert non vide de chacun des $X_i$. Soit alors $f_1$ le morphisme défini dans $U \cup \bigcup X_i$, égal à $g$ dans $U$ et à $f$ dans l'intersection de $\bigcup X_i$ et du domaine de définition de $f$. Comme $U \cup \bigcup X_i$ est dense dans $X$, $f_1$ est une extension de $f$ (prop.4) et est par suite définie en $x$.

3. \emph{Faisceau des fonctions rationnelles}.

Soit $X$ un préschéma. Pour tout ouvert $U \subset X$, désignons par $R(U)$ l'anneau des pseudo-fonctions sur $U$ (n°1, déf.1) ; c'est une $\Gamma(U, \mathcal{O}_X)$-algèbre ; en outre, si $V \subset U$ est un second ouvert de $X$, toute section de $\mathcal{O}_X$ au-dessus d'une partie ouverte partout dense de $U$ donne par restriction à $V$ une section au-dessus d'une partie ouverte partout dense de $V$, et si deux sections coïncident au-dessus d'une partie ouverte partout dense de $U$, leurs restrictions à $V$ coïncident au-dessus d'une partie ouverte partout dense de $V$. On définit donc ainsi un homomorphisme d'algèbres $\rho_{U,V} : R(U) \to R(V)$, et il est clair que si $U \supset V \supset W$ sont trois ouverts de $X$, on a $\rho_{U,W} = \rho_{V,W} \circ \rho_{U,V}$ ; les $R(U)$ définissent donc un préfaisceau d'algèbres sur $X$.

Définition 3. --- On appelle \emph{faisceau des pseudo-fonctions} sur un préschéma $X$ et on désigne par $R(X)$ le faisceau d'algèbres associé au préfaisceau formé des $R(U)$.

Lorsque $X$ est réduit, $R(U)$ s'identifie à l'anneau des fonctions rationnelles sur $U$ (n°2, prop.6) et on dit alors que $R(X)$ est le faisceau des fonctions rationnelles sur $X$.

Pour tout préschéma $X$ et tout ouvert $U \subset X$, il est clair que

\newpage

%% ===== Page 124 =====

- 89 -

le faisceau induit $\underline{R}(X)|U$ n'est autre que $\underline{R}(U)$.

Proposition 8. — Soit $X$ un préschéma n'ayant un nombre fini de composantes irréductibles $X_i$ ($1 \le i \le n$). Alors le préfaisceau défini par les $R(U)$ est égal au faisceau $\underline{R}(X)$. Ce dernier est quasi-cohérent ; pour tout ouvert $U \subset X$, $\Gamma(U, \underline{R}(X)) = R(U)$ s'identifie canoniquement au composé direct des anneaux locaux des points génériques des $X_i$ tels que $U \cap X_i \neq \emptyset$.

La dernière assertion résulte du n°1, cor. 1 de la prop. 2. On en déduit aussitôt que le préfaisceau $U \to R(U)$ satisfait à l'axiome (F 1) des faisceaux. Il satisfait aussi à l'axiome (F 2) : en effet, si $(V_\alpha)$ est un recouvrement ouvert d'un ouvert $U \subset X$, et si $s_\alpha \in R(V_\alpha)$ sont tels que les restrictions de $s_\alpha$ et $s_\beta$ à $V_\alpha \cap V_\beta$ coïncident pour tout couple d'indices, on en conclut que pour chaque indice $i$ tel que $U \cap X_i \neq \emptyset$, les composantes dans $R(X_i)$ de tous les $s_\alpha$ tels que $V_\alpha \cap X_i \neq \emptyset$ sont les mêmes ; désignant cette composante par $s_i$, il est clair que l'élément de $R(U)$ ayant les $s_i$ pour composantes (restriction $s_\alpha$ à chaque $V_\alpha$). Pour voir que $\underline{R}(X)$ est quasi-cohérent, on peut se limiter au cas où $X$ est affine, et si $x_i$ est le point générique de $X_i$ et $X = \operatorname{Spec}(A)$, il est immédiat alors que $\underline{R}(X) = \tilde{M}$, où $M$ est somme directe des $A$-modules $A_{x_i}$.

Corollaire 1. — Si $X$ est irréductible, tout $\underline{R}(X)$-module quasi-cohérent $\mathcal{F}$ est un faisceau simple.

Il suffit de montrer que tout point de $x$ admet un voisinage tel que $\mathcal{F}$ soit un faisceau simple (FAC, II, 1, 36, lemme 2), autrement dit on est ramené au cas où $X$ est affine ; on peut en outre supposer que $\mathcal{F}$ est alors le conoyau d'un homomorphisme $(\underline{R}(X))^{(I)} \to (\underline{R}(X))^{(J)}$, et tout revient donc à voir que $\underline{R}(X)$ est constant, ce qui est évident puisque $\Gamma(U, \underline{R}(X)) = R(X)$ pour tout ouvert non vide $U$.

\newpage

%% ===== Page 125 =====

- 90 -

Corollaire 2. — Si $X$ est irréductible, pour tout $\mathcal{O}_X$-module quasi-cohérent $\mathcal{F}$, $\mathcal{F} \otimes_{\mathcal{O}_X} \mathcal{R}(X)$ est un faisceau simple ; si en outre $X$ est réduit, $\mathcal{F} \otimes_{\mathcal{O}_X} \mathcal{R}(X)$ est isomorphe à un faisceau de la forme $(\mathcal{R}(X))^{(I)}$.

La seconde assertion résulte de ce que $\mathcal{R}(X)$ est alors un corps. Si $X$ est irréductible, le faisceau structural $\mathcal{O}_X$ est un sous-faisceau de $\mathcal{R}(X)$, car on définit un homomorphisme injectif $\mathcal{O}_X \to \mathcal{R}(X)$ en faisant correspondre à toute section $f$ de $\mathcal{O}_X$ au-dessus d'un ouvert $U$ (nécessairement partout dense) la pseudo-fonction $cl(f)$ correspondante. Pour tout $\mathcal{O}_X$-module $\mathcal{F}$, l'homomorphisme précédent définit donc un homomorphisme $\mathcal{F} = \mathcal{F} \otimes_{\mathcal{O}_X} \mathcal{O}_X \to \mathcal{F} \otimes_{\mathcal{O}_X} \mathcal{R}(X)$, qui sur chaque fibre n'est autre que l'homomorphisme $s \to s \otimes 1$ de $\mathcal{F}_x$ dans $\mathcal{F}_x \otimes_{\mathcal{O}_x} \mathcal{R}(X)$. Lorsque $X$ est un préschéma intègre (donc chacun des $\mathcal{O}_x$ intègre et $\mathcal{R}(X)$ égal à son corps des fractions), on dit que $\mathcal{F}$ est sans torsion si chacun des $\mathcal{O}_x$-modules $\mathcal{F}_x$ est sans torsion ; alors les homomorphismes précédents sont injectifs, autrement dit :

Corollaire 3. — Si $X$ est un préschéma intègre, tout $\mathcal{O}_X$-module quasi-cohérent sans torsion $\mathcal{F}$ est isomorphe à un sous-faisceau simple de la forme $(\mathcal{R}(X))^{(I)}$, engendré par ce sous-faisceau (considéré comme $\mathcal{R}(X)$-module).

Le cardinal de $I$ est appelé le rang de $\mathcal{F}$ ; c'est donc aussi le rang de chacun des $\mathcal{O}_x$-modules $\mathcal{F}_x$, et aussi le rang de $\Gamma(U, \mathcal{F})$ en tant que $\Gamma(U, \mathcal{O}_X)$-module sur tout ouvert affine $U \neq \emptyset$. En particulier :

Corollaire 4. — Sur un préschéma intègre $X$, tout $\mathcal{O}_X$-module quasi-cohérent sans torsion et de rang 1 est isomorphe à un sous-faisceau de $\mathcal{R}(X)$.

\newpage

%% ===== Page 126 =====

\marginpar{Archives Grothendieck sept. 59}

\setcounter{page}{91}

§ 6. Schémas d'anneaux locaux.

1. Anneaux locaux apparentés.

Pour tout anneau local $A$, nous désignerons par $\mathfrak{m}(A)$ l'idéal maximal de $A$.

Lemme 1. — Soient $A, B$ deux anneaux locaux tels que $A \subset B$ ; les conditions suivantes sont équivalentes : (i) $\mathfrak{m}(B) \cap A = \mathfrak{m}(A)$ ; (ii) $\mathfrak{m}(A) \subset \mathfrak{m}(B)$ ; (iii) $1$ n'appartient pas à l'idéal de $B$ engendré par $\mathfrak{m}(A)$.

Il est évident que (i) entraîne (ii) et (ii) entraîne (iii) ; enfin si (iii) est vérifiée, $\mathfrak{m}(B) \cap A$ contient $\mathfrak{m}(A)$ et ne contient pas $1$ donc est égal à $\mathfrak{m}(A)$.

Lorsque les conditions équivalentes du lemme 1 sont vérifiées, on dit que $B$ domine $A$. Il est clair que, dans l'ensemble des sous-anneaux locaux d'un anneau $R$, la relation de domination est une relation d'ordre.

Considérons maintenant un corps $R$. Pour tout sous-anneau $A$ de $R$, nous désignerons par $L(A)$ l'ensemble des anneaux locaux $A_{\mathfrak{p}}$, où $\mathfrak{p}$ parcourt le spectre premier de $A$ ; comme $\mathfrak{p} = (\mathfrak{p} A_{\mathfrak{p}}) \cap A$, l'application $\mathfrak{p} \to A_{\mathfrak{p}}$ de $\operatorname{Spec}(A)$ dans $L(A)$ est bijective.

Lemme 2. — Soient $R$ un corps, $A$ un sous-anneau de $R$. Pour qu'un sous-anneau local $M$ de $R$ domine un anneau $A_{\mathfrak{p}} \in L(A)$, il faut et il suffit que $A \subset M$ ; l'anneau local $M$ domine alors $A_{\mathfrak{p}}$ et correspond à $\mathfrak{p} = \mathfrak{m}(M) \cap A$.

En effet, si $M$ domine $A_{\mathfrak{p}}$, on a $\mathfrak{m}(M) \cap A_{\mathfrak{p}} = \mathfrak{p} A_{\mathfrak{p}}$ d'après le lemme 1, d'où l'unicité de $\mathfrak{p}$ ; d'autre part, si $A \subset M$, $\mathfrak{m}(M) \cap A$ est premier dans $A$, et comme $A - \mathfrak{p} \subset M - \mathfrak{m}(M)$, $A_{\mathfrak{p}} \subset M_{\mathfrak{m}(M)} = M$, donc $M$ domine $A_{\mathfrak{p}}$.

Lemme 3. — Soient $R$ un corps, $M, N$ deux sous-anneaux locaux de $R$, $P$ le sous-anneau de $R$ engendré par $M \cup N$. Les conditions suivantes

\newpage

%% ===== Page 127 =====

- 92 -

sont équivalentes :
(i) Il existe un idéal premier $\mathfrak{p}$ de $R$ tel que $\mathfrak{m}(M) = \mathfrak{p} \cap M$, $\mathfrak{m}(N) = \mathfrak{p} \cap N$.
(ii) L'idéal engendré dans $R$ par $\mathfrak{m}(M) \cup \mathfrak{m}(N)$ est distinct de $R$.
(iii) Il existe un sous-anneau local $Q$ de $R$ dominant à la fois $M$ et $N$.

Il est clair que (i) implique (ii) ; inversement, si $\notin \mathfrak{p}$, $\mathfrak{a}$ est contenu dans un idéal maximal $\mathfrak{n}$ de $R$ et comme $1 \notin \mathfrak{n}$, $\mathfrak{n} \cap M$ contient $\mathfrak{m}(M)$ et est distinct de $M$, donc $\mathfrak{n} \cap M = \mathfrak{m}(M)$, et de même $\mathfrak{n} \cap N = \mathfrak{m}(N)$. Il est clair que si $Q$ domine $M$ et $N$, $P \subset Q$, et $\mathfrak{m}(M) = \mathfrak{m}(Q) \cap M = (\mathfrak{m}(Q) \cap P) \cap M$, $\mathfrak{m}(N) = (\mathfrak{m}(Q) \cap P) \cap N$, donc (iii) entraîne (i) ; la réciproque est évidente en prenant $Q = R_{\mathfrak{p}}$.

Lorsque les conditions du lemme 3 sont satisfaites, on dit que les anneaux locaux $M$ et $N$ sont apparentés.

Proposition 1. - Soient $A, B$ deux sous-anneaux d'un corps $R$, $C$ le sous-anneau de $R$ engendré par $A \cup B$. Les conditions suivantes sont équivalentes :
(1) Pour tout anneau local $Q$ contenant $A$ et $B$, on a $A_{\mathfrak{p}} = B_{\mathfrak{q}}$, en posant $\mathfrak{p} = \mathfrak{m}(Q) \cap A$, $\mathfrak{q} = \mathfrak{m}(Q) \cap B$.
(2) Pour tout idéal premier $\mathfrak{r}$ de deux $C$, on a $A_{\mathfrak{p}} = B_{\mathfrak{q}}$ en posant $\mathfrak{p} = \mathfrak{r} \cap A$, $\mathfrak{q} = \mathfrak{r} \cap B$.
(3) Si $A \in L(A)$ et $N \in L(B)$ sont apparentés, ils sont identiques.
(4) On a $L(A) \cap L(B) = L(C)$.

Les lemmes 2 et 3 prouvent que (i) et (iii) sont équivalentes ; il est clair que (i) entraîne (ii) en l'appliquant à $Q = C_{\mathfrak{r}}$ ; inversement (ii) entraîne (i), car si $Q$ contient $A \cup B$, il contient $C$ et si $\mathfrak{r} = \mathfrak{m}(Q) \cap C$, on a $\mathfrak{p} = \mathfrak{r} \cap A$ et $\mathfrak{q} = \mathfrak{r} \cap B$ d'après le lemme 2. Il est immédiat que (iv) implique (i), car si $Q$ contient $A \cup B$, il domine un anneau local $C_{\mathfrak{r}} \in L(C)$ (lemme 2) ; on a par hypothèse $C_{\mathfrak{r}} \in L(A) \cap L(B)$, et les lemmes 1 et 2 prouvent que $C_{\mathfrak{r}} = A_{\mathfrak{p}} = B_{\mathfrak{q}}$. Prouvons enfin que (iii)

\newpage

%% ===== Page 128 =====

- 93 -

entraîne (iv). Soit $Q \in L(C)$ ; $Q$ domine un $M \in L(A)$ et un $N \in L(B)$ (lemme 2), donc $M$ et $N$, étant apparentés, sont identiques par hypothèse. Comme on a alors $C \subset M$, $M$ domine un $Q' \in L(C)$, donc $Q$ domine $Q'$, ce qui (lemme 2) entraîne nécessairement $Q=Q'=M$, donc $Q \in L(A) \cap L(B)$. Inversement, si $Q \in L(A) \cap L(B)$, on a $C \subset Q$, donc (lemme 2) $Q$ domine un $Q'' \in L(C) \subset L(A) \cap L(B)$ ; $Q$ et $Q''$ étant apparentés sont identiques, donc $Q''=Q \in L(C)$, ce qui achève la démonstration.

2. Anneaux locaux d'un schéma intègre.

Soient $X$ un préschéma intègre, $R$ son corps des fonctions rationnelles, identique à l'anneau local du point générique $\xi$ de $X$. Pour tout $x \in X$, l'anneau local $\mathcal{O}_x$ est intègre ; tout ouvert affine $U$ contenant $x$ contient $\xi$ ; si $A$ est l'anneau de $U$, $R$ est le corps des fractions de $A$ (§ 5, n°1, cor. de la prop. 1), donc $R$ est aussi l'anneau des fractions de $\mathcal{O}_x$, et on identifie canoniquement $\mathcal{O}_x$ à un sous-anneau de $R$ en faisant correspondre à tout germe de section $s \in \mathcal{O}_x$ l'unique fonction rationnelle ayant ce germe au point $x$ (toute section ayant pour germe $s$ étant en effet définie sur un ouvert dense de $X$). Pour tout $f \in R$, le domaine de définition de $f$ (§ 5, n°2, déf. 2) n'est autre que l'ensemble ouvert $\delta(f)$ des $x \in X$ tels que $f \in \mathcal{O}_x$ (avec l'identification précédente). Pour tout ouvert $U \subset X$, on a $\Gamma(U, \mathcal{O}_X) = \bigcap_{x \in U} \mathcal{O}_x$ d'après ce qui précède.

Proposition 2. - Soient $X$ un préschéma intègre, $R$ son corps des fonctions rationnelles. Pour que $X$ soit un schéma, il faut et il suffit que la relation "$\mathcal{O}_x$ et $\mathcal{O}_y$ sont apparentés" entre points $x, y$ de $X$ implique $x=y$.

Supposons cette condition vérifiée et montrons que $X$ est séparé. Soient $U$ et $V$ deux ouverts affines distincts dans $X$, $A$ et $B$ leurs anneaux, identifiés à des sous-anneaux de $R$ ; $U$ (resp. $V$) s'identi-

\newpage

%% ===== Page 129 =====

fie donc à $L(A)$ (resp. $L(B)$) et l'hypothèse entraîne (prop.1) que si $C$ est le sous-anneau de $R$ engendré par $A \cup B$, $W=U \cap V$ s'identifie à $L(A) \cap L(B) = L(C)$. En outre, on sait que tout sous-anneau $E$ de $R$ est égal à l'intersection des anneaux locaux appartenant à $L(E)$ ; $C$ s'identifie donc à l'intersection des anneaux $\mathcal{O}_z$, où $z \in W$, autrement dit à $\Gamma(W, \mathcal{O}_W)$. Considérons alors le sous-préschéma induit par $W$ sur $W$ ; à l'homomorphisme identique $\varphi: C \to \Gamma(W, \mathcal{O}_W)$ correspond (§ 2, n°2, prop.3) un morphisme $\Phi = (\varphi, \psi): W \to \operatorname{Spec}(C)$. Nous allons voir que $\Phi$ est un isomorphisme de préschémas, d'où résultera que $W$ est un ouvert affine. Soient $T \subset W$ un ouvert affine de $X$, $E = \Gamma(T, \mathcal{O}_T)$ son anneau ; comme $T$ est affine, tout idéal premier de $E$ est de la forme $\mathfrak{m}_x \cap E$, où $x \in T$, et par suite $\varphi^{-1}(\mathfrak{m}_x \cap E) = \mathfrak{m}_x \cap C$. En d'autres termes, $\psi(x) = \mathfrak{m}_x \cap C$ pour tout $x \in W$ ; si on compose $\psi$ avec l'application $\rho: \operatorname{Spec}(C) \to \operatorname{Spec}(A)$ qui, à tout idéal premier $\mathfrak{r} \subset C$, fait correspondre $\mathfrak{r} \cap A$, on voit donc que $\rho \circ \psi$ n'est autre que l'injection canonique $W \to U$, ce qui montre déjà que $\psi$ est injective ; en outre pour tout idéal premier $\mathfrak{p}$ de $C$ est de la forme $\mathfrak{m}_x \cap C$, où $x \in L(C)=W$, donc $\psi$ est bijective. De plus, $\psi_x: \mathcal{O}_{\psi(x)} \to \mathcal{O}_x$ est l'injection $\mathcal{O}_{\mathfrak{m}_x \cap C} \to \mathcal{O}_x$, si $\mathfrak{r} = \mathfrak{m}_x \cap C$ ; comme $\mathcal{O}_{\mathfrak{r}} = \mathcal{O}_x$ (prop.1), $\psi_x$ est bijective pour tout $x \in W$. Il reste à voir que $\psi$ est un homéomorphisme, autrement dit que pour toute partie fermée $F \subset W$, $\psi(F)$ est fermé dans $C$. Or $F$ est la trace sur $W$ d'un ensemble fermé de $U$, de la forme $V(\mathfrak{a})$, où $\mathfrak{a}$ est un idéal de $A$ (§ 1,n°1, cor.1 de la prop.2) ; montrons que $\psi(F) = V(\mathfrak{a} C)$, ce qui prouvera notre assertion. En effet, les idéaux premiers de $C$ contenant $\mathfrak{a} C$ sont les idéaux premiers de $C$ contenant $\mathfrak{a}$, donc les idéaux de la forme $\mathfrak{m}_x \cap C$ où $\mathfrak{a} \subset \mathfrak{m}_x$ et $x \in W$ ; comme $\mathfrak{a} \subset \mathfrak{m}_x$ équivaut à $x \in V(\mathfrak{a}) = F$ pour $x \in U$, on a bien $\psi(F) = V(\mathfrak{a} C)$.

Il est maintenant immédiat de voir que $X$ est séparé, car $U \cap V$ est affine et son anneau $C$ est engendré par la réunion $A \cup B$ des anneaux.

\newpage

%% ===== Page 130 =====

fie donc à $L(A)$ (resp. $L(B)$) et l'hypothèse entraîne (prop. 1) que si $C$ est le sous-anneau de $R$ engendré par $A \cup B$, $W=U \cap V$ s'identifie à $L(A) \cap L(B) = L(C)$. En outre, on sait que tout sous-anneau $E$ de $R$ est égal à l'intersection des anneaux locaux appartenant à $L(E)$; $C$ s'identifie donc à l'intersection des anneaux $\mathcal{O}_z$, où $z \in W$, autrement dit à $\Gamma(W, \mathcal{O}_W)$. Considérons alors le sous-préschéma induit par $X$ sur $W$ ; à l'homomorphisme identique $\varphi : C \to \Gamma(W, \mathcal{O}_W)$ correspond (§ 2, n° 2, prop. 3) un morphisme $\Phi = (\psi, \theta) : W \to \operatorname{Spec}(C)$. Nous allons voir que $\Phi$ est un isomorphisme de préschémas, d'où résultera que $W$ est un ouvert affine. Soient $T \subset W$ un ouvert affine de $X$, $E = \Gamma(T, \mathcal{O}_X)$ son anneau ; comme $T$ est affine, tout idéal premier de $E$ est de la forme $\mathfrak{m}_x \cap E$, où $x \in T$, et par suite $\varphi^{-1}(\mathfrak{m}_x \cap E) = \mathfrak{m}_x \cap C$. En d'autres termes, $\psi(x) = \mathfrak{m}_x \cap C$ pour tout $x \in W$ ; si on compose $\psi$ avec l'application $\rho : \operatorname{Spec}(C) \to \operatorname{Spec}(A)$ qui, à tout idéal premier $\mathfrak{r} \subset C$, fait correspondre $\mathfrak{r} \cap A$, on voit donc que $\rho \circ \psi$ n'est autre que l'injection canonique $W \to U$, ce qui montre déjà que $\psi$ est injective ; en outre, tout idéal premier de $C$ est de la forme $\mathfrak{m}_x \cap C$, où $x \in L(C) = W$, donc $\psi$ est bijective. De plus, $\theta_x : C_{\psi(x)} \to \mathcal{O}_x$, si $\mathfrak{r} = \mathfrak{m}_x \cap C$ ; comme $C_{\mathfrak{r}} = \mathcal{O}_x$ (prop. 1), $\theta_x$ est bijective pour tout $x \in W$. Il reste à voir que $\psi$ est un homéomorphisme, autrement dit que pour toute partie fermée $F \subset W$, $\psi(F)$ est fermé dans $C$. Or $F$ est la trace sur $W$ d'un ensemble fermé de $U$, de la forme $V(\mathfrak{a})$, où $\mathfrak{a}$ est un idéal de $A$ (§ 1, n° 1, cor. 1 de la prop. 2) ; montrons que $\psi(F) = V(\mathfrak{a}C)$, ce qui prouvera notre assertion. En effet, les idéaux premiers de $C$ contenant $\mathfrak{a}C$ sont les idéaux premiers de $C$ contenant $\mathfrak{a}$, donc les idéaux de la forme $\mathfrak{m}_x \cap C$ où $\mathfrak{a} \subset \mathfrak{m}_x$ et $x \in W$ ; comme $\mathfrak{a} \subset \mathfrak{m}_x$ équivaut à $x \in V(\mathfrak{a}) = F$ pour $x \in U$, on a bien $\psi(F) = V(\mathfrak{a}C)$.

Il est maintenant immédiat de voir que $X$ est séparé, car $U \cap V$ est affine et son anneau $C$ est engendré par la réunion $A \cup B$ des anneaux

\newpage

%% ===== Page 131 =====

- 95 -

de $U$ et $V$ (§ 3, n°6, prop. 20).

Inversement, supposons $X$ séparé, et soient $x, y$ deux points de $X$ tels que $\mathcal{O}_x$ et $\mathcal{O}_y$ soient apparentés. Soient $U$ (resp. $V$) un ouvert affine contenant $x$ (resp. $y$), d'anneau $A$ (resp. $B$) ; on sait alors que $U \cap V$ est affine, et que son anneau $C$ est engendré par $A$ et $B$ (§ 3, n°6, prop. 20). Si $\mathfrak{p} = \mathfrak{m}_x \cap A$, $\mathfrak{q} = \mathfrak{m}_y \cap B$, on a $A_\mathfrak{p} = \mathcal{O}_x$, $B_\mathfrak{q} = \mathcal{O}_y$, et comme $A_\mathfrak{p}$ et $B_\mathfrak{q}$ sont apparentés, il existe un idéal premier $\mathfrak{r}$ de $C$ tel que $\mathfrak{r} \cap A = \mathfrak{p}$, $\mathfrak{r} \cap B = \mathfrak{q}$ (n°1, lemme 3). Mais alors il existe un point $z \in U \cap V$ tel que $\mathfrak{m}_z = \mathfrak{r}$, puisque $U \cap V$ est affine, et on a évidemment $x = z$ et $y = z$, d'où $x = y$.

Corollaire 1. - Soit $X$ un schéma intègre, $x, y$ deux points de $X$. Pour que $x \in \overline{\{y\}}$, il faut et il suffit que $\mathcal{O}_y \subset \mathcal{O}_x$, autrement dit que toute fonction rationnelle définie en $x$ soit définie en $y$.

La condition est évidemment nécessaire puisque $\delta(f)$ est ouvert pour toute fonction rationnelle $f \in R$ ; montrons qu'elle est suffisante. Si $\mathcal{O}_y \subset \mathcal{O}_x$, il existe un idéal premier $\mathfrak{p}$ de $\mathcal{O}_x$ tel que $\mathcal{O}_y$ domine $(\mathcal{O}_x)_\mathfrak{p}$ (n°1, lemme 2) ; or (§ 2, n°2, prop. 4) il existe $z \in X$ tel que $x \in \overline{\{z\}}$ et $\mathcal{O}_z = (\mathcal{O}_x)_\mathfrak{p}$ ; comme $\mathcal{O}_z$ et $\mathcal{O}_y$ sont apparentés, on a $z = y$ par la prop. 2, d'où le corollaire.

Corollaire 2. - Si $X$ est un schéma intègre, l'application $x \to \mathcal{O}_x$ est injective ; autrement dit, si $x, y$ sont deux points distincts de $X$, il existe une fonction rationnelle définie en l'un de ces points et non en l'autre.

Cela résulte du cor. 1 et de l'axiome $(T_0)$.

Corollaire 3. - Soit $X$ un schéma intègre dont l'espace de base est noethérien ; lorsque $f$ parcourt le corps $R$ des fonctions rationnelles sur $X$, les $\delta(f)$ forment une base de la topologie de $X$.

En effet, toute partie fermée de $X$ est alors réunion finie d'ensembles fermés irréductibles, c'est-à-dire de la forme $\overline{\{y\}}$ (§ 2, n°1,

\newpage

%% ===== Page 132 =====

- 96 -

prop.1). Or, si $x \notin \{y\}$, il existe une fonction rationnelle $f$ définie en $x$ et non en $y$ (cor.1), autrement dit on a $x \in \delta(f)$ et $\delta(f)$ ne rencontre pas $\{y\}$. Le complémentaire de $\{y\}$ est par suite réunion d'ensembles de la forme $\delta(f)$, et en vertu de la première remarque, tout ouvert de $X$ est réunion d'intersections finies d'ensembles de la forme $\delta(f)$.

Le cor.3 montre que la topologie de $X$ est entièrement caractérisée par la donnée de la famille d'anneaux locaux $(\mathcal{O}_x)_{x \in X}$ ayant $R$ pour corps des fractions ; il revient au même d'ailleurs de dire que les parties fermées de $X$ sont définies de la façon suivante : étant donnée une partie finie $\{x_1, \dots, x_n\}$ de $X$, on considère l'ensemble des $y \in X$ tels que $\mathcal{O}_y \subset \mathcal{O}_{x_i}$ pour un indice $i$ au moins, et ces ensembles (pour tous les choix de la partie finie $\{x_1, \dots, x_n\}$ de $X$) sont les ensembles fermés de $X$. En outre, une fois connue la topologie de $X$, le faisceau structural $\mathcal{O}_X$ est aussi bien déterminé par la famille des $\mathcal{O}_x$, puisque $\Gamma(U, \mathcal{O}_X) = \bigcap_{x \in U} \mathcal{O}_x$. La famille $(\mathcal{O}_x)_{x \in X}$ détermine donc complètement le préschéma lorsque $X$ est un schéma intègre d'espace de base noethérien.

On dit qu'un morphisme $f : X \to Y$ de préschémas est dominant si $f(X)$ est dense dans $Y$.

Proposition 3. - Soient $X, Y$ deux schémas intègres, $f : X \to Y$ un morphisme dominant, $K$ (resp. $L$) le corps des fonctions rationnelles de $X$ (resp. $Y$). Alors $L$ s'identifie à un sous-corps de $K$, et pour tout $x \in X$, $\mathcal{O}_{f(x)}$ est l'unique anneau local de $Y$ dominé par $\mathcal{O}_x$.

En effet, soit $f = (\varphi, \theta)$ ; si $a$ est le point générique de $X$, on a $\varphi(a) \in \overline{\{a\}}$, et comme par thèse $\varphi(X)$ est dense dans $Y$, on voit que $Y = \overline{\{\varphi(a)\}}$, autrement dit $\varphi(a)$ est le point générique de $Y$; $\theta_a$ est par suite un monomorphisme de $L = \mathcal{O}_{\varphi(a)}$ dans $K = \mathcal{O}_a$. Comme tout ouvert affine de $Y$ contient

\newpage

%% ===== Page 133 =====

- 97 -

\textit{la}, il résulte de la prop. 3 du § 2, n$^\circ$2 que pour tout ouvert affine $U$ de $Y$, l'homomorphisme de $\Gamma(U, \mathcal{O}_Y)$ dans $\Gamma(\psi^{-1}(U), \mathcal{O}_X)$ correspondant à la restriction de $\psi$ à $\Gamma(U, \mathcal{O}_Y)$. Donc, pour tout $x \in X$, $\mathcal{O}_{x}^{\psi}$ est la restriction à $\mathcal{O}_{\psi(x)}$ de $\mathcal{O}_x$, et est par suite un monomorphisme ; on sait en outre, par définition d'un morphisme, que l'image réciproque par $\psi_x^\sharp$ de l'idéal minimal de $\mathcal{O}_{\psi(x)}$ est l'idéal minimal de $\mathcal{O}_x$, et par suite, si on identifie $L$ à un sous-corps de $K$ par $\psi_a^\sharp$, $\mathcal{O}_{\psi(x)}$ est dominé par $\mathcal{O}_x$ (n$^\circ$1, lemme 1) ; c'est d'ailleurs le seul anneau local de $Y$ dominé par $\mathcal{O}_x$, puisque deux anneaux locaux de $Y$ qui sont apparentés sont identiques (prop. 2).

\underline{Proposition 4} . - Soient $X$ un préschéma irréductible, $f : X \to Y$ une immersion locale (resp. un isomorphisme local) ; on suppose en outre le morphisme $f$ séparé. Alors $f$ est une immersion (resp. une immersion ouverte).

Soit $f = (\psi, \theta)$ ; il suffit de prouver, dans les deux cas, que $\psi$ est un homéomorphisme de $X$ sur $\psi(X)$ ; l'hypothèse entraînera alors que $\psi(X)$ est localement fermé (resp. ouvert) dans $Y$ et que $f$ est une immersion (resp. une immersion ouverte) (cf. § 3, n$^\circ$4, déf. 3). Remplaçant $f$ par $f_{\operatorname{red}}$ (§ 3, n$^\circ$4, prop. 10), on peut supposer $X$ et $Y$ réduits. En outre, par hypothèse, pour tout $x \in X$, il y a un voisinage ouvert $U$ de $x$ et un voisinage ouvert $V$ de $\psi(x)$ tel que la restriction de $f$ à $U$ s'écrive $j \circ g$, où $g$ est un isomorphisme du préschéma induit sur $U$, sur un sous-préschéma formé $Z$ du préschéma induit sur $V$ et $j$ l'injection canonique de $Z$ dans $V$. Si $Y'$ est le plus petit sous-préschéma de $Y$ ayant $\psi(X)$ comme espace de base (§ 3, n$^\circ$4, prop. 9), il est clair que $Y'$ majore $Z$, puisque $\psi(X) \cap V$ est l'espace de base de $Z$. Donc le morphisme $f$ est majoré par le morphisme d'injection $Y' \to Y$ ; en outre l'hypothèse que $f$ est une immersion ouverte et que $X$ est réduit, entraîne-

\newpage

%% ===== Page 134 =====

ne que les préschémas induits par $Y'$ sur $\psi(X) \cap V = \psi(X) \cap V$ est identique à $Z$ (§ 3, n°4, prop. 9) ; $f$ est par suite le composé d'une immersion locale (resp. d'un isomorphisme local) de $X$ dans $Y'$ et de l'injection canonique $Y' \to Y$, et on peut donc se réduire au cas où $Y=Y'$, autrement dit au cas où $f$ est dominant. Si $x$ est le point générique de $X$, $\psi(x)$ contient alors $\psi(X)$, et est donc identique à $Y$, autrement dit $Y$ est irréductible. On est ainsi ramené au cas où $X$ et $Y$ sont tous deux intègres. En outre, la question étant locale sur $Y$, on peut supposer $Y$ affine, donc un schéma ; comme $f$ est supposé séparé (tenant compte de § 3, n°6, prop. 21 et cor. 4 de la prop. 20), on en conclut que $X$ est aussi un schéma (§ 3, n°6, cor. 2 de la prop. 20), et on est ainsi dans les conditions d'application de la prop. 3. Pour tout $x \in X$, $\phi_x^\flat$ est donc injectif ; mais l'hypothèse que $f$ est une immersion locale implique que $\phi_x^\flat$ est surjectif (§ 3, n°2, prop. 4), donc $\phi_x^\flat$ est bijectif, autrement dit $\mathcal{O}_{\psi(x)} = \mathcal{O}_x$ (avec l'identification de la prop. 3). Cela implique donc que $\psi$ est une application injective (cor. 2 de la prop. 2) ; comme, avec les notations introduites plus haut, et l'hypothèse $Y'=Y$, on a vu que $\psi(U) = \psi(X) \cap V$ et que la restriction de $\psi$ à $U$ est par hypothèse un homéomorphisme de $U$ sur $\psi(U)$, on en conclut bien que $\psi$ est un homéomorphisme de $X$ sur $\psi(X)$.

3. Les schémas de Chevalley.

Soient $X$ un schéma intègre, $R$ son corps de fonctions rationnelles ; désignons par $X'$ l'ensemble des sous-anneaux locaux $\mathcal{O}_x$ de $R$ ($x \in X$). L'ensemble $X'$ vérifie les trois conditions suivantes :

(Sch 1) Pour tout $M \in X'$, $R$ est le corps des fractions de $M$.

(Sch 2) Il existe un ensemble fini de sous-anneaux noethériens $A_i$ de $R$ tels que $X' = \bigcup_i \operatorname{Spec}(A_i)$ et que, pour tout couple d'indices $i, j$, le sous-anneau $A_{ij}$ de $R$ engendré par $A_i \cup A_j$ soit de type fini sur $A_i$.

\newpage

%% ===== Page 135 =====

- 99 -

(Sch 3) Deux éléments $M, N$ de $X'$ qui sont apparentés sont identiques.

On a en effet vu au début du $n^\circ 2$ que (Sch 1) est satisfaite, et (Sch 3) résulte de la prop. 2 du $n^\circ 2$. Pour démontrer (Sch 2), il suffit de recouvrir $X$ par un nombre fini d'ouverts affines $U_i$ d'anneaux noethériens et de prendre $A_i = \Gamma(U_i, \mathcal{O}_X)$; l'hypothèse que $X$ est un schéma entraîne alors que $U_i \cap U_j$ est affine et que $\Gamma(U_i \cap U_j, \mathcal{O}_X) = A_{ij}$ (§ 3, $n^\circ 6$, prop. 20); en outre, comme $U_i \cap U_j \to U_i$ est une immersion ouverte, c'est un morphisme de type fini (§ 4, $n^\circ 2$, prop. 7) donc $A_{ij}$ est une algèbre de type fini sur $A_i$ (§ 4, $n^\circ 2$, prop. 6).

Les structures dont les axiomes sont (Sch 1), (Sch 2), (Sch 3) généralisent les "schémas" au sens de C. Chevalley, qui suppose en outre que $R$ est une extension de type fini d'un sous-corps $K$ et que les $A_i$ sont des $K$-algèbres de type fini (Sém. Cartan-Chevalley, 1955-56). Inversement, si on a une telle structure sur un ensemble $X'$, on peut lui associer un schéma intègre $X$ en utilisant les remarques suivant le cor. 3 de la prop. 2 du $n^\circ 2$ : l'espace de base de $X$ est égal à $X'$, muni de la topologie caractérisée comme il a été dit dans ces remarques, et du faisceau $\mathcal{O}_X$ tel que $\Gamma(U, \mathcal{O}_X) = \bigcap_{x \in U} \mathcal{O}_x$ pour tout ouvert $U \subset X$. Nous laissons au lecteur le soin de vérifier qu'on obtient bien ainsi un schéma intègre, dont les anneaux locaux sont les éléments de $X'$ ; nous n'utiliserons pas ce résultat par la suite.

\newpage

%% ===== Page 136 =====

\marginpar{Archives Grothendieck-Sept 59}

- 100 -

\section*{§ 7. Théorie de la dimension. Ensembles constructibles.}

Pour l'essentiel, ce paragraphe ne fait que reprendre, dans un langage plus géométrique, des faits bien connus en théorie de la dimension pour les anneaux commutatifs. Font exception les th. 2 et 3 empruntés au Séminaire Cartan-Chevalley 1955-56. Sauf au chap. III, § 2, n° et § 3, n°, nous n'utiliserons pas les résultats de ce paragraphe avant le chap. IV.

\subsection*{1. Dimension combinatoire d'un espace topologique.}

Soit $I$ un ensemble ordonné ; une chaîne d'éléments de $I$ est par définition une suite finie strictement croissante $i_0 < i_1 < \dots < i_n$ d'éléments de $I$ ; par définition, la longueur de cette suite est $n$. Si $X$ est un espace topologique, l'ensemble de ses parties fermées irréductibles est ordonné par inclusion, d'où la notion de chaîne de parties fermées irréductibles.

\textbf{Définition 1.} - Soit $X$ un espace topologique. On appelle dimension combinatoire de $X$ (ou simplement dimension de $X$ s'il n'y a pas de confusion) et on note $\dim(X)$, ou simplement $\dim(X)$, la borne supérieure des longueurs des chaînes de parties fermées irréductibles de $X$.

Pour tout $x \in X$, on appelle dimension combinatoire de $X$ en $x$, et on note $\dim_x(X)$ le nombre $\inf(\dim(U))$, où $U$ parcourt l'ensemble des voisinages ouverts de $x$. On dit que $X$ est équidimensionnel si toutes les composantes irréductibles de $X$ ont même dimension.

\textbf{Définition 2.} - Étant donnée une partie fermée irréductible $Y$ d'un espace topologique $X$, on appelle codimension combinatoire (ou simplement codimension) de $Y$ dans $X$ la borne supérieure des longueurs des chaînes de parties fermées irréductibles de $X$, dont $Y$ est le plus petit élément ; on la note $\operatorname{codim}_X(Y)$. Si $Y$ est une partie fermée quelconque de $X$, on appelle codimension combinatoire de $Y$ dans $X$ et on note encore $\operatorname{codim}_X(Y)$ la borne inférieure des codimensions dans...

\newpage

%% ===== Page 137 =====

- 101 -

\noindent \textbf{Proposition 1} .- (i) Si $Y$ est une partie d'un espace topologique $X$, $\dim(Y) \leqslant \dim(X)$.

\noindent (ii) Si un espace topologique $X$ est réunion finie de parties fermées $X_i$, alors $\dim(X) = \sup_i (\dim(X_i))$.

\noindent \rule{\linewidth}{0.4pt}
\noindent pour toute partie irréductible $Z$ fermée dans $Y$, l'adhérence $\overline{Z}$ de $Z$ dans $X$ est irréductible ; d'où (i). Si $X = \bigcup_i X_i$, les $X_i$ étant fermés, toute chaîne de parties fermées irréductibles de $X$ est nécessairement contenue dans l'un des $X_i$, ce qui prouve (ii). \rule{\linewidth}{0.4pt} On notera que si $f$ est une application continue de $X$ dans un espace topologique $Y$, on peut avoir $\dim(f(X)) > \dim(X)$ ; on en a un exemple en prenant pour $X$ un espace discret à 2 éléments $a, b$, pour $Y$ l'ensemble $\{a, b\}$ muni de la topologie dans laquelle $\emptyset$, $\{a\}$ et $\{a, b\}$ sont les seuls ensembles fermés ; si $f$ est l'application identique $X \to Y$, $\dim Y = 1$ et $\dim X = 0$.

\noindent \textbf{Corollaire} .- Soient $x \in X$, $U$ un voisinage de $x$, $Y_i$ ($1 \leqslant i \leqslant n$) des parties fermées de $U$ telles que $x \in Y_i$ pour $1 \leqslant i \leqslant n$ et que $U$ soit la réunion des $Y_i$. Alors $\dim_x X = \sup_i (\dim_x(Y_i))$.

\noindent En effet, d'après la prop. 1, $\dim_x(X) = \inf_V (\sup_i(\dim(Y_i \cap V)))$ lorsque $V$ parcourt l'ensemble des voisinages ouverts de $x$ contenus dans $U$. Le corollaire est évident si $\sup_i (\dim_x(Y_i)) = +\infty$ ; sinon, il y a un voisinage $V$ de $x$ tel que \rule{1.5cm}{0.4pt} $\dim(Y_i \cap V) = \dim(Y_i)$ pour tout $i$ et tout $V \subset V_0$.

\noindent \textbf{Proposition 2} .- (i) Soit $\Phi$ l'ensemble des parties fermées irréductibles de $X$. On a $\dim(X) = \sup_{Y \in \Phi} (\operatorname{codim}_X(Y))$.

\noindent (ii) On a $\dim(X) = \sup_{x \in X} \dim_x(X)$.

\noindent (i) découle aussitôt des déf. 1 et 2. \rule{4cm}{0.4pt}

\newpage

%% ===== Page 138 =====

- 102 -

D'autre part, si $Z_0 \subset Z_1 \subset \dots \subset Z_n$ est une chaîne de parties fermées irréductibles de $X$ et si $x \in Z_0$, pour tout ouvert $U \subset X$ contenant $x$, $U \cap Z_i$ est irréductible et fermé dans $U$, et comme $U \cap Z_i = Z_i$, les $U \cap Z_i$ sont distincts ; donc $\dim(U) \ge n$, ce qui démontre (ii).

\emph{Corollaire 1}. - Si $(U_\alpha)$ est un recouvrement ouvert de $X$, on a $\dim(X) = \sup_{\alpha}(\dim(U_\alpha))$.

\emph{Corollaire 2}. - Soit $X$ un espace noethérien vérifiant l'axiome $(T_0)$, et soit $F$ l'ensemble des points fermés de $X$. Alors $\dim(X) = \sup_{x \in F}(\dim_x(X))$.

Avec les notations de la démonstration de la prop. 2 (ii), il suffit d'observer qu'il existe dans $Z_0$ un point fermé ($\S$ 4, n$^\circ$ 1, lemme 1).

\emph{Corollaire 3}. - Soit $Y$ une partie fermée de $X$. Alors
$$(1) \qquad \dim(Y) + \operatorname{codim}_X(Y) \le \dim(X).$$
Toute chaîne de parties fermées irréductibles de $Y$ étant contenue dans une composante irréductible de $Y$, on est ramené (par la déf. 2) au cas où $Y$ est irréductible, et (1) est alors conséquence des définitions 1 et 2.

\emph{Corollaire 4}. - Soient $Y, Z, T$ trois parties fermées irréductibles de $X$ telles que $Y \subset Z \subset T$. Alors
$$(2) \qquad \operatorname{codim}_T(Y) \ge \operatorname{codim}_T(Z) + \operatorname{codim}_Z(Y).$$
C'est une conséquence immédiate de la déf. 2.

\emph{Proposition 3}. - Soit $X$ un espace noethérien vérifiant l'axiome $(T_0)$. Pour que $\dim(X) = 0$, il faut et il suffit que $X$ soit fini et discret. Pour qu'une partie fermée $Y$ de $X$ soit de codimension nulle dans $X$, il faut et il suffit que $Y$ contienne une composante irréductible de $X$.

La seconde assertion est immédiate (et indépendante de toute hypothèse sur $X$). De même, si $X$ est séparé (et a fortiori discret), il est clair que $\dim(X) = 0$, les seuls ensembles fermés irréductibles étant les points. Inversement, si $X$ est noethérien et satisfait à $(T_0)$, toute composante irréductible de $X$ contient un point fermé ($\S$ 4, n$^\circ$ 1, lemme 1).

\newpage

%% ===== Page 139 =====

- 103 -

donc doit être réduite à ce point si $\dim(X)=0$ ; le nombre de ces composantes étant fini, $X$ est fini et discret.

\emph{Corollaire 1.} — Soit $X$ un espace noethérien vérifiant $(T_0)$. Pour que $\dim_x(X)=0$, il faut et il suffit que $x$ soit isolé.

La condition est (videmment suffisante (sans hypothèse sur $X$). Elle est nécessaire, car il en résulte que $\dim(U)=0$ pour un voisinage ouvert $U$ de $x$, et comme $U$ est noethérien et vérifie $(T_0)$, $U$ est fini et discret.

\emph{Corollaire 2.} — Soit $X$ un espace irréductible. Pour qu'une partie fermée $Y$ de $X$ soit telle que $\operatorname{codim}_X(Y)=0$, il faut et il suffit que $Y=X$.

Nous dirons qu'une chaîne $Z_0 \subset Z_1 \subset \dots \subset Z_n$ de parties fermées irréductibles est \emph{saturée} s'il n'existe pas de partie fermée irréductible $Z'$ distincte des $Z_i$ et telle que $Z_k \subset Z' \subset Z_{k+1}$ pour un indice $k$.

\emph{Proposition 4.} — Soit $X$ un espace topologique tel que pour deux parties fermées irréductibles $Y \subset Z \subset X$, on ait $\operatorname{codim}_Z(Y) < +\infty$. Les conditions suivantes sont équivalentes :

a) Deux chaînes saturées de parties fermées irréductibles de $X$, ayant mêmes extrémités, ont même longueur.

b) Si $Y, Z, T$ sont trois parties fermées irréductibles de $X$ telles que $Y \subset Z \subset T$, on a
$$(3) \quad \operatorname{codim}_T(Y) = \operatorname{codim}_T(Z) + \operatorname{codim}_Z(Y) \text{.} $$

Il est immédiat que a) entraîne b). Réciproquement, supposons b) vérifié, et démontrons que si deux chaînes saturées de mêmes extrémités ont pour longueurs $n$ et $m$, on a nécessairement $n=m$. Raisonnons par récurrence sur $n$, la proposition étant évidente pour $n=1$. Supposons donc $n > 1$, $n < m$, et soit $Z_0 \subset Z_1 \subset \dots \subset Z_n$ une chaîne saturée telle qu'il existe une autre chaîne saturée d'extrémités $Z_0, Z_n$ et de longueur $m$. Comme $\operatorname{codim}_{Z_n}(Z_0) \ge n$, et que

\newpage

%% ===== Page 140 =====

- 104 -
$\operatorname{codim}_{Z}(Z_1) = 1$, il résulte de b) que $\operatorname{codim}_{Z}(Z_n) = \operatorname{codim}_{Z}(Z_0) - 1 > n - 1$ ce qui contredit l'hypothèse de récurrence.

Lorsque les conditions de la prop. 4 sont remplies, on dit que $X$ satisfait à la condition des chaînes.

\emph{Proposition 5.} - Soit $X$ un espace noethérien, de dimension finie, vérifiant l'axiome $(T_0)$. Les conditions suivantes sont équivalentes :
a) Deux chaînes maximales de parties fermées irréductibles de $X$ ont même longueur.
b) ($X$ est équidimensionnel), les points fermés de $X$ ont même codimension, et $X$ satisfait à la condition des chaînes.
c) $X$ est équidimensionnel et pour deux parties fermées irréductibles $Y \subset Z$ de $X$, on a
$$(4) \qquad \dim(Z) = \dim(Y) + \operatorname{codim}_{Z}(Y)$$
d) Les codimensions des points fermés de $X$ sont les mêmes, et pour deux parties fermées irréductibles $Y \subset Z$ de $X$, on a
$$(5) \qquad \operatorname{codim}_{X}(Y) = \operatorname{codim}_{Z}(Y) + \operatorname{codim}_{X}(Z)$$

Les hypothèses sur $X$ entraînent que les extrémités d'une chaîne maximale de parties fermées irréductibles de $X$ sont nécessairement une composante irréductible de $X$ et un point fermé ; en outre, toute chaîne saturée $(Y \subset Z)$ est contenue dans une chaîne maximale autres que ceux de la chaîne donnée dont les éléments sont ou bien contenus dans $Y$ ou bien contiennent $Z$. Ces remarques établissent aussitôt l'équivalence de a) et b), et prouvent aussi que pour toute partie fermée irréductible $Y$ de $X$, on a
$$(6) \qquad \dim(Y) + \operatorname{codim}_{X}(Y) = \dim(X)$$
et comme (3) a aussi lieu, (4) et (5) s'en déduisent aussitôt. Inversement, (4) entraîne (3), donc la condition des chaînes (prop. 4) et en outre, en appliquant (4) pour $Y$ réduit à un point

\newpage

%% ===== Page 141 =====

ferm\'e $x$)
(et $Z$ \'egal \`a une composante irr\'eductible de $X$, on obtient $\operatorname{codim}_X(x) = \dim(Z)$; on en conclut que $c)$ entra\^\i ne $b)$. De m\^eme, $(5)$ entra\^\i ne $(3)$, donc la condition des cha\^\i nes; avec le m\^eme choix de $Y$ et $Z$ que ci-dessus, on a encore, d'apr\`es $(6)$, $\operatorname{codim}_X(x) = \dim X$, donc (puisque toute composante irr\'eductible de $X$ contient un point ferm\'e), $d)$ entra\^\i ne $b)$.

Lorsque les conditions de la prop. 5 sont remplies, on dit que $X$ est \emph{strictement \'equidimensionnel}.

\emph{Corollaire.} -- Si $X$ est strictement \'equidimensionnel, il en est de m\^eme de toute r\'eunion de composantes irr\'eductibles de $X$; en outre, pour toute partie ferm\'ee $Y$ de $X$, on a la relation $(6)$.

La premi\`ere assertion d\'ecoule aussit\^ot de la prop. 5. D'autre part, $(6)$ a lieu pour toute composante irr\'eductible $Y_i$ de $Y$, donc celle des $Y_i$ pour laquelle $\dim(Y_i)$ est le plus grand est aussi celle pour laquelle $\operatorname{codim}_X(Y_i)$ est le plus petit, donc $(6)$ est vrai par d\'efinition de $\dim(Y)$ et de $\operatorname{codim}_X(Y)$.

\emph{Remarque.} -- Le lecteur remarquera que les d\'emonstrations des prop. 4 et 5 s'appliquent \`a un ensemble ordonn\'e quelconque (le fait que les \'el\'ements de cet ensemble sont des parties ferm\'ees irr\'eductibles d'un espace topologique n'intervient pas).

2. Dimension des pr\'esch\'emas.

Compte tenu des d\'ef. 1 et 2 du n$^\circ 1$, les d\'efinitions classiques de la dimension d'un anneau commutatif et du rang d'un id\'eal dans un tel anneau \'equivaut aux suivantes :

\emph{D\'efinition 3.} -- Soient $A$ un anneau commutatif. On appelle dimension de $A$ la dimension de $\operatorname{Spec}(A)$; si $I$ est un id\'eal de $A$, on appelle rang de $I$ la codimension de $V(I)$ dans $\operatorname{Spec}(A)$. On dit que $A$ satisfait \`a la condition des cha\^\i nes (resp. est \'equidimensionnel, strictement \'equidimensionnel) si $\operatorname{Spec}(A)$ satisfait \`a la condition \marginpar{et on note rg(I)}

\newpage

%% ===== Page 142 =====

- 106 -

des chaînes (resp. est équidimensionnel, strictement équidimensionnel).

On observera que dans un préschéma $X$, la relation d'inclusion entre parties fermées irréductibles est identique à la relation $x \in \{y\}$ entre points de $X$ (§ 2, n°1, prop.1).

Proposition 6. — Soient $X$ un préschéma, $Y$ une partie fermée irréductible de $X$; on a
$$(7) \quad \operatorname{codim}_X(Y) = \dim(\mathcal{O}_{X,y})$$
En effet, soit $y$ le point générique de $Y$, $U$ un ouvert affine contenant $y$; comme toute partie fermée irréductible de $X$ contenant $y$ est contenue dans $\overline{U}$, il y a correspondance biunivoque entre ces parties et leurs traces sur $U$. On peut par suite supposer $X$ affine, et la relation (7) résulte alors de la correspondance biunivoque entre idéaux premiers de $A = \Gamma(X, \mathcal{O}_X)$ contenus dans $J_y$, et idéaux premiers de $\mathcal{O}_{X,y} = \mathcal{O}_X / \mathfrak{m}_y$.

Corollaire 1. — Pour que le préschéma $X$ satisfasse à la condition des chaînes, il faut et il suffit que chacun des anneaux locaux $\mathcal{O}_{X,x}$ ($x \in X$) vérifie cette condition ; si de plus $X$ est localement noethérien, il suffit que la condition des chaînes soit vérifiée par les $\mathcal{O}_{X,y}$ correspondant aux points fermés de $X$.

En effet, par définition de la condition des chaînes (n°1) et en vertu de la correspondance biunivoque entre parties fermées irréductibles et points de $X$, on voit aussitôt que pour vérifier la condition des chaînes dans $X$, il suffit de le faire dans chacun des préschémas locaux $\operatorname{Spec}(\mathcal{O}_{X,x})$ (§ 2, n°2) ; si de plus $X$ est localement noethérien, pour tout point $x$ il existe un point fermé $y$ tel que $y \in \overline{\{x\}}$ donc $x \in \operatorname{Spec}(\mathcal{O}_{X,y})$, d'où la seconde assertion du corollaire.

Corollaire 2. — Pour tout préschéma $X$, on a $\dim(X) = \sup_{x \in X}(\dim(\mathcal{O}_{X,x}))$; si de plus $X$ est localement noethérien et si $F$ est l'ensemble des

\newpage

%% ===== Page 143 =====

- 107 -

points fermés de $X$, alors $\dim(X) = \sup_{x \in F} (\dim(\mathcal{O}_x))$.

Il suffit de remarquer que toute chaîne de parties fermées irréductibles de $X$ correspond aux points contenus dans un même schéma local (resp. dans le schéma local d'un point fermé si $X$ est localement noethérien).

Rappelons maintenant que si $A$ est un anneau local noethérien, $\dim(A)$ est égale au nombre minimum de générateurs d'un idéal de définition de $A$ ; pour tout anneau noethérien $A$, la dimension de $A$ est évidemment finie, et pour tout idéal premier $\mathfrak{p} \subset A$, le rang de $\mathfrak{p}$, égal à $\dim(A_{\mathfrak{p}})$ est aussi fini. Par suite :

Corollaire 3. — Pour tout préschéma localement noethérien $X$ et toute partie fermée irréductible $Y$, $\operatorname{codim}_X(Y)$ est finie. Si en outre $X$ est affine, $\operatorname{codim}_X(Y)$ est égal au nombre minimum de sections de $\mathcal{O}_X$ au-dessus de $X$ telles que $Y$ soit une composante irréductible de l'ensemble des $x \in X$ où ces sections "s'annulent" simultanément.

Remarque 4. 1) Un anneau local noethérien ne vérifie pas nécessairement la condition des chaînes (Nagata, On the chain problem of prime ideals, Nagoya Math. Journal, 10 (1956) ; voir aussi Nagata, Note on a chain condition of prime ideals, à paraître dans Kyoto J. Math.).

2) On dit qu'un préschéma $X$ est régulier si tous les anneaux locaux $\mathcal{O}_x$ ($x \in X$) sont réguliers ; un tel préschéma est donc réduit et ses composantes irréductibles sont intègres puisque les $\mathcal{O}_x$ sont intègres (en effet, cf. cor. 1 de la prop. 6). Un tel préschéma vérifie la condition des chaînes ; il en est ainsi de tout anneau local régulier (J.P. Serre). On dira qu'un anneau $A$ est régulier si $\operatorname{Spec}(A)$ est régulier. Comme la condition des chaînes se conserve évidemment par passage à un anneau quotient, les quotients d'anneaux réguliers vérifient cette condition. On sait que les anneaux noethériens qui interviennent en Géométrie algébrique ou en arithmétique sont tous de cette nature.

\newpage

%% ===== Page 144 =====

- 108 -

On notera qu'un anneau régulier $A$ n'est pas nécessairement strictement équidimensionnel, autrement dit (prop. 5) les codimensions des divers points fermés de $X = \operatorname{Spec}(A)$ ne sont pas nécessairement les mêmes.
Par exemple, soient $B$ un anneau de valuation discrète, $(\pi)$ son idéal maximal, $k$ son corps résiduel, $K$ le corps des fractions de $B$ ; soit $A$ l'anneau de polynômes à une indéterminée $B[T]$. Il y a dans $A$ des idéaux maximaux de rang 2, par exemple $(\pi) + (T)$ ; mais il y a aussi des idéaux maximaux de rang 1, par exemple l'idéal principal $(\pi T - 1)$, dont l'anneau local complété est $\hat{K}[[T]]$.

Si $k$ est un corps, $A$ une algèbre intègre de type fini sur $k$, $\mathfrak{m}$ un idéal maximal de $A$, on sait que la dimension de $A_{\mathfrak{m}}$ est égale au degré de transcendance $\deg.\operatorname{tr.}_k K$ sur $k$ du corps des fractions $K$ de $A$. Comme $A$ est quotient d'un anneau régulier, il vérifie la condition des chaînes et est par suite (prop. 5) strictement équidimensionnel.

Proposition 7. — Si $X$ est un préschéma algébrique sur un corps $k$, de point générique $z$, alors $X$ est strictement équidimensionnel et de dimension égale à $\deg.\operatorname{tr.}_k \kappa(z)$.

On peut en effet se borner au cas où $X$ est réduit (les dimensions de $X$ et de $X_{\text{red}}$ étant évidemment les mêmes), et appliquer la remarque précédente à tout ouvert affine $U$ de $X$ (qui contient nécessairement $z$, donc est un schéma intègre admettant $z$ comme point générique). Cette démonstration prouve en outre :

Corollaire 1. — Sous les conditions de la prop. 7, la dimension de $X$ est égale à celle de toute partie ouverte non vide de $X$.

Il convient de noter que cette dernière propriété n'est pas valable pour tous les préschémas intègres, même strictement équidimensionnels. Par exemple si $B$ est un anneau de valuation discrète, $\mathfrak{p}$

\newpage

%% ===== Page 145 =====

- 109 -

son idéal maximal, $X=\operatorname{Spec}(B)$ a deux points, $(0)$ et $\mathfrak{p}$, ce dernier étant le seul point fermé ; on a $\dim(X)=1$ mais $\dim\{(0)\}=0$.

Corollaire 2. — Soit $X$ un préschéma algébrique sur un corps $k$. Pour tout point $x \in X$, on a $\dim_x(X)=\sup_i(\dim(X_i))$, où $X_i$ parcourt l'ensemble des composantes irréductibles de $X$ contenant $x$. Si $\{x\}$ est fermé on a $\dim_x X = \operatorname{codim}_X \overline{\{x\}} = \dim(\mathcal{O}_{X,x})$.

En effet, pour tout ouvert $U$ contenant $x$, on a (n°1, cor. de la prop.1) $\dim_x(X)=\sup_i(\dim(X_i))$ et d'après le cor. 1, $\dim_x(X_i) = \dim(X_i)$. D'autre part, comme toute chaîne de parties fermées irréductibles contenant $x$ est contenue dans l'un des $X_i$, on a $\operatorname{codim}_X(\{x\}) = \sup_i(\operatorname{codim}_{X_i}(\{x\}))$, et en vertu de la formule (6) (applicable puisque $X_i$ est strictement équidimensionnel), $\operatorname{codim}_{X_i}(\{x\}) = \dim(X_i)$, d'où la seconde assertion du corollaire.

Corollaire 3. — Soient $X, Y$ deux préschémas algébriques sur un corps $k$. Pour tout $k$-morphisme $f : X \to Y$, on a $\dim(\overline{f(X)}) \leq \dim(X)$.

En considérant les composantes irréductibles de $\overline{f(X)}$, on est aussitôt ramené au cas où $Y$ est irréductible et $f$ dominant (n°1, prop.1 (ii)) ; en outre, au moins une des composantes irréductibles $X_i$ de $X$ est telle que $f(X_i)$ soit dense dans $Y$, donc on peut aussi supposer $X$ irréductible. Alors, si $x$ est le point générique de $X$, $f(x)=y$ est le point générique de $Y$, et comme $f$ correspond à un $k$-homomorphisme de $\kappa(y)$ dans $\kappa(x)$, on a $\deg.\operatorname{tr.}_k \kappa(y) \leq \deg.\operatorname{tr.}_k \kappa(x)$, d'où le corollaire.

Ici encore, ce résultat ne s'étend pas à des schémas intègres quelconques. Par exemple, soient $B$ un anneau de valuation discrète, $K$ son corps des fractions, $X=\operatorname{Spec}(K)$, $Y=\operatorname{Spec}(B)$ ; l'injection canonique $B \to K$ correspond un morphisme $f : X \to Y$ de type fini ($K$ étant en effet engendré par $1/\pi$, si $\pi$ est l'idéal maximal de $B$) ; mais on a $\dim(Y)=1$ et $\dim(X)=0$.

Corollaire 4. — Soient $X, Y$ deux préschémas algébriques sur $k$ ; on a

\newpage

%% ===== Page 146 =====

- 110 -

$$\dim(X \times_k Y) = \dim(X) + \dim(Y) \ ;$$
en particulier, pour tout corps $k$, extension de $k$, $\dim(X \otimes_k K) = \dim(X)$.

Si $(U_i)$ et $(V_j)$ sont des recouvrements ouverts affines de $X$ et $Y$, $(U_i \times V_j)$ est un recouvrement ouvert affine de $X \times_k Y$, donc on peut se borner au cas où $X$ et $Y$ sont affines ($n^\circ 1$, cor. 1 de la prop. 2). D'autre part, si $X_i$ (resp. $Y_j$) sont les composantes irréductibles de $X$ (resp. $Y$), $X \times_k Y$ est réunion des $X_i \times Y_j$, qui sont fermés (§ 2, $n^\circ 2$, prop. 5), donc ($n^\circ 1$, prop. 1 (ii)) on peut supposer $X$ et $Y$ irréductibles ; enfin, comme $(X \times_k Y)_{\text{red}}$ et $(X_{\text{red}} \times_k Y_{\text{red}})_{\text{red}}$ sont identiques (§ 3, $n^\circ 4$, cor. de la prop. 11), on peut supposer $X$ et $Y$ réduits. Soient donc $X = \operatorname{Spec}(A)$, $Y = \operatorname{Spec}(B)$, où $A$ et $B$ sont des algèbres intègres de type fini sur $k$. Si $m = \dim(X)$, $n = \dim(Y)$, il résulte du lemme de normalisation (Bourbaki, Alg. comm.) que $A$ (resp. $B$) est un module de type fini sur un sous-anneau $A_1 = k[a_1, \dots, a_m]$ (resp. $B_1 = k[b_1, \dots, b_n]$) où les $a_i$ (resp. $b_j$) sont algébriquement indépendants sur $k$ ; $C = A \otimes_k B$ est donc un module de type fini sur $C_1 = A_1 \otimes_k B_1$. Pour tout idéal premier minimal $\mathfrak{p}$ dans $C$, on voit par suite que le degré de transcendance sur $k$ du corps des fractions de $C/\mathfrak{p}$ est au plus $m+n$ ; d'autre part, en vertu du th. de Cohen-Seidenberg, il existe un idéal premier $\mathfrak{q}$ de $C$ tel que $\mathfrak{q} \cap C_1 = \{0\}$, donc il y a un des idéaux premiers minimaux dans $C$ tel que le corps des fractions de $C/\mathfrak{p}$ soit de degré fini sur le corps des fractions de $C_1$, et par suite de degré de transcendance $m+n$ sur $k$.

Proposition 8. — Soient $X, Y$ deux préschémas localement noethériens, $f: X \to Y$ un morphisme. Soient $y \in Y$, et soit $x$ le point générique d'une composante irréductible de $f^{-1}(y)$, considéré comme préschéma sur $\kappa(y)$ (§ 2, $n^\circ 7$, prop. 17). Alors $\operatorname{codim}_X \{x\} \leqslant \operatorname{codim}_Y \{y\}$ (ou, ce qui revient au même, $\dim \mathcal{O}_x \leqslant \dim \mathcal{O}_y$).

La question est locale en $x$ et en $y$, puisque toute composante

\newpage

%% ===== Page 147 =====

- 111 -

irréductible de $f^{-1}(y)$, de point générique $x$, est contenue dans l'adhérence d'un voisinage ouvert affine de $x$. On peut aussi remplacer $Y$ par le schéma local noethérien $Y = \operatorname{Spec}(B)$, où $B$ est un anneau local, d'idéal maximal $\mathfrak{m} = \mathfrak{j}_y$, et $X = \operatorname{Spec}(A)$, où $A$ est une algèbre noethérienne sur $B$ ; on peut en outre (remplaçant $B$ par son image dans $A$) supposer $B \subset A$. Alors il s'agit de voir que si $\mathfrak{p}$ est un idéal premier minimal associé à $\mathfrak{m}A$, on a $\operatorname{rg}(\mathfrak{p}) \le \operatorname{rg}(\mathfrak{m})$. Or, dans l'anneau local $A_{\mathfrak{p}}$, l'idéal maximal $\mathfrak{p}A_{\mathfrak{p}}$ est engendré par $\mathfrak{m}A_{\mathfrak{p}}$, donc admet un système de $r$ générateurs, si $r$ est le rang de $\mathfrak{m}$, ce qui démontre notre assertion.

Corollaire 1. - Soient $X, Y$ des préschémas localement noethériens, $f: X \to Y$ un morphisme, $Y'$ une partie fermée irréductible de $Y$, $X'$ une composante irréductible de $f^{-1}(Y')$, dont l'image dans $Y$ soit dense dans $Y'$. Alors $\operatorname{codim}_X(X') \le \operatorname{codim}_Y(Y')$.

Si $x$ et $y$ sont les points génériques de $X'$ et $Y'$ respectivement, l'hypothèse signifie que $f(x) = y$, et le cor. 1 n'est qu'une autre forme de la prop. 8.

Corollaire 2. - Soient $X$ un préschéma localement noethérien, $f$ une section de $\mathcal{O}_X$, qui n'est identiquement nulle dans aucune composante irréductible de $X$. Soit $Z_f$ l'ensemble des $x \in X$ tels que $f(x) = 0$. Alors toute composante irréductible de $Z_f$ est de codimension 1.

Il suffit de démontrer qu'il en est ainsi dans un voisinage ouvert affine de tout point de $Z_f$, donc on peut supposer $X = \operatorname{Spec}(A)$ affine d'anneau $A$ noethérien ; par hypothèse $f$ n'est contenu dans aucun des idéaux premiers minimaux de $A$, et tous les idéaux premiers minimaux associés à l'idéal principal $(f)$ sont donc de rang 1 ("Hauptidealsatz" de Krull ; cf. Bourbaki, Alg. comm.), ce qui prouve le corollaire. Ce dernier se déduit aussi de la prop. 8, en considérant...

\newpage

%% ===== Page 148 =====

- 112 -

dérant $f$ comme une $X$-section du préschéma $S = X \otimes_{\mathbb{Z}} \mathbb{Z}[T]$ (§ 5, n°1) ; si $G$ est le graphe de ce morphisme, $g$ sa projection sur $Y = \operatorname{Spec}(\mathbb{Z}[T])$, $Z_f$ s'identifie à $g^{-1}(y)$, où $\mathfrak{j}_y$ est l'idéal premier $(T)$ ; la prop. 8 est donc applicable et donne $\operatorname{codim}_X(Z_f) \le 1$. L'hypothèse implique $\operatorname{codim}_X(M) \le 1$ pour toute composante irréductible $M$ de $Z_f$ ; par ailleurs, on ne peut avoir $\operatorname{codim}_X(M) = 0$ en vertu de l'hypothèse (n°1, cor. 2 de la prop. 3).

Corollaire 3. — La conclusion du cor. 2 est encore valable lorsqu'on remplace $\mathcal{O}_X$ par un $\mathcal{O}_X$-module localement isomorphe à $\mathcal{O}_X$.

Cela résulte du fait que la question est locale.

3. Dimension des fibres d'un morphisme de type fini : formule de dimension, théorème de Chevalley.

Proposition 9. — Soient $A$ un anneau noethérien, $T$ une indéterminée, $\mathfrak{p}$ un idéal premier de $A$ ; alors $\mathfrak{p} = \mathfrak{p}A[T]$ est premier dans $A[T]$ et $\mathfrak{p}' \cap A = \mathfrak{p}$. Il existe une infinité d'autres idéaux premiers de $A[T]$ dont l'intersection avec $A$ est $\mathfrak{p}$ ; ces idéaux sont deux à deux sans relation d'inclusion. En outre, si $\mathfrak{q}$ est un tel idéal, on a
$$(8) \quad \operatorname{rg}(\mathfrak{q}) = \operatorname{rg}(\mathfrak{p}') + 1 = \operatorname{rg}(\mathfrak{p}) + 1$$

Dans les premières assertions, on se ramène aussitôt, en remplaçant $A$ par $A/\mathfrak{p}$, au cas où $\mathfrak{p} = (0)$ ; elles résultent alors du fait que les idéaux premiers de $A[T]$ dont l'intersection avec $A$ se réduit à $0$ sont exactement ceux ne rencontrant pas la partie multiplicative $S = A - \{0\}$ de l'anneau intègre $A$, donc correspondent biunivoquement aux idéaux premiers de $S^{-1}A[T] = K[T]$, où $K$ est le corps des fractions de $A$. La définition du rang d'un idéal premier implique alors $\operatorname{rg}(\mathfrak{q}) \ge \operatorname{rg}(\mathfrak{p}') + 1 \ge \operatorname{rg}(\mathfrak{p}) + 1$ ; il reste à prouver que $\operatorname{rg}(\mathfrak{p}) + 1 \ge \operatorname{rg}(\mathfrak{q})$ (ce qui est la seule partie de la démonstration où intervient l'hypothèse que $A$ est noethérien). Si $\operatorname{rg}(\mathfrak{p}) = r$, il suffit de trouver un système de $r+1$ éléments de $\mathfrak{q}$ engendrant un idéal de définition de

\newpage

%% ===== Page 149 =====

(A[T])_{\mathfrak{q}} ; si $K$ est le corps des fractions de $A/\mathfrak{p}$, il y a un élément $x \in \mathfrak{q}$ dont la classe mod $\mathfrak{p}$ engendre dans $K[T]$ l'idéal premier correspondant à $\mathfrak{q}/\mathfrak{p}$ ; si $x_1, \dots, x_r$ sont des éléments de $A$ formant un système de générateurs d'un idéal de définition de $A_{\mathfrak{p}}$, il est clair que les $r+1$ éléments $x, x_1, \dots, x_r$ répondront à la question.

Corollaire 1. — Pour tout anneau noethérien $A$, on a $\dim(A[T_1, \dots, T_r]) = \dim(A)+r$ ($T_i$ indéterminées).

Corollaire 2. — Pour tout préschéma localement noethérien $X$, la dimension de $X \otimes_{\mathbb{Z}} \mathbb{Z}[T_1, \dots, T_r]$ est $\dim(X)+r$.

Cela résulte du cor. 1 et du n°1, cor. 1 de la prop. 2.

Corollaire 3. — Sous les hypothèses de la prop. 9, soit $\mathfrak{q}$ un idéal premier de $A[T]$ tel que $\mathfrak{q} \cap A = \mathfrak{p}$ ; si $k$ et $k'$ sont les corps résiduels de $A_{\mathfrak{p}}$ et $(A[T])_{\mathfrak{q}}$, on a
$$(9) \qquad \dim(A_{\mathfrak{p}}) + 1 = \dim((A[T])_{\mathfrak{q}}) + \operatorname{deg.tr.}_k k'$$

Si $\mathfrak{q} = \mathfrak{p}[T]$, on a, d'après la prop. 9, $\operatorname{rg}(\mathfrak{q}) = \operatorname{rg}(\mathfrak{p})$, donc $\dim(A_{\mathfrak{p}}) = \dim((A[T])_{\mathfrak{q}})$ ; comme dans ce cas $k' = k(T)$, la formule (9) est bien vérifiée. Dans le cas contraire, $\operatorname{rg}(\mathfrak{q}) = \operatorname{rg}(\mathfrak{p}) + 1$, et comme $\mathfrak{q}$ correspond à un idéal premier $(0)$ de $k[T]$, $k'$ est une extension algébrique de $k$, et on a encore la formule (9).

Lemme 1. — Soient $A$ un anneau local intègre noethérien, $\mathfrak{m}$ son idéal maximal, $k$ son corps résiduel. Soient $B$ un anneau intègre contenant $A$, tel que $B = A[x]$, où $x \in B$, $\mathfrak{q}$ un idéal premier de $B$ tel que $\mathfrak{q} \cap A = \mathfrak{m}$, $k'$ le corps résiduel de $B_{\mathfrak{q}}$. Alors on a
$$(10) \qquad \dim(A_{\mathfrak{m}}) + \operatorname{deg.tr.}_A B_{\mathfrak{q}} \ge \dim(B_{\mathfrak{q}}) + \operatorname{deg.tr.}_k k'$$
où on a noté $\operatorname{deg.tr.}_A B_{\mathfrak{q}}$ le degré de transcendance du corps des fractions de $B$ sur celui de $A$. En outre, si $A[T]$ vérifie la condition des chaînes, les deux membres de (10) sont égaux.

Si $x$ est transcendant sur $A$, les deux membres de (10) sont égaux en vertu de (9). Dans le cas contraire, on a $B = A[T]/\mathfrak{p}$, où $\mathfrak{p}$ est

\newpage

%% ===== Page 150 =====

- 114 -

un idéal premier $\neq (0)$ de $A[T]$, tel que $\mathfrak{p} \cap A = (0)$, donc de rang 1 en vertu de la prop. 9. L'idéal $\mathfrak{q}$ de $B$ correspond à un idéal premier $\mathfrak{r} \supset \mathfrak{p}$ de $A[T]$ tel que $\mathfrak{r} \cap A = \mathfrak{m}$, et on a $B_{\mathfrak{q}} = A[T]_{\mathfrak{r}} / \mathfrak{p} A[T]_{\mathfrak{r}}$ ; on a donc, par la formule (1) du n°1
$$(11) \qquad \dim(A[T]_{\mathfrak{r}}) \ge \operatorname{rg}(\mathfrak{p}A[T]_{\mathfrak{r}}) + \dim(B_{\mathfrak{q}}) = 1 + \dim(B_{\mathfrak{q}})$$
car $\operatorname{rg}(\mathfrak{p}A[T]_{\mathfrak{r}}) = \operatorname{rg}(\mathfrak{p}) = 1$ ; en outre les deux membres de (11) sont égaux lorsque $A[T]$ (et par suite $A[T]_{\mathfrak{r}}$) vérifie la condition des chaînes (prop. 5 et 9). D'autre part, la formule (9) donne
$$\dim(A[T]_{\mathfrak{r}}) = \dim(A) + 1 - \operatorname{deg.tr.}_k k'$$
puisque les corps résiduels de $B_{\mathfrak{q}}$ et de $(A[T])_{\mathfrak{r}}$ sont les mêmes. Enfin on a alors $\operatorname{deg.tr.}_A B$, d'où finalement (10).

Théorème 1 (formule de dimension). — Soient $A$ un anneau local intègre noethérien, $B$ un anneau intègre contenant $A$ et de type fini sur $A$, $\mathfrak{q}$ un idéal premier de $B$ tel que $\mathfrak{q} \cap A$ soit l'idéal maximal de $A$, $k$ et $k'$ les corps résiduels de $A$ et de $B_{\mathfrak{q}}$. On a alors l'inégalité
$$(12) \qquad \dim(A) + \operatorname{deg.tr.}_A B \ge \dim(B_{\mathfrak{q}}) + \operatorname{deg.tr.}_k k'$$
En outre les deux membres sont égaux si toute sous-algèbre de type fini de $B$ vérifie la condition des chaînes (en particulier si $A$ est quotient d'un anneau local régulier).

Supposons que $B = A[x_1, \dots, x_n]$ ; et raisonnons par récurrence sur $n$ ; posons $C = A[x_1, \dots, x_{n-1}]$ ; et $\mathfrak{r} = \mathfrak{q} \cap C$ ; $C_{\mathfrak{r}}$ est un anneau local intègre noethérien, et si on pose $B_{\mathfrak{s}} = \mathfrak{s}^{-1}B$, $\mathfrak{q}_{\mathfrak{s}} = \mathfrak{s}^{-1}\mathfrak{q}$, où $\mathfrak{s} = C - \mathfrak{r}$, on a $B_{\mathfrak{q}} = (B_{\mathfrak{s}})_{\mathfrak{q}_{\mathfrak{s}}}$, $\mathfrak{q}_{\mathfrak{s}} \cap C_{\mathfrak{r}} = \mathfrak{r}C_{\mathfrak{r}}$ et $B_{\mathfrak{s}} = C_{\mathfrak{r}}[x_n]$ ; en outre les corps des fractions de $B_{\mathfrak{s}}$ et $C_{\mathfrak{r}}$ sont ceux de $B$ et de $C$. Si $k_1$ est le corps résiduel de $C_{\mathfrak{r}}$, on a donc, d'après (10)
$$\dim(C_{\mathfrak{r}}) + \operatorname{deg.tr.}_A B \ge \dim(B_{\mathfrak{q}}) + \operatorname{deg.tr.}_{k_1} k'$$
D'autre part, l'hypothèse de récurrence donne
$$\dim(A) + \operatorname{deg.tr.}_A C \ge \dim(C_{\mathfrak{r}}) + \operatorname{deg.tr.}_k k_1$$
d'où (12) en ajoutant membre à membre ; en outre, si

\newpage

%% ===== Page 151 =====

- 115 -

toute sous-algèbre de type fini de $C[T]$ vérifie la condition des chaînes, il en est de même de toute sous-algèbre de type fini de $C_r[T]$, ce qui achève la démonstration.

Corollaire 1. — Soient $Y$ un préschéma irréductible localement noethérien, $X$ un préschéma irréductible, $f : X \to Y$ un morphisme dominant de type fini. Soit $\xi$ (resp. $\eta$) le point générique de $X$ (resp. $Y$) et soit $e = \dim(f^{-1}(\eta)) = \operatorname{deg.tr.}_{k(\eta)} k(\xi)$ ("dimension de la fibre générique"). Pour tout $x \in X$, on a, en posant $y = f(x)$
(13)
$$\operatorname{codim}\{y\} \ge \operatorname{deg.tr.}_{k(y)} k(x) + \operatorname{codim}\{x\}$$
En outre, les deux membres sont égaux lorsque $\mathcal{O}_y$ est un quotient d'anneau local régulier.

Comme la question est locale, on peut supposer $X$ et $Y$ affines, soient $Y = \operatorname{Spec}(A)$, $X = \operatorname{Spec}(B)$, où $B$ est une algèbre intègre de type fini sur $A$; comme $f$ est dominant, $A$ s'identifie à une sous-algèbre de $B$ (§ 1, n°2, cor. 4 de la prop. 5), donc $\mathcal{O}_y$ à une sous-algèbre de $\mathcal{O}_x$. Comme $k(\eta)$ et $k(\xi)$ sont les corps de fractions de $A$ et $B$ respectivement, donc de $\mathcal{O}_y$ et $\mathcal{O}_x$ respectivement, la formule (13) résulte de (12) appliquée aux anneaux $\mathcal{O}_y$ et $\mathcal{O}_x$.

On notera d'ailleurs que $\operatorname{deg.tr.}_{k(y)} k(x)$ n'est autre que la dimension de $f^{-1}(y) \cap \{x\}$, comme il résulte de la prop. 7 du n°2, appliquée à $f^{-1}(y)$ considéré comme préschéma algébrique sur $k(y)$. En particulier :

Corollaire 2. — Les hypothèses du cor. 1 étant vérifiées, supposons en outre que $x$ soit fermé dans $f^{-1}(y)$ et que $\mathcal{O}_y$ soit un quotient d'anneau régulier. Alors
(14)
$$\operatorname{codim}\{x\} = e + \operatorname{codim}\{y\}$$

Corollaire 3. — Les hypothèses du cor. 1 étant vérifiées, on a

\newpage

%% ===== Page 152 =====

- 116 -

$$(15) \quad \dim(X) \leqslant \dim(Y) + e$$

Supposons en outre que tous les anneaux locaux $\mathcal{O}_y$ ($y \in Y$) soient quotients d'anneaux réguliers. Pour que les deux membres de (15) soient égaux, il faut et il suffit que $\dim(Y) = \sup_{x \in X} (\dim(\mathcal{O}_{f(x)}))$.

L'inégalité (15) résulte en effet de (13) appliqué à un point fermé $x$ de $X$ ($n^\circ 2$, cor. 2 de la prop. 6). Si les $\mathcal{O}_y$ sont quotients d'anneaux réguliers, l'égalité (14) montre que $\dim(X) = e + \sup_{x \in X} (\dim(\mathcal{O}_{f(x)}))$.

En particulier :

Sous les hypothèses du cor. 1,

\underline{Corollaire 4. - On a $\dim(X) = \dim(Y) + e$ dans chacun des deux cas suivants :}

$1^\circ$ $X$ et $Y$ sont des préschémas algébriques sur un corps $k$ ;

$2^\circ$ tous les anneaux $\mathcal{O}_y$ ($y \in Y$) sont quotients d'anneaux réguliers, et $f(X) = Y$.

La seconde assertion découle immédiatement du cor. 3 ainsi que du $n^\circ 2$, cor. 2 de la prop. 6 ; la première est une conséquence directe de la prop. 7 du $n^\circ 2$.

Corollaire 5. - Soient $Y$ un préschéma localement noethérien, $f: X \to Y$ un morphisme de type fini, et posons $e = \sup_{y \in Y} (\dim(f^{-1}(y)))$. Alors

$$(16) \quad \dim(X) \leqslant \dim(Y) + e$$

Si en outre tous les $\mathcal{O}_y$ sont quotients d'anneaux réguliers, et si $\dim(f^{-1}(y)) = e$ pour tout $y \in Y$, les deux membres de (16) sont égaux.

Comme la dimension de $X$ est la borne supérieure des dimensions de ses composantes irréductibles, on est ramené au cas où $X$ est irréductible ; alors $\overline{f(X)}$ est irréductible, étant l'adhérence de $f(X)$ \marginpar{; on peut aussi supposer $X$ et $Y$ réduits ($34, n^\circ 2$, cor. 3 de la prop. 9)} où $x$ est le point générique de $X$ ; comme $\dim(f(X)) \leqslant \dim(Y)$, on est ramené au cas où $X$ et $Y$ sont irréductibles et $f$ dominant, et le corollaire est alors conséquence du cor. 3.

Remarques. - Nagata a démontré que l'égalité peut être en défaut dans la formule (12), même si $B$ est entier sur $A$ et $\mathfrak{q}$ maximal dans $B$ ; autrement dit, il se peut que dans ces conditions, on ait

\newpage

%% ===== Page 153 =====

- 117 -

$\operatorname{rg}(\mathfrak{q}) = \dim(B_{\mathfrak{q}}) < \dim(A)$ (Nagata, \emph{On the chain problem of prime ideals}, Nagoya Math. Jour., 10 (1956)). On dit avec Nagata qu'un anneau local intègre noethérien $A$ vérifie la condition des chaînes stricte ("second chain condition") si la clôture intégrale de $A$ est un anneau strictement équidimensionnel (ou, ce qui revient au même, lorsqu'il en est ainsi pour toute $A$-algèbre finie sur $A$ contenue dans cette clôture). Alors on peut montrer que toute $A$-algèbre de type fini vérifie la condition des chaînes (Nagata, \emph{Note on a chain condition of prime ideals}, à paraître dans Kyoto J. Math.) ; \emph{a fortiori}, l'égalité a lieu dans (12).

\emph{Lemme 2.} - Sous les hypothèses du cor. 1 du th. 1, toutes les composantes de $f^{-1}(y)$ sont de dimension $\ge e$.

Si les deux membres de (13) sont égaux, le lemme résulte de cette égalité appliquée à un point générique $x$ d'une composante irréductible de $f^{-1}(y)$, en tenant compte de la prop. 8 du n° 2. En fait, Nagata a démontré (dans l'article du Kyoto J. Math. précité) que le lemme 2 est valable sans restriction sur les $\mathcal{O}_y$, mais sa démonstration est délicate. Au chap. III, § 4, nous obtiendrons une démonstration directe simple du th. 2 ci-dessous (et \emph{a fortiori} du lemme 2 qui en est une conséquence), comme corollaire du "main theorem" de Zariski ; jusque là, le § actuel excepté, nous n'utiliserons le lemme 2 et le th. 2 que dans les cas où les deux membres de (13) sont égaux.

\emph{Lemme 3.} - Si $A$ est un anneau, $B$ un anneau entier sur $A$, on a
$$\dim(B) \le \dim(A).$$

En effet, si $\mathfrak{p}, \mathfrak{q}$ sont deux idéaux premiers de $B$ tels que $\mathfrak{p} \subset \mathfrak{q}$ et $\mathfrak{p} \cap A = \mathfrak{q} \cap A$, on a nécessairement alors $\mathfrak{p} = \mathfrak{q}$ (Bourbaki, Alg. comm.).

\newpage

%% ===== Page 154 =====

\marginpar{Archives}
\marginpar{Jothentick Sept 59}

- 118 -

Proposition 10. -- Soient $Y$ un préschéma irréductible localement noethérien, $X$ un préschéma irréductible, $f : X \to Y$ un morphisme dominant de type fini. $\xi$ (resp. $\eta$) le point générique de $X$ (resp. $Y$).
(i) Il existe un ouvert non vide $U \subset Y$ tel que pour tout $y \in U$, la fibre $f^{-1}(y)$ soit équidimensionnelle de dimension $e = \deg.tr._{\kappa(y)}\kappa(\xi)$.
(ii) Pour tout $y \in f(X)$, toute composante de $f^{-1}(y)$ est de dimension $\ge e$.

(ii) n'est autre que le lemme 2. Pour démontrer (i), on peut évidemment supposer $X$ et $Y$ réduits (§ 4, n°2, cor.3 de la prop.9), donc intègres. On peut évidemment supposer $X$ affine, et $X$ réunion d'un nombre fini d'ouverts affines $X_i$ ; si à chaque $X_i$ correspond un ouvert non vide $U_i \subset Y$ tels que les composantes irréductibles de $X_i \cap f^{-1}(y)$ soient de dimension $e$ pour tout $y \in U_i$, l'ouvert $U = \bigcap_i U_i$ vérifiera la condition de (i). On peut donc aussi supposer $X$ affine, et la proposition sera conséquence du lemme suivant :

Lemme 4. -- Soient $X, Y$ deux préschémas intègres, $f : X \to Y$ un morphisme dominant de type fini, tel que $f^{-1}(y)$ soit affine pour tout ouvert affine $V \subset Y$ (de tels morphismes seront étudiés au chap.II, § 2, sous le nom de morphismes affines). Il existe alors un ouvert affine non vide $U$ de $Y$, tel que, si on désigne par $C$ l'anneau de $U \otimes_{\mathbb{Z}} k[T_1, \dots, T_e]$ ($T_i$ indéterminées, $e = \deg.tr._{\kappa(y)}\kappa(X)$), l'anneau de $f^{-1}(U)$ soit entier fini sur $C$.

Supposons en effet le lemme démontré. Alors, pour $y \in U$, $f^{-1}(y)$ considéré comme préschéma sur $\kappa(y)$, est un schéma affine dont l'anneau, isomorphe à $D \otimes_A \kappa(y)$ (B anneau de $U$) est entier fini sur $\kappa(y)[T_1, \dots, T_e]$; le lemme 3 et le cor. 1 de la prop.9 montrent alors que $\dim(f^{-1}(y)) \le e$ ce qui, tenant compte de (ii), achève de prouver (i).

Pour établir le lemme 4, on peut évidemment supposer $Y$ affine d'anneau $A$, donc $X$ affine d'anneau $B$, où $B$ est une $A$-algèbre intègre.

\newpage

%% ===== Page 155 =====

\marginpar{pour}
(contenant $A$ de type fini sur $A$). Soit $K$ le corps des fractions de $A$ ; $B \otimes_A K$ est une $K$-algèbre intègre de type fini ; en vertu du lemme de normalisation (sur $K$), il existe des éléments algébriquement indépendants $c_1, \dots, c_e$ de $B \otimes_A K$ tels que $B \otimes_A K$ soit un anneau entier fini sur l'anneau $B' = K[c_1, \dots, c_e]$. On peut écrire $c_i = b_i/a$, où $a \in A$, $b_i \in B$ ; d'autre part, si $(x_j)$ est un système fini de générateurs de $B$ sur $A$, les $x_j$ sont entiers sur $B'$, donc sur un anneau de la forme $A_g[c_1, \dots, c_e]$ ($g \in A$). Si on pose $g = af$, on voit que les $x_j$, et a fortiori les $x_j/g$, sont entiers sur $A_g[b_1, \dots, b_e]$ ; comme les $b_i$ sont algébriquement indépendants sur le corps des fractions $K$ de $A$, il suffit de prendre $U = D(g)$ pour répondre à la question.

Théorème 2. — Soient $Y$ un préschéma noethérien $f: X \to Y$ un morphisme de type fini. Pour tout entier $n$, l'ensemble des $x \in X$ tels que $\dim_x(f^{-1}(f(x))) \ge n$ est fermé.

La question étant locale sur $Y$, on peut supposer $Y$ noethérien. Si $\mathcal{E}$ est l'ensemble des parties fermées $Y' \subset Y$ telles que le théorème soit vrai pour $Y'$ et $f^{-1}(Y')$, il suffit de prouver que le complémentaire $E'$ de $\mathcal{E}$ dans l'ensemble des parties fermées de $Y$ est vide. Dans le cas contraire $\mathcal{E}'$ admet un élément minimal $Y_0$, puisque $Y$ est noethérien ; remplaçant $Y$ par $Y_0$, on est ramené à démontrer le théorème en supposant qu'il est vrai pour toute partie fermée $Y' \subsetneq Y$ de $Y$. D'autre part, si $X_i$ ($1 \le i \le m$) sont les composantes irréductibles de $X$, on a $F_n(X) = \bigcup_i F_n(X_i)$ (n°1, cor. de la prop. 1), et on peut donc supposer $X$ irréductible ; enfin, comme on peut supposer $X$ et $Y$ réduits (§ 4, n°2, cor. 3 de la prop. 9), on peut remplacer $Y$ par $f(X)$ (§ 3, n°4, cor. de la prop. 9), donc on peut supposer $f$ dominant et $Y$ irréductible. En vertu de la prop. 10 (ii), on a $F_n(X) = X$ pour $n \le e$, on peut par suite supposer $n > e$. Mais alors, en vertu de la prop. 10 (i), il existe une partie fermée $Y' = Y - U$ distincte

\newpage

%% ===== Page 156 =====

- 120 -

de $Y$ telle que $F_n(X) \subset f^{-1}(Y')$, ce qui termine la démonstration.

\emph{Corollaire 1.} - L'ensemble des $x \in X$ tels que $x$ soit isolé dans $f^{-1}(f(x))$ est ouvert dans $X$.

C'est en effet le complémentaire de $F_1(X)$.

\emph{Corollaire 2.} - Sous les hypothèses du th. 2, supposons de plus que l'application $f$ soit fermée. Alors, pour tout entier $n$, l'ensemble des $y \in Y$ tels que $\dim(f^{-1}(y)) \geq n$ est fermé.

En effet, cet ensemble n'est autre que $f(F_n(X))$ ($n^\circ 2, \text{prop.} 2 \text{ (ii)}$).

\emph{Remarque.} - On notera que la prop. 10 (et par suite le lemme 2) sont conséquences du th. 2 : en effet, sous les hypothèses de la prop. 10, l'ensemble des $x$ tels que $\dim(f^{-1}(f(x))) < e$ est ouvert, et comme il ne peut contenir $\xi$, il est vide ; de même, l'ensemble des $x$ tels que $\dim(f^{-1}(f(x))) = e$ est ouvert et non vide, puisqu'il contient $\xi$.

Signalons enfin le résultat plus facile suivant :

\emph{Proposition 11.} - Soit $Y$ un préschéma dont l'espace de base est quasi compact. Si $f: X \to Y$ est un morphisme de type fini, il existe un entier $n$ tel que, pour tout $y \in Y$, $\dim(f^{-1}(y)) \leq n$.

Comme il y a un recouvrement affine fini $(V_i)$ de $Y$ tel que $f^{-1}(V_i)$ soit réunion finie d'ouverts affines, on est ramené aussitôt au cas où $X=\operatorname{Spec}(B)$ et $Y=\operatorname{Spec}(A)$ sont affines ($n^\circ 1, \text{prop.} 1 \text{ (ii)}$). Si $B$ est une $A$-algèbre engendrée par $n$ éléments, pour tout $y \in Y$, l'anneau du schéma $f^{-1}(y)$ sur $\kappa(y)$ est une algèbre ayant $n$ générateurs, donc $\dim(f^{-1}(y)) \leq n$ ($\text{cor.} 1 \text{ de la prop.} 9$).

\emph{4. Ensembles constructibles.}

\emph{Définition 4 (Chevalley).} - Soit $X$ un espace topologique. On appelle ensemble constructible dans $X$ une réunion finie de parties localement fermées de $X$.

\newpage

%% ===== Page 157 =====

- 121 -

Définition 5. — Soit $h$ une application d'un espace topologique $X$ dans un ensemble $T$. On dit que $h$ est constructible si $h^{-1}(t)$ est constructible pour tout $t \in T$, et vide sauf pour un nombre fini de valeurs de $t$ ; pour toute partie $S$ de $T$, $h^{-1}(S)$ est alors constructible.

On peut dire que l'ensemble des parties constructibles de $X$ est la plus petite partie de $\mathcal{P}(X)$ contenant les ensembles fermés et stable par intersection finie et passage au complémentaire (ce qui entraîne qu'elle est aussi stable par réunion finie). L'image réciproque d'un ensemble constructible par une application continue est constructible. Si $Y$ est une partie constructible de $X$, les parties de $Y$ qui sont constructibles en tant que sous-espaces de $Y$ sont identiques à celles qui sont constructibles en tant que sous-espaces de $X$.

Proposition 12. — Soit $X$ un espace topologique noethérien. Pour qu'une partie $E$ de $X$ soit constructible, il faut et il suffit que pour toute partie fermée irréductible $Y$ de $X$, $E \cap Y$ soit rare dans $Y$ ou contienne un ouvert non vide de $Y$.

La nécessité de la condition résulte de ce que $E \cap Y$ doit être une partie constructible de $Y$, et du

Lemme 5. — Soient $X$ un espace topologique irréductible, $E$ une partie constructible de $X$. Pour que $E$ soit dense dans $X$, il faut et il suffit que $E$ contienne un ouvert non vide de $X$. Si $X$ admet un point générique (c'est-à-dire si $X = \{\bar{x}\}$), il revient au même de dire que $x \in E$.

En effet, comme $X$ est irréductible, tout ouvert non vide est dense. Inversement, supposons que $E = \bigcup_i (U_i \cap F_i)$, où les $U_i$ sont ouverts non vides dans $X$ et les $F_i$ fermés dans $X$ ; on a donc $E \subset \bigcup_i F_i$, et par suite si $E \neq X$, $X$ est égal à un des $F_i$, d'où $E \subset F_i$. Si $X = \{\bar{x}\}$ et si $x \in E$, il est clair que $E$ est dense dans $X$ ; la réciproque résulte de

\newpage

%% ===== Page 158 =====

- 122 -

ce qui précède et de ce que tout ouvert non vide contient $x$.

Pour démontrer que la condition de la prop. 12 est suffisante, considérons l'ensemble $F$ des parties fermées $Y$ de $X$ telles que $E \cap Y$ soit constructible (par rapport à $Y$ ou par rapport à $X$, ce qui revient au même). Il s'agit de montrer que le complémentaire $F'$ de $F$ dans l'ensemble de toutes les parties fermées de $X$ est vide. Dans le cas contraire $F'$ contient un élément minimal $Y_0$ ; remplaçant $X$ par $Y_0$, on voit qu'on peut supposer que pour toute partie fermée $Y \neq X$ de $X$, $E \cap Y$ est constructible. Si $X$ est réductible et si $X_i$ ($1 \leq i \leq n$) sont ses composantes irréductibles, les $E \cap X_i$ sont constructibles par hypothèse, donc aussi $E$ qui est leur réunion (finie). Si $X$ est irréductible, ou bien $E$ est rare, donc [rayé: $E$ est constructible] $E \cap X = E$ est constructible ; ou bien $E$ contient un ouvert non vide $U$, donc est réunion de $U \cap E$ et de $E \cap (X-U)$ ; comme l'ensemble fermé $X-U$ est distinct de $X$, $E \cap (X-U)$ est constructible, et il en est par suite de même de $E$.

Corollaire. - Soient $X$ un espace noethérien, $(E_{\alpha})$ une famille filtrante croissante de parties constructibles de $X$, telle que :

1° $E = \bigcup_{\alpha} E_{\alpha}$.

2° Toute partie fermée irréductible de $X$ est contenue dans l'adhérence d'un des $E_{\alpha}$.

Alors il existe $\alpha$ tel que $X = \overline{E_{\alpha}}$.

Lorsque toute partie fermée irréductible de $X$ admet un point générique, l'hypothèse 2° peut être supprimée.

\newpage

%% ===== Page 159 =====

les $E_{\alpha}$. Soit $M$ l'ensemble des parties fermées de $X$ contenues dans l'un des $E_{\alpha}$ ; il s'agit de montrer que le complémentaire $M^c$ de $M$ dans l'ensemble de toutes les parties fermées de $X$ est vide. Sinon, $M^c$ contient un élément minimal ; on est donc ramené à prouver le corollaire en supposant que toute partie fermée $Y \neq X$ de $X$ est contenue dans un des $E_{\alpha}$. La proposition est alors évidente si $X$ est réductible, car chacune des composantes irréductibles $X_i$ de $X$ est contenue dans un $E_{\alpha_i}$, et il existe un $E_\beta$ contenant tous les $E_{\alpha_i}$. Supposons donc $X$ irréductible. Alors par hypothèse il existe un indice $\beta$ tel que $X = E_{\beta}$, donc (lemme 5) $E_{\beta}$ contient un ouvert non vide $U$. Mais alors l'ensemble fermé $X-U$ est contenu dans un $E_{\gamma}$, et il suffit de prendre $E_{\delta}$ contenant $E_{\beta}$ et $E_{\gamma}$. Lorsque toute partie fermée irréductible $Y$ de $X$ admet un point générique $y$, il existe $\alpha$ tel que $y \in E_{\alpha}$, donc $\{y\} \subset E_{\alpha}$, et la condition $2^{\circ}$ est conséquence de $1^{\circ}$.

\underline{Proposition 13}.-- Soient $X$ un espace noethérien, $x$ un point de $X$, $E$ une partie constructible de $X$. Pour que $E$ soit un voisinage de $x$, il faut et il suffit que pour toute partie fermée irréductible $Y$ de $X$ contenant $x$, $E \cap Y$ soit dense dans $Y$ (s'il existe un point générique de $Y$, cela signifie aussi (lemme 5) que $y \in E$).

La condition est évidemment nécessaire ; prouvons qu'elle est suffisante. Soit $M$ l'ensemble des parties fermées de $X$ contenant $x$ et tel que $E \cap Y$ soit un voisinage de $x$ dans $Y$ ; en considérant un ensemble minimal dans le complémentaire de $M$ par rapport à l'ensemble de toutes les parties fermées de $X$ contenant $x$, on se ramène au cas où pour toute partie fermée $Y \neq X$ de $X$ contenant $x$, $E \cap Y$ est un voisinage de $x$ dans $Y$. Si $X$ est réductible, chacune des composantes irréductibles $X_i$ de $X$ contenant $x$ est distincte de $X$, donc $E \cap X_i$ est un voisinage de $x$ par rapport à $X_i$, et par suite $E$ est un voisinage de $x$.

\newpage

%% ===== Page 160 =====

- 124 -

dans la réunion (ouverte) des composantes irréductibles de $X$ contenant $x$, et a fortiori dans $X$. Si $X$ est irréductible, $E$ contient une partie ouverte non vide $U$ ; la proposition est évidente si $x \in U$ ; sinon, par hypothèse $x$ est intérieur à $E \cap (X-U)$ par rapport à $X-U$, donc l'adhérence dans $X$ de $X-E$ ne contient pas $x$, ce qui achève la démonstration.

\emph{Proposition 14}. — Soit $h$ une application d'un espace noethérien $X$ dans un ensemble $T$. Pour que $h$ soit constructible, il faut et il suffit que pour toute partie fermée irréductible $Y$ de $X$, il existe dans $Y$ un ouvert non vide $U$ dans lequel $h$ est constante.

La condition est nécessaire : en effet, par hypothèse, $h$ ne prend dans $Y$ qu'un nombre fini de valeurs $t_i$, et chacun des ensembles $Y \cap h^{-1}(t_i)$ est constructible. Comme ils ne peuvent être tous rares, l'un d'entre eux contient un ensemble ouvert non vide (lemme 5). Pour voir que la condition est suffisante, considérons l'ensemble $M$ des parties fermées $Y$ de $X$ telles que la restriction de $h$ à $Y$ soit constructible ; en considérant un ensemble minimal dans le complémentaire de $M$ par rapport à l'ensemble de toutes les parties fermées de $X$, on se ramène au cas où pour toute partie fermée $Y \neq X$ de $X$, la restriction de $h$ à $Y$ est constructible. Si $X$ est réductible, la restriction de $h$ à chacune des composantes irréductibles $X_i$ de $X$ est constructible, d'où suit aussitôt (déf. 2) que $h$ est constructible. Si $X$ est irréductible, il existe par hypothèse une partie ouverte non vide $U$ de $X$ dans laquelle $h$ est constante ; d'autre part, la restriction de $h$ à $X-U$ est constructible par hypothèse, et il résulte aussitôt que $h$ est constructible.

\emph{Corollaire}. — Soit $X$ un espace noethérien dans lequel toute partie fermée irréductible admet un point générique. Si $h$ est une application de $X$ dans un ensemble $T$ telle que pour tout $t \in T$, $h^{-1}(t)$ soit

\newpage

%% ===== Page 161 =====

- 125 -

constructible , alors $h$ est constructible .

En effet , si $Y$ est une partie fermée irréductible de $X$ et $y$ son point générique , $h^{-1}(h(y))$ est constructible et contient $y$ , donc (lemme 5) contient une partie ouverte non vide de $Y$ , et il suffit d'appliquer la prop.14 .

Théorème 3 (Chevalley) . — Soient $Y$ un préschéma noethérien , $f: X \to Y$ un morphisme de type fini . Alors l'image par $f$ de toute partie constructible de $X$ est une partie constructible de $Y$ .

Il suffit de prouver que $f(Z)$ est constructible pour toute partie localement fermée $Z$ de $X$. On peut supposer $X$ et $Y$ réduits en considérant un sous-préschéma de $X$ ayant $Z$ pour espace de base (§ 3, n°4, prop.9) et la restriction de $f$ à $Z$ (§ 4, n°2, prop.7 et 8) on est ramené à prouver que $f(X)$ est constructible . Appliquons le critère de la prop.12 ; pour toute partie fermée irréductible $T$ de $Y$ , il faut prouver que $T \cap f(X) = f(f^{-1}(T))$ est rare dans $T$ ou contient un ouvert non vide de $T$ ; en prenant un sous-préschéma de $Y$ ayant $T$ pour espace de base (§ 3, n°4, prop 9) , et considérant son image réciproque par $f$ (§ 3, n°2, cor.3 de la prop.5) , on voit finalement qu'on est ramené à prouver que si on suppose $Y$ irréductible et $f(X)$ dense dans $Y$ , $f(X)$ contient un ouvert non vide de $Y$ . On peut supposer en outre $Y$ affine ; alors $X$ est par hypothèse réunion finie d'ouverts affines $V_i$ et comme $f(X)$ est dense dans $Y$ , un au moins des $f(V_i)$ est dense dans $Y$ (l'intersection d'un nombre fini d'ouverts non vides de $Y$ étant non vide) . On est ainsi ramené au cas où en outre $X$ et $Y$ sont affines ; enfin , si $(X_j)_{1 \le j \le n}$ est la famille des composantes irréductibles de $X$ , on voit comme ci-dessus que l'un des $f(X_j)$ est dense dans $Y$ ; on peut donc supposer $X$ irréductible . Alors $X = \operatorname{Spec}(B)$ , $Y = \operatorname{Spec}(A)$ , où $A$ est un anneau intègre noethérien , $B$ un anneau $A$-algèbre intègre de type fini , contenant $A$ . Nous allons voir qu'il existe $g \in A$ tel que

\marginpar{($\S4, n^\circ2$, cor.3 de la prop.9)}

\newpage

%% ===== Page 162 =====

- 126 -

pour tout $x \in D(g)$ (c'est-à-dire tel que $g \notin \mathfrak{j}_x$), l'idéal $\mathfrak{j}_x$ soit l'intersection de $A$ et d'un idéal premier de $B$, ce qui prouvera que $D(g) \subset f(X)$. On est donc ramené à démontrer le

Lomme 6. - Soient $A$ un anneau intègre, $B$ une $A$-algèbre intègre de type fini contenant $A$. Il existe $g \in A$ tel que tout homomorphisme de $A$ dans un corps algébriquement clos $\Omega$, non nul en $g$, se prolonge en un homomorphisme de $B$ dans $\Omega$.

Or, il s'agit là d'un résultat classique d'algèbre commutative (Sém. Cartan-Chevalley 1955, 3-01).

Corollaire. - Sous les hypothèses du th. 3, les fonctions $y \to \dim(f^{-1}(y))$ sur $Y$ et $x \to \dim(f^{-1}(f(x)))$ sur $X$ sont constructibles.

En vertu du th. 2 du n° 3, pour tout $n$, l'ensemble des $x$ tels que $\dim(f^{-1}(f(x))) = n$ est différence de deux ensembles fermés, donc constructible, et l'assertion relative à $x \to \dim(f^{-1}(f(x)))$ résulte donc du cor. de la prop. 14. Comme $f(E_n)$ est l'ensemble des $y$ tels que $\dim(f^{-1}(y)) = n$, il est constructible d'après le th. 3 et on conclut de même que la fonction $y \to \dim(f^{-1}(y))$ est constructible.

5. Applications ouvertes et équidimensionnalité.

Proposition 15. - Soient $Y$ un préschéma localement noethérien, $f: X \to Y$ un morphisme de type fini, $x$ un point de $X$, $y = f(x)$. Les conditions suivantes sont équivalentes :

(i) L'image par $f$ de tout voisinage de $x$ est un voisinage de $y$.

(ii) Pour tout $y' \in Y$ dont $y$ est une spécialisation, il existe $x' \in X$ tel que $f(x') = y'$ et que $x$ soit une spécialisation de $x'$.

(iii) Pour toute partie fermée irréductible $Y'$ de $Y$ contenant $y$, il existe une partie fermée irréductible $X'$ de $X$ contenant $x$ et telle que $f(X')$ contienne $Y'$ et soit dense dans $Y'$.

\newpage

%% ===== Page 163 =====

- 127 -

Pour prouver que (i) entraîne (iii), soit $Y' = \{y'\}$, et soit $F$ la réunion des composantes irréductibles de $f^{-1}(Y')$ qui ne contiennent pas $x$. $X - F = U$ est un voisinage de $x$ dans $X$ ; en effet, si $V$ est un voisinage ouvert affine de $y$, d'anneau noethérien, l'intersection de $f^{-1}(V)$ et d'une composante irréductible de $f^{-1}(Y')$ est une composante irréductible de $f^{-1}(V)$, donc ces composantes sont en nombre fini, et notre assertion résulte de ce que $f^{-1}(V)$ est un voisinage de $x$ dans $X$. Cela étant, par hypothèse $f(U)$ est un voisinage de $y$ dans $Y$, donc il existe $x' \in U$ tel que $f(x') = y'$ ; par définition $x'$ appartient à une des composantes irréductibles de $f^{-1}(Y')$ contenant $x$, et comme $y' \in f(X')$, $f(X')$ est dense dans $Y'$. D'ailleurs, si $X'$ est irréductible et fermé, contient $x$ et est tel que $f(X') \subset Y'$ est dense dans $Y'$, et si $x''$ est le point générique de $X'$, on a nécessairement $f(x'') = y'$, donc (iii) entraîne (ii). Enfin, (ii) entraîne (i) : en effet, si $U$ est un ensemble ouvert dans $X$ contenant $x$, $U$ contient tout $x'$ dont $x$ est une spécialisation, donc $f(U)$ contient par hypothèse tout $y'$ dont $y$ est une spécialisation. Comme $f(U)$ est constructible (nº 4, th. 3), il résulte de la prop. 13 du nº 4 que $f(U)$ est un voisinage de $x$.

\subsection*{Sur $X, Y, f$}
Corollaire 1. — Les hypothèses étant celles de la prop. 15, les conditions suivantes sont équivalentes :
(i) $f$ est une application ouverte.
(ii) Pour tout $x \in X$ et tout $y' \in Y$ tel que $f(x) \in \{y'\}$, il existe $x' \in X$ tel que $y' = f(x')$ et $x \in \{x'\}$.
(iii) Pour toute partie fermée irréductible $Y'$ de $Y$, il existe une composante irréductible $X'$ de $f^{-1}(Y')$ telle que $f(X')$ soit dense dans $Y'$.

\newpage

%% ===== Page 164 =====

\setcounter{page}{127}

pour prouver que (i) entraîne (iii), soit $Y'=\{y'\}$, et soit $F$ la réunion des composantes irréductibles de $f^{-1}(Y')$ qui ne contiennent pas $x$. $X-F=U$ est un voisinage de $x$ dans $X$ ; en effet, si $V$ est un voisinage ouvert affine de $y$, d'anneau noethérien, l'intersection de $f^{-1}(V)$ et d'une composante irréductible de $f^{-1}(y')$ est une composante irréductible de $f^{-1}(V \cap Y')$, donc ces composantes sont en nombre fini, et notre assertion résulte de ce que $f^{-1}(V)$ est un voisinage de $x$ dans $X$. Cela étant, par hypothèse $f(U)$ est un voisinage de $y$ dans $Y$, donc il existe $x' \in U$ tel que $f(x')=y'$ ; par définition $x'$ appartient à une des composantes irréductibles de $f^{-1}(Y')$ contenant $x$, et comme $y' \in f(X')$, $f(X')$ est dense dans $Y'$. D'ailleurs, si $X'$ est irréductible et fermé, contient $x$ et est tel que $f(X') \subset Y'$ est dense dans $Y'$, et si $x''$ est le point générique de $X'$, on a nécessairement $f(x'')=y'$, donc (iii) entraîne (ii). Enfin, (ii) entraîne (i) : en effet, si $U$ est un ensemble ouvert dans $X$ contenant $x$, $U$ contient tout $x'$ dont $x$ est une spécialisation, donc $f(U)$ contient par hypothèse tout $y'$ dont $y$ est une spécialisation. Comme $f(U)$ est constructible (n$^o$ 4, th. 3), il résulte de la prop. 13 du n$^o$ 4 que $f(U)$ est un voisinage de $x$.

\section*{Sur $X, Y, f$}

Corollaire 1. — Les hypothèses étant celles de la prop. 15, les conditions suivantes sont équivalentes :

(i) $f$ est une application ouverte.

(ii) pour tout $x \in X$ et tout $y' \in Y$ tel que $f(x) \in \overline{\{y'\}}$, il existe $x' \in X$ tel que $y'=f(x')$ et $x \in \overline{\{x'\}}$.

(iii) pour toute partie fermée irréductible $Y'$ de $Y$, il existe une composante irréductible $X'$ de $f^{-1}(Y')$ telle que $f(X')$ soit dense dans $Y'$.

\newpage

%% ===== Page 165 =====

\begin{tikzcd}
X' \arrow[r] \arrow[d, "f^{Y'}"] & X'' \arrow[d, "f^{Y''}"] \\
Y' \arrow[r, "g"] & Y''
\end{tikzcd}

est commutatif, les flèches horizontales étant les injections canoniques. Cela étant, tout ouvert $U'$ de $X'$ est de la forme $U' = U \cap X'$, où $U$ est un ouvert dans $X''$; comme $(f^{Y''})^{-1}(Y') = X'$ (§ 1, $n^\circ 2$ formule (5)), on a $f^{Y'}(U' \cap X') = f^{Y''}(U) \cap Y'$, et comme $f^{Y''}(U)$ est ouvert dans $Y''$ par hypothèse, on voit que $f^{Y'}(U')$ est ouvert dans $Y'$.

On verra au chap. IV que si $Y'$ est localement noethérien, pour tout morphisme $g: Y' \to Y$ (non nécessairement de type fini), $f^{Y'}$ est encore ouvert.

Corollaire 3. — Les hypothèses sur $X, Y, f$ étant celles de la prop. 15, les conditions suivantes sont équivalentes :

(i) $f$ est universellement ouvert au-dessus de $y \in Y$.

(ii) Pour tout morphisme de type fini $g: Y' \to Y$, tout $y' \in Y'$ tel que $g(y') = y$, et tout $x' \in (f^{Y'})^{-1}(y')$, alors, pour tout $y'_1 \in Y'$ dont $y'$ est spécialisation, il existe $x'_1$ dont $x'$ est spécialisation et tel que $f^{Y'}(x'_1) = y'_1$.

(iii) Pour tout préschéma irréductible $Y'$, tout morphisme de type fini $g: Y' \to Y$ et tout $y' \in Y'$ tel que $g(y') = y$, tout point de $(f^{Y'})^{-1}(y')$ appartient à une composante irréductible de $X^{Y'}$ dont l'image par $f^{Y'}$ est dense dans $Y'$.

Cela résulte aussitôt de la prop. 15.

Corollaire 4. — Soit $X$ un préschéma localement noethérien irréductible, $\eta$ son point générique, $f: X \to Y$ un morphisme de type fini universellement ouvert (resp. universellement ouvert au-dessus de $y \in Y$)

\newpage

%% ===== Page 166 =====

- 130 -

Soient $Y'$ un préschéma irréductible de point générique $\eta'$, et
$g : Y' \to Y$ un morphisme dominant de type fini. Alors les points génériques
des composantes irréductibles de $X^{Y'}$ sont les points génériques
des composantes irréductibles de $(f^{Y'})^{-1}(\eta') = f^{-1}(\eta) \times_Y Y'$
(resp. pour tout $y' \in Y'$ tel que $g(y') = y$, tout point de $(f^{Y'})^{-1}(y')$
appartient à une composante irréductible de $X^{Y'}$ dont l'image par $f^{Y'}$
est dense dans $Y'$). En particulier, si $g(\eta') = f(\eta)$ et si $X$ est irré-
ductible, alors $X^{Y'}$ est irréductible (resp. $X^{Y'}$ a une seule composante
irréductible dont l'image par $f^{Y'}$ est dense dans $Y'$, et
$(f^{Y'})^{-1}(y')$ est contenu dans cette composante).

Observons que si $f$ (resp. $f^{Y'}$) est ouvert et $Y$ (resp. $Y'$) irréductible,
la restriction de $f$ (resp. $f^{Y'}$) à toute composante irré-
ductible de $X$ (resp. $X^{Y'}$) est un morphisme dominant. Lorsque $f$
est universellement ouverte, les assertions du corollaire résultent
du

Lemme 7. - Soient $X, Y$ deux préschémas noethériens, $f : X \to Y$ un mor-
phisme. Supposons $Y$ irréductible de point générique $\eta$ ; alors les
points génériques des composantes irréductibles de $X$ dont l'image
par $f$ est dense dans $Y$ sont les points génériques des composantes
irréductibles de $f^{-1}(\eta)$.

En effet, les composantes irréductibles de $X$ qui rencontrent $f^{-1}(\eta)$
sont celles pour lesquelles $f(Z)$ est dense dans $Y$ ; si $\zeta$ est le point
générique de $Z$, et que $f(\zeta) = \eta$, il est clair que $Z \cap f^{-1}(\eta)$ est l'unique
composante irréductible de $f^{-1}(\eta)$ contenant $\zeta$, puisque c'est l'adhérence de $\{\zeta\}$
dans $f^{-1}(\eta)$ et que toute composante irréductible de $f^{-1}(\eta)$ est conte-
nue dans une composante irréductible de $X$ ; d'où le lemme.

Pour appliquer le lemme au cas où $f$ est universellement ouverte, il
suffit de remarquer qu'on peut remplacer $Y$ par un ouvert affine

\newpage

%% ===== Page 167 =====

- 131 -

donc supposer $Y$ (et par suite $Y'$ et $X$) noethérien ; cela entraîne la première assertion . Pour la seconde, le lemme montre que $f^{-1}(\eta)$ est irréductible, donc aussi $(f^{Y'})^{-1}(\eta') = f^{-1}(\eta)$, et une nouvelle application du lemme prouve que $X^{Y'}$ est irréductible .

Quant aux assertions relatives au cas où $f$ est universellement ouverte au-dessus de $y$, la première a été énoncée dans le cor.3 ; la seconde résulte comme ci-dessus du lemme 7 .

\medskip
\noindent
Proposition 16 . - Soient $Y$ un préschéma localement noethérien irréductible de point générique $\eta$, $X$ un préschéma irréductible, $f : X \to Y$ un morphisme dominant de type fini . Soient $x \in X$, $y=f(x)$ ; supposons que pour toute partie ouverte $U$ de $X$ qui rencontre une des composantes irréductibles de $f^{-1}(y)$ contenant $x$, $f(U)$ soit un voisinage de $y$. Alors les composantes irréductibles de $f^{-1}(y)$ contenant $x$ ont une dimension égale à $\dim(f^{-1}(\eta))$.

On peut évidemment se borner au cas où $x$ est point générique d'une des composantes irréductibles de $f^{-1}(y)$, et l'hypothèse entraîne alors que l'image par $f$ de tout voisinage de $x$ est un voisinage de $y$. En vertu de la prop.15, pour toute chaîne $Y_0 = \{y\} \subset Y_1 \subset \dots \subset Y_r$ de parties fermées irréductibles de $Y$, il existe une chaîne $X_0 = \{x\} \subset X_1 \subset \dots \subset X_r$ de parties fermées irréductibles de $X$ telle que $f(X_i)$ soit contenue dans $Y_i$ et dense dans $Y_i$ : il suffit en effet d'appliquer la prop.15 par récurrence descendante sur $i$. On en conclut par définition (n°2, prop.6) que $\dim(O_y) \le \dim(O_x)$, et par suite (n°2, prop.8) $\dim(O_x) = \dim(O_y)$ ; le cor.1 du th.1 donne alors $\deg.\operatorname{tr.}_{K(y)} K(x) \le e$, et on conclut à l'aide de la prop.10 (ii) du n°3 .

\medskip
\noindent
Corollaire . - Si $f$ est ouvert, pour tout $y \in f(X)$, toutes les composantes de $f^{-1}(y)$ sont de dimension $e$.

Au chap. III nous verrons inversement que lorsque $Y$ est normale et $f$ est universellement ouvert au-dessus de $y$.

\marginpar{la conclusion de la prop. 16 implique que}

\newpage

%% ===== Page 168 =====

\marginpar{Archives Grothendieck-sept 59}

71, before prop. 10 , add :

Corollaire 3. -- Si $f : X \to Y$ est de type fini, il en est de même de $f_{red}$.

On a en effet le diagramme commutatif
$$
\begin{tikzcd}
X_{red} \arrow[r, "f_{red}"] \arrow[d, "j"] & Y_{red} \arrow[d, "j'"] \\
X \arrow[r, "f"] & Y
\end{tikzcd}
$$
d'après la prop. 8, $f \circ j$ est de type fini, donc il en est de même de $j' \circ f_{red}$ ; mais $j'$ étant une immersion fermée, est de type fini (prop. 7), et la conclusion résulte du cor. 2.

after the proof of prop. 9
p. 52 insert a new corollary :

Corollaire XX. -- Soient $X$ un préschéma réduit, $f : X \to Y$ un morphisme, $Z$ le sous-préschéma fermé réduit de $Y$ ayant pour espace de base l'adhérence $\overline{f(X)}$ ; alors $f$ se factorise en $X \xrightarrow{g} Z \xrightarrow{i} Y$, où $i$ est le morphisme d'injection.

On se ramène aussitôt au cas où $Y = \operatorname{Spec}(B)$ est affine ; soit $\varphi$ l'homomorphisme $\Gamma(Y, \mathcal{O}_Y) \to \Gamma(X, \mathcal{O}_X)$ correspondant à $f$ (§ 2, n$^\circ$2, prop. 3). Désignons par $\mathfrak{a}$ le noyau de $\varphi$, et soit $Z$ le sous-préschéma fermé de $Y$ défini par le faisceau quasi-cohérent $\tilde{\mathfrak{a}}$ ; comme $\varphi$ se factorise en $B \to B/\mathfrak{a} \to \Gamma(X, \mathcal{O}_X)$, on est ramené au cas où $\mathfrak{a} = 0$, $Z = Y$, et tout revient à démontrer dans ce cas que $Y$ est réduit et $f(X)$ partout dense dans $Y$. Soit $(V_\alpha)$ un recouvrement de $X$ par des ouverts affines, et soit $\varphi_\alpha$ l'homomorphisme $\Gamma(Y, \mathcal{O}_Y) \to \Gamma(V_\alpha, \mathcal{O}_{V_\alpha})$ composé de $\varphi$ et de l'homomorphisme de restriction ; si $b$ est le noyau de $\varphi_\alpha$, on a $\bigcap_\alpha \mathfrak{b}_\alpha = (0)$. Notons d'autre part que comme $X$ est réduit, il n'y a pas d'élément nilpotent dans $\Gamma(V_\alpha, \mathcal{O}_{V_\alpha})$ ; donc l'image par $\varphi_\alpha$ d'un élément nilpotent de $B$ est nécessairement 0, donc $\mathfrak{a} \subset \mathfrak{b}_\alpha$ pour tout $\alpha$, autrement dit $\mathfrak{b} = 0$. D'autre part, dire que $y \in f(X)$ signifie que $j_y$ est de la forme $\varphi_\alpha^{-1}(j_{x})$ pour un $\alpha$ et un $x \in V_\alpha$ ; comme l'interse

\newpage

%% ===== Page 169 =====

tion des $j_x$ pour $x \in V_\alpha$ est réduite à $0$, l'intersection des $\varphi_\alpha^{-1}(j_x)$ pour $x \in V_\alpha$ est $\mathfrak{a}$ ; on en conclut que $\overline{f(X)} \cap \mathfrak{a} = (0)$, d'où $Y = \overline{f(X)}$.

P. $91$, after the definition of $\operatorname{Spec}(A)$, add :

Pour que $\operatorname{Spec}(A)$ soit vide, il faut et il suffit que l'anneau $A$ soit réduit à $0$.

P. $27$, replace the last $5$ lines of the proof of prop. $4$ by :

ma affine ; dans ce cas, la proposition résulte du :

Lemme $2$ bis. — Soient $A$ un anneau, $S$ une partie multiplicative de $A$, $\varphi$ l'homomorphisme canonique $A \to S^{-1}A$ ; alors $^a\varphi$ est un homomorphisme de $\operatorname{Spec}(S^{-1}A)$ sur le sous-espace de $\operatorname{Spec}(A)$ formé des $x$ tels que $j_x \cap S = \emptyset$ ; et pour tout $y \in \operatorname{Spec}(S^{-1}A)$, si on pose $x = ^a\varphi(y)$, $\varphi_y$ est un isomorphisme de $\mathcal{O}_x$ sur $\mathcal{O}_y$.

La première partie a déjà été démontrée ($\S$ $1$, n$^{\text{o}} 2$, cor. $3$ de la prop. $5$)

La seconde est un résultat classique d'algèbre commutative (chap. $0, \S$ $1$, n$^{\text{o}}$ ).

P. $41$ bis (end of $\S$ $2$, n$^{\text{o}} 7$), replace Prop. $19$ by :

Proposition $19$. — Soient $f : X \to Y$ un morphisme, $y$ un point de $Y$, $Z$ le préschéma $\operatorname{Spec}(\mathcal{O}_y)$, $p = (y, \psi)$ la projection $X \times_Y \operatorname{Spec}(\mathcal{O}_y) \to X$ ; $\psi$ est un homomorphisme de l'espace de base de $X \times_Y Z$ sur le sous-espace $f^{-1}(y)$, et pour tout $t \in X \times_Y Z$, si on pose $z = \psi(t)$, $\psi_t$ est un isomorphisme de $\mathcal{O}_z$ sur $\mathcal{O}_t$.

Comme l'image canonique de $Z$ dans $Y$ (identifiée à $Z$) est contenue dans tout ouvert affine contenant $y$ (n$^{\text{o}} 2$), on se ramène encore, comme dans la prop. $17$, au cas où $X = \operatorname{Spec}(A)$ et $Y = \operatorname{Spec}(B)$ sont affines, $A$ étant une algèbre sur $B$. Alors $X \times_Y Z$ est le spectre premier de

\newpage

%% ===== Page 170 =====

- xxii -

P.58, cor. de prop.18 devient cor.1 ; ajouter après :

Corollaire 2. — Si $X, Y$ sont deux $S$-préschémas tels que $Y$ soit séparé au-dessus de $S$, la projection $X \times_S Y \to X$ est un morphisme séparé.

Il suffit, dans la prop.18, de remplacer $Y$ par $S$ et $S'$ par $Y$ (§ 2, n°5, cor. de la prop.8).

P.62, après cor.3 de prop.22, ajouter :

Corollaire 4. — Soient $X_1, X_2$ deux $S$-préschémas séparés au-dessus de $S$ ; soient $f: Y \to X_1$, $g: Y \to X_2$ deux $S$-morphismes. Si $f$ ou $g$ est une immersion fermée, $(f,g)_S$ est une immersion fermée.

En effet, on sait (n°6, cor.2 de la prop.18) que les projections $p_1, p_2$ sont des morphismes séparés ; le corollaire résulte donc du cor.1.

P.65, après déf.1, ajouter :

Remarque. — Si $X$ est localement noethérien, le faisceau structural $\mathcal{O}_X$ est un faisceau cohérent d'anneaux (FAC, I, 2, 15) : en effet, il suffit de montrer que le noyau d'un homomorphisme $\mathcal{O}_X^n \to \mathcal{O}_X$ est un $\mathcal{O}_X$-module de type fini ; la question étant locale, on est ramené au cas où $X = \operatorname{Spec}(A)$, $A$ étant un anneau noethérien, et tout revient à remarquer que le noyau d'un homomorphisme $A^n \to A$ est de type fini (§ 1, n°3, cor.4 un th.1). On en déduit en particulier que tout [marginpar: quasi-cohérent sous-faisceau un faisceau cohérent sur $X$ est cohérent, et que si $F$ est un $\mathcal{O}_X$-Module pour lequel tout $x \in X$ a un voisinage $U$ tel que $F|_U$ soit isomorphe au conoyau d'un homomorphisme $\mathcal{O}_U^p \to \mathcal{O}_U^q$, $F$ est cohérent (FAC, loc. cit.).]

\newpage

%% ===== Page 171 =====

- xxiii -

P. 67, replace def. 2 by :

\emph{Définition 2. — On dit qu'un préschéma est artinien s'il est affine et si son anneau est artinien.}

P. 74, replace the line before cor. 1 by :

$X$ est un schéma fini sur $K$, de rang $[A:K]$, que l'on note $rg_K(X)$ ; pour deux préschémas finis sur $K$, $X, Y$, on a 
$$rg_K(X+Y) = rg_K(X) + rg_K(Y), \quad rg_K(X \times_K Y) = rg_K(X) rg_K(Y).$$

P. 78, after cor. to prop. 16, add :

\emph{Remarque. — Si $\theta^b$ est surjectif, il résulte de la démonstration de la prop. 16 (i) que l'on peut toujours supposer que l'ensemble ouvert $U$ est de la forme $\phi^{-1}(V)$, où $V$ est ouvert dans $Y$.}

P. 1, after the introduction of $f(x)$, add :

Plus généralement, lorsque $(X, \mathcal{O}_X)$ est un espace annelé tel que tout anneau $\mathcal{O}_x$ soit un anneau local ; pour toute section $f \in \Gamma(X, \mathcal{O}_X)$, nous désignerons par $f(x)$ la classe du germe $f_x$ de $f$ au point $x \in X$, modulo l'idéal maximal de $\mathcal{O}_x$.

(This of course will ultimately go into chap. 0 when that chapter is rewritten).

P. 54, at the end of n° 4, add :

Nous dirons qu'un préschéma $X$ est intègre s'il est irréductible et réduit. Il revient au même de dire que l'espace de base $X$ est irréductible et que le faisceau structural $\mathcal{O}_X$ est intègre (i.e. que chaque fibre $\mathcal{O}_x$ est intègre), comme il résulte de la prop. 4 du § 1, n° 1.

\newpage

%% ===== Page 172 =====

- xxiii -

P.67, replace def.2 by :

\emph{Définition 2} .— On dit qu'un préschéma est artinien s'il est affine et si son anneau est artinien .

P.74, replace the line before cor.1 by :

$X$ est un schéma fini sur $K$, de rang $[A:K]$, que l'on note $r_K(X)$ ; pour deux préschémas finis sur $K$, $X, Y$, on a 
$$r_K(X+Y)=r_K(X)+r_K(Y) \quad , \quad r_K(X \times Y)=r_K(X)r_K(Y) \quad .$$

P.78, after cor. to prop.16, add :

\emph{Remarque} .— Si $f^b$ est surjectif, il résulte de la démonstration de la prop.16 (i) que l'on peut toujours supposer que l'ensemble ouvert $U$ est de la forme $\varphi^{-1}(V)$, où $V$ est ouvert dans $Y$.

P.1, after the introduction of $f(x)$, add :

Plus généralement, lorsque $(X,\mathcal{O}_X)$ est un espace annelé tel que tout anneau $\mathcal{O}_x$ soit un anneau local, pour toute section $f \in \Gamma(X, \mathcal{O}_X)$, nous désignerons par $f(x)$ la classe du germe $f_x$ de $f$ au point $x \in X$, modulo l'idéal maximal de $\mathcal{O}_x$.

(This of course will ultimately go into chap.0 when that chapter is rewritten).

P. 54, at the end of n°4, add :

Nous dirons qu'un préschéma $X$ est intègre s'il est irréductible et réduit. Il revient au même de dire que l'espace de base $X$ est irréductible et que le faisceau structural $\mathcal{O}_X$ est intègre (i.e. que chaque fibre $\mathcal{O}_x$ est intègre), comme il résulte de la prop.4 du § 1, n°1.

\newpage

%% ===== Page 173 =====

- II-ii -

1. II-22 , after the definition of $A(X)$ , add :
Si $U$ est une partie ouverte de $S$, on a donc $A(f^{-1}(U)) = A(X)|U$ .

P. II-29 , after cor.3 , add :
Corollaire 4. -- Soit $X$ un préschéma affine au-dessus de $S$ ; pour que $X$ soit de type fini sur $S$, il faut et il suffit que la $\mathcal{O}_S$-algèbre $A(X)$ soit de type fini.

On est aussitôt ramené au cas où $S$ est affine ; alors $X$ est un schéma affine (cor.2) et si $S = \operatorname{Spec}(A)$, $X = \operatorname{Spec}(B)$, $A(X)$ est le faisceau associé à la $A$-algèbre $B$ ; [deleted: xxxxxxxx] le corollaire résulte donc des définitions et de I,4,2, prop.6 .

P. II-33 , at the end of prop.12 , add :
(vi) Si $f : X \to Y$ est affine, il en est de même de $f_{\operatorname{red}}$.

P. II-33 , at the end of the proof of prop.12 , add :
Enfin, pour démontrer (vi), on remarque que le diagramme
$$\begin{tikzcd}
X_{\operatorname{red}} \arrow[r, "f_{\operatorname{red}}"] \arrow[d] & Y_{\operatorname{red}} \arrow[d] \\
X \arrow[r, "f"] & Y
\end{tikzcd}$$
est commutatif, les flèches verticales étant des immersions fermées, donc des morphismes séparés ; la conclusion résulte alors de (ii) et de (v).

P. II-36 , at the end of prop.15 , add :
(vi) Si $f : X \to Y$ est un morphisme fini (resp. entier), il en est de même de $f_{\operatorname{red}}$.

P. II-37 , at the end of the proof of prop.15 , add :
Enfin, (vi) se déduit de (ii) et (v) par le même raisonnement formel que dans la prop.12.

P. II-38 , before n°8 , add :
Proposition 17. -- Soit $f : X \to Y$ un morphisme fini. Alors, pour tout $y \in Y$, la fibre $f^{-1}(y)$ est un espace discret fini.

\newpage

%% ===== Page 174 =====

- II-iii -

En effet, en tant que $\kappa(y)$-préschéma, $f^{-1}(y)$ est identifié à $X \times_{\operatorname{Spec}(\kappa(y))} \operatorname{Spec}(\kappa(y))$ (I, 2, 7, prop. 17), qui est fini au-dessus de $\operatorname{Spec}(\kappa(y))$ (prop. 15 (iii)) ; c'est donc un schéma affine dont l'anneau est une algèbre de dimension finie sur $\kappa(y)$ (prop. 14). Le corollaire résulte par suite de I, 4, 2, prop. 12.

P. II-3, in the statement of prop. 2, suppress the assumption on $Y$, and replace the proof by :

Si $\mathcal{G}$ est quasi-cohérent (resp. cohérent), pour tout $x \in X$, il y a un voisinage ouvert $V$ de $f(x)$ tel que $\mathcal{G}|V$ soit isomorphe au conoyau d'un homomorphisme $\mathcal{O}_Y^m|V \to \mathcal{O}_Y^n|V$ (resp. $\mathcal{O}_Y^m|V \to \mathcal{O}_Y^n|V$). Comme $f^*(\mathcal{O}_Y) = \mathcal{O}_X$ et que $f^*$ est exact à droite et commute aux sommes directes (chap 0, $\S$ 2, $n^o$ ), on voit que si $U=f^{-1}(V)$, $f^*(\mathcal{G})|U$ est isomorphe au conoyau d'un homomorphisme $\mathcal{O}_X^m|U \to \mathcal{O}_X^n|U$ (resp. $\mathcal{O}_X^m|U \to \mathcal{O}_X^n|U$) ; d'où la proposition, compte tenu de ce que $X$ est localement noethérien, lorsque $\mathcal{G}$ est supposé cohérent.

P. II-4, in the statement about the tensor product $\mathcal{F}_1 \otimes \dots \otimes \mathcal{F}_n$, delete the assumption that the $X_i$ are locally noetherian.

P. II-8, add a footnote concerning the "set" $\mathfrak{m}$ :

(*) Pour que $\mathfrak{m}$ soit un ensemble, il faut modifier légèrement la définition donnée ci-dessus. Convenons d'appeler \emph{hauteur} d'un faisceau $\mathcal{H}$ sur un espace $V$ le plus petit cardinal supérieur aux cardinaux des fibres de $\mathcal{H}$. Étant donné un préschéma $X$, et un cardinal $\mathfrak{m}$, il y a un ensemble $\mathfrak{N}$ de faisceaux quasi-cohérents définis sur des ouverts de $X$ tel que, pour tout couple $(V, \mathcal{H})$ formé d'un ouvert $V \subset X$ et d'un faisceau quasi-cohérent $\mathcal{H}$ sur $V$, de hauteur $\le \mathfrak{m}$, il y ait un faisceau $\mathcal{H}' \in \mathfrak{N}$ tel que $\mathcal{H}'$ et $\mathcal{H}$ soient isomorphes (cela est dû à ce que pour un anneau donné $A$, tout $A$-module de cardinal $\le \mathfrak{m}$ est isomorphe à un quotient de $A^{(\mathfrak{m})}$). D'autre part si $X$ est affine noethérien et si $\mathcal{F}$ est donné sur $U$, de hauteur $\mathfrak{n}$, la méthode de I, 1, 4, th. 2 donne un faisceau quasi-cohérent sur $X$, induisant $\mathcal{F}$ et dont la hauteur est $\le \mathfrak{n}^{\operatorname{Card}(U)}$. Dans le cas général, on prend pour $\mathfrak{m}$ un cardinal infini supérieur à $\mathfrak{n}^{\operatorname{Card}(X)}$; on considère

\newpage

%% ===== Page 175 =====

- II-iv -

re l'ensemble $\mathbb{N}$ correspondant, que l'on peut supposer être tel que deux faisceaux distincts de $\mathbb{N}$ ne sont jamais isomorphes. On ordonne alors $\mathbb{N}$ par la relation "il existe une immersion ouverte de $H$ dans $H'$", et le reste de la démonstration est inchangée (on peut évidemment supposer que $F \in \mathbb{N}$).

From p. II-11 on, the word "divisoriel" should everywhere be replaced by "inversible". P. II-12, in the remark after def. 1, add: si $\mathcal{O}_X$ est un faisceau d'anneaux cohérent.

P. II-12, replace the proof of prop. 11 by:

Plus généralement, si $L$ est localement libre de rang fini, on a un isomorphisme canonique fonctoriel $\operatorname{Hom}_{\mathcal{O}_X}(L, \mathcal{O}_X) \otimes_{\mathcal{O}_X} L' \to \operatorname{Hom}_{\mathcal{O}_X}(L, L')$ : pour tout ouvert $U$ tel que $L|_U$ soit isomorphe à $\mathcal{O}_U^n$, à tout couple $(u, t)$ où $u \in \Gamma(U, \operatorname{Hom}_{\mathcal{O}_X}(L, \mathcal{O}_X))$ et $t \in \Gamma(U, L')$, on fait correspondre l'élément $s \mapsto u(s)t$ de $\Gamma(U, \operatorname{Hom}_{\mathcal{O}_X}(L, L'))$. Pour voir qu'on a un isomorphisme, on peut supposer $U=X$, et comme $\operatorname{Hom}_{\mathcal{O}_X}(\mathcal{O}_X^n, \mathcal{O}_X)$ est canoniquement isomorphe à $\mathcal{O}_X^n$ et $\operatorname{Hom}_{\mathcal{O}_X}(\mathcal{O}_X^n, L')$ à $(L')^n$, on est ramené au cas $L=\mathcal{O}_X$, qui est immédiat.

P. II-15, line 7, add:
l'associativité de cette multiplication se vérifie aussitôt.

P. II-15 and 16, replace subsection D) by:

Proposition 13. - Soient $(X, \mathcal{O}_X)$ un espace annelé tel que les $\mathcal{O}_{X,x} = \mathcal{O}_x$ soient des anneaux locaux, dont on désigne par $\mathfrak{m}_x$ l'idéal maximal. Soit $L$ un $\mathcal{O}_X$-module inversible, et soit $f \in \Gamma(X, L)$. Pour tout $x \in X$, les conditions suivantes sont équivalentes : a) $f_x$ est un générateur de $\mathcal{O}_x$ ; b) $f_x \notin \mathfrak{m}_x L_x$ ; c) il existe une section $g$ de $L^{-1}$ au-dessus d'un voisinage $U$ de $x$ telle que l'image canonique de $f \otimes g$ dans $\Gamma(U, \mathcal{O}_X)$ (prop. 10) soit l'unité.

On peut évidemment se borner au cas où $L=\mathcal{O}_X$ ; l'équivalence de a) et b) est évidente, et il est clair que c) entraîne b) ; d'autre

\newpage

%% ===== Page 176 =====

- II-v -

part, si $f_x \notin \mathfrak{m}_x$, $f_x$ est inversible dans $\mathcal{O}_x$, soit $f_x g_x = 1_x$. Par définition des germes de sections, cela veut dire qu'il existe un voisinage $U$ de $x$ et une section $g$ de $\mathcal{O}_X$ au-dessus de $U$ telle que $fg = 1$ dans $U$, d'où c).

Corollaire 1. - L'ensemble $X_f$ des $x$ vérifiant les conditions équivalentes a), b), c) est ouvert dans $X$.

Cela est évident pour la condition c).

Corollaire 2. - Soient $L, L'$ deux faisceaux inversibles sur $X$ (les $\mathcal{O}_X$ étant supposés être des anneaux locaux) ; si $f \in \Gamma(X, L)$, $g \in \Gamma(X, L')$ on a
$$(1) \qquad X_f \cap X_g = X_{f \otimes g}$$

Par définition, on peut en effet se ramener au cas où $L=L'=\mathcal{O}_X$ ; alors $f \otimes g$ s'identifie au produit $fg$, et le corollaire est évident.

Par abus de langage, on dit parfois (lorsqu'aucune confusion n'en résulte) que $X_f$ est l'ensemble des $x \in X$ où $f$ ne s'annule pas.

Remarque. - Supposons que $X$ soit un schéma local (I,2,2) ; si $a$ est l'unique point fermé de $X$, tout ouvert affine contenant $a$ est nécessairement $X$ entier ; il s'ensuit que tout faisceau inversible sur $X$ est nécessairement isomorphe à $\mathcal{O}_X$ (ou, comme on dit encore, est trivial). Cette propriété n'a pas lieu en général pour un schéma affine quelconque $\operatorname{Spec}(A)$ ; on verra au chap. IV que si $A$ est un anneau normal, elle équivaut au fait que $A$ est factoriel.

p. II-18, in prop. 14 (former prop. 15), instead of $J(x)=0$, read :
telle que le support de $F$ soit contenu dans le support de $\mathcal{O}_X / J$.

p. II-18, line -4, instead of "contenu dans $Y$" read "égal à $Y$".

\newpage

%% ===== Page 177 =====

- II-VI -

P. II-20, replace lines 3-14 by :
montrons que pour tout $z \in Y' \cap Y''$, l'homomorphisme $F(z) \to F'(z) \oplus F''(z)$ est bijectif. En effet, par définition de $J$, on a $J \cap J' = J$ ; si $z \notin Y''$, on a donc $J'(z) = J(z)$ et par suite pour tout germe de section $\tilde{s}(z)$ de $F(z)$, on a $\tilde{s}(z) \otimes 1_z = \tilde{s}(z)$ ; autrement dit $F'(z) = F(z)$, et comme $F''(z) = 0$, notre assertion est établie dans ce cas ; on raisonne de même si $z \notin Y'$. Par suite, le noyau et le conoyau de $F \to F' \oplus F''$ ont leurs supports dans $Y' \cap Y''$, et sont dans $\mathcal{K}'$ par hypothèse ; il en est de même de $F'$ et $F''$ pour la même raison, donc $F' \oplus F''$ est aussi dans $\mathcal{K}'$, $\mathcal{K}'$ étant exacte. La conclusion résulte alors de la considération des deux suites exactes
$$ 0 \to \operatorname{Im} \varphi \to F'' \to \operatorname{Coker} \varphi \to 0 $$
$$ 0 \to \operatorname{Ker} \varphi \to F \to \operatorname{Im} \varphi \to 0 $$
et de l'hypothèse que $\mathcal{K}'$ est exacte.

P. II-21, last line of corollary, instead of "remplacée par $G_y \oplus 0$" read "supprimée".

P. II-52, replace Prop. 7 by :
Proposition 7. - Si l'anneau $S$ est intègre, le schéma $\operatorname{Proj}(S)$ est réduit et son espace de base est irréductible et non vide. Inversement, si $\operatorname{Proj}(S)$ possède ces propriétés, et si en outre $S_+$ ne contient pas d'élément nilpotent $\neq 0$ et $S_0$ ne contient pas de diviseur de $0$ dans $S$, alors $S$ est intègre.
Si $S$ est intègre, $(0)$ est un idéal premier gradué, donc l'irréductibilité de l'espace de base $\operatorname{Proj}(S)$ résulte du cor. 6 de la prop. 5. D'autre part, pour $f$ homogène $\neq 0$ dans $S_+$, $S_f$ et a fortiori $S_{(f)}$ n'a pas d'élément nilpotent et n'est pas réduit à $0$, donc $D_+(f)$ est un schéma réduit non vide ; il en est par suite de même de $X = \operatorname{Proj}(S)$ (1, 3, 4, déf. 5). Réciproquement, si $S_+$ n'a pas d'élément nilpotent $\neq 0$, pour tout $f$ homogène $\neq 0$ dans $S$, $D_+(f) \neq \emptyset$ (cor. 1 de

\newpage

%% ===== Page 178 =====

- II-vii -

la prop.4) ; l'hypothèse que $\operatorname{Proj}(S)$ est irréductible entraîne $D_+(f) \cap D_+(g) \neq \emptyset$ pour $f,g$ homogènes et $\neq 0$ dans $S_+$, donc $fg \neq 0$ d'après (7), ce qui montre que $S_+$ n'a pas de diviseur de $0$ ; comme $S_0$ ne contient pas non plus de diviseur de $0$ dans $S$, on en conclut aussitôt que $S$ est intègre.

P. II-60, après la ligne 8, add :
(puisque $\mathcal{O}_X(n)$ est localement libre)
Il est clair que $\Gamma_*$ est un foncteur covariant additif exact à gauche ; en particulier, si $J$ est un faisceau d'idéaux sur $X$, $\Gamma_*(J)$ s'identifie à un idéal gradué de $\Gamma_*(\mathcal{O}_X)$.

P. II-62, avant $n^\circ 4$, add :

\emph{Remarque}. - Avec les notations antérieures, les homomorphismes composés
$$(20) \quad \tilde{M} \xrightarrow{\alpha} (\Gamma_*(\tilde{M}))^{\sim} \xrightarrow{\beta} \tilde{M}$$
et
$$(21) \quad \Gamma_*(F) \xrightarrow{\alpha} \Gamma_*((\Gamma_*(F))^{\sim}) \xrightarrow{\Gamma_*(\beta)} \Gamma_*(F)$$
sont les isomorphismes identiques. En effet, la vérification de (20) est locale, et dans un ouvert affine $D_+(f)$, elle résulte aussitôt des définitions et de ce que $\beta$, appliqué à des faisceaux quasi-cohérents, est déterminé par son action sur les sections au-dessus de $D_+(f)$ (I,1,3,cor.2 du th.1). La vérification de (21) se fait pour chaque degré séparément : si on pose $M = \Gamma_*(F)$, on a $M_n = \Gamma(X, F(n))$ et $(\Gamma_*(\tilde{M}))_n = \Gamma(X, (\tilde{M}(n)))^{\sim} = \Gamma(X, (M(n))^{\sim})$. Or, si $f \in S_1$ et $z \in M_n$, $(z)$ est l'élément $z/1$ de $(M(n))_{(f)}$, égal à $(f/1)^n (z/f^n)$ ; il lui correspond par $\alpha(f)$ la section $((a_f(f))^n / D_+(f))((z/D_+(f))((a_f(f))^n D_+(f))^{-1})$ au-dessus de $D_+(f)$, c'est-à-dire la restriction de $z$ à $D_+(f)$, ce qui vérifie (21).

\newpage

%% ===== Page 179 =====

- II-vii -
la prop.4) ; l'hypothèse que $\operatorname{Proj}(S)$ est irréductible entraîne
$D_+(f) \cap D_+(g) \neq \emptyset$ pour $f,g$ homogènes et $\neq 0$ dans $S_+$, donc $fg \neq 0$ d'après
(7), ce qui montre que $S_+$ n'a pas de diviseur de $0$ ; comme $S_0$ ne con-
tient pas non plus de diviseur de $0$ dans $S$, on en conclut aussitôt
que $S$ est intègre.

P. II-60, after line 8, add :
(puisque $\mathcal{O}_X(n)$ est localement libre)
Il est clair que $\Gamma_*$ est un foncteur covariant additif exact à gauche ; en particulier, si $\mathcal{J}$ est un faisceau d'idéaux sur $X$, $\Gamma_*(\mathcal{J})$
s'identifie à un idéal gradué de $\Gamma_*(\mathcal{O}_X)$.

P. II-62, before n°4, add :

Remarque .- Avec les notations antérieures, les homomorphismes composés
$$(20) \quad \tilde{M} \xrightarrow{\alpha} (\Gamma_*(\tilde{M}))^{\sim} \xrightarrow{\beta} \tilde{M}$$
et 
$$(21) \quad \Gamma_*(F) \xrightarrow{\alpha} \Gamma_*((\Gamma_*(F))^{\sim}) \xrightarrow{\Gamma_*(\beta)} \Gamma_*(F)$$
sont les isomorphismes identiques. En effet, la vérification de
(20) est locale, et dans un ouvert affine $D_+(f)$, elle résulte aussitôt
des définitions et de ce que $\beta$, appliqué à des faisceaux quasi-cohérents, est déterminé par son action sur les sections au-dessus
de $D_+(f)$ (I,1.3, cor.2 du th.1). La vérification de (21) se fait
pour chaque degré séparément : si on pose $M=\Gamma_*(F)$, on a $M_n = \Gamma(X, F(n))$ et $(\Gamma_*(\tilde{M}))_n = \Gamma(X, (\tilde{M}(n))^{\sim})$. Or, si $f \in S_1$ et
$z \in M_n$, $\beta(z)$ est l'élément $z/1$ de $(M(n))_{(f)}$, égal à $(f/1)^n(z/f^n)$ ;
il lui correspond par $\alpha$ la section
$((\alpha_f(f))^n/D_+(f))((z/D_+(f))(\alpha_f(f))^{-n})$ au-dessus de $D_+(f)$,
c'est-à-dire la restriction de $z$ à $D_+(f)$, ce qui vérifie (21).

\newpage

%% ===== Page 180 =====

part , si $f_x \notin \mathfrak{m}_x$ , $f_x$ est inversible dans $\mathcal{O}_x$ , soit $f_x g_x = 1_x$ . Par
définition des germes de sections , cela veut dire qu'il existe un
voisinage $U$ de $x$ et une section $g$ de $\mathcal{O}_X$ au-dessus de $U$ telle que $fg$
$= 1$ dans $U$ , d'où c) .

Corollaire 1 . — L'ensemble $X_f$ des $x$ vérifiant les conditions équi-
valentes a), b), c) est ouvert dans $X$ .

Cela est évident pour la condition c) .

Corollaire 2 . — Soient $L, L'$ deux faisceaux inversibles sur $X$ (les $\mathcal{O}_X$
étant supposés être des anneaux locaux) ; si $f \in \Gamma(X, L)$, $g \in \Gamma(X, L')$
on a
$$(1) \quad f_x \otimes g_x = f \otimes g$$
Par définition , on peut en effet se ramener au cas où $L=L'=\mathcal{O}_X$ ; a-
lors $f \otimes g$ s'identifie au produit $fg$ , et le corollaire est évident .
Par abus de langage , on dit parfois (lorsqu'aucune confusion n'en
résulte) que $X_f$ est l'ensemble des $x \in X$ où $f$ ne s'annule pas .

Remarque . — Supposons que $X$ soit un schéma local (1,2,2) ; si $x$ est
l'unique point fermé de $X$ , tout ouvert affine contenant $x$ est néces-
sairement $X$ tout entier ; il s'ensuit que tout faisceau inversible
sur $X$ est nécessairement isomorphe à $\mathcal{O}_X$ (ou , comme on dit encore ,
est trivial) . Cette propriété n'a pas lieu en général pour un schéma
affine quelconque $\operatorname{Spec}(A)$ ; on verra au chap. IV que si $A$ est un an-
neau normal , elle équivaut au fait que $A$ est factoriel .

p. II-18 , in prop. 14 (former prop. 15) , au lieu de $J(x)$
$=0$ , read :
telle que le support de $F$ soit contenu dans le support de $\mathcal{O}_X/J$ .

p. II-18 , line -4 , instead of "contenu dans $Y$" read "égal à $Y$" .

\newpage

%% ===== Page 181 =====

\begin{itemize}
\item[-- II-v --]
\end{itemize}

part, si $f_x \in \mathfrak{m}_x$, $f_x$ est inversible dans $\mathcal{O}_x$, soit $f_x g_x = 1_x$. Par définition des germes de sections, cela veut dire qu'il existe un voisinage $U$ de $x$ et une section $g$ de $\mathcal{O}_X$ au-dessus de $U$ telle que $fg = 1$ dans $U$, d'où $c$.

Corollaire 1. -- L'ensemble $X_f$ des $x$ vérifiant les conditions équivalentes a), b), c) est ouvert dans $X$.

Cela est évident pour la condition c).

Corollaire 2. -- Soient $L, L'$ deux faisceaux inversibles sur $X$ (les $\mathcal{O}_x$ étant supposés être des anneaux locaux) ; si $f \in \Gamma(X, L)$, $g \in \Gamma(X, L')$ on a
$$(1) \quad X_f \cap X_g = X_{f \otimes g}$$

Par définition, on peut en effet se ramener au cas où $L=L'=\mathcal{O}_X$ ; alors $f \otimes g$ s'identifie au produit $fg$, et le corollaire est évident.

Par abus de langage, on dit parfois (lorsqu'aucune confusion n'en résulte) que $X_f$ est l'ensemble des $x \in X$ où $f$ ne s'annule pas.

Remarque. -- Supposons que $X$ soit un \emph{schéma local} (I, 2, 2) ; si $x$ est l'unique point fermé de $X$, tout ouvert affine contenant $x$ est nécessairement $X$ tout entier ; il s'ensuit que tout faisceau inversible sur $X$ est nécessairement \emph{isomorphe} à $\mathcal{O}_X$ (ou, comme on dit encore, est trivial). Cette propriété n'a pas lieu en général pour un schéma affine quelconque $\operatorname{Spec}(A)$ ; on verra au chap. IV que si $A$ est un anneau normal, elle équivaut au fait que $A$ est factoriel.

P. II-18, in prop. 14 (former prop. 15), instead of $J(x) = 0$, read :

telle que le support de $F$ soit contenu dans le support de $\mathcal{O}_X/J$.

P. II-18, line -4, instead of "contenu dans $Y$" read "égal à $Y$".

\newpage

%% ===== Page 182 =====

\marginpar{Archives Grothendieck-sept. 59}

\begin{center}
--- xxiv ---
\end{center}

P.112, before n°3, insert :

\emph{Remarque} .— Rappelons enfin la définition de la dimension d'un anneau local noethérien $A$ par son polynôme caractéristique, que nous utiliserons au chap. III : si $\mathfrak{q}$ est un idéal de définition de $A$, la longueur de $A/\mathfrak{q}^n$ est un polynôme en $n$ pour les valeurs assez grandes de $n$, et le degré de ce polynôme est la dimension de $A$ (th. de Hilbert-Samuel) .

P.23-24, after prop. 1, add :

Lorsque $y \in \{x\}$, on dit que $y$ est \underline{spécialisation de $x$}.

\newpage

%% ===== Page 183 =====

- 0-ii -

$\Gamma_U(f) : \Gamma(U, \mathcal{F}_\alpha) \to \Gamma(U, \mathcal{G}_\alpha)$ forment un système projectif d'homomorphismes, les diagrammes
$$\begin{tikzcd}
\Gamma(U, \mathcal{F}_\alpha) \arrow[r] \arrow[d] & \Gamma(U, \mathcal{G}_\alpha) \arrow[d] \\
\Gamma(V, \mathcal{F}_\alpha) \arrow[r] & \Gamma(V, \mathcal{G}_\alpha)
\end{tikzcd}$$
étant en outre commutatifs pour $U \supset V$ ; on en conclut que si $f_U = \varprojlim \Gamma_U(f_\alpha)$, les homomorphismes $f_U$ définissent un homomorphisme de faisceaux $f : \varprojlim \mathcal{F}_\alpha \to \varprojlim \mathcal{G}_\alpha$, que l'on désigne par $\varprojlim f_\alpha$.

On peut donc dire que $\varprojlim$ est un foncteur de la catégorie des systèmes projectifs de faisceaux de groupes abéliens sur $X$ dans la catégorie des faisceaux de groupes abéliens sur $X$. Ce foncteur est \emph{exact à gauche} : autrement dit, si $(\mathcal{F}'_\alpha)$, $(\mathcal{F}_\alpha)$, $(\mathcal{F}''_\alpha)$ sont trois systèmes projectifs de faisceaux de groupes abéliens sur $X$, et si pour chaque $\alpha$, on a une suite exacte avec commutativité des diagrammes
$$\begin{tikzcd}
0 \arrow[r] & \mathcal{F}'_\alpha \arrow[r, "f_\alpha"] \arrow[d] & \mathcal{F}_\alpha \arrow[r, "g_\alpha"] \arrow[d] & \mathcal{F}''_\alpha \arrow[d] \\
0 \arrow[r] & \mathcal{F}'_\beta \arrow[r, "f_\beta"] & \mathcal{F}_\beta \arrow[r, "g_\beta"] & \mathcal{F}''_\beta
\end{tikzcd}$$
pour $\alpha \leq \beta$, alors en posant $f = \varprojlim f_\alpha$, $g = \varprojlim g_\alpha$, la suite
$$0 \to \varprojlim \mathcal{F}'_\alpha \to \varprojlim \mathcal{F}_\alpha \to \varprojlim \mathcal{F}''_\alpha$$
est exacte. Cela résulte aussitôt de ce que le foncteur $\Gamma$ est exact à gauche, et de la propriété correspondante pour le foncteur $\varprojlim$ dans la catégorie des systèmes projectifs de groupes abéliens.

Soit $\varphi$ une application continue de $X$ dans un espace topologique $Y$. Si $(\mathcal{F}_\alpha)$ est un système projectif de faisceaux de groupes abéliens sur $X$, on a
$$\varphi_*(\varprojlim \mathcal{F}_\alpha) = \varprojlim \varphi_*(\mathcal{F}_\alpha) \,.$$
En effet, pour tout ouvert $V \subset Y$, on a par définition
$$\Gamma(V, \varprojlim \varphi_*(\mathcal{F}_\alpha)) = \varprojlim \Gamma(V, \varphi_*(\mathcal{F}_\alpha)) = \varprojlim \Gamma(\varphi^{-1}(V), \mathcal{F}_\alpha)$$
$$= \Gamma(\varphi^{-1}(V), \varprojlim \mathcal{F}_\alpha) = \Gamma(V, \varphi_*(\varprojlim \mathcal{F}_\alpha))$$

\newpage

%% ===== Page 184 =====

- 0-iii -

En particulier, si $X$ est une partie fermée de $Y$, et si on identifie un faisceau sur $Y$ de support contenu dans $X$ au faisceau qu'il induit sur $X$, on voit en prenant pour $\varphi$ l'injection canonique $X \to Y$ [unclear: XXXXXX] qu'une limite projective de tels faisceaux a encore son support contenu dans $X$.

Si maintenant $(G_{\alpha}, v_{\alpha\beta})$ est un système projectif de faisceaux de groupes abéliens sur $Y$, $(\varphi^*(G_{\alpha}))$ est un système projectif de faisceaux de groupes abéliens sur $X$, pour les homomorphismes $\varphi^*(v_{\alpha\beta})$; à toute section $s \in \Gamma(U, \varprojlim G_{\alpha})$, où $G = \varprojlim G_{\alpha}$, on peut associer la famille des sections $s_{\alpha} = v_{\alpha} \in \Gamma(U, \varphi^*(G_{\alpha}))$, et on définit ainsi un homomorphisme canonique $\varphi^*(\varprojlim G_{\alpha}) \to \varprojlim \varphi^*(G_{\alpha})$, mais cet homomorphisme n'est pas nécessairement bijectif.

Les considérations qui précèdent s'appliquent sans changement aux faisceaux d'anneaux et aux faisceaux de modules.

2. La condition de Mittag-Leffler.

Soit $(A_{\alpha}, u_{\alpha\beta})$ un système projectif de groupes abéliens. On appelle condition de Mittag-Leffler la condition suivante :

(ML) Pour tout indice $\alpha$, il existe $\beta \geqslant \alpha$ tel que, pour tout $\gamma \geqslant \beta$, on ait $u_{\alpha\gamma}(A_{\gamma}) = u_{\alpha\beta}(A_{\beta})$.

On notera que cette condition est satisfaite si les $u_{\alpha\beta}$ sont surjectives; il en est de même lorsque les $A_{\alpha}$ sont des modules artiniens sur un système projectif d'anneaux (les $u_{\alpha\beta}$ étant des homomorphismes), car alors les $u_{\alpha\beta}(A_{\beta})$, pour $\beta \geqslant \alpha$, forment un ensemble filtrant (pour $\supset$) de sous-modules de $A_{\alpha}$, et ont par suite un plus petit élément.

Lemme 1. - Soit $0 \to A_{\alpha}' \to A_{\alpha} \to A_{\alpha}'' \to 0$ une suite exacte de systèmes projectifs de groupes abéliens.

(i) Si $(A_{\alpha}')$ vérifie (ML), il en est de même de $(A_{\alpha}'')$.

(ii) Si $(A_{\alpha}')$ et $(A_{\alpha}'')$ vérifient (ML), il en est de même de $(A_{\alpha})$.

\newpage

%% ===== Page 185 =====

- 0-iv -

[crossed out: ] Soit $(u_{\alpha\beta})$ le système d'homomorphismes définissant le système projectif $(A_{\alpha})$; nous désignerons son noyau par $H_{\alpha\beta}$; le système d'homomorphismes $(u'_{\alpha\beta})$ définissant $(A'_{\alpha})$ est formé des restrictions [crossed out: aux] $u_{\alpha\beta}|_{A'_{\beta}}$; et le système d'homomorphismes $(u''_{\alpha\beta})$ définissant $(A''_{\alpha})=(A_{\alpha}/A'_{\alpha})$ [crossed out: est] obtenu en passant aux quotients à partir de $(u_{\alpha\beta})$; le noyau de $u''_{\alpha\beta}$ est donc $(H_{\alpha\beta}+A'_{\beta})/A'_{\beta}$.

(i) Si on a $u_{\alpha\beta}(A_{\beta})=u_{\alpha\beta}(A'_{\beta})$, il résulte du fait que $g_{\beta}$ et $g_{\gamma}$ sont surjectives, que $u''_{\alpha\beta}(A''_{\beta})=g_{\alpha}(u_{\alpha\beta}(A_{\beta}))=u''_{\alpha\gamma}(A''_{\gamma})$.

(ii) Soit $x_{\alpha} \in A_{\alpha}$ tel que $x_{\alpha}=u_{\alpha\beta}(x_{\beta})$ pour un $x_{\beta} \in A_{\beta}$; on en tire $g_{\alpha}(x_{\alpha})=u''_{\alpha\beta}(g_{\beta}(x_{\beta}))$; supposons qu'il existe $x_{\gamma} \in A_{\gamma}$ tel que $g_{\alpha}(x_{\alpha})=u''_{\alpha\gamma}(g_{\gamma}(x_{\gamma}))$; cela s'écrit $g_{\alpha}(x_{\alpha}-u_{\alpha\gamma}(x_{\gamma}))=0$, ou encore $u''_{\alpha\beta}(g_{\beta}(x_{\beta}-u_{\beta\gamma}(x_{\gamma})))=0$. Posons $z_{\beta}=x_{\beta}-u_{\beta\gamma}(x_{\gamma})$; on a par hypothèse $g_{\beta}(z_{\beta}) \in (H_{\alpha\beta}+A'_{\beta})/A'_{\beta}$, donc $z_{\beta} \in H_{\alpha\beta}+A'_{\beta}$; il existe par suite [crossed out: ] $x'_{\beta} \in A'_{\beta}$ tel que [crossed out: ] $x_{\beta}-u_{\beta\gamma}(x_{\gamma})=u_{\alpha\beta}(z_{\beta})=u_{\alpha\beta}(x'_{\beta})$. Supposons qu'il existe [crossed out: ] $x'_{\gamma} \in A'_{\gamma}$ tel que $u_{\alpha\beta}(x'_{\beta})=u_{\alpha\gamma}(x'_{\gamma})$; on en conclut $x_{\alpha}=u_{\alpha\gamma}(x_{\gamma}+x'_{\gamma})$, ce qui achève la démonstration.

Proposition 1. - Soit $I$ un ensemble ordonné filtrant admettant une partie cofinale dénombrable. Soit $0 \to A'_{\alpha} \to A_{\alpha} \to A''_{\alpha} \to 0$ une suite exacte de systèmes projectifs de groupes abéliens, ayant $I$ pour ensemble d'indices. Si $(A'_{\alpha})$ satisfait à la condition (ML), la suite
$$0 \to \varprojlim A'_{\alpha} \to \varprojlim A_{\alpha} \to \varprojlim A''_{\alpha} \to 0$$
est exacte.

Vu l'hypothèse, on peut se borner au cas où $I=\mathbb{N}$; par extraction d'une suite d'indices, on peut aussi supposer que, si $u_{nm}$ désigne l'homomorphisme [crossed out: ] $A_m \to A_n$ pour $n \leq m$, [crossed out: ] tous les [crossed out: ] sous-groupes $u_{nm}(A'_{\beta})$ sont égaux pour $n \leq m$; on désignera ce sous-groupe par $B_n$. Si on pose, pour simplifier, $v_n=u_{n,n+1}$, on a évidemment $v_n(B_{n+1})=B_n=v_n(A'_{n+1})$ et par suite $A'_{n+1}=B_{n+1}+(H_n \cap A'_{n+1})$, en désignant par $H_n$ le noyau de $v_n$. Tout revient à prouver que l'homomorphisme

\newpage

%% ===== Page 186 =====

- 0 - v -

$\varprojlim \varepsilon_n : \varprojlim A_n \to \varprojlim A''_n$ est surjective. Étant donnée une suite $(x''_n)$ telle que $x''_n \in A''_n$ et $x''_n = u''_n(x''_{n+1})$ pour tout $n$, nous allons construire par récurrence une suite $(x_n)$ telle que $x_n \in A_n$, $x''_n = \varepsilon_n(x_n)$ et $x_n = v_n(x_{n+1})$ pour tout $n$. Supposons construite une suite $(x_0, \dots, x_{n-1}, y_n)$ telle que $x_i \in A_i$ pour $i \leq n-1$ et $y_n \in A_n$, $x''_i = \varepsilon_i(x_i)$ pour $i \leq n-1$ et $x''_n = \varepsilon_n(y_n)$ et $x_i = v_i(x_{i+1})$ pour $i \leq n-2$, $x_{n-1} = v_{n-1}(y_n)$. Par hypothèse il existe $z_{n+1} \in A_{n+1}$ tel que $\varepsilon_{n+1}(z_{n+1}) = x''_{n+1}$, donc $\varepsilon_n(v_n(z_{n+1})) = x''_n = \varepsilon_n(y_n)$, autrement dit $v_n(z_{n+1}) - y_n \in A'_n = B_n + (H_{n-1} \cap A'_n)$. Il existe par suite $z'_n \in B_{n+1}$ tel que $y_n - v_n(z_{n+1} + z'_n) \in H_{n-1}$. Posons $y_{n+1} = z_{n+1} + z'_n$, et $x_n = v_n(y_{n+1})$; il est clair que $v_{n-1}(x_n) = v_{n-1}(y_n) = x_{n-1}$ ; et comme $y_n - x_n \in A'_n$, $\varepsilon_n(x_n) = x''_n$; enfin, comme $z'_n \in A'_{n+1}$, on a $\varepsilon_{n+1}(y_{n+1}) = x''_{n+1}$ ; et la récurrence peut se poursuivre, d'où la proposition.

Corollaire. — Les hypothèses sur $I$ étant les mêmes, soit $(K_\alpha)_{\alpha \in I}$ un système projectif de complexes $K = (K^n)_{n \in \mathbb{Z}}$ dont l'opérateur différentiel est de degré $+1$. Pour chaque $n$ il existe un homomorphisme canonique

$$ h_n : H^n(\varprojlim K_\alpha) \to \varprojlim H^n(K_\alpha) $$

Si, pour tout degré $n$, le système projectif de groupes abéliens $(K^n_\alpha)_{\alpha \in I}$ vérifie (ML), alors tous les homomorphismes $h_n$ sont surjectifs. Si en outre, pour un degré $n$, le système projectif $(H^{n-1}(K_\alpha))_{\alpha \in I}$ vérifie (ML), l'homomorphisme $h_n$ est bijectif.

Désignons par $Z^n$ et $B^n$ les groupes des cycles et des cobords, et posons $K^n = \varprojlim K_\alpha$; l'existence des homomorphismes $h_n$ provient de la commutativité des diagrammes

$$
\begin{tikzcd}
\dots \arrow[r] & K^{n-1} \arrow[d] \arrow[r] & K^n \arrow[d] \arrow[r] & K^{n+1} \arrow[d] \arrow[r] & \dots \\
\dots \arrow[r] & K^{n-1}_\alpha \arrow[r] & K^n_\alpha \arrow[r] & K^{n+1}_\alpha \arrow[r] & \dots
\end{tikzcd}
$$

Considérons les suites exactes

$$ (*_n) \quad 0 \to B^n(K_\alpha) \to Z^n(K_\alpha) \to H^n(K_\alpha) \to 0 $$

\newpage

%% ===== Page 187 =====

$$(**_n) \quad 0 \to Z^{n-1}(K_\alpha) \to K^{n-1} \to B^n(K_\alpha) \to 0$$
Comme $B^n(K_\alpha)$ est un quotient de $K^{n-1}$, l'hypothèse et le lemme 1 (i) montrent que $(B^n(K_\alpha))_{\alpha \in I}$ vérifie (ML) ; d'après la prop. 1, la suite
$$0 \to \varprojlim B^n(K_\alpha) \to \varprojlim Z^n(K_\alpha) \to \varprojlim H^n(K_\alpha) \to 0$$
est exacte. Comme $\varprojlim B^n(K_\alpha)$ s'identifie à un sous-groupe de $B^n(K)$ et $\varprojlim Z^n(K_\alpha)$ à un sous-groupe de $Z^n(K)$, $h_n$ est surjective. Si maintenant on suppose en outre que le système projectif $(H^{n-1}(K_\alpha))_{\alpha \in I}$ vérifie (ML), les suites exactes $(**_{n-1})$ et le lemme 1 (ii) montrent qu'il en est de même du système projectif $(Z^{n-1}(K_\alpha))_{\alpha \in I}$ ; mais alors la prop. 1 appliquée aux suites exactes $(**_n)$, prouve que la suite
$$0 \to \varprojlim Z^{n-1}(K_\alpha) \to K^{n-1} \to \varprojlim B^n(K_\alpha) \to 0$$
est exacte ; comme $\varprojlim Z^{n-1}(K_\alpha) \subset Z^{n-1}(K)$ et $\varprojlim B^n(K_\alpha) \supset B^n(K)$, les inclusions précédentes sont nécessairement des égalités, et $h_n$ est bijective.

\emph{Proposition 2.} — Soient $X$ un espace topologique, $(F_k)_{k \in \mathbb{N}}$ un système projectif de faisceaux de groupes abéliens sur $X$, $F = \varprojlim F_k$. On suppose vérifiées les conditions suivantes :
(i) Il existe une base $\mathcal{B}$ de la topologie de $X$ telle que, pour tout $U \in \mathcal{B}$ et tout $i > 0$, le système projectif $(H^i(U, F_k))_{k \in \mathbb{N}}$ vérifie (ML).
(ii) Pour tout $x \in X$ et tout $i > 0$, on a $\varprojlim (\varinjlim H^i(U, F_k)) = 0$, lorsque $U$ parcourt l'ensemble des voisinages de $x$.
(iii) Les homomorphismes $u_{hk}: F_k \to F_h$ définissant le système projectif $(F_k)$ sont surjectifs.

Dans ces conditions, pour tout $i \ge 0$, l'homomorphisme canonique
$$h_i : H^i(X, F) \to \varprojlim H^i(X, F_k)$$
est surjectif ; si en outre le système projectif $(H^{i-1}(X, F_k))_{k \in \mathbb{N}}$ vérifie (ML) pour une valeur de $i$, $h_i$ est bijectif.

a) Nous allons d'abord supposer que les $F_k$ sont flasques ainsi

\newpage

%% ===== Page 188 =====

- 0-vii -

que les noyaux des $u_{hk}$, et si (iii) est vérifiée, on a nécessairement $H^i(X, \mathcal{F}) = 0$ pour $i > 0$. Il suffira de prouver que pour tout ouvert $U$ de $X$ et tout recouvrement ouvert $\mathcal{U}$ de $U$, on a $\check{H}^i(\mathcal{U}, \mathcal{F}) = 0$ pour $i > 0$; il en résultera en effet d'abord que pour la cohomologie de Čech, on a $\check{H}^i(U, \mathcal{F}) = 0$ pour tout $i > 0$, puis en vertu de G,II,5.9.2 appliqué à l'ensemble de tous les ouverts de $X$ que $H^i(X, \mathcal{F}) = 0$ pour tout $i > 0$. Comme les $\mathcal{F}_k$ sont flasques, on a $H^i(U, \mathcal{F}_k) = 0$ pour $i > 0$ (G,II,5.2.3); considérons pour chaque $k$ le complexe $C^*(\mathcal{U}, \mathcal{F}_k)$ des cochaînes du nerf du recouvrement $\mathcal{U}$ (G,II,5.1), qui forment évidemment un système projectif de complexes. Montrons que toutes les applications $C^*(\mathcal{U}, \mathcal{F}_k) \to C^*(\mathcal{U}, \mathcal{F}_h)$ ($h \le k$) sont surjectives : il suffit évidemment par définition de montrer que pour tout ouvert $V$ de $X$, $\Gamma(V, \mathcal{F}_k) \to \Gamma(V, \mathcal{F}_h)$ est surjectif ; mais cela résulte de ce que si $N_{hk}$ est le noyau de $u_{hk}$, la suite $0 \to N_{hk} \to \mathcal{F}_k \to \mathcal{F}_h \to 0$ est exacte par hypothèse et donne donc la suite exacte de cohomologie
$$0 \to \Gamma(V, N_{hk}) \to \Gamma(V, \mathcal{F}_k) \to \Gamma(V, \mathcal{F}_h) \to H^1(V, N_{hk}) \to \dots$$
où $H^1(V, N_{hk}) = 0$ puisque $N_{hk}$ est flasque par hypothèse. Le système projectif $(C^*(\mathcal{U}, \mathcal{F}_k))_{k \in \mathbb{N}}$ vérifie (ML), et il en est de même de $(H^i(\mathcal{U}, \mathcal{F}_k))_{k \in \mathbb{N}}$ pour tout $i \ge 0$, car cela est trivial pour $i > 0$, et comme $H^0(\mathcal{U}, \mathcal{F}_k) = \Gamma(U, \mathcal{F}_k)$ (d'après ce qui précède), la condition (ML) est remplie aussi pour $i = 0$. On peut donc appliquer le cor. de la prop. 1, donc $H^i(\mathcal{U}, \mathcal{F}) = \varprojlim H^i(\mathcal{U}, \mathcal{F}_k) = 0$ pour tout $i > 0$.

b) Passons au cas général, et considérons la résolution canonique $C^*(\mathcal{X}, \mathcal{F}_k) = (C^i(\mathcal{X}, \mathcal{F}_k))_{i \ge 0}$ de $\mathcal{F}_k$ par des faisceaux flasques (G,II,4.3). Pour tout $i \ge 0$, considérons le système projectif $(C^i(\mathcal{X}, \mathcal{F}_k))_{k \in \mathbb{N}}$ de faisceaux flasques ; il satisfait aux conditions de a). En effet, si $N_{hk}$ est le noyau de $u_{hk}$ pour $h \le k$, la suite $0 \to N_{hk} \to \mathcal{F}_k \to \mathcal{F}_h \to 0$ est exacte et d'après (iii) notre assertion résulte de ce que le foncteur \marginpar{(G,II, 5.2.2)}

\newpage

%% ===== Page 189 =====

- 0-VIII -

$A \to \mathcal{C}^i(X, A)$ est exact (G,II,4.3). Soit $\mathcal{G}^i = \varprojlim \mathcal{C}^i(X, \mathcal{F}_k)$ ; on a donc $H^j(X, \mathcal{G}^i) = 0$ pour $j > 0$ et $i \ge 0$. Nous allons montrer que $\mathcal{G}^* = (\mathcal{G}^i)$ est une résolution du faisceau $\mathcal{F}$ ; de ce que $H^j(X, \mathcal{G}^*) = 0$ pour $j > 0$, il résultera alors que l'on pourra calculer la cohomologie $H^n(X, \mathcal{F})$ qui sera égale à $H^n(\Gamma(\mathcal{G}^*))$ (G,II,4.7.1).

Il est clair que, par passage à la limite projective, on déduit des suites exactes
$$ 0 \to \mathcal{F}_k \to \mathcal{C}^0(X, \mathcal{F}_k) \to \mathcal{C}^1(X, \mathcal{F}_k) \to \dots $$
un faisceau différentiel $\mathcal{G}^*$ :
$$ 0 \to \mathcal{F} \to \mathcal{G}^0 \to \mathcal{G}^1 \to \dots $$
Pour prouver notre assertion, il faut établir que $H^i(\mathcal{G}^*) = 0$ pour $i > 0$. Ce faisceau est engendré par le préfaisceau $U \mapsto H^i(\Gamma(U, \mathcal{G}^*))$ (du système projectif (G,II,4.1) ; or, le complexe $\Gamma(U, \mathcal{G}^*)$ est la limite projective des complexes $(\Gamma(U, \mathcal{C}^*(X, \mathcal{F}_k)))_{k \in \mathbb{N}}$. Pour chaque $i \ge 0$, les applications $\Gamma(U, \mathcal{C}^i(X, \mathcal{F}_k)) \to \Gamma(U, \mathcal{C}^i(X, \mathcal{F}_h))$ ($h \le k$) sont surjectives ; on a vu en effet ci-dessus que l'homomorphisme de faisceaux $\mathcal{C}^i(X, \mathcal{F}_k) \to \mathcal{C}^i(X, \mathcal{F}_h)$ est surjectif et que son noyau est flasque ; la suite exacte de cohomologie prouve alors notre assertion, comme dans a). Compte tenu de l'hypothèse (i) (où on notera que $H^i(U, \mathcal{F}_k) = H^i(\Gamma(U, \mathcal{C}^*(X, \mathcal{F}_k)))$, la résolution canonique $\mathcal{C}^*(U, \mathcal{F}_k|U)$ étant induite sur $U$ par $\mathcal{C}^*(X, \mathcal{F}_k)$), on peut appliquer le cor. de la prop. 1 au système projectif de complexes $(\Gamma(U, \mathcal{C}^*(X, \mathcal{F}_k)))_{k \in \mathbb{N}}$ ; et on a donc $H^i(\Gamma(U, \mathcal{G}^*)) = \varprojlim H^i(U, \mathcal{F}_k)$ pour tout $i \ge 0$. L'hypothèse (ii) prouve bien alors, par définition, que les faisceaux $H^i(\mathcal{G}^*)$ sont nuls pour $i > 0$.

On a donc pour tout $i \ge 0$, $H^i(X, \mathcal{F}) = H^i(\Gamma(X, \mathcal{G}^*))$, et $\Gamma(X, \mathcal{G}^*) = \varprojlim \Gamma(X, \mathcal{C}^*(X, \mathcal{F}_k))$. On vient de voir ci-dessus que les applications $\Gamma(X, \mathcal{C}^i(X, \mathcal{F}_k)) \to \Gamma(X, \mathcal{C}^i(X, \mathcal{F}_h))$ sont toutes surjectives ; la conclusion résulte donc encore du cor. de la prop. 1.

\newpage

%% ===== Page 190 =====

\marginpar{Archives Grothendieck sept 59}

\begin{center}
-- II-1 --
\end{center}

\section*{CHAPITRE II}
\section*{ETUDE ELEMENTAIRE DE QUELQUES CLASSES DE MORPHISMES}

Les morphismes définis et étudiés dans ce chapitre le sont sans utilisation des méthodes cohomologiques ; une étude plus poussée, utilisant ces dernières, en sera faite au chap. III.

\subsection*{§ 1. Compléments sur les faisceaux quasi-cohérents.}

\subsubsection*{1. Opérations sur les faisceaux quasi-cohérents.}

\textbf{Proposition 1.} -- Soit $f : X \to Y$ un morphisme de préschémas. On suppose qu'il existe un recouvrement $(Y_\alpha)$ de $Y$ par des ouverts affines tel que chacun des $f^{-1}(Y_\alpha)$ admette un recouvrement fini $(X_{\alpha i})$ par des ouverts affines contenus dans $f^{-1}(Y_\alpha)$, tel que chacune des intersections $X_{\alpha i} \cap X_{\alpha j}$ soit elle-même réunion finie d'ouverts affines. Dans ces conditions, pour tout faisceau quasi-cohérent $\mathcal{F}$ sur $X$, $f_*(\mathcal{F})$ est un faisceau quasi-cohérent sur $Y$.

La question étant locale (sur $Y$), on peut supposer $Y$ égal à l'un des $Y_\alpha$, donc supprimer les indices $\alpha$.

a) Supposons d'abord que les $X_i \cap X_j$ soient eux-mêmes des ouverts affines. Posons $\mathcal{F}_i = \mathcal{F}|_{X_i}$, $\mathcal{F}_{ij} = \mathcal{F}|_{X_i \cap X_j}$, et soient $\mathcal{F}_i$ et $\mathcal{F}_{ij}$ les images de $\mathcal{F}$ et $\mathcal{F}_{ij}$ respectivement par les restrictions respectives de $f$ en $X_i$ et $X_i \cap X_j$ ; on sait que les $\mathcal{F}_i$ et $\mathcal{F}_{ij}$ sont quasi-cohérents (I, 1, 6, cor. 2 de la prop. 8). Posons $\mathcal{G} = \bigoplus_i \mathcal{F}_i$, $\mathcal{H} = \bigoplus_{ij} \mathcal{F}_{ij}$ ; $\mathcal{G}$ et $\mathcal{H}$ sont quasi-cohérents (I, 1, 3, cor. 4 du th. 1). Nous voulons définir un homomorphisme $u : \mathcal{G} \to \mathcal{H}$ tel que $f_*(\mathcal{F})$ soit le noyau de $u$ ; il en résultera que $f_*(\mathcal{F})$ est quasi-cohérent (I, 1, 3, cor. 4 du th. 1). Il suffit de définir $u$ comme homomorphisme de préfaisceaux (G, I, 1, 9) ; tenant compte de la définition de $\mathcal{G}$ et $\mathcal{H}$, il suffit donc, pour tout ensemble ouvert $W \subset Y$, de définir un homomorphisme

$$u_W : \bigoplus_i \Gamma(f^{-1}(W) \cap X_i, \mathcal{F}) \to \bigoplus_{ij} \Gamma(f^{-1}(W) \cap X_i \cap X_j, \mathcal{F})$$

de façon à satisfaire aux conditions de compatibilité usuelles lorsque $W$ varie. Si, pour toute section $s_i \in \Gamma(f^{-1}(W) \cap X_i, \mathcal{F})$, on dé-

\newpage

%% ===== Page 191 =====

-- II - 2 --
signe par $s_{i|j}$ la restriction de $s_i$ à $f^{-1}(W) \cap X_i \cap X_j$, on posera
$$u_W((s_i)) = (s_{i|j} - s_{j|i})$$
et les conditions de compatibilité sont évidemment remplies. Pour prouver que le noyau $R$ de $u$ est $f_*(F)$, définissons un homomorphisme de $f_*(F)$ dans $R$ en faisant correspondre à toute section $s \in \Gamma(f^{-1}(W), F)$ la famille $(s_i)$, où $s_i$ est la restriction de $s$ à $f^{-1}(W) \cap X_i$ ; les axiomes (F 1) et (F 2) des faisceaux (G,II,1.1) entraînent que cet homomorphisme est bijectif, ce qui termine la démonstration dans ce cas.

b) Dans le cas général, le même raisonnement s'applique une fois que l'on a établi que les faisceaux $F_{ij}$ sont quasi-cohérents. Or, par hypothèse $X_i \cap X_j$ est réunion finie d'ouverts affines $X_{ijk}$ ; et comme les $X_{ijk}$ sont des ouverts affines dans un schéma, l'intersection de deux quelconques d'entre eux est encore un ouvert affine (I,3, prop. et cor. 1 de la prop.). On est donc ramené au premier cas, et la prop.1 est démontrée.

\emph{Corollaire.} - La conclusion de la prop.1 est valable dans chacun des cas suivants :

(i) $f$ est séparé, et $Y$ réunion d'une famille $(Y_\alpha)$ d'ouverts affines tels que chacun des $f^{-1}(Y_\alpha)$ soit réunion finie d'ouverts affines.

(ii) $f$ est séparé et de type fini.

(iii) $Y$ est réunion d'une famille $(U_\alpha)$ d'ouverts tels que $f^{-1}(U_\alpha)$ soit un espace noethérien (ce qui a lieu par exemple si $X$ est noethérien).

Dans le cas (i) les $X_i \cap X_j$ sont affines (I,3, prop.). Le cas (ii) est un cas particulier de (i). Enfin, dans le cas (iii), on peut se ramener au cas où $X$ est noethérien ; alors $Y$ affine $X$ admet un recouvrement ouvert affine fini $(X_i)$, et les $X_i \cap X_j$, étant quasi-compacts, sont réunions finies d'ouverts affines.

\emph{Proposition 2.} - Soit $f : X \to Y$ un morphisme de préschémas. Pour

\newpage

%% ===== Page 192 =====

- II-3 -

\noindent \emph{Proposition 2.} — Soit $f: X \to Y$ un morphisme de préschémas. Pour tout faisceau quasi-cohérent $\mathcal{G}$ sur $Y$, $f^*(\mathcal{G})$ est un faisceau quasi-cohérent sur $X$. Si de plus $X$ et $Y$ sont localement noethériens, pour tout faisceau cohérent $\mathcal{G}$ sur $Y$, $f^*(\mathcal{G})$ est un faisceau cohérent sur $X$.

La question est locale (sur $X$) ; soit $x \in X$, et soit $V$ un ouvert affine de $Y$ contenant $f(x)$ ; il y a un ouvert affine $U$ contenant $x$ et contenu dans $f^{-1}(V)$. Comme la restriction de $f^*(\mathcal{G})$ à $U$ est l'image réciproque de $\mathcal{G}|_V$ par la restriction de $f$ à $U$, on est ramené au cas où $X$ et $Y$ sont affines, et la proposition a été démontrée dans ce cas (I,1,6, cor. 2 de la prop. 8). Si en outre $X$ et $Y$ sont localement noethériens, on peut supposer les anneaux de $U$ et $V$ noethériens ; tenant compte de (I,1,6, prop. 8) et de la caractérisation des faisceaux cohérents sur un schéma affine (I,1,5, th. 3), la seconde assertion de la proposition revient à voir que si $N$ est un $B$-module de type fini, $M \otimes_B A$ est un $A$-module de type fini, ce qui est évident.

\noindent \emph{Proposition 3.} — Soit $X$ un préschéma (resp. un préschéma localement noethérien) ; soient $\mathcal{F}$ et $\mathcal{G}$ deux faisceaux quasi-cohérents (resp. cohérents) sur $X$ ; alors $\mathcal{F} \otimes_{\mathcal{O}_X} \mathcal{G}$ est quasi-cohérent (resp. cohérent). Si $\mathcal{F}$ est cohérent et $\mathcal{G}$ quasi-cohérent (resp. cohérent), alors $\underline{\operatorname{Hom}}_{\mathcal{O}_X}(\mathcal{F}, \mathcal{G})$ est quasi-cohérent (resp. cohérent).

La question étant locale, on peut supposer $X$ affine (resp. affine noethérien) ; en outre, si $\mathcal{F}$ est cohérent, on peut supposer qu'il est le conoyau d'un homomorphisme $\mathcal{O}_X^m \to \mathcal{O}_X^n$. L'assertion relative aux faisceaux quasi-cohérents résulte alors de (I,1,3, cor. 6 du th. 1) ; l'assertion relative aux faisceaux cohérents résulte de (I,1,5, th. 3) et du fait que si $M$ et $N$ sont des $A$-modules de type fini, il en est de même de $M \otimes_A N$ et de $\operatorname{Hom}_A(M, N)$ lorsque $A$ est noethérien.

\noindent \emph{Définition 1.} — Soient $X, Y$ deux $S$-préschémas, $p, q$ les projections de $X \times_S Y$, $\mathcal{F}$ (resp. $\mathcal{G}$) un faisceau quasi-cohérent sur $X$ (resp. $Y$). On

\newpage

%% ===== Page 193 =====

- II-4 -

appelle produit tensoriel de $F$ et $G$ sur $\mathcal{O}_S$ (ou sur $S$) et on note $F \otimes_S G$ (ou $F \otimes G$) le produit tensoriel $p^*(F) \otimes_{\mathcal{O}_{X \times_S Y}} q^*(G)$.

Ce produit tensoriel est donc un faisceau quasi-cohèrent sur $X \times_S Y$ ; il est cohèrent si $F$ et $G$ le sont, et si $X, Y$ et $X \times_S Y$ sont localement noethériens. On notera que si on prend $X=Y=S$, on retrouve la notion usuelle de produit tensoriel de faisceaux de $\mathcal{O}_S$-modules sur $S$. Si on prend $Y=S$, $F=\mathcal{O}_X$ et si on désigne par $f$ le morphisme structural $X \to Y$, on obtient le faisceau $f^*(G)$ : le produit tensoriel ordinaire et l'image réciproque apparaissent donc comme des cas particuliers du produit tensoriel général.

La déf. 1 entraîne immédiatement que, pour $X$ et $Y$ fixés, $F \otimes_S G$ est un bifoncteur covariant additif et exact à droite en $F$ et $G$.

Le préschéma $S$ étant fixé, on peut considérer les couples $(X, F)$, où $X$ est un $S$-préschéma et $F$ un faisceau quasi-cohèrent sur $X$, comme formant une catégorie : on définira un morphisme $(X, F) \to (Y, G)$ comme un couple $(u, v)$, où $u=(f, \phi)$ est un morphisme $X \to Y$, et $v$ un homomorphisme $G \to u_*(F)$ (cf. chap. 0, § 1, n°5) ; on sait (loc. cit.) qu'il revient au même de se donner un homomorphisme $u^*(G) \to F$.

Proposition 4. - Soient $f : T \to X$, $g : T \to Y$ deux $S$-morphismes, $F$ (resp. $G$) un faisceau quasi-cohèrent sur $X$ (resp. $Y$). On a alors
$$(f, g)^*(F \otimes_S G) = f^*(F) \otimes_{\mathcal{O}_T} g^*(G)$$
Si $p, q$ sont les projections de $X \times_S Y$, la formule résulte en effet des relations $(f, g)p = f$ et $(f, g)q = g$ et du fait qu'une image réciproque d'un produit tensoriel de faisceaux algébriques est le produit tensoriel des images réciproques de ces faisceaux (chap. 0, § 2 n°5).

Corollaire 1. - Soient $f : X \to X'$, $g : Y \to Y'$ deux $S$-morphismes, $F'$ (resp. $G'$) un faisceau quasi-cohèrent sur $X'$ (resp. $Y'$). On a alors

\newpage

%% ===== Page 194 =====

- II-5 -

$$(f \times_S g)^*(F' \otimes_{S'} G') = f^*(F') \otimes_S g^*(G')$$

Cela résulte de la prop. 4 et du fait que $f \times_S g = (f \circ p, g \circ q)_S$, $p$ et $q$ étant les projections de $X \times_S Y$.

\noindent \textit{Corollaire 2.} -- Soient $X, Y, Z$ trois $S$-préschémas, $F$ (resp. $G, H$) un faisceau quasi-cohérent sur $X$ (resp. $Y, Z$) ; le faisceau $F \otimes_S G \otimes_S H$ est l'image réciproque de $(F \otimes_S G) \otimes_S H$ par l'isomorphisme canonique de $X \times_S Y \times_S Z$ sur $(X \times_S Y) \times_S Z$.

En effet, cet isomorphisme s'écrit $(p_1, p_2)_S \times_S p_3$, en désignant par $p_1, p_2, p_3$ les projections de $X \times_S Y \times_S Z$.

\noindent \textit{Corollaire 3.} -- Si $X$ est un $S$-préschéma, tout faisceau quasi-cohérent $F$ sur $X$ est l'image réciproque de $F \otimes_S \mathcal{O}_S$ par l'isomorphisme canonique de $X$ sur $X \times_S S$ (I, 2, 5, prop. 8).

En effet, cet isomorphisme est $(1_X, \varphi)_S$, où $\varphi$ est le morphisme structural $X \to S$, et le corollaire résulte de la prop. 4 et du fait que $\varphi^*(\mathcal{O}_S) = \mathcal{O}_X$.

Soient $X$ un $S$-préschéma, $F$ un faisceau quasi-cohérent sur $X$, $S' \to S$ un morphisme ; on désigne par $F'$, ou $F^{S'}$, le faisceau quasi-cohérent $F \otimes_S \mathcal{O}_{S'}$ sur $X \times_S S' = X^{S'}$ ; donc $F^{S'} = p^*(F)$ ($p$ projection de $X^{S'}$ dans $X$).

\noindent \textit{Proposition 5.} -- Soit $\varphi' : S'' \to S'$ un morphisme. Pour tout faisceau quasi-cohérent $F$ sur un $S$-préschéma $X$, $(F^{S'})^{\varphi'}$ est l'image réciproque de $F^{S''}$ par l'isomorphisme canonique de $(X^{S'})^{S''}$ sur $X^{S''}$.

Cela résulte aussitôt des définitions et de (I, 2, 5, prop. 9).

Cela s'écrit encore $(F \otimes_S \mathcal{O}_{S'}) \otimes_{S'} \mathcal{O}_{S''} = F \otimes_S \mathcal{O}_{S''}$.

\noindent \textit{Proposition 6.} -- Soient $Y$ un $S$-préschéma, $f : X \to Y$ un $S$-morphisme. Pour tout faisceau quasi-cohérent $G$ sur $Y$ et tout morphisme $S' \to S$, on a $(f^*G)^{S'} = (f^{S'})^*(G^{S'})$.

Cela résulte aussitôt de la commutativité du diagramme
\begin{tikzcd}
X^{S'} \arrow[r, "f^{S'}"] \arrow[d] & Y^{S'} \arrow[d] \\
X \arrow[r, "f"] & Y
\end{tikzcd}

\newpage

%% ===== Page 195 =====

-- II-6 --

\emph{Corollaire} .- Soient $X, Y$ deux $S$-préschémas, $F$ (resp. $G$) un faisceau quasi-cohérent sur $X$ (resp. $Y$). L'image réciproque du faisceau $F^S \otimes G^S$ par l'isomorphisme canonique $(X \times_S Y)^S \to X^S \times_S Y^S$ (I, 2, 5, cor. 2 de la prop. 9) est égale à $(F \otimes G)^S$.

Si $p, q$ sont les projections de $X \times_S Y$, l'isomorphisme en question n'est autre que $(p^S, q^S)_S$; le corollaire résulte des prop. 4 et 6.

\emph{Proposition 7} .- Soient $S, X, Y$ trois schémas affines d'anneaux respectifs $A, B, C$. Soit $\tilde{M}$ (resp. $\tilde{N}$) un $B$-module (resp. $C$-module), $F = \tilde{M}$ (resp. $G = \tilde{N}$) le faisceau quasi-cohérent associé ; alors $F \otimes_S G$ est isomorphe au faisceau associé au $(B \otimes_A C)$-module $M \otimes_A N$.

En effet, en vertu de (I, 1, 6, prop. 3), $F \otimes_S G$ est isomorphe au faisceau associé au $(B \otimes_A C)$-module
$$(M \otimes_B (B \otimes_A C)) \otimes_{B \otimes_A C} ((B \otimes_A C) \otimes_C N)$$
et en raison des isomorphismes canoniques entre produits tensoriels, ce dernier est isomorphe à $M \otimes_B (B \otimes_A C) \otimes_C N = (M \otimes_B B) \otimes_A (C \otimes_C N) = M \otimes_A N$.

2. Prolongement des faisceaux quasi-cohérents. Application aux sous-préschémas.

\emph{Lemme 1} .- Soit $X$ un espace topologique, $F$ un faisceau d'ensembles sur $X$, $U$ une partie ouverte de $X$, $G$ un sous-faisceau de $F|U$. Soit $\overline{G}$ le sous-faisceau de $F$ tel que, pour tout ouvert $V$ de $X$, $\Gamma(V, \overline{G})$ soit formé des sections $s \in \Gamma(V, F)$ telles que la restriction de $s$ à $U \cap V$ soit une section de $G$ au-dessus de $U \cap V$. Alors on a $\overline{G}|U = G$, et $\overline{G}$ est le plus grand sous-faisceau de $F$ induisant sur $U$ un sous-faisceau de $G$.

Le fait que $\overline{G}$ soit un sous-faisceau de $F$ et les propriétés énoncées dans le lemme résultent aussitôt des définitions.

On notera que dans l'énoncé du lemme 1, on peut remplacer "ensembles" par "groupes", "modules", "anneaux", etc. (en général par des objets d'une catégorie admettant des limites inductives).

\emph{Lemme 2} .- Soit $X = \operatorname{Spec}(A)$ un schéma affine, $F = \tilde{M}$ le faisceau quasi-

\newpage

%% ===== Page 196 =====

- II-7 -

cohèrent sur $X$ associé à un $A$-module $M$, $\mathcal{G}$ le sous-faisceau de $F|D(f)$ associé à un sous-module $N'$ de $M_f$. Alors $\mathcal{G}$ est associé au sous-module de $M$ image réciproque de $N'$ par l'application canonique $i_f: A \to A_f$.

Il suffit de vérifier que pour tout $g \in A$, les sections du faisceau associé à $i_f^{-1}(N')$ sont les sections de $F$ induisant sur $D(f) \cap D(g) = D(fg)$ une section de $\mathcal{G}$ ; d'après (I, 1, 3, th. 1 et prop. 6) cela revient à montrer que tout élément $z/g^n \in M_g$ ($z \in M$) dont l'image canonique dans $M_{fg}$ appartient à $N'_{fg} = (N')_g$ ($g'$ image de $g$ dans $A_f$) est un élément de $N_g$. On a par hypothèse $z/g^n = z_1/(fg)^p$ avec $z_1 \in N'$, et on peut évidemment supposer $n=p$ ; cela signifie qu'il existe un entier $n$ tel que $(fg)^n (f^n g^n z - g^n z_1) = 0$, ou encore que dans $M_f$, on a $g^{n+m} z = f^{n+m} z_1/f^m \in N'$ ; et on conclut que $g^{n+m} z \in N$, et par suite $z/g^n = g^{n+m} z / g^{n+m} \in N_g$.

Théorème 1. - Soit $X$ un préschéma localement noethérien.

(i) Tout faisceau quasi-cohèrent (resp. cohèrent) sur un ouvert $U$ de $X$ est la restriction à $U$ d'un faisceau quasi-cohèrent (resp. cohèrent) sur $X$.

(ii) Soient $F$ un faisceau quasi-cohèrent (resp. cohèrent) sur $X$, $U$ un ouvert de $X$, $\mathcal{G}$ un sous-faisceau quasi-cohèrent de $F|U$. Alors le plus grand sous-faisceau de $F$ induisant sur $U$ un sous-faisceau de $F|U$ contenu dans $\mathcal{G}$ (lemme 1) est quasi-cohèrent (resp. cohèrent).

(iii) Soient $F$ un faisceau quasi-cohèrent sur $X$, $U$ un ouvert de $X$, $\mathcal{G}$ un sous-faisceau cohèrent de $F|U$. Alors il existe un sous-faisceau cohèrent de $F$ induisant $\mathcal{G}$ sur $U$.

(i) Si $X$ est affine, il est noethérien (I, 4, prop.) et tout ouvert de $X$ est quasi-compact, donc on est ramené au cas traité dans (I, 1, 4, cor. du th. 2) et (I, 1, 5, cor. du th. 3). Dans le cas général, considérons l'ensemble $U$ des couples $(V, H)$ où $V$ est

\newpage

%% ===== Page 197 =====

un ouvert de $X$ contenant $U$, et $H$ un faisceau quasi-cohérent (resp. cohérent) sur $V$ induisant le faisceau donné $F$ sur $U$. L'ensemble $\mathcal{U}$ est inductif quand on l'ordonne par la relation de restriction : en effet, si $\mathcal{L}$ est une partie totalement ordonnée de $\mathcal{U}$, $V_1$, la réunion des ensembles $V$ tels que $(V, H) \in \mathcal{L}$, il est clair qu'il y a sur $V_1$ un unique faisceau $H_1$ dont la restriction à tout $V$ tel que $(V, H) \in \mathcal{L}$ est le faisceau $H$ ; tout revient à voir que $H_1$ est quasi-cohérent (resp. cohérent). La question étant locale, on remarque que pour tout $x \in V_1$, il y a un voisinage ouvert affine $W \subset V_1$ de $x$ qui soit un espace noethérien ; comme il y a donc un ouvert maximal parmi les $W \cap V$ tels que $(V, H) \in \mathcal{L}$, on conclut qu'il y a un $(V, H) \in \mathcal{L}$ pour lequel $W \subset V$, et cela entraîne notre assertion. Cela étant, soit $(V_0, H_0)$ un élément maximal de $\mathcal{U}$ ; montrons que $V_0 = X$. Supposons le contraire et soit $x \in X$ non dans $V_0$ ; soit $T$ un voisinage ouvert affine de $x$. Comme on l'a vu plus haut, il y a un faisceau quasi-cohérent $G$ sur $T$ tel que $G|_{V_0 \cap T} = H_0|_{V_0 \cap T}$ ; si $H'$ désigne le faisceau sur $V' = V_0 \cup T$ qui coïncide avec $H_0$ sur $V_0$ et avec $G$ sur $T$, on aurait $(V', H') \in \mathcal{U}$ contrairement à la définition de $(V_0, H_0)$, ce qui achève de démontrer (i).

(ii) D'après la définition de $\overline{G}$ (lemme 1), il est clair que pour tout ouvert $V$ de $X$, on a $\overline{G}|_V = \overline{G'}$, en désignant par $G'$ le faisceau $G|_{U \cap V}$, et par $\overline{G'}$ le faisceau correspondant sur $V$. Cette remarque montre que l'on peut se limiter au cas où $X$ est affine et noethérien. Alors $U$ est réunion finie d'ouverts affines $D(f_i)$ ; la définition de $\overline{G}$ montre que ce faisceau est l'intersection des sous-faisceaux $\overline{G_i}$ de $F$, où $G_i = G|_{D(f_i)}$. Or, on sait (lemme 2) que les $\overline{G_i}$ sont quasi-cohérents ; tout revient donc à voir que l'intersection de deux sous-faisceaux $M, N$ d'un faisceau $\tilde{F}$ ($M, N$ étant des $A$-modules) est quasi-cohérent, ce qui résulte aussitôt de ce que $M \cap N$ est le noyau de l'homomorphisme $M \to P/N$ (1.1.3).

\newpage

%% ===== Page 198 =====

-- II-9 --

cor.4 du th.1) . Cela établit l'assertion relative au cas où $F$ est quasi-cohérent ; l'assertion relative au cas où $F$ est cohérent est conséquence du lemme suivant :

Lemme 3. -- Si $X$ est localement noethérien, tout sous-faisceau (resp. tout faisceau quotient) quasi-cohérent d'un faisceau cohérent $F$ est un faisceau cohérent.

La question étant locale, on peut supposer $X$ affine, et l'anneau $A(X)=A$ noethérien ; $F$ est alors isomorphe à un faisceau $\widetilde{M}$ associé à un $A$-module $M$ de type fini (I,1,5,th.3). Tout sous-faisceau (resp. faisceau quotient) quasi-cohérent de $\widetilde{M}$ est alors isomorphe à un faisceau associé à un sous-module (resp. module quotient) de $M$ (I,1,3, cor.4 du th.1) ; un tel sous-module (resp. module quotient) étant de type fini, le lemme en résulte (I,1,5,th.3).

(iii) Supposons d'abord $X$ affine et l'anneau $A=A(X)$ noethérien ; alors d'après (ii) le faisceau $G$ est quasi-cohérent, donc on peut le supposer de la forme $\widetilde{M}$, où $M$ est un $A$-module. Or, $M$ est limite inductive d'un système inductif $(M_\lambda)$ des $A$-modules de type fini. D'autre part, $U$ étant quasi-compact est réunion d'un nombre fini d'ouverts affines $D(f_i)$ avec $f_i \in A$ ($1 \le i \le n$) ; le faisceau $\widetilde{M}|D(f_i)$ est associé au $A_{f_i}$-module $M_{f_i}$, et ce dernier est limite inductive des $(M_\lambda)_{f_i}$ (chap.0, § 1, n°3). Or, $A_{f_i}$ est noethérien, et $\widetilde{G}|D(f_i)=G|D(f_i)$ est cohérent par hypothèse, donc (I,1,5,th.3) $M_{f_i}$ est de type fini ; les $(M_\lambda)_{f_i}$ étant aussi de type fini, il existe un indice $\lambda_i$ tel que $(M_\lambda)_{f_i} = M_{f_i}$ pour $\lambda \ge \lambda_i$. Soit alors $\mu$ un indice $\ge \lambda_i$ pour $1 \le i \le n$ ; on a donc $(M_\mu)_{f_i} = M_{f_i}$ pour $1 \le i \le n$, et on conclut que le sous-faisceau cohérent $\widetilde{M}_\mu$ induit $G|(D(f_i))$ sur chacun des $D(f_i)$, et par suite induit $G$ sur $U$.

Pour passer au cas général, on considère de nouveau l'ensemble $\mathcal{M}$ des couples $(V,H)$, où $V$ est un ouvert de $X$ contenant $U$, et $H$ un sous-faisceau cohérent de $F|V$ induisant ...

\newpage

%% ===== Page 199 =====

- II - 10 -

on montre alors par le même raisonnement que dans (i) que $M$ est inductif et que si $(V, M)$ en est un élément maximal, on a nécessairement $V = X$ (en utilisant le résultat démontré pour $X$ affine d'anneau noethérien). C.Q.F.D.

Corollaire. - Sur un préschéma localement noethérien $X$, tout faisceau quasi-cohérent $\mathcal{F}$ est limite inductive de ses sous-faisceaux cohérents.

Pour tout ouvert affine $U$ d'anneau noethérien, $\mathcal{F}|U$ est associé à un $A(U)$-module $M_U$, lequel est limite inductive de sous-modules $M_{U,\lambda}$ de type fini $(I, 1, 3, \text{cor.} 4 \text{ du th. } 1)$; $\mathcal{F}|U$ est par suite limite inductive des sous-faisceaux cohérents $\widetilde{M}_{U,\lambda}$ $(I, 1, 5, \text{th.} 3)$. D'après le th. 1 (iii) $\widetilde{M}_{U,\lambda}$ est induit sur $U$ par un sous-faisceau cohérent $\mathcal{G}_{\lambda}$ de $\mathcal{F}$. Les sommes finies de sous-faisceaux $\mathcal{G}_{\lambda}$ sont des sous-faisceaux cohérents de $\mathcal{F}$ $(I, 2, 13, \text{cor.} \text{ du th.} 2)$, et il est clair que $\mathcal{F}$ est limite inductive de ces sous-faisceaux.

Proposition 8. - Soient $X$ un préschéma localement noethérien, $Y$ un sous-préschéma de $X$. Il existe un plus petit sous-préschéma $\overline{Y}$ de $X$ majorant $(I, 3, 1)$ $Y$, dont l'espace de base est l'adhérence de l'espace de base de $Y$, et dont le faisceau structural induit $\mathcal{O}_Y$ sur $Y$.

Soit $U$ un ouvert de $X$ tel que $Y$ soit un sous-préschéma fermé du préschéma induit sur $U$; $\mathcal{O}_Y$ est donc la restriction à $Y$ d'un faisceau $(\mathcal{O}_X|U)/\mathcal{J}$, où $\mathcal{J}$ est un faisceau quasi-cohérent d'idéaux; mais comme $X$ est localement noethérien et que tout idéal d'un anneau noethérien est de type fini, $\mathcal{J}$ est en fait un faisceau cohérent $(I, 1, 3, \text{th.} 1$ et $I, 1, 5, \text{th.} 3)$. Les sous-préschémas fermés de $X$ majorant $Y$ correspondent bijectivement aux sous-faisceaux cohérents $\mathcal{J}'$ d'idéaux sur $X$, induisant sur $U$ un faisceau contenu dans $\mathcal{J}$; parmi ces faisceaux il y en a un plus grand $\overline{\mathcal{J}}$ (th. 1, (ii) appliqué à $\mathcal{O}_X$ et $\mathcal{J}$), le sous-préschéma fermé défini par $\overline{\mathcal{J}}$ est le...

\newpage

%% ===== Page 200 =====

- II-11 -

sous-préschéma cherché ; il est clair par définition de $\mathcal{J}$ (lemme 1) que le support de $\mathcal{O}_X/\mathcal{J}$ est contenu dans l'adhérence de $Y$, et comme ce support est fermé, il est égal à cette adhérence. Comme $Y$ est localement fermé, il est ouvert dans $\overline{Y}$, et il résulte du lemme 1 que $Y$ est le sous-préschéma induit par $\overline{Y}$.

Corollaire 1. - Toute section de $\mathcal{O}_{\overline{Y}}$ sur un ouvert $V$ de $\overline{Y}$ qui est nulle dans $V \cap Y$ est nulle.

En effet, soit $s$ une telle section ; pour tout $x \in \overline{Y}$ il existe un voisinage ouvert $W$ de $x$ dans $X$ et une section $s'$ de $\mathcal{O}_X$ sur $W$ tels que l'image de $s'$ dans $\mathcal{O}_X/\mathcal{J}$ coïncide avec $s$ dans $W \cap \overline{Y}$ ; l'hypothèse signifie que $s'(y) \in \mathcal{J}(y)$ pour tout $y \in W \cap Y$, et par définition de $\mathcal{J}$ (lemme 1) on a donc aussi $s'(z) \in \mathcal{J}(z)$ pour tout $z \in W$, d'où en particulier $s(x)=0$.

Corollaire 2. - Pour tout ouvert $V$ de $X$, le sous-préschéma $\overline{Y \cap V}$ est égal au préschéma induit par $\overline{Y}$ sur l'ensemble ouvert $\overline{Y \cap V}$.

En effet, si $\mathcal{J}'$ est le faisceau $\mathcal{J}|(U \cap V)$, il résulte du lemme 1 que $\mathcal{J}'$ et $\mathcal{J}$ induisent le même faisceau sur $V$.

Corollaire 3. - Soient $X$ un $S$-préschéma, $Z$ un préschéma séparé au-dessus de $S$. Soient $Y$ un sous-préschéma de $X$ ; si deux $S$-morphismes $f, g$ de $X$ dans $Z$ ont même restriction à $Y$, ils sont identiques.

Soit $h=(f,g)_S : X \to Z \times_S Z$ ; comme la diagonale $T=\Delta_Z(Z)$ est un sous-préschéma fermé de $Z \times_S Z$, $Y'=h^{-1}(T)$ est un sous-préschéma fermé de $Y$ (I,3,2, cor. 3 de la prop. 5). Si $u:Y \to Z$ est la restriction commune de $f$ et $g$ à $Y$, la restriction de $h$ à $Y$ est $h'=(u,u)_S$, qui se factorise en $\Delta_Z \circ u$ ; comme $\Delta_Z^{-1}(T)=Z$, on a $h'^{-1}(T)=Y$, et par suite $Z'$ est un sous-préschéma fermé de $X$ induisant $Y$ ; en d'autres termes $Z'$ majore $Y$, donc aussi $\overline{Y}$ par définition. Autrement dit, $h^{-1}(T)=\overline{Y}$, et par suite $h$ se factorise en $\overline{\Delta}_Z \circ v$, où $v$ est un morphisme $\overline{Y} \to Z$, ce qui entraîne $f=g=v$.

3. Faisceaux divisoriels.

\newpage

%% ===== Page 201 =====

- II-12 -

Dans ce n°, $X$ désigne un espace annelé, non nécessairement un préschéma.

\textbf{Définition 1.} - On dit qu'un faisceau algébrique $\mathcal{F}$ sur $X$ est localement libre (resp. localement libre de rang $n$) si pour tout $x \in X$, il existe un voisinage ouvert $U$ de $x$ tel que $\mathcal{F}|U$ soit isomorphe à un faisceau de la forme $\mathcal{O}_X^{(I)}|U$ (resp. à $\mathcal{O}_X^n|U$). Un faisceau localement libre de rang $1$ est encore appelé faisceau divisoriel. Il est clair qu'un faisceau localement libre est quasi-cohérent, un faisceau localement libre de rang fini cohérent.

A) Opérations tensorielles.

\textbf{Proposition 9.} - Si $L, L'$ sont deux faisceaux divisoriels, il en est de même de $L \otimes_{\mathcal{O}_X} L'$ et de $\mathcal{H}om_{\mathcal{O}_X}(L, \mathcal{O}_X)$, noté aussi $L^{-1}$.

La proposition est évidente, puisque $\mathcal{O}_X \otimes_{\mathcal{O}_X} \mathcal{O}_X$ et $\mathcal{H}om_{\mathcal{O}_X}(\mathcal{O}_X, \mathcal{O}_X)$ s'identifient canoniquement à $\mathcal{O}_X$.

On écrit $L^{\otimes n}$ pour le produit tensoriel de $n$ faisceaux identiques à un faisceau divisoriel $L$ ($n \ge 1$) ; on pose par convention $L^{\otimes 0} = \mathcal{O}_X$ et $L^{(-n)} = (L^{-1})^{\otimes n}$ pour $n \ge 1$. Alors :

\textbf{Proposition 10.} - Quels que soient les entiers rationnels $n, m$, il existe un isomorphisme canonique fonctoriel $L^{\otimes n} \otimes_{\mathcal{O}_X} L^{\otimes m} \sim L^{\otimes (n+m)}$.

Il suffit évidemment de définir un isomorphisme $L \otimes_{\mathcal{O}_X} L^{-1} \sim \mathcal{O}_X$ : si $s \in \Gamma(U, L)$ et $u \in \Gamma(U, L^{-1})$, on fait correspondre à $s \otimes u$ la section $u(s) \in \Gamma(U, \mathcal{O}_X)$, ce qui définit l'isomorphisme cherché.

\textbf{Proposition 11.} - Soient $L, L'$ deux faisceaux divisoriels ; il existe un isomorphisme canonique fonctoriel $L^{-1} \otimes L' \sim \mathcal{H}om_{\mathcal{O}_X}(L, L')$.

Pour tout couple $(u, t)$, où $u \in \Gamma(U, L^{-1})$ et $t \in \Gamma(U, L')$, on fait correspondre à $u \otimes t$ l'élément $s \to u(s)t$ de $\Gamma(U, \mathcal{H}om_{\mathcal{O}_X}(L, L'))$. Pour voir qu'on a un isomorphisme, on peut supposer que $L|U$ et $L'|U$ sont isomorphes à $\mathcal{O}_X|U$ ; autrement dit on peut supposer que $L = L' = \mathcal{O}_X$, et alors c'est immédiat.

B) Groupe de Picard.

Se donner un faisceau divisoriel revient à se donner un recouvrement...

\newpage

%% ===== Page 202 =====

- II-13 -

ment $U = (U_\lambda)_{\lambda \in L}$ de $X$ par des ouverts, pour tout couple $(\lambda, \mu)$ d'indices un automorphisme $\theta_{\lambda\mu}$ de $\mathcal{O}_X|_{U_\lambda \cap U_\mu}$ de sorte que si on désigne par $\theta_{\lambda\mu, \nu}$ les restrictions de $\theta_{\lambda\mu}$ à $U_\lambda \cap U_\mu \cap U_\nu$, on ait $\theta_{\lambda\mu} \cdot \theta_{\mu\nu} \cdot \theta_{\nu\lambda} = 1$ ; cela permet en effet de "recoller" les faisceaux $\mathcal{O}_X|_{U_\lambda}$ au moyen des $\theta_{\lambda\mu}$ (FAG, I, 1, 4, prop. 4) ce qui donne évidemment un faisceau divisoriel. Si on observe qu'un automorphisme $\theta_{\lambda\mu}$ de $\mathcal{O}_X|_{U_\lambda \cap U_\mu}$ équivaut à la donnée d'une section du faisceau des groupes d'éléments inversibles $\mathcal{O}_X^*$ dont les fibres sont les groupes d'éléments inversibles $\mathcal{O}_x^*$ des anneaux $\mathcal{O}_x$, au-dessus de $U_\lambda \cap U_\mu$ on voit que tout cocycle de $U$ à valeurs dans $\mathcal{O}_X^*$ définit un faisceau divisoriel (noter que les cochaînes sont ici notées multiplicativement) ; par passage à un recouvrement plus fin $V$, le cocycle correspondant donne le même faisceau divisoriel ; en outre, deux cocycles cohomologues donnent deux faisceaux divisoriels isomorphes. On obtient ainsi une application surjective du groupe de cohomologie $H^1(X, \mathcal{O}_X^*)$ (égal, comme on sait, au groupe de cohomologie de Cech $\check{H}^1(X, \mathcal{O}_X^*)$) (G, II, 5.9, cor. du th. 5.9.1)) dont les valeurs sont les classes de faisceaux divisoriels isomorphes sur $X$. On voit donc tout d'abord que ces classes forment un ensemble, que nous noterons $\operatorname{Pic}(X)$. De plus, l'application $H^1(X, \mathcal{O}_X^*) \to \operatorname{Pic}(X)$ est en fait bijective, car par définition tout isomorphisme d'un faisceau divisoriel $L$ sur un autre $L'$ se représente, pour un recouvrement $U$ assez fin, par un 1-cobord de $U$ à valeurs dans $\mathcal{O}_X^*$.

Notons maintenant que si, pour deux éléments $l, l'$ de $\operatorname{Pic}(X)$, on désigne par $ll'$ la classe du faisceau divisoriel $L \otimes_{\mathcal{O}_X} L'$, où $L \in l$, $L' \in l'$, on définit sur $\operatorname{Pic}(X)$ une structure de groupe abélien, appelé groupe de Picard de $X$. Pour cette structure de groupe, la bijection $H^1(X, \mathcal{O}_X^*) \to \operatorname{Pic}(X)$ est un isomorphisme. En effet, soient $L, L'$ deux faisceaux divisoriels, $U = (U_\alpha)$ un recouvrement ouvert de $X$ tel que les sections de $L, L'$ au-dessus de $U_\alpha$ soient de la forme

\newpage

%% ===== Page 203 =====

- II - 14 -

Proposition 12. — Soit $f : X \to Y$ un morphisme d'espaces annelés. Pour tout faisceau divisoriel $L$ sur $Y$, l'image réciproque $f^*(L)$ est un faisceau divisoriel sur $X$. En outre, si $L, L'$ sont deux faisceaux divisoriels sur $Y$, on a $f^*(L \otimes_{\mathcal{O}_Y} L') = f^*(L) \otimes_{\mathcal{O}_X} f^*(L')$ et $f^*(L^{-1}) = (f^*(L))^{-1}$.

La première assertion résulte de ce que les images réciproques par $f$ de deux faisceaux localement isomorphes sont localement isomorphes et de ce que $f^*(\mathcal{O}_Y) = \mathcal{O}_X$ (chap. 0, § 2, n° ); les autres assertions ont été démontrées au chap. 0, § 2, n°.

Corollaire. — Par passage aux classes de faisceaux divisoriels isomorphes, l'application $L \to f^*(L)$ définit un homomorphisme $f^* : \operatorname{Pic}(Y) \to \operatorname{Pic}(X)$.

On peut d'ailleurs voir aisément que si on identifie $\operatorname{Pic}(X)$ et $\operatorname{Pic}(Y)$ à $H^1(X, \mathcal{O}_X^*)$ et $H^1(Y, \mathcal{O}_Y^*)$ au moyen de l'isomorphisme canonique défini ci-dessus, l'homomorphisme $f^*$ s'identifie à l'homomorphisme $H^1(Y, \mathcal{O}_Y^*) \to H^1(X, \mathcal{O}_X^*)$ déduit de l'homomorphisme $\theta : f^*(\mathcal{O}_Y^*) \to \mathcal{O}_X^*$ dans la cohomologie de Čech (on a posé $f = (\varphi, \theta)$) : il suffit de remarquer que si $U_\alpha, U_\beta$ sont deux ouverts de $Y$, $s_\alpha, s_\beta$ des sections de $\mathcal{O}_Y$ au-dessus de $U_\alpha$ et $U_\beta$ respectivement telles que $s_\alpha = \varepsilon_{\alpha\beta} s_\beta$ au-dessus de $U_\alpha \cap U_\beta$, avec $\varepsilon_{\alpha\beta} \in \Gamma(U_\alpha \cap U_\beta, \mathcal{O}_Y^*)$, les sections correspondantes $(s_\alpha \circ \varphi) \otimes 1$ et $(s_\beta \circ \varphi) \otimes 1$ de $\mathcal{O}_X$ au-dessus de $\varphi^{-1}(U_\alpha)$ et $\varphi^{-1}(U_\beta)$ sont telles que, au-dessus de $\varphi^{-1}(U_\alpha \cap U_\beta)$, on ait
$$(s_\alpha \circ \varphi) \otimes 1 = ((\varepsilon_{\alpha\beta} \circ \varphi) (s_\beta \circ \varphi)) \otimes 1 = (\theta(\varepsilon_{\alpha\beta} \circ \varphi)) \cdot ((s_\beta \circ \varphi) \otimes 1).$$

\newpage

%% ===== Page 204 =====

- II-15 -

C) \emph{Anneaux et modules gradués construits à l'aide d'un faisceau divisoriel.}

Soit $L$ un faisceau divisoriel sur $X$. On désigne par $\Gamma_*(L)$ le groupe abélien somme directe $\bigoplus_{n \in \mathbb{Z}} \Gamma(X, L^{\otimes n})$ ; on le munit d'une structure d'anneau gradué, en faisant correspondre au couple $(s_n, s_m)$, avec $s_n \in \Gamma(X, L^{\otimes n})$, $s_m \in \Gamma(X, L^{\otimes m})$ la section de $L^{\otimes (n+m)}$ sur $X$ qui correspond canoniquement (prop. 10) à $s_n \otimes s_m$. Il est clair que, pour $X$ fixé, $\Gamma_*(L)$ est un foncteur covariant par rapport à $L$. En outre, si $f: X \to Y$ est un morphisme d'espaces annelés ($f = (\phi, \theta)$), pour tout faisceau divisoriel $L$ sur $Y$, l'application $s_n \mapsto (s_n \phi) \otimes 1$ est un homomorphisme $\psi_n : \Gamma(Y, L^{\otimes n}) \to \Gamma(X, f^*(L^{\otimes n}))$, et la prop. 12 montre que ces homomorphismes définissent un homomorphisme d'anneaux gradués $f^* : \Gamma_*(L) \to \Gamma_*(f^*(L))$.

Si maintenant $F$ est un faisceau algébrique quelconque sur $X$, $L$ un faisceau divisoriel sur $X$, on pose
$$\Gamma_*(L, F) = \bigoplus_{n \in \mathbb{Z}} \Gamma(X, F \otimes L^{\otimes n})$$
On munit ce groupe abélien d'une structure de module gradué sur l'anneau gradué $\Gamma_*(L)$ de la façon suivante : au couple $(s_n, u)$, où $s_n \in \Gamma(X, L^{\otimes n})$ et $u \in \Gamma(X, F \otimes L^{\otimes m})$, on fait correspondre la section de $F \otimes L^{\otimes (n+m)}$ qui correspond canoniquement (prop. 10) à $s_n \otimes u$. Il est clair que, pour $X$ et $L$ fixé, $\Gamma_*(L, F)$ est un foncteur covariant dans la catégorie des $\Gamma_*(L)$-modules gradués ; pour $X$ et $F$ fixés, c'est un foncteur covariant en $L$ à valeurs dans la catégorie des groupes abéliens. Enfin, si $f: Y \to X$ est un morphisme d'espaces annelés (de degré 0), on définit comme ci-dessus un \emph{di-homomorphisme} de modules gradués $f^* : \Gamma_*(L, F) \to \Gamma_*(f^*(L), f^*(F))$.

D) \emph{Ouverts définis par une section d'un faisceau divisoriel.}

Supposons désormais que tous les anneaux $\mathcal{O}_X(x)$ soient des anneaux locaux ; soit $L$ un faisceau divisoriel sur un ouvert $V \subset X$, et soit $f$ une section de $L$ au-dessus de $V$. Nous désignerons par

\newpage

%% ===== Page 205 =====

\marginpar{Archives Grothendieck sept. 59}
- II-16 -

l'ensemble des $x \in V$ ayant la propriété suivante : il existe un voisinage ouvert $W \subset V$ de $x$ et un isomorphisme $\theta : L|_W \to \mathcal{O}_X|_W$ tel que la section $\theta(f|_W) = s$ de $\mathcal{O}_X$ au-dessus de $W$ soit telle que $s(x)$ n'appartienne pas à l'idéal maximal de $\mathcal{O}_X(x)$. On notera que cette propriété est alors valable pour tout voisinage ouvert $W' \subset W$ de $x$ et tout isomorphisme $\theta' : L|_{W'} \to \mathcal{O}_X|_{W'}$, car au-dessus de $W \cap W'$, la section $s' = \theta'(f|_{W'})$ est de la forme $y \mapsto \epsilon(y)s(y)$, où $\epsilon(y)$ est inversible dans $\mathcal{O}_X(y)$.

Proposition 13. — Si $X$ est un préschéma, $L$ un faisceau divisoriel sur $X$, l'ensemble $V_f$ est ouvert dans $X$ pour toute partie ouverte $V$ de $X$ et toute section $f$ de $L$ au-dessus de $V$.

Comme pour tout ouvert $W \subset V$, on a évidemment $W \cap V_f = V_{f|_W}$, il suffit de considérer le cas où $V$ est affine et où $L = \mathcal{O}_X$ considéré comme module sur lui-même; mais alors, $X_f$ n'est autre que l'ensemble $D(f)$ défini dans 1,1,1.

Proposition 14. — Soient $L, L'$ deux faisceaux divisoriels sur un espace annelé $X$; si $f \in \Gamma(X, L)$, $g \in \Gamma(X, L')$, on a
$$(1) \quad X_f \cap X_g = X_{f \otimes g}$$

Par définition, on peut en effet se ramener au cas où $L = L' = \mathcal{O}_X$, et alors $f \otimes g$ s'identifie au produit $fg$, et la proposition est évidente.

Par abus de langage, on dit parfois (lorsqu'aucune confusion n'en résulte) que $X_f$ est l'ensemble des $x \in X$ où $f$ ne s'annule pas (ce qui revient à remplacer en chaque point $f(x)$ par sa classe dans le corps résiduel de $\mathcal{O}_X(x)$).

4. Prolongement de sections de faisceaux quasi-cohérents.

Théorème 2. — Soit $X$ un préschéma dont l'espace de base est noethérien ou un schéma dont l'espace de base est quasi-compact. Soient $L$ un faisceau divisoriel sur $X$, $f$ une section de $L$ au-dessus de $X$,

\newpage

%% ===== Page 206 =====

- II-17 -

$F$ un faisceau quasi-cohérent sur $X$.

(i) Si $s \in \Gamma(X, F)$ est telle que $s|_{X_f} = 0$, alors il existe un entier $n > 0$ tel que $f^n \otimes s = 0$.

(ii) Pour toute section $s \in \Gamma(X_f, F)$, il existe un entier $n > 0$ tel que $s \otimes f^n$ se prolonge en une section de $F \otimes L^{\otimes n}$ au-dessus de $X$.

(i) Comme l'espace de base de $X$ est quasi-compact, donc réunion finie d'ouverts affines, on est aussitôt ramené au cas où $X$ est affine et $L = \mathcal{O}_X$. Dans ce cas $f$ s'identifie à un élément de $A(X)$ et on a $X_f = D(f)$ ; $s$ s'identifie à un élément de $A(X)$-module $M$ et $s|_{X_f}$ à l'élément correspondant de $M_f$, et le résultat est trivial.

(ii) De nouveau $X$ est réunion finie d'ouverts affines $U_i$ ($1 \leq i \leq n$) tels que $L|_{U_i} = \mathcal{O}_X|_{U_i}$, et $s \otimes f^n$ s'identifie sur chaque $U_i$ à $s \cdot f^n$. On sait alors (I, 1, 4, th. 2) qu'il existe $n$ tel que pour chaque $i$, $s|_{U_i \cap X_f}$ se prolonge en une section $s_i$ de $F \otimes L^{\otimes n}$ au-dessus de $U_i$. Si $s_{ij}$ est la restriction de $s_i - s_j$ à $U_i \cap U_j$, on a $s_{ij} = 0$ dans $X_f \cap U_i \cap U_j$. Or, si $X$ est un espace noethérien, $U_i \cap U_j$ est quasi-compact ; si $X$ est un schéma, $U_i \cap U_j$ est affine, donc encore quasi-compact. En vertu de (i), il existe donc un entier $m$ (qu'on peut supposer indépendant de $i$ et $j$) tel que $(s_i|_{U_j} - s_j|_{U_i}) \otimes f^m = 0$. On conclut aussitôt qu'il existe une section $s'$ de $F \otimes L^{\otimes (n+m)}$ au-dessus de $X$, induisant $s \otimes f^{\otimes (n+m)}$ sur $X_f$.

Les corollaires qui suivent donnent une interprétation du th. 2 en langage plus algébrique :

Corollaire 1. — Les hypothèses étant celles du th. 2, considérons l'anneau gradué $A_\bullet = \Gamma_\bullet(L)$ et le $A_\bullet$-module gradué $M_\bullet = \Gamma_\bullet(L, F)$. Si $f \in A_n$, où $n \in \mathbb{Z}$, alors on a $\Gamma(X_f, F) \cong (M_\bullet)_f$ (sous-groupe du module de fractions $(M_\bullet)_f$ formé des éléments de degré 0).

\newpage

%% ===== Page 207 =====

-- II-18 --

Corollaire 2. - On suppose vérifiées les hypothèses du th. 2 et en outre on suppose que $L = \mathcal{O}_X$. Alors si on pose $A = \Gamma(X, \mathcal{O}_X)$, $M = \Gamma(X, F)$, le $A$-module $\Gamma(X_f, F)$ est isomorphe à $M_f$.

Proposition 15. - Soient $X$ un préschéma noethérien, $F$ un faisceau cohérent sur $X$, $J$ un faisceau quasi-cohérent d'idéaux de $\mathcal{O}_X$, tels que $J(x) = (0)$ pour tout $x$ appartenant au support de $F$. Alors il existe un entier $n$ tel que $J^n F = 0$.

Comme $X$ est réunion finie d'ouverts affines dont les anneaux sont noethériens, on peut supposer $X$ affine d'anneau $A$ noethérien ; alors $F = \tilde{M}$, où $M$ est un $A$-module de type fini, et $J = \tilde{I}$, où $I$ est un idéal de $A$ (I, 1.3, th. 1 et 4, th. 2). Comme $A$ est noethérien, $I$ admet un système fini de générateurs $f_i$ ($1 \le i \le m$) ; par hypothèse, toute section de $F$ au-dessus de $X$ est nulle dans chacun des $D(f_i)$ ; si $s_j$ ($1 \le j \le q$) sont des sections de $F$ engendrant $M$, il existe donc un entier $h > 0$ (indépendant de $i$ et $j$) tel que $f_i^h s_j = 0$ (I, 1.4, th. 2), et par suite $I^h s = 0$ pour tout $s \in M$. On en conclut que si $n = mh$, on a $I^n = 0$, et par suite le faisceau quasi-cohérent correspondant $J^n F = (I^n)^{\sim}$ est nul.

5. Le lemme de dévissage.

Définition 2. - Soit $\mathcal{C}$ une catégorie abélienne. On dit qu'une sous-classe $\mathcal{C}'$ de $\mathcal{C}$ est exacte si, pour toute suite exacte $0 \to A' \to A \to A'' \to 0$ dans $\mathcal{C}$ dont deux termes sont dans $\mathcal{C}'$, le troisième est dans $\mathcal{C}'$.

Théorème 3. - Soit $X$ un préschéma noethérien ; on désigne par $\mathcal{K}$ la catégorie abélienne des faisceaux cohérents sur $X$. Soient $\mathcal{K}'$ une sous-classe exacte de $\mathcal{K}$, $X'$ une partie fermée de l'espace de base de $X$. On suppose que pour toute partie fermée irréductible $Y$ de $X'$ de point générique $y$, il existe un faisceau $G \in \mathcal{K}'$ de support contenu dans $Y$ et tel que $G_y$ soit un module simple sur l'anneau local $\mathcal{O}_y = \mathcal{O}_{X, y}$. Alors tout faisceau cohérent de support contenu dans $X'$ appartient à $\mathcal{K}'$ (et en particulier, si $X'=X$, $\mathcal{K}' = \mathcal{K}$).

\newpage

%% ===== Page 208 =====

\noindent -- II-19 --

Soit $E$ l'ensemble des parties fermées $Y$ de $X'$ ayant la propriété suivante : tout faisceau cohérent sur $X$ de support contenu dans $Y$ appartient à $\underline{K}'$ ; nous allons montrer que $E$ est l'ensemble de toutes les parties fermées de $X'$, ou encore que l'ensemble $E'$ des parties fermées de $X'$ n'appartenant pas à $E$ est vide. Comme $X'$ est un espace noethérien, $E'$ admet un élément minimal, et on est donc ramené à prouver la proposition suivante : si $Y$ est une partie fermée de $X'$ telle que toute partie fermée de $Y$ distincte de $Y$, appartient à $E$, alors $Y$ appartient à $E$.

Soit donc $F \in \underline{K}$ à support contenu dans $Y$, et prouvons que $F \in \underline{K}'$. Soit $J$ le plus grand faisceau d'idéaux quasi-cohérent de $\mathcal{O}_X$ tel que $Y$ soit le support de $\mathcal{O}_X/J$ (I, 3, 4, prop. 9) ; la restriction de $\mathcal{O}_X/J$ à $Y$ définit le plus petit sous-préschéma de $X$ ayant pour support $Y$, et on sait que ce sous-préschéma est réduit ; nous le désignerons par $Y_{\text{red}}$. D'après la prop. 15 du n$^\circ$ 4, il existe un entier $n$ tel que $J^n F = 0$ ; on a donc pour $1 \le k \le n$, une suite exacte
$$0 \to J^{k-1}F/J^k F \to F/J^k F \to F/J^{k-1}F \to 0$$
de faisceaux cohérents sur $X$ ; comme la catégorie $\underline{K}'$ est exacte, on voit, par récurrence sur $k$, qu'il suffit de prouver que chacun des faisceaux $F_k = J^{k-1}F/J^k F$ est dans $\underline{K}'$. On est donc ramené à prouver que $F \in \underline{K}'$ sous l'hypothèse supplémentaire que $JF = 0$ ; il revient au même de dire que $F|Y$ est un faisceau cohérent sur le préschéma $Y_{\text{red}}$ (FGA, I, 2, 17). Distinguons deux cas :

a) $Y$ est réductible. Soit $Y = Y' \cup Y''$, $Y'$ et $Y''$ étant des parties fermées de $Y$ distinctes de $Y$ ; désignons par $J'$ (resp. $J''$) le plus grand faisceau d'idéaux quasi-cohérent de $\mathcal{O}_X$ tel que $Y'$ (resp. $Y''$) soit le support de $\mathcal{O}_X/J'$ (resp. $\mathcal{O}_X/J''$), et posons $F' = F \otimes (\mathcal{O}_X/J')$, $F'' = F \otimes (\mathcal{O}_X/J'')$. On définit un homomorphisme $F \to F'$ (resp. $F \to F''$) en faisant correspondre à toute section $s \in F(U)$ la section $s \otimes 1_U$ (resp. $s \otimes 1''_U$), $1'_U$ (resp. $1''_U$) étant la section de $\mathcal{O}_X/J'$ (resp. $\mathcal{O}_X/J''$)

\newpage

%% ===== Page 209 =====

\marginpar{II-20}
\noindent égale à l'élément unité de la fibre correspondant en tout point $x \in U$ ; d'où un homomorphisme $F \to F' \oplus F''$. La suite
$$(2) \quad 0 \to F \to F' \oplus F''$$
est exacte. En effet, par définition de $J$, on a $J' \cap J''=J$, donc $O_X/J$ s'identifie à un sous-faisceau de $(O_X/J') \oplus (O_X/J'')$ ; si en un point $z \in Y$ un germe de section $\tilde{s}(z)$ de $F(z)$ est tel que $\tilde{s}(z) \otimes 1'_z = \tilde{s}(z) \otimes 1''_z = 0$, on en déduit donc $\tilde{s}(z) = 0$, d'où notre assertion. En outre, si $z \notin Y''$, on a $J'(z)=J(z)$, donc $\tilde{s}(z) \otimes 1'_z = \tilde{s}(z)$ ; autrement dit $F'(z)=F(z)$, et comme $F''(z)=0$, $F(z) \to F'(z) \oplus F''(z)$ est surjectif. Il en est de même si $z \notin Y'$ ; d'autres termes, le conoyau de l'homomorphisme $F \to F' \oplus F''$ a son support dans $Y' \cap Y''$. Or par hypothèse $F'$ et $F''$ sont dans $\underline{K}'$, donc il en est de même de $F' \oplus F''$ puisque $\underline{K}'$ est exacte ; l'hypothèse entraînant aussi que le conoyau de $F \to F' \oplus F''$ est dans $\underline{K}'$, on en conclut que $F \in \underline{K}'$, $\underline{K}'$ étant exacte.

b) $Y$ est irréductible ; soit $y$ un point générique, et soit $V$ un schéma affine de $X$ contenant $y$ ; $Y_{\text{red}} \cap V$ est irréductible et réduit, dont l'anneau est $A(V)/J(Y \cap V)$, et on a $J(Y \cap V) = j_y$ dans $A(V)$, donc $J(y)$ est l'idéal maximal $j_y O_X(y)$ de $O_X(y)$, et $O_{Y_{\text{red}}}(y) = O_X(y)/J(y)$ est le corps résiduel $\kappa(y)$. Comme $F|Y$ est un faisceau cohérent sur $Y_{\text{red}}$, $F_y$ est un espace vectoriel de dimension finie sur $\kappa(y)$. Cela étant, il y a par hypothèse un faisceau $G \in \underline{K}'$ dont le support est contenu dans $Y$ et tel que $G_y$ soit un $O_X(y)$-module simple, ou, ce qui revient au même, un $\kappa(y)$-espace vectoriel de dimension $1$. Il y a par suite un $\kappa(y)$-isomorphisme $(G_y)^n \to F_y$, et comme $G^n$ et $F$ sont des faisceaux cohérents, il existe un voisinage ouvert $W$ de $y$ dans $X$ et un isomorphisme $G^n|W \to F|W$ (\text{FAC}, I, 2, 14, prop. 5). Soit $H$ le graphe de cet isomorphisme, identifié à un sous-faisceau de $(G^n \times F)|W$, et qui est un faisceau cohérent sur $W$, isomorphe à $G^n|W$ et à $F|W$ ; il existe par suite un sous-faisceau cohérent $H_0$ de $G^n \times F$, induisant $H$ sur $W$ et $0$ sur $Y$ (no. 4, ...

\newpage

%% ===== Page 210 =====

- II-21 -

1) . Les restrictions $u: H_0 \to \mathbb{G}^m$ et $v: H_0 \to F$ à $H_0$ des projections de $\mathbb{G}^m \times F$ sont alors des homomorphismes de faisceaux cohérents, qui, dans $W$ et dans $Y$, se réduisent à des isomorphismes ; autrement dit, les noyaux et conoyaux de $u$ et $v$ ont leur support dans l'ensemble fermé $Y - Y \cap W$, distinct de $Y$, et appartiennent par suite à $\underline{K}'$ ; d'autre part $\mathbb{G}^m \in \underline{K}'$ et par suite aussi $\mathbb{G}^m \in \underline{K}'$ puisque $\underline{K}'$ est exacte. On en conclut tout d'abord que $H_0 \in \underline{K}'$, puis que $F \in \underline{K}'$, puisque $\underline{K}'$ est exacte. C.Q.F.D.

Corollaire . — Supposons que la sous-classe exacte $\underline{K}'$ de $\underline{K}$ ait en outre la propriété que tout facteur direct d'un faisceau cohérent $M \in \underline{K}'$ appartient encore à $\underline{K}'$. Dans ces conditions, la conclusion du th. 3 subsiste encore lorsque la condition "$G_y$ est un module simple sur $\mathcal{O}_y$" est remplacée par "$G_y \neq 0$".

En effet, le raisonnement ne doit être modifié que dans le cas b) : $G_y$ est alors un espace vectoriel sur $\kappa(y)$, de dimension $q > 0$, et on a donc cette fois un isomorphisme $(G_y)^m \to (F_y)^q$ ; la fin du raisonnement du th. 3 prouve alors que $F^q \in \underline{K}'$, et l'hypothèse additionnelle sur $\underline{K}'$ entraîne alors $F \in \underline{K}'$.

\newpage

%% ===== Page 211 =====

\marginpar{Archives Grothendieck sept 59}

\setcounter{page}{22}
\section*{\textsuperscript{-- II-22 --} § 2. Morphismes affines.}

Les résultats de ce paragraphe sont les contreparties "globales" de ceux du chap. I, § 1 ; ils ne sont donc pas essentiellement nouveaux et fournissent simplement un langage commode pour la suite.

1. Compléments sur les images directes de faisceaux de modules.

Soient $X$ un $S$-préschéma, $f : X \to S$ son morphisme structural ; on sait (chap. 0, § 2, n\textsuperscript{o}5) que l'image directe $f_*(\mathcal{O}_X)$ est une $\mathcal{O}_S$-algèbre, que nous noterons $\underline{A}(X)$ ; de même, pour tout faisceau algébrique $\mathcal{F}$ sur $X$, nous écrirons $\underline{A}(\mathcal{F})$ l'image directe $f_*(\mathcal{F})$, qui est un $\underline{A}(X)$-module (et non seulement un $\mathcal{O}_S$-module).

Soient $Y$ un second $S$-préschéma, $g : Y \to S$ son morphisme structural, $u : X \to Y$ un $S$-morphisme ; on a donc le diagramme commutatif
\begin{center}
\begin{tikzcd}
X \arrow[dr, "f"'] \arrow[r, "u"] & Y \arrow[d, "g"] \\
& S
\end{tikzcd}
\end{center}
(1)

On a par définition $u = (\varphi, \vartheta)$, où $\vartheta : \mathcal{O}_Y \to u_*(\mathcal{O}_X)$ est un homomorphisme de faisceaux d'anneaux : on en déduit (chap. 0, § 2, n\textsuperscript{o}5) un homomorphisme de $\mathcal{O}_S$-algèbres $g_*(\vartheta) : g_*(\mathcal{O}_Y) \to g_*(u_*(\mathcal{O}_X))$, et comme $g_*(u_*(\mathcal{O}_X)) = f_*(\mathcal{O}_X)$ d'après (1), on obtient ainsi un homomorphisme $\underline{A}(u) : \underline{A}(Y) \to \underline{A}(X)$ de $\mathcal{O}_S$-algèbres.

Soient maintenant $\mathcal{F}$ un faisceau algébrique sur $X$, $\mathcal{G}$ un faisceau algébrique sur $Y$, et $v : \mathcal{G} \to u_*(\mathcal{F})$ un homomorphisme de $\mathcal{O}_Y$-modules : $g_*(v) : g_*(\mathcal{G}) \to g_*(u_*(\mathcal{F})) = f_*(\mathcal{F})$ est alors non seulement un homomorphisme de $\mathcal{O}_S$-modules, que nous noterons $\underline{A}(v)$, mais encore le couple $(\underline{A}(u), \underline{A}(v))$ constitue un bi-homomorphisme de $\underline{A}(Y)$-module $\underline{A}(\mathcal{G})$ dans le $\underline{A}(X)$-module $\underline{A}(\mathcal{F})$.

Le préschéma $S$ étant fixé, on peut considérer les couples $(X, \mathcal{F})$, où $X$ est un $S$-préschéma et $\mathcal{F}$ un faisceau algébrique sur $X$, comme formant une catégorie, en définissant un morphisme $(X, \mathcal{F}) \to (Y, \mathcal{G})$ comme un couple $(u, v)$, où $u : X \to Y$ un $S$-morphisme et $v : \mathcal{G} \to u_*(\mathcal{F})$ un $\mathcal{O}_Y$-homomorphisme. On peut alors dire que $(\underline{A}(X), \underline{A}(\mathcal{F}))$ est un fonc-

\newpage

%% ===== Page 212 =====

teur contravariant à valeurs dans la catégorie dont les objets sont
les couples formés d'une $\mathcal{O}_S$-algèbre et d'un module sur cette algèbre
et les morphismes les di-homomorphismes.

Dans ce qui suit, nous dirons pour abréger que le $S$-préschéma $X$ est
adéquat lorsque pour tout faisceau quasi-cohérent $\mathcal{F}$ sur $X$, son image
directe $\underline{A}(\mathcal{F})$ est un faisceau quasi-cohérent sur $S$; en particulier
alors $\underline{A}(X)$ est un faisceau quasi-cohérent sur $S$. On a vu au §1, n°1, cor.
de la prop.1, des conditions suffisantes pour que $X$ soit adéquat.
Proposition 1. - Soient $Y$ un préschéma, $\mathcal{B}$ une $\mathcal{O}_Y$-algèbre qui
est un faisceau quasi-cohérent de $\mathcal{O}_Y$-modules; $\mathcal{B}$ définit sur $Y$ une
structure d'espace annelé $(Y, \mathcal{B})$, et tout $\mathcal{B}$-module est muni d'une
structure de $\mathcal{O}_Y$-module. Pour qu'un $\mathcal{B}$-module $\mathcal{F}$ soit un faisceau
quasi-cohérent sur l'espace annelé $(Y, \mathcal{B})$, il faut et il suffit que
$\mathcal{F}$ soit quasi-cohérent sur le préschéma $(Y, \mathcal{O}_Y)$.

La condition est nécessaire : en effet, si $\mathcal{F}$ est quasi-cohérent sur
$(Y, \mathcal{B})$, c'est le conoyau d'un homomorphisme $\mathcal{B}^{(I)} \to \mathcal{B}^{(J)}$ ; comme cet
homomorphisme est aussi un $\mathcal{O}_Y$-homomorphisme de $\mathcal{O}_Y$-modules, et que
$\mathcal{B}^{(I)}$ et $\mathcal{B}^{(J)}$ sont quasi-cohérents sur $(Y, \mathcal{O}_Y)$, puisqu'il en est ainsi
de $\mathcal{B}$ (I, 1, 3, cor. 4 du th. 1), $\mathcal{F}$ est quasi-cohérent sur $(Y, \mathcal{O}_Y)$ (I, 1, 3,
cor. 4 du th. 1).

La condition est suffisante ; on peut en effet se borner au cas où
$Y$ est un schéma affine d'anneau $A$ ; alors on a $\mathcal{B} = \tilde{B}$, où $B$ est une $A$-
algèbre. En effet, la donnée d'une structure de $\mathcal{O}_Y$-algèbre sur un
$\mathcal{O}_Y$-module revient à la donnée d'un $\mathcal{O}_Y$-homomorphisme $\mathcal{B} \otimes_{\mathcal{O}_Y} \mathcal{B} \to \mathcal{B}$ satis-
faisant aux conditions d'associativité et de commutativité qui se
traduisent par la commutativité des diagrammes
$$
\begin{tikzcd}
\mathcal{B} \otimes \mathcal{B} \otimes \mathcal{B} \arrow[r, "1 \otimes f"] \arrow[d, "f \otimes 1"'] & \mathcal{B} \otimes \mathcal{B} \arrow[d, "f"] \\
\mathcal{B} \otimes \mathcal{B} \arrow[r, "f"] & \mathcal{B}
\end{tikzcd}
\qquad
\begin{tikzcd}
\mathcal{B} \otimes \mathcal{B} \arrow[r, "f"] \arrow[d, "\sigma"'] & \mathcal{B} \\
\mathcal{B} \otimes \mathcal{B} \arrow[ur, "f"'] & 
\end{tikzcd}
$$
où $\sigma$ est la symétrie $b \otimes b' \to b' \otimes b$ en chaque point. Notre assertion
résulte alors des propriétés fonctorielles de $M \to \tilde{M}$ (I, 1, 3).

\newpage

%% ===== Page 213 =====

- 11-24 -

De même, on a $F = \tilde{M}$, où $M$ est un $A$-module, et la donnée d'une structure de $B$-Module sur $F$ revient à la donnée d'un $\mathcal{O}_Y$-homomorphisme :
$$ B \otimes_{\mathcal{O}_Y} \tilde{M} \to \tilde{M} $$
avec commutativité du diagramme
\begin{tikzcd}
B \otimes_{\mathcal{O}_Y} (B \otimes_A M) \arrow{r}{1 \otimes g} \arrow{d}{f \otimes 1} & B \otimes_A M \arrow{d}{g} \\
B \otimes_A M \arrow{r}{g} & M
\end{tikzcd}
Donc $M$ est muni d'une structure de $B$-module ; on en conclut que $M$ (en tant que $B$-module) est isomorphe au conoyau d'un $B$-homomorphisme $B^{(I)} \to B^{(J)}$ ; en tant que $A$-module, il est aussi isomorphe au conoyau de cet homomorphisme, donc $F$ est isomorphe au conoyau de l'homomorphisme $B^{(I)} \to B^{(J)}$, non seulement en tant que $\mathcal{O}_Y$-module, mais aussi en tant que $B$-Module ; d'où la proposition.

Proposition 2. - Soient $Y$ un préschéma localement noethérien, $B$ un faisceau quasi-cohérent d'$\mathcal{O}_Y$-algèbres. On suppose que pour tout $y \in Y$, il existe un voisinage $V$ de $y$ dans $Y$ tel que $\Gamma(V, B)$ soit une algèbre de type fini sur $\Gamma(V, \mathcal{O}_Y)$. Alors $B$ est un faisceau cohérent d'anneaux (FAC, I, 2, 15).

On peut de nouveau se limiter au cas où $Y$ est un schéma affine d'anneau noethérien $A$ ; on a vu dans la démonstration de la prop. 1 que $B = \tilde{B}$, où $B$ est une $A$-algèbre. L'hypothèse permet en outre de supposer que $B$ est une $A$-algèbre de type fini, et par suite est un anneau noethérien. Cela étant, il faut prouver que le conoyau d'un $B$-homomorphisme $B^I \to B$ est un $B$-module de type fini ; ce noyau est isomorphe à $\tilde{N}$, où $N$ est le noyau de l'homomorphisme correspondant $B^I \to B$ ; mais comme cet homomorphisme n'est seulement un $A$-homomorphisme, mais aussi un $B$-homomorphisme, $N$ est un sous-module de $B^I$, donc un $B$-module de type fini, ce qui achève la démonstration.

Corollaire. - Sous les hypothèses de la prop. 2, pour qu'un $B$-module $N$ soit cohérent, il faut et il suffit qu'il soit un $\mathcal{O}_Y$-module quasi-cohérent et un $B$-module de type fini. S'il en est ainsi, et si $N$ est un sous-$B$-Module ou un $B$-Module quotient de $M$, pour que $N$ soit un $B$-module cohérent, il faut et il suffit qu'il soit un...

\newpage

%% ===== Page 214 =====

- II-25 -

Les conditions sur $M$ sont évidemment nécessaires, compte tenu de la prop.1 ; montrons qu'elles sont suffisantes. On peut se limiter au cas où $Y$ est affine d'anneau noethérien $A$ et où il existe un $B$-homomorphisme $B^n \to M \to 0$. On a alors $B = \tilde{B}$, $M = \tilde{M}$, où $B$ est une $A$-algèbre de type fini et $M$ un $B$-module de type fini ; le noyau $P$ de l'homomorphisme $B^n \to M$ est un sous-module de $B^n$, donc est de type fini. On en conclut que $M$ est le conoyau d'un $B$-homomorphisme $B^m \to B^n$, et est donc cohérent puisque $B$ est un faisceau cohérent d'anneaux (FAC, I, 2, 15 prop. 7). Le même raisonnement montre qu'un sous-$B$-module de $M$ est de type fini, d'où la seconde partie du corollaire.

2. \emph{Préschémas affines sur un préschéma}.

Définition 1. - On dit qu'un $S$-préschéma $X$ est affine au-dessus de $S$ s'il existe un recouvrement $(S_\alpha)$ de $S$ par des ouverts affines tels que, pour tout $\alpha$, le préschéma induit par $X$ sur l'ensemble ouvert $f^{-1}(S_\alpha)$ ($f$ morphisme structural de $X$) soit affine.

Exemple. - Tout sous-préschéma fermé de $S$ est un $S$-préschéma affine au-dessus de $S$ (I, 3, 2, cor. 2 de la prop. 4).

Remarque. - Un préschéma $X$ affine au-dessus de $S$ n'est pas nécessairement un schéma affine, comme le montre l'exemple $X=S$. D'autre part, si un schéma $X$ est un $S$-préschéma, il n'est pas nécessairement affine au-dessus de $S$. Rappelons toutefois que si $S$ est un schéma, un $S$-préschéma qui est un schéma affine est affine au-dessus de $S$ (I, 3, 7, cor. 3 de la prop. 22).

Proposition 3. - Tout $S$-préschéma qui est affine au-dessus de $S$ est séparé au-dessus de $S$.

Cela résulte aussitôt de (I, 3, 6, cor. 2 de la prop. 20) et de (I, 3, 6, prop. 19).

Proposition 4. - Soient $X$ un $S$-préschéma affine au-dessus de $S$, $f$ le morphisme structural $X \to S$. Pour tout ouvert $U \subset S$, $f^{-1}(U)$ est affine.

\newpage

%% ===== Page 215 =====

- II-26 -

ne au-dessus de $U$.

Par définition, on est ramené au cas où $S$ est affine et $X$ affine ; alors $f$ est de la forme $(\varphi, \tilde{\varphi})$, où $\varphi$ est un homomorphisme $A \to B$, en désignant par $A$ et $B$ les anneaux de $S$ et $X$ respectivement. Comme les $D(g)$, où $g \in A$, forment une base de $S$, on est ramené au cas où $U = D(g)$ ; mais on sait alors que $f^{-1}(U) = D(\varphi(g))$ $(I, 1, 2, \text{prop.} 5)$, d'où la proposition.

Proposition 5 . - Un $S$-préschéma $X$ qui est affine au-dessus de $S$ est $S$-adéquat ($n^\circ 1$).

Cela résulte de la définition et du cor. de la prop. 1 du $\S 1, n^\circ 1$, compte tenu de la prop. 3.

En particulier, le faisceau d'algèbres $\mathcal{A}(X)$ ($n^\circ 1$) est alors quasi-cohérent sur $S$.

Proposition 6 . - Soient $X, Y$ deux $S$-préschémas ; on suppose $X$ affine au-dessus de $S$ et $Y$ $S$-adéquat ($n^\circ 1$). Alors l'application $u \to A(u)$ de l'ensemble $\operatorname{Hom}_S(Y, X)$ dans l'ensemble $\operatorname{Hom}(\mathcal{A}(X), \mathcal{A}(Y))$ ($n^\circ 1$) est bijective.

Soit $(s_\alpha)$ un recouvrement de $S$ par des ouverts affines tels que $f^{-1}(s_\alpha)$ soit affine ($f$ morphisme structural $X \to S$). Si $g$ est le morphisme structural $Y \to S$, pour tout $S$-morphisme $u : Y \to X$, la restriction $u_\alpha$ de $u$ à $g^{-1}(s_\alpha)$ est un $S$-morphisme $g^{-1}(s_\alpha) \to f^{-1}(s_\alpha)$. Inversement, si pour tout $\alpha$, $v_\alpha$ est un $S$-morphisme $g^{-1}(s_\alpha) \to f^{-1}(s_\alpha)$, tel que les restrictions de $v_\alpha$ et $v_\beta$ à $g^{-1}(s_\alpha \cap s_\beta)$ coïncident, les $v_\alpha$ sont les restrictions aux $s_\alpha$ d'un $S$-morphisme $Y \to X$. Cela étant, il y a correspondance biunivoque entre les $S$-morphismes $g^{-1}(s_\alpha) \to f^{-1}(s_\alpha)$ et les homomorphismes de $\Gamma(s_\alpha, \mathcal{O}_S)$-algèbres $\Gamma(f^{-1}(s_\alpha), \mathcal{O}_X) \to \Gamma(g^{-1}(s_\alpha), \mathcal{O}_Y)$ $(I, 2, 2, \text{prop.} 3)$. Par définition, on a $\Gamma(f^{-1}(s_\alpha), \mathcal{O}_X) = \Gamma(s_\alpha, \mathcal{A}(X))$ et $\Gamma(g^{-1}(s_\alpha), \mathcal{O}_Y) = \Gamma(s_\alpha, \mathcal{A}(Y))$ ; comme $\mathcal{A}(X)$ et $\mathcal{A}(Y)$ sont des $\mathcal{O}_S$-modules quasi-cohérents (en vertu de l'hypothèse et de la prop. 5), à tout $\Gamma(s_\alpha, \mathcal{O}_S)$ homomor-

\newpage

%% ===== Page 216 =====

- II-27 -

phisme $\Gamma(S_\alpha, \mathcal{A}(X)) \to \Gamma(S_\alpha, \mathcal{A}(Y))$ correspond canoniquement un homomorphisme de $\mathcal{O}_S$-module et réciproquement (I, 1, 3, cor. 2 du th. 1); en outre, comme la propriété d'être un homomorphisme d'anneaux se traduit par la commutativité du diagramme

$$\begin{tikzcd}
\Gamma(S_\alpha, \mathcal{A}(X)) \otimes_{\Gamma(S_\alpha, \mathcal{O}_S)} \Gamma(S_\alpha, \mathcal{A}(X)) \arrow[r] \arrow[d] & \Gamma(S_\alpha, \mathcal{A}(Y)) \otimes_{\Gamma(S_\alpha, \mathcal{O}_S)} \Gamma(S_\alpha, \mathcal{A}(Y)) \arrow[d] \\
\Gamma(S_\alpha, \mathcal{A}(X)) \arrow[r] & \Gamma(S_\alpha, \mathcal{A}(Y))
\end{tikzcd}$$

les homomorphismes correspondants de $\mathcal{O}_S$-modules rendent le diagramme correspondant commutatif (I, 1, 2, cor. 6 du th. 1) et par suite sont des homomorphismes de $\mathcal{O}_S$-algèbres. Il y a donc correspondance biunivoque entre les $S$-morphismes $Y \to X$ et les familles $(w_\alpha)$ d'homomorphismes de $\mathcal{O}_S$-algèbres $\mathcal{A}(X)|S_\alpha \to \mathcal{A}(Y)|S_\alpha$ telles que $w_\alpha$ et $w_\beta$ coïncident dans $S_\alpha \cap S_\beta$; mais ces familles correspondent biunivoquement aux homomorphismes de $S$-algèbres $\mathcal{A}(X) \to \mathcal{A}(Y)$, ce qui achève la démonstration.

Corollaire. — Soient $X, Y$ deux $S$-préschémas qui sont affines au-dessus de $S$. Alors l'application $u \to \mathcal{A}(u)$ de $\operatorname{Hom}_S(Y, X)$ dans $\operatorname{Hom}(\mathcal{A}(X), \mathcal{A}(Y))$ est bijective, et pour que $u$ soit un isomorphisme, il faut et il suffit que $\mathcal{A}(u)$ le soit.

Cela résulte des prop. 5 et 6.

3. Préschéma affine au-dessus de $S$ associé à une $\mathcal{O}_S$-algèbre.

Proposition 7. — Pour toute $\mathcal{O}_S$-algèbre quasi-cohérente $\mathcal{B}$, il existe un préschéma $X$ affine au-dessus de $S$, et un seul, à un $S$-isomorphisme près, tel que $\mathcal{A}(X) = \mathcal{B}$.

Pour tout ouvert affine $U \subset S$, soit $X_U$ le préschéma $\operatorname{Spec}(\Gamma(U, \mathcal{B}))$; comme $\Gamma(U, \mathcal{B})$ est une $\Gamma(U, \mathcal{O}_S)$-algèbre, $X_U$ est un $S$-préschéma (I, 2, 3); comme $\mathcal{B}$ est quasi-cohérent, la $\mathcal{O}_S$-algèbre $\mathcal{A}(X_U)$ s'identifie à

\newpage

%% ===== Page 217 =====

- II-28 -
$B|_U$ (I, 1, 3, th. 1 et 1, 5). Soit $X_{U,V}$ le préschéma induit par $X_U$ sur l'image réciproque de $U \cap V$ par le morphisme structural $f_U : X_U \to S$; $X_{U,V}$ et $X_{V,U}$ sont affines sur $U \cap V$ (n° 2, prop. 4), et par définition, $\mathcal{A}(X_{U,V})$ et $\mathcal{A}(X_{V,U})$ s'identifient canoniquement à $B|_{U \cap V}$. Donc (n° 2, prop. 6), il y a un isomorphisme canonique $\theta_{U,V} : X_{V,U} \to X_{U,V}$; en outre, si $W$ est un troisième ouvert affine de $S$, et si on désigne par $\theta_{U,V}, \theta_{V,W}$ et $\theta_{U,W}$ les restrictions de $\theta_{U,V}, \theta_{V,W}$ et $\theta_{U,W}$ aux images réciproques de $U \cap V \cap W$ dans $X_U, X_V$ et $X_W$ respectivement par les morphismes structuraux, on a $\theta_{V,W} \circ \theta_{U,V} = \theta_{U,W}$. Il existe donc un préschéma $X$, un recouvrement $(T_U)$ de $X$ par des ouverts affines, et pour chaque $U$ un isomorphisme $\varphi_U$ de $X_U$ sur le préschéma induit $T_U$, de sorte que $\varphi_U$ applique $f_U^{-1}(U \cap V)$ sur $T_U \cap T_V$ et que l'on ait $\theta_{U,V} = \varphi_U^{-1} \circ \varphi_V$ (I, 3, 8). Le morphisme $f_U \circ \varphi_U^{-1}$ fait de $T_U$ un $S$-préschéma, et les morphismes $f_U$ et $f_V$ coïncident dans $T_U \cap T_V$; donc $X$ est un $S$-préschéma. Il est clair que $X$ est affine au-dessus de $S$ et que $\mathcal{A}(X) = B$, donc $\mathcal{A}(X) = B$. L'unicité à isomorphisme près résulte aussitôt de la prop. 6 du n° 2.

On dit que le préschéma $X$ ainsi défini est associé à la $\mathcal{O}_S$-algèbre $B$.

Corollaire 1. — Soit $X$ un préschéma affine au-dessus de $S$, et $f : X \to S$ le morphisme structural. Pour tout ouvert affine $U \subset S$, le préschéma induit sur $f^{-1}(U)$ est un schéma affine.

Comme on peut supposer $X$ associé à une $\mathcal{O}_S$-algèbre, le corollaire résulte de la construction de $X$ décrite dans la démonstration précédente.

Corollaire 2. — Si $S$ est un schéma affine, tout préschéma affine au-dessus de $S$ est un schéma affine.

Corollaire 3. — Soient $X$ un préschéma affine au-dessus de $S$, $X'$ ...

\newpage

%% ===== Page 218 =====

- 29 -

un préschéma affine au-dessus de $X$ ; alors $X'$ est affine au-dessus de $S$.

On est aussitôt ramené à démontrer que si on suppose en outre que $S$ est un schéma affine, alors $X'$ est un schéma affine ; mais cela résulte du cor. 2.

Soit $X$ un préschéma affine au-dessus de $S$. Pour définir un préschéma $X'$ affine au-dessus de $X$, il revient au même, d'après le cor. 3 de la prop. 7, de se donner un préschéma $X'$ affine au-dessus de $S$, et un $S$-morphisme $f': X' \to X$ ; en d'autres termes (prop. 7), cela revient à se donner une $\mathcal{O}_S$-algèbre quasi-cohérente $\mathcal{B}' = \mathcal{A}(X')$ et un homomorphisme $\mathcal{A}(X) \to \mathcal{B}'$ de $\mathcal{O}_S$-algèbres (que l'on peut envisager aussi comme définissant sur $\mathcal{B}'$ une structure de $\mathcal{B}$-algèbre). D'ailleurs, si $f$ est le morphisme structural $X \to S$, on a $\mathcal{B}' = f_*(\mathcal{O}_{X'})$, puisque $(f \circ g)_* = f_* \circ g_*$.

4. \emph{Faisceaux quasi-cohérents sur un préschéma affine au-dessus de $S$.}

Nous laissons au lecteur les démonstrations des propositions suivantes, qui sont tout à fait analogues à celles des prop. 6 et 7 respectivement :

Proposition 8. - Soient $X, Y$ deux $S$-préschémas ; on suppose $X$ affine au-dessus de $S$, et $Y$ $S$-adéquat ($n^\circ 1$). Soit $\mathcal{F}$ (resp. $\mathcal{G}$) un $\mathcal{O}_Y$-module (resp. un $\mathcal{O}_X$-module) quasi-cohérent. Alors l'application $(u, \alpha) \to (\mathcal{A}(u), \mathcal{A}(\alpha))$ de l'ensemble des morphismes $(Y, \mathcal{G}) \to (X, \mathcal{F})$ dans l'ensemble des $\mathcal{A}$-homomorphismes $\mathcal{A}(\mathcal{F}) \to \mathcal{A}(\mathcal{G})$ ($n^\circ 1$) est bijective.

Proposition 9. - Pour tout couple $(\mathcal{B}, \mathcal{M})$ formé d'une $\mathcal{O}_S$-algèbre quasi-cohérente et d'un $\mathcal{B}$-module $\mathcal{M}$ quasi-cohérent (en tant que $\mathcal{O}_S$-module ou en tant que $\mathcal{B}$-module, ce qui est équivalent $n^\circ 1$, prop. 11), il existe un couple $(X, \mathcal{F})$ formé d'un préschéma $X$ affine au-dessus de $S$ et d'un $\mathcal{O}_X$-module quasi-cohérent $\mathcal{F}$, tels que $\mathcal{A}(X) = \mathcal{B}$ et $\mathcal{A}(\mathcal{F}) = \mathcal{M}$ ; en outre ce couple est unique à un isomorphisme près.

\newpage

%% ===== Page 219 =====

-- II-30 --

Lemme 1. — Soient $B$ une $\mathcal{O}_S$-algèbre quasi-coh\=erente, $M, N$ deux $B$-modules quasi-coh\=erents. Alors le $B$-module $M \otimes_B N$ est quasi-coh\=erent. Il en est de m\^eme du $B$-module $\operatorname{Hom}_B(M, N)$ lorsque pour tout $s \in S$, il existe un voisinage $U$ de $s$ et deux entiers $p, q$ tels que $M|_U$ soit isomorphe au conoyau d'un homomorphisme $(B|_U)^p \to (B|_U)^q$.

En vertu de la prop. 1 du n$^\circ$ 1, il revient au m\^eme de prouver que $M \otimes_B N$ et $\operatorname{Hom}_B(M, N)$ sont quasi-coh\=erents en tant que $\mathcal{O}_S$-modules. On peut donc supposer $S$ affine d'anneau $A$, avec $B = \tilde{B}$, $M = \tilde{M}$, $N = \tilde{N}$, o\`u $B$ est une $A$-algèbre, $M$ et $N$ des $B$-modules. Si $p$ est l'homomorphisme $A \to B$, on voit en raisonnant comme dans I, 1, 3, cor. 6 du th. 1, que tout revient \`a prouver que, pour tout \=el\=ement $f \in A$, si on pose $f' = \varphi(f)$, le $A_f$-module $(M \otimes_B N)_{[f']}$ est isomorphe \`a $(M_{f'} \otimes_{B_{f'}} N_{f'})$ et (lorsque $M$ est conoyau d'un homomorphisme $B^p \to B^q$) que $(\operatorname{Hom}_B(M, N))_{[f']}$ est isomorphe \`a $(\operatorname{Hom}_{B_{f'}}(M_{f'}, N_{f'}))$ ; mais cela a \=et\=e vu au chap. 0, \S 1, n$^\circ$ 5.7.

Proposition 10. — Soient $Y$ un pr\=esch\=ema affine au-dessus de $S$, $X, X'$ deux pr\=esch\=emas affines au-dessus de $Y$ (donc aussi au-dessus de $S$) (cor. 3 de la prop. 7). Soient $B, A, A'$ les $\mathcal{O}_S$-alg\=ebres correspondant \`a $Y, X$ et $X'$ respectivement. Alors $X \times_Y X'$ est affine au-dessus de $Y$ (donc de $S$), et la $\mathcal{O}_S$-alg\=ebre correspondante est $A \otimes_B A'$.

En vertu du lemme 1, $A \otimes_B A'$ est une $\mathcal{O}_S$-alg\=ebre quasi-coh\=erente, donc correspond \`a un pr\=esch\=ema $Z$ affine au-dessus de $S$ (n$^\circ$ 3, prop. 7) ; en outre, les $B$-homomorphismes $A \to A \otimes_B A'$ et $A' \to A \otimes_B A'$ correspondent \`a des $Y$-morphismes $Z \to X$ et $Z \to X'$. \marginpar{(prop 6)} Pour voir que le triplet form\=e de $Z$ et de ces deux morphismes est un produit $X \times_Y X'$, on peut se borner au cas o\`u $S$ est un sch\=ema affine d'anneau $C$. Mais alors $Y, X$ et $X'$ sont des sch\=emas affines (n$^\circ$ 3, cor. 2 de la prop. 7) dont les anneaux $B, A, A'$ sont des $C$-alg\=ebres telles que $B = \tilde{B}$, $A = \tilde{A}$, $A' = \tilde{A'}$. La d\=emonstration du lemme 1 montre alors que $A \otimes_B A' = (A \otimes_B A')^\sim$, donc l'anneau du sch\=ema affine $Z$

\newpage

%% ===== Page 220 =====

s'identifie à $A \otimes_B A'$, et les morphismes $Z \to X$, $Z \to X'$ sont les morphismes $u, v$ correspondant aux homomorphismes $u: A \to A \otimes_B A'$, $v: A' \to A \otimes_B A'$. La proposition résulte alors de I, 2, 4, prop. 6.

Corollaire. — Soit $F$ (resp. $F'$) un $\mathcal{O}_X$-module (resp. un $\mathcal{O}_{X'}$-module) quasi-cohérent ; on a $A(F \otimes_Y F') = A(F) \otimes_{A(Y)} A(F')$.

On sait que $F \otimes_Y F'$ est quasi-cohérent ($\S 1, n^\circ 1$); on se ramène aussitôt au cas où $S$ (et par suite $Y, X, X'$) sont des schémas affines, et où $F = \tilde{M}$, $F' = \tilde{M}'$, $M$ (resp. $M'$) étant un $A$-module (resp. un $A'$-module), avec les notations de la démonstration de la prop. 10. Alors $F \otimes_Y F'$ s'identifie au faisceau associé au $(A \otimes_B A')$-module $M \otimes_B M'$ ($\S 1, n^\circ 1$, prop. 7), d'où le corollaire ($X$ étant affine au-dessus de $S$).

En particulier, lorsque $X=X'=Y$, on voit que si $F, G$ sont des $\mathcal{O}_X$-modules quasi-cohérents, on a $A(F \otimes_X G) = A(F) \otimes_{A(X)} A(G)$. Si de plus pour tout $x \in X$, il y a un voisinage $V$ de $x$ et deux entiers $p,q$ tels que $F|_V$ soit isomorphe au conoyau d'un homomorphisme $\mathcal{O}_X^p|_V \to \mathcal{O}_X^q|_V$, on voit comme dans la démonstration du lemme 1 que $A(\mathcal{H}om_{\mathcal{O}_X}(F, G)) = \operatorname{Hom}_{A(X)}(A(F), A(G))$.

5. Changement du préschéma de base.

Proposition 11. — Soit $X$ un préschéma affine au-dessus de $S$. Pour toute extension $S' \to S$ du préschéma de base, $X' = X \times_S S'$ est affine au-dessus de $S'$. On a en outre $A(X') = g^*(A(X))$.

Pour démontrer la proposition, on peut se limiter au cas où $S$ et $S'$ sont affines ; il suffit alors de prouver que $X'$ est un schéma affine. Mais alors $X$ est un schéma affine, et si $A, A'$ et $B$ sont les anneaux de $S, S'$ et $X$ respectivement, on sait que $X'$ est alors un schéma affine d'anneau $A' \otimes_A B$ (I, 2, 4, prop. 6) ; en outre, on a alors $A(X) = \tilde{B}$, et $g^*(\tilde{B})$ est le $\tilde{A}'$-module associé à $A' \otimes_A B$ (I, 1, 6, prop. 8).

Corollaire 1. — Sous les hypothèses de la prop. 11, soit $f$ le morphisme structural $X \to S$, $f', g'$ les projections $X' \to S'$, $X' \to X$, de sorte

\newpage

%% ===== Page 221 =====

\marginpar{Archives Grothendieck-IHES-59}
\begin{center}
- II-32 -
\end{center}

que le diagramme
$$\begin{tikzcd}
X \arrow[d, "f"] \arrow[r, "\tilde{g}"] & X' \arrow[d, "f'"] \\
S \arrow[r, "g"] & S'
\end{tikzcd}$$
est commutatif. Pour tout $\mathcal{O}_X$-module quasi-cohérent $F$, il existe un isomorphisme canonique de $\mathcal{O}_{S'}$-modules
$$u : g^*(f_*(F)) \xrightarrow{\sim} f'_*(g'^*(F)) \text{ .}$$

On sait (chap. 0, \S 2, n$^\circ$5) que pour définir un homomorphisme $g^*(f_*(F)) \to f'_*(g'^*(F))$, il revient au même de définir un homomorphisme $f_*(F) \to g_*(f'_*(g'^*(F))) = f_*(g'_*(g'^*(F)))$; on prendra l'homomorphisme $f_*(\lambda)$, où $\lambda$ est l'homomorphisme canonique $F \to g'_*(g'^*(F))$ (loc. cit.). Pour prouver qu'il s'agit d'un isomorphisme, on peut, comme dans la prop. 11, se ramener au cas où $S$ et $S'$ sont affines; on a alors $F = \widetilde{M}$, où $M$ est un $B$-module, et on constate aussitôt que $g^*(f_*(F))$ et $f'_*(g'^*(F))$ sont tous deux égaux au $\mathcal{O}_{S'}$-module associé au $A'$-module $A' \otimes_A M$ (M considéré comme $A$-module), et que $u$ est l'homomorphisme associé à l'identité.

Remarque. — On se gardera de croire que le cor. 1 soit vrai lorsque $X$ n'est pas supposé affine au-dessus de $S$, même lorsque $S' = \operatorname{Spec}(\kappa(s))$ ($s \in S$), et que $S' \to S$ est le morphisme canonique (I, 2, 2) — auquel cas $X'$ n'est autre que la fibre $f^{-1}(s)$ (I, 2, 7, prop. 17). En d'autres termes, lorsque $X$ n'est pas affine au-dessus de $S$, l'opération "image directe de faisceaux quasi-cohérents" ne permute pas à l'opération "passage aux fibres". Nous verrons cependant au chap. III (\S, n$^\circ$) un résultat dans ce sens, de nature "asymptotique", valable pour les faisceaux cohérents sur $X$ lorsque $f$ est propre et $S$ noethérien.

Corollaire 2. — Soit $X$ un $S$-préschéma, $S'$ un préschéma affine au-dessus de $S$; alors $X' = X_{S'}$ est un préschéma affine au-dessus de $X$. Lorsque $S = \operatorname{Spec}(A)$, $S' = \operatorname{Spec}(A')$ sont affines, $X'$ est associé à la $\mathcal{O}_X$-algèbre $\mathcal{O}_X \otimes_A A'$ ($A$ et $A'$ étant identifié à des faisceaux simples

\newpage

%% ===== Page 222 =====

- 11-33 -

sur $X$).

Il suffit d'intervertir les rôles de $X$ et de $S'$ dans la prop. 11.

\subsection*{6. Morphismes affines}

Nous dirons qu'un morphisme $f: X \to Y$ de préschémas est affine s'il définit $X$ comme préschéma affine au-dessus de $Y$. Les propriétés des préschémas affines au-dessus d'un autre se traduisent comme suit dans cette terminologie :

\noindent Proposition 12. --- (i) Une immersion fermée est affine.
(ii) Le composé de deux morphismes affines est affine.
(iii) Si $f: X \to Y$ est un $S$-morphisme affine, $f^{S'}: X^{S'} \to Y^{S'}$ est affine pour toute extension $S' \to S$ de la base.
(iv) Si $f: X \to Y$ et $g: X' \to Y'$ sont deux $S$-morphismes affines, $f \times_S g: X \times_S X' \to Y \times_S Y'$ est affine.
(v) Si $f: X \to Y$ et $g: Y \to Z$ sont deux morphismes tels que $g \circ f$ soit affine et $g$ séparé, alors $f$ est affine.

En effet, (i) n'est autre que l'exemple suivant la déf. 1 du $n^\circ 2$, et (ii) le cor. 3 de la prop. 7 du $n^\circ 3$; (iii) découle de la prop. 11 du $n^\circ 5$, puisque $X^{S'}$ s'identifie au produit $X \times_Y Y'$ $(I, 2, 5, \text{cor. } 2 \text{ de la prop. } 9)$; (iv) se déduit de (iii) en factorisant $f \times_S g$ de la façon habituelle $X \times_S X' \to Y \times_S X' \to Y \times_S Y'$, et appliquant (ii). Enfin, pour démontrer (v), on peut se ramener au cas où $Z$ est un schéma affine; alors $X$ est un schéma affine et comme $Y$ est un schéma, (v) n'est autre que la propriété déjà signalée en Remarque suivant la déf. 1 du $n^\circ 2$, que l'on peut aussi énoncer comme

\noindent Corollaire. --- Si $X$ est un schéma affine et $Y$ un schéma, tout morphisme $X \to Y$ est affine.

\subsection*{7. Préschémas entiers sur un autre}

Etant donné un anneau $A$, rappelons qu'une $A$-algèbre $B$ est dite entière sur $A$ si tout élément de $B$ est racine dans $B$ d'un polynôme unitaire à coefficients dans $A$; il revient au même de dire que tout élément...

\newpage

%% ===== Page 223 =====

- II-34 -

ment de $B$ est contenu dans une sous-algèbre de $B$ qui est un $A$-module de type fini . Lorsqu'il en est ainsi , la sous-algèbre de $B$ engendrée par une partie finie de $B$ est un $A$-module de type fini ; pour que $B$ soit une algèbre entière et de type fini sur $A$, il faut et il suffit donc que $B$ soit un $A$-module de type fini . On notera que dans tout ceci , on ne suppose pas que l'homomorphisme $A \to B$ définissant la structure de $A$-algèbre soit injectif .

Définition 2. — Soient $X$ un $S$-préschéma , $f: X \to S$ son morphisme structural . On dit que $X$ est entier au-dessus de $S$, ou que $f$ est un morphisme entier, s'il existe un recouvrement $(S_\alpha)$ de $S$ par des ouverts affines tels que pour tout $\alpha$, le préschéma induit $f^{-1}(S_\alpha)$ soit un schéma affine dont l'anneau $B_\alpha$ est une algèbre entière sur l'anneau $A_\alpha$ de $S_\alpha$. On dit que $X$ est entier fini (ou simplement fini) au-dessus de $S$, ou que $f$ est un morphisme entier (ou simplement fini) si $X$ est entier sur $S$ et de type fini sur $S$.

Il est clair que si $X$ est entier au-dessus de $S$, il est affine au-dessus de $S$. Pour qu'un préschéma $X$ affine au-dessus de $S$ soit entier (resp. entier fini) au-dessus de $S$, il faut et il suffit que la $\mathcal{O}_S$-algèbre quasi-cohérente associée $\mathcal{A}(X)$ soit telle qu'avec les notations de la déf. 2, $\Gamma(S_\alpha, \mathcal{A}(X))$ soit une algèbre entière (resp. entière et de type fini) sur $\Gamma(S_\alpha, \mathcal{O}_S)$. Une $\mathcal{O}_S$-algèbre quasi-cohérente ayant cette propriété est dite entière (resp. entière finie, ou simplement finie) sur $\mathcal{O}_S$ si elle possède cette propriété. Se donner un préschéma entier (resp. fini) au-dessus de $S$ revient donc à se donner une $\mathcal{O}_S$-algèbre quasi-cohérente entière (resp. finie) sur $\mathcal{O}_S$. On notera qu'une $\mathcal{O}_S$-algèbre quasi-cohérente est finie si et seulement si c'est un $\mathcal{O}_S$-module de type fini.

Proposition 13. — Soit $S$ un préschéma localement noethérien. Pour qu'un préschéma $X$ affine au-dessus de $S$ soit fini au-dessus de $S$

\newpage

%% ===== Page 224 =====

- II-35 -

$S$, il faut et il suffit que la $\mathcal{O}_S$-algèbre $\mathcal{A}(X)$ soit cohérente.
D'après la remarque précédente, tout revient à remarquer que si $S$ est localement noethérien, les $\mathcal{O}_S$-modules quasi-cohérents de type fini ne sont autres que les $\mathcal{O}_S$-modules cohérents, ce qui résulte de I,1,5, th.3.

Proposition 14. - Soient $X$ un préschéma entier (resp. fini) au-dessus de $S$, $f: X \to S$ le morphisme structural. Alors, pour tout ouvert affine $U \subset S$, d'anneau $A$, $f^{-1}(U)$ est un schéma affine dont l'anneau $B$ est une $A$-algèbre entière (resp. finie) sur $A$.

On sait déjà ($n^\circ 3$, cor. 1 de la prop. 7) que $f^{-1}(U)$ est affine. Soit $\varphi: A \to B$ l'homomorphisme ; par hypothèse, il existe un recouvrement fini de $U$ par des ouverts $D(g_i)$ ($g_i \in A$) tels que, si on pose $h_i = \varphi(g_i)$, $B_{h_i}$ soit une algèbre entière (resp. finie) sur $A_{g_i}$. Considérons d'abord le cas où on suppose toutes les $B_{h_i}$ finies sur les $A_{g_i}$ : autrement dit, en tant que $A_{g_i}$-module, chaque $B_{h_i}$ admet un système de générateurs fini, qu'on peut supposer de la forme $(b_{ij} / h_i^n)$, $n$ étant le même pour tous les indices $i$. Montrons que les $b_{ij}$ forment un système de générateurs de $B$ en tant que $A$-module. En effet, soit $B'$ le sous-$A$-module de $B$ engendré par ces éléments. Par hypothèse, pour chaque $i$, il existe des $a_{ij} \in A$ et un entier $m$ (qu'on peut supposer indépendant de $i$) tel que, dans $B_{h_i}$, on ait $b/1 = (\sum_j a_{ij} b_{ij}) / h_i^m$ ; cela entraîne qu'il existe un entier $r \geqslant m$ tel que pour tout $i$, on ait $h_i^r b \in B'$. Or, comme les $D(g_i)$ recouvrent $U$, l'idéal de $A$ engendré par les $g_i$ est égal à $A$, autrement dit il existe des éléments $a_i \in A$ tels que $\sum_i a_i g_i = 1$ ; comme $g_i^r b = h_i^r b$ par définition, on a $b = \sum_i (a_i g_i^r) b = \sum_i a_i (h_i^r b) \in B'$, d'où notre assertion.

Supposons seulement maintenant chaque $B_{h_i}$ entière sur $A_{g_i}$. Pour tout $b \in B$, soit $C$ la sous-algèbre de $B$ engendrée par $b$. Comme...

\newpage

%% ===== Page 225 =====

- II-36 -

tout $i$, $C_{h_i}$ est l'algèbre sur $A_{g_i}$ engendrée par $b/1$ dans $B_{h_i}$, $C_{h_i}$
est par hypothèse un $A_{g_i}$-module de type fini, donc le raisonnement
précédent prouve que $C$ est un $A$-module de type fini, ce qui achève
la démonstration.

Corollaire .- Si $X$ est un préschéma entier (resp. fini) au-dessus de
$S$, alors, pour tout ouvert $U \subset S$, $f^{-1}(U)$ est entier (resp. fini)
au-dessus de $U$.

Cela résulte de la prop. 14 et du fait que les ouverts affines forment
une base de $S$.

Proposition 15 .- (i) Une immersion fermée est finie (et a fortiori
entière).

(ii) Le composé de deux morphismes finis (resp. entiers) est fini
(resp. entier).

(iii) Si $f : X \to Y$ est un $S$-morphisme fini (resp. entier), $f^{S'} : X^{S'} \to Y^{S'}$
est fini (resp. entier) pour toute extension $S' \to S$ de la base.

(iv) Si $f : X \to Y$ et $g : X' \to Y'$ sont deux $S$-morphismes finis (resp. en
tiers), il en est de même de $f \times_S g : X \times_S X' \to Y \times_S Y'$.

(v) Si $f : X \to Y$ et $g : Y \to Z$ sont deux morphismes tels que $g \circ f$ soit
fini (resp. entier) et $g$ séparé, alors $f$ est fini (resp. entier).

Pour démontrer qu'une immersion fermée $X \to S$ est finie, on peut se
borner au cas où $S = \operatorname{Spec}(A)$, et tout revient à remarquer qu'un
anneau quotient $A/I$ est un $A$-module monogène. Pour démontrer que le
composé de deux morphismes finis (resp. entiers) $X \to Y, Y \to Z$ est fini
(resp. entier), on peut encore supposer $Z$ (et par suite $X, Y$) affines
et alors l'assertion équivaut à dire que si $B$ est une $A$-algèbre finie
(resp. entière), et $C$ une $B$-algèbre finie (resp. entière), alors
$C$ est une $A$-algèbre finie (resp. entière). Pour démontrer (iii), on
peut se borner au cas où $S = Y$, puisque $X^{S'}$ s'identifie à $X \times_Y S'$ ;
on peut en outre supposer $S = \operatorname{Spec}(A)$, $S' = \operatorname{Spec}(A')$; alors $X$ est affine
d'anneau $B$, $X^{S'}$ est affine d'anneau $A' \otimes_A B$, et tout

\newpage

%% ===== Page 226 =====

- II-37 -

revient à remarquer que si $B$ est une $A$-algèbre finie (resp. entière), $A' \otimes_A B$ est une $A'$-algèbre finie (resp. entière). (iv) se déduit de (ii) et (iii) en factorisant $X \times_S X' \to Y \times_S X' \to Y \times_S Y'$. Enfin, (v) résulte de (i), (ii) et (iv) de la façon habituelle : on factorise $f : X \to X \times_Y Z \xrightarrow{p_2} Y$, $p_2$ s'identifiant à $(g \circ f) \times_Z Y$ ; d'après (iv), $p_2$ est fini (resp. entier), et comme $\Gamma_f$ est une immersion fermée (I, 3.7, prop. 22), il résulte de (i) que $\Gamma_f$ est un morphisme fini ; on conclut à l'aide de (ii).

Remarque. - L'hypothèse que $g$ est séparé est essentielle pour la validité de (v) : en effet, si $Y$ n'est pas séparé sur $Z$, l'identité $Y \to Y$ est le morphisme composé $Y \xrightarrow{\Delta_Y} Y \times_Z Y \xrightarrow{p_1} Y$, mais $\Delta_Y$ n'est pas un morphisme entier, comme il résulte de la prop. 16 qui suit.

Proposition 16. - Si $f : X \to Y$ est un morphisme entier, pour toute partie fermée $Z$ de $X$, $f(Z)$ est fermé dans $Y$.

Comme pour toute partie fermée $Z$ de $X$, il existe un sous-préschéma de $X$ ayant pour support $Z$ (I, 3.4, prop. 9), il résulte de la prop. 15, (i) et (ii) qu'on peut se borner à prouver que $f(X)$ est fermé dans $Y$. En vertu de la prop. 15 (vi), on peut supposer $X$ et $Y$ réduits ; en outre, si $T$ est le sous-préschéma réduit de $Y$ ayant pour support de base $\overline{f(X)}$ (I, 3.4, prop. 9), on sait que $f$ se factorise en $X \xrightarrow{g} T \xrightarrow{j} Y$, où $j$ est le morphisme d'injection (I, 3.4, cor. de la prop. 9), et comme $j$ est séparé (I, 3.6, prop. 16), il résulte de la prop. 15 (v) que $g$ est un morphisme entier. On peut donc supposer que $f(X)$ est dense dans $Y$. Enfin, la question étant locale sur $Y$, on peut se borner au cas où $Y$ est un spectre premier $\operatorname{Spec}(A)$. Alors $Y = \operatorname{Spec}(B)$, où $B$ est une $A$-algèbre entière sur $A$ ; en outre, l'hypothèse que $f(X)$ est dense dans $Y$ et que $A$ n'a pas d'élément nilpotent (I, 3.4, prop. 8) entraîne que l'homomorphisme $\varphi : A \to B$ est injectif (I, 1.2, cor. 4 de la prop. 5).

\newpage

%% ===== Page 227 =====

\centerline{-- II-38 --}

Dire que dans ces conditions $f(X)=Y$ signifie que tout idéal premier de $A$ est la trace sur $A$ d'un idéal premier de $B$, ce qui n'est outre le premier th. de Cohen-Seidenberg (Bourbaki, Alg. comm., chap. V, § 8).

\section*{8. Fibré vectoriel associé à un faisceau de modules}

Soient $A$ un anneau, $E$ un $A$-module. Rappelons qu'on appelle \emph{algèbre symétrique} sur $E$ et qu'on note $S(E)$ l'algèbre quotient de l'algèbre tensorielle $T(E)$ par l'idéal bilatère engendré par les éléments $x \otimes y - y \otimes x$, où $x, y$ parcourent $E$. L'algèbre $S(E)$ est solution du problème universel suivant : si $\sigma$ est l'application canonique $E \to S(E)$ (obtenue par composition de $E \to T(E)$ et de l'application canonique $T(E) \to S(E)$), toute application $A$-linéaire $E \to B$, où $B$ est une $A$-algèbre commutative, se factorise en $E \xrightarrow{\sigma} S(E) \xrightarrow{g} B$, où $g$ est un $A$-homomorphisme d'algèbres. On déduit aussitôt de cette caractérisation que pour deux $A$-modules $E, F$, on a $S(E \oplus F) \cong S(E) \otimes S(F)$ à un isomorphisme canonique près. En outre, $S(E)$ est un foncteur covariant de la catégorie des $A$-modules dans celle des $A$-algèbres. Par abus de langage, un produit $\sigma(x_1)\sigma(x_2) \dots \sigma(x_n)$, où les $x_i \in E$, se note souvent $x_1 x_2 \dots x_n$ si aucune confusion n'en résulte. L'algèbre $S(E)$ est graduée de façon évidente ; l'algèbre $S(A)$ est canoniquement isomorphe à l'algèbre de polynômes à une indéterminée $A[T]$ ; l'algèbre $S(A^n)$ à l'algèbre de polynômes à $n$ indéterminées $A[T_1, \dots, T_n]$.

Soit $\varphi: A \to B$ un homomorphisme d'anneaux. Si $F$ est un $B$-module, l'application canonique $F \to S(F)$ donne une application canonique $F_{[\varphi]} \to S(F_{[\varphi]})$ qui se factorise donc en $F \to S(F_{[\varphi]}) \to S(F)_{[\varphi]}$ ; on notera que l'homomorphisme canonique $S(F_{[\varphi]}) \to S(F)_{[\varphi]}$ n'est pas bijectif en général. Si $E$ est un $A$-module, tout $A$-homomorphisme $E \to F_{[\varphi]}$ (c'est-à-dire tout homomorphisme $E \to F$) donne donc canoniquement un homomorphisme d'algèbres $S(E) \to S(F_{[\varphi]}) \to S(F)_{[\varphi]}$, c'est-à-dire un $A$-homomorphisme d'algèbres $S(E) \to S(F)$. Avec les mêmes notations, il résulte aussitôt des définitions et de la propriété analogue pour $T(E)$.

\newpage

%% ===== Page 228 =====

- II-39 -

que pour tout $A$-module $E$, $S(E \otimes_A B)$ s'identifie canoniquement à l'algèbre $S(E) \otimes_A B$.

Enfin, soit $R$ une partie multiplicative de $A$ ; appliquant ce qui précède à l'anneau $B=R^{-1}A$, et se rappelant que $R^{-1}E = R^{-1}A \otimes_A E$, on voit que l'on a $S(R^{-1}E) = R^{-1}S(E)$ à un isomorphisme canonique près. En outre, si $R \supset R'$ est une seconde partie multiplicative de $A$, le diagramme

\begin{tikzcd}
R^{-1}E \arrow[r] & S(R^{-1}E) \\
\downarrow & \downarrow \\
{R'}^{-1}E \arrow[r] & S({R'}^{-1}E)
\end{tikzcd}

est commutatif.

Cela étant, considérons maintenant un espace annelé $(S, \mathcal{A})$, et soit $E$ un $\mathcal{A}$-module sur $S$. Si, à tout ouvert $U \subset S$, on associe le $\Gamma(U, \mathcal{A})$-module $S(\Gamma(U, E))$, on définit ainsi (vu le caractère fonctoriel de $S$) un préfaisceau d'algèbres ; on dit que le faisceau associé, que l'on note $S(E)$, est l'algèbre symétrique du $\mathcal{A}$-module $E$ ; il résulte aussitôt de ce qui précède que $S(E)$ est solution d'un problème universel : tout homomorphisme de $\mathcal{A}$-modules $E \to B$, où $B$ est une $\mathcal{A}$-algèbre, se factorise en $E \to S(E) \to B$, la seconde flèche étant un homomorphisme de $\mathcal{A}$-algèbres. Il y a donc correspondance biunivoque entre homomorphismes de $\mathcal{A}$-modules $E \to B$ et homomorphismes de $\mathcal{A}$-algèbres $S(E) \to B$. En particulier, tout homomorphisme $u: E \to F$ de $\mathcal{A}$-modules définit un homomorphisme $S(u): S(E) \to S(F)$ de $\mathcal{A}$-algèbres, et $S(E)$ est donc un foncteur covariant. Si $E, F$ sont deux $\mathcal{A}$-modules, $S(E \oplus F)$ s'identifie canoniquement à $S(E) \otimes_{\mathcal{A}} S(F)$, comme le montre la considération des préfaisceaux correspondants.

Soit $(T, \mathcal{B})$ un second espace annelé, $f$ un morphisme $(S, \mathcal{A}) \to (T, \mathcal{B})$. Si $F$ est un $\mathcal{B}$-module, $S(f^*(F))$ s'identifie à $f^*(S(F))$; en effet, si $f=(\varphi, f^\#)$ [unclear: et U est un ouvert de S], une section de $S(f^*(F))$ sur $U$ s'identifie au voisinage de chaque point de $U$...

\newpage

%% ===== Page 229 =====

-- II-40 --

on a $\underline{S}(f^*(\mathcal{F})) \simeq \underline{S}(f^*(\mathcal{F})) \otimes_{f^*(\mathcal{B})} f^*(\mathcal{B}) \simeq \underline{S}(\underline{f}^*(\mathcal{F})) \otimes_{f^*(\mathcal{B})} f^*(\mathcal{B})$ ; pour tout ouvert $U$ de $S$ et toute section $u$ de $\underline{S}(f^*(\mathcal{E}))$ au-dessus de $U$, $u$ coïncide, dans un voisinage $V$ de chaque point $s \in U$, avec un élément de $\underline{S}(f^*(\mathcal{E}))(V, f^*(\mathcal{E}))$; par suite, il y a un voisinage $W$ de $f(s)$ dans $T$, une section $u'$ de $\underline{S}(\mathcal{F})$ au-dessus de $W$, et un voisinage $V' \subset V \cap f^{-1}(W)$ de $s$ tels que $u$ coïncide avec $u' \circ f$ au-dessus de $V'$, d'où notre assertion.

On notera enfin que $\underline{S}(\mathcal{E})$ est somme directe infinie des $\underline{S}_n(\mathcal{E})$, où le faisceau de modules $\underline{S}_n(\mathcal{E})$ est associé au préfaisceau des $\underline{S}_n(\Gamma(U, \mathcal{E}))$.

Proposition 18. — Soient $A$ un anneau, $S = \operatorname{Spec}(A)$ son spectre premier, et $\mathcal{E}$ le $\mathcal{O}_S$-module associé à un $A$-module $M$ ; alors la $\mathcal{O}_S$-Algèbre $\underline{S}(\mathcal{E})$ est associée à la $A$-algèbre $S(M)$.

En effet, pour tout $f \in A$, $\underline{S}(\mathcal{E})_f = (S(M))_f$, et la proposition résulte donc des définitions.

Corollaire. — Si $S$ est un préschéma, $\mathcal{E}$ un $\mathcal{O}_S$-Module quasi-cohérent, la $\mathcal{O}_S$-Algèbre $\underline{S}(\mathcal{E})$ est quasi-cohérente. Si en outre $S$ est localement noethérien, et si $\mathcal{E}$ est cohérent, alors chacun des $\mathcal{O}_S$-Modules $\underline{S}_n(\mathcal{E})$ est cohérent.

La première assertion est conséquence immédiate de la prop. 18 et de I, 1, 4, th. 2 ; la seconde résulte de ce que si $\mathcal{E}$ est un $A$-module de type fini, $\underline{S}_n(\mathcal{E})$ est un $A$-module de type fini ; on applique alors I, 1, 5, th. 3.

Définition 3. — Soit $\mathcal{E}$ un $\mathcal{O}_S$-Module quasi-cohérent. On appelle fibré vectoriel défini par $\mathcal{E}$ et on note $\mathbb{V}(\mathcal{E})$ le préschéma affine au-dessus de $S$ associé ($n^o 3$) à la $\mathcal{O}_S$-Algèbre $\underline{S}(\mathcal{E})$.

En vertu de la prop. 6 du $n^o 2$, pour tout $S$-préschéma adéquat, il y a correspondance biunivoque canonique entre les $S$-morphismes $X \to \mathbb{V}(\mathcal{E})$ et les $\mathcal{O}_S$-homomorphismes d'algèbres $\underline{S}(\mathcal{E}) \to \mathcal{A}(X)$, et par suite aussi entre ces morphismes et les $\mathcal{O}_S$-homomorphismes de Module

\newpage

%% ===== Page 230 =====

- II-41 -

$\underline{E \to A(X) = f_*(O_X)}$. En particulier :

1$^\circ$ Prenons pour $X$ un sous-préschéma induit par $S$ sur un ouvert $U \subset S$. Alors, si $p$ est le morphisme structural $\underline{V}(E) \to S$, un $S$-morphisme $u : U \to \underline{V}(E)$ est par définition un morphisme tel que $p \circ u = 1_U$, autrement dit une section du morphisme $p$ au-dessus de $U$, ou, comme nous dirons encore, une $S$-section de $\underline{V}(E)$ au-dessus de $U$. D'après ce qu'on vient de voir, ces $S$-sections correspondent biunivoquement aux homomorphisme de $\mathcal{O}_S$-Module $E \to f_*(O_U)$, ce qui revient au même (chap. 0, § 2, n$^\circ$ 5) aux $\mathcal{O}_S$-homomorphismes $f^*(E) \to O_U$ ($f$ injection canonique $U \to S$). En d'autres termes, le faisceau des germes de $S$-sections de $\underline{V}(E)$ s'identifie canoniquement au faisceau $\underline{\operatorname{Hom}}_{\mathcal{O}_S}(E, \mathcal{O}_S)$ (qu'on peut appeler le dual du faisceau $E$).

2$^\circ$ Prenons pour $X$ le spectre $\{s\}$ d'un corps $K$, le morphisme structural $f : X \to S$ correspond donc à un monomorphisme $\kappa(s) \to K$, où $s = f(s)$ (I, 2, 2, cor. de la prop. 5); les $S$-morphismes $\{s\} \to \underline{V}(E)$ ne sont donc autres que les points géométriques de $\underline{V}(E)$ à valeurs dans l'extension $K$ de $\kappa(s)$, points qui sont localisés aux points de la fibre $p^{-1}(s)$. L'ensemble de ces points, qu'on peut encore appeler fibre géométrique rationnelle sur $K$ de $\underline{V}(E)$ au-dessus du point $s$, s'identifie d'après ce qui précède à l'ensemble des $\mathcal{O}_S$-homomorphismes $E \to f_*(O_X)$ (chap. 0, § 2, n$^\circ$ 5), ou, ce qui revient au même, à l'ensemble des $\mathcal{O}_S$-homomorphismes $f^*(E) \to K$. Mais on a par définition $f^*(E) = E \otimes_{\mathcal{O}_S} K = E^s \otimes_{\kappa(s)} K$, en posant $E^s = E_s / m_s E_s$, $m_s$ étant l'idéal maximal de $\mathcal{O}_S$ ; la fibre géométrique de $\underline{V}(E)$ au-dessus de $s$ s'identifie donc au dual de l'espace vectoriel $E^s \otimes_{\kappa(s)} K$ ; si $E^s$ ou $K$ est de dimension finie sur $\kappa(s)$, ce dual s'identifie aussi à $(E^s)^\vee \otimes_{\kappa(s)} K$, où $(E^s)^\vee$ est le dual du $\kappa(s)$-espace vectoriel $E^s$.

Proposition 12. - (i) $\underline{V}(E)$ est un foncteur contravariant de la catégorie des $\mathcal{O}_S$-Modules quasi-cohérents dans la catégorie des préschémas affines au-dessus de $S$.

\newpage

%% ===== Page 231 =====

- II-42 -

(ii) Si $\mathcal{E}$ est un $\mathcal{O}_S$-module de type fini, $\mathbf{V}(\mathcal{E})$ est de type fini sur $S$.
(iii) Si $\mathcal{E}$ et $\mathcal{F}$ sont deux $\mathcal{O}_S$-modules, $\mathbf{V}(\mathcal{E} \oplus \mathcal{F})$ s'identifie canoniquement à $\mathbf{V}(\mathcal{E}) \times_S \mathbf{V}(\mathcal{F})$.
(iv) Soient $g: S' \to S$ un morphisme ; pour tout $\mathcal{O}_S$-module $\mathcal{E}$, $g^*(\mathbf{V}(\mathcal{E}))$ s'identifie canoniquement à $(\mathbf{V}(\mathcal{E}))^{S'} = \mathbf{V}(\mathcal{E}) \times_S S'$.

(i) est conséquence immédiate de la prop. 6 du n° 2, compte tenu de ce que tout homomorphisme $\mathcal{E} \to \mathcal{F}$ de $\mathcal{O}_S$-modules définit canoniquement un homomorphisme $S(\mathcal{E}) \to S(\mathcal{F})$ de $\mathcal{O}_S$-algèbres. (ii) découle aussitôt des définitions et du fait que $S(\mathcal{E})$ est une $\mathcal{A}$-algèbre de type fini si $\mathcal{E}$ est un $\mathcal{A}$-module de type fini. Pour démontrer (iii), on peut supposer $S = \operatorname{Spec}(A)$, $\mathcal{E} = \widetilde{M}$, $\mathcal{F} = \widetilde{N}$, où $M$ et $N$ sont des $A$-modules ; l'assertion résulte alors de ce que $S(M \oplus N) = S(M) \otimes_A S(N)$. De même, pour démontrer (iv), on peut supposer $S, S'$ affines, soit $S = \operatorname{Spec}(A)$, $S' = \operatorname{Spec}(A')$, et $\mathcal{E} = \widetilde{M}$, $M$ étant un $A$-module, alors $g^*(\mathcal{E}) = \widetilde{N}$, où $N = M \otimes_A A'$ (1.1, 6, prop. 8), et le résultat découle de ce que $S(N) = S(M) \otimes_A A'$.

Si on prend en particulier $\mathcal{E} = \mathcal{O}_S$, le préschéma $\mathbf{V}(\mathcal{O}_S)$, qu'on note aussi $\mathbb{A}_S$, est appelé le préschéma structural au-dessus de $S$ ; c'est le préschéma affine au-dessus de $S$ associé à la $\mathcal{O}_S$-algèbre $S(\mathcal{O}_S)$, qui s'identifie à la $\mathcal{O}_S$-algèbre $\mathcal{O}_S[T] = \mathcal{O}_S \otimes_{\mathbb{Z}} \mathbb{Z}[T]$, où $T$ est une indéterminée, $\mathbb{Z}$ et $\mathbb{Z}[T]$ des faisceaux simples sur $S$. Comme on l'a vu ci-dessus, le faisceau des germes de $S$-sections de $\mathbb{A}_S$ s'identifie à $\mathcal{O}_S$ lui-même. Pour tout morphisme $g : S' \to S$, on a $g^*(\mathcal{O}_S) = \mathcal{O}_{S'}$, donc $\mathbb{A}_{S'}$ s'identifie canoniquement à $(\mathbb{A}_S)^{S'} = \mathbb{A}_S \times_S S'$ en vertu de la prop. 19.

\newpage

%% ===== Page 232 =====

\marginpar{Archives Grothendieck sept 59}

\centerline{-- II-43 --}

\section*{§ 3. Spectres premiers homogènes.}

\subsection*{1. Généralités sur les anneaux gradués.}

\noindent \textbf{Notations.} --
Soit $S$ un anneau gradué, à degrés positifs ; on désignera par $S_n$ la partie de $S$ formée des éléments homogènes de degré $n$ ($n \ge 0$) ; par $S_+$ la somme (directe) des $S_n$ tels que $n > 0$ ; on a $1 \in S_0$, et $S_+$ est un idéal gradué de $S$ et $S_0$ un sous-anneau de $S$. Pour tout élément homogène $f$ de degré $d > 0$, on désigne par $S_{(f)}$ le sous-anneau (évidemment gradué) de l'anneau des fractions $S_f$ formé des éléments homogènes de degré $0$ : tout élément de $S_{(f)}$ est donc de la forme $x/f^n$ avec $x \in S_{nd}$ ($n \ge 0$).

Soit $M$ un module gradué sur $S$ (à degrés positifs ou négatifs) ; on désignera par $M_n$ le sous-module de $M$ formé des éléments de degré $n$. Avec les mêmes hypothèses sur $f$, on désigne par $M_{(f)}$ le sous-module (évidemment gradué) du module de fractions $M_f$, formé des éléments homogènes de degré $0$ ; les éléments de $M_{(f)}$ sont donc de la forme $z/f^n$ avec $z \in M_{nd}$ ($n \ge 0$). Il est clair que $M_{(f)}$ est un $S_{(f)}$-module.

On notera que pour tout entier $k > 0$, l'anneau $S_{(f)}$ et le module $M_{(f)}$ sont entièrement déterminés par la connaissance des éléments des modules $S_{nd}$ et $M_{nd}$, car si $x \in S_{nd}$ : $x/f^n = (x f^{k-1)n}) / f^{kn}$ et de même dans $M_{(f)}$.

\noindent \textbf{Lemme 1.} -- Soient $d, d'$ deux entiers $> 0$, $f \in S_d$, $g \in S_{d'}$. Il existe alors un isomorphisme canonique d'anneaux
$$S_{(fg)} \simeq (S_{(f)})_{(g/f^{d'})} ;$$
si on identifie ces deux anneaux, il existe un isomorphisme canonique de modules
$$M_{(fg)} \simeq (M_{(f)})_{(g/f^{d'})}$$
En effet, les éléments de $M_{(fg)}$ sont de la forme $x/(fg)^{nd''}$, où $x \in S_{nd''d''}$ ($d''$ étant un entier arbitraire $> 0$) ; à cet élément on fait correspondre l'élément $(x/f^{nd''}) / (g/f^{d'})^{nd''}$ de l'anneau

\newpage

%% ===== Page 233 =====

- II-44 -

$ (S_{(f)})_{g^d/f^d} $. Cette application est bien définie; s'il existe $k \geqslant 0$ tel que $(fg)^k x = 0$ dans $S$, on a aussi $(g^d/f^d)^k (x/f^{nd}) = 0$ dans $S_{(f)}$, car cet élément est égal à $(fg)^{kd} x / f^{(k+nd)d} = 0$ par hypothèse. Réciproquement, comme tout élément de $(S_{(f)})_{g^d/f^d}$ s'écrit $(x/f^{nd}) / (g^d/f^d)^{n'}$ pour un $n$ convenable et un $x \in S_{nd'}$, on fait correspondre à cet élément l'élément $x/(fg)^{nd d'}$ de $S_{(fg)}$; on a encore une application bien définie, car si on a $(g^d/f^d)^k (x/f^{nd}) = 0$, cela s'écrit $g^{kd} x / f^{kd+nd} = 0$, donc il existe $h$ tel que $f^h g^{kd} x = 0$ dans $S$, et cela entraîne $(fg)^{h+kd} x = 0$. On procède de la même façon pour $M$, et il est immédiat de vérifier que les applications ainsi définies sont des isomorphismes (resp. des isomorphismes de modules) inverses l'une de l'autre.

Avec les mêmes notations et hypothèses, il est clair que l'homomorphisme canonique $S_g \to S_{fg}$, transformant $x/f^n$ en $g^k x / (fg)^k$ est de degré 0, donc donne par restriction un homomorphisme canonique d'anneaux $S_{(f)} \to S_{(fg)}$. Le lemme 1 entraîne alors :

Lemme 2. — L'anneau $S_{(fg)}$ (resp. le module $M_{(fg)}$) est engendré par la réunion des images canoniques de $S_{(f)}$ et $S_{(g)}$ (resp. de $M_{(f)}$ et $M_{(g)}$).

Il suffit en effet de voir que $1/(g^d/f^d) = f^d/g^d$ appartient à l'image canonique de $S_{(g)}$, ce qui est évident par définition.

Pour tout entier $d > 0$, nous désignerons par $S^{(d)}$ l'anneau gradué somme directe des $S_{nd}$ ($n \geqslant 0$), par $M^{(d)}$ le $S^{(d)}$-module gradué somme directe des $M_{nd}$ ($n \in \mathbb{Z}$).

Proposition 1. — Soit $d$ un entier $> 0$, et soit $f \in S_d$. Alors il existe un isomorphisme canonique d'anneaux $S_{(f)} \cong S^{(d)} / (f-1)S^{(d)}$; si on identifie ces deux anneaux, il existe un isomorphisme canonique de $M_{(f)} \cong M^{(d)} / (f-1)M^{(d)}$.

Le premier de ces isomorphismes se définit en faisant correspondre

\newpage

%% ===== Page 234 =====

-- II-45 --

à $x/f^n$, où $x \in S_{nd}$, l'élément $\overline{x}$, classe de $x \bmod (f-1)S^{(d)}$ ; cette application est bien définie, car si $f^n x = 0$, on peut écrire $x = (1-f^n)x = (f-1)(f^{n-1} + f^{n-2} + \dots + 1)x \in (f-1)S^{(d)}$. On remarque d'autre part que si $x \in S_{nd}$ est tel que $x = (f-1)y$ avec $y = y_{hd} + y_{(h+1)d} + \dots + y_{kd}$, $y_{jd} \in S_{jd}$, $y_{hd} \neq 0$, on a nécessairement $h=n$ et $x = -y_{hd}$, ainsi que les relations $y_{(j+1)d} = f y_{jd}$ pour $h \le j \le k-1$, $f y_{kd} = 0$, ce qui donne finalement $f^{k-n}x = 0$ ; à toute classe $\overline{x} \bmod (f-1)S^{(d)}$ d'un élément $x \in S_{nd}$ on peut donc faire correspondre l'élément $x/f^n$, car la remarque qui précède prouve que cette application est bien définie. Il est immédiat que les deux applications ainsi définies sont des homomorphismes d'anneaux, inverses l'une de l'autre. On procède exactement de même pour $M$.

Soit $T$ une partie multiplicative de $S_+$, formée d'éléments homogènes ; $T_0 = T \cup \{1\}$ est alors une partie multiplicative de $S$ ; l'anneau $T^{-1}S$ est encore gradué de façon évidente (les éléments de $T_0$ étant homogènes), et on désigne par $S_{(T)}$ le sous-anneau de $T^{-1}S$ formé des éléments de degré $0$ ; ce sont donc les éléments de la forme $x/h$, où $h \in T$ et $x$ est homogène de degré égal à celui de $h$. On sait que $T^{-1}S$ s'identifie canoniquement à la limite inductive des anneaux $S_f$, où $f$ parcourt $T$ (pour les homomorphismes canoniques $S_f \to S_g$, où $g$ est un multiple de $f$) ; comme cette identification respecte les degrés, elle identifie $S_{(T)}$ à la limite inductive des $S_{(f)}$ pour $f \in T$. Pour tout $S$-module gradué $M$, on définit de même le sous-module $M_{(T)}$ (sur l'anneau $S_{(T)}$) formé des éléments de degré $0$ de $T^{-1}M$, et on voit que c'est la limite inductive des $M_{(f)}$ où $f \in T$.

Soit $\mathfrak{p}$ un idéal premier gradué de $S$, qui est donc somme directe des sous-groupes $\mathfrak{p}_n = \mathfrak{p} \cap S_n$ ; supposons que $\mathfrak{p}$ ne contienne pas $S_+$. Alors si $f \in S_+$ n'appartient pas à $\mathfrak{p}$, la relation $f^n x \in \mathfrak{p}$ est équivalente à $x \in \mathfrak{p}$ ; en particulier $f^n x \in \mathfrak{p}_m$ équivaut à $x \in \mathfrak{p}_{m-nd}$.

\newpage

%% ===== Page 235 =====

— II-46 —

Proposition 2. — Soit $n_0$ un entier $> 0$ ; pour tout $n \ge n_0$, soit $\mathfrak{p}_n$ un sous-groupe de $S_n$. Pour qu'il existe un idéal premier gradué $\mathfrak{p}$ de $S$, ne contenant pas $S_+$, et tel que $\mathfrak{p} \cap S_n = \mathfrak{p}_n$ pour tout $n \ge n_0$ : il faut et il suffit que les conditions suivantes soient vérifiées : $1^\circ$ $S_m \mathfrak{p}_n \subset \mathfrak{p}_{m+n}$ pour tout $m \ge 0$ et tout $n \ge n_0$ ; $2^\circ$ pour $m \ge n_0$, $n \ge n_0$, $f \in S_m$, $g \in S_n$, la relation $fg \in \mathfrak{p}_{m+n}$ entraîne $f \in \mathfrak{p}_m$ ou $g \in \mathfrak{p}_n$ ; $3^\circ$ $\mathfrak{p}_n \neq S_n$ pour un $n \ge n_0$ au moins. En outre l'idéal premier gradué $\mathfrak{p}$ est alors unique.

En effet, si $\mathfrak{p} \not\supset S_+$, il existe au moins un $k > 0$ tel que $\mathfrak{p} \cap S_k \neq S_k$, puisque $\mathfrak{p}$ est gradué ; si $f \in S_n$ n'appartient pas à $\mathfrak{p}$, la relation $\mathfrak{p} \cap S_{n-mk} = S_{n-mk}$ entraîne $\mathfrak{p} \cap S_n = S_n$ d'après les remarques précédentes ; donc si $\mathfrak{p} \cap S_n = S_n$ à partir d'une certaine valeur de $n$, on aurait, nécessairement $\mathfrak{p} \supset S_+$. Les conditions de l'énoncé sont donc nécessaires. Inversement, si elles sont satisfaites, et si pour un entier $d \ge n_0$, $f \in S_d$ n'appartient pas à $\mathfrak{p}$, alors, si $\mathfrak{p}$ existe, $\mathfrak{p}_m$ pour $m < n_0$, est nécessairement égal à l'ensemble des $z \in S_m$ tels que $f^r z \in \mathfrak{p}_{m+rd}$ presque pour tout $r$. Cela prouve déjà que si $\mathfrak{p}$ existe, il est unique. Reste à montrer que si on définit les $\mathfrak{p}_m$ pour $m < n_0$ par la condition précédente, $\mathfrak{p} = \bigoplus \mathfrak{p}_m$ est un idéal premier. Notons d'abord qu'en vertu des hypothèses, pour $n \ge n_0$, $\mathfrak{p}_n$ est aussi défini comme l'ensemble des $x \in S_n$ tels que $f^r x \in \mathfrak{p}_{n+rd}$ pour $r$ assez grand. Cela étant si $g \in S_m$, $x \in \mathfrak{p}_n$, on a $f^r gx \in \mathfrak{p}_{m+n+rd}$ pour $r$ assez grand, et par suite $gx \in \mathfrak{p}_{m+n}$. D'autre part, si $x \in S_m$, $y \in S_n$ sont tels que $xy \in \mathfrak{p}_{m+n}$ et $y \notin \mathfrak{p}_n$, alors $f^r xy \in \mathfrak{p}_{m+n+rd}$ pour $r$ assez grand, $f^{2r} xy \in \mathfrak{p}_{m+n+2rd}$, donc on en conclut $f^r x \in \mathfrak{p}_{m+rd}$ pour $r$ assez grand, donc $x \in \mathfrak{p}_m$, ce qui achève la démonstration.

Nous dirons qu'une partie $J$ de $S_+$ est un idéal de $S_+$ si c'est l'intersection avec $S_+$ un idéal de $S$ ; on dit que c'est un idéal premier gradué de $S_+$ s'il est l'intersection de $S_+$ et d'un idéal premier gradué de $S$.

\newpage

%% ===== Page 236 =====

- II-47 -

ne contenant pas $S_+$ (cet idéal premier étant d'ailleurs unique d'après la prop.2). Si $J$ est un idéal de $S_+$, la racine de $J$ dans $S_+$ est l'ensemble des éléments de $S_+$ dont une puissance appartient à $J$, autrement dit c'est l'ensemble $r_+(J) = r(J) \cap S_+$.

Si $\mathfrak{p}$ est un idéal premier gradué de $S_+$, on désignera par $S_{(\mathfrak{p})}$ et $M_{(\mathfrak{p})}$ l'anneau $S_{(T)}$ et le module $M_{(T)}$, où $T$ désigne l'ensemble des éléments homogènes de $S_+$ qui n'appartiennent pas à $\mathfrak{p}$.

\subsection*{2. Spectre premier homogène d'un anneau gradué.}

Etant donné un anneau gradué $S$, on appelle spectre premier homogène de $S$ et on désigne par $\operatorname{Proj}(S)$ l'ensemble des idéaux premiers gradués de $S_+$ ; ou, ce qui revient au même, l'ensemble des idéaux premiers gradués de $S$ ne contenant pas $S_+$ ; nous allons définir sur $\operatorname{Proj}(S)$ une structure de schéma.

Pour toute partie $E$ de $S$, soit $V_+(E)$ l'ensemble des idéaux premiers gradués de $S_+$ ne contenant pas $S_+$ et contenant $E$ ; c'est donc la partie $V(E) \cap \operatorname{Proj}(S)$ de $\operatorname{Spec}(S)$. De 1.1.1, prop.1 on déduit donc :

(1) $$V_+(0) = \operatorname{Proj}(S) \quad ; \quad V_+(S) = V_+(S_+) = \emptyset$$

(2) $$V_+(\bigcup E_\alpha) = \bigcap V_+(E_\alpha)$$

(3) $$V_+(EE') = V_+(E) \cup V_+(E')$$

On ne change pas $V_+(E)$ en remplaçant $E$ par l'idéal gradué engendré par $E$ ; en outre, si $J$ est un idéal gradué de $S$, il résulte du n°1, que l'on a

(4) $$V_+(J) = V_+(\bigcup_{n \ge 0} J \cap S_n)$$

pour tout $n > 0$ : car si $\mathfrak{p} \in V_+(J)$ contient les éléments homogènes de $J$ de degré $\ge n$, comme par hypothèse il existe un $f$ homogène non contenu dans $\mathfrak{p}$, pour tout $x \in J \cap S_m$, on a $f^r x \in J_{m+rd}$ pour tout $r > 0$ (si $f \in S_d$), donc $f^r x \in \mathfrak{p}_{m+rd}$ pour $r$ assez grand, d'où $x \in \mathfrak{p}_m$. Enfin, on a

(5) $$V_+(J) = V_+(r(J))$$

\newpage

%% ===== Page 237 =====

--- II-48 ---

Par définition, les $V_+(E)$ sont les parties fermées de $X = \operatorname{Proj}(S)$ pour la topologie induite par la topologie spectrale de $\operatorname{Spec}(S)$, et qu'on appellera encore topologie spectrale sur $\operatorname{Proj}(S)$. On pose, pour tout $f \in S$
$$(6) \qquad D_+(f) = D(f) \cap \operatorname{Proj}(S) = \operatorname{Proj}(S) - V_+(f)$$
et on a par suite $(1, 1, 1, \text{formule } (3))$
$$(7) \qquad D_+(fg) = D_+(f) \cap D_+(g)$$

\emph{Proposition 3. --} Lorsque $f$ parcourt l'ensemble des éléments homogènes de $S_+$, les $D_+(f)$ forment une base de la topologie de $X = \operatorname{Proj}(S)$.

On a vu en effet ci-dessus que toute partie fermée de $X$ est intersection d'ensembles de la forme $V_+(f)$, où $f$ est homogène de degré $>0$.

Soit $f$ un élément homogène de $S_+$, de degré $d > 0$ ; pour tout idéal premier gradué $\mathfrak{p}$ de $S$, ne contenant pas $f$, on sait que l'ensemble des $x/f^n$, où $x \in \mathfrak{p}_n$ et $n \ge 0$, un idéal premier de l'anneau des fractions $S_{(f)}$ ; sa trace sur $(S_f)_0 = S_{(f)}$ est donc un idéal premier de cet anneau, que nous désignerons par $\psi_f(\mathfrak{p})$ ; c'est donc l'ensemble des $x/f^n$, où $x \in \mathfrak{p}_{nd}$ et $n \ge 0$. On a ainsi défini une application $\psi_f : D_+(f) \to \operatorname{Spec}(S_{(f)})$; si $g \in S_d$, est un second élément homogène de $S_+$, on a un diagramme commutatif
$$(8) \qquad \begin{tikzcd} D_+(fg) \arrow[r, "\psi_{fg}"] \arrow[d, hook] & \operatorname{Spec}(S_{(fg)}) \arrow[d, "\omega"] \\ D_+(f) \arrow[r, "\psi_f"] & \operatorname{Spec}(S_{(f)}) \end{tikzcd}$$
où la flèche verticale de gauche est l'inclusion, et celle de droite est l'application $\omega_{fg,f}$ déduite de l'homomorphisme canonique $\omega = \omega_{fg,f} : S_{(f)} \to S_{(fg)}$ ; en effet, si $x/f^n \in \omega^{-1}(\psi_{fg}(\mathfrak{p}))$, où $fg \notin \mathfrak{p}$, on a par définition $g^n x / (fg)^n \in \psi_{fg}(\mathfrak{p})$, donc $g^n x \in \mathfrak{p}$, et par suite $x \in \mathfrak{p}$, et la réciproque est évidente.

\emph{Proposition 4. --} L'application $\psi_f$ est un homéomorphisme de $D_+(f)$ sur $\operatorname{Spec}(S_{(f)})$.

En premier lieu, $\psi_f$ est continue ; car si $h \in S_{nd}$ est tel que...

\newpage

%% ===== Page 238 =====

- II-49 -

et $h/f^n \in \varphi_f(\mathfrak{p})$, on a par définition $h \in \mathfrak{p}$, donc $\varphi_f^{-1}(D(h/f^n)) = D_+(hf)$, et notre assertion résulte de la formule (7) ; en outre, les $D_+(hf)$, où $h$ parcourt les ensembles $S_{nd}$, forment une base de la topologie de $D_+(f)$ en vertu de la prop.3 et de la formule (7) ; ce qui prouve donc, compte tenu de l'axiome $(T_0)$, que $\varphi_f$ est injective et que l'application inverse $\varphi_f^{-1} : \operatorname{Spec}(S_{(f)}) \to D_+(f)$ est continue. Enfin, pour voir que $\varphi_f$ est surjective, on remarque que, si $\mathfrak{q}_0$ est un idéal premier de $S_{(f)}$ et si, pour tout $n > 0$, on désigne par $\mathfrak{p}_n$ l'ensemble des $x \in S_n$ tels que $x^d/f^n \in \mathfrak{q}_0$, les $\mathfrak{p}_n$ vérifient les conditions de la prop.2 (en effet, si $x \in S_n, y \in S_n$ tels que $x^d/f^n \in \mathfrak{q}_0$ et $y^d/f^n \in \mathfrak{q}_0$, on a $(x+y)^{2d}/f^{2n} \in \mathfrak{q}_0$, d'où $(x+y)^d/f^n \in \mathfrak{q}_0$, puisque $\mathfrak{q}_0$ est premier). Si $\mathfrak{p}$ est l'idéal premier gradué de $S$ ainsi défini, on a bien $\varphi_f(\mathfrak{p}) = \mathfrak{q}_0$, car si $x \in S_{nd}$, les relations $x/f^n \in \mathfrak{q}_0$ et $x^d/f^n \in \mathfrak{q}_0$ sont équivalentes.

\marginpar{premier gradué}

Corollaire 1. - Pour que $D_+(f) = \emptyset$, il faut et il suffit que $f$ soit nilpotent.

En effet, pour que $\operatorname{Spec}(S_{(f)}) = \emptyset$, il faut et il suffit que $S_{(f)}$ soit réduit à 0, ou encore que $f/f = 0$ dans $S_{(f)}$, ce qui signifie que $f$ est nilpotent.

Corollaire 2. - Soit $E$ un sous-groupe gradué de $S_+$. Les conditions suivantes sont équivalentes :

a) $V_+(E) = X$.

b) Tout élément de $E$ est nilpotent.

c) Les composantes homogènes de tout élément de $E$ sont nilpotentes.

Il est clair que c) entraîne b) et que b) entraîne a). Si $J$ est l'idéal gradué de $S$ engendré par $E$, a) équivaut à $V_+(J) = X$ ; a fortiori a) entraîne que tout élément homogène $f \in E \subset J$ est tel que $V_+(f) = X$, donc $f$ est nilpotent en vertu du cor. 1.

\newpage

%% ===== Page 239 =====

\noindent \underline{Corollaire 3} .— Si $J$ est un idéal gradué de $S_+$, $r(J)$ est l'intersection des idéaux premiers gradués de $S$ contenant $J$.

En considérant l'anneau gradué $S/J$, on peut se ramener au cas où $J=(0)$. Il faut montrer que si $f \in S_+$ n'est pas nilpotent, il y a un idéal premier gradué ne contenant pas $f$ ; or, une au moins des composantes homogènes de $f$ n'est pas nilpotente, et on est ainsi ramené au cas où $f$ est homogène ; l'assertion résulte alors du cor. 1.

Pour toute partie $Y$ de $X$, soit $j_+(Y)$ l'ensemble des $f \in S_+$ tels que $Y \subset V_+(f)$ ; il revient au même de dire que $j_+(Y)=j(Y) \cap S_+$; $j_+(Y)$ est donc un idéal de $S_+$, égal à sa racine dans $S_+$.

\underline{Proposition 5} .— a) Pour toute partie $E$ de $S_+$, $j_+(V_+(E))$ est la racine de l'idéal gradué de $S$ engendré par $E$.

b) Pour toute partie $Y$ de $X$, $V_+(j_+(Y))=Y$, adhérence de $Y$ dans $X$.

c) Si $J$ est l'idéal gradué de $S_+$ engendré par $E$, on a $V_+(E)=V_+(J)$, et l'assertion résulte alors du cor. 3 de la prop. 4.

b) Comme $V_+(J)=\bigcap_{f \in J} V_+(f)$, la relation $Y \subset V_+(f)$ entraîne $Y \subset V_+(f)$ pour tout $f \in J$, et par suite $J \supset j_+(Y)$, d'où $V_+(j_+(Y)) \subset V_+(J)$, ce qui prouve b) par définition des ensembles fermés.

\underline{Corollaire 1} .— Les parties fermées $Y$ de $X=\operatorname{Proj}(S)$ et les idéaux gradués de $S_+$ égaux à leur racine dans $S_+$ se correspondent bijectivement par les applications décroissantes $Y \to j_+(Y)$, $J \to V_+(J)$ ; à la réunion $Y_1 \cup Y_2$ de deux parties fermées correspond $j_+(Y_1) \cap j_+(Y_2)$, et à l'intersection d'une famille quelconque $(Y_\lambda)$ de parties fermées correspond la racine dans $S_+$ de la somme des $j_+(Y_\lambda)$.

\underline{Corollaire 2} .— Soit $J$ un idéal gradué de $S_+$; pour que $V_+(J)=\emptyset$, il faut et il suffit que tout élément de $S_+$ ait une puissance dans $J$.

Ce dernier corollaire s'exprime aussi sous la forme équivalente :

\underline{Corollaire 3} .— Soit $(f_\alpha)$ une famille d'éléments homogènes de $S_+$. Pour que les $D_+(f_\alpha)$ forment un recouvrement de $X=\operatorname{Proj}(S)$, il faut et il suffit que tout élément de $S_+$ ait une puissance dans l'idéal engendré par les $f_\alpha$.

\newpage

%% ===== Page 240 =====

\marginpar{II-51}

Corollaire 4. — Soient $(f_\alpha)$ une famille d'éléments homogènes de $S_+$, $f$ un élément de $S_+$. Les relations suivantes sont équivalentes :
(i) $D_+(f) \subset \bigcup_\alpha D_+(f_\alpha)$ ; (ii) $\bigcap_\alpha V_+(f_\alpha) \subset V_+(f)$ ; (iii) une puissance de $f$ appartient à l'idéal engendré par les $f_\alpha$.

Corollaire 5. — Pour que $X = \operatorname{Proj}(S)$ soit vide, il faut et il suffit que tout élément de $S_+$ soit nilpotent.

Corollaire 6. — Dans la correspondance biunivoque décrite dans le cor. 1, aux parties fermées irréductibles de $X$ correspondent les idéaux premiers gradués dans $S_+$.

En effet, si $Y = Y_1 \cup Y_2$, où $Y_1$ et $Y_2$ sont fermés et distincts de $Y$, on a $j_+(Y) = j_+(Y_1) \cap j_+(Y_2)$, les idéaux $j_+(Y_1)$ et $j_+(Y_2)$ étant distincts de $j_+(Y)$, donc $j_+(Y)$ n'est pas premier. Inversement, si $J$ est un idéal gradué non premier de $S_+$, il existe deux éléments $f, g$ de $S_+$ tels que $f \notin J$, $g \notin J$, $fg \in J$ ; alors $V_+(f) \not\subset V_+(J)$, $V_+(g) \not\subset V_+(J)$ mais $V_+(J) \subset V_+(f) \cup V_+(g)$ par (3) ; on en conclut que $V_+(J)$ est réunion des ensembles fermés $V_+(f) \cap V_+(J)$ et $V_+(g) \cap V_+(J)$, qui sont distincts de $V_+(J)$.

Soient $f, g$ deux éléments homogènes de $S_+$; considérons les schémas affines $Y_f = \operatorname{Spec}(S_{(f)})$, $Y_g = \operatorname{Spec}(S_{(g)})$ et $Y_{fg} = \operatorname{Spec}(S_{(fg)})$. En vertu du lemme 1, le morphisme $\psi_{fg,f} : Y_{fg} \to Y_f$, correspondant à l'homomorphisme canonique $\varphi_{fg,f} : S_{(f)} \to S_{(fg)}$, est une immersion ouverte (I, 1, 3, prop. 6). Au moyen de l'homomorphisme réciproque de $\psi_{f} : D_+(f) \to Y_f$ (prop. 4), on peut transporter à $D_+(f)$ la structure de schéma affine de $Y_f$; en vertu de la commutativité du diagramme (8), le schéma affine $D_+(fg)$ s'identifie alors au schéma induit sur l'ensemble ouvert $D_+(fg)$ de l'espace de base par le schéma $D_+(f)$. Il est clair alors que $\operatorname{Proj}(S)$ est muni d'une unique structure de préschéma, dont la restriction à chaque $D_+(f)$ est le préschéma qui vient d'être défini. En outre :

\newpage

%% ===== Page 241 =====

\noindent -- II-52 --

\noindent \emph{Proposition 6. --- Le préschéma $\operatorname{Proj}(S)$ est un schéma.}

\noindent Il suffit (I,3,6,prop.20) de montrer que quels que soient $f,g$ homogènes dans $S_+$, $D_+(f) \cap D_+(g) = D_+(fg)$ est affine et que son anneau est engendré par les images canoniques des anneaux de $D_+(f)$ et $D_+(g)$ ; le premier point est évident par définition, et le second résulte de ce que $S_{(fg)}$ est engendré par les images canoniques de $S_{(f)}$ et $S_{(g)}$ ($n^\circ 1$ lemme 2).

\noindent Lorsqu'on parlera du spectre premier homogène $\operatorname{Proj}(S)$ comme d'un schéma, il s'agira toujours de la structure qui vient d'être définie.

\noindent \emph{Proposition 7. --- Si l'anneau $S$ est intègre, le schéma $\operatorname{Proj}(S)$ est réduit et son espace de base est irréductible.}

\noindent Comme $(0)$ est alors un idéal premier gradué, l'irréductibilité de l'espace de base $\operatorname{Proj}(S)$ résulte du cor. 6 de la prop. 5. D'autre part pour $f$ homogène dans $S_+$, $S_{(f)}$ n'a \emph{a fortiori} pas d'élément nilpotent, donc $D_+(f)$ est un schéma réduit ; il en est de même de $X = \operatorname{Proj}(S)$ (I,3,4, déf. 5).

\noindent Etant donné un anneau commutatif $A$, rappelons qu'on dit qu'un anneau gradué $S$ est une $A$-algèbre graduée s'il est muni d'une structure de $A$-algèbre telle que chacun des sous-groupes $S_n$ ($n \ge 0$) soit un $A$-module ; il suffit d'ailleurs pour cela que $S_0$ soit une $A$-algèbre, autrement dit on définit sur $S$ une structure de $A$-algèbre graduée en définissant une structure de $A$-algèbre sur $S_0$ et en définissant pour $a \in A$, $x \in S_n$, $a \cdot x = (a \cdot 1) x$.

\noindent \emph{Proposition 8. --- Supposons que $S$ soit une $A$-algèbre graduée. Alors sur $X = \operatorname{Proj}(S)$, le faisceau $\mathcal{O}_X$ est un faisceau de $A$-algèbres ($A$ étant considéré comme faisceau simple sur $X$) ; d'autres termes, $X$ est un schéma au-dessus de $Y = \operatorname{Spec}(A)$.}

\noindent Il suffit de noter que pour tout $f$ homogène dans $S_+$, $S_{(f)}$ est une

\newpage

%% ===== Page 242 =====

-- II-53 --

algèbre sur $A$, et que le diagramme
$$\begin{tikzcd}
S_{(f)} \arrow[r] & S_{(fg)} \\
& A \arrow[u] \arrow[ul]
\end{tikzcd}$$
est commutatif pour $f, g$ homogènes dans $S_+$.

Remarque. — Soit $d$ un entier $> 0$ ; le sous-anneau $S^{(d)}$ de $S$ peut être considéré comme anneau gradué en prenant $S_{nd}$ comme ensemble des éléments homogènes de degré $n$ dans cet anneau. L'application $\mathfrak{p} \to \mathfrak{p} \cap S^{(d)}$ est une bijection de l'ensemble $\operatorname{Proj}(S)$ sur $\operatorname{Proj}(S^{(d)})$. En effet, supposons donné un idéal premier gradué $\mathfrak{p}' \in \operatorname{Proj}(S^{(d)})$, et posons $\mathfrak{p}_{nd} = \mathfrak{p}' \cap S_{nd}$. Pour tout $n$, définissons $\mathfrak{p}_n$ comme l'ensemble des $x \in S_n$ tels que $x^d \in \mathfrak{p}_{nd}$ ; si $x \in \mathfrak{p}_n, y \in \mathfrak{p}_n$, on a $(x+y)^{2d} \in \mathfrak{p}_{2nd}$, donc $(x+y)^d \in \mathfrak{p}_{nd}$ puisque $\mathfrak{p}'$ est un idéal premier. On en conclut aussitôt que les $\mathfrak{p}_n$ vérifient les conditions de la prop. 2 du $n^\circ 1$, et par suite il existe un seul idéal premier $\mathfrak{p} \in \operatorname{Proj}(S)$ tel que $\mathfrak{p} \cap S^{(d)} = \mathfrak{p}'$. De plus, comme pour tout $f$ homogène dans $S_+$, on a $V_+(f) = V_+(f^d)$, on voit que $\operatorname{Proj}(S)$ et $\operatorname{Proj}(S^{(d)})$ s'identifient en tant qu'espaces topologiques ; enfin, $S_{(f)}$ et $S_{(f^d)}$ s'identifient également en vertu du lemme 1 du $n^\circ 1$, donc $\operatorname{Proj}(S)$ et $\operatorname{Proj}(S^{(d)})$ s'identifient en tant que schémas.

Désignons par $S'$ l'anneau dont le groupe additif est la somme directe de $\mathbb{Z}$ et de $S_+$, la multiplication dans $S_+$ étant la même que dans l'anneau $S$, et le produit $n.x$ pour $x \in S_n$ étant la somme de $n$ termes égaux à $x$ ; il est clair que $S'$ est un anneau gradué avec $S'_0 = \mathbb{Z}$ et $S'_n = S_n$ pour $n > 0$. Si à tout idéal premier $\mathfrak{p} \in \operatorname{Proj}(S)$ on fait correspondre l'idéal premier $\mathfrak{p}' \in \operatorname{Proj}(S')$ tel que $\mathfrak{p}'_n = \mathfrak{p}_n$ pour tout $n > 0$ (prop. 2, $n^\circ 1$), on voit aussitôt qu'on définit un isomorphisme canonique identifiant les

\newpage

%% ===== Page 243 =====

schémas $\operatorname{Proj}(S)$ et $\operatorname{Proj}(S')$.

De même, si $S$ est une $A$-algèbre graduée, on peut remplacer $S$ par l'algèbre graduée somme directe de $A$ et des $S_n$ ($n>0$) avec $(S'_A)_0 = A$ et $(S'_A)_n = S_n$, et on voit comme ci-dessus que les schémas $\operatorname{Proj}(S)$ et $\operatorname{Proj}(S'_A)$ s'identifient canoniquement.

3. \emph{Faisceau associé à un module gradué}.

Soit $M$ un $S$-module gradué. Pour tout $f$ homogène dans $S_+$, $M_{(f)}$ est un $S_{(f)}$-module, et il lui correspond donc un faisceau quasi-cohérent $(M_{(f)})^{\sim}$ associé (I, 1, 3) de modules sur le schéma affine $\operatorname{Spec}(S_{(f)})$, identifié à $D_+(f)$.

Proposition 9. -- Il existe sur $X = \operatorname{Proj}(S)$ un $\mathcal{O}_X$-module quasi-cohérent et un seul tel que pour tout $f$ homogène dans $S_+$, on ait $\Gamma(D_+(f), \widetilde{M}) = M_{(f)}$, l'homomorphisme de restriction $\Gamma(D_+(f), \widetilde{M}) \to \Gamma(D_+(fg), \widetilde{M})$ pour $f, g$ homogènes dans $S_+$, étant l'homomorphisme canonique $M_{(f)} \to M_{(fg)}$ ($n^{\circ}1$).

En effet, la restriction à $D_+(fg)$ du faisceau $(M_{(f)})^{\sim}$ sur $D_+(f)$ s'identifie canoniquement ($n^{\circ}1$, lemme 1) au faisceau associé au module $(M_{(f)})_{(g/f^d)}$ ; puisque $D_+(fg)$ s'identifie au spectre premier de $(S_{(f)})_{(g/f^d)}$ (I, 1, 3, prop. 6) ; mais d'après le lemme 1 du $n^{\circ}1$, on voit que la restriction à $D_+(fg)$ de $(M_{(f)})^{\sim}$ s'identifie canoniquement à $(M_{(fg)})^{\sim}$ ; on en conclut qu'il existe un isomorphisme $\theta_{g,f}$ de la restriction à $D_+(fg)$ du faisceau $(M_{(f)})^{\sim}$ sur la restriction à $D_+(fg)$ du faisceau $(M_{(g)})^{\sim}$ ; en outre, si $h$ est un troisième élément homogène de $S_+$, on a $\theta_{g,f} \circ \theta_{f,h} \circ \theta_{h,g} = 1$ dans $D_+(fgh)$. On en conclut (FAC, I, 1, 4) qu'il existe sur $X$ un $\mathcal{O}_X$-module $F$ et pour chaque $f$ homogène dans $S_+$, un isomorphisme $\eta_f$ de $F|_{D_+(f)}$ sur $(M_{(f)})^{\sim}$, tels que $\theta_{g,f} \circ \eta_g \circ \eta_f^{-1} = 1$. Si alors on considère le faisceau associé au préfaisceau formé des $\Gamma(D_+(f)) = M_{(f)}$, avec les homomorphismes de restriction canoniques $M_{(f)} \to M_{(fg)}$ comme homomorphismes de restriction, ce qui précède prouve que $F$ et $G$ sont isomorphes, et par suite (compte

\newpage

%% ===== Page 244 =====

-- 55 --

Définition 1. — On dit que le faisceau $\widetilde{M}$ défini dans la prop. 9 est le faisceau associé au $S$-module gradué $M$.

Rappelons que les $S$-modules gradués forment une catégorie lorsqu'on restreint les homomorphismes de modules à ceux qui sont de degré 0. Avec cette convention :

Proposition 10. — Le foncteur $M \to \widetilde{M}$ est un foncteur covariant additif exact de la catégorie des $S$-modules gradués dans la catégorie des $\mathcal{O}_X$-modules, qui commute aux limites inductives et aux sommes directes.

En effet, ces propriétés étant locales, il suffit de les vérifier sur les faisceaux $\widetilde{M}|_{D_+(f)} = M_{(f)}$. Or les foncteurs $M \to M_{\mathfrak{p}}$, $M \to (M)_{\mathfrak{p},0} = M_{(f)}$ et $M \to (M)_{(f)}^\sim$ ont tous trois les propriétés d'exactitude et de permutabilité aux limites inductives et sommes directes (I, 1, 3, cor. 3 et 4 du th. 1) ; d'où la proposition.

Proposition 11. — Pour tout $\mathfrak{p} \in X = \operatorname{Proj}(S)$, on a $\widetilde{M}_{\mathfrak{p}} = M_{(\mathfrak{p})}$.

On a en effet $\widetilde{M}_{\mathfrak{p}} = \varinjlim \Gamma(D_+(f), \widetilde{M})$, où $f$ parcourt les éléments homogènes de $S_+ \setminus \mathfrak{p}$ ; comme $\Gamma(D_+(f), \widetilde{M}) = M_{(f)}$, la proposition résulte de la définition de $M_{(\mathfrak{p})}$ (n°1).

Proposition 12. — On suppose $S$ engendré par l'ensemble $S_1$ des éléments homogènes de degré 1. Pour que $\widetilde{M}=0$, il faut et il suffit que pour tout $z \in M$ et tout $f \in S_+$, il existe une puissance de $f$ annulant $z$.

En effet, la condition $\widetilde{M}=0$ équivaut à $M_{(f)}=0$ pour tout $f$ homogène dans $S_+$, et la condition est donc suffisante (sans hypothèse sur $S_1$).

Inversement, si $\widetilde{M}=0$ pour tout $f \in S_1$ homogène, pour tout $z \in M$ il existe par définition une puissance $f^n$ telle que $f^n z=0$. L'hypothèse sur $S_1$ entraîne alors que cette condition est aussi sa-

\newpage

%% ===== Page 245 =====

-- 56 --

Si $M$ est un $S$-module gradué, et $n$ un entier rationnel, nous désignerons par $M(n)$ le $S$-module gradué défini par $(M(n))_k = M_{n+k}$ pour tout $k \in \mathbb{Z}$. On définit ainsi en particulier, pour tout $n \in \mathbb{Z}$, un $S$-module gradué $S(n)$ (il faut prendre $S_m=(0)$ pour $m < 0$). On dit qu'un $S$-module gradué est libre s'il est isomorphe, en tant que module gradué, à une somme directe de modules de la forme $S(n)$ ; il revient au même de dire que $M$ admet une base sur $S$ formée d'éléments homogènes.

Proposition 13. — Soient $d$ un entier $> 0$, $f \in S_d$. Alors les faisceaux $(S(nd))^{\sim} | D_+(f)$ et $\mathcal{O}_X | D_+(f)$ sont isomorphes pour tout $n \in \mathbb{Z}$.

En effet, la multiplication par l'élément inversible $(f^n/1)$ est une bijection de $S_f$ sur lui-même, donc une bijection de $S(f) = (S_f)_{(0)}$ sur $(S(nd))_{(f)} = (S(nd))_{(f)}$ et $S(nd)_{(f)}$ sont par suite isomorphes, d'où la proposition.

Corollaire 1. — Le faisceau $(S(nd))^{\sim}$ est un faisceau inversible sur l'ensemble ouvert $\bigcup_{f \in S_d} D_+(f)$.

Corollaire 2. — Si $S_+$ est engendré par l'ensemble $S_1$ des éléments homogènes de rang $1$, les faisceaux $(S(n))^{\sim}$ sont inversibles pour tout $n \in \mathbb{Z}$.

Il suffit d'appliquer le cor. 1 avec $d=1$, en remarquant que $X = \bigcup_{f \in S_1} D_+(f)$.

Nous poserons dans ce qui suit
$$(9) \quad \mathcal{O}_X(n) = (S(n))^{\sim} \text{ pour tout } n \in \mathbb{Z},$$
et pour tout faisceau algébrique $\mathcal{F}$ sur une partie ouverte $U$ de $X$,
$$(10) \quad \mathcal{F}(n) = \mathcal{F} \otimes_{\mathcal{O}_X} (\mathcal{O}_X(n)|U) \text{ pour tout } n \in \mathbb{Z}.$$

Soient $M, N$ deux $S$-modules gradués ; rappelons que $M \otimes_S N$ est un $S$-module gradué par la définition
$$(M \otimes_S N)_q = \sum_{m+n=q} (M_m \otimes_{S_0} N_n)$$

\newpage

%% ===== Page 246 =====

- 11-57 -

où $\varphi_{mn}$ est le produit tensoriel des injections $M_n \to M$, $N_n \to N$. Pour tout $f \in S_d$ ($d>0$), on définit un homomorphisme canonique fonctoriel de $S_{(f)}$-modules
$$(11) \quad \lambda_f : M_{(f)} \otimes_{S_{(f)}} N_{(f)} \to (M \otimes_S N)_{(f)}$$
en composant l'homomorphisme $M_{(f)} \otimes_{S_{(f)}} N_{(f)} \to M_f \otimes_{S_f} N_f$ (provenant des injections $M_{(f)} \to M_f$, $N_{(f)} \to N_f$ et $S_{(f)} \to S_f$) et l'isomorphisme canonique $M_f \otimes_{S_f} N_f \to (M \otimes_S N)_f$ (chap. 0, § 1, n°3) et en notant que ce dernier conserve les degrés ; pour $x \in M_{nd}$, $y \in N_{n'd}$, on a donc $\lambda_f(x/f^n \otimes y/f^{n'}) = (x \otimes y) / f^{n+n'}$.

Il résulte aussitôt de cette définition que si $g \in S_{d'}$ ($d'>0$), le diagramme
$$\begin{tikzcd} M_{(f)} \otimes_{S_{(f)}} N_{(f)} \arrow[r, "\lambda_f"] \arrow[d] & (M \otimes_S N)_{(f)} \arrow[d] \\ M_{(fg)} \otimes_{S_{(fg)}} N_{(fg)} \arrow[r, "\lambda_{fg}"] & (M \otimes_S N)_{(fg)} \end{tikzcd}$$
où la flèche verticale de droite est l'homomorphisme canonique, et la flèche verticale de gauche provient des homomorphismes canoniques, est commutatif. On en conclut que les $\lambda_f$ définissent un homomorphisme canonique fonctoriel de $\mathcal{O}_X$-modules
$$(12) \quad \lambda : \tilde{M} \otimes_{\mathcal{O}_X} \tilde{N} \to (M \otimes_S N)^\sim.$$

Rappelons que $\operatorname{Hom}_S(M, N)$ désigne l'ensemble des homomorphismes (de tous les degrés) de modules gradués $M \to N$, et qu'il devient un $S$-module gradué lorsqu'on prend pour $(\operatorname{Hom}_S(M, N))_n$ l'ensemble des homomorphismes de degré $n$. Pour tout $f \in S_d$ ($d>0$), on définit un homomorphisme canonique fonctoriel de $S_{(f)}$-modules
$$(13) \quad \rho_f : (\operatorname{Hom}_S(M, N))_{(f)} \to \operatorname{Hom}_{S_{(f)}}(M_{(f)}, N_{(f)})$$
en faisant correspondre à $u/f^n$, où $u$ est un homomorphisme de degré $nd$, l'homomorphisme $M_{(f)} \to N_{(f)}$ qui transforme $x/f^k$ en $u(x)/f^{k+n}$. On a encore, pour $g \in S_{d'}$ ($d'>0$), un diagramme commutatif

\marginpar{$\psi_{mn}$}

\newpage

%% ===== Page 247 =====

\marginpar{Archives Grothendieck sept 59}

$$- II-58 -$$

$$\begin{tikzcd}
(\operatorname{Hom}_S(M,N))_{(f)} \arrow[d] \arrow[r, "\rho_f"] & \operatorname{Hom}_{S_{(f)}}(M_{(f)}, N_{(f)}) \arrow[d] \\
(\operatorname{Hom}_S(M,N))_{(fg)} \arrow[r, "\rho_{fg}"] & \operatorname{Hom}_{S_{(fg)}}(M_{(fg)}, N_{(fg)})
\end{tikzcd}$$

(la flèche de gauche étant l'homomorphisme canonique, celle de droite provenant des homomorphismes canoniques). On en conclut que les $\rho_f$ définissent un homomorphisme fonctoriel de $\mathcal{O}_X$-modules
(14) $\quad \rho : (\underline{\operatorname{Hom}}_S(M,N))^{\sim} \to \underline{\operatorname{Hom}}_{\mathcal{O}_X}(\tilde{M}, \tilde{N})$.

\textbf{Proposition 14.} - Supposons $S$ engendré par $S_1$. Alors $\lambda$ est un isomorphisme ; il en est de même de $\mu$ si on suppose en outre que $M$ soit le conoyau d'un homomorphisme d'un $S$-module gradué libre de rang fini dans un autre (ce qui sera le cas lorsque $S$ est noethérien et $M$ un $S$-module gradué de type fini).

Comme $X$ est réunion des $D_+(f)$, où $f \in S_+$, on est ramené à prouver que $\lambda_f$ et $\mu_f$ sont des isomorphismes, sous les hypothèses envisagées, lorsque $f$ est homogène de degré $1$. Or, on définit alors une application $\mathbb{Z}$-bilinéaire $M_n \times N_m \to M_{(f)} \otimes_{S_{(f)}} N_{(f)}$ en faisant correspondre à $(x,y)$ l'élément $(x/f^n) \otimes (y/f^m)$ ; ces applications définissent donc une application $\mathbb{Z}$-linéaire $M \otimes_S N \to M_{(f)} \otimes_{S_{(f)}} N_{(f)}$, et si $s \in S_h$, cette application transforme $(sx) \otimes y$ en $(s/f^h)((x/f^n) \otimes (y/f^m))$ (si $x \in M_n, y \in N_m$). On en déduit donc un homomorphisme $\rho_f : M \otimes_S N \to M_{(f)} \otimes_{S_{(f)}} N_{(f)}$ relatif à l'homomorphisme d'anneaux $S \to S_{(f)}$ qui transforme $s$ en $s/f^h$. Supposons en outre que pour un élément $\sum (x_i \otimes y_i)$ de $M \otimes N$ ($x_i, y_i$ homogènes), on ait $f^r \sum (x_i \otimes y_i) = 0$, autrement dit $\sum (f^r x_i \otimes y_i) = 0$. Si $x_i$ (resp. $y_i$) est de degré $n_i$ (resp. $p_i$), on en déduit $\sum (f^r x_i / f^{n_i+r}) \otimes (y_i / f^{p_i}) = 0$, c'est-à-dire $\rho_f (\sum (x_i \otimes y_i)) = 0$. Par suite $\rho_f$ se factorise en $M \otimes_S N \to (M \otimes_S N)_{(f)} \xrightarrow{\rho'_f} M_{(f)} \otimes_{S_{(f)}} N_{(f)}$ ; si $\lambda'_f$ est la restriction de $\rho'_f$ à $(M \otimes_S N)_{(f)}$, on vérifie aussitôt que $\lambda_f$ et $\lambda'_f$ sont des $S_{(f)}$-homomorphismes réciproques, d'où la première partie de la proposition.

Pour démontrer la seconde partie, supposons que $M$ soit le conoyau

\newpage

%% ===== Page 248 =====

- 11-59 -

d'un homomorphisme de modules gradués $P \to Q$, où $P$ et $Q$ sont des sommes directes de modules de la forme $S(n)$ ; utilisant l'exactitude à droite des foncteurs $L \to \operatorname{Hom}_S(L, N)$ et $M \to M_{(f)}$ ; on est aussitôt ramené à démontrer que $\mu_f$ est un isomorphisme lorsque $M=S(n)$. Or, pour tout $z \in N$, soit $u_z$ l'homomorphisme de $S(n)$ dans $N$ tel que $u_z(1)=z$ ; on voit aussitôt que $\sigma : z \to u_z$ est un isomorphisme homogène de degré 0 de $N(-n)$ sur $\operatorname{Hom}_S(S(n), N)$. Il lui correspond un isomorphisme $\sigma_f : (N(-n))_{(f)} \xrightarrow{\sim} (\operatorname{Hom}_S(S(n), N))_{(f)}$.
Soit d'autre part $\varphi_f$ l'isomorphisme $N_{(f)} \xrightarrow{\sim} \operatorname{Hom}_{S_{(f)}}((S(n))_{(f)}, N_{(f)})$ qui à tout $z' \in N_{(f)}$ fait correspondre l'homomorphisme $v_{z'}$, tel que $v_{z'}(s/f^k) = sz'/f^{n+k}$ ($s \in S_{n+k} = (S(n))_k$) ; on constate aisément que l'application composée
$$(N(-n))_{(f)} \xrightarrow{\sigma_f} (\operatorname{Hom}_S(S(n), N))_{(f)} \xrightarrow{\varphi_f^{-1}} N_{(f)}$$
est l'isomorphisme $z/f^k \to z/f^{k-n}$ de $(N(-n))_{(f)}$ sur $N_{(f)}$, donc $\mu_f$ est un isomorphisme.

Corollaire 1. - Quels que soient $n, m$ dans $\mathbb{Z}$, on a
$$(15) \quad \mathcal{O}_X(n) \otimes \mathcal{O}_X \mathcal{O}_X(m) = \mathcal{O}_X(n+m)$$
à un isomorphisme canonique près.

Cela résulte de la prop. 14 et de l'existence de l'isomorphisme canonique de degré 0 $S(m) \otimes_S S(n) \to S(n+m)$ qui à l'élément $1 \otimes 1$ ($1 \in (S(m))_{-m}$, $1 \in (S(n))_{-n}$) fait correspondre l'élément $1 \in (S(n+m))_{-(n+m)}$.
En particulier, pour tout $n \in \mathbb{Z}$, on a
$$(15) \quad \mathcal{O}_X(n) = (\mathcal{O}_X(1))^{\otimes n}$$
avec la convention introduite au $\S$ 1, n° 3, lorsque $n \leq 0$.

Corollaire 2. - Pour tout $S$-module gradué $M$ et tout $n \in \mathbb{Z}$, on a
$$(16) \quad (M(n))^{\sim} = \tilde{M}(n)$$
à un isomorphisme canonique près.

Cela résulte de la prop. 14, des formules (9) et (10), et de l'isomorphisme canonique de degré 0, $M(n) \to M \otimes_S S(n)$ qui à tout $m \in (M(n))_h = M_{n+h}$ fait correspondre $m \otimes 1 \in (M \otimes_S S(n))_h$.

\newpage

%% ===== Page 249 =====

- II-60 -

Dans toute la fin de ce n° nous supposons $S_+$ engendré par l'ensemble $S_1$ des éléments homogènes de degré 1, de sorte que $\mathcal{O}_X(1)$ est un faisceau inversible (cor. 2 de la prop. 13). Pour tout $\mathcal{O}_X$-module $F$, on posera (§ 1, n°3)
$$(17) \qquad \Gamma_*(F) = \Gamma_*(O_X(1), F) = \bigoplus_{n \in \mathbb{Z}} \Gamma(X, F(n))$$;
rappelons en particulier que $\Gamma_*(\mathcal{O}_X)$ est muni d'une structure d'anneau gradué (compte tenu de (15)) et $\Gamma_*(F)$ d'une structure de module gradué sur $\Gamma_*(\mathcal{O}_X)$.

Soit $M$ un $S$-module gradué ; pour tout $f \in S_d$ ($d > 0$), $x \to x/1$ est un homomorphisme de groupes abéliens $M_0 \to M_{(f)}$ ; et comme $M_{(f)}$ s'identifie canoniquement à $\Gamma(D_+(f), \widetilde{M})$, on obtient ainsi un homomorphisme de groupes abéliens $\alpha_f : M_0 \to \Gamma(D_+(f), \widetilde{M})$. Il est clair que le diagramme

\begin{tikzcd}
& M_0 \arrow[dl, "\alpha_f"'] \arrow[dr, "\alpha_{fg}"] & \\
\Gamma(D_+(f), \widetilde{M}) \arrow[rr] & & \Gamma(D_+(fg), \widetilde{M})
\end{tikzcd}

est commutatif pour tout $g \in S_{d'}$ ($d' > 0$) ; cela signifie que pour tout $x \in M_0$, les sections $\alpha_f(x)$ et $\alpha_{fg}(x)$ coïncident dans $D_+(f) \cap D_+(g)$ et par suite il existe une section unique $\alpha_0(x) \in \Gamma(X, \widetilde{M})$ dont la restriction à chaque $D_+(f)$ est $\alpha_f(x)$. On a ainsi défini un homomorphisme de groupes abéliens $\alpha_0 : M_0 \to \Gamma(X, \widetilde{M})$. Appliquant ce résultat au $S$-module gradué $M(n)$ ($n \in \mathbb{Z}$), on obtient pour chaque $n$ un homomorphisme de groupes abéliens $\alpha_n : M_n = (M(n))_0 \to \Gamma(X, \widetilde{M}(n))$ (compte tenu du cor. 2 de la prop. 14) ; d'où un homomorphisme (de degré 0) de groupes abéliens gradués :
$$(18) \qquad \alpha : M \to \Gamma_*(\widetilde{M})$$
qui dans chaque $M_n$ coïncide avec $\alpha_n$.

Si on prend en particulier $M=S$, on vérifie aussitôt (compte tenu de la définition de la multiplication dans $\Gamma_*(\mathcal{O}_X)$) que $\alpha : S \to \Gamma_*(\mathcal{O}_X)$ est un homomorphisme (de degré 0) d'anneaux.

\newpage

%% ===== Page 250 =====

-- II-61 --

\emph{neaux gradués}, et pour tout $S$-module gradué $M$, (18) est un di-
homomorphisme de modules gradués.

\noindent Proposition 15. — Pour tout $f \in S_d$ ($d>0$), $D_+(f)$ est identique à l'ensemble des $p \in X$ où la section $\alpha_d(f)$ de $\mathcal{O}_X(d)$ ne s'annule pas (§ 1, n°3)

Comme $X = \bigcup_{g \in S_1} D_+(g)$, il suffit de vérifier que, pour tout $g \in S_1$, l'ensemble des $p \in D_+(g)$ où $\alpha_d(f)$ ne s'annule pas est identique à $D_+(fg)$. Or, la restriction de $\alpha_d(f)$ à $D_+(g)$ est la section correspondante à l'élément $f/1$ de $S_{(g)}$; par l'isomorphisme canonique $S_{(g)} \cong S_{(g^{(d)})}$ (prop. 13), cette section de $\mathcal{O}_X(d)$ au-dessus de $D_+(g)$ s'identifie à la section de $\mathcal{O}_X$ au-dessus de $D_+(g)$ qui correspond à l'élément $f/g^d$ de $S_{(g)}$; dire que cette section s'annule en $p \in D_+(g)$ signifie que $f/g^d \in \mathfrak{q}$, où $\mathfrak{q}$ est l'idéal premier de $S_{(g)}$ correspondant à $p$ (n°2, prop. 4); par définition cela veut dire que $f \in \mathfrak{p}$, d'où la proposition.

Soit maintenant $F$ un $\mathcal{O}_X$-module, et posons $M = \Gamma_*(F)$; en vertu de l'existence de l'homomorphisme $\alpha: S \to \Gamma_*(\mathcal{O}_X)$, $M$ peut être considéré comme un $S$-module gradué. Pour tout $f \in S_d$ ($d>0$), il résulte de la prop. 15 que la restriction à $D_+(f)$ de la section $\alpha_d(f)$ de $\mathcal{O}_X(d)$ est inversible; il en est de même par suite de même de la restriction à $D_+(f)$ de la section $\alpha_{nd}(f^n)$ de $\mathcal{O}_X(nd)$, pour tout $n>0$. Soit alors $z \in M_{nd} = \Gamma(X, F(nd))$; s'il existe un entier $k \ge 0$ tel que la restriction à $D_+(f)$ de $f^k z$, c'est-à-dire la section $(z|_{D_+(f)})(\alpha_{dk}(f^k)|_{D_+(f)})$ de $F((n+k)d)$ soit nulle, alors en raison de la remarque précédente, on a aussi $z|_{D_+(f)} = 0$. Cela montre qu'on définit un $S_{(f)}$-homomorphisme $\beta_f : M_{(f)} \to \Gamma(D_+(f), F)$ en faisant correspondre à l'élément $z/f^n$ la section $(z|_{D_+(f)})(\alpha_{nd}(f^n)|_{D_+(f)})^{-1}$ de $F$ au-dessus de $D_+(f)$. On vérifie en outre aussitôt que le diagramme

\begin{tikzcd}
M_{(f)} \arrow[r, "\beta_f"] \arrow[d] & \Gamma(D_+(f), F) \arrow[d] \\
M_{(fg)} \arrow[r, "\beta_{fg}"] & \Gamma(D_+(fg), F)
\end{tikzcd}

\newpage

%% ===== Page 251 =====

\begin{center}
-- II-62 --
\end{center}

est commutatif ($g \in S_d, d > 0$). Si l'on se rappelle que $M_{(f)}$ s'identifie canoniquement à $\Gamma(D_+(f), \tilde{M})$ et que les $D_+(f)$ forment une base de la topologie de $X$ (n$^\circ$2, prop.3), on voit que les $\beta_f$ proviennent d'un unique homomorphisme canonique de $\mathcal{O}_X$-modules
$$(19) \quad \beta : (\Gamma_*(F))^\sim \to F$$
qui est évidemment fonctoriel.

\subsection*{4. Conditions de finitude.}

\textbf{Lemme 3.} -- Soit $S$ un anneau gradué à degrés positifs. Pour que $S_+$ soit un idéal de type fini, il faut et il suffit que $S$ soit une $S_0$-algèbre de type fini.

Si $S$ est une $S_0$-algèbre de type fini, on peut écrire $S = S_0[f_1, \dots, f_n]$ où on peut supposer les $f_i$ homogènes et de degrés $> 0$ ; il est clair alors que les $f_i$ forment un système de générateurs du $S$-module $S_+$. Inversement, si $S_+$, en tant que $S$-module, admet un système fini de générateurs $g_j$ ($1 \le j \le m$), on peut supposer les $g_j$ homogènes de degrés $> 0$ ; montrons alors que $S = S_0[g_1, \dots, g_m]$. En effet, il suffit de voir, par récurrence sur le degré $d$ d'un élément homogène $f \in S_+$, que $f$ est un polynôme par rapport aux $g_j$ à coefficients dans $S_0$. Or, on a par hypothèse $f = \sum c_j g_j$ ; où $c_j \in S$, et dans cette relation, on peut évidemment remplacer $c_j$ par sa composante homogène de degré $d - \deg(g_j)$ ; l'hypothèse de récurrence entraîne alors le résultat.

\textbf{Lemme 4.} -- Soit $S$ un anneau gradué à degrés positifs, tel que $S_+$ soit un idéal de type fini. Dans ces conditions :

(i) les $S_n$ ($n \ge 0$) sont des $S_0$-modules de type fini ;

(ii) pour tout $S$-module de type fini $M$, et tout $d > 0$, $M^{(d)}$ est un $S^{(d)}$-module de type fini ;

(iii) pour tout $d > 0$, $S^{(d)}$ est une $S_0$-algèbre de type fini ;

(iv) pour tout $n > 0$, il existe un entier $m_0$ tel que $S_m \subset S_+^n$ pour $m \ge m_0$.

Soit $(f_i)_{1 \le i \le p}$ un système de générateurs homogènes de l'idéal $S_+$, qui est

\newpage

%% ===== Page 252 =====

-- II-63 --

(lemme 3), aussi un système de générateurs de la $S_0$-algèbre $S$ ; il est clair que si $n_i$ est le degré de $f_i$, les produits $f_1^{r_1} \dots f_p^{r_p}$ tels que $\sum_{i=1}^p r_i n_i = d$ forment un système de générateurs du $S_0$-module $S_d$, d'où (i).
Soient $x_j$ ($1 \le j \le n$) des générateurs du $S$-module $M$, qu'on peut supposer homogènes, et désignons par $m_j$ le degré de $x_j$ ; tout élément de $M^{(d)}$ est par définition combinaison linéaire à coefficients dans $S_0$ d'éléments de la forme $g_1^{r_1} \dots g_p^{r_p} x_j$, où $r_i < d$ ($1 \le i \le p$), et $m_j + \sum_{i=1}^p r_i n_i$ est un multiple de $d$ ; il n'y a qu'un nombre fini de systèmes $(r_i)$ vérifiant ces conditions, d'où (ii). De même tout élément de $S^{(d)}$ est une combinaison linéaire à coefficients dans $S_0$ de monômes dont chacun est produit d'un monôme par rapport aux $f_i$ et d'un monôme $f_1^{r_1} \dots f_p^{r_p}$, où $r_i < d$ ($1 \le i \le p$) et $\sum_{i=1}^p r_i n_i$ est un multiple de $d$ ; ces derniers monômes (en nombre fini) et les $f_i^d$ forment donc un système de générateurs de la $S_0$-algèbre $S^{(d)}$. Enfin, il n'y a qu'un nombre fini de systèmes $(r_i)$ tels que $\sum_{i=1}^p r_i n_i < n$, donc si $m_0$ est la plus grande valeur des sommes $\sum_{i=1}^p r_i n_i$ pour ces systèmes, on voit que l'on a $S_m \subset S_m^n$ pour tout $m \ge m_0$.

\underline{Proposition 16.} — (i) Si $S$ est un anneau gradué noethérien, $X = \operatorname{Proj}(S)$ est un schéma noethérien.

(ii) Si $S$ est une $A$-algèbre graduée de type fini, $X = \operatorname{Proj}(S)$ est un schéma de type fini au-dessus de $Y = \operatorname{Spec}(A)$.

(i) Si $S$ est noethérien, l'idéal $S_+$ admet un système fini de générateurs homogènes $f_i$ ($1 \le i \le n$), donc (n°2, cor. 3 de la prop. 5), l'espace de base $X$ est réunion des $D_+(f_i) = \operatorname{Spec}(S_{(f_i)})$, et tout revient à voir que chacun des $S_{(f_i)}$ est noethérien ; en vertu de la prop. 1 du n°1, on est ramené à prouver que les anneaux $S^{(d)}$ ($d > 0$) sont noethériens. Or l'anneau $S = S/S_+$ est noethérien, et $S^{(d)}$ est une $S_0$-algèbre de type fini (lemme 4), d'où la conclusion.
(ii) L'hypothèse entraîne que $S_0$ est une $A$-algèbre de type fini, et $S$ une $S_0$-algèbre de type fini, donc $S_+$ est un idéal de type fini.

\newpage

%% ===== Page 253 =====

— II-64 —

(lemme 3) ; comme donc (i), on est ramené à prouver que les $S^{(d)}$ sont des $A$-algèbres de type fini (I, 4, 2, déf. 2) ; mais cela résulte du lemme 4 et du fait que $S_0$ est une $A$-algèbre de type fini.

Nous considérerons dans ce qui suit les conditions de finitude suivantes pour un $S$-module gradué $M$ :

(TF) Il existe un entier $n$ tel que le sous-module $\bigoplus_{k \ge n} M_k$ soit un $S$-module de type fini.

(TN) Il existe un entier $n$ tel que $M_k = 0$ pour $k \ge n$.

\emph{Proposition 17.} — Soient $S$ un anneau gradué noethérien, $M$ un $S$-module gradué.

(i) Si $M$ vérifie la condition (TF), le $\mathcal{O}_X$-module $\widetilde{M}$ est cohérent.

(ii) Supposons que $M$ vérifie (TF) ; pour que $\widetilde{M} = 0$, il faut et il suffit que $M$ vérifie (TN).

La condition (TN) implique par définition $M_{(f)} = 0$ pour tout $f$ homogène dans $S_+$, donc $\widetilde{M} = 0$. Si $M$ vérifie (TF), le sous-module gradué $M' = \bigoplus_{k \ge n} M_k$, qui est par hypothèse de type fini, est tel que $M/M'$ satisfasse (TN) ; on a par suite $(M/M')^{\sim} = 0$, et l'exactitude du foncteur $M \mapsto \widetilde{M}$ (n° 3, prop. 10) entraîne $\widetilde{M} = \widetilde{M'}$; pour prouver que $\widetilde{M}$ est cohérent, on est donc ramené au cas où $M$ est de type fini. La question étant locale, il suffit, puisque $S_{(f)}$ est noethérien (lemme 4 et n° 1, prop. 1) de prouver que $M_{(f)}$ est un $S_{(f)}$-module de type fini (I, 1, 5, th. 3) ; mais en vertu des lemmes 3 et 4, $M^{(d)}$ est un $S^{(d)}$-module de type fini, et la prop. 1 du n° 1 établit alors notre assertion. Supposons en outre que $\widetilde{M} = 0$ (M vérifiant (TF)) ; alors on a aussi $\widetilde{M'} = 0$ et la condition (TN) pour $M'$ est identique à la condition (TN) pour $M$, donc pour prouver que $\widetilde{M} = 0$ implique que $M$ vérifie (TN), on peut encore se limiter au cas où $M$ est engendré par un nombre fini d'éléments homogènes $x_i$ ($1 \le i \le p$) ; soit d'autre part $(f_j)$ ($1 \le j \le q$) un système de générateurs homogènes de l'idéal $S_+$. On a par hypothèse $M_{(f_j)} = 0$ pour tout $j$, donc il existe un entier $n$ tel que $f_j^n x_i = 0$.

\newpage

%% ===== Page 254 =====

- II-65 -

quels que soient $i$ et $j$. Soit $n_j$ le degré de $f_j$, et soit $m$ la plus grande valeur de $\sum r_j n_j$ pour les systèmes $(r_j)$ en nombre fini tels que $\sum r_j \le n$; il est clair alors que si $k > n$ on a $S_k x_i = 0$ pour tout $i$; si $h$ est le plus grand des degrés des $x_i$, on en conclut que $M_k = 0$ pour $k > h + m$, ce qui termine la démonstration.

Corollaire. - Soit $S$ un anneau gradué noethérien; pour que $X = \operatorname{Proj}(S) = \emptyset$, il faut et il suffit qu'il existe $n$ tel que $S_k = 0$ pour $k \ge n$. La condition $X = \emptyset$ est en effet équivalente à $S_+ = 0$, et $S$ est un $S$-module monogène.

Théorème 1. - On suppose que $S$ est noethérien et engendré par l'ensemble $S_1$ des éléments homogènes de degré 1. Alors, pour tout $\mathcal{O}_X$-module quasi-cohérent $\mathcal{F}$, l'homomorphisme canonique $\beta : (\Gamma_*(\mathcal{F}))^{\sim} \to \mathcal{F}$ (formule (19)) est un isomorphisme.

En effet, l'espace de base $X$ est alors noethérien (prop. 16), et on a par suite, pour tout $f \in S_d$, un isomorphisme de $(\Gamma_*(\mathcal{F}))_{(f)}$ sur $\Gamma(D_+(f), \mathcal{F})$ (§ 1, n° 4, cor. 1 du th. 2), qui coïncide avec l'homomorphisme $\beta_f$ défini au n° 3 (grâce en particulier à la prop. 15), puisque $(\Gamma_*(\mathcal{F}))_{(f)}$ a été identifié à $(\mathcal{F}(D_+(f)))$; cela démontre le théorème.

Remarque. - Le cor. 1 du th. 2 du § 1, n° 4 est aussi valable si $X$ est un schéma dont l'espace de base est quasi-compact; le th. 1 est donc valable sous la seule hypothèse que $S_+$ est engendré par un nombre fini d'éléments $f_i \in S_1$, car alors $X$ est réunion des $S(f_i)$ qui sont quasi-compacts, et on sait d'autre part que $X$ est un schéma (n° 2, prop. 6).

Corollaire 1. - Tout faisceau quasi-cohérent $\mathcal{F}$ sur $X$ est isomorphe à un faisceau de la forme $\widetilde{M}$, où $M$ est un $S$-module gradué.

Corollaire 2. - Tout faisceau cohérent $\mathcal{F}$ sur $X$ est isomorphe à un faisceau de la forme $\widetilde{M}$, où $M$ est un $S$-module gradué de type fini.

\marginpar{On peut supposer $\mathcal{F} = \widetilde{M}$ (cor. 1).}
Soit $(f_\lambda)_{\lambda \in L}$ un système de générateurs homogènes de $S_+$, et

\newpage

%% ===== Page 255 =====

- 66 -

pour toute partie finie $H$ de $L$, soit $M_H$ le sous-module gradué de $M$ engendré par les $f_\lambda$ tels que $\lambda \in H$ ; il est clair que $M$ est limite inductive de ses sous-modules $M_H$ ; donc $F$ est limite inductive de ses sous-faisceaux $\widetilde{M}_H$ ($n^\circ 3$, prop. 10) . Mais comme $F$ est cohérent et $X$ noethérien (prop. 16), donc quasi-compact, il existe une partie finie $H$ de $L$ telle que $F = \widetilde{M}_H$ (chap. 0, § 2, $n^\circ$).

Corollaire 3. - Soit $F$ un faisceau cohérent sur $X$. Il existe alors un entier $n_0$ tel que, pour tout $n \ge n_0$, $F(n)$ soit isomorphe à un quotient d'un faisceau de la forme $\mathcal{O}_X^k$ ($k > 0$ dépendant de $n$).

En vertu du cor. 2, on peut supposer $F = \widetilde{M}$, où $M$ est un $S$-module gradué engendré par un nombre fini d'éléments homogènes $x_i$ ; si $x_i$ est de degré $m_i$, $M$ est donc quotient d'une somme directe finie de module $S(m_i)$ ; par suite ($n^\circ 3$, prop. 10), on peut se borner au cas où $M = S(n)$ donc $F(n) = (S(m+n))^{\sim}$ ($n^\circ 3$, cor. 1 de la prop. 14). On peut alors prendre $n_0 = -m$ si $m < 0$, $n_0 = 0$ si $m \ge 0$ ; en effet, on a alors $m+n \ge 0$ pour $n \ge n_0$, et l'idéal $S_{m+n} + S_{m+n+1} + \dots$ de $S$ est de type fini, autrement dit il existe un entier $k$ et un homomorphisme du $S$-module $S^k$ sur cet idéal. Si on considère cet homomorphisme comme un homomorphisme $S^k \to S(m+n)$, on voit que son conoyau vérifie la condition (TN), et on en conclut (prop. 17 et $n^\circ 3$, prop. 10) que l'homomorphisme correspondant $(S^k)^{\sim} \to (S(m+n))^{\sim}$ est surjectif.

Corollaire 4. - Soit $F$ un faisceau cohérent sur $X$. Il existe un entier $n_0$ tel que pour tout $n \ge n_0$, $F$ soit isomorphe à un quotient d'un faisceau $(\mathcal{O}_X(-n))^k$ ($k$ dépendant de $n$).

Cela résulte du cor. 3 ainsi que du $n^\circ 3$, cor. 1 de la prop. 14.

Corollaire 5. - Soit $F$ un faisceau cohérent sur $X$. Il existe un entier $n_0$ tel que pour $n \ge n_0$, $F(n)$ soit engendré par un nombre fini de ses sections au-dessus de $X$.

5. Comportements fonctoriels.

Soient $S, S'$ deux anneaux gradués à degrés positifs, $\phi : S' \to S$ un

\newpage

%% ===== Page 256 =====

- 67 -

\marginpar{1, 2, 3}
\marginpar{1, $\phi_f$}

homomorphisme d'anneaux gradués. Nous désignerons par $G(\phi)$ la partie ouverte de $X = \operatorname{Proj}(S)$ complémentaire de $V_+(\phi(S'_+))$; autrement dit, c'est la réunion des $D_+(f)$ où $f$ parcourt l'ensemble des éléments homogènes de $\phi(S'_+)$ (car l'intersection de $S_n$ et de l'idéal gradué de $S$ engendré par $\phi(S'_+)$ est le $S_0$-module engendré par $\phi(S'_n)$). La restriction à $G(\phi)$ de l'application continue ${}^a\phi$ de $\operatorname{Spec}(S)$ dans $\operatorname{Spec}(S')$ est donc une application continue de $G(\phi)$ dans $\operatorname{Proj}(S')$, que nous noterons encore ${}^a\phi$ par abus de langage. Si $f' \in S'_+$ et $f = \phi(f')$, on a
$$(28) \quad {}^a\phi^{-1}(D_+(f')) = D_+(f)$$
d'après la remarque précédente sur $G(\phi)$, la formule (7) du n°2 et I, 1, 2, formule (5). D'autre part, l'homomorphisme $\phi$ définit canoniquement un homomorphisme $S'_{f'} \to S_f$ d'anneaux gradués, et par suite un homomorphisme $S'(f') \to S(f)$ par restriction aux éléments de degré 0; nous désignerons ce dernier homomorphisme par $\phi_f$; il lui correspond un morphisme $({}^a\phi_f, \phi_f) : \operatorname{Spec}(S(f)) \to \operatorname{Spec}(S'(f'))$ de schémas affines. Si on identifie canoniquement $\operatorname{Spec}(S(f))$ au schéma induit par $\operatorname{Proj}(S)$ sur $D_+(f)$ (n°2), on a défini un morphisme $D_+(f) \to D_+(f')$, et ${}^a\phi$ s'identifie à la restriction de ${}^a\phi$ à $D_+(f)$. D'autre part, il est immédiat que si $g'$ est un second élément homogène de $S'_+$, et $g = \phi(g')$, le diagramme
$$\begin{tikzcd}
D_+(f) \arrow[r] & D_+(f') \\
D_+(fg) \arrow[u] \arrow[r] & D_+(f'g') \arrow[u]
\end{tikzcd}$$
est commutatif, en raison de la commutativité du diagramme
$$\begin{tikzcd}
S(f') \arrow[d, "\phi_f"] \arrow[r, "\omega_{f', g'}"] & S(f'g') \arrow[d, "\phi_{fg}"] \\
S(f) \arrow[r, "\omega_{f, g}"] & S(fg)
\end{tikzcd}$$
Par suite, compte tenu de la définition de $G(\phi)$ et de la formule (7) du n°2 :

Proposition 18. - Étant donné un homomorphisme $\phi : S' \to S$ d'anneaux

\newpage

%% ===== Page 257 =====

- II-68 -

gradués, il existe un morphisme et un seul $(\phi, \varphi)$ du préschéma induit $\mathcal{G}(\varphi)$ dans $\operatorname{Proj}(S')$, tel que pour tout élément homogène $f' \in S'_+$, la restriction de ce morphisme à $D_+(f)$ (où $f = \varphi(f')$) coïncide avec le morphisme associé à l'homomorphisme $\varphi_{f'} : S'_{(f')} \to S_{(f)}$.

Soient $S''$ un troisième anneau gradué à degrés positifs, $\varphi' : S'' \to S'$ un homomorphisme d'anneaux gradués, et posons $\varphi'' = \varphi \circ \varphi'$. En raison de (20) et de la formule $\varphi''^* = (\varphi'^*).(\varphi^*)$ on vérifie aussitôt que l'on a $\mathcal{G}(\varphi'') \subset \mathcal{G}(\varphi)$ et que, si $\phi, \phi'$ et $\phi''$ sont les morphismes associés à $\varphi, \varphi'$ et $\varphi''$, on a $\phi'' = \phi \circ \phi''|_{\mathcal{G}(\varphi'')}$.

Supposons que $S$ (resp. $S'$) soit une $A$-algèbre graduée (resp. une $A'$-algèbre graduée), et soit $\psi$ un homomorphisme $A' \to A$ tel que le diagramme
$$\begin{array}{ccc} A' & \xrightarrow{\psi} & A \\ & \searrow & \downarrow \\ & & S' \xrightarrow{\varphi} S \end{array}$$
(23)
soit commutatif. On peut alors considérer $\mathcal{G}(\varphi)$ et $\operatorname{Proj}(S')$ comme des préschémas au-dessus de $\operatorname{Spec}(A)$ et $\operatorname{Spec}(A')$ respectivement ; si $\Phi$ et $\Psi$ sont les morphismes correspondant respectivement à $\varphi$ et $\psi$, le diagramme
$$\begin{array}{ccc} \mathcal{G}(\varphi) & \xrightarrow{\Phi} & \operatorname{Proj}(S') \\ \downarrow & & \downarrow \\ \operatorname{Spec}(A) & \xrightarrow{\Psi} & \operatorname{Spec}(A') \end{array}$$
est alors commutatif : il suffit en effet de le démontrer pour la restriction de $\Phi$ à $D_+(f)$, où $f = \varphi(f')$, $f'$ étant homogène dans $S'_+$; et alors cela résulte de la commutativité du diagramme
$$\begin{array}{ccc} A' & \xrightarrow{\psi} & A \\ \downarrow & & \downarrow \\ S'_{(f')} & \xrightarrow{\varphi_{f'}} & S_{(f)} \end{array}$$
Soient maintenant $M$ un $S$-module gradué, et considérons le $S'$-module $M_{[\varphi]}$ qui est évidemment gradué. Soit $f'$ homogène dans $S'_+$, et soit $f = \varphi(f')$; on sait (chap. 0, § 1, n°5) qu'il y a un isomorphisme canonique $(M_{[\varphi]})_{f'} \simeq (M_f)_{f'}$ et il est immédiat

\newpage

%% ===== Page 258 =====

que cet isomorphisme conserve les degrés, donc il donne un isomorphisme canonique $(M_{[\varphi]})(f') \simeq (M(f))[\varphi_f]$. Il correspond à cet isomorphisme un isomorphisme du faisceau $(M_{[\varphi]})\sim|D_+(f')$ sur le faisceau $(\Phi_+)_* (\tilde{M}|D_+(f))$ (n°3 et I, 1, 6, prop. 7). En outre, si $g'$ est un second élément homogène de $S'_+$ et $g = \varphi(g')$, le diagramme
$$\begin{tikzcd}
(M_{[\varphi]})(f') \simeq (M(f))[\varphi_f] \\
\downarrow \\
(M_{[\varphi]})(f'g') \simeq (M(fg))[\varphi_{fg}]
\end{tikzcd}$$
est commutatif, d'où on conclut aussitôt que l'isomorphisme $(M_{[\varphi]})\sim|D_+(f'g') \simeq (\Phi_+)_{fg*} (\tilde{M}|D_+(fg))$ est la restriction à $D_+(f'g')$ de $(M_{[\varphi]})\sim|D_+(f') \simeq (\Phi_+)_{f*} (\tilde{M}|D_+(f))$. Comme $\Phi_+$ est la restriction à $D_+(f)$ du morphisme $\Phi$, on voit donc, compte tenu de (22) :

\textsc{Proposition 19}. — Il existe un isomorphisme canonique fonctoriel du faisceau $(M_{[\varphi]})\sim$ sur le faisceau image directe $\Phi_*(\tilde{M}(\varphi))$.

On en déduit aussitôt un isomorphisme canonique fonctoriel de l'ensemble des $\varphi$-homomorphismes $M' \to M$ d'un $S'$-module gradué $M'$ dans un $S$-module gradué $M$ (correspondant à l'homomorphisme $\varphi : S' \to S$), sur l'ensemble des $\mathcal{O}_X$-homomorphismes $\tilde{M}' \to \Phi_*(\tilde{M}(\varphi))$ (où on a posé $X' = \operatorname{Proj}(S'))$.

Soient $A, A'$ deux anneaux, $\psi : A' \to A$ un homomorphisme d'anneaux, définissant un morphisme $\Psi : X = \operatorname{Spec}(A) \to Y' = \operatorname{Spec}(A')$. Soit $S'$ une $A'$-algèbre graduée à degrés positifs, et posons $S = S' \otimes_{A'} A$, qui est graduée de façon évidente ; $\varphi : S' \to S$ est alors un homomorphisme d'anneaux gradués rendant commutatif le diagramme (23). Comme $\varphi$ engendre ici $\varphi(S'_+)S_+$, on a $G(\varphi) = \operatorname{Proj}(S) = X$ ; d'où, en posant $X' = \operatorname{Proj}(S')$, un diagramme commutatif
$$(24) \quad \begin{tikzcd}
X \arrow[r, "\Phi"] \arrow[d, "\Psi"] & X' \arrow[d, "\Psi'"] \\
Y \arrow[r, "\psi"] & Y'
\end{tikzcd}$$
Soit maintenant $M'$ un $S'$-module gradué, et posons...

\newpage

%% ===== Page 259 =====

- II-70 -

qui est un $S$-module gradué de façon évident par les $M'_n \otimes_{A'} A$. Dans ces conditions :

Proposition 20. — Le diagramme (24) identifie le schéma $X$ au produit $X' \times_{Y'} Y$ ; le faisceau $\widetilde{M}$ s'identifie alors à $\Phi^* (\widetilde{M'}) \otimes_{\mathcal{O}_{X'}} \mathcal{O}_X$ (§ 1, n°1).

La première assertion sera démontrée si on prouve que pour tout $f'$ homogène dans $S'$, en posant $f = \varphi(f')$, les restrictions de $\Phi$ et $p$ à $D_+(f)$ identifient ce schéma au produit $D_+(f') \times_{Y'} Y$ (I, 2, 4, lemme 4) ; autrement dit, il suffit de prouver que $S_{(f)}$ s'identifie canoniquement à $S'_{(f')} \otimes_{A'} A$, ce qui est immédiat en vertu de l'existence de l'isomorphisme canonique $S_f \simeq S'_{f'} \otimes_{A'} A$ (chap. 0, § 1, n°5) conservant les degrés. De même, $M_f$ s'identifie canoniquement à $M'_f \otimes_{A'} A$ par un isomorphisme conservant les degrés (loc. cit.), donc $M_{(f)}$ s'identifie canoniquement à $M'_{(f')} \otimes_{A'} A$, ce qui établit la seconde assertion.

Proposition 21. — Soit $S$ un anneau gradué et soit $X = \operatorname{Proj}(S)$.

(i) Si $\varphi : S \to S'$ est un homomorphisme surjectif d'anneaux gradués, le morphisme correspondant $\Phi$ est défini dans $\operatorname{Proj}(S')$ tout entier et est une immersion fermée dans $X$.

(ii) Supposons en outre $S$ noethérien et engendré par l'ensemble $S_1$ des éléments homogènes de degré $1$. Soit $X'$ un sous-préschéma fermé de $X$, défini par un faisceau cohérent d'idéaux $\mathcal{J}$. Soit $I$ l'idéal gradué de $S$, image réciproque de $\Gamma_*(X', \mathcal{J})$ (n°3), et posons $S' = S/I$. Alors $X'$ est le sous-préschéma correspondant à l'immersion fermée $\operatorname{Proj}(S') \to X$ définie par l'homomorphisme canonique $S \to S'$.

(i) Comme $\varphi(S)_+$ engendre $S'_+$, $G(\varphi) = \operatorname{Proj}(S')$ ; montrons que si $I$ est le noyau de $\varphi$, $\Phi$ est une immersion fermée dont le sous-préschéma associé est défini par le faisceau $\widetilde{I}$ (qui est cohérent puisque $S$ est noethérien (n°4, prop. 17)). La question étant locale sur

\newpage

%% ===== Page 260 =====

- II-71 -

$X$, soit $f \in S$, et posons $f' = \varphi(f)$ ; comme $\varphi$ est un homomorphisme d'anneaux gradués, on vérifie aussitôt que $\varphi_f : S_{(f)} \to S_{(f')}$ est surjectif et que son noyau est $I_{(f)}$, d'où notre assertion (I, 3, 1).

(ii) D'après ce qui précède, on est ramené à prouver que l'homomorphisme de $\mathcal{O}_X$-modules $\tilde{I} \to \mathcal{O}_X$ déduit de l'injection $I \to S$ est un isomorphisme de $\tilde{I}$ sur $J$. Or, en vertu du th. 1 du n° 4 et de son cor. 2, il existe un sous-$S$-module gradué de type fini $I'$ de $\Gamma_*(J)$ tel que l'isomorphisme canonique $\beta : (\Gamma_*(J))^\sim \to J$ donne par restriction à $\tilde{I}'$ un isomorphisme $\tilde{I}' \to J$ ; d'autre part, comme $I$ est de type fini, puisque $S$ est noethérien, on peut supposer que $I' \supset \alpha(I)$. De même il existe un sous-$S$-module gradué $T'$ de $\Gamma_*(\mathcal{O}_X)$ tel que l'isomorphisme canonique $\beta : (\Gamma_*(\mathcal{O}_X))^\sim \to \mathcal{O}_X$, restreint à $\tilde{T}'$, donne un isomorphisme $\tilde{T}' \to \mathcal{O}_X$ ; on peut en outre supposer que $T'$ contient l'image $\alpha(S)$. Alors (n° 3, Remarque) le morphisme $\tilde{S} \to \tilde{T}'$ est un isomorphisme. En vertu de la prop. 10 du n° 3 et de la prop. 17 du n° 4, comme $T'$ vérifie (TF), on en déduit que $T'_n / \alpha(S)_n$ vérifie (TN), autrement dit $(S)_n = T'_n$ à partir d'un certain rang ; en outre, si $N$ est le noyau de $S \to T'$, $N$ est un idéal de $S$, donc de type fini, et vérifie donc aussi (TN) d'après la prop. 17 du n° 4, donc $S_n \to T'_n$ est bijectif à partir d'un certain rang. Cela étant, dans le diagramme commutatif

\begin{tikzcd}
I_n \arrow[r] \arrow[d] & S_n \arrow[d] \\
I'_n \arrow[r] & T'_n
\end{tikzcd}

$u$ et $u'$ sont des injections et par définition $I_n$ est l'image réciproque par $\alpha$ de $u'(I'_n)$ ; si $S_n \to T'_n$ est bijectif, on en conclut que $I_n \to I'_n$ est bijectif, autrement dit le noyau et le conoyau de $I \to I'$ vérifient (TN). Donc (prop. 10 du n° 3 et prop. 17 du n° 4), $I \to I'$ est un isomorphisme, ce qui achève la démonstration.

Corollaire 1. - Les hypothèses et notations étant celles de la prop.

\newpage

%% ===== Page 261 =====

- II-72 -

21 (ii) : pour que le sous-préschéma fermé $X'$ soit réduit et que son espace de base soit irréductible, il suffit que l'idéal $I$ soit premier.

Cela résulte en effet des prop. 21 et 7.

Corollaire 2. - Soient $A$ un anneau, $M$ un $A$-module, $S$ une $A$-algèbre graduée engendrée par $S_1$. Supposons qu'il existe un $A$-homomorphisme surjectif $M \to S_1$. Si $S_A(M)$ désigne l'algèbre symétrique du $A$-module $M$, graduée de la façon usuelle, $\operatorname{Proj}(S)$ est isomorphe à un sous-préschéma fermé de $\operatorname{Proj}(S_A(M))$.

En effet, l'homomorphisme de $A$-modules $M \to S_1$ se prolonge canoniquement en un $A$-homomorphisme d'algèbres graduées $S_A(M) \to S$, et par hypothèse cet homomorphisme est surjectif.

\newpage

%% ===== Page 262 =====

\marginpar{Archive\\ Grothendieck - sept 59}

- II-73 -

§ 4. Morphismes projectifs et quasi-projectifs.

1. Préschéma associé à un faisceau gradué d'algèbres.

Soient $Y$ un préschéma, $S$ une $\mathcal{O}_Y$-algèbre graduée quasi-cohérente, à degrés positifs. Soient $U$ un ouvert affine de $Y$, $A$ son anneau ; par hypothèse, $S|_U$ est isomorphe à l'algèbre graduée $S_A$ où $S = \Gamma(U, S)$ est une $A$-algèbre graduée ; posons $X_U = \operatorname{Proj}(T(U, S))$. Soit $U' \subset U$ un second ouvert affine de $Y$, $A'$ son anneau, $j$ l'injection canonique $U' \to U$, qui correspond canoniquement à l'homomorphisme de restriction $A \to A'$; on a $S|_{U'} = j^*(S|_U)$, et par suite $S = \Gamma(U', S)$ s'identifie canoniquement à $S \otimes_A A'$ (1, 1, 6, prop. 8). On en conclut (cf. § 3, n° 5, prop. 20) que $X_{U'}$ s'identifie canoniquement à $X_U \times_U U'$, et par suite aussi à $f_U^{-1}(U')$, en désignant par $f_U$ le morphisme structural $X_U \to U$ (I, 3, 2, cor. 3 de la prop. 5). Nous désignerons par $\sigma_{U', U}$ l'isomorphisme canonique $f_U^{-1}(U') \xrightarrow{\sim} X_{U'}$, ainsi défini, par $\sigma_{U', U}$ l'immersion ouverte $X_{U'} \to X_U$ obtenue en composant $\sigma_{U', U}^{-1}$ et l'injection canonique $f_U^{-1}(U') \to X_U$.
Il est clair que si $U'' \subset U'$ est un troisième ouvert affine de $Y$, on a $f_{U''}^{-1}(U') \circ \sigma_{U', U} \circ \sigma_{U'', U'} = \sigma_{U'', U}$.

Proposition 1. -- Soit $Y$ un préschéma. Pour toute $\mathcal{O}_Y$-algèbre graduée quasi-cohérente $S$ à degrés positifs, il existe un préschéma $X$ au-dessus de $Y$, et un seul à un isomorphisme près, ayant la propriété suivante : si $f$ est le morphisme structural $X \to Y$, pour tout ouvert affine $U$ de $Y$, il existe un isomorphisme du préschéma induit $f^{-1}(U)$ sur $X_U = \operatorname{Proj}(\Gamma(U, S))$ tel que, si $V$ est un second ouvert affine de $Y$ contenu dans $U$, le diagramme
$$
\begin{tikzcd}
f^{-1}(V) \ar[d] \ar[r] & X_V \ar[d] \\
f^{-1}(U) \ar[r] & X_U
\end{tikzcd}
$$
(1) soit commutatif.

Avec les notations du début de ce n°, soit $X_{U, V}$ le préschéma induit...

\newpage

%% ===== Page 263 =====

- II-74 -

par $X_U$ sur $f^{-1}(U \cap V)$, pour deux ouverts affines quelconques $U, V$ de $Y$ ; nous allons définir un $Y$-isomorphisme $\theta_{U,V} : X_U \xrightarrow{\sim} X_V$. Pour cela, considérons un ouvert affine $W \subset U \cap V$ : en composant les isomorphismes $f^{-1}(W) \xrightarrow{\sigma_{W,U}} X_U$, $X_W \xrightarrow{\sigma_{W,V}^{-1}} f^{-1}(W)$, on obtient un isomorphisme $\theta_W$, et on vérifie aussitôt que si $W' \subset W$ est un ouvert affine, $\theta_{W'}$ est la restriction de $\theta_W$ à $f^{-1}(W')$ ; les $\theta_W$ sont donc bien les restrictions d'un $Y$-isomorphisme $\theta_{U,V}$. En outre, si $U, V, W$ sont trois ouverts affines, $\theta_{j,V}, \theta_{V,W}, \theta_{W,U}$ les restrictions de $\theta_{U,V}, \theta_{V,W}, \theta_{W,U}$ aux images réciproques de $U \cap V \cap W$ dans $X_U, X_W, X_U$ respectivement, il résulte des définitions précédentes que l'on a $\theta_{V,W} \circ \theta_{U,V} = \theta_{U,W}$. L'existence de $X$ vérifiant les propriétés énoncées se déduit donc du I, 3, 8 ; son unicité à un $Y$-isomorphisme près est triviale, compte tenu de (1).

Nous dirons que le préschéma $X$ ainsi défini est associé à la $\mathcal{O}_Y$-algèbre graduée $S$, et nous le noterons $\operatorname{Proj}(S)$. Il est immédiat que $\operatorname{Proj}(S)$ est séparé au-dessus de $Y$ (§ 3, n°2, prop. 6 et I, 3, 6, prop. 19), et de type fini en $Y$ si $S$ est une $\mathcal{O}_Y$-algèbre de type fini.

Nous désignerons par $S_d$ le $\mathcal{O}_Y$-module quasi-cohérent formé des éléments homogènes de degré $d > 0$ de $S$, par $S^{(d)}$ la $\mathcal{O}_Y$-algèbre graduée quasi-cohérente somme directe des $S_{nd}$ ($d > 0, n \in \mathbb{N}$).

Proposition 2. — Soit $d > 0$ et soit $f \in \Gamma(Y, S_d)$. Il existe une partie ouverte $X_f$ de l'espace de base de $X = \operatorname{Proj}(S)$ ayant la propriété suivante : pour tout ouvert affine $U$ de $Y$, $X_f \cap f^{-1}(U) = D_+(f|U)$ dans $f^{-1}(U)$ identifié à $\operatorname{Proj}(\Gamma(U, S))$ (f désignant le morphisme structural $X \to Y$). En outre, le préschéma induit sur $X_f$ est canoniquement isomorphe au préschéma associé au faisceau d'algèbres $S^{(d)} / (f-1)S^{(d)}$.

On a $f|U \in \Gamma(U, S_d) = (\Gamma(U, S))_d$. Si $U, U'$ sont deux ouverts affines de $Y$ tels que $U' \subset U$, $f|U'$ est l'image canonique de $f|U$ par l'homomorphisme de restriction $\Gamma(U, S) \to \Gamma(U', S)$ ; donc $D_+(f|U')$ s'identi-

\newpage

%% ===== Page 264 =====

- II - 75 -

est égal au préschéma induit sur l'image réciproque $\phi^{-1}(U, (D_+(f/U))$ dans $X_U$, (§ 3, n°5) ; d'où la première assertion.

En outre, le préschéma induit sur $D_+(f/U)$ par $X_U$ s'identifie canoniquement à $\operatorname{Spec}((\Gamma(U, S))_{(f|U)})$, ces identifications étant compatibles avec les homomorphismes de restriction ; la seconde assertion résulte donc du § 3, n°1, prop. 1 et de la Remarque suivant la prop. 5 du § 3 (par abus de langage).

On dira encore $X_f$ (en tant qu'ouvert l'espace de base $X$) est l'ensemble des $x \in X$ où $f$ ne s'annule pas.

Corollaire 1. - Si $f \in \Gamma(Y, S_d)$, $g \in \Gamma(Y, S_e)$ ; on a
$$(2) \qquad X_{fg} = X_f \cap X_g$$

Il suffit en effet de considérer l'intersection des deux membres avec un ensemble $\phi^{-1}(U)$, où $U$ est ouvert affine dans $Y$, et d'appliquer la formule (7) du § 3, n°2.

Corollaire 2. - Soit $(f_\alpha)$ une famille de sections de $S$ au-dessus de $Y$ telle que $f_\alpha \in \Gamma(Y, S_{d_\alpha})$ ; si le faisceau d'idéaux engendré par cette famille contient tous les $S_n$ à partir d'un certain rang, l'espace de base $X$ est réunion des $X_{f_\alpha}$.

En effet, pour tout ouvert affine $U$ de $Y$, $\phi^{-1}(U)$ est réunion des $X_{f_\alpha} \cap \phi^{-1}(U)$ (§ 3, n°2, cor. 3 de la prop. 5).

Corollaire 3. - Soit $J$ un faisceau quasi-cohérent d'idéaux sur $Y$, et soit $A = \mathcal{O}_Y / J$. Posons $S = A[T]$, où $T$ est une indéterminée. Alors $\operatorname{Proj}(S)$ s'identifie canoniquement au sous-préschéma fermé de $Y$ défini par $J$.

En appliquant le cor. 2 à l'unique section égale à $T$ en chaque point de $X = \operatorname{Proj}(S)$, on voit que $X_T = X$. On a ici $d=1$, et $S^{(1)} / (f-1)S^{(1)} = S / (f-1)S$ est canoniquement isomorphe à $A$, d'où le corollaire.

\newpage

%% ===== Page 265 =====

-- II-76 --

pour la première assertion. D'autre part, supposons $Y$ irréductible et soit $(U_\alpha)$ une base de la topologie de $Y$ formée d'ouverts affines (nécessairement irréductibles); si on a démontré que chacun des $\varphi^{-1}(U_\alpha)$ est irréductible (donc non vide) ($\varphi$ étant le morphisme structural $X \to Y$), il s'en suivra que $\varphi^{-1}(U_\alpha) \cap \varphi^{-1}(U_\beta)$, qui contient un $\varphi^{-1}(U_\gamma)$, est non vide, et comme cela est vrai pour tout $\beta$, on en conclut que $\varphi^{-1}(U_\alpha)$ est dense dans $X$, et par suite que $X$ est irréductible et $\varphi(X)$ dense dans $Y$.

Soit donc $Y = \operatorname{Spec}(A)$ et $S = \tilde{S}$, où $S$ est une $A$-algèbre graduée. Soit $x \in X$, $y = \varphi(x)$ ; l'anneau local $\mathcal{O}_{X,x} = \mathcal{O}_{X,y}$ est isomorphe à une fibre du faisceau structural $\mathcal{O}_Z$ du préschéma $Z = X \times_Y \operatorname{Spec}(\mathcal{O}_y)$ ($\S 2.8, \operatorname{prop}. 19$) ; mais $Z$ s'identifie à $\operatorname{Proj}(S \otimes_A A_y) = \operatorname{Proj}(S_y)$ ($\S 3, n^\circ 5, \operatorname{prop}. 20$) et par hypothèse $S_y = S(y)$ est un anneau intègre. Donc le préschéma $Z$ est réduit ($\S 3, n^\circ 2, \operatorname{prop}. 7$) et $\mathcal{O}_{X,x}$ est par suite sans élément nilpotent, ce qui prouve la première assertion.

Si $S = \bigoplus_n A_n$ est un sous-anneau de $S_y$ pour tout $y \in Y$, donc le schéma $Y$ est réduit et chacun des $A_y$ est intègre. Si de plus $Y$ est irréductible, l'anneau $A$ est intègre ($I, 1.1, \operatorname{prop}. 4$). Montrons que l'anneau $S$ est intègre, d'où résultera la proposition ($\S 3, n^\circ 2, \operatorname{prop}. 7$). En premier lieu, si $f \neq 0$ est un élément de $A$, $f$ n'est pas diviseur de $0$ dans $S$ ; en effet, $fg = 0$ avec $g \in S$ entraîne $(f/1)(g/1) = 0$ dans tout $S_y$ ($y \in Y$) ; comme $A$ est intègre, $f/1 \neq 0$ et comme $S_y$ est intègre on tire $g/1 = 0$. Cette relation ayant lieu pour tout $y$, on a $g=0$ ($I, 1.1, \operatorname{th}. 1$). On en conclut que l'homomorphisme canonique $S \to S_y$ est injectif, et comme $S_y$ est intègre, il en est de même de $S$.

\underline{Corollaire} : Supposons $Y$ réduit et irréductible, et que le faisceau quasi-cohérent d'algèbres graduées $S$ soit tel que $X = \operatorname{Proj}(S)$ soit réduit et irréductible et que $\varphi(X)$ soit dense dans $Y$. Soit $S'$ le fais-

\newpage

%% ===== Page 266 =====

\noindent -- II-77 --

\noindent de $\mathcal{O}_Y$ des algèbres graduées, somme directe de $\mathcal{O}_Y$ et des $S_d$ ($d > 0$), et soit $\mathcal{N}'$ le faisceau d'idéaux (gradué) de $S'$ tel que pour tout $y \in Y$, $\mathcal{N}'(y)$ soit l'idéal de $S_y$ formé des éléments nilpotents. Alors $\mathcal{N}'$ est quasi-cohérent, le faisceau $S'' = S'/\mathcal{N}'$ est intègre et $\operatorname{Proj}(S'')$ s'identifie canoniquement à $X$.

\noindent On sait que l'idéal des éléments nilpotents d'un anneau gradué est un idéal gradué ; l'existence du faisceau $\mathcal{N}'$ et le fait qu'il est quasi-cohérent se démontrent comme dans I, 3, 4, prop. 8. On sait déjà (prop. 3) que $\operatorname{Proj}(S')$ et $\operatorname{Proj}(S)$ s'identifient canoniquement, et on peut (puisque $\varphi(X)$ est partout dense) donc supposer que $S'=S$. La question étant locale sur $Y$, on peut de plus supposer que $Y = \operatorname{Spec}(A)$ est affine et $S = \tilde{S}$, où $S$ est une algèbre graduée telle que $S_0 = A$ ; soit $N$ l'idéal gradué des éléments nilpotents de $S$. Soit $S = S/N$ et désignons par $x \mapsto \overline{x}$ l'homomorphisme canonique $S \to \overline{S}$. Pour tout $f \in S_d$ ($d > 0$), on vérifie aussi que le noyau de l'homomorphisme surjectif $S_{(f)} \to \overline{S}_{(\overline{f})}$, défini par $x/f^n \mapsto \overline{x}/\overline{f}^n$ ($x \in S_{nd}$), est formé des éléments nilpotents de $S_{(f)}$ ; mais par hypothèse le préschéma $D_+(f) = \operatorname{Spec}(S_{(f)})$ est réduit, donc l'homomorphisme précédent est un isomorphisme, et comme le diagramme

$$
\begin{tikzcd}
S_{(f)} \arrow[r] \arrow[d] & \overline{S}_{(\overline{f})} \arrow[d] \\
S_{(fg)} \arrow[r] & \overline{S}_{(\overline{fg})}
\end{tikzcd}
$$

\noindent est commutatif ($g \in S_d$ et non nilpotent) ; on en conclut que $X$ est canoniquement isomorphe à $\operatorname{Proj}(\overline{S})$. Comme par hypothèse $A = S_0$ est intègre, on a $\overline{S}_0 = A$, et la conclusion résulte donc de la prop. du § 3, n° 2.

\newpage

%% ===== Page 267 =====

\subsection*{Proposition 3.} — Soit $S$ une $\mathcal{O}_Y$-algèbre graduée quasi-cohérente à degrés positifs.
(i) Il existe un isomorphisme canonique de $\operatorname{Proj}(S)$ sur $\operatorname{Proj}(S^{(d)})$, pour tout $d > 0$.
(ii) Soit $S'$ la $\mathcal{O}_Y$-algèbre graduée somme directe de $\mathcal{O}_Y$ et des $S_n$ ($n > 0$), $\operatorname{Proj}(S)$ et $\operatorname{Proj}(S')$ sont canoniquement isomorphes.
(iii) Soit $L$ un $\mathcal{O}_Y$-module inversible (§ 1, n°3), et soit $S(L)$ la $\mathcal{O}_Y$-algèbre graduée somme directe des $S_n \otimes L^{\otimes n}$, alors $\operatorname{Proj}(S)$ et $\operatorname{Proj}(S(L))$ sont isomorphes.

Dans chacun des cas il suffit de définir l'isomorphisme localement sur $Y$, avec les opérations de restriction étant triviales. On peut donc supposer $Y$ affine, et alors (i) et (ii) résultent du § 3, n°2, Remarque ; pour (iii), on peut de plus supposer $L$ isomorphe à $\mathcal{O}_Y$, et la proposition est alors évidente.

2. Faisceau sur $\operatorname{Proj}(S)$ associé à un faisceau de modules gradués.

Soient $Y$ un préschéma, $S$ une $\mathcal{O}_Y$-algèbre graduée quasi-cohérente, à degrés positifs, $M$ un $S$-module gradué quasi-cohérent (sur $(Y, \mathcal{O}_Y)$, ou sur l'espace annelé $(Y, S)$, ce qui revient au même (§ 2, n°1, prop.1)). Avec les notations du n°1, désignons par $\widetilde{M}_U$ le faisceau de modules $(\Gamma(U, M))^{\sim}$ sur $X_U = \operatorname{Proj}(\Gamma(U, S))$ ; comme $\Gamma(U, M)^{\sim}$ s'identifie canoniquement à $\Gamma(U, M) \otimes_A A'$; on a $\widetilde{M}_U = \phi_U^*(\widetilde{M}_U)$ (§ 3, n°5, prop.20).

\subsection*{Proposition 5.} — Il existe sur $\operatorname{Proj}(S) = X$ un $\mathcal{O}_X$-module quasi-cohérent $\widetilde{M}$ un seul tel que, pour tout ouvert affine $U$ de $Y$, l'isomorphisme $\phi^{-1}(U) \cong \operatorname{Proj}(\Gamma(U, S))$ (morphisme structural $X \to Y$) transforme $\widetilde{M}|_{\phi^{-1}(U)}$ en $(\Gamma(U, M))^{\sim}$.

Comme $\phi_{U', U}$ s'identifie au morphisme d'injection $\phi^{-1}(U') \to \phi^{-1}(U)$, la proposition résulte aussitôt de la relation $\widetilde{M}_{U'} = \phi_{U', U}^*(\widetilde{M}_U)$ et du principe de recollement des faisceaux (FAC, I, 1, 4).

\newpage

%% ===== Page 268 =====

On dit que $\widetilde{M}$ est le $\mathcal{O}_X$-module associé au $S$-module gradué $M$.

Proposition 6. — Soit $M$ un $S$-module gradué quasi-cohérent, et soit $f \in \Gamma(Y, S_d)$ ($d > 0$). L'isomorphisme canonique $\phi$ de $X_f$ sur le préschéma affine au-dessus de $Y$ associé au faisceau d'algèbres $S^{(d)} / (f-1)S^{(d)}$ transforme $\widetilde{M}|_{X_f}$ en le faisceau (sur ce préschéma) associé au faisceau de modules $M^{(d)} / (f-1)M^{(d)}$ ($\S$ 2, n°4, prop. 9).

La question étant locale sur $Y$, on est aussitôt ramené à la prop. 1 du $\S$ 3, n°1, compte tenu de la remarque suivant la prop. 20 du $\S$ 3, n°5.

Proposition 7. — Le foncteur $M \mapsto \widetilde{M}$ est un foncteur additif exact et covariant de la catégorie des $S$-modules gradués quasi-cohérents dans la catégorie des $\mathcal{O}_X$-modules quasi-cohérents qui commute aux limites inductives et aux sommes directes.

La question étant locale sur $Y$, se ramène à la prop. 10 du $\S$ 3, n°3.

Proposition 8. — Soit $d > 0$. Sur l'ensemble ouvert $X_f$, le faisceau $(S(nd))^{\sim}$ est inversible ; en particulier, si la $\mathcal{O}_Y$-algèbre $S$ est engendrée par $S_1$, les faisceaux $(S(n))^{\sim}$ sont inversibles pour tout $n \in \mathbb{Z}$.

Il s'agit encore de questions locales sur $Y$, et on est ramené à la prop. 13 du $\S$ 3, n°3 et à ses corollaires.

On posera encore, pour tout $n \in \mathbb{Z}$
$$(3) \qquad \mathcal{O}_X(n) = (S(n))^{\sim}$$
et pour tout faisceau algébrique $F$ sur $X$
$$(4) \qquad F(n) = F \otimes_{\mathcal{O}_X} \mathcal{O}_X(n)$$

Proposition 9. — Supposons $S$ engendrée par $S_1$. Alors, si $M$ et $N$ sont deux $S$-modules gradués quasi-cohérents, il existe un isomorphisme canonique fonctoriel
$$(5) \qquad \lambda : \widetilde{M} \otimes_{\mathcal{O}_X} \widetilde{N} \xrightarrow{\sim} (\widetilde{M \otimes_S N})^{\sim}$$
Si de plus $M$ est, en tout point de $Y$, localement isomorphe au co-

\newpage

%% ===== Page 269 =====

— II-80 —

noyau d'un homomorphisme $S^q \to S^q / U$, il existe un isomorphisme canonique fonctoriel

(6) $\mu : (\underline{\operatorname{Hom}}_S(M,N))^{\sim} \xrightarrow{\sim} \underline{\operatorname{Hom}}_{\mathcal{O}_X}(\tilde{M}, \tilde{N})$ .

Les isomorphismes $\lambda$ et $\mu$ ont été définis au § 3, n°3 (prop. 14), lorsque $Y$ est affine ; ces définitions étant locales se transportent aussitôt au cas général considéré ici.

Corollaire 1. — Quels que soient $n,m$ dans $\mathbb{Z}$, on a

(7) $\mathcal{O}_X(m) \otimes_{\mathcal{O}_X} \mathcal{O}_X(n) = \mathcal{O}_X(m+n)$

à un isomorphisme canonique près. En particulier

(8) $\mathcal{O}_X(n) = (\mathcal{O}_X(1))^{\otimes n}$

Corollaire 2. — Pour tout $S$-module gradué $M$ et tout $n \in \mathbb{Z}$, on a

(9) $(\widetilde{M(n)}) \simeq \tilde{M}(n)$

à un isomorphisme canonique près.

Cela résulte des propriétés correspondantes pour $Y$ affine (§ 3, n°3, cor. 1 et 2 de la prop. 14).

Dans toute la fin de ce n°, nous supposons que $S$ est engendrée par ses éléments de degré 1.

Soit $p$ le morphisme structural $X = \operatorname{Proj}(S) \to Y$. Pour tout $\mathcal{O}_X$-module $F$, nous poserons

(10) $\Gamma_*(F) = \bigoplus_{n \in \mathbb{Z}} p_*(F(n))$

et en particulier

(11) $\Gamma_*(\mathcal{O}_X) = \bigoplus_{n \in \mathbb{Z}} p_*(\mathcal{O}_X(n))$

On sait (chap. 0, § 2, n°) qu'il existe un homomorphisme canonique $p_*(F) \otimes_{\mathcal{O}_Y} p_*(G) \to p_*(F \otimes_{\mathcal{O}_X} G)$ pour deux $\mathcal{O}_X$-modules $F, G$ ; on déduit de (7) que $\Gamma_*(\mathcal{O}_X)$ est muni d'une structure de $\mathcal{O}_Y$-algèbre graduée, et de même (4) définit sur $\Gamma_*(F)$ une structure de module gradué sur $\Gamma_*(\mathcal{O}_X)$.

Soit $M$ un $S$-module gradué quasi-cohérent. Pour tout ouvert affine $U$ de $Y$, on a défini au § 3, n°3, un homomorphisme de groupes abéliens $\alpha_{U,M} : \Gamma(U, M) \to \Gamma(p^{-1}(U), \tilde{M})$. Il est immédiat que ces hom-

\newpage

%% ===== Page 270 =====

- II-81 - \hfill (20)
(\S 3, n°5, Remarque 2 suivant prop. 20), et définissent donc un homomorphisme de faisceaux de groupes abéliens
$\alpha : \underline{M} \to p_*(\underline{\widetilde{M}})$. Appliquant ce résultat à chacun des $M_n = \widetilde{M(n)}$, et tenant compte de (9), on définit un homomorphisme canonique fonctoriel (de degré 0) de faisceaux de groupes abéliens gradués
$$(12) \quad \alpha : \underline{\widetilde{M}} \to \underline{\Gamma}_*(\underline{\widetilde{M}})$$

En prenant en particulier $M = S$, on vérifie que $\alpha : \underline{S} \to \underline{\Gamma}_*(\mathcal{O}_X)$ est un homomorphisme de $\mathcal{O}_Y$-algèbres graduées, et que (12) est un $\alpha$-homomorphisme de modules gradués, relatif à cet homomorphisme.

\emph{Proposition 10}. — Pour toute section $f \in \Gamma(Y, S_d)$, $X_f$ est identique à l'ensemble des points de $X$ où $\alpha_d(f)$ (considérée comme section de $\mathcal{O}_X(d)$) ne s'annule pas (\S 1, n°3).

(Par définition, $\alpha_d(f)$ est une section de $p_*(\mathcal{O}_X(d))$ au-dessus de $Y$, mais une telle section est aussi par définition une section de $\mathcal{O}_X(d)$ au-dessus de $X$. La définition de $X_f$ (n°1, prop.2) ramène au cas où $Y$ est affine, qui a été traité au \S 3, n°3, prop. 15).

Soit maintenant $\mathcal{F}$ un faisceau algébrique sur $X$, ayant la propriété que chacun des $p_*(\mathcal{F}(n))$ est quasi-cohérent sur $Y$. Alors $\Gamma_*(\mathcal{F}) = \bigoplus_{n \in \mathbb{Z}} p_*(\mathcal{F}(n))$ est quasi-cohérent (I, 1, 3, cor. 4 du th. 1) et par suite le faisceau $(\Gamma_*(\mathcal{F}))^{\sim}$ est défini et est un $\mathcal{O}_X$-module quasi-cohérent.
(I, 1, 3, cor. 4 du th. 1 et Pour tout ouvert affine $U$ de $Y$, on a (\S 3, n°2, prop. 10)
$$(\Gamma(U, \bigoplus_{n \in \mathbb{Z}} p_*(\mathcal{F}(n))))^{\sim} = (\bigoplus_{n \in \mathbb{Z}} \Gamma(U, p_*(\mathcal{F}(n))))^{\sim}$$
$$= (\bigoplus_{n \in \mathbb{Z}} \Gamma(p^{-1}(U), \mathcal{F}(n)))^{\sim} = (\bigoplus_{n \in \mathbb{Z}} \Gamma(p^{-1}(U), \mathcal{F}(n)))^{\sim} =$$
$$= (\Gamma_*( \mathcal{F}|_{p^{-1}(U)} ))^{\sim}$$
et par suite (\S 3, n°3) on a un homomorphisme canonique
$$\rho_U : (\Gamma(U, \bigoplus_{n \in \mathbb{Z}} p_*(\mathcal{F}(n))))^{\sim} \to \mathcal{F}|_{p^{-1}(U)} \text{ .}$$
En outre, ces homomorphismes sont compatibles avec les opérations de restriction (\S 3, n°5, Remarque 2 suivant Prop. 20), et on en déduit donc un homomorphisme canonique fonctoriel

\newpage

%% ===== Page 271 =====

(13) $\beta : (\Gamma_*(F))^{\sim} \to F$
pour les faisceaux algébriques $F$ sur $X$ satisfaisant à la condition précédente.

3. Conditions de finitude.

Proposition 11. — Soient $Y$ un préschéma localement noethérien, $S$ une $\mathcal{O}_Y$-algèbre graduée quasi-cohérente engendrée par $S_1$ ; on suppose en outre $S_1$ cohérent. Alors $X = \operatorname{Proj}(S)$ est de type fini au-dessus de $Y$ et est localement noethérien ; si en outre $Y$ est noethérien, $X$ est noethérien.

Tout revient à prouver la première assertion (I, 4.2, prop. 10). La question étant locale sur $Y$, on peut supposer $Y$ affine [noethérien en anneau $A$]. On a alors $S \simeq S'$, où $S' = \Gamma(Y, S_1)$ et par hypothèse $S$ est engendré par $S_1 = \Gamma(Y, S_1)$, et $S_1$ est un $A$-module de type fini (I, 1.5, th. 3). Alors $S$ est une $A$-algèbre de type fini, et la proposition résulte du § 3, n° 4, prop. 16.

Soient $S$ une $\mathcal{O}_Y$-algèbre graduée quasi-cohérente, $M$ un $S$-module gradué quasi-cohérent. Nous considérerons les conditions de finitude suivantes :

(TF') Il existe un recouvrement $(U_\alpha)$ de $Y$ formé d'ouverts affines, et pour chaque $\alpha$ un entier $n_\alpha$ tel que le module $\bigoplus_{k \ge n_\alpha} (M|_{U_\alpha})$ soit de type fini.

(TN') Il existe un recouvrement $(U_\alpha)$ de $Y$ formé d'ouverts affines, et pour chaque $\alpha$ un entier $n_\alpha$ tel que $M|_{U_\alpha} = 0$ pour $k \ge n_\alpha$.

Proposition 12. — Soient $Y$ un préschéma localement noethérien, $S$ une $\mathcal{O}_Y$-algèbre graduée quasi-cohérente engendrée par $S_1$, $S_1$ étant supposé cohérent. Soit $M$ un $S$-module gradué quasi-cohérent.

(i) Si $M$ vérifie la condition (TF'), $\widetilde{M}$ est cohérent.

(ii) Supposons que $M$ vérifie (TF') ; pour que $\widetilde{M} = 0$, il faut et il suffit que $M$ vérifie (TN').

Les questions étant locales sur $Y$, on est ramené au cas où $Y$ est

\newpage

%% ===== Page 272 =====

- II-83 -

Remarque. - Si $Y$ est noethérien, on peut supposer dans (TF') et (TF') que le recouvrement $(U_\alpha)$ est fini et que les $n_\alpha$ sont tous égaux à un même $n$; ces conditions signifient alors respectivement que $\bigoplus_{k \ge n} M_k$ est un $S$-module de type fini, et que $M_k = 0$ pour $k \ge n$.

Théorème 1. - Soient $Y$ un préschéma localement noethérien, $S$ une $\mathcal{O}_Y$-algèbre graduée quasi-cohérente engendrée par $S_1$, $S_1$ étant supposé cohérent. Pour tout faisceau quasi-cohérent $\mathcal{F}$ sur $X = \operatorname{Proj}(S)$, l'homomorphisme canonique $\beta$ (n°2) est défini et est un isomorphisme.

En effet, le morphisme structural $p : X \to Y$ est séparé et de type fini (prop. 11); chacun des faisceaux $\mathcal{F}(n)$ étant quasi-cohérent (§ 1, n°1, prop. 3), il en est de même de $p_*(\mathcal{F}(n))$ (§ 1, n°1, cor. de la prop. 1), donc $\beta$ est défini. Pour voir que c'est un isomorphisme, on se ramène au cas où $Y$ est affine d'anneau noethérien $A$, $S = \tilde{S}$, où $S$ est une $A$-algèbre graduée engendrée par $S_1$, $S_1$ étant de type fini. La définition de $\beta$ et le th. 1 du § 3, n°4 prouvent alors le théorème.

Corollaire 1. - Tout faisceau quasi-cohérent $\mathcal{F}$ sur $X$ est isomorphe à un faisceau de la forme $\tilde{M}$, où $M$ est un $S$-module gradué.

C'est immédiat en prenant $M = \Gamma_*(\mathcal{F})$.

Corollaire 2. - On suppose en outre $Y$ noethérien. Si $\mathcal{F}$ est cohérent, on peut supposer $M$ de type fini.

En effet, on peut supposer $\mathcal{F} = \tilde{M}$ (cor. 1); $M$ est limite inductive de ses sous-faisceaux cohérents (§ 1, n°2, cor. du th. 1); on peut en outre supposer que ces sous-faisceaux $M_\lambda$ sont gradués, comme il résulte de la démonstration du corollaire cité. La prop. 7 du n°2 montre que $\tilde{M}$ est limite inductive des $\tilde{M}_\lambda$, donc $\mathcal{F}$ limite inductive des $\beta(\tilde{M}_\lambda)$. Mais comme $X$ est noethérien (prop. 11) et $\mathcal{F}$ cohérent, $\mathcal{F}$ est égal à un

\newpage

%% ===== Page 273 =====

des $\beta(\tilde{N}_\lambda)$ (chap. 0, $\S$ 2, $n^\circ$ 1).

On suppose $Y$ noethérien.

Corollaire 3. — Soit $\mathcal{F}$ un faisceau cohérent sur $X$. Il existe $n_0$ tel que pour $n \geq n_0$, l'homomorphisme canonique $p^*(p_*(\mathcal{F}(n))) \to \mathcal{F}(n)$ (chap. 0, $\S$ 2, $n^\circ$ 1) soit surjectif.

En effet, pour tout $x \in X$, soit $U$ un voisinage ouvert affine de $y=p(x)$ dans $Y$, dont l'anneau soit noethérien. Il existe un entier $n_0(U)$ tel que pour $n \geq n_0(U)$, $\mathcal{F}(n)|p^{-1}(U)$ soit engendré par un nombre fini de ses sections au-dessus de $p^{-1}(U)$ ($\S$ 3, $n^\circ$ 4, cor. 5 du th. 1) ; mais ces dernières sont par définition images canoniques de sections de $p^*(p_*(\mathcal{F}(n)))$ au-dessus de $p^{-1}(U)$, donc $\mathcal{F}(n)|p^{-1}(U)$ est égal à l'image canonique de $p^*(p_*(\mathcal{F}(n)))|p^{-1}(U)$. Enfin, comme $Y$ est noethérien, il existe un recouvrement fini de $Y$ par des ouverts affines $U_i$ d'anneaux noethériens, et en prenant pour $n_0$ le plus grand des $n_0(U_i)$, on achève la démonstration.

Nous verrons au chap. III, $\S$ , que sous les conditions du cor. 3, $p_*(\mathcal{F}(n))$ est un faisceau cohérent. Mais nous avons déjà remarqué ci-dessus que $p_*(\mathcal{F}(n))$ est quasi-cohérent, donc ($\S$ 1, $n^\circ$ 2, cor. du th. 1) limite inductive de ses sous-faisceaux cohérents. On en déduit (chap. 0, $\S$ 2, $n^\circ$ 1) que $\mathcal{F}(n)$ est l'image canonique d'un faisceau $p^*(\mathcal{G})$, où $\mathcal{G}$ est un sous-faisceau cohérent de $p_*(\mathcal{F}(n))$. En d'autres termes (compte tenu de (4) et (7)) :

Corollaire 4. — On suppose $Y$ noethérien. Si $\mathcal{F}$ est un faisceau cohérent, il existe un entier $n_0$ tel que pour tout $n \geq n_0$, $\mathcal{F}$ soit isomorphe à un quotient d'un faisceau de la forme $p^*(\mathcal{G})(-n)$, où $\mathcal{G}$ est un faisceau cohérent sur $Y$ (dépendant de $n$).

4. Comportements fonctoriels.

Soient $Y$ un préschéma, $\mathcal{S}, \mathcal{S}'$ deux $\mathcal{O}_Y$-algèbres graduées quasi-cohérentes à degrés positifs ; posons $X = \operatorname{Proj}(\mathcal{S})$, $X' = \operatorname{Proj}(\mathcal{S}')$, et soient $p, p'$ les morphismes structuraux de $X$ et $X'$. Soit $\varphi : \mathcal{S} \to \mathcal{S}'$

\newpage

%% ===== Page 274 =====

- II-85 -

un $\mathcal{O}_Y$-homomorphisme d'Algèbres graduées, de degré $0$. Pour tout ouvert affine $U$ de $Y$, posons $S_U = \Gamma(U, \mathcal{S})$, $S'_U = \Gamma(U, \mathcal{S}')$; l'homomorphisme $\varphi$ définit un homomorphisme $\varphi_U : S_U \to S'_U$ de $A_U$-algèbres graduées (où $A_U = \Gamma(U, \mathcal{O}_Y)$). Il lui correspond dans $p^{-1}(U)$ un ensemble ouvert $G(\varphi_U)$ et un morphisme $\Phi_U : G(\varphi_U) \to p^{-1}(U)$ (§ 3, n°5). En outre, en vertu de la commutativité du diagramme
$$\begin{tikzcd}
S_U \arrow[r, "\varphi_U"] \arrow[d] & S'_U \arrow[d] \\
S_V \arrow[r, "\varphi_V"] & S'_V
\end{tikzcd}$$
(où $V$ est un ouvert affine contenu dans $U$), on vérifie aussitôt que $G(\varphi_V) = G(\varphi_U) \cap p^{-1}(V)$ et que $\Phi_V$ est la restriction de $\Phi_U$ à $G(\varphi_U)$. On a ainsi défini une partie ouverte $G(\varphi)$ de $X'$ telle que $G(\varphi) \cap p^{-1}(U) = G(\varphi_U)$ pour tout ouvert affine $U \subset Y$, et un morphisme $\Phi : G(\varphi) \to X$, qui sera dit associé à $\varphi$ et que nous noterons $\operatorname{Proj}(\varphi)$. Lorsque localement $\varphi(S_+)$ engendre $S'_+$ (en tant que $\mathcal{O}_Y$-module) on a $G(\varphi)=X'$.

Soit maintenant $M'$ un $S'$-module gradué quasi-cohérent, qui donne un $S$-module gradué $M'_{[\varphi]}$.

\textbf{Proposition 13.} — Il existe un isomorphisme canonique fonctoriel du faisceau $(M'_{[\varphi]})^{\sim}$ sur le faisceau image directe $\Phi_*(M'^{\sim} | G(\varphi))$.

On définit en effet aussitôt cet isomorphisme lorsque $Y$ est affine (tenant compte du § 3, n°5, prop. 19), et dans le cas général il suffit de vérifier que ces isomorphismes ainsi définis dans les ouverts $p^{-1}(U)$ ($U$ ouvert affine de $Y$) sont compatibles avec les opérations de restriction, ce qui est immédiat.

\textbf{Proposition 14.} — Soient $Y, Y'$ deux préschémas, $\varphi : Y' \to Y$ un morphisme, $\mathcal{S}$ une $\mathcal{O}_Y$-Algèbre graduée quasi-cohérente, $\mathcal{S}' = \varphi^*(\mathcal{S})$ son image réciproque; le $Y'$-préschéma $\operatorname{Proj}(\mathcal{S}')$ s'identifie canoniquement à $\operatorname{Proj}(\mathcal{S}) \times_Y Y'$. En outre si $\mathcal{M}$ est un $\mathcal{S}$-module gradué quasi-cohérent, le faisceau $(\varphi^*(\mathcal{M}))^{\sim}$ sur $X'$ s'identifie alors à $\tilde{\mathcal{M}} \otimes_{\mathcal{O}_X} \mathcal{O}_{X'}$.

\newpage

%% ===== Page 275 =====

- II-86 -
Notons d'abord que $\psi^*(S)$ et $\psi^*(M)$ sont quasi-cohérents (§ 1, n°1, prop.2). Si $U$ est un ouvert affine de $Y$, $U' \subset \psi^{-1}(U)$ un ouvert affine de $Y'$, $A, A'$ les anneaux de $U$ et $U'$, on a $S|_U = \tilde{S}$, où $S$ est une $A$-algèbre graduée, et $S'|_{U'}$ s'identifie alors à $(S \otimes_A A')^{\sim}$ (I, 1, 6, prop. 8) ; la première assertion résulte alors de la prop. 20 du § 3, n°5, la vérification de la compatibilité de l'identification ainsi faite pour chaque couple $U, U'$ avec les opérations de restriction étant immédiate. Soient $q : \operatorname{Proj}(S) \to Y$, $q' : \operatorname{Proj}(S') \to Y'$ les morphismes structuraux ; $q'^{-1}(U')$ s'identifie alors à $q^{-1}(U) \times_U U'$, et les deux faisceaux $(\psi^*(M))^{\sim}|_{q'^{-1}(U')}$ et $(M \otimes_{\mathcal{O}_Y} \mathcal{O}_{Y'})^{\sim}|_{q'^{-1}(U')}$ s'identifient alors canoniquement tous deux à $(M \otimes_A A')^{\sim}$, où $M = \Gamma(U, M)$ ; d'où la seconde assertion, les compatibilités avec les opérations de restriction se vérifiant comme d'ordinaire.

Enfin, si on pose $M' = \psi^*(M)$, l'homomorphisme canonique $\alpha' : M' \to \Gamma_*(\tilde{M'})$ (n°2) n'est autre que $\psi^*(\alpha)$. De même, soit la projection $\operatorname{Proj}(S') \to \operatorname{Proj}(S)$ pour un faisceau algébrique $F$ sur $\operatorname{Proj}(S)$ posons $F' = \psi^*(F)$ ; si les faisceaux $\tilde{M}$ et $q_*(F(n))$ et $q'_*(F'(n))$ sont quasi-cohérents, l'homomorphisme canonique $\beta' : (\Gamma_* (F'))^{\sim} \to F'$ n'est autre que $\psi^*(\beta)$. La vérification de ces assertions se ramène encore au cas où $Y$ et $Y'$ sont affines, et alors on applique la Remarque 2 du § 3, n°5.

Proposition 15. - Soient $Y$ un préschéma, $S$ une $\mathcal{O}_Y$-algèbre graduée quasi-cohérente.

(i) Pour tout sous-faisceau gradué quasi-cohérent d'idéaux dans $S$, $X = \operatorname{Proj}(S/J)$ s'identifie canoniquement à un sous-préschéma fermé de $\operatorname{Proj}(S)$.

(ii) Supposons en outre $Y$ localement noethérien, $S$ engendrée par $S_1$ et $S_1$ cohérent. Soit $X$ un sous-préschéma fermé de $X = \operatorname{Proj}(S)$, dé-

\newpage

%% ===== Page 276 =====

- II-87 -

fini par un faisceau cohérent $\mathcal{I}$ d'idéaux. Soit $J$ le faisceau gradué d'idéaux dans $S$, image réciproque de $\mathcal{I}'(\mathcal{I})$ par l'homomorphisme $\alpha : S \to \Gamma_*(\mathcal{O}_X)$. Alors $X'$ s'identifie canoniquement à $\operatorname{Proj}(S/J)$.

(i) L'application canonique $S \to S'=S/J$ étant surjective, on a vu au début de ce $n^\circ$ qu'il lui correspond un $Y$-morphisme $\operatorname{Proj}(S') \to \operatorname{Proj}(S)$. On est immédiatement ramené à démontrer la proposition dans le cas où $Y$ est affine, et elle découle alors aussitôt de la prop. 21 (i) du $\S 3, n^\circ 5$. On notera en outre que cela prouve que le sous-préschéma fermé de $\operatorname{Proj}(S)$ auquel s'identifie $\operatorname{Proj}(S/J)$ est défini par le faisceau quasi-cohérent d'idéaux $\mathcal{J}$ sur $\operatorname{Proj}(S)$.

(ii) On est ramené à démontrer que l'homomorphisme $\mathcal{J} \to \mathcal{O}_X$ déduit de l'injection canonique $I \to S$, est un isomorphisme de $\mathcal{J}$ sur le faisceau $\mathcal{I}$, et comme la question est locale sur $Y$, on peut supposer $Y$ affine d'anneau noethérien, ce qui entraîne $S = \tilde{S}$, où $S$ est une $A$-algèbre graduée engendrée par $S_1$ de type fini, et par suite $S$ est noethérien. Il suffit alors d'appliquer la prop. 21 (ii) du $\S 3, n^\circ 5$.

Corollaire 2. — Soient $Y$ un préschéma, $S$ une $\mathcal{O}_Y$-algèbre graduée quasi-cohérente engendrée par $S_1$, $M$ un $\mathcal{O}_Y$-module quasi-cohérent. Supposons qu'il existe un $\mathcal{O}_Y$-homomorphisme surjectif $M \to S_1$. Si $S_{\mathcal{O}_Y}(M)$ est la $\mathcal{O}_Y$-algèbre symétrique de $M$ ($\S 2, n^\circ 8$), $\operatorname{Proj}(S)$ est isomorphe à un sous-préschéma fermé de $\operatorname{Proj}(S_{\mathcal{O}_Y}(M))$.

En effet, l'homomorphisme $M \to S_1$ se prolonge en un homomorphisme de $\mathcal{O}_Y$-algèbres graduées $S_{\mathcal{O}_Y}(M) \to S$, qui est par hypothèse surjectif.

5. \emph{Fibrés projectifs}. Morphismes dans un fibré projectif.

Définition 1. — Soient $Y$ un préschéma, $E$ un $\mathcal{O}_Y$-module quasi-cohérent. On appelle fibré projectif sur $Y$ défini par $E$, et on note $\mathbb{P}(E)$, le $Y$-préschéma $\operatorname{Proj}(S_{\mathcal{O}_Y}(E)) = P$. Le $\mathcal{O}_P$-module $\mathcal{O}_P(1)$ ($n^\circ 2$) est appelé le faisceau fondamental sur $P$.

\newpage

%% ===== Page 277 =====

\marginpar{Archives \\ Grothendieck sept 59}

II-80 --

Soient $\mathcal{E}, \mathcal{F}$ deux $\mathcal{O}_Y$-modules quasi-cohérents, $u : \mathcal{E} \to \mathcal{F}$ un $\mathcal{O}_Y$-homomorphisme ; il lui correspond canoniquement un homomorphisme $\bar{u} : \mathcal{S}_{\mathcal{O}_Y}(\mathcal{E}) \to \mathcal{S}_{\mathcal{O}_Y}(\mathcal{F})$ de $\mathcal{O}_Y$-Algèbres graduées. Si $u$ est surjectif, il en est de même de $\bar{u}$, et par suite ($n^\circ 4$, prop. 15) $\operatorname{Proj}(\bar{u})$ est une immersion fermée $\mathbb{P}(\mathcal{F}) \to \mathbb{P}(\mathcal{E})$, que nous noterons $\mathbb{P}(u)$. On pose pour abréger $\mathcal{S}=\mathcal{S}_{\mathcal{O}_Y}(\mathcal{E})$, $\mathcal{S}'=\mathcal{S}_{\mathcal{O}_Y}(\mathcal{F})$, et si on désigne par $\mathcal{J}$ le noyau de $\bar{u}$, qui est un faisceau gradué d'idéaux, il est clair que $\mathcal{J}(n)$ est le noyau de l'homomorphisme correspondant $\mathcal{S}(n) \to \mathcal{S}'(n)$, et par suite $\mathbb{P}=\mathbb{P}(\mathcal{E})$, $\mathbb{P}'=\mathbb{P}(\mathcal{F})$ est isomorphe à $\mathcal{O}_{\mathbb{P}}(n)/\mathcal{J}(n)$ ($n^\circ 2$, prop. 7 et cor. 2 de la prop. 9), et ce dernier n'est autre que le faisceau $j_*(\mathcal{O}_{\mathbb{P}'}(n))$, en posant $j=\mathbb{P}(u)$ ; autrement dit, on a
$$(14) \quad j_*(\mathcal{O}_{\mathbb{P}'}(n)) = \mathcal{J}_n(\mathcal{O}_{\mathbb{P}}(n))$$

Soit maintenant $\psi : Y' \to Y$ un morphisme, et posons $\mathcal{E}'=\psi^*(\mathcal{E})$ ; on a alors $\mathcal{S}_{\mathcal{O}_{Y'}}(\mathcal{E}') = \psi^*(\mathcal{S}_{\mathcal{O}_Y}(\mathcal{E}))$ ($\S 2, n^\circ 8$) ; par suite ($n^\circ 4$, prop. 14), on a
$$(15) \quad \mathbb{P}(\psi^*(\mathcal{E})) = \mathbb{P}(\mathcal{E}) \times_Y Y'$$
à un isomorphisme canonique près. En outre, $\psi^*(\mathcal{S}_{\mathcal{O}_Y}(\mathcal{E})(n)) = \mathcal{S}_{\mathcal{O}_{Y'}}(\mathcal{E}')(n)$ ; donc ($n^\circ 4$, prop. 14)
$$(16) \quad \mathcal{O}_{\mathbb{P}'}(n) = \mathcal{O}_{\mathbb{P}}(n) \otimes_Y Y'$$

Posons $\mathbb{P}=\mathbb{P}(\mathcal{E})$, et désignons par $p$ le morphisme structural $\mathbb{P} \to Y$.
Comme $\mathcal{S}=(\mathcal{S}_{\mathcal{O}_Y}(\mathcal{E}))_1$, par définition, on a un homomorphisme canonique $\alpha_1 : \mathcal{E} \to p_*(\mathcal{O}_{\mathbb{P}}(1))$ ($n^\circ 2$), et par suite un homomorphisme canonique (chap. 0, $\S 2, n^\circ$)
$$(17) \quad \alpha : p^*(\mathcal{E}) \to \mathcal{O}_{\mathbb{P}}(1)$$

Lorsque $Y=\operatorname{Spec}(A)$ est affine, on explicite aisément cet homomorphisme $\alpha$ : remontant aux définitions ($\S 3, n^\circ 3$) : on a alors $\mathcal{E}=\tilde{E}$, où $E$ est un $A$-module, et $\mathcal{S}_{\mathcal{O}_Y}(\mathcal{E})=(\mathcal{S}_A(E))^{\sim}$, d'où $\mathbb{P}=\operatorname{Proj}(\mathcal{S}_A(E))$. Posons pour simplifier $\mathcal{S}=\mathcal{S}_A(E)$, et notons que lorsque $f$ parcourt $E$, les $D_+(f)$ forment un recouvrement de $\mathbb{P}$ puisque $E$ engendre $\mathcal{S}$ ; on voit alors aussitôt que la restriction de (17) à $D_+(f)$ correspond canoniquement ($I, 1, 3$, cor. 2 du th. 1) à l'homomorphisme de $E \otimes_A S(f)$ dans

\newpage

%% ===== Page 278 =====

$(s(1))_{(f)}$ qui, à tout élément $x \otimes 1$, fait correspondre $x/1$, et par suite, à tout élément $x \otimes (x_1 x_2 \dots x_k / f^k)$ ($x, x_i \in E$) fait correspondre $xx_1 x_2 \dots x_k / f^k$. Ceci montre que lorsque $Y$ est affine, l'homomorphisme (17) est surjectif, et comme cette propriété est locale sur $Y$, elle subsiste dans le cas général.

Soit maintenant $X$ un $Y$-préschéma et $r : X \to P$ un $Y$-morphisme, de sorte que si $q$ est le morphisme structural $X \to Y$, on a le diagramme commutatif :

$$\begin{tikzcd}
P \arrow[d, "p"] & X \arrow[l, "r"] \arrow[dl, "q"] \\
Y
\end{tikzcd}$$

Comme le foncteur $r^*$ est exact à droite, on déduit de l'homomorphisme surjectif (17) un homomorphisme surjectif $r^*(\alpha_1) : r^*(p^*(E)) \to r^*(\mathcal{O}_P(1))$. Mais $r^*(p^*(E)) = q^*(E)$, et $r^*(\mathcal{O}_P(1))$ est localement isomorphe à $r^*(\mathcal{O}_P) = \mathcal{O}_X$, autrement dit c'est un faisceau inversible $L_r$ sur $\mathcal{O}_X$, et on a ainsi défini un $\mathcal{O}_X$-homomorphisme surjectif :

(18) $$\varphi_r : q^*(E) \to L_r$$

Lorsque $Y = \operatorname{Spec}(A)$ est affine, on explicite encore cet homomorphisme de la façon suivante (avec les notations ci-dessus) : étant donné $f \in E$, soit $V$ un ouvert affine de $X$ contenu dans $r^{-1}(D_+(f))$, et soit $B$ son anneau, qui est une $A$-algèbre ; la restriction de $r$ à $V$ correspond à un $A$-homomorphisme $\omega : S_{(f)} \to B$, on a $q^*(E) = (E \otimes_A B)^\sim$, et $L_r = L_f \otimes_{S_{(f)}} B$ ; $L_f$ est un $B$-module libre ayant pour base $(f/1) \otimes 1$ (cf. $\S 3, n^\circ 3, \text{prop.} 13$), et $\varphi_r$ correspond à l'homomorphisme $E \otimes_A B \to L_f \otimes_{S_{(f)}} B$ qui transforme $x \otimes 1$ en $(f/1) \otimes \omega(x/f)$. En outre, si $j \in V$ est un idéal de $B$, le point $p = r(j) \in P$ se détermine de la façon suivante : on l'identifie canoniquement ($\S 3, n^\circ 2, \text{prop.} 4$) à l'idéal $\omega^{-1}(j)$ de $S_{(f)}$ ; ce dernier est l'ensemble des $s/f^n$ ($s \in S_n, n > 0$ arbitraire) tels que $\omega(s/f^n) \in j$.

\marginpar{monogène}

\newpage

%% ===== Page 279 =====

- II-90 -

Or la puissance tensorielle $L^{\otimes n}$ s'identifie canoniquement à la puissance symétrique $S_n(L)$ et il correspond donc à $u$ un homomorphisme $u_n=S_n(u) : S_n(E \otimes_A B) = S_n(A_B) \to L^{\otimes n}$ qui transforme $s \otimes 1$ en $(f/1)^{\otimes n} \otimes \omega(s/f^n)$ ; comme $p$ est l'ensemble des $s \in S_n$ tels que $s/f^n \in \omega^{-1}(j)$, on voit finalement que $p$ est l'ensemble des éléments homogènes $s \in S_n$ (identifiés à des sections de $S_n(E)$) qui, par l'homomorphisme $S_n(\varphi_r)$, sont transformés en sections de $L^{\otimes n}$ s'annulant au point $j$ (i.e. dont le germe au point $j$ est dans l'idéal maximal de l'anneau local $B_j$ correspondant).

Cette description montre que la donnée de $u$ détermine complètement l'homomorphisme $\omega$ et la restriction de $r$ à $V$. Donnons inversement, nous maintenant de façon générale ($X, Y$ et $E$ étant fixés) un faisceau inversible $L$ sur $X$, et un homomorphisme surjectif $\varphi : q^*(E) \to L$ ; nous allons en déduire canoniquement un morphisme \marginpar{$r=(\varphi, \theta)$} $r : X \to P$. L'homomorphisme $\varphi$ de $\mathcal{O}_X$-modules définit un $\mathcal{O}_X$-homomorphisme

(19) $$\mathcal{S}(\varphi) : S_{\mathcal{O}_X}(q^*(E)) \to \bigoplus_{n \geq 0} L^{\otimes n}$$

Soit $t$ un point de $X$ ; pour définir $\varphi(t) \in P$, on peut supposer d'abord que $Y=\operatorname{Spec}(A)$ est affine, donc $E=\tilde{E}$ et $P=\operatorname{Proj}(S)$ avec $S=S_A(E)$ ; pour tout $n > 0$, soit $p_n$ l'ensemble des sections $s \in S_n$ ayant la propriété suivante : l'image par $S_n(\varphi)$ de la section $s$ s'annule au point $t$. Il est immédiat que $p_n$ est en raison de l'hypothèse de surjectivité et du fait que $X$ est le support de $L$ ; les $p_n$ vérifient donc les conditions de la prop. 2 et nous définirons $p(t)$ comme l'idéal premier gradué de $S$ tel que $p \cap S_n = p_n$ pour tout $n > 0$.

Pour passer au cas général, il suffit de vérifier que lorsqu'on passe d'un ouvert affine $U \subset Y$ à un ouvert affine $V \subset U$, la définition...

\newpage

%% ===== Page 280 =====

on de $\psi(t)$ ne change pas, en supposant $q(t) \in V \subset U$ ; mais cela est immédiat, compte tenu de la définition de $\operatorname{Proj}(S)$ (n°1).

Pour définir $r$ dans un voisinage de $t \in X$, considérons encore en premier lieu le cas où $X$ et $Y$ sont affines. Soit $f$ un élément de $E_1$ non dans $\mathfrak{p} = \psi(t)$ ; on peut d'autre part supposer que $L$ est isomorphe à $\mathcal{O}_X$, où $L = \tilde{L}$, où $L$ est un $B$-module libre de rang 1. Soit $c$ un générateur de $L$ ; l'homomorphisme $\psi$ provient d'un homomorphisme $E \otimes_A B \to L$ qu'on peut composer avec l'homomorphisme canonique $E \to E \otimes_A B$ pour obtenir un homomorphisme de $A$-modules $E \to B$, où $v_c$ est un homomorphisme de $A$-modules $E \to B$. On en déduit canoniquement un homomorphisme de $A$-algèbres $v_c : S_A(E) \to B$, et par suite un homomorphisme de $A$-algèbres $S_A(E) \to B_g$, qui à tout élément $x_1 x_2 \dots x_n / f^n$ ($z_i \in E$) fait correspondre $v_c(z_1) \dots v_c(z_n)/g^n$, où $g = v_c(f)$. Cet homomorphisme, et l'anneau $B_g$, sont indépendants du générateur $c$ de $L$ considéré, car si on remplace $c$ par un autre générateur $c' = ac$ où $a$ est un élément inversible dans $B$, on a $v_{c'}(x) = a^{-1} v_c(x)$. Il correspond à $\psi$ un morphisme $r_f : D(g) \to D_+(f)$ et on vérifie aussitôt (compte tenu de l'identification canonique de $\operatorname{Spec}(S_{(f)})$ et de $D_+(f)$, §3, n°2, prop. 4) que pour tout $t' \in D(g)$, on a $r_f(t') = \psi(t')$, ce qui prouve la continuité de $\psi$ ; il est clair par ailleurs que $r_f$ ne dépend pas du point $D(g)$ initialement considéré. Enfin, si $f'$ est un second élément de $E$ non dans $\mathfrak{p}$, et $g' = v_c(f')$, $r_f$ est la restriction à $D(gg')$ de $r_f$ et $r_{f'}$. Cela montre l'existence du morphisme $r$ lorsque $X$ et $Y$ sont affines, et on passe de là aisément, d'abord au cas où $Y$ est affine, $X$ quelconque, puis au cas général en tenant compte de la définition de $\operatorname{Proj}(S)$ (n°1). En outre il y a un isomorphisme canonique $L_r \to L$ tel que $\psi = \varepsilon \circ p_r$ ; pour le définir, on peut se placer dans les conditions locales envisagées plus haut, et pour tout $x \in E$, faire correspondre à l'élément

\newpage

%% ===== Page 281 =====

-- 92 --

$$(x/1) \otimes 1 \text{ de } L_x \text{ en l'élément } v_c(x) \in L \text{ ; vu l'hypothèse de surjectivité, il suffit de vérifier que si } x/1 = 0 \text{ dans } S_{(f)}(1) \text{, on a } v_c(x) = 0 \text{ dans } B \text{ ; mais la première relation signifie que } f^n x = 0 \text{ dans } S_{n+1} \text{ pour un certain } n \text{, et on en conclut que } v_c(f^n x) = g^n v_c(x) = 0 \text{ dans } B \text{, d'où notre assertion.}$$

Nous désignerons par $r_{L, \varphi}$ le morphisme $X \to P$ ainsi défini par $L$ et $\varphi$.
On vérifie aussitôt que l'on a $r_{L, \varphi} = r$, et nous avons donc démontré le

\textsc{Théorème 2.} -- Les applications $r \mapsto (L, \varphi)$ et $(L, \varphi) \mapsto r_{L, \varphi}$ mettent en correspondance biunivoque l'ensemble des $Y$-morphismes $r : X \to P$ et l'ensemble des classes de couples $(L, \varphi)$, où $L$ est un $\mathcal{O}_X$-module inversible et un homomorphisme surjectif $\varphi : (q^*(E)) \to L$, deux couples $(L, \varphi)$, $(L', \varphi')$ étant équivalents s'il existe un $\mathcal{O}_X$-isomorphisme $\tau : L \to L'$ tel que $\varphi' = \tau \circ \varphi$.

\textsc{Corollaire 1.} -- Soit $U$ un ouvert de $Y$ ; l'ensemble des $Y$-sections canoniques de $P(E)$ au-dessus de $U$ est en correspondance biunivoque avec l'ensemble des sous-faisceaux quasi-cohérents $F$ de $E_{|U}$ tels que $(E_{|U})/F$ soit un faisceau inversible sur $U$.

Il suffit de prendre $X = U$ et de remarquer (avec les notations du th. 2) que $(L, \varphi)$ et $(L', \varphi')$ sont équivalents si et seulement si $\varphi$ et $\varphi'$ ont même noyau.

\textsc{Corollaire 2.} -- Supposons que tout $\mathcal{O}_Y$-module inversible soit isomorphe à $\mathcal{O}_Y$. Soient $V$ le groupe $\operatorname{Hom}_{\mathcal{O}_Y}(E, \mathcal{O}_Y)$, considéré comme module sur $A = \Gamma(Y, \mathcal{O}_Y)$, et soit $V^*$ le sous-ensemble de $V$ formé des homomorphismes surjectifs. Alors l'ensemble des $Y$-sections de $P(E)$ au-dessus de $Y$ s'identifie à $V^* / A^*$, où $A^*$ est le groupe des unités de $A$.

En particulier :
1° Le cor. 2 s'applique lorsque $Y$ est un schéma local (§ 1, n°3, Remarque). Si $Y$ est un préschéma quelconque, et un point $y$ de $Y$, $Y' = \operatorname{Spec}(k(y))$, la fibre de $P(E)$ au-dessus de $y$ est, d'après $q^{-1}(y)$

\newpage

%% ===== Page 282 =====

- II-93 -

(15), identifiée à $\operatorname{Proj}(S^\bullet)$, où $S^\bullet$ est l'algèbre symétrique de $\mathcal{E}_y \otimes \kappa(y) = E_y / \mathfrak{m}_y E_y = E_y^y$. Plus généralement, si $K$ est une extension de $\kappa(y)$, $q^{-1}(y) \otimes_{\kappa(y)} K$ s'identifie à $\operatorname{Proj}(S^\bullet \otimes_{\kappa(y)} K)$. Le cor. 2 montre donc que les points géométriques de $\mathbb{P}(E)$ à valeurs dans l'extension $K$ de $\kappa(y)$, ou encore les points géométriques de $q^{-1}(y)$ à valeurs dans $K$, s'identifient aux éléments de l'espace projectif associé au dual $\operatorname{Hom}_K(E_y \otimes_{\kappa(y)} K, K)$ du $K$-espace vectoriel $E_y \otimes_{\kappa(y)} K$.

2° Supposons que $Y$ soit affine d'anneau $A$, tel en outre que tout $\mathcal{O}_Y$-module inversible soit trivial (§ 1, n°3, Remarque) ; prenons en outre $E=\mathcal{O}_Y^n$ ; alors, dans le cor. 2, $V$ s'identifie à $A^n$ (I, 1.3, cor. 2 du th. 1), $V^*$ à l'ensemble des systèmes $(f_i)_{1 \le i \le n}$ engendrant l'idéal $A$ ; deux tels systèmes définissent la même $Y$-section de $\mathbb{P}_A^{n-1}$, autrement dit le même point de $\mathbb{P}_A^{n-1}$ à valeurs dans $A$ si et seulement si l'un se déduit de l'autre par multiplication par un élément inversible de $A$.

Ces propriétés justifient la terminologie de "fibré projectif" pour $\mathbb{P}(E)$.

Remarques. — 1) Le cor. 1 (pour $U=Y$ entraine inversement le th. 1, car les $Y$-morphismes de $X$ dans $\mathbb{P}$ s'identifient aux $X$-morphismes de $X$ dans $\mathbb{P} \times_Y X$, entraînant dit aux $X$-sections de $\mathbb{P} \times_Y X$ en vertu de la définition du produit.

2) Nous verrons au chap. IV, § que si $Y$ est localement noethérien et connexe, et si $E$ est cohérent et localement libre, alors tout faisceau inversible $\mathcal{L}$ sur $\mathbb{P}(E)$ est isomorphe à un faisceau de la forme $\mathcal{L}' \otimes_{\mathcal{O}_Y} \mathcal{O}_{\mathbb{P}}(m)$, où $\mathcal{L}'$ est un faisceau inversible sur $Y$, bien déterminé à un isomorphisme près, et $m$ est un entier bien déterminé. En d'autres termes, $H^1(\mathbb{P}, \mathcal{O}_{\mathbb{P}}^*)$ est isomorphe à $\mathbb{Z} \times H^1(Y, \mathcal{O}_Y^*)$ (§ 1, n°3). Nous verrons aussi que $p_*(\mathcal{L}) = 0$ si $m < 0$ et que $p_*(\mathcal{L})$ est isomorphe à $\mathcal{L}' \otimes_{\mathcal{O}_Y} S^m(E)$ si $m \ge 0$ ($S^m$ désignant ici

\newpage

%% ===== Page 283 =====

- II-94 -

le $\mathcal{O}_Y$-module puissance $n$-ème symétrique de $E$. Si $F$ est un $\mathcal{O}_Y$-module quasi-cohérent, tout $Y$-morphisme $P(E) \to P(F)$ est donc déterminé par la donnée d'un faisceau inversible $L'$ sur $Y$, d'un entier $m \ge 0$ et d'un $\mathcal{O}_Y$-homomorphisme $\psi : F \to L' \otimes_{\mathcal{O}_Y} S^m(E)$ tel que l'homomorphisme correspondant $\psi^\sharp$ (de faisceaux sur $P(E)$) soit surjectif. Nous verrons aussi que si le $Y$-morphisme considéré est un isomorphisme, alors $m=1$ et $F$ est isomorphe à $E \otimes_{\mathcal{O}_Y} L'^{\vee}$ (réciproque de la prop. 16). Ceci permettra aussi de déterminer le faisceau des germes d'automorphismes de $P(E)$ comme le quotient du faisceau de groupes $\underline{\operatorname{Aut}}(E)$ (localement isomorphe à $GL(n, \mathcal{O}_Y)$) par $\mathcal{O}_Y^*$.

Soient $E, F$ deux $\mathcal{O}_Y$-modules quasi-cohérents, posons $P_1 = P(E)$, $P_2 = P(F)$ et désignons par $p_1, p_2$ les morphismes structuraux $P_1 \to Y, P_2 \to Y$. Soit $Q = P_1 \times_Y P_2$, et soient $q_1, q_2$ les projections $Q \to P_1, Q \to P_2$ ; le faisceau $\mathcal{O}_Q(1) \otimes_Y \mathcal{O}_Q(1) = q_1^*(\mathcal{O}_{P_1}(1)) \otimes \mathcal{O}_Q q_2^*(\mathcal{O}_{P_2}(1))$ sur $Q$ est inversible puisque $q_1^*(\mathcal{O}_{P_1}(1))$ est localement isomorphe à $q_1^*(\mathcal{O}_{P_1}) = \mathcal{O}_Q$ $(i=1,2)$. D'autre part, si $p = p_1 \circ q_1 = p_2 \circ q_2$ est le morphisme structural $Q \to Y$, on a $p^*(E \otimes_{\mathcal{O}_Y} F) = q_1^*(p_1^*(E)) \otimes_{\mathcal{O}_Q} q_2^*(p_2^*(F))$ ; les homomorphismes surjectifs canoniques (17), $p_1^*(E) \to \mathcal{O}_{P_1}(1)$, $p_2^*(F) \to \mathcal{O}_{P_2}(1)$ donnent donc, par produit tensoriel, un homomorphisme
$$(20) \quad \nu : p^*(E \otimes_{\mathcal{O}_Y} F) \to \mathcal{O}_Q(1) \otimes \mathcal{O}_Q \mathcal{O}_Q(1)$$
évidemment surjectif ; on en déduit un morphisme canonique, dit morphisme de Veronese :
$$\nu : P(E) \times_Y P(F) \to P(E \otimes_{\mathcal{O}_Y} F)$$

Proposition 16. — Le morphisme de Veronese est une immersion fermée.

La question étant locale sur $Y$, on peut supposer $Y$ affine d'anneau $A$, de sorte que $E = \tilde{E}, F = \tilde{F}, E$ et $F$ étant deux $A$-modules, et $E \otimes_A F = (\widetilde{E \otimes_A F})$ ; posons $S^{(1)} = S_A(E)$, $S^{(2)} = S_A(F)$, $S = S_A(E \otimes_A F)$. Soient $f \in E, g \in F$ ; nous allons montrer que la restriction de $\nu$ à $D_+(f) \times_Y D_+(g)$ est une immersion fermée dans $D_+(f \otimes g)$.

\newpage

%% ===== Page 284 =====

- II-95 -

L'homomorphisme (20) correspond à l'homomorphisme $x \otimes y \to (x/f) \otimes (y/f)$ de $E \otimes_A F$ dans $L = (S^{(1)})_{(f)} \otimes_A (S^{(2)})_{(g)}$. Or, l'élément $c = (f/1) \otimes (g/1)$ est un générateur de $L$ considéré comme $S^{(1)} \otimes_A S^{(2)}$-module ($\S 3, n^\circ 3, \text{prop.} 13$); avec les notations de la démonstration du th. 2, on a donc $v_c(f \otimes g) = 1$, et l'application $\omega$ de $S^{(1 \otimes 2)}$ dans $S^{(1)} \otimes S^{(2)}$ est telle que $\omega((x \otimes y)/(f \otimes g)) = (x/f) \otimes (y/g)$, est un homomorphisme surjectif, ce qui établit notre assertion. En outre, $D_+(f) \times D_+(g)$ est l'image réciproque par $v$ de $D_+(f \otimes g)$ ; en effet, si un point $y \in \operatorname{Proj}(S^{(1)}) \times \operatorname{Proj}(S^{(2)})$ n'appartient pas à $D_+(f) \times D_+(g)$, au moins une des projections sur $\operatorname{Proj}(S^{(1)})$ et $\operatorname{Proj}(S^{(2)})$, n'appartient pas à $D_+(f)$ (resp. $D_+(g)$), et par suite, ou bien $p_1(f)$ s'annule au point $p_1(y)$, ou bien $p_2(g)$ s'annule au point $p_2(y)$ ; on en conclut que si $y \in D_+(f^c) \times D_+(g^c)$, $(f/1) \otimes (g/1)$ (dans $(S^{(1)})_{(f^c)} \otimes (S^{(2)})_{(g^c)}$) s'annule au point $y$ ; comme d'autre part par définition $f \otimes g$ ne s'annule en aucun point de $D_+(f \otimes g)$, aucun point de $D_+(f \otimes g)$ ne peut être l'image de $y$, ce qui prouve notre assertion, et établit la prop. 18, puisque les $D_+(f \otimes g)$ forment une base de la topologie de $\underline{\operatorname{Proj}}(E \otimes_A F)$ et les $D_+(f) \times D_+(g)$ une base de la topologie de $\underline{\operatorname{Proj}}(E) \times \underline{\operatorname{Proj}}(F)$.

Proposition 19. — Le morphisme de Veronese est fonctoriel en $E$ et $F$ pour les morphismes surjectifs de $\mathcal{O}_Y$-modules.

Il faut montrer que si $E \to E'$ est un $\mathcal{O}_Y$-homomorphisme surjectif, le diagramme
$$\begin{tikzcd}
\underline{\operatorname{Proj}}(E') \times \underline{\operatorname{Proj}}(F) \arrow[d, "j \times 1"] \arrow[r] & \underline{\operatorname{Proj}}(E' \otimes F) \arrow[d] \\
\underline{\operatorname{Proj}}(E) \times \underline{\operatorname{Proj}}(F) \arrow[r] & \underline{\operatorname{Proj}}(E \otimes F)
\end{tikzcd}$$
est commutatif ; $j$ étant l'immersion fermée $\underline{\operatorname{Proj}}(E') \to \underline{\operatorname{Proj}}(E)$. Comme $j \times 1$ est une immersion fermée ($I, 3, 2, \text{prop.} 5$) et que, si l'on pose $p_1 = \underline{\operatorname{Proj}}(E')$, on a $(j \times 1)^* (\mathcal{O}_{p_1}(1) \otimes \mathcal{O}_{p_2}(1)) = j^*(\mathcal{O}_{p_1}(1)) \otimes \mathcal{O}_{p_2}(1) = \mathcal{O}_{p_1}(1) \otimes \mathcal{O}_{p_2}(1)$ on voit que notre assertion découle de (14) et du $\S 1, n^\circ 1, \text{cor.} 1$ de la prop. 4.

\newpage

%% ===== Page 285 =====

- II-96 -

sertion résulte de la prop. 17 (i) et (ii).

\textsc{Proposition 20.} — Soit $\psi: Y' \to Y$ un morphisme. Si $E' = \psi^*(E)$, $F' = \psi^*(F)$, le morphisme de Veronese $\mathbb{P}(E') \times_Y \mathbb{P}(F') \to \mathbb{P}(E' \otimes F')$ s'identifie à $\psi^Y$.

Posons $\mathbb{P}_1 = \mathbb{P}(E')$, $\mathbb{P}_2 = \mathbb{P}(F')$; on sait (formule (15)) que $\mathbb{P}_i' \simeq \mathbb{P}_i \times_Y Y' (i=1,2)$, donc le morphisme $\mathbb{P}_1' \times_{Y'} \mathbb{P}_2' \to Y'$ s'identifie à $\mathbb{P}_1 \times_Y \mathbb{P}_2 \to Y$. D'autre part, $E' \otimes F'$ s'identifie à $\psi^*(E \otimes F)$, donc $\mathbb{P}(E' \otimes F') \simeq (\mathbb{P}(E \otimes F))_Y \times Y'$ (formule (15)). Enfin $\mathcal{O}_{\mathbb{P}_1'}(1) \otimes \mathcal{O}_{\mathbb{P}_2'}(1) \simeq L' \otimes_{Y'} \mathcal{O}_{Y'}$ s'identifie à $L \otimes_Y Y'$, en vertu de la formule (16) de ce n° et du cor. de la pro. 6 du § 1, n°1. L'homomorphisme canonique $\psi^*(E \otimes F) \to L'$ s'identifie donc à $\psi^Y$, et la proposition résulte de la prop. 17, (iii).

\textsc{6. Morphismes projectifs.}

\textsc{Définition 2.} — On dit qu'un $Y$-préschéma $X$ est projectif si $X$ est $Y$-isomorphe à un sous-préschéma fermé de $\mathbb{P}(E)$, où $E$ est un $\mathcal{O}_Y$-module quasi-cohérent de type fini. On dit qu'un morphisme $X \to Y$ est projectif s'il fait de $X$ un $Y$-préschéma projectif.

On notera que cette définition n'est pas de nature locale sur $Y$; [les contre-exemples de Nagata (Illinois J. Math.) et Hironaka ( ] montrent que, même si $X$ et $Y$ sont des schémas algébriques non singuliers sur un corps algébriquement clos, $X$ peut être localement projectif sur $Y$ sans être projectif sur $Y$.

\textsc{Proposition 21.} — (i) Soit $S$ une $\mathcal{O}_Y$-algèbre graduée quasi-cohérente à degrés positifs. Supposons qu'il existe un recouvrement ouvert fini $(U_i)$ de $Y$ tel que chacune des algèbres graduées $\Gamma(U_i, S)$ soit de type fini sur $\Gamma(U_i, \mathcal{O}_Y)$. Alors $\operatorname{Proj}(S)$ est projectif sur $Y$.
(ii) Inversement, si $Y$ est noethérien, tout préschéma $X$ projectif sur $Y$ est isomorphe à un préschéma $\operatorname{Proj}(S)$, où $S$ est une $\mathcal{O}_Y$-algèbre graduée quasi-cohérente à degrés positifs, engendrée par $S_1$, et $S_1$ est de type fini (et par suite cohérent).

\newpage

%% ===== Page 286 =====

(i) résulte en effet de la prop. 3 (i) du n° 1 et de son corollaire, ainsi que du cor. de la prop. 15 du n° 4 ; (ii) résulte de la prop. 15 (ii) du n° 4 .

\emph{Proposition 22}. -- Un morphisme projectif est séparé et de type fini.

On sait en effet qu'il en est ainsi du morphisme structural $P(\mathcal{E}) \to Y$ (n° 1) et une immersion fermée est séparée (I, 3, 6, prop. 16) et de type fini (I, 4, 2, prop. 7).

\emph{Théorème 3}. -- (i) Une immersion fermée est un morphisme projectif.
(ii) Si $f: X \to Y$ et $g: Y \to Z$ sont des morphismes projectifs et si $Z$ est noethérien, $g \circ f$ est projectif.
(iii) Si $f: X \to Y$ est un $S$-morphisme projectif, $f^{S'}: X^{S'} \to Y^{S'}$ est projectif.
(iv) Si $f: X \to Y$ et $g: X' \to Y'$ sont deux $S$-morphismes projectifs, $f \times_S g: X \times_S X' \to Y \times_S Y'$ est projectif.
(v) Soient $f: X \to Y$ et $g: Y \to Z$ deux morphismes tels que $g \circ f$ soit projectif. Si $g$ est séparé et $Y$ noethérien, $f$ est projectif.
(vi) Si $f: X \to Y$ est projectif et si $Y_{\text{red}}$ est noethérien, $f_{\text{red}}$ est projectif.

(1) résulte de la prop. 21 (i) et du cor. 3 de la prop. 2 du n° 1. Pour prouver (iii), on identifie $X^{S'}$ à $X \times_Y Y'$ (I, 2, 5, cor. 2 de la prop. 9) et la proposition résulte de la formule (15) du n° 5 lorsque $X = P(\mathcal{E})$ ; si $X$ est un sous-préschéma fermé de $P(\mathcal{E})$, $X^{S'}$ est un sous-préschéma fermé de $(P(\mathcal{E}))^{S'}$ (I, 3, 2, cor. 1 de la prop. 5), d'où (iii) en général. De même, pour démontrer (iv), on est aussitôt ramené au cas où $X = P(\mathcal{E}), X' = P(\mathcal{E}')$ ; soient $p, p'$ les projections de $X \times_S X' = T$ dans $Y$ et $Y'$ respectivement ; d'après la formule (15) du n° 5, on a $P(p^*(\mathcal{E})) = P(\mathcal{E}) \times_Y T$ et $P(p'^*(\mathcal{E}')) = P(\mathcal{E}') \times_{Y'} T$ ; d'où
$$P(p^*(\mathcal{E})) \times_T P(p'^*(\mathcal{E}')) = (P(\mathcal{E}) \times_Y T) \times_T (T \times_{Y'} P(\mathcal{E}')) = P(\mathcal{E}) \times_Y (T \times_{Y'} P(\mathcal{E}'))$$
$$= P(\mathcal{E}) \times_Y (Y \times_S P(\mathcal{E}')) = P(\mathcal{E}) \times_S P(\mathcal{E}') \text{ . Or, } p^*(\mathcal{E}) \text{ et } p'^*(\mathcal{E}') \text{ sont de ty-}$$

\newpage

%% ===== Page 287 =====

- II-98 -

pe fini sur $T$, donc il en est de même de $p^*(\underline{E}) \otimes_T p'^*(\underline{E}')$ ; comme $P(p^*(\underline{E}) \otimes_T p'^*(\underline{E}'))$ s'identifie à un sous-préschéma fermé de $P(p^*(\underline{E}) \otimes_T p'^*(\underline{E}'))$ d'après la prop.18, cela achève de prouver (iv). (v) et (vi) se déduisent des autres assertions par les raisonnements habituels (tenant compte de ce que $X_{\text{red}} \to X$ est une immersion fermée). Reste donc à prouver (ii). Comme $Y$ est de type fini sur $Z$ (prop.22) et que $Z$ est supposé noethérien, $Y$ est noethérien. Par hypothèse, $X$ est isomorphe à un sous-préschéma fermé de $P(\underline{E})$, où $\underline{E}$ est un $\mathcal{O}_Y$-module quasicohérent de type fini, et par suite cohérent. On en conclut ($n^\circ 4$, cor.4 du th.1) que $\underline{E}$ est isomorphe au quotient d'un faisceau de la forme $g^*(\mathcal{F})(n)$, où $\mathcal{F}$ est un $\mathcal{O}_Z$-module cohérent. Donc ($n^\circ 5$) $P(\underline{E})$ est isomorphe à un sous-préschéma fermé de $P(g^*(\mathcal{F})(n))$. Or, comme $g^*(\mathcal{F})(n) = g^*(\mathcal{F}) \otimes_{\mathcal{O}_Y} \mathcal{O}_Y(n)$ et que $\mathcal{O}_Y(n)$ est inversible, $P(g^*(\mathcal{F})(n))$ est isomorphe à $P(g^*(\mathcal{F}))$, et $P(g^*(\mathcal{F}))$ est d'autre part isomorphe à $P(\mathcal{F}) \times_Z Y$ ($n^\circ 5$, formule (15)). En vertu de (iv) et de l'hypothèse, $P(\underline{E}) \times_Y Y$ est projectif sur $Z$, donc il en est de même de $X$, qui s'identifie à un sous-préschéma fermé de $P(\underline{E}) \times_Z Y$.

Théorème 4. - Soit $Y$ un préschéma localement noethérien. Tout morphisme projectif $f: X \to Y$ est une application fermée de l'espace de base de $X$ dans celui de $Y$.

La question étant locale sur $Y$ (cor. du th.3), on peut supposer $Y$ affine, donc noethérien ; il est clair en outre qu'on peut supposer $X = P(\underline{E}) = \operatorname{Proj}(S)$, où $S = \mathcal{S}_{\mathcal{O}_Y}(\underline{E})$, $\underline{E}$ étant de type fini sur $\mathcal{O}_Y$, donc cohérent. [le texte ici est barré : En outre, toute partie fermée de l'espace de base de $X$ est support d'un sous-préschéma fermé (I, 3, 4, prop.9) ; mais un tel sous-préschéma est lui-même isomorphe à un préschéma $P(\mathcal{F})$ (prop.21 (ii)).] Finalement on est donc ramené à prouver que $f(X)$ est fermé dans $Y$. Or,

\newpage

%% ===== Page 288 =====

- 11-99 -

pour tout $y \in Y$, la fibre $f^{-1}(y)$ s'identifie, en tant que préschéma, à $\operatorname{Proj}(\underline{S}) \times_Y \operatorname{Spec}(k(y))$ (I, 2, 7, prop. 17), donc à $\mathbb{P}(\underline{E}_y / \mathfrak{m}_y \underline{E}_y)$, fibré projectif sur $\operatorname{Spec}(k(y))$ (n° 5, formule (15)). Cette fibre est vide si et seulement si $S_{k(y)}(\underline{E}_y / \mathfrak{m}_y \underline{E}_y)$ satisfait à la condition (TN) (n° 3, prop. 12 (ii)), et comme $k(y)$ est un corps, cela signifie que l'espace vectoriel $\underline{E}_y / \mathfrak{m}_y \underline{E}_y$ est réduit à $0$ ; mais comme $\underline{E}_y$ est un module de type fini sur $\mathcal{O}_Y$, cela signifie que $\underline{E}_y = 0$, en vertu du lemme de Nakayama. Or, comme $\underline{E}$ est un $\mathcal{O}_Y$-module cohérent, l'ensemble des $y \in Y$ tels que $\underline{E}_y = 0$ est ouvert (FAC, I, 24, 12, prop. 2). On voit par suite que le complémentaire de $f(X)$ dans $Y$ est ouvert.

\emph{Corollaire}. - Soient $X$ un $Y$-préschéma projectif, $Z$ un $Y$-préschéma noethérien et séparé sur $Y$. Tout $Y$-morphisme de $X$ dans $Z$ est une application fermée.

Si $f: X \to Z$ est un tel morphisme et $g: Z \to Y$ le morphisme structural de $Y$. Il suffit de remarquer que $g \circ f$ est projectif et $g$ séparé par hypothèse ; on peut donc appliquer le th. 3 (v) qui prouve que $f$ est projectif.

7. Morphismes quasi-projectifs.

\emph{Définition 3}. - On dit qu'un $Y$-préschéma $X$ est quasi-projectif s'il est isomorphe à un sous-préschéma d'un $Y$-préschéma projectif. Un morphisme $f: X \to Y$ est dit quasi-projectif s'il fait de $X$ un $Y$-préschéma quasi-projectif.

Il est clair qu'un morphisme projectif est quasi-projectif.

\emph{Proposition 23}. - Un morphisme quasi-projectif $X \to Y$ est séparé ; il est de type fini lorsque $Y$ est localement noethérien.

En effet, si $X$ est un sous-préschéma d'un $Y$-préschéma projectif $Z$, $X$ est séparé sur $Z$ (I, 2, 6, prop. 16) ; d'autre part $Z$ est de type fini sur $Y$ (prop. 22), donc localement noethérien si $Y$ l'est ; dans ce cas, $X$ est donc de type fini sur $Z$ (I, 4, 2, prop. 7), d'où la conclusion (I, 4, 2, prop. 8).

\newpage

%% ===== Page 289 =====

\marginpar{Archives \\ Grothendieck sept 59}

$$- \text{II-}100 -$$

\textsc{Proposition} 24. — Soit $Y$ un préschéma localement noethérien. Pour qu'un $Y$-préschéma $X$ soit quasi-projectif, il faut et il suffit que $X$ soit isomorphe à un sous-préschéma induit sur un ouvert $U$ d'un $Y$-préschéma projectif $Z$. On peut même supposer que $U$ est dense dans $Z$.

La condition est évidemment suffisante ; elle est nécessaire, car si $X$ est sous-préschéma d'un $Y$-préschéma projectif $Z$, $Z$ est localement noethérien (puisque de type fini sur $Y$), donc $X$ est le préschéma induit par un sous-préschéma fermé $X'$ de $Z$, et $X$ est ouvert dans $X'$ ($\S$ 1, n°2, prop. 8) ; comme $X'$ est projectif sur $Y$, la conclusion en résulte.

\textsc{Proposition} 25. — (i) Un morphisme d'immersion est quasi-projectif.
(ii) Si $f: X \to Y$ et $g: Y \to Z$ sont des morphismes quasi-projectifs, et si $Y$ est noethérien, $g \circ f$ est quasi-projectif.
(iii) Si $f: X \to Y$ est un $S$-morphisme quasi-projectif, $f^{S'}: X^{S'} \to Y^{S'}$ est quasi-projectif.
(iv) Si $f: X \to Y$ et $g: X' \to Y'$ sont deux $S$-morphismes quasi-projectifs, $f \times_S g : X \times_S X' \to Y \times_S Y'$ est quasi-projectif.
(v) Soient $f: X \to Y$ et $g: Y \to Z$ deux morphismes tels que $g \circ f$ soit quasi-projectif. Si $g$ est séparé et $Y$ noethérien, $f$ est quasi-projectif.
(vi) Si $f: X \to Y$ est quasi-projectif et si $Y_{\text{red}}$ est noethérien, $f_{\text{red}}$ est quasi-projectif.

(i), (iii) et (iv) sont des conséquences immédiates du th. 3 (i), (iii) et (iv), compte tenu de ce que si $h: X \to Z$ est une immersion, $h^{S'}: X^{S'} \to Z^{S'}$ est une immersion (I, 3, 2, cor. 1 de la prop. 5) et si $h': X' \to Z'$ est une seconde immersion, $h \times h'$ est une immersion (I, 3, 2, prop. 5). (v) et (vi) se déduisent comme d'ordinaire des autres assertions (compte tenu de ce que $X_{\text{red}}$, étant un sous-préschéma

\newpage

%% ===== Page 290 =====

\noindent -- II-101 --

d'un $Y$-préschéma quasi-projectif, est quasi-projectif sur $Y$. Reste donc à démontrer (ii). On peut évidemment se ramener au cas où $X = \mathbb{P}(\mathcal{E})$, où $\mathcal{E}$ est un $\mathcal{O}_Y$-module cohérent. D'autre part, en vertu de la prop. 24, on peut identifier $X$ au sous-préschéma induit sur un ouvert d'un $Z$-préschéma projectif $Y'$. En vertu du § 1, n°2, th. 1 ; $\mathcal{E}$ est la restriction à $Y$ d'un $\mathcal{O}_{Y'}$-module cohérent $\mathcal{E}'$, donc, par définition (n°1, prop. 1), $\mathbb{P}(\mathcal{E})$ s'identifie à un sous-préschéma induit sur une partie ouverte de $\mathbb{P}(\mathcal{E}')$. En vertu du th. 3 (ii), $\mathbb{P}(\mathcal{E}')$ est projectif sur $Z$, donc $\mathbb{P}(\mathcal{E})$ est quasi-projectif sur $Z$.

\medskip

\noindent Proposition 26. — Soit $\mathcal{E}$ un faisceau quasi-cohérent sur un préschéma $Y$. Il existe une immersion canonique du fibré vectoriel $\mathbb{V}(\mathcal{E})$ dans le fibré projectif $\mathbb{P}(\mathcal{E} \oplus \mathcal{O}_Y)$.

\medskip

\noindent Posons $\mathcal{E}' = \mathcal{E} \oplus \mathcal{O}_Y$, $\mathcal{S} = \mathcal{S}(\mathcal{E})$, $\mathcal{S}' = \mathcal{S}(\mathcal{E}')$, $\mathbb{P} = \mathbb{P}(\mathcal{S}')$, et soit $f$ la section de $\mathcal{E}' \to S$ au-dessus de $Y$ qui en chaque point $y \in Y$ est égale à l'élément unité de $\mathcal{O}_y$. Le sous-préschéma induit sur la partie ouverte $\mathbb{P}_f$ de $\mathbb{P}$ définie par la section $f$ est affine au-dessus de $Y$ et canoniquement isomorphe au préschéma associé au faisceau d'algèbres $\mathcal{S}'/(f-1)\mathcal{S}'$ (n°1, prop. 2) ; mais $\mathcal{S}'$ est canoniquement isomorphe à $\mathcal{S} \otimes_{\mathcal{O}_Y} \mathcal{S}(\mathcal{O}_Y) = \mathcal{S} \otimes_{\mathcal{O}_Y} \mathcal{O}_Y[T]$ ($T$ indéterminée), $f$ s'identifiant à $1 \otimes T$ ; on en conclut aussitôt que $\mathcal{S}'/(f-1)\mathcal{S}'$ s'identifie canoniquement à $\mathcal{S}$, d'où la conclusion. On dit que $\mathbb{P}(\mathcal{E} \oplus \mathcal{O}_Y)$ est la fermeture projective de $\mathbb{V}(\mathcal{E})$.

\medskip

\noindent Proposition 27. — Soit $Y$ un préschéma noethérien. Tout morphisme affine de type fini $f : X \to Y$ est quasi-projectif.

\medskip

\noindent En effet, $\mathcal{A} = f_*(\mathcal{O}_X)$ est un faisceau quasi-cohérent d'algèbres sur $Y$ (§ 2, n°2, prop. 5) qui est de type fini par hypothèse. Comme $Y$ est noethérien, il y a donc un recouvrement fini $(U_i)$ de $Y$ par des ouverts affines tels que chaque algèbre $\mathcal{A}(U_i)$ soit engendrée par un nombre fini de ses éléments $s_{ij}$ ; soit $\mathcal{E}_i$ le faisceau quasi-cohérent

\newpage

%% ===== Page 291 =====

-- 102 --
cohérent de modules sur $U_i$ engendré par les $s_{ij}$, qui est un sous-faisceau de $A|U_i$ (considéré comme faisceau de modules) ; comme $Y$ est noethérien et $E_i$ de type fini, $E_i$ est cohérent, donc est la restriction à $U_i$ d'un sous-faisceau de modules cohérent $E'$ de $A$ (§ 1, n°2, th. 1 (iii)). Soit $E$ le $\mathcal{O}_Y$-module cohérent somme des $E'_i$ ; il résulte aussitôt des définitions que l'injection $E \to A$ se prolonge en un homomorphisme surjectif $S_{\mathcal{O}_Y}(E) \to A$, autrement dit (§ 2, n°3, prop. 7 et I, 3, 1, prop. 1) $X$ est isomorphe à un sous-préschéma fermé de $\mathbb{V}(E)$ ; la proposition résulte donc de la prop. 26.

\medskip

Proposition 28. -- Tout morphisme fini $f: X \to Y$ est projectif.

Avec les notations de la prop. 27, $A = f_*(\mathcal{O}_X)$ est ici un $\mathcal{O}_Y$-module de type fini et quasi-cohérent ; l'application identique $A \to A$ se prolonge en un homomorphisme surjectif $S_{\mathcal{O}_Y}(A) \to A$ d'algèbres, donc $X$ s'identifie à un sous-préschéma fermé de $\mathbb{P}(A)$ (prop. 25), qui lui-même s'identifie à un sous-préschéma induit sur un ouvert de $\mathbb{P}(\mathcal{O}_Y)$. D'autre part, comme $X$ est fini sur $Y$, il est à fortiori fini sur $P$ puisque $P$ est séparé sur $Y$ (§ 2, n°7, prop. 15 (v)) ; donc l'immersion $X \to P$ est fermée (§ 2, n°7, prop. 16), d'où la proposition.

\medskip

Proposition 29. -- Soient $Y$ un schéma affine noethérien, $X$ un schéma quasi-projectif sur $Y$. Pour tout $\mathcal{O}_X$-module cohérent $F$, il existe un entier $n > 0$ et un $\mathcal{O}_X$-module inversible $L$ tels que $F$ soit isomorphe à un quotient de $\mathcal{O}_X^n \otimes L = L^n$.

En effet, $X$ peut être identifié à un préschéma induit sur un ouvert d'un $Y$-schéma projectif $X'$ ; le faisceau cohérent $F$ est alors restriction à $X$ d'un $\mathcal{O}_{X'}$-module cohérent $F'$ (§ 1, n°2, th. 1 (i)) puisque $X'$ est noethérien ; mais on sait alors que $F'$ est isomorphe à un quotient de faisceau de la forme $p^* (E) \otimes_{\mathcal{O}_{X'}} \mathcal{O}_{X'}(m)$, où $m$ est un entier rationnel et $E$ un $\mathcal{O}_Y$-module cohérent (n°4, cor. 4 du th. 1). Comme $Y$ est affine, il existe $n > 0$ tel que $E$ soit isomorphe...

\newpage

%% ===== Page 292 =====

-- II-103 --

phe à un quotient d'un faisceau $\mathcal{O}_Y^n$ ($I, 1, 5, \text{th.} 3$ et $I, 1, 3, \text{cor.} 4$ du th. 1) ; comme $f^*$ est exact à droite, $f^*(\mathcal{E})$ est un quotient de $f^*(\mathcal{O}_Y^n) = \mathcal{O}_X^n$, d'où la proposition.

Corollaire. -- Sous les hypothèses de la prop. 29, pour qu'un $X$-préschéma $X'$ soit projectif (resp. quasi-projectif), il faut et il suffit qu'il soit isomorphe à un sous-schéma fermé (resp. à un sous-schéma) d'un fibré projectif de la forme $\mathbb{P}_X^n$.

En effet, la prop. 29 montre que $\mathbb{P}(\mathcal{F})$ est isomorphe à un sous-schéma fermé d'un schéma $\mathbb{P}_X^{n-1}$ ($n^\circ 5$ et prop. 16 du $n^\circ 5$).

On ignore si la prop. 29 et son corollaire sont valables pour un préschéma noethérien quelconque $X$.

8. Faisceaux très amples et faisceaux amples.

Définition 4. -- Soit $X$ un $Y$-préschéma, $f : X \to Y$ son morphisme structural. On dit qu'un faisceau inversible $\mathcal{L}$ sur $X$ est très ample pour $f$ (ou simplement très ample, si aucune confusion n'en résulte) s'il existe un $\mathcal{O}_Y$-module quasi-cohérent $\mathcal{E}$ de type fini et une $Y$-immersion $i$ de $X$ dans $\mathbb{P}(\mathcal{E})$ tels que $\mathcal{L} = i^*(\mathcal{O}_{\mathbb{P}(\mathcal{E})}(1))$. On dit qu'un faisceau inversible $\mathcal{L}$ sur $X$ est ample pour $f$ (ou simplement ample) s'il existe un entier $n > 0$ tel que $\mathcal{L}^{\otimes n}$ soit très ample pour $f$.

Il est clair qu'un faisceau très ample est ample. Pour qu'il existe sur $X$ un faisceau très ample, il faut et il suffit que $X$ soit quasi-projectif sur $Y$.

Proposition 30. -- Soient $\mathcal{L}, \mathcal{L}'$ deux faisceaux inversibles sur $X$. Si $\mathcal{L}$ et $\mathcal{L}'$ sont très amples (resp. amples) pour un morphisme $f : X \to Y$, il en est de même de $\mathcal{L} \otimes \mathcal{L}'$.

En effet, soient $\mathcal{E}, \mathcal{E}'$ deux $\mathcal{O}_Y$-modules quasi-cohérents de type fini, $i, i'$ des immersions de $X$ dans $\mathbb{P}(\mathcal{E})$ et $\mathbb{P}(\mathcal{E}')$ respectivement, telles que $\mathcal{L} = i^*(\mathcal{O}_{\mathbb{P}(\mathcal{E})}(1))$, $\mathcal{L}' = i'^*(\mathcal{O}_{\mathbb{P}(\mathcal{E}')}(1))$. Posons $\mathcal{Q} = \mathcal{E} \otimes \mathcal{E}'$, et considérons le morphisme de Veronese $v : \mathbb{P}(\mathcal{E}) \times_Y \mathbb{P}(\mathcal{E}') \to \mathbb{P}(\mathcal{Q})$ ; comme $i$ et $i'$ sont...

\newpage

%% ===== Page 293 =====

- II-104 -

des immersions, il en est de même de $(r, r')_Y : X \to \mathbb{P}_Y \times \mathbb{P}'$, (I, 3, 5, cor 3 de la prop. 15), et comme $v$ est une immersion (n° 5, prop. 18), $r'' : X \xrightarrow{(r, r')_Y} \mathbb{P} \times_Y \mathbb{P}' \to \mathbb{Q}$ est aussi une immersion. D'autre part (n° 5, cor. de la prop. 18), $v^*(\mathcal{O}_{\mathbb{Q}}(1))$ est isomorphe à $\mathcal{O}_{\mathbb{P}}(1) \otimes_Y \mathcal{O}_{\mathbb{P}'}(1)$, donc ($§$ 1, n° 1, prop. 4), $r''^*(\mathcal{O}_{\mathbb{Q}}(1))$ est isomorphe à $L \otimes_Y L'$, ce qui démontre la proposition pour les faisceaux très amples. Comme $(L \otimes_Y L')^{\otimes n} = L^{\otimes n} \otimes_Y L'^{\otimes n}$, on en déduit la proposition pour les faisceaux amples.

Remarque. - Dans la démonstration de la prop. 30, pour prouver que $(r, r')_Y$ est une immersion, il suffit que l'un des morphismes $r, r'$ soit une immersion ; on voit donc qu'on peut affirmer que $L \otimes_Y L'$ est très ample lorsque $L$ est très ample et qu'il existe un homomorphisme surjectif $q^*(E') \to L'$, où $q : X \to Y$ est le morphisme structural, et $E'$ un $\mathcal{O}_Y$-module quasi-cohérent de type fini.

Proposition 31. - (i) Soit $\psi : Y' \to Y$ un morphisme. Si $L$ est un faisceau sur $X$ très ample (resp. ample) pour $f : X \to Y$, $L' = L \otimes_Y \mathcal{O}_{Y'}$ est un faisceau très ample (resp. ample) sur $X' = X \times_Y Y'$ pour $f' : X' \to Y'$.

(ii) Soient $K$ un faisceau inversible sur $Y$, $L$ un faisceau inversible sur $X$. Pour que $L$ soit très ample (resp. ample) pour $f : X \to Y$, il faut et il suffit que $L \otimes f^*(K)$ soit très ample (resp. ample) pour $f$.

(iii) Soient $X, X'$ deux $Y$-préschémas, $L$ et $L'$ des faisceaux très amples (resp. amples) sur $X$ et $X'$ respectivement ; alors $L \otimes_Y L'$ est très ample (resp. ample) sur $X \times_Y X'$.

(iv) Soient $f : X \to Y$, $g : Y \to Z$ deux morphismes, $Z$ étant supposé noethérien. Soit $L$ un faisceau très ample sur $X$ pour le morphisme $f$, $K$ un faisceau ample sur $Y$ pour le morphisme $g$. Alors il existe un entier $n > 0$ tel que $L \otimes f^*(K^{\otimes n})$ soit très ample pour $g \circ f$.

(v) résulte de ce que, si $r : X \to \mathbb{P}(E) = \mathbb{P}$ est une immersion, il en est de même de $r' : X' \to \mathbb{P}(E') = \mathbb{P}'$, où $E' = \psi^*(E)$ (I, 3, 2, cor. 1 de la ...

\newpage

%% ===== Page 294 =====

\section*{II-104 bis -}

pro. 5), et on a $(r_Y')^* (\mathcal{O}_{\mathbb{P}'}(1)) = \mathbb{L}'$ ($n^\circ 5$, prop. 17), et pour prouver (ii), notons que d'après la remarque suivant la prop. 30, si $\mathbb{L}$ est ample il en est de même de la réciproque $L \otimes f^* (K) \otimes f^* (K^{-1})$ ; enfin, comme $(L \otimes f^* (K))^{\otimes n} = L^{\otimes n} \otimes f^* (K^{\otimes n})$, l'assertion (ii) en résulte pour les faisceaux amples. La démonstration de (iii) est tout à fait analogue à celle de la prop. 30, utilisant le fait que si on a des immersions $r : X \to \mathbb{P}(E) = P$, $r' : X' \to \mathbb{P}(E') = P'$ alors $r \times r'$ est une immersion de $X \times X'$ dans $\mathbb{P}(E) \times \mathbb{P}(E') = P \times P'$ (I, 3.2, prop. 5), et de ce que l'image réciproque de $\mathcal{O}_{\mathbb{P}}(1) \otimes \mathcal{O}_{\mathbb{P}'}(1)$ par cette immersion est $\mathbb{L} \otimes_Y \mathbb{L}'$ ($\S 1$, $n^\circ 1$, cor. 1 de la prop. 4). Enfin, pour établir (iv), remarquons que l'on peut en remplaçant $K$ par $K^{\otimes n}$, supposer que $K$ est très ample ; on a alors des immersions $r : X \to \mathbb{P}(E) = P$, $r' : Y \to \mathbb{P}(G) = P'$ avec $\mathbb{L} = r^* (\mathcal{O}_{\mathbb{P}}(1))$, $\mathbb{K} = r'^* (\mathcal{O}_{\mathbb{P}'}(1))$; en outre, en vertu des démonstrations du th. 3 du $n^\circ 6$ et de la prop. 25 du $n^\circ 7$, il existe une immersion $s$ de $\mathbb{P}(E) = P$ dans $\mathbb{P}(F) = P''$ (où $F$ est un $\mathcal{O}_Z$-module cohérent) telle que $s^* (\mathcal{O}_{\mathbb{P}''}(1)) = \mathcal{O}_{\mathbb{P}}(1)$. En outre, $\mathbb{P}(F)$ s'identifie à $\mathbb{P}(F) \times_Z Y$, et $\mathcal{O}_{\mathbb{P}''}(1)$ à $\mathcal{O}_{\mathbb{P}(F)}(1) \otimes_{\mathcal{O}_Z} \mathcal{O}_Y$ ($n^\circ 5$, formule (16)), donc $\mathbb{L}$ est l'image réciproque de ce dernier faisceau par $s \circ r$ ; $\mathbb{L} \otimes f^* (K)$ est donc l'image réciproque par $s \circ r$ de $\mathcal{O}_{\mathbb{P}(F)}(1) \otimes_{\mathcal{O}_Z} \mathbb{K}$. Mais $r'$ définit une immersion $1 \times r'$ de $\mathbb{P}(F) \times_Z Y$ dans $\mathbb{P}(F) \times_Z \mathbb{P}(G)$, et $\mathcal{O}_{\mathbb{P}(F)}(1) \otimes_{\mathcal{O}_Z} \mathbb{K}$ est l'image réciproque par $1 \times r'$ de $\mathcal{O}_{\mathbb{P}(F)}(1) \otimes_{\mathcal{O}_Z} \mathcal{O}_{\mathbb{P}(G)}(1)$ ; par suite $\mathbb{L} \otimes f^* (K)$ est l'image réciproque de ce dernier faisceau par une immersion de $X$ dans $\mathbb{P}(F) \times_Z \mathbb{P}(G)$ ; il suffit enfin d'utiliser le morphisme de Veronese pour ce dernier proschéma pour établir que $\mathbb{L} \otimes f^* (K)$ est très ample pour $g \circ f$.

Corollaire 1. - Soient $X$ un $Y$-proschéma, $f : X \to Y$ son morphisme structural, $g : X \to Y'$ un $Y$-morphisme. Si un faisceau inversible $\mathbb{L}$ sur

\newpage

%% ===== Page 295 =====

- II-104 ter -

$X$ est très ample (resp. ample) pour $f$, il est très ample (resp. ample) pour $g$.

En effet, le morphisme $(1,g) : X \to X \times_Y Y'$ n'est autre que le morphisme graphe de $g$, donc est une immersion (I, 3, 5, cor. 3 de la prop. 15), et $L$ est l'image réciproque de $L' = L \otimes_Y \mathcal{O}_{Y'}$ par cette immersion ; en vertu de la prop. 31 (i), $L'$ est l'image réciproque de $\mathcal{O}_{P'}(1)$ par une immersion de $X \times_Y Y'$ dans $P' = \mathbb{P}(\mathcal{E}')$, d'où le corollaire.

Corollaire 2. — Soient $f : X \to Y$, $g : Y \to Z$ deux morphismes, $Z$ étant supposé noethérien. Si $L$ est un faisceau ample sur $X$ pour $f$, $K$ un faisceau ample sur $Y$ pour $g$, $L \otimes f^*(K)$ est ample sur $X$ pour $g \circ f$.

Comme $(L \otimes f^*(K))^{\otimes n} = L^{\otimes n} \otimes f^*(K^{\otimes n})$, cela résulte de la prop. 31 (iv).

Corollaire 3. — Si, dans le cor. 2, on suppose de plus $f$ fini, alors $f^*(K)$ est ample pour $g \circ f$.

En effet, en vertu du cor. 2, il suffit de montrer que $\mathcal{O}_X$ est alors très ample pour $f$. Or, pour l'immersion de $X' = \mathbb{P}(\mathcal{E})$ dans $\mathbb{P}(\mathcal{E} \oplus \mathcal{O}_Y) = \mathbb{P}$ définie dans la prop. 26, on constate sans peine que l'image réciproque de $\mathcal{O}_P(1)$ n'est autre que $\mathcal{O}_X$. D'autre part, si $\mathcal{A} = f_*(\mathcal{O}_X)$, on a défini dans la prop. 28 une immersion de $X$ dans $\mathbb{P}(\mathcal{A})$, et l'image réciproque de $\mathcal{O}_{\mathbb{P}(\mathcal{A})}(1)$ par cette immersion est $\mathcal{O}_X$ ; d'où notre assertion.

9. Caractérisation des faisceaux amples.

Soient $g : X \to Y$ un morphisme de préschémas, $\mathcal{S}$ une $\mathcal{O}_Y$-algèbre graduée en degrés positifs ; $g^*(\mathcal{S})$ est alors une $\mathcal{O}_X$-algèbre graduée (chap. 0, §9). Considérons sur $X$ un faisceau inversible $L$, et la $\mathcal{O}_Y$-algèbre graduée $S' = \bigoplus_{n \ge 0} L^{\otimes n}$ ; supposons donné un $\mathcal{O}_X$-homomorphisme $\phi : g^*(\mathcal{S}) \to S'$.

\newpage

%% ===== Page 296 =====

- II-104-quater -

de $\mathcal{O}_X$-algèbres graduées $\varphi : q^*(S) \to S'$, tel que $\varphi_n : q^*(S_n) \to L^{\otimes n}$ soit surjectif lorsque $n$ est assez grand. Nous allons voir qu'on peut déduire de $\varphi$ un morphisme $\psi : X \to \operatorname{Proj}(S)$ de façon canonique, en généralisant le procédé décrit au $n^\circ 5$ (que les considérations qui suivent sont destinées à remplacer).

Soit $x \in X$ ; pour définir $\psi(x) \in \operatorname{Proj}(S)$, on peut supposer d'abord que $Y = \operatorname{Spec}(A)$ est affine et $S = \tilde{S}$, où $S$ est une $A$-algèbre graduée. Pour tout $n > 0$, soit $p_n$ l'ensemble des $s \in S_n$ tels que $\varphi_n(s)$ s'annule au point $x$. L'hypothèse de surjectivité et le fait que $X$ est le support de chacun des $L^{\otimes n}$ impliquent que $p_n \cdot S_n$ pour $n$ assez grand ; les $p_n$ vérifient par suite les conditions de la prop. 2 du § 3, $n^\circ 1$, et on définit $p = \psi(x)$ comme l'idéal premier gradué de $S$ tel que $p \cap S_n = p_n$ pour tout $n > 0$. Pour passer au cas général, il suffit de vérifier que lorsqu'on passe d'un ouvert affine $U \subset Y$ à un ouvert affine $V \subset U$, la définition de $\psi(x)$ ne change pas, en supposant que $q(x) \in V \subset U$ ; ce qui est immédiat, compte tenu de la définition de $\operatorname{Proj}(S)$ ($n^\circ 1$).

Pour définir $\psi$ dans le voisinage de $x \in X$, considérons encore en premier lieu le cas où $X$ et $Y$ sont affines, et soit $X = \operatorname{Spec}(B)$ ; on peut d'autre part supposer que $L$ est isomorphe à $\mathcal{O}_X$, soit $L = \tilde{L}$, où $L$ est un $B$-module libre de rang 1. Soit $e$ un générateur de $L$, de sorte que $e^n$ est un générateur de $L^{\otimes n}$ ; on peut alors écrire pour tout $n$, $\varphi_n(s \otimes 1) = v_n(s) e^n$, où $v_n$ est un homomorphisme $S_n \to B$ de $A$-modules ; on outre on a par hypothèse, pour $s \in S_n, s' \in S_m : v_{n+m}(ss') = v_n(s)v_m(s')$.

Soit $f \in S_d$ ($d > 0$) ; nous allons définir un homomorphisme $\omega$ de $A$-algèbres $S_{(f)} \to B_f$ ; il suffit, pour $s \in S_{nd}$, de poser $\omega(s/f^n) = v_{nd}(s)/v_d(f)^n$. En outre, cet homomorphisme et l'anneau $B_f$ sont indépendants du générateur $e$ choisi, car tout autre générateur de

\marginpar{ $\varphi_n(s \otimes 1)$ où $1$ est la section de $q^*(S)$ correspondant à $s$}

\newpage

%% ===== Page 297 =====

-- 105 --
$L$ est de la forme $\varepsilon \sigma$, où $\varepsilon$ est inversible dans $B$ ; les homomorphismes $v_n'$ correspondants sont donc tels que $v_n'(s) = \varepsilon^n v_n(s)$ et on a $g' = \varepsilon^d g$, notre assertion en résulte. Il correspond donc à $\omega$ un morphisme $r_\varphi : D_+(f) \to D_+(f')$ et on vérifie aussitôt (compte tenu de l'identification canonique de $\operatorname{Spec}(S_{(f)})$ et de $D_+(f)$, cf. §3, n°2, prop. 4) que pour tout $x' \in D(g)$, on a $r_\varphi(x') = \psi(x')$; il est clair par ailleurs que $r_\varphi$ ne dépend pas du point $x \in D(g)$ initialement considéré. Enfin, si $f' \in S_{d'}$ ($d' > 0$) est un second élément homogène, et $g' = v_{d'}(f')$, on vérifie aussitôt que $r_{f'f}$ est la restriction de $r_f$ et de $r_{f'}$, à $D(gg')$. Cela montre l'existence du morphisme $r$ lorsque $X$ et $Y$ sont affines ; on passe de là aisément d'abord au cas où $Y$ est affine, $X$ quelconque, puis au cas général en tenant compte de la définition de $\operatorname{Proj}(S)$ (n°1).

Proposition 32. -- Soient $Y$ un préschéma localement noethérien, $X$ un $Y$-préschéma de type fini sur $Y$, $q : X \to Y$ le morphisme structural. Soient $\mathcal{S}$ une $\mathcal{O}_Y$-algèbre graduée, $\mathcal{L}$ un $\mathcal{O}_X$-module inversible, $\varphi : q^*(\mathcal{S}) \to \bigoplus_{n \geqslant 0} \mathcal{L}^{\otimes n}$, surjectif dans les degrés assez grands ; soit $r : X \to \operatorname{Proj}(\mathcal{S})$ le morphisme correspondant. Les conditions suivantes sont équivalentes :

(i) $r$ est un homéomorphisme de $X$ sur le sous-espace $r(X)$ de $\mathcal{P}$.

(ii) Pour tout couple de points $x, y$ de $X$ tels que $x \notin \{y\}$ et $q(x) = q(y)$, il existe un voisinage $U$ de $q(x)$ dans $Y$, un entier $n > 0$ et une section $s$ de $\mathcal{S}_n$ au-dessus de $U$, telle que si $s'$ est la section correspondante de $q^*(\mathcal{S}_n)$ au-dessus de $q^{-1}(U)$, l'image de $s'$ par $\varphi$ est une section de $\mathcal{L}^{\otimes n}$ nulle en $y$ et non nulle en $x$.

Il est immédiat que (i) entraîne (ii), car on a dans les hypothèses de (ii), $r(x) \notin \{r(y)\}$, et comme on peut supposer $Y = \operatorname{Spec}(A)$ affine, avec $\mathcal{S} = \tilde{S}$, il existe $f \in S_d$ (pour un $d > 0$) tel que $r(x) \in D_+(f)$ et $r(y) \notin D_+(f)$, et la section $s = f$ répond à la question.

Pour prouver que (ii) entraîne (i), il suffit de montrer que, pour...

\newpage

%% ===== Page 298 =====

- II-105 bis -

tout $x \in X$ et tout voisinage ouvert $V$ de $x$, il existe un voisinage ouvert $V'$ de $r(x)$ dans $P$ tel que $r^{-1}(V')$ soit contenu dans $V$. La question étant locale sur $Y$, on peut supposer $Y$ affine ; donc noethérien, et $X$ est alors aussi noethérien ; l'ensemble fermé $X-V$ est donc réunion de ses composantes en nombre fini, qui sont de la forme $\overline{\{y\}}$ ($y \in X$) ; on est par suite ramené au cas où $V=X-\overline{\{y\}}$ avec $x \notin \overline{\{y\}}$. Si $q(x) \notin \overline{\{q(y)\}}$, il suffit de prendre pour $V'$ l'image réciproque dans $P$ de $Y-\overline{\{q(y)\}}$. Sinon, on peut appliquer la condition (ii) et prendre pour $V'$ l'ensemble ouvert $P_s$ des $z$ où $s$ ne s'annule pas ($n^{\circ}1$, prop.2) ; il résulte de la définition de $r$ que $r^{-1}(V')$ ne contient pas $y$.

Corollaire 1. - Supposons que la $\mathcal{O}_Y$-algèbre $S$ soit de type fini. Pour que $r_{L, \varphi}$ soit une immersion, il faut et il suffit que les conditions équivalentes de la prop. 32 soient vérifiées, et en outre que l'on ait la propriété suivante : pour tout $x \in X$ et tout $\alpha \in \mathcal{O}_x$, il existe un entier $n > 0$ et deux sections $u, v$ de $S_n$ au-dessus d'un voisinage de $q(x)$ telles que, si $u', v'$ sont les images par $\varphi$ des sections correspondantes de $q^*(S_n)$ au-dessus de $q^{-1}(U)$, on ait $v'(x) \neq 0$, et que $\alpha$ soit le germe de $u'/v'$ au point $x$.

Notons que $P = \operatorname{Proj}(S)$ est alors de type fini sur $Y$, donc $r$ est un $Y$-morphisme de type fini (I, 4, 2, cor. 2 de la prop. 9) ; comme on suppose que $r$ est un homéomorphisme de $X$ sur $r(X)$, pour prouver que c'est une immersion, il suffit de montrer que c'est une immersion locale, et par suite de montrer que pour tout $x \in X$, $\varphi_x$ est un épimorphisme de $\mathcal{O}_{r(x)}$ sur $\mathcal{O}_x$ (I, 4, 3, prop. 16 (i)). On se ramène donc aussitôt au cas où $X$ et $Y$ sont affines, et où $\alpha$ est le germe d'une section définie au-dessus de $X$ tout entier ; la définition de $r$ donnée ci-dessus prouve alors que $\alpha$ est bien l'image par $\varphi_x$ d'un

\newpage

%% ===== Page 299 =====

\hfill \text{--- II-105 ter ---}

du germe de section $(u/v)_{\psi(x)}$ lorsque la condition de l'énoncé est remplie, et la réciproque est immédiate.

Corollaire 2. --- Supposons que $S$ soit de type fini et que $S$ soit de type fini et que pour tout ouvert affine $U \subset Y$, il existe un entier $m > 0$ tel que, pour tout entier $n > 0$ et toute section $t$ de $\mathcal{L}^{\otimes n}$ au-dessus de $q^{-1}(U)$, il existe une section $s$ de $S_m$ au-dessus de $U$ telle que $t = f_m(s')$, où $s'$ est la section correspondante de $q^*(S_m)$ au-dessus de $q^{-1}(U)$. Alors si $r$ est un homéomorphisme de $X$ sur $r(X)$, $r$ est une immersion de $X$ dans $P$.

Il faut vérifier la condition du cor. 1. Prenons pour $U$ un ouvert affine contenant $q(x)$; $\alpha$ est le germe d'une section de $\mathcal{O}_X$ au-dessus d'un voisinage ouvert $V \subset q^{-1}(U)$ de $x$, et en vertu de la prop. 32, on peut supposer qu'il existe un voisinage ouvert affine $W \subset U$ de $q(x)$ et une section $v$ de $S_m$ au-dessus de $W$ tels que, si $v'$ est la section correspondante de $\mathcal{L}^{\otimes m}$ au-dessus de $q^{-1}(W)$, $V$ soit l'ensemble des points de $X$ où $v'$ ne s'annule pas. Il existe alors un entier $h > 0$ tel que $v'^h$ se prolonge en une section de $\mathcal{L}^{\otimes mh}$ au-dessus de $q^{-1}(W)$ (correspondant à l'ouvert affine $W$) (§ 1, n$^\circ$4, th. 2), on peut évidemment supposer $h$ de la forme $nm$, et la section prolongée est donc de la forme $u'$, correspondant à une section $u$ de $S_{mh}$ au-dessus de $W$; comme $\alpha = (u'/v'^h)_x$, la condition du cor. 1 est satisfaite.

Nous utiliserons surtout les critères précédents lorsque $S = S_{\mathcal{O}_Y}(\mathcal{E})$, $\mathcal{E}$ étant un $\mathcal{O}_Y$-Module quasi-cohérent; tout épimorphisme $\varphi : q^*(\mathcal{E}) \to \mathcal{L}$ définit alors un épimorphisme (encore noté $\varphi$) $S \to S' = \bigoplus_{n \geq 0} \mathcal{L}^{\otimes n}$, et le morphisme $r_{\varphi}$ est alors celui qui a été noté de cette façon au n$^\circ$5.

\newpage

%% ===== Page 300 =====

- II-105 \emph{ter} -

du genre de section $(u/v)\psi(x)$ lorsque la condition de l'énoncé est remplie, et la réciproque est immédiate.

Corollaire 2. - Supposons que $S$ soit de type fini et que, pour tout ouvert affine $U \subset S$, il existe un entier $m > 0$ tel que, pour tout entier $n > 0$ et toute section $t$ de $\mathbb{L}^{\otimes nm}$ au-dessus de $q^{-1}(U)$, il existe une section $s$ de $\mathbb{S}_{nm}$ au-dessus de $U$ telle que $t = f_{nm}(s)$, où $s'$ est la section correspondante de $\mathbb{Q}^{\ast}(\mathbb{S}_{nm})$ au-dessus de $q^{-1}(U)$. Alors si $r$ est un homéomorphisme de $X$ sur $r(X)$, $r$ est une immersion de $X$ dans $P$.

Il faut vérifier la condition du cor. 1. Prenons pour $U$ un ouvert affine contenant $q(x)$; $x$ est le germe d'une section $\alpha$ de $\mathcal{O}_X$ au-dessus d'un voisinage ouvert $V \subset q^{-1}(U)$ de $x$, et en vertu de la prop. 32, on peut supposer qu'il existe un voisinage ouvert affine $W \subset U$ de $q(x)$ et une section $v$ de $\mathbb{S}_k$ au-dessus de $W$ tels que, si $v'$ est la section correspondante de $\mathbb{L}^{\otimes k}$ au-dessus de $q^{-1}(W)$. $V$ soit l'ensemble des points de $X$ où $v'$ ne s'annule pas. Il existe alors un entier $h > 0$ tel que $v'^h$ se prolonge en une section de $\mathbb{L}^{\otimes kh}$ au-dessus de $q^{-1}(W)$ (en correspondant à l'ouvert affine $W$) (§ 1, n°4, th. 2). On peut évidemment supposer $h$ de la forme $mn$, et la section prolongée est donc de la forme $u'$, correspondant à une section $u$ de $\mathbb{S}_{kh}$ au-dessus de $W$ ; comme $\alpha = (u'/v'^h)_x$, la condition du cor. 1 est satisfaite.

Nous utiliserons surtout les critères précédents lorsque $S = \mathcal{S}_Y(E)$, $E$ étant un $\mathcal{O}_Y$-module quasi-cohérent ; tout épimorphisme $\varphi : q^*(E) \to \mathbb{L}$ définit alors un épimorphisme (encore noté $\varphi$) $S \to S' = \bigoplus_{n \geq 0} \mathbb{L}^{\otimes n}$, et le morphisme $\pi_{\varphi}$ est alors celui qui a été noté de cette façon au n°5.

\newpage

%% ===== Page 301 =====

- 11-105 quinquiès -

de tous les éléments intervenant de cette façon. Notons d'autre part que l'on peut supposer les $W_i$ choisis de façon que $\mathcal{L}|V_i$ soit de la forme $\tilde{L}_i$, où $L_i$ est un module monogène libre sur $A(V_i)$ ; comme $q^*(S_n)|V_i \to \mathcal{O}_V^n|V_i$ est surjectif par hypothèse, il existe un élément $h_i$ de $E_i$ dont l'image par $\varphi_n$ (qui est une section de $\mathcal{O}_V^n$ au-dessus de $q^{-1}(U_i)$) donne par restriction à $V_i$ un générateur de $\mathcal{O}_V^n$ (I, 3, cor. 2 et 4 du th. 1) ; soit $R_i' = R_i \cup \{h_i\}$. Le sous-faisceau $E_i'$ de $S_n|U_i$ engendré par les éléments de $R_i'$ est de type fini, donc cohérent puisque $U_i$ est noethérien ; il existe par suite un sous-faisceau cohérent $E'$ de $S_n$ induisant $E_i'$ sur $U_i$ (§ 1, n° 2, th. 2) ; si $E$ est le sous-faisceau de $S_n$ somme des $E_i'$, il est clair que $E$ est cohérent et que la restriction $\varphi'$ de $\varphi_n$ à $q^*(E)$ est un épimorphisme sur $\mathcal{O}_V^n$. Considérons alors dans $\mathbb{P}(E)$ les ouverts affines $W_i'$ définis par les sections $f_i$ de $E$ au-dessus des $U_i$ ; le choix des $E_i'$ et la description du morphisme $r' : r^{-1}(\mathcal{O}_V^n) \to \mathcal{O}_V^n$ montrent que la restriction de $r'$ à $V_i$ correspond à un homomorphisme surjectif $A(W_i') \to A(V_i)$. On en conclut que cette restriction est une immersion fermée ; pour voir que $r'$ est une immersion, on peut se limiter au cas où $Y$ est affine, la question étant locale sur $Y$ ; mais alors, comme $W_i$ et $W_i'$ sont définis par la même section $f_i$, il résulte de la définition de $r'$ et de la prop. 2 du n° 1 que $r'^{-1}(W_i') = V_i$, ce qui achève la démonstration.

\newpage

%% ===== Page 302 =====

- 11-106 -

\sout{est cohérent et que la restriction de $\varphi'$ à $q^*(\mathbb{E}')$ est un épimorphisme}
\sout{Cela étant, si on considère dans $\mathbb{P}(\mathbb{E}')$ les ouverts affines $W'_i$ définis}
\sout{par les sections $f_i$ de $\mathbb{E}'$ au-dessus des $U_i$, il est clair que $A(W'_i)$}
\sout{s'identifie à $A(W_i)$ et la description du morphisme $r' = r_{L', \varphi'}$ montre}
\sout{que la restriction de $r'$ à $V_i$ est un morphisme de $V_i$ dans $W'_i$ qui}
\sout{correspond au même homomorphisme $A(W'_i) \to A(V_i)$ que la restric-}
\sout{tion de $r$ à $V_i$. On en conclut que la restriction de $r'$ à $V_i$}
\sout{est une immersion fermée, donc que $r'$ est encore une immersion de}
\sout{$X$ dans $\mathbb{P}(\mathbb{E}')$.}

Les résultats précédents nous permettent de caractériser les faisceaux amples et très amples par des propriétés intrinsèques.

Proposition 34. - Soient $X$ un préschéma noethérien, $q : X \to Y$ un morphisme de type fini, $L$ un faisceau inversible sur $X$. Les conditions suivantes sont équivalentes :

(i) $L$ est très ample.

(ii) L'homomorphisme canonique $\mu : q^*(q_*(L)) \to L$ est surjectif (autrement dit $L$ est engendré par ses sections au-dessus des $q^{-1}(U)$, où $U$ parcourt les ouverts affines de $Y$), et le morphisme correspondant $X \to \mathbb{P}(q_*(L))$ est une immersion.

(iii) Il existe un sous-faisceau cohérent $\mathbb{E}$ de $q_*(L)$ tel que la restriction à $q^*(\mathbb{E})$ de l'homomorphisme canonique $q^*(q_*(L)) \to L$ soit surjective et que le morphisme correspondant $X \to \mathbb{P}(\mathbb{E})$ soit une immersion.

La prop. 33 montre que (ii) entraîne (iii), $q_*(L)$ étant quasi-cohérent (§ 1, $n^\circ 1$, cor. de la prop. 1) ; d'autre part, la prop. 32 et son cor. 1 prouvent que (iii) entraîne (ii), car l'injection $\mathbb{E} \to q_*(L)$ donne canoniquement un homomorphisme $S_{\mathcal{O}_X}(q^*(\mathbb{E})) \to S_{\mathcal{O}_X}(q^*(q_*(L)))$, qui, composé avec l'homomorphisme $S_{\mathcal{O}_X}(q^*(M)) \to \bigoplus_{n \geq 0} L^{\otimes n}$, donne l'homomor...

\marginpar{P9:}

\newpage

%% ===== Page 303 =====

$\varphi' :$ phisme $S_{\mathcal{O}_X}(q^*(E)) \to \bigoplus_{n \ge 0} L^{\otimes n}$ qui définit l'immersion $X \to \mathbb{P}(E)$ ; comme les critères de la prop. 32 et de son cor. 1 reposent sur l'existence d'images de sections par $\varphi$ ayant des propriétés données, il suffira de prendre les images de sections par $\varphi'$ ayant ces propriétés pour répondre à la question. Il est immédiat que (iii) entraîne (i) puisqu'on a vérifié au [th. 2] (chap. 0, $\S$ 2) que le morphisme \text{=P} $\mathbb{P}(E) \to \mathbb{P}(E)$ l'image de $X \to \mathbb{P}(E)$ du faisceau $\mathcal{O}_{\mathbb{P}}(1)$ est isomorphe au faisceau inversible $L$ qui a servi à définir ce morphisme. Enfin (i) entraîne (iii) : supposons en effet que $L = f^*(\mathcal{O}_{\mathbb{P}}(1))$, où $f$ est une immersion $X \to \mathbb{P}(F)$, $F$ étant un $\mathcal{O}_Y$-Module cohérent ; on sait ($n^\circ 5$) qu'on en déduit canoniquement un homomorphisme $\theta : q^*(F) \to L$, et par suite aussi un homomorphisme $\theta' : F \to q_*(L)$. Soit $E$ l'image de $F$ par $\theta$, et comme $E$ est un faisceau quasi-cohérent de type fini, donc cohérent puisque $Y$ est noethérien, et que $\theta$ se factorise en $F \to E \to q_*(L)$, et par suite $\theta$ se factorise en $q^*(F) \to q^*(E) \to L$ ; en outre, comme $q^*$ est exact à droite, $q^*(\theta)$ est surjectif, donc il lui correspond une immersion fermée $\mathbb{P}(E) \to \mathbb{P}(F)$ ($n^\circ 4$, prop. 15). A la factorisation précédente de $\theta$ correspond ($n^\circ 5$, prop. 17) une factorisation $X \to \mathbb{P}(E) \to \mathbb{P}(F)$ et comme $X \to \mathbb{P}(F)$ est une immersion ainsi que $\mathbb{P}(E) \to \mathbb{P}(F)$, il en est de même de $X \to \mathbb{P}(E)$.

Corollaire. — Pour que $L$ soit très ample (resp. ample), il faut et il suffit qu'il existe un recouvrement ouvert $(U_\alpha)$ de $Y$ tel que chacun des $L|_{q^{-1}(U_\alpha)}$ soit très ample (resp. ample).

C'est immédiat pour les faisceaux très amples, la condition (ii) de la prop. 34 étant locale en $Y$. D'autre part, comme on peut supposer le recouvrement $(U_\alpha)$ fini ($Y$ étant quasi-compact), l'assertion relative aux faisceaux amples se déduit de celle relative aux faisceaux très amples par la définition des faisceaux amples.

Proposition 35. — Les hypothèses et notations étant les mêmes que dans la prop. 34, les conditions suivantes sont équivalentes :

\newpage

%% ===== Page 304 =====

- II-106 ter -

(i) $\underline{L}$ est ample .
(ii) Pour tout couple de points $x, y$ de $X$ tels que $x \notin \overline{\{y\}}$ et $q(x) \in \overline{\{q(y)\}}$, il existe un entier $n > 0$, un voisinage ouvert $U$ de $q(x)$ et une section $s$ de $\underline{L}^{\otimes n}$ au-dessus de $q^{-1}(U)$, nulle en $y$ et non nulle en $x$.
(iii) Pour tout $\mathcal{O}_X$-module cohérent $\underline{F}$, l'homomorphisme canonique $\mu : q^*(q_*(\underline{F} \otimes \underline{L}^{\otimes n})) \to \underline{F} \otimes \underline{L}^{\otimes n}$ est surjectif dès que $n$ est assez grand .
(iv) Pour tout faisceau cohérent d'idéaux $\underline{J}$ de $\mathcal{O}_X$, l'homomorphisme canonique $\mu : q^*(q_*(\underline{J} \otimes \underline{L}^{\otimes n})) \to \underline{J} \otimes \underline{L}^{\otimes n}$ est surjectif dès que $n$ est assez grand .

(i) implique (iii). En effet : il suffit de le démontrer lorsque $\underline{L}$ est très ample, donc de la forme $i^*(\mathcal{O}_P(1))$, où $i : X \to P(\underline{E})$ est une immersion et $\underline{E}$ un $\mathcal{O}_X$-module cohérent. Alors $X$ s'identifie à un sous-préschéma induit sur un ouvert de $\overline{X}$, plus petit sous-préschéma fermé de $P$ majorant $X$ (§ 1, n°2, prop. 8), et le faisceau cohérent $\underline{F}$ sur $X$ est induit par un faisceau cohérent $\underline{\overline{F}}$ sur $\overline{X}$ (§ 1, n°2, th. 1). On est donc ramené à démontrer notre assertion lorsque $X$ est un sous-préschéma fermé de $P$, et en remplaçant $\underline{F}$ par $i_*(\underline{F})$, on est même ramené au cas où $X=P$. Mais alors le résultat a déjà été démontré (n°3, cor. 3 du th. 1).

Il est clair que (iii) entraîne (iv). Prouvons que (iv) entraîne (ii). Prenons en effet pour $\underline{J}$ le faisceau cohérent d'idéaux définissant la partie fermée $\overline{\{y\}}$. Soit $t$ une section de $\underline{J}$ au-dessus d'un voisinage ouvert $V$ de $x$, ne s'annulant pas en $x$ ; en la multipliant par une section de $\underline{L}^{\otimes n}$ au-dessus d'un voisinage de $x$, égale à un générateur de $\underline{L}^{\otimes n}$ au point $x$, on obtient une section $t'$ de $\underline{J} \otimes \underline{L}^{\otimes n}$ au-dessus d'un voisinage de $x$, ne s'annulant pas en $x$. Par hypothèse, il existe un voisinage ouvert $U$ de $x$ et une section $s$ de $\underline{J} \otimes \underline{L}^{\otimes n}$ au-

\newpage

%% ===== Page 305 =====

- 11-106 quater -

dessus de $q^{-1}(U)$ dont la restriction à un voisinage de $x$ est égale à $t'$ ; l'image de cette section dans $\underline{L}^{\otimes n}$ répond à la question .

Montrons enfin que (ii) entraîne (i) . Considérons l' $\mathcal{O}_Y$-algèbre graduée quasi-cohérente $S = \bigoplus_{n \geq 0} q_*(\underline{L}^{\otimes n})$ ; comme une section de $\underline{L}^{\otimes n}$ qui est non nulle en un point $x$ engendre $(\underline{L}^{\otimes n})_x$, la condition (ii) entraîne en particulier que l'homomorphisme canonique $q^*(q_*(\underline{L}^{\otimes n})) \to \underline{L}^{\otimes n}$ est surjectif pour un $n=d$ ; il en est alors de même de $q^*(q_*(\underline{L}^{\otimes hd})) \to \underline{L}^{\otimes hd}$ pour tout $h > 0$ ; donc en remplaçant $\underline{L}$ par $\underline{L}^{\otimes d}$, on peut supposer que l'homomorphisme canonique $q^*(S) \to \bigoplus_{n \geq 0} \underline{L}^{\otimes n}$ est surjectif, et définit donc un morphisme de $X$ dans $\operatorname{Proj}(S)$. La prop. 32 montre que ce morphisme est un homomorphisme de $X$ sur son image $r(X)$ ; en outre, la définition de $S$ montre que la condition de la cor. 2 de la prop. 32 est trivialement satisfaite, donc $r$ est une immersion. Enfin, il résulte de la prop. 33 la conclusion.

Corollaire 1. — Soient $\underline{F}$ un $\mathcal{O}_Y$-module quasi-cohérent, $r: X \to P(\underline{F})$ un morphisme qui soit un homomorphisme de $X$ sur $r(X)$ ; alors $\underline{L} = r^*(\mathcal{O}_P(1))$ est ample.

En effet, la condition (ii) de la prop. 34 est remplie comme conséquence de la condition (ii) de la prop. 32.

Corollaire 2. — Pour que $\underline{L}$ soit ample il faut et il suffit que son image réciproque sur $X_{\operatorname{red}}$ le soit.

En effet, $X_{\operatorname{red}}$ est un sous-préschéma fermé de $X$ et $\underline{L}'$ vérifie les conditions du cor. 1.

Remarques. — 1) Nous verrons au chap. III, § 2, n° 1 une généralisation du cor. 2 : si $X$ est propre sur $Y$, $X'$ fini sur $X$ et dominant $X$, pour que $\underline{L}$ soit ample il faut et il suffit que son image réciproque sur $X'$ soit ample.

\newpage

%% ===== Page 306 =====

\marginpar{Archives Grothendieck - sept 59}

- II-106 -

dessus de $q^{-1}(U)$ dont la restriction à un voisinage de $x$ est égale à $t'$ ; l'image de cette section dans $L^{\otimes n}$ répond à la question.

Montrons enfin que (ii) entraîne (i). Considérons l' $\mathcal{O}_Y$-algèbre graduée quasi-cohérente $S = \bigoplus_{n \ge 0} q_* (L^{\otimes n})$ ; comme une section de $L^{\otimes n}$ qui est non nulle en un point $x'$ engendre $(L^{\otimes n})_x$, la condition (ii) entraîne en particulier que l'homomorphisme canonique 
$q^* (q_* (L^{\otimes n})) \to L^{\otimes n}$ est surjectif pour un $n=d$ ; il en est alors de même de $q^* (q_* (L^{\otimes hd})) \to L^{\otimes hd}$ pour tout $h > 0$ ; donc en remplaçant $L$ par $L^d$ on peut supposer que l'homomorphisme canonique $q^* (\bigoplus_{n \ge 0} L^{\otimes n}) \to \bigoplus_{n \ge 0} L^{\otimes n}$ est surjectif, et définit donc un morphisme de $X$ dans $\operatorname{Proj}(S)$. La prop. 32 montre que ce morphisme est un homomorphisme de $X$ sur son image $r(X)$ ; en outre, la définition de $S$ montre que la condition du cor. 2 de la prop. 32 est trivialement satisfaite, donc est une immersion. Enfin, il résulte de la prop. 33 la conclusion.

Corollaire 1. - Soient $L$ un $\mathcal{O}_X$-module quasi-cohérent, $r$ un morphisme $X \to P(L)$ qui soit un homomorphisme de $X$ sur $r(X)$ ; alors $L=r^*(\mathcal{O}_P(1))$ est ample.

En effet, la condition (ii) de la prop. 34 est remplie, comme conséquence de la condition (ii) de la prop. 32.

Corollaire 2. - Pour que $L$ soit ample, il faut et il suffit que son image réciproque sur $X_{\text{red}}$ le soit.

En effet, $X_{\text{red}}$ est un sous-préschéma fermé de $X$ et $L'$ vérifie les conditions du cor. 1.

Remarques. - 1) Nous verrons au chap. III, §2, n° 1 une généralisation du cor. 2 : si $X$ est propre sur $Y$, $X'$ fini sur $Y$ et dominant $X$, pour que $L$ soit ample il faut et il suffit que son image réciproque sur $X'$ soit ample.

\newpage

%% ===== Page 307 =====

- II-107 bis -

(v) Si $f : X \to Y$ et $g : Y \to Z$ sont deux morphismes tels que $g \circ f$ soit propre et $g$ séparé, $f$ est propre.

(vi) Si $f : X \to Y$ est propre, il en est de même de $f_{\text{red}}$.

Comme d'ordinaire, on n'a à démontrer que les propriétés (i), (ii) et (iii), les autres en découlant formellement. Dans chacun de ces trois cas, la vérification de la condition (a) de la déf. 1 découle de résultats antérieurs (I, 3, 6, prop. 16, 17 et 18 et I, 4, 2, prop. 7, 8 et 9) ; reste à vérifier la condition (b). C'est immédiat dans le cas (i), car si $X \to Y$ est une immersion fermée, il en est de même de $X \times_Y Y' \to Y \times_Y Y' = Y'$ (I, 3, 2, prop. 5). Pour démontrer (ii), considérons un morphisme $Z' \to Z$, où $Z'$ est localement noethérien ; comme $Y \to Z$ est de type fini, il en est de même de $Y \times_Z Z' \to Z' = Y \times_Z Z'$ (I, 4, 2, prop. 9), donc $Y \times_Z Z'$ est localement noethérien (I, 4, 2, prop. 10). Or, on peut écrire $X \times_Z Z' = X \times_Y (Y \times_Z Z')$ (I, 2, 5, formule (1)), et par suite la projection $X \times_Z Z' \to Z'$ se factorise en $X \times_Y (Y \times_Z Z') \to Y \times_Z Z' \to Z'$. Compte tenu de la remarque précédente, (ii) résulte donc du fait que le composé de deux applications fermées est fermée. Enfin, pour toute extension $S' \to S$, $X^{S'}$ s'identifie à $X \times_Y Y^{S'}$, donc pour tout morphisme $Z \to Y^{S'}$, où $Z$ est localement noethérien, on peut écrire $X^{S'} \times_{Y^{S'}} Z = (X \times_Y Y^{S'}) \times_{Y^{S'}} Z = X \times_Y Z$, et comme $X \times_Y Z$ est fermée par hypothèse, cela prouve (iii).

Corollaire 1. - Soient $f : X \to Y$, $g : Y \to Z$ deux morphismes. Si $g \circ f$ est propre, $f$ surjectif, $f$ séparé et de type fini, alors $g$ est propre.

Il n'y a à vérifier que la condition (b) de la déf. 1. Or pour tout morphisme $Z' \to Z$ tel que $Z'$ soit localement noethérien, le diagramme
$$\begin{tikzcd}
X \times_Z Z' \arrow[r, "f \times 1"] \arrow[d] & Y \times_Z Z' \arrow[d, "p'"] \\
X \times_Z Z' \arrow[r, "p"] & Z'
\end{tikzcd}$$
(où $p$ et $p'$ sont les projections) est commutatif ; en outre $f \times 1$

\newpage

%% ===== Page 308 =====

- II-108 -

est surjectif puisqu'il en est ainsi de $f$ (I, 2.7, prop. 11), et $p'$ est une application fermée. Toute partie fermée $F$ de $Y \times_Z Z'$ est alors l'image par $f \times 1$ d'une partie fermée $E$ de $X \times_Z Z'$, donc $p(F) = p'(E)$ est fermé dans $Z'$, d'où le corollaire.

Corollaire 2. — Soit $f: X \to Y$ un morphisme séparé de type fini. Supposons que $X$ soit réunion finie d'ensembles fermés $X_i$ ($1 \le i \le n$) et désignons par $X_i$ un sous-préschéma fermé de $X$ ayant $X_i$ pour espace de base ; soit $j_i$ l'injection canonique $X_i \to X$. Supposons d'autre part que pour chaque $i$, il existe un sous-préschéma fermé $Y_i$ de $Y$ tel que l'on ait $f \circ j_i = h_i \circ f_i$, où $f_i$ est un morphisme $X_i \to Y_i$ et $h_i$ l'injection canonique $Y_i \to Y$. Alors, pour que $f$ soit propre, il faut et il suffit que chacun des $f_i$ le soit.

En effet, si $f$ est propre, il en est de même de $f \circ j_i$ puisque $j_i$ est une immersion fermée (prop. 2 (i) et (ii)), donc, comme $h_i$ est une immersion fermée, et par suite est séparé, $f_i$ est propre (prop. 2 (v)). Inversement, supposons les $f_i$ propres et considérons le préschéma $Z = \sum X_i$, et soit $u$ le morphisme $Z \to X$ qui se réduit à $j_i$ dans chaque $X_i$. La restriction de $f \circ u$ à chaque $X_i$ est égale à $f \circ j_i = h_i \circ f_i$, donc est propre puisqu'il en est ainsi de $h_i$ et de $f_i$ (prop. 2 (i) et (ii)) ; il résulte aussitôt de la déf. 1 que $u$ est alors propre. Mais comme $u$ est surjectif par hypothèse, et $f$ séparé et de type fini, on conclut que $f$ est propre par application du cor. 1.

En particulier :

Corollaire 3. — Si $f: X \to Y$ est séparé et de type fini, et si $f_{red}$ est propre, alors $f$ est propre.

Il suffit en effet d'appliquer le cor. 2 avec $n=1$, $X_1 = X_{red}$, $Y_1 = Y_{red}$.

Si $X$ et $Y$ sont noethériens et $f: X \to Y$ un morphisme séparé de type fini, les cor. 2 et 3 permettent de se ramener au cas de [texte coupé].

\newpage

%% ===== Page 309 =====

- II-109 -

\text{préschémas intègres pour vérifier que } f \text{ est propre. En effet, on peut alors supposer } X \text{ et } Y \text{ réduits par le cor.3 ; soient } X_i \text{ } (1 \le i \le n) \text{ les composantes irréductibles de } X, Y_j \text{ } (1 \le j \le m) \text{ celles de } Y \text{ ; désignons encore par } X_i \text{ (resp. } Y_j \text{) le sous-préschéma fermé réduit ayant } X_i \text{ (resp. } Y_j \text{) pour espace de base (I,3,4,prop.9) ; si } x_i \text{ est le point générique de } X_i \text{, et } y_j \text{ une des composantes irréductibles de } Y \text{ contenant } f(x_i) \text{, } f(x_i) \text{ est contenu dans l'adhérence de } \{f(x_i)\} \text{, donc dans } Y_j \text{ . Comme le sous-espace } f^{-1}(Y_j) \text{ contient } X_i \text{, le sous-préschéma } f^{-1}(Y_j) \text{ de } X \text{ majore le sous-préschéma } X_i \text{, et par suite (I,3,2) } f \circ j_i \text{ } (j_i \text{ injection } X_i \to X) \text{ se factorise en } h_j \circ f_i \text{ où } h_j \text{ est l'injection } Y_j \to Y \text{ et } f_i \text{ un morphisme } X_i \to Y_j \text{ ; on est donc dans les conditions d'application du cor.2 et il suffit que les } f_i \text{ soient propres pour que } f \text{ le soit.}

\text{Corollaire 4. - Soient } X, Y \text{ des S-préschémas séparés de type fini sur } S \text{, et soit } f: X \to Y \text{ un S-morphisme. Pour que } f \text{ soit propre, il faut et il suffit que pour tout S-préschéma } S' \text{ localement noethérien, le morphisme } f \times 1: X \times_S S' \to Y \times_S S' \text{ soit une application fermée.}

\text{Notons d'abord que si } \varphi: X \to S \text{ et } \psi: Y \to S \text{ sont les morphismes structuraux, on a par définition } \psi \circ f = \varphi \text{, donc } f \text{ est séparé et de type fini (I,3,6,cor.4 de la prop.20 et I,4,2,cor.2 de la prop.9). Si maintenant } Y' \text{ est un } Y\text{-préschéma localement noethérien, } Y' \text{ peut aussi être considéré comme un } S\text{-préschéma, et comme } Y \to S \text{ est séparé (I,3,6,lemme 1), } X \times_Y Y' \text{ s'identifie à un sous-préschéma fermé de } X \times_S Y' \text{. Cela étant, si } f \text{ est propre, il en est de même de } f \times 1: X \times_S S' \to Y \times_S S' \text{ (prop.2 (ii)), et comme } Y \times_S S' \text{ est de type fini sur } S' \text{, donc localement noethérien, } f \times 1 \text{ est une application fermée (prop.1). Inversement, d'après la remarque faite ci-dessus, il suffit de montrer que } p: X \times_Y Y' \to Y' \text{ est une application fermée. Or, si } h \text{ est le morphisme } Y' \to Y \text{, le morphisme graphe } \Gamma_h \text{ est une immersion fermée de } Y' \text{ dans}

\newpage

%% ===== Page 310 =====

- II-110 -

est commutatif ; comme $f \times 1$ est une application fermée, il est clair qu'il en est de même de $p$.

\emph{Proposition 3. — Tout morphisme projectif est propre.}

On sait en effet qu'un tel morphisme $X \to Y$ est séparé et de type fini ($\S$ 4, $n^{\circ}6$, prop. 22), et si $Y'$ est un $Y$-préschéma localement noethérien, $X \times_Y Y' \to Y'$ est projectif ($\S$ 4, $n^{\circ}6$, th. 3 (iv)), et comme $Y'$ est localement noethérien, le morphisme précédent est une application fermée ($\S$ 4, $n^{\circ}6$, th. 4).

\emph{Corollaire 1. — Soit $Y$ un préschéma localement noethérien. Pour qu'un morphisme $f : X \to Y$ soit projectif, il faut et il suffit qu'il soit quasi-projectif et propre.}

La condition est nécessaire en vertu de la prop. 3. Inversement, supposons-la remplie ; il existe alors un $Y$-préschéma $P$ projectif sur $Y$ et une immersion $g : X \to P$ telle que $f = h \circ g$, où $h$ est le morphisme $P \to Y$ ; comme $h$ est séparé, on en conclut que $g$ est propre (prop. 2 (v)), et comme $P$ est de type fini sur $Y$, donc localement noethérien, $g$ est une immersion fermée (prop. 1), donc $f$ est projectif.

\emph{Corollaire 2. — Tout morphisme fini est propre.}

Il est en effet projectif ($\S$ 4, $n^{\circ}7$, prop. 28).

2. \emph{Le lemme de Chow.}

\emph{Théorème 1 (lemme de Chow). — Soit $S$ un préschéma noethérien, $X$ un $S$-préschéma de type fini et séparé sur $S$.}

(1) \emph{Il existe un $S$-préschéma $X'$ quasi-projectif sur $S$ et un S-morphisme $f : X' \to X$, projectif et surjectif.}

\newpage

%% ===== Page 311 =====

- II-111 -

(ii) On peut prendre $X'$ et $f$ tels qu'il existe un ouvert $U \subset X$ tel que $U' = f^{-1}(U)$ soit dense dans $X'$ et que la restriction de $f$ à $U'$ soit un isomorphisme $U' \xrightarrow{\sim} U$.

(iii) Si $X$ est réduit (resp. irréductible, intègre), on peut supposer $X'$ réduit (resp. irréductible, intègre).

La démonstration se fait en plusieurs étapes.

A) On peut d'abord supposer $X$ irréductible ; en effet, comme $X$ est noethérien, ses composantes irréductibles $X_i$ ($1 \leq i \leq n$) sont en nombre fini. Si le théorème est démontré pour chacune d'elles, et si $X'_i$ et $f_i : X'_i \to X_i$ sont le préschéma et le morphisme correspondant à $X_i$, le préschéma $X'$ somme des $X'_i$ et le morphisme $f : X' \to X$ dont la restriction à $X'_i$ est $f_i$, répondront à la question. Il est immédiat en effet que la condition (iii) est vérifiée par $X'$ si elle l'est par les $X'_i$ ; il en est de même de (ii), en prenant pour $U$ la réunion des $U_i \cap X_i$ (non vides). Enfin, comme les $X'_i$ sont quasi-projectifs sur $S$, il en est de même de $X'$ (§ 4, n°7, Remarque) ; de même, les $X'_i \to X_i$ sont projectifs par le th.3 (i) et (ii) du § 4, n°6, donc $X' \to X$ est projectif (§ 4, n°6, Remarque) et est évidemment surjectif par définition.

B) $X$ est réunion d'une famille finie $(X_k)_{1 \leq k \leq n}$ d'ouverts affines non vides ; comme $S$ est un schéma, ces derniers sont affines sur $S$ (I, 3, 7, cor. 3 de la prop. 22). Il existe par suite pour chaque $k$ une immersion ouverte $\varphi_k : X_k \to P_k$, où $P_k$ est projectif sur $S$ (§ 4, n°7, prop. 24 et 27). Soit $U = \bigcap X_k$ ; comme $X$ est irréductible et les $X_k$ non vides, $U$ est non vide, et par suite dense dans $X$ ; les restrictions des $\varphi_k$ à $U$ définissent un morphisme $\varphi : U \to P$, où

\newpage

%% ===== Page 312 =====

- II-112 -

$$
\begin{tikzcd}
U \arrow[r, "\varphi"] \arrow[d, "j_k"] & P \arrow[d, "p_k"] \\
X_k \arrow[r, "\varphi_k"] & P_k
\end{tikzcd}
$$
(1)

soient commutatifs, $j_k$ désignant l'injection canonique $U \to X_k$ et $p_k$ la projection $P \to P_k$ (I, 2, 4). Si $j$ est l'injection canonique $U \to X$, le morphisme $\psi = (j, \varphi) : U \to X \times_S P$ est une immersion (I, 3, 5, cor. 4 de la prop. 15). Soit $X'$ le plus petit sous-préschéma fermé de $X \times_S P$ majorant l'immersion $\psi$ (§ 1, n°2, prop. 3) : $\psi$ peut donc se factoriser en
(2)
$$ \psi : U \xrightarrow{\psi'} X' \xrightarrow{h} X \times_S P $$
où $\psi'$ est une immersion ouverte et $h$ une immersion fermée. Si $q_1 : X \times_S P \to X$ et $q_2 : X \times_S P \to P$ sont les projections, poserons $f : X' \to X$ et $g : X' \to X \times_S P \to P$. Nous allons voir que $X'$ et $f$ répondent à la question.

C) Montrons d'abord que $f$ est projectif et surjectif et que la restriction de $f$ à $f^{-1}(U) = U'$ est un isomorphisme de $U'$ sur $U$. Comme les $P_k$ sont projectifs sur $S$, il en est de même de $P$ (§ 4, n°6, th. 3 (iv)), donc $X \times_S P$ est projectif sur $X$ (§ 4, n°6, th. 3 (iii)) et il en est de même par définition de $X'$ qui est un sous-préschéma fermé de $X \times_S P$. D'autre part, on a $f \circ \psi' = q_1 \circ (h \circ \psi') = q_1 \circ \psi = j$, donc $f(X')$ contient l'ensemble ouvert partout dense $U$ de $X$ ; mais $f$ est une application fermée de $X'$ dans $X$ puisque $f$ est projectif (§ 4, n°6, 4), donc $f(X') = X$. Enfin, remarquons que $U \times_S P$ s'identifie à un préschéma induit sur un ouvert de $X \times_S P$, l'injection canonique de $X'$ s'identifiant à $j \times 1$, et le sous-préschéma $X' \cap (U \times_S P)$ contient évidemment le préschéma induit sur $U$. Mais d'autre part, l'immersion $\psi$ se factorise en $U \xrightarrow{(\psi', j \times 1)} U \times_S P \xrightarrow{h'} X \times_S P$, et comme $X$ est séparé sur $S$, le morphisme graphe $\Gamma_f$ est une immersion fermée (I, 3, 7, prop. 22), donc $f(U)$ est un sous-préschéma fermé de $U \times_S P$ ; il est donc égal à $X' \cap (U \times_S P)$, et com-

\newpage

%% ===== Page 313 =====

- II-113 -

me par ailleurs $U'$ majore évidemment $\psi(U)$, on voit que l'on a $\psi(U) = U' \cap (U \times_S P)$, et comme $\psi$ est une immersion, la restriction de $f$ à $U'$ est bien un isomorphisme sur $U$ ; par ailleurs, en raison de la définition de $X'$, $U'$ est dense dans $X'$.

D) Prouvons maintenant que $g$ est une immersion, ce qui établira que $X'$ est quasi-projectif sur $S$, puisque $P$ est projectif sur $S$. Posons:
$$
\begin{aligned}
V_k &= p_k(X_k) \quad (\text{ouvert de } P_k) \\
W_k &= p_k^{-1}(V_k) \quad (\text{ouvert de } P) \\
X'_k &= f^{-1}(X_k) \\
X''_k &= g^{-1}(W_k)
\end{aligned} \quad (\text{ouverts de } X')
$$

Il est clair que les $X'_k$ forment un recouvrement de $X'$ ; nous allons d'abord voir qu'il en est de même des $X''_k$, en prouvant que $X'_k \subset X''_k$.
Il suffira pour cela de montrer que le diagramme
(3)
$$
\begin{tikzcd}
X'_k \arrow[r, "g|_{X'_k}"] \arrow[d] & P \arrow[d, "p_k"] \\
X_k \arrow[r, "\psi_k"] & P_k
\end{tikzcd}
$$
est commutatif. Or $X'_k$ s'identifie au préschéma induit $X' \cap (X_k \times_S P)$ et comme $U' \subset X'_k$, $X'_k$ est le plus petit sous-préschéma fermé de lui-même majorent $U'$ (§ 1, n°2, cor. 2 de la prop. 8). Pour montrer la commutativité (3), il suffit donc de prouver qu'en le composant avec l'injection $U' \to X'_k$, ou ce qui revient au même, avec $\psi$, on obtient un diagramme commutatif ($P_k$ étant séparé au-dessus de $S$ ; § 1, n°2, cor. 3 de la prop. 8) ; mais par définition le diagramme qu'on obtient ainsi n'est autre que (1), d'où notre assertion.

On voit donc que les $W_k$ forment un recouvrement ouvert de $g(X')$ ; pour prouver que $g$ est une immersion, il suffit donc de prouver que chacune des restrictions $g|_{X''_k}$ est une immersion dans $W_k$ (I, 2, 2, cor. 1 de la prop. 4). Pour cela, considérons le morphisme $u_k : W_k \xrightarrow{p_k} V_k \xrightarrow{j_k} X_k \to X$ ; comme $X$ est séparé sur $S$, le morphisme graphe $\Gamma_{u_k} : W_k \to X \times_S W_k$ est une immersion fermée, donc le gra-

\newpage

%% ===== Page 314 =====

\noindent - II-114 -

\noindent par $u_k(W_k)$ est un sous-préschéma fermé de $X \times_S W_k$; si nous montrons que ce sous-préschéma majore $U'$, il majorera aussi $X'_k \cap (X \times_S W_k)$ (§ 1, $n^\circ$2, cor. 2 de la prop. 8), et ce dernier sous-préschéma n'est autre que $X''_k$ par définition. Comme la restriction de $q_2$ à $\Gamma_{u_k}(W_k)$ est un isomorphisme sur $W_k$, la restriction de $g$ à $X''_k$ sera une immersion dans $W_k$, et notre assertion sera démontrée. Or, pour prouver que $\Gamma_{u_k}(W_k)$ majore $U'$, il suffit, comme on le voit aisément à partir de la définition du produit, de vérifier que le diagramme
$$\begin{tikzcd}
& & W_k \arrow[d, "u_k"] \\
U \arrow[r, "\psi"] & X \times_S W_k \arrow[ur, "q_2"] \arrow[r, "q_1"] & X
\end{tikzcd}$$
est commutatif ; or, par définition de $u_k$, et les relations $q_1 \circ \psi = j$, $q_2 \circ \psi = \psi$, on est ramené à la commutativité du diagramme (1).

\noindent E) Il est clair que comme $U$, et par suite $U'$, est irréductible, il en est de même de $X'$ dans la construction précédente. Si en outre $X$ est réduit, il en est de même de $U'$, donc $X'$ est aussi réduit (§ 1, $n^\circ$2, cor. 4 de la prop. 8). Cela achève la démonstration.

\noindent \emph{Corollaire 1.} — Soient $S$ un schéma noethérien, $X$ un schéma séparé et de type fini sur $S$. Pour que $X$ soit propre sur $S$, il faut et il suffit qu'il existe un schéma $X'$ projectif sur $S$ et un $S$-morphisme surjectif $f : X' \to X$. Lorsqu'il en est ainsi, on peut en outre choisir $f$ de sorte qu'il existe un ouvert dense $U$ de $X$ tel que la restriction de $f$ à $f^{-1}(U)$ soit un isomorphisme sur $U$, $f^{-1}(U)$ étant dense en $X'$. Si en outre $X$ est irréductible (resp. réduit), on peut supposer qu'il en est de même de $X'$.

\noindent La condition est suffisante en vertu de la prop. 3 et du cor. 1 de la prop. 2. Elle est nécessaire, car avec les notations du th. 1, si $X$ est propre sur $S$, $X'$ est propre sur $S$, car il est projectif, donc propre sur $X$ (prop. 3), d'où notre assertion (prop. 2 (ii)); comme en outre $X'$ est quasi-projectif sur $S$, il est projectif sur $S$

\newpage

%% ===== Page 315 =====

- II-115 -

(n°1, cor. 1 de la prop. 3).

\emph{Remarque}. — Le $S$-morphisme surjectif $f : X' \to X$ est un morphisme birationnel lorsque $X$ et $X'$ sont irréductibles (I, 4, 3).

\emph{Corollaire 2}. — Soient $S$ un schéma localement noethérien, $X$ un $S$-préschéma séparé et de type fini sur $S$, $f_0 : X \to S$ son morphisme structural. Pour que $X$ soit propre sur $S$, il faut et il suffit que pour tout morphisme de type fini $S' \to S$, $f_0^{S'} : X \times_S S' \to S'$ soit une application fermée.

La condition est évidemment nécessaire (n°1, déf. 1) ; prouvons qu'elle est suffisante. La question étant locale sur $S$, on peut supposer $S$ noethérien ; si on reprend les notations de la démonstration du th. 1, il résulte de l'hypothèse que la projection $q_2 : X \times_S P \to P$ est une application fermée. Comme l'application $g : X' \to P$ est composée de $q_2$ et de l'immersion fermée $h : X' \to X \times_S P$ et que $g$ est une immersion, $g$ est une immersion fermée. On en conclut que $g$ est propre (n°1, prop. 2 (i)), et comme le morphisme $f : X' \to X$ est surjectif et que l'on a $g = f \circ f$, on conclut du cor. 1 de la prop. 2 du n°1 que $f_0$ est propre.

3. \emph{Sections rationnelles d'un préschéma propre au-dessus d'un autre}.

\emph{Définition 2}. — On dit qu'un point $x$ d'un préschéma $X$ est régulier (resp. normal) si l'anneau local $\mathcal{O}_x$ est régulier (resp. normal). On dit que $X$ est régulier (resp. normal) si chaque point de $X$ est régulier (resp. normal).

Un préschéma régulier ou normal est réduit. Si $X$ est localement noethérien, dire qu'un point $x \in X$ est normal signifie que $\mathcal{O}_x$ est intègre et intégralement clos ; comme un anneau local régulier a ces propriétés, tout point régulier est alors normal. Pour un point fermé $x$ d'un préschéma localement noethérien $X$, il revient au même de dire que $x$ est régulier, ou normal, ou que $\mathcal{O}_x$ est un anneau de valuation.

\newpage

%% ===== Page 316 =====

- II-116 -

\emph{discrète}.

Théorème 2. — Soit $S = \operatorname{Spec}(A)$ le spectre d'un anneau de valuation discrète $A$ ; soit $X \to S$ le morphisme structural ; si $s$ est le point générique de $S$, toute section de $f^{-1}(s)$ au-dessus de $s$ se prolonge d'une seule manière en une $S$-section de $X$. 
Il revient au même (I, 5, 2) de dire que toute $S$-section rationnelle de $X$ est partout définie, ou encore (I, 2, 6) que les points de $X$ à valeurs dans $A$ sont les mêmes que les points de $X$ à valeurs dans le corps des fractions $K$ de $A$.

Comme $S$ et $X$ sont localement noethériens, une section $g$ de $f^{-1}(s)$ au-dessus de $s$ équivaut à la donnée d'un isomorphisme $\theta : \mathcal{O}_x \to \mathcal{O}_s$, où $x = g(s) \in f^{-1}(s)$, sur $\mathcal{O}_s$. Soit $Y$ l'adhérence de $\{x\}$ dans $X$ ; comme $S$ est réduit, le morphisme $g$ se factorise en un morphisme de $\{s\}$ dans le sous-préschéma fermé réduit de $X$ ayant $Y$ pour espace de base, suivi de l'injection canonique $Y \to X$ (I, 3, 4, cor. de la prop. 9). On peut évidemment remplacer $X$ par $Y$ ($n^\circ 1$, prop. 2, (i) et (ii)), et par suite supposer $X$ intègre.
Le Lemme de Chow donne alors un préschéma intègre $X'$ et un morphisme surjectif, projectif et birationnel $f' : X' \to X$ tel que $f \circ f' : X' \to S$ soit projectif. Si $x'$ est le point générique de $X'$, il correspond à un isomorphisme $\theta'$ de $\mathcal{O}_x$ sur $\mathcal{O}_{x'}$ ; en composant l'isomorphisme réciproque $\theta'^{-1}$ avec $\theta$, on obtient une section $g'$ de $f'^{-1}(s)$ au-dessus de $s$, et on a évidemment $g = f' \circ g'$. Il suffira donc de prouver le théorème pour $g'$ et $X'$, autrement dit, on se ramène au cas où $X$ est projectif sur $A$, donc isomorphe à un sous-préschéma fermé d'un $S$-schéma $\mathbb{P}_A^n$ (§ 4, $n^\circ 6$, prop. 21) ; si $j$ est l'injection canonique $X \to \mathbb{P}_A^n$, $j \circ g$ est une section de $h^{-1}(s)$ au-dessus de $s$, en désignant par $h$ le morphisme structural $\mathbb{P}_A^n \to S$ ; si on peut le prolonger en une $S$-section $g_1$ de $\mathbb{P}_A^n$, on aura $g_1(S) \subset X$ puisque $X = \{x\}$ et $g_1(s) = x$ ; on sait alors, comme $X$ est réduit, que $g_1$ se factorise en une

\newpage

%% ===== Page 317 =====

$j \circ g$, où $\bar{g}$ est la $S$-section de $X$ cherchée (I,3,4,cor. de la prop.9).
On est finalement ramené ainsi au cas où $X=\mathbb{P}^n_A$. Comme $A$ est un anneau local, tout faisceau inversible sur $S$ est isomorphe à $\mathcal{O}_S$ ; donc (§ 4, n°5, cor.2 du th.2) un point de $X$ à valeurs dans $K$ s'identifie à une classe d'éléments $(e_0, e_1, \dots, e_n)$ de $K$, où $e_i \neq 0$ et les $e_i$ sont des éléments non tous nuls de $K$. Si $u$ est une uniformisante de $A$, en multipliant les $e_i$ par une puissance convenable de $u$, on peut supposer qu'ils appartiennent à $A$, et que l'un au moins est inversible. Mais alors (§ 4, n°5, cor.2 du th.2), le système $(e_0, e_1, \dots, e_n)$ définit aussi un point de $X$ à valeurs dans $A$, ce qui achève la démonstration.

Corollaire 1. — Soit $Y$ un préschéma localement noethérien (donc réduit) et normal, et soit $X$ un schéma propre au-dessus de $Y$. Soit $f$ une $Y$-section rationnelle de $X$, et soit $Y'$ l'ensemble des points $y \in Y$ où $f$ n'est pas définie. Alors $\operatorname{codim}_Y Y' \ge 2$.

Cela signifie (I,7,1,déf.2 et I,7,2,prop.6) que pour toute composante irréductible $Z$ de $Y'$, de point générique $z$, on a $\dim(\mathcal{O}_z) \ge 2$.

Dans le cas contraire, on aurait $\dim(\mathcal{O}_z)=1$ et par hypothèse $O_z$ (intègre et intégralement clos), $\mathcal{O}_z$ serait un anneau de valuation discrète (Bourbaki, Alg. comm.). Si alors $f$ est définie dans $U$, $U \cap \operatorname{Spec}(\mathcal{O}_z)$ n'est pas vide (I,5,2,lemme 2) et la restriction $f'$ de $f$ à $U \cap \operatorname{Spec}(\mathcal{O}_z)$ est une application rationnelle $\operatorname{Spec}(\mathcal{O}_z) \to X$ ; par hypothèse $f$ ne serait pas un morphisme (I,5,2,prop.7), ce qui contredit le th.2.

Corollaire 2. — Soient $S$ un préschéma localement noethérien, $X$ et $Y$ deux $S$-préschémas de type fini, tels que $Y$ soit normal (donc réduit) et $X$ propre sur $S$. Soit $f$ une $S$-application rationnelle de $Y$ dans $X$, et soit $Y'$ l'ensemble des $y \in Y$ où $f$ n'est pas définie. Alors

\newpage

%% ===== Page 318 =====

$\operatorname{codim}_Y(Y') \ge 2$.

Pour tout ouvert $U \subset Y$, $\epsilon \to (\epsilon, 1)_S$ est une application biunivoque de l'ensemble des $S$-morphismes $g$ de $U$ dans $X$ sur l'ensemble des $U$-sections de $X \times_S U$ (ou de $X \times_S Y$) ; il s'ensuit que l'on peut identifier l'ensemble des $S$-applications rationnelles de $Y$ dans $X$ à l'ensemble des sections rationnelles de $X^Y$ ; comme $X^Y$ est propre sur $Y$ ($n^\circ 1$, prop. 2 (iii)), on peut appliquer le cor. 1, d'où la conclusion.

\newpage

%% ===== Page 319 =====

\marginpar{Archives Grothendieck-sept 59}

\section*{-- II-105 quater -- (new version)}

Proposition 33. -- Soient $Y$ un préschéma noethérien, $X$ un $Y$-préschéma de type fini sur $Y$, $q: X \to Y$ le morphisme structural, $S$ une $\mathcal{O}_Y$-algèbre graduée en degrés positifs, $L$ un $\mathcal{O}_Y$-module inversible, $\varphi$ un homomorphisme de $\mathcal{O}_X$-algèbres graduées $q^*(S) \to \bigoplus_{n \ge 0} L^{\otimes n}$, surjectif dans les degrés assez grands ; soit $r = r_{L, \varphi}: X \to \operatorname{Proj}(S)$ le morphisme correspondant, et supposons que $r$ soit une immersion (cor. 1 de la prop. 32). Alors, il existe un entier $n > 0$, un sous-$\mathcal{O}_Y$-module cohérent $E_n$ de $S_n$ tel que [le morphisme composé $\varphi_n \to q^*(S_n) \otimes L^{\otimes n}$] soit un épimorphisme et que le morphisme correspondant de $X$ dans $\mathbb{P}(E_n)$ soit une immersion ; en particulier $L$ est ample.

Soit $W$ un ouvert de $\operatorname{Proj}(S)$ tel que $r(X)$ soit fermé dans $W$ (I, 3, 1) ; par définition $W$ est réunion d'une famille $(W_\alpha)$ d'ouverts affines de la forme $D_+(f_\alpha)$, où $f_\alpha$ est une section de $S_{n_\alpha}$ au-dessus d'un ouvert affine $U_\alpha$ de $Y$. Par hypothèse $X$ est noethérien ($q$ étant de type fini donc quasi-compact, et il en est de même de $r(X)$ homéomorphe à $X$) ; il existe par suite une sous-famille finie $(W_i)_{1 \le i \le n}$ de $(W_\alpha)$ qui recouvre $r(X)$ ; en outre, comme $D_+(f_i^n) = D_+(f_i)$, on peut supposer les $f_i$ de même degré $n$, $n$ étant $E_n$ de plus choisi assez grand pour que $\varphi_n$ soit surjectif. Soit $V_i = r^{-1}(W_i)$ ; le sous-préschéma induit par $X$ sur $V_i$ est isomorphe à un sous-préschéma fermé de $W_i$, donc est affine et son anneau $A(V_i)$ est isomorphe à un anneau quotient de $A(W_i)$, ce dernier n'étant autre que $(S_i^{(n)})_{(f_i)}$, où $S_i = \Gamma(U_i, S)$ est une algèbre graduée sur $A_i = \Gamma(U_i, \mathcal{O}_Y)$. Comme $q$ est de type fini, $A(V_i)$ est une $A_i$-algèbre de type fini (I, 4, 2, prop. 6) ; considérons un système fini de générateurs de $A(V_i)$ sur $A_i$ ; comme $A(V_i)$ est une algèbre quotient de $(S_i^{(n)})_{(f_i)}$, ces générateurs sont images d'éléments de la forme $s_{ij}/f_i^{m_i}$ où les $s_{ij}$ sont des éléments homogènes de degré $nm_i$ ($1 \le j \le \lambda_i$) ; en multipliant les $s_{ij}$ par des puissances convenables des $f_i$, on

\newpage

%% ===== Page 320 =====

- II-105 quinquies - (new version)

peut supposer tous les $m_i$ égaux à un même entier $m$, et finalement, en remplaçant $n$ par $nm$, on peut supposer les $m_i$ égaux à $1$, les $f_i$ et les $g_{ij}$ appartenant au même module $E_i = \Gamma(U_i, S_n)$ ($1 \le i \le \mu, 1 \le j \le \lambda_i$) soit $R_i$ la partie finie de $E_i$ formée de ces éléments. Notons d'autre part que l'on peut supposer les $W_i$ choisis de façon que $L|V_i$ soit de la forme $\tilde{L}_i$, où $L_i$ est un module monogène libre sur $A(V_i)$ ; comme $q^*(S_n)|V_i \to L^{\otimes n}|V_i$ est surjectif par hypothèse, il existe un élément $h_i$ de $E_i$ dont l'image par $\varphi_n$ (qui est une section de $L^{\otimes n}$ au-dessus de $q^{-1}(U_i)$) donne par restriction à $V_i$ un générateur de $L_i^{\otimes n}$ ($I, 1, 3$, cor. 2 et 4 du th. 1) ; soit $R'_i = R_i \cup \{h_i\}$. Le sous-faisceau algébrique $E'_i$ de $S_n|U_i$ engendré par les éléments de $R'_i$ est de type fini, donc cohérent puisque $U_i$ est noethérien ; il existe par suite un sous-faisceau cohérent $E_i$ de $S_n$ induisant $E'_i$ sur $U_i$ (§ $1, n^\circ 2, th. 2$) ; si $E$ est le sous-faisceau de $S_n$ somme des $E_i$, il est clair que $E$ est cohérent et que le composé $q^*(E) \to q^*(S_n) \xrightarrow{\varphi_n} L^{\otimes n}$ est un épimorphisme. Considérons alors dans $\mathbb{P}(E)$ les ouverts affines $W'_i$ définis par les sections $f_i$ de $E$ au-dessus des $U_i$ ; le choix des $E_i$ et la description du morphisme $r' = r_{\varphi_n}$ montrent que la restriction de $r'$ à $V_i$ correspond à un homomorphisme surjectif $A(W'_i) \to A(V_i)$. On en conclut que cette restriction est une immersion fermée ; pour voir que $r'$ est une immersion, on peut se limiter au cas où $Y$ est affine, la question étant locale sur $Y$ ; mais alors, comme $W'_i$ et $V_i$ sont définis par la même section $f_i$, il résulte de la définition de $r'$ et de la prop. 2 du $n^\circ 1$ que $r'^{-1}(W'_i) = V_i$, ce qui achève la démonstration.

Corollaire. — Supposons en outre que $S = S_{\mathcal{O}_Y}(E)$, où $E$ est un $\mathcal{O}_Y$-module quasi-cohérent et que l'homomorphisme $\varphi_1 : q^*(E) \to L$ soit déjà surjectif ; alors la conclusion de la prop. 33 est déjà vraie pour $n=1$.

On notera ici que l'on peut supposer au début de la démonstration

\newpage

%% ===== Page 321 =====

\begin{itemize}
    \item[— II-105 sextiès —]
\end{itemize}

que les $f_i$ sont des sections de $S_1 = \underline{\underline{E}}$ ; les $g_{ij}$ sont alors des polynômes à coefficients dans $A_i$ par rapport à un nombre fini d'éléments de $\Gamma(U_i, \underline{\underline{E}})$ ; c'est ce dernier module qu'on dénote ici par $E_i$, et on prend pour $R_i$ l'ensemble formé de tous les éléments de $E_i$ au moyen desquels s'expriment les $g_{ij}$ ($1 \le j \le \lambda_i$). La fin de la démonstration est alors inchangée (avec $n=1$), d'où la conclusion.

\newpage

%% ===== Page 322 =====

\text{-- II-105 --}

sont des sections de $S_1 = E$ ; les $g_{ij}$ sont alors des polynômes à coefficients dans $A_i$ par rapport à un nombre fini d'éléments de $\Gamma(U_i, E)$ ; c'est ce dernier module qu'on dénote ici par $E_i$, et on prend pour $R_i$ l'ensemble formé de tous les éléments de $E_i$ au moyen desquels s'expriment les $g_{ij}$ ($1 \le j \le \lambda_i$). La fin de la démonstration est alors inchangée (avec $n=1$), d'où la conclusion.

\newpage

%% ===== Page 323 =====

-- ii-ix --

P. II-75, before prop.4, bring back prop.3 from p.II-78, and add:

Corollaire .- Supposons qu'il existe un recouvrement ouvert $(U_i)$ de $Y$ tel que chacune des algèbres graduées $\Gamma(U_i, S)$ soit de type fini sur $\Gamma(U_i, \mathcal{O}_Y)$. Alors il existe $d>0$ tel que $S^{(d)}$ soit engendrée par $S_d$.

Le corollaire résultera du lemme suivant :

Lemme 1 .- Soit $S$ une $A$-algèbre graduée de type fini. Il existe un entier $m > 0$ tel que $S_{mn} = (S_m)^n$ pour tout entier $n > 0$.

En effet, ce lemme étant supposé démontré, il correspondra à chaque algèbre graduée $\Gamma(U_i, S)$ un entier $m_i$, et il suffira de prendre pour $d$ le p.p.c.m. des $m_i$.

Pour démontrer le lemme, on peut supposer $S$ engendrée par un nombre fini d'éléments homogènes $f_i$, de degré $h_i$ ($1 \le i \le r$) ; soit $h$ le p.p.c.m. des $h_i$ et posons $g_i = f_i^{h/h_i}$, de sorte que tous les $g_i$ sont de degré $h$. Désignons par $R_k$ le $A$-module engendré par les monômes de degré $k$ par rapport aux $g_i$ ; il est clair que $R_p R_q = R_{p+q}$, et d'autre part, tout monôme de degré $jh$ par rapport aux $f_i$ peut s'écrire comme produit d'un monôme par rapport aux $g_i$ et d'un monôme de la forme $f_1^{t_1} f_2^{t_2} \dots f_r^{t_r}$, avec $t_i < h/h_i$. Ces monômes sont en nombre fini, désignons-les par $z_j$ ($1 \le j \le e$) ; $z_j$ est de degré multiple de $h$, soit $e_j h$. Soit $p > \max(e_j)$, et prenons $m = ph$ ; pour tout entier $k > 0$, on a $S_{mk} = R_{mk} + R_{mk-e_1 h} z_1 + \dots + R_{mk-e_j h} z_j$, d'où $S_{mk} = R_{(k-1)m} S_m = R_{m}^{k-1} S_m = S_m^k$, ce qui prouve que $S_{mn} = S_m^n$.

Nous dirons qu'une $\mathcal{O}_Y$-algèbre $S$ est intègre si pour chaque $y \in Y$, l'anneau local $S_y$ est intègre.

P. II-76, last line, before "Soit $S'$", insert :

\newpage

%% ===== Page 324 =====

\noindent Supposons de plus que pour tout $y \in Y$, l'annulateur de $(S_y)_+$ dans $\mathcal{O}_y$ soit réduit à $0$. Soit $S'$ le faisceau.

\noindent p. II-77, replace the last sentence by :
On peut supposer que $S=S_+$. Remarquons maintenant que l'annulateur de $S_+$ dans $A$ est réduit à $0$ ; s'il contenait en effet un élément $a \neq 0$, comme on a $a/1 \neq 0$ dans $A_y$ (puisque $A$ est intègre), $a/1$ appartiendrait à l'annulateur de $(S_y)_+=(S_+)_y$, contrairement à l'hypothèse. Nous allons en déduire que si $a \neq 0$ dans $A$, la relation $af=0$ pour un élément homogène $f \in S_+$ entraîne $f=0$. En effet, il existe $g \in S_+$ tel que $g^k a \neq 0$ pour tout $k > 0$, car $g^k a = 0$ entraîne $(ga)^k = 0$ et par suite $ga = 0$, et on sait qu'il y a au moins un $g$ homogène ne satisfaisant pas à cette condition. Cela entraîne que $a/1 \neq 0$ dans $S_{(g)}$ ; comme le préschéma $D_+(g) = \operatorname{Spec}(S_{(g)})$ est non vide et irréductible par hypothèse, $S_{(g)}$ est intègre, donc la relation $(a/1)(f^d/g^d) = 0$ dans $S_{(g)}$ entraîne $f^d/g^d = 0$ ($d, d'$ degrés de $f$ et $g$ respectivement), donc $g^{k'}f^d = 0$ pour un $k' > 0$, et comme $S_+$ ne contient pas d'élément nilpotent $f=0$ ; mais cela signifierait que $D_+(fg) = \emptyset$, et par suite $D_+(f) \cap D_+(g) = \emptyset$, avec $D_+(f) \neq \emptyset, D_+(g) \neq \emptyset$, contrairement à l'hypothèse que $X$ est irréductible. Le même raisonnement montre qu'il n'y a pas de diviseur de $0$ dans $S_+$; comme $S$ est intègre, il en est de même de $S$, et par suite aussi chaque $S_y$ est intègre, ce qui achève la démonstration.

\noindent p. II-82, before $n^\circ 3$, insert :

\noindent Remarque. - Les homomorphismes composés
$$\begin{tikzcd}
\tilde{\mathcal{M}} \ar[r, "\alpha"] \ar[d, phantom, ""{coordinate, name=Y}] & (\Gamma(\tilde{\mathcal{M}}))^\sim \ar[d, bend left=40, "\Gamma(f)"] \\
\Gamma(\mathcal{F}) \ar[r, "\Gamma(f)"] & \Gamma((\Gamma(\mathcal{F}))^\sim) \ar[u, phantom, ""{coordinate, name=X}] \ar[r, "\Gamma(f)"] & \Gamma(\mathcal{F})
\end{tikzcd}$$
sont les isomorphismes identiques (lorsqu'ils sont définis). La question est en effet locale sur $Y$, et on est donc ramené aux résultats

\newpage

%% ===== Page 325 =====

- II-xi -

correspondants du § 3 ($n^\circ 3$, formules (20) et (21)), compte tenu des propriétés usuelles des faisceaux quasi-cohérents (I, 1, 3).

P. II-80, before line -3, insert:

En vertu de la prop. 8, et de l'exactitude à gauche du foncteur $\Gamma_*$, le foncteur covariant $\Gamma_*$ est additif et exact à gauche. En particulier, si $\mathcal{J}$ est un faisceau d'idéaux dans $\mathcal{O}_X$, $\Gamma_*(\mathcal{J})$ s'identifie à un sous-faisceau gradué de $\Gamma_*(\mathcal{O}_X)$.

P. II-88, replace lines 5 to 12 by:

On peut donc dire que $\underline{P}(\mathcal{E})$ est un foncteur contravariant lorsqu'on convient de ne prendre comme morphismes pour les $\mathcal{O}_Y$-modules quasi-cohérents que les homomorphismes surjectifs. En outre, si on pose $P=\underline{P}(\mathcal{E})$, $P'=\underline{P}(\mathcal{F})$ et $j=\underline{P}(u)$, on a
(14) $$ j^*(\mathcal{O}_P(n)) \cong \mathcal{O}_{P'}(n) $$
comme il résulte du § 3, $n^\circ 5$, cor. de la prop. 20.

P. II-88, after equation (16), insert:

Lorsque l'on prend $\mathcal{E}=\mathcal{O}_Y^n$, on écrit $\mathbb{P}_Y^{n-1}$ au lieu de $\underline{P}(\mathcal{E})$; si en outre $Y$ est affine d'anneau $A$, on écrit aussi $\mathbb{P}_A^{n-1}$.

P. II-91, between lines -7 and -8, insert:

est défini de même comme un morphisme $D_+(gg') \to D_+(ff')$, à partir de l'homomorphisme $S_{(ff')} \to S_{gg'}$ qui fait correspondre à $x_1 \dots x_{2n}/(ff')^n$ l'élément $v_c(x_1) \dots v_c(x_{2n})/(gg')^n$ ; on vérifie aussitôt que $r_{ff'}$

\newpage

%% ===== Page 326 =====

\noindent P. II-88 , before line -12 , insert :

\noindent Proposition 16 .-- Soit $L$ un faisceau inversible sur $Y$ . Pour tout $\mathcal{O}_Y$-module quasi-cohérent $E$ , il existe un $Y$-isomorphisme de $\mathbb{P}(E)$ sur $\mathbb{P}(E \otimes L)$ .

\noindent Remarquons d'abord que si $A$ est un anneau , $E$ un $A$-module , $L$ un $A$-module monogène, on définit un homomorphisme de $A$-modules $S_n(E \otimes L) \to S_n(E) \otimes L^{\otimes n}$ en faisant correspondre à $(x_1 \otimes y_1) \dots (x_n \otimes y_n)$ l'élément $(x_1 x_2 \dots x_n) (y_1 \otimes y_2 \otimes \dots \otimes y_n)$ ($x_i \in E$, $y_i \in L$ pour $1 \le i \le n$) ; il est immédiat que cet homomorphisme est en fait un isomorphisme, en se ramenant au cas où $L=A$. On en conclut un isomorphisme canonique de modules gradués $S_\bullet(E \otimes L) \cong \bigoplus_{n \ge 0} S_n(E) \otimes L^{\otimes n}$.

\noindent En revenant aux conditions de la prop.16 , les remarques qui précèdent définissent un isomorphisme canonique $S_\bullet(E \otimes L) \cong \bigoplus_{n \ge 0} S_n(E) \otimes_{\mathcal{O}_Y} L^{\otimes n}$, en définissant cet isomorphisme comme isomorphisme de préfaisceaux. La proposition résulte alors de la prop.3 (iii) du n°1 .

\marginpar{et tenant compte de I,1, 3, cor.4 du th.1}

\noindent P. II-94 , after line 11 , insert :

\noindent Proposition 17 .-- (i) Soit $u : X' \to X$ un morphisme : si le $Y$-morphisme $r : X \to \mathbb{P}$ correspond à l'homomorphisme $\varphi : q^*(E) \to L$, le $Y$-morphisme $r'$ correspond à $u^*(\varphi) : u^*(q^*(E)) \to u^*(L)$ .

\noindent (ii) Soient $E, F$ deux $\mathcal{O}_Y$-modules quasi-cohérents, $v : E \to F$ un homomorphisme surjectif ; $j : \mathbb{P}(F) \hookrightarrow \mathbb{P}(E)$ l'immersion fermée correspondante. Si le $Y$-morphisme $r : X \to \mathbb{P}(F)$ correspond à l'homomorphisme $\varphi : q^*(F) \to L$, le $Y$-morphisme $j \circ r$ correspond à $q^*(E) \xrightarrow{q^*(v)} q^*(F) \to L$.

\noindent (iii) Soit $\psi : Y' \to Y$ un morphisme, et posons $E' = \psi^*(E)$. Si le $Y$-morphisme $r : X \to \mathbb{P}$ correspond à l'homomorphisme $\varphi : q^*(E) \to L$, le $Y'$-morphisme $r' : X' \to \mathbb{P}(E')$ correspond à $\varphi' = \varphi \otimes 1_{Y'} : q^*(E) \otimes_{\mathcal{O}_Y} \mathcal{O}_{Y'} \to L \otimes_{\mathcal{O}_Y} \mathcal{O}_{Y'}$.

\noindent (i) découle immédiatement des définitions ; il en est de même de (ii)

\newpage

%% ===== Page 327 =====

- II-xiii -

compte tenu de (14) . D'autre part, d'après (15), on a le diagramme commutatif
\begin{tikzcd}
Y' \arrow[r, "p'"] \arrow[d, "\psi"] & P' = P \times_Y Y' \arrow[r, "r'"] \arrow[d, "j'"] & X' \arrow[d, "u"] \\
Y \arrow[r, "p"] & P \arrow[r, "r"] & X
\end{tikzcd}

En vertu de (16), on a $(r^{Y'})^*(\mathcal{O}_{P'}(1)) = (r')^*(\mathcal{O}_P(1)) = u^*(r^*(\mathcal{O}_P(1))) = u^*(\mathcal{L}) = \mathcal{L} \otimes \mathcal{O}_{Y'}$, et d'autre part (n°4) $v^*(\alpha)$ n'est autre que l'homomorphisme canonique $\alpha': (p^{Y'})^*(\mathcal{E}') \to \mathcal{O}_{P'}(1)$ ; d'où (iii) .

P. II-56, replace 3 last lines and first line of p. II-57 by :
Le foncteur $F \to F(n)$ est exact, car la question est locale, et pour un voisinage convenable $U$ de $x \in X$, $\mathcal{O}_X(n)|_U$ est isomorphe à $\mathcal{O}_X|_U$.

Soient $M, N$ deux $S$-modules gradués ; on définit sur $M \otimes N$ une structure de $S$-module gradué de la façon suivante. Sur le produit tensoriel $M \otimes N$, on peut définir une structure de $\mathbb{Z}$-module gradué en prenant $(M \otimes N)_q = \bigoplus_{m+n=q} (M_m \otimes N_n)$ où $\bigoplus_{m+n=q}$ est le produit tensoriel des injections $M_m \to M$, $N_n \to N$. Comme $M \otimes N = (M \otimes N)/P$, où $P$ est le sous-module de $M \otimes N$ formé des éléments $(xs) \otimes y - x \otimes (sy)$, où $x \in M, y \in N, s \in S$, et que $P$ est évidemment un sous-module gradué de $M \otimes N$, on obtient de cette façon une graduation sur $M \otimes N$.

P. II-57, after formula (12), replace the four next lines by :
Rappelons que pour des modules gradués, $\operatorname{Hom}_S(M,N)$ désigne, non l'ensemble de tous les homomorphismes de $M$ dans $N$, mais la somme (directe) des sous-groupes $(\operatorname{Hom}_S(M,N))_n$, où cette notation désigne l'ensemble des homomorphismes de degré $n$. Si $u$ est un homomorphisme de degré $n$ et $s \in S_k$, $s.u$ est de degré $k+n$, donc les $(\operatorname{Hom}_S(M,N))_n$ définissent sur $\operatorname{Hom}_S(M,N)$ une structure de $S$-module gradué.

P. II-69, first line after prop. 19, read "homomorphisme" instead of "isomorphisme" .

\newpage

%% ===== Page 328 =====

- II-xiv -

P. II-95 , before prop.19 , insert :
Corollaire .- Si on pose $P=\mathbb{P}(E \oplus F)$, $v^*(O_P(1))$ est canoniquement isomorphe à $O_{P_1}(1) \otimes O_{P_2}(1) = L$.
En effet , avec les notations de la démonstration de la prop.18 , le faisceau induit par $O_P(1)$ sur $D_+(f \otimes g)$ correspond canoniquement au $S(f \otimes g)$-module monogène libre ayant pour générateur $(f \otimes g)/1$, donc le faisceau induit par $v^*(O_P(1))$ sur $D_+(f) \times D_+(g)$ correspond au $(S(f) \otimes S(g))$-module monogène ayant le même générateur ; en lui faisant correspondre le module monogène de générateur $(f/1) \otimes (g/1)$, on définit localement l'isomorphisme de l'énoncé ; il est immédiat de vérifier que ces isomorphismes locaux sont compatibles avec les opérations de restriction .

P. II-96 , before n°6 , insert :
Remarque .- Le préschéma somme de $\mathbb{P}(E)$ et $\mathbb{P}(F)$ est de même canoniquement isomorphe à un sous-préschéma fermé de $\mathbb{P}(E \oplus F)$. En effet les homomorphismes surjectifs $E \oplus F \to E$ et $E \oplus F \to F$ correspondent à des immersions fermées $\mathbb{P}(E) \to \mathbb{P}(E \oplus F)$ et $\mathbb{P}(F) \to \mathbb{P}(E \oplus F)$ ; tout revient à voir que les espaces de base des sous-préschémas formés de $\mathbb{P}(E \oplus F)$ ainsi obtenus sont sans point commun. La question étant locale sur $Y$, on peut supposer que $Y=\operatorname{Spec}(A)$, $E=\tilde{E}$, $F=\tilde{F}$, $E$ et $F$ étant des $A$-modules ; or $S_n(E)$ et $S_n(F)$ s'identifient à des sous-modules (d'intersection réduite à $0$) de $S_n(E \oplus F)$, et si $p$ est un idéal premier gradué de $S_A(E)$ tel que $p_n \supset S_n(E)$, il lui correspond dans $S_A(E \oplus F)$ un idéal premier gradué dont la trace sur $S_A(E)$ est $p_n$, mais qui contient $S_+(F)$, comme on le voit aussitôt ; deux points de $\operatorname{Proj}(S_A(E))$ et $\operatorname{Proj}(S_A(F))$ respectivement ne peuvent donc avoir même image dans $\operatorname{Proj}(S_A(E \oplus F))$.

\newpage

%% ===== Page 329 =====

\noindent - II-XV -

\vspace{1em}

\noindent P. II-98, before th.4, insert :

\noindent Corollaire. — Soit $f : X \to Y$ un morphisme projectif. Pour tout ouvert $U \subset Y$, la restriction de $f$ au sous-préschéma induit $f^{-1}(U)$ est un morphisme projectif.

\noindent En effet, cette restriction s'identifie à $f^U : X \times_Y U \to U$, et il suffit d'appliquer le th.3 (iii).

\noindent Remarque. — Soient $f : X \to Y$, $f' : X' \to Y$ deux morphismes projectifs, et soient $Z$ le préschéma somme de $X$ et $X'$, $g$ le morphisme $Z \to Y$ coïncidant avec $f$ dans $X$, avec $f'$ dans $X'$. Si $f$ et $f'$ sont projectifs, il en est de même de $g$. En effet, si $X$ (resp. $X'$) s'identifie à un sous-préschéma fermé de $\mathbb{P}(\underline{E})$ (resp. $\mathbb{P}(\underline{E}')$), où $\underline{E}$ et $\underline{E}'$ sont des $\mathcal{O}_Y$-modules quasi-cohérents de type fini, on a vu (n$^\circ$5, Remarque) que $Z$ s'identifie à un sous-préschéma fermé de $\mathbb{P}(\underline{E} \oplus \underline{E}')$.

\vspace{1em}

\noindent P. II-101, before prop.26, insert :

\noindent Remarque. — Si $f : X \to Y$, $f' : X' \to Y$ sont deux morphismes quasi-projectifs, $Z$ le préschéma somme de $X$ et $X'$, le morphisme $g : Z \to Y$ coïncidant avec $f$ dans $X$ et avec $f'$ dans $X'$, est quasi-projectif, comme il résulte aussitôt de la Remarque du n$^\circ$5.

\vspace{1em}

\noindent P. II-104, before prop.31, insert :

\noindent Remarque. — Dans la démonstration de la prop.30, pour prouver que $(r, r')_Y$ est une immersion, il suffit que l'un des morphismes $r, r'$ soit une immersion ; on voit donc qu'on peut affirmer que $L \otimes L'$ est très ample lorsque $L$ est très ample et qu'il existe un homomorphisme surjectif $q^* (\underline{E}') \to L'$, où $q : X \to Y$ est le morphisme structural, et $\underline{E}'$ un $\mathcal{O}_Y$-module quasi-cohérent de type fini.

\newpage

%% ===== Page 330 =====

- II-xvi -

P. II-109, replace last 3 lines and first lines of p. II-110, until prop. 3, by :

d'après la remarque faite ci-dessus, dans le diagramme commutatif
\begin{tikzcd}
X \times_Y Y' \arrow[r] \arrow[d] & Y \times_Y Y' = Y' \arrow[d] \\
X \times_{S'} Y' \arrow[r] & Y \times_{S'} Y'
\end{tikzcd}
les flèches verticales sont des immersions fermées ; il en résulte immédiatement que si $X \times_{S'} Y' \to Y \times_{S'} Y'$ est une application fermée, il en est de même de $X \times_Y Y' \to Y'$.

P. II-11, before n° 3, add :

Corollaire 4. -- Si $Y$ est réduit, il en est de même de $\overline{Y}$.

Sinon, en vertu de I, 3, 4, prop. 9 appliqué aux sous-espaces $Y$ et $\overline{Y}$, il existerait un sous-faisceau $\mathcal{J}$ de $\mathcal{O}_Y$ strictement plus grand que $\mathcal{J}$, induisant $\mathcal{J}$ sur $\overline{Y}$, contrairement à la définition de $\overline{Y}$.

P. II-38, before Prop. 17, add :

Proposition 16 bis. -- Si $X \xrightarrow{f} Y$ est un morphisme entier, on a $\dim(X) = \dim(f(X)) \le \dim(Y)$.

En utilisant la prop. 16, et remplaçant $Y$ par le sous-préschéma fermé réduit ayant pour base $\overline{f(X)}$ (I, 3, 4, prop. 9 et cor., et ci-dessus, prop. 15 (v)) on se ramène au cas où $Y = f(X)$. Si $X_i$ est une composante irréductible de $X$, il résulte de la prop. 16 que $f(X_i)$ est irréductible de point générique $x_i$ ; on peut donc supposer $X$ et $Y$ irréductibles, avec $Y = f(X)$, et aussi que $X$ et $Y$ sont réduits (prop. 15 (vi)) ; enfin, en remplaçant $Y$ par un ouvert affine de $Y$, on peut supposer $Y$ (et par suite $X$) affine. Si $X = \operatorname{Spec}(B)$, $Y = \operatorname{Spec}(A)$, $B$ est alors un anneau entier sur $A$, donc il résulte du premier th. de Cohen-Seidenberg que $\dim(B) = \dim(A)$.

\newpage

%% ===== Page 331 =====

\begin{center}
-- II-xvii --
\end{center}

p. II-71-72, cor. 1 und 2 become cor. 2 and 3; insert before :

\noindent Corollaire 1. -- Sous les hypothèses de la prop. 21 [unclear: (i)], on a
$$\Phi((s(n)))\tilde{} = (s'(n))\tilde{} \quad \text{pour tout } n \in \mathbb{Z}, \text{ et par suite } \Phi(\mathcal{F}(n)) = \Phi(\mathcal{F})(n) \quad \text{pour tout } \mathcal{O}_X\text{-module } \mathcal{F}.$$

\noindent Soit en effet $f \in S$, et $f' = \varphi(f)$; on déduit de $\varphi$ un homomorphisme
$$(s(n))(f) \otimes_{S(f)} S'(f') \to (s'(n))(f')$$
en faisant correspondre à $(x_{n+k}/f^k) \otimes 1$ ($x_{n+k} \in S_{n+k}$) l'élément
$x'_{n+k}/f'^k$ où $x'_{n+k} = \varphi(x_{n+k})$; cet homomorphisme est surjectif puisque
$\varphi$ l'est ; en outre, tout élément de $(s(n))_{(f)} \otimes_{S_{(f)}} S'_{(f')}$ peut se
mettre sous la forme $(x_{n+k}/f^k) \otimes 1$; si on a $f^h x_{n+k} = 0$
pour un $h$, on en conclut que $f'^h x'_{n+k} \in I$, d'où $f'^h x'_{n+k}/f'^n+h+k = 0$
dans $S_{(f')}$, et comme $(x_{n+k}/f^k) \otimes 1 = (f^h x_{n+k}/f^{n+h+k}) \otimes 1$,
on a bien $(x_{n+k}/f^k) \otimes 1 = 0$, ce qui démontre le corollaire.

\noindent p. II-72, after cor. 3, add :

\noindent Remarque. -- Supposons que l'homomorphisme $\varphi : S \to S'$ soit tel que
la restriction $S_n \to S'_n$ soit un isomorphisme lorsque $n$ est assez grand.
Comme on peut supposer que $S \to S'$ est aussi un isomorphisme (en remplaçant $S_0$ et $S'_0$ par $\mathbb{Z}$, cf. n° 3, Remarque finale), on en conclut
que pour $d$ assez grand, $S_{(d)} \to S'_{(d)}$ est un isomorphisme, donc (n°
2, Remarque) $\operatorname{Proj}(S)$ et $\operatorname{Proj}(S')$ s'identifient par $\varphi$, et cet
isomorphisme identifie $\mathcal{O}_X(n)$ et $\mathcal{O}_{X'}(n)$ pour tout $n$.

\noindent p. II-87, cor. becomes cor. 2; insert before :

\noindent Corollaire 1. -- Sous les conditions de la prop. 15 (i), et en supposant $S$ engendré par $S_1$ : $\mathcal{O}_X(n)$ s'identifie à l'image réciproque de $\mathcal{O}_X(n)$ par l'injection canonique.

\noindent Comme la question est locale, on peut supposer $Y$ affine, et
on est ramené au cor. 1 de la prop. 21 du § 3, n° 5.

\newpage

%% ===== Page 332 =====

-- II-xviii --

The last lines of p.II-76 have been destroyed by error. Read :

\emph{Proposition 4}. — Si le faisceau quasi-cohérent d'algèbres graduées $S$ est intègre, alors $X = \operatorname{Proj}(S)$ est réduit. Si de plus $Y$ est irréductible et $S_0 = \mathcal{O}_Y$, alors $X$ est irréductible et son image dans $Y$ est partout dense.

On peut se ramener au cas où $Y$ est un schéma affine : c'est évident.

\newpage

%% ===== Page 333 =====

-- II-xix --

p. II-21, after the corollary, add :

Remarque .— Les démonstrations du th. 3 et de son corollaire montrent que ce dernier est encore valable sous les conditions plus faibles suivantes sur $K'$ :

1) La somme directe de deux faisceaux de $K'$ est dans $K'$.

2) Si $F^k \in K'$ pour un entier $k$, alors $F \in K'$.

3) Soit $u : F \to G$ un homomorphisme de faisceaux cohérents sur $X$, et supposons que $G$ (ainsi que $\ker u$ et $\operatorname{coker} u$) soient dans $K'$, et en outre que les dimensions des supports de $\ker u$ et $\operatorname{coker} u$ soient strictement inférieures à celle du support de $G$ (resp. $F$) (resp. $G \in K'$). On observera en effet que dans le cas a) du th. 3, on peut prendre pour $Y'$ et $Y''$ des réunions de composantes irréductibles de $Y$, et alors la dimension de $Y' \cap Y''$ est strictement inférieure à celle de $Y$ ; et dans le cas b), comme $Y$ est irréductible, toute partie fermée de $Y$, distincte de $Y$, a une dimension strictement inférieure à celle de $Y$.

p. II-43, add a new section :

6. Cônes projetants .

Soit $S$ un anneau gradué à degrés positifs. On dit que le schéma affine $X' = \operatorname{Spec}(S)$ est le cône projetant du spectre premier homogène $X = \operatorname{Proj}(S)$ ; l'ensemble $V(S_+)$ (fermé) des idéaux premiers de $S$ contenant $S_+$ est appelé le sommet de $X'$ ; il est réduit à un point lorsque $S_+$ est maximal, c'est-à-dire lorsque $S_0$ est un corps. L'ensemble ouvert $X'' = X' - V(S_+)$ est réunion des $S_f$, lorsque $f$ parcourt l'ensemble des éléments homogènes de $S_+$, car un idéal premier ne contenant pas $S_+$ a au moins un élément homogène de $S_+$ n'appartenant pas à $\mathfrak{p}$. Pour tout $f$ homogène dans $S_+$, l'injection canonique $S_{(f)} \to S_f$

\newpage

%% ===== Page 334 =====

- II-xx -

définit un morphisme $u_f : D(f) \to D_+(f)$, qui est surjectif d'après (6) ; en outre, si $g$ est un second élément homogène de $S_+$, l'application canonique $S_f \to S_{fg}$ transforme $S_{(f)}$ en $S_{(fg)}$ ; donc $u_{fg}$ est la restriction de $u_f$ à $D(fg)$, ce qui prouve que les $u_f$ sont les restrictions d'un morphisme canonique surjectif $u : X'' \to X$.

\emph{Proposition 22.} -- Pour tout $x \in X$, la fibre $u^{-1}(x)$ est un schéma affine sur $\kappa(x)$ isomorphe au spectre de $\kappa(x)[T]$, où $T$ est une indéterminée.

Observons d'abord que pour tout idéal premier gradué $\mathfrak{p} \subset S_+$, $u^{-1}(\mathfrak{p})$ est l'ensemble des idéaux premiers $\mathfrak{q}$ de $S$ tels que $\mathfrak{q} \cap S_n = \mathfrak{p} \cap S_n$ pour tout $n > 0$. On a donc $u^{-1}(D_+(f)) \subset D(f)$ pour tout $f$ homogène dans $S_+$, et comme les $D_+(f)$ recouvrent $X$, $u^{-1}(x)$ est isomorphe à un préschéma de la forme $D(f) \times_{D_+(f)} \operatorname{Spec}(\kappa(x))$, où $x \in D_+(f)$. Or, on a $S_f = S_{(f)}[f/1]$, et il est immédiat que $f/1$ et ses puissances sont des éléments libres sur $S_{(f)}$ ($n^\circ 3$, prop. 13), donc $S_f$ est une $S_{(f)}$-algèbre isomorphe à $S_{(f)}[T]$, et la proposition en résulte aussitôt.

Si $S_0$ est un corps $k$, il revient au même de dire que $X$ est intègre ou que son cône projetant $X''$ est intègre ($n^\circ 3$, prop. 7).

\newpage

%% ===== Page 335 =====

\v - II-xxi -

p. II-104 ter, the cor.2 is incorrect, and should be replaced by :

Corollaire 2. - Soient $f : X \to Y$, $g : Y \to Z$ deux morphismes, $Z$ étant supposé noethérien. Si $\mathcal{L}$ est un faisceau ample sur $X$ pour $f$, $\mathcal{K}$ un faisceau ample sur $Y$ pour $g$, il existe un entier $n > 0$ tel que $\mathcal{L} \otimes f^*(\mathcal{K}^{\otimes n})$ soit ample sur $X$ pour $g \circ f$.

En effet $(\mathcal{L} \otimes f^*(\mathcal{K}^{\otimes n}))^{\otimes m} = \mathcal{L}^{\otimes m} \otimes f^*(\mathcal{K}^{\otimes nm})$, et on peut d'abord prendre $n$ assez grand pour que $\mathcal{K}^{\otimes n}$ soit très ample, puis $m$ assez grand pour que $\mathcal{L}^{\otimes m}$ soit très ample ; le corollaire résulte alors de la prop.31 (iv).

p. II-105 bis, in cor.1, suppress the first sentence ; the beginning of the proof should then read :

Notons que $r$ est un $Y$-morphisme de type fini, car son composé avec $\operatorname{Proj}(S) \to Y$ est de type fini (I, 4, 2, cor.2 de la prop.9) ; comme on suppose...

p. II-105 ter, suppress in cor.2 the assumption that $S$ is of finite type.

p. II-105 quater and 105 quinquies are rewritten (see further on).

p. II-106 ter and 106-quater, from line -6 of p.106 ter to line 3 of p.106-quater, replace "soit t.." by :

Par hypothèse, pour $n$ assez grand, il existe un voisinage ouvert affine $U$ de $q(x)$ dans $Y$ tel que $(\mathcal{J} \otimes \mathcal{L}^n)|q^{-1}(U)$ soit engendré par ses sections ; il existe donc une section $s$ de ce faisceau qui ne s'annule par en $x$ ; comme $x \notin \{y\}$, l'image $s$ de cette section dans $\mathcal{L}^n$ répond à la question.

p. II-106, in the statement of prop.34 (iii), read :

...tel que l'homomorphisme composé $q^*(E) \to q^*(q_*(L)) \to \mathcal{L}$ soit surjectif...

\newpage

%% ===== Page 336 =====

\marginpar{Archives Grothendieck Sept 59}

\begin{center}
- II-xix -
\end{center}

P. II-21, after the corollary, add :

\emph{Remarque}. — Les démonstrations du th.3 et de son corollaire montrent que ce dernier est encore valable sous les conditions plus faibles suivantes sur $K'$ :

1) La somme directe de deux faisceaux de $K'$ est dans $K'$.

2) Si $F^k \in K'$ pour un entier $k$, alors $F \in K'$.

3) Soit $u : F \to G$ un homomorphisme de faisceaux cohérents sur $X$, et supposons que $G$ (ainsi que $\ker u$ et $\operatorname{coker} u$) soient dans $K'$, et en outre que les dimensions des supports de $\ker u$ et $\operatorname{coker} u$ soient strictement inférieures à celle du support de $G$ (resp. $F$), alors $F \in K'$. 

On observera en effet que dans le cas a) du th.3, on peut prendre pour $Y'$ et $Y''$ des réunions de composantes irréductibles de $Y$, et alors la dimension de $Y' \cap Y''$ est strictement inférieure à celle de $Y$ ; et dans le cas b), comme $Y$ est irréductible, toute partie fermée de $Y$, distincte de $Y$, a une dimension strictement inférieure à celle de $Y$.

P. II-43, add a new section :

6. \emph{C\^{o}nes projetants}.

Soit $S$ un anneau gradué à degrés positifs. On dit que le schéma affine $X' = \operatorname{Spec}(S)$ est le cône projetant du spectre premier homogène $X = \operatorname{Proj}(S)$ ; l'ensemble fermé $V(S_+)$ des idéaux premiers de $S$ contenant $S_+$ est appelé le sommet de $X'$ ; il est réduit à un point lorsque $S_+$ est maximal, c'est-à-dire lorsque $S_0$ est un corps. L'ensemble ouvert $X'' = X' - V(S_+)$ est réunion des $D_f$, lorsque $f$ parcourt l'ensemble des éléments homogènes de $S_+$, car pour tout idéal premier $\mathfrak{p}$ ne contenant pas $S_+$, il y a au moins un élément homogène de $S_+$ n'appartenant pas à $\mathfrak{p}$. Pour tout $f$ homogène dans $S_+$, l'injection canonique $S_{(f)} \to S_f$

\newpage

%% ===== Page 337 =====

- II-xx -

définit un morphisme $u_f : D(f) \to D_+(f)$, qui est surjectif d'après (6) ; en outre, si $g$ est un second élément homogène de $S_+$, l'application canonique $S_f \to S_{fg}$ transforme $S_{(f)}$ en $S_{(fg)}$, donc $u_g$ est la restriction de $u_f$ à $D(fg)$, ce qui prouve que les $u_f$ sont les restrictions d'un morphisme canonique surjectif $u : X'' \to X$.

Proposition 22. - Pour tout $x \in X$, la fibre $u^{-1}(x)$ est un schéma affine sur $\kappa(x)$ isomorphe au spectre de $\kappa(x)[T]$, où $T$ est une indéterminée.

Observons d'abord que pour tout idéal premier gradué $\mathfrak{p} \subset S_+$, $u^{-1}(\mathfrak{p})$ est l'ensemble des idéaux premiers $\mathfrak{q}$ de $S$ tels que $\mathfrak{q} \cap S_n = \mathfrak{p} \cap S_n$ pour tout $n > 0$. On a donc $u^{-1}(D_+(f)) \subset D(f)$ pour tout $f$ homogène dans $S_+$, et comme les $D_+(f)$ recouvrent $X$, $u^{-1}(x)$ est isomorphe à un préschéma de la forme $D(f) \times_{D_+(f)} \operatorname{Spec}(\kappa(x))$, où $x \in D_+(f)$. Or, on a $S_f = S_{(f)}[f/1]$, et il est immédiat que $f/1$ et ses puissances sont des éléments libres sur $S_{(f)}$ ($n^\circ 3$, prop. 13), donc $S_f$ est une $S_{(f)}$-algèbre isomorphe à $S_{(f)}[T]$, et la proposition en résulte aussitôt.

Si $S_0$ est un corps $k$, il revient au même de dire que $X$ est intègre ou que son cône projetant $X''$ est intègre ($n^\circ 3$, prop. 7).

\newpage

%% ===== Page 338 =====

- II-xx -

définit un morphisme $u_f : D(f) \to D_+(f)$, qui est surjectif d'après (6) ; en outre, si $g$ est un second élément homogène de $S_+$, l'application canonique $S_f \to S_{fg}$ transforme $S_{(f)}$ en $S_{(fg)}$, donc $u_{fg}$ est la restriction de $u_f$ à $D(fg)$, ce qui prouve que les $u_f$ sont les restrictions d'un morphisme canonique surjectif $u : X^n \to X$.

Proposition 22. - Pour tout $x \in X$, la fibre $u^{-1}(x)$ est un préschéma affine sur $\kappa(x)$ isomorphe au spectre de $\kappa(x)[T]$, où $T$ est une indéterminée.

Observons d'abord que pour tout idéal premier gradué $\mathfrak{p} \subset S_+$, $u^{-1}(\mathfrak{p})$ est l'ensemble des idéaux premiers $\mathfrak{q}$ de $S$ tels que $\mathfrak{q} \cap S_n = \mathfrak{p} \cap S_n$ pour tout $n > 0$. On a donc $u^{-1}(D_+(f)) \subset D(f)$ pour tout $f$ homogène dans $S_+$, et comme les $D_+(f)$ recouvrent $X$, $u^{-1}(x)$ est isomorphe à un préschéma de la forme $D(f) \times_{D_+(f)} \operatorname{Spec}(\kappa(x))$, où $x \in D_+(f)$. Or, on a $S_f = S_{(f)}[f/1]$, et il est immédiat que $f/1$ et ses puissances sont des éléments libres sur $S_{(f)}$ ($n^\circ 3$, prop. 13), donc $S_f$ est une $S_{(f)}$-algèbre isomorphe à $S_{(f)}[T]$, et la proposition en résulte aussitôt.

Si $S_0$ est un corps $k$, il revient au même de dire que $X$ est intègre ou que son cône projetant $X^n$ est intègre ($n^\circ 3$, prop. 7).

\newpage

%% ===== Page 339 =====

\marginpar{Archives Grothendieck-sept 59}

\section*{III-1 - CHAPITRE III}

\section*{COHOMOLOGIE DES FAISCEAUX ALGÉBRIQUES COHÉRENTS. APPLICATIONS.}

\subsection*{§ 1. Cohomologie des schémas affines.}

\subsection*{1. Rappels sur le complexe de l'algèbre extérieure.}

Soient $A$ un anneau, $f = (f_i)_{1 \le i \le n}$ un système de $n$ éléments de $A$.
Le complexe de l'algèbre extérieure $K_*(f)$ correspondant à $f$ est un complexe de chaînes se définissant de la façon suivante : le $A$-module gradué $K_*(f)$ est égal à l'algèbre extérieure $\bigwedge(A^n)$, graduée de la façon usuelle, et l'opérateur bord est la multiplication intérieure par $i_f$ considéré comme élément du dual $A^n$ ; on rappelle que $i_f$ est une antidérivation de $\bigwedge(A^n)$, définie par la condition que, si $(e_i)_{1 \le i \le n}$ est la base canonique de $A^n$, $i_f(e_i) = f_i$ ; la vérification de la relation $i_f^2 = 0$ est immédiate.

Une définition équivalente est la suivante : pour chaque $i$, on considère un complexe de chaînes $K_*(f_i)$ défini de la façon suivante : $K_0(f_i) = K_1(f_i) = A$, $K_n(f_i) = 0$ pour $n > 1$ ; l'opérateur bord est défini par la condition que $d_0 : A \to A$ est la multiplication par $f_i$. On prend alors pour $K_*(f)$ le produit tensoriel $K_*(f_1) \otimes \dots \otimes K_*(f_n)$ (G.I, 2.7) muni de son degré total ; la vérification de l'isomorphisme de ce complexe et du complexe $K_*(f)$ défini ci-dessus est immédiate.

Pour tout $A$-module $M$, on définit le complexe de chaînes
$$(1) \qquad K_*(f, M) = K_*(f) \otimes_A M$$
et le complexe de cochaînes
$$(2) \qquad K^*(f, M) = \operatorname{Hom}_A(K_*(f), M)$$
Si $g$ est une $k$-cochaîne de ce complexe, il résulte des définitions précédentes que si on pose $g(i_1, \dots, i_k) = g(e_{i_1} \wedge \dots \wedge e_{i_k})$, $g$ est alternée et l'on a
$$(3) \qquad dg(i_1, \dots, i_{k+1}) = \sum_{j=1}^{k+1} (-1)^j f_{i_j} g(i_1, \dots, \hat{i}_j, \dots, i_{k+1})$$
Des complexes précédents, on déduit comme d'ordinaire les modules dérivés d'homologie et de cohomologie
$$(4) \qquad H_*(f, M) = H_*(K_*(f, M))$$

\newpage

%% ===== Page 340 =====

(5) $$ H^*(f, M) = H^*(K^*(f, M)) $$
On définit d'ailleurs un $A$-isomorphisme $K_*(f, M) \cong K^*(f, M)$ en faisant correspondre à toute chaîne $z = \sum (e_{i_1} \wedge \dots \wedge e_{i_k}) \otimes z_{i_1 \dots i_k}$ la cochaîne $\varepsilon_z$ telle que $\varepsilon_z(j_1, \dots, j_{n-k}) = \nu z_{j_1 \dots j_{n-k}}$, où $(j_h)_{1 \le h \le n-k}$ est la suite complémentaire de $(i_h)_{1 \le h \le k}$ dans $\{1, \dots, n\}$ et $\nu$ le nombre d'inversions de la permutation $i_1, \dots, i_k, j_1, \dots, j_{n-k}$ ; on vérifie que $\delta \varepsilon_z = d(\varepsilon_z)$ et par suite les modules d'homologie et de cohomologie (4) et (5) sont isomorphes. Dans ce chapitre, on étudiera surtout les groupes de cohomologie $H^n(f, M)$.

Il est immédiat que $H^*(f, M)$ est un foncteur cohomologique (T, II, 2.1) de la catégorie des $A$-modules dans celle des $A$-modules gradués, nul pour les degrés $< 0$ et $> n$. En outre, on a
(6) $$ H^0(f, M) = \operatorname{Hom}_A(A/(f), M) $$
où $(f)$ désigne l'idéal de $A$ engendré par $f_1, \dots, f_r$ ; ce module n'est autre que le sous-module de $M$ annulé par $(f)$. On a aussi
(7) $$ H^n(f, M) = (A/(f)) \otimes_A M = M / (\sum_{i=1}^n f_i M) $$
d'après l'isomorphisme entre $H^n$ et $H_0$ signalé ci-dessus.

Soit $g = (g_i)_{1 \le i \le r}$ une seconde suite de $r$ éléments de $A$ ; posons $fg = (f_i g_i)_{1 \le i \le r}$. On peut définir un homomorphisme canonique de complexes
$$ \rho_g : K_*(fg) \to K_*(f) $$
comme l'extension canonique à l'algèbre extérieure $\bigwedge(A^r)$ de l'application $A$-linéaire $(x_1, \dots, x_r) \to (g_1 x_1, \dots, g_r x_r)$ de $A^r$ dans $A^r$ ; pour voir qu'on a bien un homomorphisme de complexes, il suffit de remarquer, de façon générale, que si $u: E \to F$ est une application $A$-linéaire, $x \in F'$ et $y = {}^t u(x) \in E'$, alors on a la formule
(8) $$ (\wedge u) i_y = i_x (\wedge u) $$
en effet, les deux membres sont des antidérivations de $\bigwedge F$, et il suffit de vérifier qu'ils coïncident dans $F$, ce qui résulte aussitôt des définitions. Lorsqu'on identifie $K_*(f)$ au produit tensoriel

\newpage

%% ===== Page 341 =====

- III-3 -

$K_*(\underline{f}_i)$, $\varphi_{\underline{g}}$ est le produit tensoriel des $\varphi_{g_i}$, où $\varphi_{g_i}$ se réduit
à l'identité en degré $0$ et à la multiplication par $g_i$ en degré $1$.
En particulier, si $m, n$ sont deux entiers tels que $m \geq n \geq 0$, on a
des homomorphismes de complexes
$$\varphi_{\underline{f}^{m-n}} : K_*(\underline{f}^n) \to K_*(\underline{f}^m)$$
et par suite des homomorphismes (notés par le même symbole)
$$\varphi_{\underline{f}^{m-n}} : K^*(\underline{f}^n, M) \to K^*(\underline{f}^m, M)$$
$$H^*(\underline{f}^n, M) \to H^*(\underline{f}^m, M)$$

Ces derniers homomorphismes vérifient évidemment la condition de
transitivité $\varphi_{\underline{f}^{q-n}} = \varphi_{\underline{f}^{q-m}} \circ \varphi_{\underline{f}^{m-n}}$ pour $n \leq m \leq q$ ; ils définissent
donc deux systèmes inductifs ; on posera
(9)
$$C^*(\underline{f}, M) = \varinjlim K^*(\underline{f}^n, M)$$
$$H^*(\underline{f}, M) = H^*(C^*(\underline{f}, M)) = \varinjlim H^*(\underline{f}^n, M)$$

la dernière relation provenant du fait que le passage à la limite
inductive permute au foncteur $H^*$; on verra ultérieurement que
$H^*(\underline{f}, M)$ ne dépend en fait que de l'idéal $(\underline{f})$ (et même de la topologie $(\underline{f})$-adique sur $A$), ce qui justifie les notations.

Il est clair que $C^*(\underline{f}, M)$ est un foncteur $A$-linéaire exact en $M$ et
$H^*(\underline{f}, M)$ un foncteur cohomologique en $M$.

Soit $\underline{f}=(f_i) \in A^r$, $\underline{g}=(g_i) \in A^r$; désignons par
$e_{\underline{g}}$ la multiplication extérieure à gauche par le vecteur $\underline{g} \in A^r$ dans
l'algèbre extérieure $\bigwedge(A^r)$; on sait que l'on a la formule d'homotopie
$$i_{\underline{f}} e_{\underline{g}} + e_{\underline{g}} i_{\underline{f}} = \langle \underline{g}, \underline{f} \rangle 1$$
dans le $A$-module $A^r$ ($1$ désignant l'automorphisme identique de $A^r$);
cette relation signifie aussi que, dans le complexe $K_*(\underline{f})$, on a
$$d e_{\underline{g}} + e_{\underline{g}} d = \langle \underline{g}, \underline{f} \rangle 1$$
Si l'idéal $(\underline{f})=A$, il existe $\underline{g} \in A^r$ tel que $\langle \underline{g}, \underline{f} \rangle = \sum g_i f_i = 1$. Par suite
$(G, I, 2.4)$

\newpage

%% ===== Page 342 =====

\begin{itemize}
\item[-- III-4 --]
\end{itemize}

\underline{Proposition 1} .-- Supposons que l'idéal $(f)$ engendré par les $f_i$ soit
$A$. Alors le complexe $K_*(f)$ est homotopiquement équivalent à $0$, et
il en est de même des complexes $K_*(f, M)$ et $K^*(f, M)$ pour tout $A$-module
$M$.

\underline{Corollaire} .-- Si $(f)=A$, on a $H^*(f, M)=0$ ou $H^*((f), M)=0$ pour tout $A$-
module $M$.

On a en effet alors $(f^n)=A$.

2. Cohomologie de \v{C}ech d'un recouvrement ouvert .

\underline{Notations} .-- Dans ce qui suit, on désignera par :

$X$ un préschéma ;

$\mathcal{F}$ un faisceau quasi-cohérent sur $X$ ;

$A=\Gamma(X, \mathcal{O}_X)$, $M=\Gamma(X, \mathcal{F})$ ;

$f=(f_i)_{1 \le i \le n}$ un système fini d'éléments de $A$ ;

$U_i = X_{f_i}$ l'ensemble ouvert (II,1,3,cor.1 de la prop.13) des
$x \in X$ où la section $f_i$ "ne s'annule pas" ;

$U = \bigcup_{i=1}^n U_i$ ;

$\mathcal{U}$ le recouvrement $(U_i)_{1 \le i \le n}$ de $U$ .

Supposons que $X$ soit, ou bien noethérien, ou bien un schéma dont
l'espace de base est quasi-compact . On sait alors (cor.2 du
(II,1,4,th.2) que
l'on a $\Gamma(U_i, \mathcal{F})=M_{f_i}$. Nous poserons $U_{i_0 \dots i_p} = U_{i_0} \cap \dots \cap U_{i_p} = X_{f_{i_0} \dots f_{i_p}}$
(II,1,3,cor.2 de la prop.13) ; on a donc aussi
\begin{equation}
(10) \quad \Gamma(U_{i_0 \dots i_p}, \mathcal{F}) = M_{f_{i_0} \dots f_{i_p}}
\end{equation}

Or (chap.0, § 1, n°6) $M_{f_{i_0} \dots f_{i_p}}$ s'identifie à la limite inductive
$\varinjlim M_{i_0 \dots i_p}^{(n)}$, où le système inductif est formé des $M_{i_0 \dots i_p}^{(n)} = M$, l'ho-
momorphisme $\varphi_m : M_{i_0 \dots i_p}^{(n)} \to M_{i_0 \dots i_p}^{(m)}$ étant la multiplication par
$(f_{i_0} \dots f_{i_p})^{m-n}$ pour $m \le n$. Désignons par $C^p(M)$ l'ensemble des
applications alternées de $[1, n]^{p+1}$ dans $M$ (pour tout $n$) ; ces $A$-nodu-
les forment encore un système inductif pour les $\varphi_m$. Si $C^p(\mathcal{U}, \mathcal{F})$ est

\newpage

%% ===== Page 343 =====

- III - 5 -

le groupe des $p$-cochaînes alternées de Čech, relativement au recouvrement $\mathcal{U}$, et à coefficients dans $\mathcal{F}$, il résulte de ce qui précède que l'on peut écrire
$$(11) \qquad C^p(\mathcal{U}, \mathcal{F}) = \varinjlim_n C_n^p(M) .$$
Or, avec les notations du n°1, $C_n^p(M)$ s'identifie à $K^{p+1}((\underline{f}), M)$, et l'application $\varphi_{nm}$ s'identifie à l'application $\varphi_{n-m}$ définie au n°1. On a donc, pour tout $p \ge 0$, un isomorphisme canonique fonctoriel en $M$
$$(12) \qquad C^p(\mathcal{U}, \mathcal{F}) \simeq K^{p+1}((\underline{f}), M) .$$
La formule (3) du n°1 montre en outre que ces isomorphismes sont compatibles avec les opérateurs cobord. Il en résulte :

Proposition 2. — Si $X$ est noethérien, ou un schéma dont l'espace de base est quasi-compact, il existe pour tout $p \ge 1$ un isomorphisme canonique fonctoriel en $M$
$$(13) \qquad H^p(\mathcal{U}, \mathcal{F}) \simeq H^{p+1}((\underline{f}), M) .$$
On a en outre une suite exacte fonctorielle en $M$
$$(14) \qquad 0 \to H^0((\underline{f}), M) \to M \to H^0(\mathcal{U}, \mathcal{F}) \to H^1((\underline{f}), M) \to 0 .$$
Les relations (13) sont en effet évidentes d'après ce qui précède. D'autre part, $M = C^0((\underline{f}), M)$ et $C^0(\mathcal{U}, \mathcal{F}) = C^1((\underline{f}), M)$, donc $H^0(\mathcal{U}, \mathcal{F})$ s'identifie au sous-groupe des 1-cocycles de $C^1((\underline{f}), M)$ ; la suite exacte (3) [celle qui résulte de] n'est autre alors que la définition des groupes de cohomologie $H^0((\underline{f}), M)$ et $H^1((\underline{f}), M)$.

Corollaire. — Supposons qu'il existe des $a_i \in \Gamma(\mathcal{U}, \mathcal{O}_X)$ tels que $\sum a_i f_i = 1$. Alors, pour tout faisceau quasi-cohèrent $\mathcal{F}$ sur $\mathcal{U}$, on a $H^i(\mathcal{U}, \mathcal{F}) = 0$ pour $i > 0$ ; si en outre $X=\mathcal{U}$, l'homomorphisme canonique $M \to H^0(X, \mathcal{F})$ est bijectif. (les $X_{f_i}=U_i$ étant par hypothèse quasi-compacts)
On peut en effet supposer $X=\mathcal{U}$, et alors $H^i((\underline{f}), M) = 0$ pour tout $i \ge 0$ (n°1, cor. de la prop.1) ; l'assertion découle alors aussitôt de (13) et (14). On observera que, puisque $H^0(\mathcal{U}, \mathcal{F}) = H^0(\mathcal{U}, \mathcal{F})$, on redémontre ainsi le

\newpage

%% ===== Page 344 =====

- III-6 -

th.1 de I,1,3 comme cas particulier.

Remarque. — Supposons que $X$ soit un schéma affine et que les $U_i = X_{f_i}$ soient des ouverts affines ; il en résulte alors que les $U_{i_0 i_1 \dots i_p}$ sont des ouverts affines (I,3,6, prop.20 ; mais par contre on notera que $U$ n'est pas nécessairement affine). Dans ce cas, les foncteurs $\Gamma(U, \mathcal{F})$ et $\Gamma(U_{i_0 i_1 \dots i_p}, \mathcal{F})$ sont exacts en $\mathcal{F}$ (I,3, cor.5 du th.1). On en déduit que si on a une suite exacte de faisceaux quasi-cohérents $0 \to \mathcal{F}' \to \mathcal{F} \to \mathcal{F}'' \to 0$, la suite de complexes
$$0 \to C^*(U, \mathcal{F}') \to C^*(U, \mathcal{F}) \to C^*(U, \mathcal{F}'') \to 0$$
est exacte, et donne donc lieu à une suite exacte de cohomologie
$$\dots \to H^i(U, \mathcal{F}') \to H^i(U, \mathcal{F}) \to H^i(U, \mathcal{F}'') \xrightarrow{\partial} H^{i+1}(U, \mathcal{F}') \dots$$
D'autre part, comme le foncteur $C^*((\underline{f}), M)$ est exact, on a de même une suite exacte de cohomologie
$$\dots \to H^i((\underline{f}), M') \to H^i((\underline{f}), M) \to H^i((\underline{f}), M'') \xrightarrow{\partial} H^{i+1}((\underline{f}), M') \dots$$
où on a posé $M' = \Gamma(X, \mathcal{F}')$, $M'' = \Gamma(X, \mathcal{F}'')$, la suite $0 \to M' \to M \to M'' \to 0$ étant exacte par hypothèse. Cela étant, comme le diagramme
$$\begin{tikzcd}
0 \to C^*(U, \mathcal{F}') \to C^*(U, \mathcal{F}) \to C^*(U, \mathcal{F}'') \to 0 \\
0 \to C^*((\underline{f}), M') \to C^*((\underline{f}), M) \to C^*((\underline{f}), M'') \to 0
\end{tikzcd}$$
est commutatif, on en conclut que les diagrammes
$$\begin{tikzcd}
H^i(U, \mathcal{F}'') \xrightarrow{\partial} H^{i+1}(U, \mathcal{F}') \\
\downarrow \cong \downarrow \cong \\
H^{i+1}((\underline{f}), M'') \xrightarrow{\partial} H^{i+2}((\underline{f}), M')
\end{tikzcd}$$
sont aussi commutatifs (G,I,2.1.1).

3. Cohomologie d'un schéma affine.

Théorème 1. — Soit $X$ un schéma affine. Pour tout faisceau quasi-cohérent $\mathcal{F}$ sur $X$, on a $H^i(X, \mathcal{F}) = 0$ pour $i > 0$.

Soit $U$ un recouvrement fini de $X$ par des ouverts affines $X_{f_i}$ ($1 \le i \le r$) ; on sait qu'alors l'idéal engendré dans $A = \Gamma(X, \mathcal{O}_X)$ par les $f_i$ est égal à $A$. On conclut donc du cor. de la prop.2 ($n^\circ 2$) que l'on a $H^i(U, \mathcal{F}) = 0$ pour $i > 0$. Comme il y a de tels recouvrements

\newpage

%% ===== Page 345 =====

- III-7 -
arbitrairement fine, la définition de la cohomologie de Čech montre
que l'on a aussi $H^n(X, \mathcal{F})=0$ pour $n>0$. Mais ceci s'applique aussi à
tout ouvert qui est affine $X_\rho$, donc $H^n(X_\rho, \mathcal{F})=0$ pour $n>0$. Mais comme l'intersection $X_f \cap X_g = X_{fg}$, on en déduit que l'on a aussi $H^n(X, \mathcal{F})=0$ pour $n>0$ (G,II,5.9.2).

Corollaire 1. - Soit $f: X \to Y$ un morphisme affine. Pour tout faisceau quasi-cohérent $\mathcal{F}$ sur $X$, on a $R^q f_*(\mathcal{F})=0$ pour $q>0$.

En effet, on sait par définition que $R^q f_*(\mathcal{F})$ est le faisceau sur $Y$ associé au préfaisceau $U \mapsto H^q(f^{-1}(U), \mathcal{F})$, où $U$ parcourt les ensembles ouverts de $Y$. Mais les ouverts affines forment une base de la topologie de $Y$, et pour un tel ouvert $U$, $f^{-1}(U)$ est affine (II,2,3, cor.1 de la prop.7), donc $H^q(f^{-1}(U), \mathcal{F})=0$ d'après le th.1, ce qui démontre le corollaire.

Corollaire 2. - Soit $f: X \to Y$ un morphisme affine. Pour tout faisceau quasi-cohérent $\mathcal{F}$ sur $X$, $H^q(X, \mathcal{F})$ est canoniquement isomorphe à $H^q(Y, f_*(\mathcal{F}))$ pour tout $q \ge 0$.

On sait en effet que $H^q(X, \mathcal{F})$ est l'aboutissement d'une suite spectrale dont le terme $E_2$ est donné par $E_2^{pq} = H^p(Y, R^q f_*(\mathcal{F}))$ (suite spectrale de Leray; G,II,4.17.1). D'après le cor.1, $E_2^{pq}=0$ pour $q>0$; on sait alors que $H^p(X, \mathcal{F})$ est isomorphe à $E_2^{p0} = H^p(Y, f_*(\mathcal{F}))$ (G,I,4.4.1).

4. Application à la cohomologie des préfaisceaux quelconques.

Proposition 3. - Soient $X$ un schéma, $(U_\alpha) = \mathcal{U}$ un recouvrement de $X$ par des ouverts affines. Pour tout faisceau quasi-cohérent $\mathcal{F}$ sur $X$, les modules de cohomologie $H^n(X, \mathcal{F})$ et $H^n(\mathcal{U}, \mathcal{F})$ sont canoniquement isomorphes.

En effet, toute intersection finie $V$ d'ouverts $U_\alpha$ est un ouvert affine puisque $X$ est un schéma (I,3,6, prop.20), donc $H^q(V, \mathcal{F})=0$ pour $q \ge 1$. La proposition résulte alors du th. de Leray (G,II,

\newpage

%% ===== Page 346 =====

- III-8 -

5.4.1) .

Corollaire .— Soient $X$ un schéma dont l'espace de base est quasi-compact, $A = \Gamma(X, \mathcal{O}_X)$, $\mathfrak{f} = (f_i)_{1 \le i \le r}$ une famille finie d'éléments de $A$ tels que les $X_{f_i}$ soient affines, $U$ la réunion des $X_{f_i}$. Soient $F$ un faisceau quasi-cohérent sur $X$, $M = \Gamma(X, F)$. Alors, avec les notations du $n^\circ 1$, on a un isomorphisme fonctoriel
$$H^q(U, F) = H^{q+1}(\mathfrak{f}, M) \quad \text{pour } q \ge 1$$
et une suite exacte fonctorielle
$$0 \to H^0(\mathfrak{f}, M) \to M \to H^0(U, F) \to H^1(\mathfrak{f}, M) \to 0 .$$
Cela résulte en effet de la prop. 3 et de la prop. 2 du $n^\circ 2$.

La remarque suivant la prop. 2 du $n^\circ 2$ montre en outre que si $X$ est affine, les diagrammes
\begin{tikzcd}
H^q(U, F') \arrow[r] \arrow[d] & H^{q+1}(\mathfrak{f}, M') \arrow[d] \\
H^{q+1}(\mathfrak{f}, M) \arrow[r] & H^{q+2}(\mathfrak{f}, M'')
\end{tikzcd}
correspondant à une suite exacte $0 \to F' \to F \to F'' \to 0$ de faisceaux quasi-cohérents, sont commutatifs.

Proposition 4 .— Soit $u: X \to Y$ un morphisme séparé. On suppose que $Y$ est réunion d'une famille $(V_\alpha)$ d'ouverts affines tels que tout $u^{-1}(V_\alpha)$ soit réunion finie d'ouverts affines dans $X$ (ce qui est par exemple le cas si $u$ est de type fini). Alors, pour tout faisceau quasi-cohérent $F$ sur $X$, les $R^q u_*(F)$ sont quasi-cohérents sur $Y$.

La question étant locale sur $Y$, on peut supposer $Y$ affine d'anneau $A$ et $X$ réunion finie d'ouverts affines $U_i$ ($1 \le i \le r$) ; soit $\underline{U}$ le recouvrement $(U_i)$. Pour tout $f \in A$, les ensembles $U_{if} = u^{-1}(Y_f) \cap U_i$ sont affines (I, 3, 7, cor. 3 de la prop. 22) ; désignons par $\underline{U}_f$ le recouvrement $(U_{if})$ de $u^{-1}(Y_f)$. Pour toute suite $(i_k)_{0 \le k \le p}$ de indices de $[1, r]$, posons comme d'ordinaire $U_{i_0 i_1 \dots i_p} = \bigcap_{k=0}^p U_{i_k}$ ; comme $X$ est un schéma (I, 3, 6, prop. 17) les $U_{i_0 \dots i_p}$ sont des ou-

\newpage

%% ===== Page 347 =====

- III-9 -

verts affines (I,3,6,prop.20) ; nous écrirons pour simplifier $U_{\underline{s}}$ au lieu de $U_{i_0, \dots, i_p}$ en posant $\underline{s} = (i_0, \dots, i_p)$. Soit $B_{\underline{s}}$ l'anneau de $U_{\underline{s}}$, et soit $\varphi_{\underline{s}}$ l'homomorphisme $A \to B_{\underline{s}}$ correspondant à la restriction de $u$ à $U_{\underline{s}}$. Si on pose $M_{\underline{s}} = \Gamma(U_{\underline{s}}, F)$, le module $\Gamma(U_{\underline{s}} \cap u^{-1}(Y_f), F)$ s'identifie canoniquement à $(M_{\underline{s}})_f$ (I,1,3,prop.6 et I,1,2, prop.5). Si on considère $M_{\underline{s}}$ comme un $A$-module au moyen de l'homomorphisme $\varphi_{\underline{s}} : A \to B_{\underline{s}}$, $(M_{\underline{s}})_f$ s'écrit aussi $(M_{\underline{s}})_f$ et est un $A_f$-module. On en conclut que $C^p(U_f, F)$, en tant que $A_f$-module, s'identifie à $(C^p(U, F))_f$, où $C^p(U, F)$ est considéré comme $A$-module ; en outre, le cobord $C^p(U_f, F) \to C^{p+1}(U_f, F)$ s'identifie à l'homomorphisme $d_f$, en désignant par $d$ le cobord $C^p(U, F) \to C^{p+1}(U, F)$. Comme le foncteur $N \to N_f$ est exact dans la catégorie des $A$-modules, on a donc $H^*(C^*(U_f, F)) = (H^*(C^*(U, F)))_f = (H^*(U, F))_f$. En vertu de la prop.3, on voit donc qu'on a défini un isomorphisme canonique de $A_f$-modules

$$\rho_f : H^*(u^{-1}(Y_f), F) \cong T(Y_f, (H^*(X, F))^{\sim})$$

En outre, si $g$ est un second élément de $A$, on vérifie aussitôt que le diagramme

\begin{tikzcd}
H^*(u^{-1}(Y_f), F) \cong \Gamma(Y_f, (H^*(X, F))^{\sim}) \\
H^*(u^{-1}(Y_{fg}), F) \cong \Gamma(Y_{fg}, (H^*(X, F))^{\sim})
\end{tikzcd}

est commutatif. On en conclut l'existence d'un isomorphisme canonique fonctoriel

$$R^*u_*(F) \cong (H^*(X, F))^{\sim}$$

ce qui démontre la prop.6. On a prouvé en même temps :

Corollaire. - Sous les hypothèses précédentes, on a, pour tout ouvert affine $U \subset Y$

$$\Gamma(U, R^q u_*(F)) = H^q(u^{-1}(U), F) \quad (q \ge 0).$$

5. Le critère de Serre.

Le th.1 du n°3 admet une réciproque lorsque $X$ est noethérien :

\newpage

%% ===== Page 348 =====

- III-10 -

avant de l'énoncer, nous commencerons par caractériser les schémas affines parmi la catégorie des espaces annelés. Soit $(X, \mathcal{O}_X)$ un espace annelé, tels que les fibres $\mathcal{O}_x$ soient des anneaux locaux ; nous désignerons par $\mathfrak{m}_x$ l'idéal maximal de $\mathcal{O}_x$. Soit $A = \Gamma(X, \mathcal{O}_X)$, et considérons le schéma affine $X' = \operatorname{Spec}(A)$; on définit canoniquement un morphisme d'espaces annelés $\varphi = (\varphi, \theta) : (X, \mathcal{O}_X) \to (X', \mathcal{O}_{X'})$ de la façon suivante : pour tout $x \in X$, soit $\varphi(x)$ l'ensemble des $f \in A$ tels que $f_x \in \mathfrak{m}_x$ ; il est immédiat que $\varphi(x)$ est un idéal premier de $A$, donc un point de $X'$. En outre, cette définition montre que pour tout $f \in A$, on a
$$(15) \qquad \varphi^{-1}(D(f)) = X_f$$
et comme $X_f$ est ouvert (II, 1, 3, cor. 1 de la prop. 13) et que les $D(f)$ forment une base de la topologie de $X'$, $\varphi$ est continue. D'autre part, $\mathcal{O}_{\varphi(x)} = \mathcal{O}_{\varphi(x)}$; dans l'homomorphisme canonique $f \to f_x$ de $A$ dans $\mathcal{O}_x$ qui, à toute section $f$ fait correspondre son germe $f_x$ en $x$, il est clair par définition que la relation $f \in \varphi(x)$ entraîne $f_x \in \mathfrak{m}_x$ ; donc que $f_x$ est inversible dans $\mathcal{O}_x$ ; on en conclut que l'homomorphisme $f \to f_x$ se factorise en $A_{\varphi(x)} \to \mathcal{O}_x$. En outre les homomorphismes $\omega_x$ sur les fibres proviennent d'un homomorphisme $\omega : \varphi^*(\mathcal{O}_{X'}) \to \mathcal{O}_X$ de faisceaux ; il suffit en effet de définir un homomorphisme $\theta : \mathcal{O}_{X'} \to \varphi_*(\mathcal{O}_X)$ et de montrer que $\omega_x = \theta_x$. Or, pour tout $f \in A$, les sections de $\mathcal{O}_{X'}$ au-dessus de $D(f)$ s'identifient canoniquement aux éléments $a/f^n$ (I, 1, 3, prop. 6 et th. 1) ; en vertu de (15), $f$ est inversible au-dessus de $X_f = \varphi^{-1}(D(f))$, donc on peut faire correspondre à $a/f^n$ la section $(a|_{\varphi^{-1}(D(f))})(f|_{\varphi^{-1}(D(f))})^{-n}$ de $\mathcal{O}_X$ au-dessus de $\varphi^{-1}(D(f))$ et on a ainsi défini un homomorphisme $\Gamma(D^{-1}(f), \mathcal{O}_{X'}) \to \Gamma(D^{-1}(f), \varphi_*(\mathcal{O}_X))$ ; il est immédiat que ces homomorphismes sont compatibles avec les opérations de restriction de $D(f)$ à $D(ff')$, et

\newpage

%% ===== Page 349 =====

\marginpar{Archives Grothendieck - sept 59}

\section*{-- III-21 --}

\section*{§ 2. Etude cohomologique des morphismes projectifs.}

\subsection*{1. Calculs explicites de certains groupes de cohomologie.}

Soient $X$ un préschéma, $\mathcal{L}$ un faisceau inversible sur $X$ ; considérons l'anneau gradué
$$(1) \quad S = \Gamma_*(X, \mathcal{L}) = \bigoplus_{n \in \mathbb{Z}} \Gamma(X, \mathcal{L}^{\otimes n}) \, .$$
Soit $(f_i)_{1 \le i \le p}$ une famille finie d'éléments homogènes de $S$, soit $f_i \in S_{d_i}$ ; posons $U_i = X_{f_i}$, $U = \bigcup_i U_i$, et désignons par $\mathcal{U}$ le recouvrement $(U_i)$ de $U$. Pour tout faisceau quasi-cohérent $\mathcal{F}$ sur $X$, on posera
$$(2) \quad \mathcal{F}^{(*)} = \mathcal{F} \otimes \mathcal{L} \otimes \bigoplus_{n \in \mathbb{Z}} \mathcal{L}^{\otimes n}$$
$$(3) \quad H^q(\mathcal{U}, \mathcal{F}^{(*)}) = \bigoplus_{n \in \mathbb{Z}} H^q(U, \mathcal{F} \otimes \mathcal{L}^{\otimes n})$$
$$(4) \quad H^q(\mathcal{U}, \mathcal{F}^{(*)}) = \bigoplus_{n \in \mathbb{Z}} H^q(U, \mathcal{F} \otimes \mathcal{L}^{\otimes n})$$

Il est clair que le faisceau (2) est un $S$-module gradué, quand on identifie $S$ au faisceau simple correspondant. Les groupes abéliens (3) et (4) sont bigradués, en prenant $(H^q(\mathcal{U}, \mathcal{F}^{(*)}))_{mn} = H^q(U, \mathcal{F} \otimes \mathcal{L}^{\otimes n})$, et de même pour (4) ; pour le deuxième degré, il est clair que (4) est un $S$-module gradué, et il en est de même pour (3), comme le voit en considérant une résolution canonique de $\mathcal{F}$.

Considérons maintenant le $S$-module gradué (II, 1, 3)
$$M = \Gamma_*(U, \mathcal{F}^{(*)}) = H^0(U, \mathcal{F}^{(*)}) = \bigoplus_{n \in \mathbb{Z}} \Gamma(U, \mathcal{F} \otimes \mathcal{L}^{\otimes n}) \, .$$
Si $X$ est noethérien, ou séparé et a un espace de base quasi-compact il résulte de II, 1, 4, th. 2 qu'en posant comme d'ordinaire $U_{i_0 i_1 \dots i_p} = \bigcap_{k=0}^p U_{i_k}$, on a, à un isomorphisme canonique près
$$\Gamma(U_{i_0 i_1 \dots i_p}, \mathcal{F}^{(*)}) = M_{f_{i_0} \dots f_{i_p}}$$
On peut encore (avec les notations du § 1, n° 2), identifier $M_{f_{i_0} \dots f_{i_p}}$ à $\lim_{\longrightarrow} M_n^{(f_{i_0} \dots f_{i_p})}$. Cette identification est un isomorphisme de $S$-modules gradués, si on définit le degré d'un élément

\newpage

%% ===== Page 350 =====

- III-22 -

$z \in \varinjlim M_{i_0 \dots i_p}^{(n)}$ de la façon suivante : si $z$ est l'image canonique
d'un élément homogène $x \in M_{i_0 \dots i_p}^{(n)} = M$, de degré $n$, son degré est égal
à $n - (d_i + \dots + d_p)$. Tenant compte de la définition des homomorphismes
$M_{i_0 \dots i_p}^{(n)} \to M_{i_0 \dots i_p}^{(k)}$ (§ 1, n°2), on voit aussitôt que cette définition ne dépend pas du "représentant" $x$ de $z$ choisi. Les S-modules $\varinjlim G^p(M)$ sont gradués de la même façon, et les raisonnements du § 1, n°2 montrent que l'on a

(5) $$ \check{H}^p(\mathcal{U}, \mathcal{F}(*)) = \varinjlim \check{H}^p(M) $$

L'isomorphisme des deux membres respectant les degrés. On a donc aussi

(6) $$ \check{H}^p(\mathcal{U}, \mathcal{F}(*)) = \check{H}^{p+1}((\underline{f}), M) = \varinjlim \check{H}^{p+1}(\underline{f}, M) $$

l'isomorphisme conservant les degrés : en vertu de ce qui précède, le degré d'un élément du second membre image canonique d'une cochaîne $g \in K^{p+1}(\underline{f}, M)$, dont les valeurs $g(i_0, \dots, i_p)$ sont dans la même composante homogène $M$, est $n - (d_i + \dots + d_p)$, et il est indépendant du choix de cette cochaîne comme représentant de l'élément considéré.

Comme les isomorphismes précédents sont compatibles avec les opérateurs cobords, on en conclut, comme dans la prop. 2 du § 1, n°2 :

Proposition 1. — Soit $X$ un préschéma noethérien, ou un schéma dont l'espace de base est quasi-compact. Pour tout $q \ge 1$, il existe un isomorphisme canonique fonctoriel en $\mathcal{F}$

(7) $$ \check{H}^q(\mathcal{U}, \mathcal{F}(*)) \simeq H^{q+1}((\underline{f}), M) $$

On a en outre une suite exacte fonctorielle en $\mathcal{F}$

(8) $$ 0 \to H^0((\underline{f}), M) \to M \to \check{H}^0(\mathcal{U}, \mathcal{F}(*)) \to H^1((\underline{f}), M) \to 0 $$

De plus, tous les homomorphismes introduits respectent les structures de groupes abéliens bigradués et de S-modules gradués.

Corollaire. — Si $X$ est un schéma dont l'espace de base est quasi-com-

\newpage

%% ===== Page 351 =====

\subsection*{III-23}

pact et si les $U_i = X_{f_i}$ sont affines, il existe un isomorphisme canonique fonctoriel en $\mathcal{F}$

$$(9) \qquad H^q(U, \mathcal{F}(*)) \simeq H^{q+1}((\underline{f}), M) \quad \text{pour } q \ge 1$$

et une suite exacte fonctorielle en $\mathcal{F}$

$$(10) \qquad 0 \to H^0((\underline{f}), M) \to H^0(U, \mathcal{F}(*)) \to H^1((\underline{f}), M) \to 0$$

Il suffit d'appliquer la prop. 3 du § 1, n$^\circ$4.

La proposition "locale" analogue à la prop. 1 est la suivante :

\textbf{Proposition 2.} — Soient $S$ un anneau gradué à degrés positifs, $f_i$ ($1 \le i \le r$) un élément homogène de $S$ de degré $d_i$, $M$ un $S$-module gradué. Soient $X = \operatorname{Proj}(S)$ le spectre premier homogène de $S$, $\mathcal{F} = \tilde{M}$ le faisceau quasi-cohérent sur $X$ associé à $M$, et posons $U_i = D_+(f_i)$, $U = \cup U_i$. Prenons en outre $L = \mathcal{O}_X(1) = \mathcal{S}(1)^\sim$, de sorte que $\mathcal{F}(*) = \bigoplus_{n \in \mathbb{Z}} \mathcal{F}(n) = \bigoplus_{n \in \mathbb{Z}} (M(n))^\sim$. Il existe alors des isomorphismes canoniques fonctoriels en $M$

$$(11) \qquad H^q(U, \mathcal{F}(*)) \simeq H^{q+1}((\underline{f}), M) \quad \text{pour } q \ge 1$$

et une suite exacte fonctorielle en $M$

$$(12) \qquad 0 \to H^0((\underline{f}), M) \to H^0(U, \mathcal{F}(*)) \to H^1((\underline{f}), M) \to 0$$

En effet, on a $\Gamma(U_{i_0 \dots i_p}, (\tilde{M}(n))^\sim) = (M_{f_{i_0} \dots f_{i_p}})_n$ par définition (II, 3, 3, prop. 9); comme, sur un schéma affine, $\Gamma$ est un foncteur qui commute aux sommes directes (I, 1, 3, cor. 4 du th. 1), on a donc $\Gamma(U_{i_0 \dots i_p}, \mathcal{F}(*)) = M_{f_{i_0} \dots f_{i_p}}$. Le reste du raisonnement est alors le même que pour la prop. 1, tenant compte de ce que les $U_{i_0 \dots i_p}$ sont affines (§ 1, n$^\circ$3, Remarque suivant la prop. 3).

\textbf{Remarques.} — 1) Dans les conditions de la prop. 2, les foncteurs $\Gamma(U_{i_0 \dots i_p}, \tilde{M}(*))$ sont exacts en $M$; on en conclut, comme dans la Remarque du § 1, n$^\circ$2, que si $0 \to M' \to M \to M'' \to 0$ est une suite exacte, les diagrammes

\begin{tikzcd}
H^q(U, \tilde{M}'(*)) \arrow[r] \arrow[d] & H^{q+1}((\underline{f}), M') \arrow[d] \\
H^{q+1}(U, \tilde{M}'(*)) \arrow[r] & H^{q+2}((\underline{f}), M')
\end{tikzcd}

\newpage

%% ===== Page 352 =====

--- III-24 ---

2) La prop. 2 sera surtout intéressante lorsque $S$ sera une $A$-algèbre engendrée par un nombre fini d'éléments de degré 1, $A$ étant supposé noethérien ; en effet, lorsqu'il en est ainsi, tout faisceau quasi-cohérent est de la forme $\widetilde{M}$ (II, 3, 4, th. 1).

Nous allons appliquer la prop. 2 dans le cas où $S = A[T_0, \dots, T_n]$, où $A$ est un anneau quelconque, les $T_i$ des indéterminées, avec $M=S$, et $f_i = T_i$. On est donc essentiellement ramené à calculer $H^q((T), S)$ où $T = (T_i)_{0 \leq i \leq r}$.

Nous allons pour cela utiliser la proposition connue suivante, dont, pour être complet, nous donnerons une démonstration :

\emph{Proposition 3. — Soient $S$ un anneau, $f = (f_i)_{1 \leq i \leq r}$ une famille finie d'éléments de $S$, $M$ un $S$-module. Si, pour tout $i$, l'homothétie définie par $f_i$ dans $M / (f_1, \dots, f_{i-1})M$ est injective, on a $H^i(f, M) = 0$ pour $i \neq r$ et $H^r(f, M) \simeq M / (\sum_{i=1}^r f_i M) = M_f$.}

La seconde assertion est vraie quel que soit $M$ (§ 1, n° 1, formule (7)).

Pour démontrer la première, il suffira de prouver que $H_i(f, M) = 0$ pour tout $i > 0$. Raisonnons par récurrence sur $r$, le cas $r=0$ étant trivial. Posons $f' = (f_i)_{1 \leq i \leq r-1}$ ; il est clair qu'elle satisfait aux conditions de la prop. 3, donc, si $L = K_.(f', M)$, on a $H_i(L) = 0$ pour $i > 0$ et $H_0(L) = M_{f'}$. Or, posons pour abréger $K = K.(f_r) = K_0 \oplus K_1$, où $K_0 = K_1 = S$, $d_0$ étant la multiplication par $f_r$ ; on a par définition $K_.(f, M) = K \otimes L$. Mais on a le lemme suivant :

\emph{Lemme 1. — Soit $K$ un complexe de $S$-modules qui soit $A$-libre et nul en dimensions $\neq 0, 1$. Pour tout complexe $L$ de $A$-modules, on a une suite exacte}

$$0 \to H_0(K \otimes_A (L)) \to H_p(K \otimes_A L) \to H_p(K \otimes_{p-1} L) \to 0$$

C'est un cas particulier d'une suite exacte de termes de bas degré

\newpage

%% ===== Page 353 =====

- III-25 -

de la suite spectrale de Künneth ; mais on peut le démontrer directement de la façon suivante. Considérons $K_0$ et $K_1$ comme des complexes (nuls en dimensions $\neq 0$ et $\neq 1$ respectivement) ; on a alors une suite exacte de complexes
$$0 \to K_0 \otimes L \to K \otimes L \to K_1 \otimes L \to 0$$
à laquelle nous pouvons appliquer la suite exacte d'homologie
$$\dots \to H_{p+1}(K_1 \otimes L) \xrightarrow{\partial} H_p(K_0 \otimes L) \to H_p(K \otimes L) \to H_p(K_1 \otimes L) \to H_{p-1}(K_0 \otimes L) \to \dots$$
Mais il est évident que $H_p(K_0 \otimes L) = K_0 \otimes H_p(L)$ et $H_p(K_1 \otimes L) = K_1 \otimes H_{p-1}(L)$ pour tout $p$ ; en outre, on vérifie aussitôt que l'opérateur $\partial : K_1 \otimes H_p(L) \to K_0 \otimes H_p(L)$ n'est autre que $d_0 \otimes 1$ ; le lemme résulte donc de la définition de $H_p(K \otimes L)$ et de $H_p(K_1 \otimes L)$ et de la suite exacte précédente.

Ce lemme étant établi, la fin de la prop.3 est immédiate : l'hypothèse de récurrence et la suite exacte du lemme 1 donnent $H_p(K \otimes L) = 0$ pour $p \geq 2$ ; la même suite exacte donne aussi $H_1(K \otimes L) = 0$ dès que l'on aura montré que $H_1(K, H_0(L)) = 0$ ; mais par définition ce groupe n'est autre que le noyau de l'homothétie définie par $f_r$ dans $M_{r-1}$ ; et ce noyau est nul par hypothèse, ce qui achève la démonstration.

Proposition 4. — Si $S = A[T_0, \dots, T_r]$ ; on a, avec $T = (T_i)_{0 \leq i \leq r}$
$$(13) \quad H^i(\underline{T}, S) = 0 \quad \text{si } i \neq r+1$$
$$(14) \quad H^{r+1}(\underline{T}, S) = S/(\underline{T}^n)$$
Le $A$-module $H^{r+1}(\underline{T}, S)$ a donc une base sur $A$ formée des classes mod. $(\underline{T}^n)$ des monômes $T_0^{p_0} \dots T_r^{p_r}$ avec $p = (p_0, \dots, p_r)$, $0 \leq p_i < n$.
Ce est un corollaire immédiat de la prop.3, dont les hypothèses sont trivialement vérifiées.

Passons à la limite inductive sur $n$ ; les relations (13) donnent $H^i(\underline{T}, S) = 0$ pour $i \neq r+1$. Pour $i = r+1$, le système inductif est formé des $S/(\underline{T}^n)$, l'homomorphisme $\varphi_{nm} : S/(\underline{T}^n) \to S/(\underline{T}^m)$ ($0 \leq n \leq m$) étant la multiplication par $(T_0 \dots T_r)^{m-n}$. Si, pour $n > \sup(p_i)_{0 \leq i \leq r}$,

\newpage

%% ===== Page 354 =====

- III-26 -

on désigne par $\xi_{\underline{p}}^{(n)} = \xi_{p_0 \dots p_r}^{(n)}$ la classe de $T_0^{n-p_0} \dots T_r^{n-p_r} \pmod{(T^n)}$,
on a $\varphi_{nm} (\xi_{\underline{p}}^{(n)}) = \xi_{\underline{p}}^{(m)}$, et ces éléments ont donc même image canonique
$\xi_{\underline{p}} = \xi_{p_0 \dots p_r}$ dans la limite inductive $H^{r+1}(\mathbb{T}, S)$; en vertu de la dé-
finition du degré donnée au début de ce n°, le degré de $\xi_{\underline{p}}$ est donc
égal à $-|\underline{p}| = -(p_0 + \dots + p_r)$. Il est clair que les $\xi_{\underline{p}}^{(n)}$ pour $0 \le p_i \le n$
forment une base de $S/(T^n)$. On déduit donc de la prop. 4 :

Corollaire. - Avec les notations de la prop. 4, on a
$$H^i((\mathbb{T}), S) = 0 \text{ pour } i \neq r+1$$
et $H^{r+1}((\mathbb{T}), S)$ est un $A$-module libre dont une base est formée des
éléments $\xi_{p_0 \dots p_r}$ tels que $p_i > 0$ ($0 \le i \le r$).

De ce corollaire et de la prop. 2, on déduit :

Proposition 5. - Soient $A$ un anneau, $r$ un entier $\ge 0$, et soit $X = \mathbb{P}_A^r$ (II, 4, 5). Alors :
(i) On a $H^i(X, \mathcal{O}_X(*)) = 0$ pour $i \neq 0, r$.
(ii) L'homomorphisme canonique (II, 3, 3) $\alpha : S \to H^0(X, \mathcal{O}_X(*))$ est bijec-
tif.
(iii) $H^r(X, \mathcal{O}_X(*))$ est un $A$-module libre ayant une base formée d'élé-
ments $\xi_{p_0 \dots p_r}$, où $p_i > 0$ pour $0 \le i \le r$ ; $\xi_{p_0 \dots p_r}$ étant de degré $-|\underline{p}| =
-(p_0 + \dots + p_r)$.

Il suffit en effet de remarquer que $H^i((\mathbb{T}), S) = 0$ dans la suite exac-
te (12) appliquée à $S = A[T_0, \dots, T_r]$, et que la prop. 2 est applica-
ble à $U=X$, puisque $X$ est la réunion des $D_+(T_i)$ ; l'identification
de l'application $S \to H^0(X, \mathcal{O}_X(*))$ de la suite exacte (12) avec l'appli-
cation $\alpha$ définie en II, 3, 3 résulte de l'identification canonique
de $H^0(U_i, \mathcal{O}_X(*))$ et de $H^0(U_i, \mathcal{O}_X(*))$ et du fait que sur un schéma affine
le foncteur $\Gamma$ commute aux sommes directes infinies.

Corollaire 1. - Les seules valeurs de $(i, n)$ pour lesquelles on puis-
se avoir $H^i(X, \mathcal{O}_X(n)) \neq 0$ sont les suivantes : $i=0$ et $n \ge 0$, $i=r$ et
$n \le -(r+1)$.

\newpage

%% ===== Page 355 =====

- III-27 -

On notera que si $A \neq 0$, on a effectivement $H^i(X, \mathcal{O}_X(n)) \neq 0$ pour les valeurs indiquées de $(i,n)$ ; cela résulte de la prop. 5, puisque $S_n$ est alors $\neq 0$ pour tous les degrés.

Dans les applications qui seront faites dans ce chapitre, nous n'utiliserons en fait que le résultat moins précis :

Corollaire 2. — Les modules $H^i(X, \mathcal{O}_X(n))$ sont de type fini ; si $i > 0$, ils sont nuls pour $n \gg 0$.

Indiquons encore que nous donnerons plus tard le théorème "global" dont la forme locale est la prop. 5.

2. Le théorème fondamental des morphismes projectifs.

Théorème 1 (Serre). — Soient $Y$ un préschéma noethérien, $S$ une $\mathcal{O}_Y$-algèbre quasi-cohérente graduée en degrés positifs, et engendrée par le $\mathcal{O}_Y$-module $S_1$. On suppose qu'il existe un entier $r \ge 0$ tel que tout point de $Y$ admette un voisinage ouvert $U$ pour lequel il existe un homomorphisme surjectif $\mathcal{O}_U^{r+1} \to S|_U$. Soient $X = \operatorname{Proj}(S)$ le $Y$-préschéma projectif défini par $S$, $f : X \to Y$ son morphisme structural. Pour tout $\mathcal{O}_X$-module cohérent $\mathcal{F}$ :

(i) Les $R^i f_* (\mathcal{F})$ sont des faisceaux cohérents.

(ii) $R^i f_* (\mathcal{F}) = 0$ pour $i > r$.

(iii) Il existe un entier $N$ tel que $n \ge N$ et $i > 0$ entraîne $R^i f_* (\mathcal{F}(n)) = 0$.

(iv) Il existe un entier $N$ tel que $n \ge N$ entraîne que l'homomorphisme canonique $f^* (f_* (\mathcal{F}(n))) \to \mathcal{F}(n)$ est surjectif.

La forme locale de ce théorème est la suivante :

Corollaire 1. — Supposons remplies les conditions du th. 1, et supposons en outre que $Y = \operatorname{Spec}(A)$ soit un schéma affine. Alors, pour tout faisceau cohérent $\mathcal{F}$ sur $X$ :

(i) Les $A$-modules $H^i(X, \mathcal{F})$ sont de type fini.

(ii) $H^i(X, \mathcal{F}) = 0$ pour $i > r$.

\newpage

%% ===== Page 356 =====

\marginpar{Contenant compte de G, cor. du th. 4.9.2}

(iii) Il existe un entier $N$ tel que $n \ge N$ et $i > 0$ entraîne $H^i(X, \mathcal{F}(n)) = 0$.

(iv) Il existe un entier $N$ tel que pour $n \ge N$, $\mathcal{F}(n)$ soit engendré par ses sections au-dessus de $X$.

Comme $Y$ est quasi-compact, le cor. 1 entraîne le th. 1, car en recouvrant $Y$ par un nombre fini d'ouverts affines, il suffira de prendre pour $N$ dans (iii) et (iv) le plus grand des entiers correspondant (pour la même raison) à ces ouverts. On peut en outre supposer qu'il existe un homomorphisme surjectif $\mathcal{O}_Y^{r+1} \to S_1$. On sait alors (II, 4.4, cor. de la prop. 15) que $X$ s'identifie à un sous-préschéma fermé de $\mathbb{P}_A^r$ ; en outre, si $j$ est l'injection canonique $X \to \mathbb{P}_A^r$ (FAC, I, 2.17), $j_*(\mathcal{F})$ est cohérent, et $j_*(\mathcal{F}(n)) = (j_*(\mathcal{F}))(n)$ en vertu de la relation (14) de II, 4.5. On est donc ramené à prouver le cor. 1 dans le cas particulier où $X = \mathbb{P}_A^r$, $S = \tilde{S}$, où $S = A[T_0, \dots, T_r]$. Comme tout faisceau cohérent $\mathcal{F}$ sur $X$ est alors de la forme $\tilde{M}$, où $M$ est un $S$-module gradué de type fini, (ii) résulte de la formule (11) du n$^{\text{o}}$ 1, car $X$ est alors recouvert par les $D_+(T_i)$ en nombre $r+1$, et on sait (§ 1; n$^{\text{o}}$ 1) que $H^i((f), M) = 0$ lorsque $f$ a moins de $k$ éléments. Notons d'autre part que (iv) a déjà été démontré (II, 3.4, cor. 5 du th. 1). Nous allons démontrer simultanément (i) et (iii). Notons que ces assertions sont vraies si $\mathcal{F} = \mathcal{O}_X(m)$ (n$^{\text{o}}$ 1, cor. 2 de la prop. 5) ; elles le sont donc aussi lorsque $\mathcal{F}$ est somme directe d'un nombre fini de faisceaux $\mathcal{O}_X(m_j)$. D'autre part, (i) et (iii) sont vraies trivialement pour $i > r$ en vertu de (ii). Nous allons alors procéder par récurrence descendante sur $i$. On sait que $\mathcal{F}$ est isomorphe à un quotient d'une somme directe d'un nombre fini de faisceaux $\mathcal{O}_X(m_j)$ (II, 3.4, cor. 4 du th. 1), autrement dit, on a une suite exacte $0 \to R \to \mathcal{E} \to \mathcal{F} \to 0$, où $R$ est cohérent (puisque $\mathcal{F}$ et $\mathcal{E}$ le sont) et où $\mathcal{E}$ vérifie (i) et (iii). Comme $\mathcal{F} \to \mathcal{F}(n)$ est un foncteur exact, on a aussi une suite exacte

\newpage

%% ===== Page 357 =====

$$ 0 \to \underline{R}(n) \to \underline{G}(n) \to \underline{F}(n) \to 0 $$
pour tout $n \in \mathbb{Z}$. La suite exacte de cohomologie donne donc la suite exacte
$$ H^{i-1}(X, \underline{G}(n)) \to H^{i-1}(X, \underline{F}(n)) \to H^i(X, \underline{R}(n)) $$
Comme $\underline{G}(n)$ est somme directe des $\mathcal{O}_X(n+m_j)$, $H^{i-1}(X, \underline{G}(n))$ est de type fini, et il en est de même de $H^{i-1}(X, \underline{R}(n))$ par l'hypothèse de récurrence, donc $H^{i-1}(X, \underline{F}(n))$ est de type fini pour tout $n$, et en particulier pour $n=0$. D'autre part, il existe un $N$ tel que pour $n \ge N$ on ait $H^i(X, \underline{R}(n))=0$ (par l'hypothèse de récurrence), et $H^{i-1}(X, \underline{G}(n))=0$ puisque $\underline{G}$ vérifie (iii) ; on a donc aussi $H^{i-1}(X, \underline{F}(n))=0$, ce qui achève la démonstration.

Corollaire 2. — Sous les hypothèses du th. 1, soit $\underline{F} \to \underline{G} \to \underline{H}$ une suite exacte de faisceaux cohérents sur $X$. Il existe alors un entier $N$ tel que pour $n \ge N$, la suite
$$ f_*(\underline{F}(n)) \to f_*(\underline{G}(n)) \to f_*(\underline{H}(n)) $$
soit exacte.

Soit $\underline{F}'$ le noyau et $\underline{G}'$ le conoyau de $\underline{F} \to \underline{G}$ ; soit $\underline{G}''$ le noyau de $\underline{G} \to \underline{H}$, soit $\underline{H}'$ le conoyau de cet homomorphisme. Comme $\underline{F} \to \underline{F}(n)$ est un foncteur exact, il suffit de prouver que si $n$ est assez grand, chacune des suites
$$ 0 \to \underline{F}'(n) \to \underline{F}(n) \to \underline{G}'(n) \to 0, \quad 0 \to \underline{G}'(n) \to \underline{G}(n) \to \underline{H}'(n) \to 0, \quad 0 \to \underline{G}''(n) \to \underline{H}(n) \to \underline{H}'(n) \to 0 $$
est exacte ; autrement dit, on peut supposer que $0 \to \underline{F} \to \underline{G} \to \underline{H} \to 0$ est exacte. On a alors la suite exacte de cohomologie
$$ 0 \to f_*(\underline{F}(n)) \to f_*(\underline{G}(n)) \to f_*(\underline{H}(n)) \to R^1 f_*(\underline{F}(n)) \to \dots $$
et la conclusion résulte du th. 1, (iii).

Corollaire 3. — Sous les hypothèses du th. 1, soit $M$ un $S$-module vérifiant la condition $(T^0)$ (II, 4, 3). Alors il existe un entier $N$ tel que, pour $n \ge N$, l'homomorphisme canonique (II, 4, 2)
$$ \alpha_n : M_n \to f_*(\underline{M}(n)) $$

\newpage

%% ===== Page 358 =====

- III-30 -

est un isomorphisme. En d'autres termes, le noyau et le conoyau
de l'homomorphisme canonique
$$\alpha : M \to \Gamma_*( \tilde{M} )$$
vérifient la condition (TN') (ou, comme on dit encore, $\alpha$ est
un (TN')-isomorphisme).

Comme $Y$ est noethérien, la question est locale et on peut donc
supposer $Y = \operatorname{Spec}(A)$ et $M = \tilde{M}$, $A$ étant noethérien, $S$ un $A$-module de type fini et $M$ un $S$-module gradué de type fini. On a alors
une suite exacte
$$L' \to L \to M \to 0$$
où $L$ et $L'$ sont des modules gradués libres de type fini sur $S$. Supposons le résultat prouvé pour le cas où $M=S$ ; il est alors vrai pour $L$ et $L'$. Considérons le diagramme commutatif
\begin{tikzcd}
L'_n \to L_n \to M_n \to 0 \\
f_*(L'(n)) \to f_*(\tilde{L}(n)) \to f_*(\tilde{M}(n)) \to 0
\end{tikzcd}
La deuxième ligne est exacte d'après le cor. 2, dès que $n$ est assez
grand ; comme il en est de même de la première et que les deux premiers homomorphismes $\alpha_n$ sont des isomorphismes, il en est de même
du troisième.

Cela étant, si $S = A[T_0, \dots, T_r]$, le corollaire (pour $M=S$) résulte
de la prop. 5 (ii) du n° 1. Dans le cas général, $S$ s'identifie à un
anneau quotient d'un anneau $S' = A[T_0, \dots, T_r]$ et $X$ à un sous-proschéma
fermé de $X' = \mathbb{P}^r_A$. Si $j$ est l'injection canonique $X \to X'$, $j_*(\tilde{S}(n))$ est
un $\mathcal{O}_{X'}$-module vérifiant (TF'), donc l'homomorphisme canonique
$\alpha_n : S_n \to \Gamma(X', j_*(\tilde{S}(n)))$ est bijectif pour $n$ assez grand en vertu
de ce qui précède, d'où le corollaire, puisque $\Gamma(X', j_*(\tilde{S}(n))) =
= \Gamma(X, \tilde{S}(n))$.

Corollaire 4. - Pour tout faisceau cohérent $\mathcal{F}$ sur $X$, $\Gamma_*( \mathcal{F} )$ est

\newpage

%% ===== Page 359 =====

- III-31 -

un $S$-module vérifiant la condition $(TF')$.

On sait en effet que $\mathcal{F}$ est isomorphe à un faisceau de la forme $\widetilde{M}$, où $M$ est un $S$-module de type fini (II,4,3, cor.2 du th.1); il suffit alors d'appliquer le cor.3.

\emph{Scholie.} - On peut donc dire que la donnée d'un faisceau algébrique cohérent sur $X$ est équivalente à la donnée d'un $S$-module gradué quasi-cohérent satisfaisant à la condition $(TF)$, deux tels $S$-modules qui ne diffèrent que par un nombre fini de composantes homogènes correspondant au même faisceau cohérent sur $X$.

On peut aussi dire qu'il y a équivalence entre la catégorie des faisceaux cohérents sur $X$, et la catégorie quotient (T,I,11) de la catégorie des $S$-modules gradués quasi-cohérents vérifiant $(TF')$, par la sous-catégorie épaisse des $S$-modules gradués quasi-cohérents vérifiant $(TN')$.

Les propriétés (iii) et (iv) du th.1 se généralisent de façon immédiate quand $\mathcal{O}_X(1)$ est remplacé par un faisceau inversible ample (II, 4,6) :

\emph{Proposition 5.} - Soient $Y$ un préschéma noethérien, $X$ un $Y$-préschéma projectif, $f$ de morphisme structural $f$, $\mathcal{L}$ un faisceau inversible sur $X$, ample pour $f$, $\mathcal{F}$ un faisceau cohérent sur $X$. Il existe alors un entier $N$ tel que, pour $n \ge N$, on ait $R^i f_*(\mathcal{F} \otimes \mathcal{L}^{\otimes n}) = 0$ pour $i > 0$, et que l'homomorphisme canonique $f^*(f_*(\mathcal{F} \otimes \mathcal{L}^{\otimes n})) \to \mathcal{F} \otimes \mathcal{L}^{\otimes n}$ soit surjectif.

Par définition, il existe un entier $k$ tel que $\mathcal{L}^k$ soit très ample pour $f$ (par définition $X$ s'identifie à un sous-préschéma de $\mathbb{P}(\mathcal{E})$, où $\mathcal{E}$ est un $\mathcal{O}_Y$-module cohérent, et on a $\mathcal{L}^k \simeq j^*(\mathcal{O}_P(1))$, en désignant par $j$ l'injection canonique $X \to P$ ; en outre, comme $f$ est un morphisme projectif par hypothèse, $X$ est un sous-préschéma fermé de $P$ (II,4, prop. ). Cela étant,

\newpage

%% ===== Page 360 =====

- III-32 -

soit $p$ le morphisme structural $P \to Y$ ; pour tout entier $m$ tel que $0 \le m \le k-1$, il résulte du th.1 qu'il existe un entier $N_m$ tel que pour $r \ge N_m$, $R^i f_* (\mathcal{F} \otimes \mathcal{L}^{\otimes(m+kr)}) = R^i p_* (j_* (\mathcal{F} \otimes \mathcal{L}^{\otimes(m+kr)})) = 0$ pour $i > 0$, car on a $j_* (\mathcal{G} \otimes \mathcal{L}^{\otimes kr}) = j_*(\mathcal{G}) \otimes \mathcal{L}^{\otimes kr}$ pour tout $\mathcal{O}_X$-module cohérent $\mathcal{G}$, $j$ étant une immersion fermée, et que pour la même raison $R^i f_* (\mathcal{G}) = R^i p_* (j_* (\mathcal{G}))$ (II, 4, cor. du th. 4.9.2). De même, pour $r \ge N_m$, l'homomorphisme canonique $p^* (p_* (j_* (\mathcal{F} \otimes \mathcal{L}^{\otimes(m+kr)}))) \to j_* (\mathcal{F} \otimes \mathcal{L}^{\otimes(m+kr)})$ est surjectif ; comme le foncteur $j^*$ est exact ($j$ étant une immersion), et que $j^* (j_* (\mathcal{F} \otimes \mathcal{L}^{\otimes(m+kr)})) = \mathcal{F} \otimes \mathcal{L}^{\otimes(m+kr)}$, on en conclut que $f^* (f_* (\mathcal{F} \otimes \mathcal{L}^{\otimes(m+kr)})) \to \mathcal{F} \otimes \mathcal{L}^{\otimes(m+kr)}$ est surjectif. Il suffit alors de prendre pour entier $N$ le plus grand des entiers $m+kN_m$ pour répondre à la question.

Corollaire. - Sous les hypothèses de la prop. 6, soit $\mathcal{F} \to \mathcal{G} \to \mathcal{H}$ une suite exacte de faisceaux cohérents sur $X$. Il existe alors un entier $N$ tel que, pour $n \ge N$, la suite
$$f_* (\mathcal{F} \otimes \mathcal{L}^{\otimes n}) \to f_* (\mathcal{G} \otimes \mathcal{L}^{\otimes n}) \to f_* (\mathcal{H} \otimes \mathcal{L}^{\otimes n})$$
soit exacte.

Le raisonnement est exactement le même que dans le cor. 2 du th. 1.

Remarques. - 1) La seconde assertion de la prop. 6 a déjà été démontrée (II, 4, [unclear]) sous des conditions plus larges, savoir lorsque $X$ est seulement supposé quasi-projectif sur $Y$. Mais il faut noter que dans ce cas la propriété $R^i f_* (\mathcal{F} \otimes \mathcal{L}^{\otimes n}) = 0$ pour $i > 0$, à partir d'une certaine valeur de $n$, n'est plus nécessairement vérifiée. Soit en effet $A$ un anneau noethérien, et prenons $Y = \operatorname{Spec}(A)$, $X' = \operatorname{Spec}(A[T_0, \dots, T_r])$, et soit $g$ le morphisme structural $X' \to Y$ ; alors $\mathcal{O}_{X'}(n)$ est très ample (pour $g$) pour tout $n$ (II, 4, 8, [unclear]) ; pour tout sous-schéma $X$ de $X'$, $\mathcal{O}_X(n)$ est donc aussi très ample (pour la restriction de $g$ à $X$). Or, prenons pour $X$ la réunion des ouverts affines $D(T_i)$ de $X'$ ($0 \le i \le r$) ; il résulte du § 1, n°4, cor. de la prop. 3 et (II, 4, 6, [unclear])

\newpage

%% ===== Page 361 =====

- III-33 -

l'anneau $\Gamma(X, \mathcal{O}_X)$ s'identifie à l'intersection des anneaux de fractions $(A[T_0, \dots, T_r])_{T_i}$ pour $0 \le i \le r$ (I, 6, 2), donc à $A[T_0, \dots, T_r]$; on conclut alors du § 1, n°4, cor. de la prop. 3 et de la formule (7) du § 1, n°1, que l'on a $H^r(X, \mathcal{O}_X) = H^r(X, \mathcal{O}_X) = A$ pour tout $n$.

3. Application aux faisceaux gradués d'algèbres .

Soit $Y$ un préschéma ; nous dirons qu'une $\mathcal{O}_Y$-algèbre graduée (à degrés positifs) $\mathcal{S}$ est spéciale si $\mathcal{S}_1$ est un $\mathcal{O}_Y$-module de type fini et si, pour $n$ assez grand, $\mathcal{S}_n = \mathcal{S}_1^n$.

Proposition 8. — Soit $Y$ un préschéma noethérien, $f : X \to Y$ le morphisme structural;
(i) Soient $X$ un $Y$-préschéma projectif et $\mathcal{L}$ un $\mathcal{O}_X$-module très ample; alors $\mathcal{S} = \bigoplus_{n \ge 0} f_*(\mathcal{L}^{\otimes n})$ est une $\mathcal{O}_Y$-algèbre graduée spéciale, $X$ est isomorphe à $\operatorname{Proj}(\mathcal{S})$ et $\mathcal{L}$ est isomorphe à $\mathcal{O}_{X}(1)$.

(ii) Inversement, soit $\mathcal{S}'$ une $\mathcal{O}_Y$-algèbre graduée spéciale, et soient $X = \operatorname{Proj}(\mathcal{S}')$, $\mathcal{L} = \mathcal{O}_X(1)$; alors l'homomorphisme canonique $\mathcal{S}' \to \bigoplus_{n \ge 0} f_*(\mathcal{O}_X(n))$ est un $(TN^0)$-isomorphisme.

De façon imagée, on peut dire que les couples $(X, \mathcal{L})$, où $X$ est un $Y$-préschéma projectif et $\mathcal{L}$ un $\mathcal{O}_X$-module très ample, peuvent s'interpréter comme des $\mathcal{O}_Y$-algèbres graduées spéciales, modulo $(TN^0)$. L'assertion (ii) est un cas particulier du cor. 3 du th. 2 du n° 2; prouvons (i). Par définition d'un faisceau très ample (II, 4, 8, déf. 4) il existe une $\mathcal{O}_Y$-algèbre graduée spéciale $\mathcal{S}'$ et une immersion $r$ de $X$ dans $P = \operatorname{Proj}(\mathcal{S}')$ telle que $\mathcal{L} = r^*(\mathcal{O}_P(1))$; comme $X$ est projectif sur $Y$ et $P$ séparé sur $Y$, $r$ est un morphisme projectif (II, 4, 6, th. 3 (v)), donc une application fermée (II, 4, 6, th. 4), autrement dit $X$ s'identifie à un sous-préschéma fermé de $P$. Remplaçant $\mathcal{S}'$ par une algèbre quotient par un faisceau gradué d'idéaux convenable, on peut même supposer que $X=P$ (II, 4, 4, prop. 15 et cor. 1 de la prop. 15). On a alors $\mathcal{S} = f_*(\mathcal{O}_X)$, et le cor. 3 du th. 1 du n° 2 montre que l'homomorphisme canonique $\alpha_n : \mathcal{S}_n \to f_*(\mathcal{O}_X(n))$ est un isomorphisme.

\newpage

%% ===== Page 362 =====

\marginpar{Archives Grothendieck sept. 59}

- III-34 -

pour $n$ assez grand. D'autre part, chacun des $S_n$ est un $\mathcal{O}_Y$-module cohérent en vertu du th. 1 (1) du n° 2, donc $S$ est une $\mathcal{O}_Y$-algèbre spéciale (puisque $S$ en est une). En outre (II, 4.4, Remarque finale), le morphisme déduit de $\alpha: S \to S$ est alors un isomorphisme de $\operatorname{Proj}(S)$ sur $\operatorname{Proj}(S')=X$, qui transforme le faisceau $S(1)\sim$ en $\mathcal{O}_X(1)$, ce qui achève la démonstration.

Remarques. - Si on suppose seulement que $L$ est un $\mathcal{O}_X$-module ample, $L^{\otimes d}$ est très ample pour $d$ assez grand, et comme les préschémas $\operatorname{Proj}(S^{(d)})$ et $\operatorname{Proj}(S)$ sont canoniquement $Y$-isomorphes (II, 4.1, prop. 3 (1)), on voit que $X$ est encore isomorphe à $\operatorname{Proj}(S)$; en outre le th. 1 du n° 2 s'applique encore à $X$ en remplaçant $S$ par $S^{(d)}$ (mais bien entendu $\mathcal{O}_X(n)$ doit ici être prise isomorphe à $S^{(d)}(n)\sim$). Il serait intéressant de savoir si $S$ est alors une $\mathcal{O}_Y$-algèbre de type fini, et quels sont alors les énoncés qui doivent remplacer le th. 1 du n° 2 et ses corollaires, ainsi que la prop. 8 ci-dessus. D'une façon générale, le rôle joué par les faisceaux inversibles dans la théorie des morphismes projectifs n'est pas éclairci. Des progrès dans ce sens permettraient peut-être de résoudre des questions telles que la suivante : si $Y$ est un préschéma noethérien, $X$ un $Y$-préschéma séparé et de type fini, $X'$ un $Y$-préschéma projectif, $X' \to X$ un morphisme entier surjectif, est-il vrai que $X$ soit un $Y$-préschéma projectif ? On l'ignore même lorsque $Y=\operatorname{Spec}(k)$, où $k$ est un corps algébriquement clos.

Proposition 9. - Soient $Y$ un préschéma noethérien, $f: X \to Y$ un morphisme projectif, $L$ un $\mathcal{O}_X$-module inversible. Les conditions suivantes sont équivalentes :

(i) $L$ est ample.

(ii) Pour tout $\mathcal{O}_X$-module cohérent $F$, il existe un entier $n_0$ tel que pour $n \ge n_0$, on ait $R^qf_*(F \otimes L^n) = 0$ pour tout $q > 0$.

\newpage

%% ===== Page 363 =====

- III-35 -

(iii) Pour tout faisceau cohérent $J$ d'idéaux dans $\mathcal{O}_X$, il existe un entier $n_0$ tel que, pour $n \ge n_0$, on ait $R^1f_*(J \otimes L^{\otimes n}) = 0$.

On a déjà démontré que (i) entraîne (ii) (n°2, prop.7), et (ii) implique trivialement (iii). Pour démontrer que (iii) entraîne (i), on peut se borner au cas où $Y$ est affine (II, 4,9, cor. de la prop.34); nous allons appliquer le critère de II, 4,9, prop.35 (iv). Soit donc $J$ un faisceau cohérent d'idéaux dans $\mathcal{O}_X$, et soit $x$ un point fermé de $X$; désignons par $K$ le faisceau cohérent d'anneaux sur $X$ dont la fibre est nulle aux points $\neq x$, et la fibre en $x$ égale à $k(x)$. Le raisonnement du critère de Serre (§ 1, n°5, th.2) montre que l'on a une suite exacte $0 \to J \to J \otimes K \to 0$, où $J'$ est un faisceau cohérent d'idéaux, et comme $L$ est localement libre, on en déduit une suite exacte
$$0 \to J' \otimes L^{\otimes n} \to J \otimes L^{\otimes n} \to J \otimes K \otimes L^{\otimes n} \to 0$$
pour tout $n$. L'hypothèse (iii) entraîne que pour $n \ge n_0$, $H^1(X, J' \otimes L^{\otimes n}) = 0$, et on en conclut, par la suite exacte de cohomologie, que $\Gamma(X, J \otimes L^{\otimes n}) \to \Gamma(J \otimes K \otimes L^{\otimes n})$ est surjective. Donc (loc. cit.) il existe un nombre fini de sections $h_i$ de $J$ au-dessus de $X$ ($1 \le i \le n$) et un voisinage ouvert $U(x)$ de $x$ tels que les restrictions des $h_i$ à $U(x)$ engendrent $J \otimes L^{\otimes n} | U(x)$. Comme $X$ est noethérien, il y a un nombre fini de points fermés $x_i$ ($1 \le i \le m$) tels que les $U(x_i)$ forment un recouvrement fermé de $X$ (I, 4,1, lemme 1). On en conclut aussitôt que $J$ est engendré par un nombre fini de ses sections au-dessus de $X$; $a$ fortiori le critère de II, 4,9, prop.35 (iv est vérifié.

\marginpar{$J \otimes L^{\otimes n}$}
\marginpar{$J \otimes L^{\otimes n}$}

\newpage

%% ===== Page 364 =====

-III-36 -

\textbf{Proposition 10.} — Soient $Y$ un préschéma noethérien, $f: X \to Y$ un morphisme projectif, $L$ un $\mathcal{O}_X$-module inversible, $g: X' \to X$ un morphisme fini de $X$-finant. Pour que $L$ soit ample, il faut et il suffit que son image réciproque $L' = g^*(L) = L \otimes_{\mathcal{O}_X} \mathcal{O}_{X'}$, le soit.

On sait déjà que la condition est nécessaire (II, 4, 8, cor. 3 de la prop. 31). Comme $g_{\text{red}}$ est fini (II, 2, 7, prop. 15) et $f_{\text{red}}$ projectif (II, 4, 6, th. 3), on peut supposer $X$ et $X'$ réduits en vertu do II, 4, 9, cor. 2 de la prop. 35. Par le raisonnement de récurrence habituel dans les espaces noethériens, on peut supposer que pour toute partie fermée $Z \subset X$, en désignant par $j$ l'injection dans $X$ du sous-préschéma réduit dont $Z$ est l'espace sous-jacent, $j^*(L)$ est ample. Nous allons en déduire tout d'abord que si un $\mathcal{O}_X$-module cohérent $\mathcal{F}$ a son support contenu dans une telle partie $Z$, il satisfait à la condition (ii) de la prop. 9. En effet, si $J$ est le faisceau cohérent d'idéaux définissant le sous-préschéma réduit $Z$, on a alors $J^n\mathcal{F}=0$ pour un entier $n > 0$ (II, 1, 4, prop. 14); filtrant $\mathcal{F}$ par les sous-faisceaux $J^k\mathcal{F}$ et utilisant la suite exacte de cohomologie, on est ramené au cas où $J\mathcal{F}=0$, autrement dit au cas où $\mathcal{F}$ est de la forme $j_*(\mathcal{G})$, où $\mathcal{G}$ est un $\mathcal{O}_Z$-module cohérent. Comme l'injection $j$ est affine, on sait (§ 1, n$^\circ$3, cor. 3 du th. 1) que
$$ R^q(f \circ j)_*(\mathcal{G} \otimes j^*(L^{\otimes m})) = R^qf_*(\mathcal{F} \otimes L^{\otimes m}) $$
pour tout $q > 0$ et tout $m$ (II, 2, 4, cor. de la prop. 9); mais comme par hypothèse $j^*(L)$ est ample, on a $R^q(f \circ j)_*(\mathcal{G} \otimes j^*(L^{\otimes m})) = 0$ pour $m$ assez grand, ce qui établit notre assertion.

Cela étant, posons $A = \mathcal{O}_X$, $B = g_*(\mathcal{O}_{X'})$; il existe deux entiers $k, m$ et un homomorphisme de $\mathcal{O}_X$-modules $A^k \to B \otimes L^{\otimes m}$ ayant la propriété suivante : il existe un ensemble ouvert non vide $V \subset X$ tel que sur $V$ soit un isomorphisme. La démonstration de ce fait suit pas celle de l'assertion analogue dans la prop. 6 du § 1, n$^\circ$5; il faut seulement utiliser le fait que si

\newpage

%% ===== Page 365 =====

- III-37 -

$U$ est un voisinage ouvert affine de $x \in X$, $x_j$ ($1 \le j \le n$) les points de la fibre (discrète) $g^{-1}(x)$, $s$ une section de $\mathcal{O}_{X'}$ au-dessus de $g^{-1}(U)$, il existe pour chaque $j$ une section $f_j$ de $\underline{L}'^m$ au-dessus de $X'$ telle que $x_j \in X'_j \subset g^{-1}(U)$ et que $s \otimes f_j$ se prolonge en une section de $\mathcal{O}_X \otimes \underline{L}'^m$ au-dessus de $X'$. La possibilité de réaliser la première condition (en prenant $m$ assez grand) résulte de ce que $\underline{L}'$ est ample et du critère (ii) de EGA II, 4.9, prop. 35 : on considère les composantes irréductibles (en nombre fini) de $g^{-1}(U)$, et on applique le critère cité à $x_j$ et au point générique de chacune de ces composantes. Quant à la possibilité de réaliser la seconde condition elle résulte de EGA II, 1.4, th. 2.

Cela étant, supposons d'abord que $\underline{F}$ soit un $\underline{B}$-module cohérent sur $X$ et prouvons alors que $R^q f_* (\underline{F} \otimes \underline{L}'^n) = 0$ pour tout $q > 0$ dès que $n$ est assez grand. Comme $g$ est affine, il existe un $\mathcal{O}_X$-module cohérent $\underline{G}$ tel que $\underline{F} = g_*(\underline{G})$ (EGA II, 2.4, prop. 9), et on a $R^q f_* (\underline{F} \otimes \underline{L}'^n) = R^q (f \circ g)_* (\underline{G} \otimes \underline{L}'^n)$ pour tout $q > 0$ et tout $n$ (§ 1, n° 3, cor. 3 du th. 1) ; l'hypothèse que $\underline{L}'$ est ample entraîne donc la conclusion pour $n$ assez grand en vertu de la prop. 9.

Abordons maintenant le cas général ; l'homomorphisme $u$ défini plus haut donne un homomorphisme

$$v : \underline{\operatorname{Hom}}_{\underline{A}}(\underline{B} \otimes \underline{L}'^m, \underline{F}) \to \underline{\operatorname{Hom}}_{\underline{A}}(\underline{A}^k, \underline{F}) = \underline{F}^k$$

dont la restriction à $V$ est un isomorphisme. Mais on a
$\underline{\operatorname{Hom}}_{\underline{A}}(\underline{B} \otimes \underline{L}'^m, \underline{F}) = \underline{\operatorname{Hom}}(\underline{L}'^m, \underline{\operatorname{Hom}}(\underline{B}, \underline{F})) = \underline{\operatorname{Hom}}(\underline{B}, \underline{F}) \otimes \underline{L}'^{(-m)}$
(EGA II, 1.3, prop. 11), et par suite, pour tout $n > 0$, on déduit de $v$ un homomorphisme

$$w : \underline{F}' \otimes \underline{L}'^n \to (\underline{F} \otimes \underline{L}'^{(n+m)})^k$$

où on a posé $\underline{F}' = \underline{\operatorname{Hom}}_{\underline{A}}(\underline{B}, \underline{F})$ ; en outre la restriction de $w$ à $V$ est un isomorphisme, autrement dit les supports de $\operatorname{Ker} w$ et de $\operatorname{Coker} w$

\newpage

%% ===== Page 366 =====

- III-38 -

sont contenus dans une partie fermée de $X$ distincte de $X$. Comme $F'$ est un $B$-module cohérent, on a $R^q f_* (F' \otimes L^n) = 0$ pour $q > 0$ et $n$ assez grand d'après ce qui a été vu plus haut. On conclut comme dans la démonstration de la prop. 6 du $\S 1, n^\circ 5$, en appliquant deux fois la suite exacte de cohomologie.

Corollaire. — Soit $(X_i)_{1 \leq i \leq n}$ une famille finie de sous-préschémas fermés de $X$, telle que $X$ soit réunion des espaces de base $X_i$. Pour que $L$ soit ample, il faut et il suffit que pour chaque $i$, $g_i^*(L)$ soit ample, en désignant par $g_i$ l'injection canonique $X_i \to X$.

Il suffit de considérer le préschéma $X'$ somme des $X_i$, et le morphisme surjectif $g : X' \to X$ dont la restriction à chaque $X_i$ est $g_i$ ; ce morphisme étant fini ($\S 1, n^\circ 5$, cor. 1 de la prop. 6), on peut lui appliquer la prop. 10.

4. Préschémas éclatés.

Soit $Y$ un préschéma, $R(Y)$ le faisceau des pseudo-fonctions sur $Y$ (1,5,3) ; nous dirons pour abréger qu'un sous-faisceau de $R(Y)$ est un faisceau d'idéaux fractionnaires sur $Y$. On dira qu'un faisceau d'idéaux (ou d'idéaux fractionnaires) $J$ sur $Y$ est localement principal si pour tout $y \in Y$, il existe un voisinage ouvert $U$ de $y$ tel que $J|_U$ soit engendré par une de ses sections au-dessus de $U$ ; si $Y$ est intègre et $J_y \neq 0$ pour tout $y \in Y$, un tel faisceau $J$ est inversible.

Définition 1. — Soit $J$ un faisceau quasi-cohérent d'idéaux (ou d'idéaux fractionnaires) sur $Y$, et soit $S = \bigoplus_{n \geq 0} J^n$ (où on pose $J^0 = \mathcal{O}_Y$), qui est une $\mathcal{O}_Y$-Algèbre graduée quasi-cohérente. On dit que le $Y$-préschéma $X = \operatorname{Proj}(S)$ est obtenu en faisant éclater le faisceau d'idéaux (resp. d'idéaux fractionnaires) $J$.

Lorsque $J$ est un faisceau d'idéaux quasi-cohérent, qui définit un

\newpage

%% ===== Page 367 =====

--- III-39 ---

sous-préschéma fermé $Y'$ de $Y$, on dit encore que $\operatorname{Proj}(S)$ est obtenu en faisant éclater le sous-préschéma $Y'$. Lorsque $J$ est de type fini, $X = \operatorname{Proj}(S)$ est un $Y$-schéma projectif.

Si on remplace $J$ par $J^d$ ($d > 0$), ou par $JL$, où $L$ est un sous-faisceau inversible de $\mathcal{O}_Y$ (resp. de $\mathcal{R}(Y)$ si $J$ est un faisceau d'idéaux fractionnaires), $\operatorname{Proj}(S)$ est remplacé par un $Y$-préschéma isomorphe (II, 4,1, prop. 3). En particulier, si $J$ est inversible, $X$ s'identifie à $Y$.

Proposition 11. — Soient $Y$ un préschéma intègre, $J$ un faisceau quasi-cohérent non nul d'idéaux entiers ou fractionnaires sur $Y$. Alors le préschéma $X$ obtenu en faisant éclater $J$ est intègre, et le morphisme structural $f : X \to Y$ est birationnel.

La première assertion résulte de ce que $S = \bigoplus_{n \ge 0} J^n$ est alors intègre (II, 4, 1, prop. 4). Soient $x, y$ les points génériques de $X$ et $Y$ ; on a $f(x) = y$ et il faut prouver que $k(x)$ est de rang $1$ sur $k(y)$. Or $x$ est aussi le point générique de la fibre $f^{-1}(y)$ ; si $\psi$ est le morphisme canonique $\xi \to Y$, où $\xi = \operatorname{Spec}(k(y))$, le préschéma $f^{-1}(y)$ s'identifie à $\operatorname{Proj}(S')$, où $S' = \psi^*(S)$ (II, 4, 4, prop. 14). Mais il est clair que $S' = \bigoplus_{n \ge 0} (J_y)^n$, et comme $J_y \neq 0$ est un idéal de $\mathcal{O}_y = k(y)$, on a $J_y = k(y)$ et $\operatorname{Proj}(S')$ s'identifie par suite à $\xi$, d'où la seconde assertion.

Proposition 12. — Soient $Y$ un préschéma intègre noethérien, $X$ un préschéma intègre, $f : X \to Y$ un morphisme projectif birationnel. Il existe alors un faisceau cohérent $J \subset \mathcal{R}(Y)$ d'idéaux fractionnaires tel que $X$ soit isomorphe au préschéma obtenu en faisant éclater $J$. En outre, si $U$ est un ouvert de $Y$ tel que la restriction de $f$ à $f^{-1}(U)$ soit un isomorphisme $f^{-1}(U) \to U$ (I, 4, 3, prop. 16), on peut supposer que $J|_U$ est localement principal.

Si $L$ est un faisceau très ample sur $X$, $X$ s'identifie à $\operatorname{Proj}(S)$, où $S = \bigoplus_{n \ge 0} f_* (L^n)$ (nº 3, prop. 8) ; il est clair en outre que $f_*(L)|_U$

\newpage

%% ===== Page 368 =====

- III-40 -

est inversible, et il en est de même de chacun des $f_*(\mathcal{L}^n)|_U$ ($n \ge 0$). Comme les $f_*(\mathcal{L}^n)$ sont cohérents ($n^\circ 2, th. 1$), la propriété précédente entraîne que les $f_*(\mathcal{L}^n)$ sont sans torsion, donc il en est de même du $\mathcal{O}_Y$-module $\mathcal{S} \otimes_{\mathcal{O}_Y} \mathcal{R}(Y)$ ($I, 5, 3, cor. 3$ de la $prop. 8$). Ce dernier est connu lorsqu'on connaît sa restriction à un ouvert non vide, par exemple à un ouvert non vide $U' \subset U$ tel que $\mathcal{L}|_{f^{-1}(U')}$ soit isomorphe à $\mathcal{O}_X|_{f^{-1}(U')}$. On voit donc que $\mathcal{S} \otimes_{\mathcal{O}_Y} \mathcal{R}(Y)$ est une $\mathcal{O}_Y$-Algèbre isomorphe à $\mathcal{R}(Y)[T]$, où $T$ est une indéterminée, et $\mathcal{S}$ s'identifie à la sous-$\mathcal{O}_Y$-Algèbre graduée engendrée par l'image de $f_*(\mathcal{L})$ dans $\mathcal{S} \otimes_{\mathcal{O}_Y} \mathcal{R}(Y)$ identifié à $\mathcal{R}(Y)[T]$ ; il est clair que cette image est de la forme $J \cdot T$, où $J$ est un sous-faisceau cohérent de $\mathcal{R}(Y)$, donc $\mathcal{S}$ est isomorphe à $\bigoplus_{n \ge 0} J^n$, ce qui achève la démonstration.

Corollaire 1. - Sous les hypothèses de la $prop. 12$, supposons en outre que, pour tout sous-faisceau cohérent $J$ de $\mathcal{R}(Y)$, il existe un $\mathcal{O}_Y$-Module inversible $\mathcal{L}$ tel que $\mathcal{T}(Y, \mathcal{L} \otimes_{\mathcal{O}_Y} \text{Hom}(J, \mathcal{O}_Y)) \neq 0$ ; alors, dans l'énoncé de la $prop. 12$, on peut supposer que $J$ est un sous-faisceau de $\mathcal{O}_Y$. Cette condition supplémentaire est toujours vérifiée si $Y$ est un schéma quasi-projectif sur un anneau $A$. ($II, 1, 3, prop. 11$)

En effet, on a $\mathcal{L} \otimes \text{Hom}(J, \mathcal{O}_Y) = \text{Hom}(J^{-1}, \text{Hom}(J, \mathcal{O}_Y)) = \text{Hom}(J \otimes \mathcal{L}^{-1}, \mathcal{O}_Y)$ ; l'hypothèse signifie donc qu'il existe un homomorphisme non nul $u$ de $J \otimes \mathcal{L}^{-1}$ dans $\mathcal{O}_Y$. Comme pour tout $y \in Y$, $(J \otimes \mathcal{L}^{-1})_y$ s'identifie à un sous-$\mathcal{O}_{Y,y}$-module du corps des fractions $\mathcal{K}(Y)_y$, $u_y$ est nécessairement injectif, donc $u$ est un isomorphisme de $J \otimes \mathcal{L}^{-1}$ sur un sous-faisceau $J'$ de $\mathcal{O}_Y$. On peut par suite remplacer $J$ par $J'$ sans changer $\text{Proj}(\mathcal{S})$ à un isomorphisme près ($II, 4, 1, prop. 3$). La dernière assertion du corollaire résulte du

Lemme 2. - Si $Y$ est un schéma quasi-projectif sur un anneau $A$, pour tout faisceau cohérent $\mathcal{F} \neq 0$ sur $Y$, il existe un $\mathcal{O}_Y$-module inversible $\mathcal{L}$ quasi-

\newpage

%% ===== Page 369 =====

- III-41 -

$L$ tel que $T(Y, L \otimes F) \neq 0$.

En effet, on peut identifier $Y$ à un sous-préschéma induit sur un ouvert partout dense d'un $A$-schéma projectif $Z$ (II, 4, 7, prop. 24), et $F$ est alors induit sur $Y$ par un $\mathcal{O}_Z$-module quasi-cohérent $\underline{G}$ (II, 1, 2, th. 1). On peut donc se ramener au cas où $Y$ est projectif sur $A$ ; en outre, comme $F$ est limite inductive de ses sous-faisceaux cohérents (II, 1, 2, cor. du th. 1) on peut se borner au cas où $F$ est cohérent et $\neq 0$. Mais alors $F(n) = F \otimes \mathcal{O}_Y(n)$ est $\neq 0$ pour tout $n$, et il est engendré par ses sections au-dessus de $Y$ lorsque $n$ est assez grand ($n^\circ 2$, cor. 1 du th. 1), d'où le lemme.

En particulier :

Corollaire 2. — Soient $X$ et $Y$ deux schémas projectifs intègres sur un corps $k$, et soit $f: X \to Y$ un morphisme birationnel. Alors $X$ est isomorphe à un $Y$-préschéma obtenu en faisant éclater un sous-schéma fermé $Y'$ (non nécessairement réduit) de $Y$.

En effet, $f$ est projectif (II, 4, 6, th. 3 (v)), et comme $Y$ est projectif sur un corps, la condition supplémentaire du cor. 1 est vérifiée ; il suffit alors de considérer le sous-schéma fermé $Y'$ de $Y$ défini par le faisceau d'idéaux $J$ du cor. 1.

Remarque. — Au chap. IV, en étudiant la notion de diviseur, nous verrons que si, dans l'énoncé de la prop. 12, on suppose $Y$ non singulier, alors $X$ peut se déduire de $Y$ en faisant éclater un sous-schéma $Y'$ de $Y$ dont l'espace sous-jacent est contenu dans $Y - U$.

\marginpar{/ que les $\mathcal{O}_y$ ($y \in Y$) sont factoriels (ce qui est le cas par exemple si ...)}

5. Caractéristique d'Euler-Poincaré et polynôme de Hilbert.

Soient $A$ un anneau artinien, $S$ une $A$-algèbre graduée spéciale, $X = \operatorname{Proj}(S)$ son spectre premier homogène. Si $F$ est un $\mathcal{O}_X$-module cohérent, les $H^i(X, F)$ ($i \ge 0$) sont des $A$-modules de type fini ($n^\circ 2$, cor. 1 du th. 1), et par suite de longueur finie puisque $A$ est artinien.

\newpage

%% ===== Page 370 =====

- III-42 -

En outre, on sait que $H^i(X, \mathcal{F}) = 0$ sauf pour un nombre fini de valeurs de $i$ (loc.cit.) ; le nombre entier
$$(15) \quad \chi_A(\mathcal{F}) = \sum_{i=0}^{\infty} (-1)^i \operatorname{long}(H^i(X, \mathcal{F}))$$
est donc défini pour tout $\mathcal{O}_X$-module cohérent $\mathcal{F}$ ; on l'appelle la caractéristique d'Euler-Poincaré de $\mathcal{F}$ (par rapport à l'anneau $A$). Pour $\mathcal{F} = \mathcal{O}_X$, on dit que $\chi_A(\mathcal{O}_X)$ est le genre arithmétique de $X$.

Proposition 13. - Soit $0 \to \mathcal{F}' \to \mathcal{F} \to \mathcal{F}'' \to 0$ une suite exacte de $\mathcal{O}_X$-modules cohérents. On a alors
$$(16) \quad \chi_A(\mathcal{F}) = \chi_A(\mathcal{F}') + \chi_A(\mathcal{F}'').$$

On a en effet, par la suite exacte de cohomologie, une suite exacte
$$0 \to H^0(X, \mathcal{F}') \to H^0(X, \mathcal{F}) \to H^0(X, \mathcal{F}'') \to H^1(X, \mathcal{F}') \to \dots$$
$$\dots \to H^r(X, \mathcal{F}') \to H^r(X, \mathcal{F}) \to H^r(X, \mathcal{F}'') \to 0$$
les modules de cohomologie d'indice $> r$ étant nuls pour $\mathcal{F}, \mathcal{F}'$ et $\mathcal{F}''$. Or on sait que dans une suite exacte de modules de longueur finie, ayant des $0$ aux deux extrémités, la somme alternée des longueurs est nulle ; appliquant ce résultat, on trouve immédiatement (16).

Théorème 2. - (i) La fonction $n \to \chi_A(\mathcal{F}(n))$ est un polynôme en $n$.
(ii) Pour $n$ assez grand, on a $\chi_A(\mathcal{F}(n)) = \operatorname{long} \Gamma(X, \mathcal{F}(n))$ ; si $\mathcal{F} = \widetilde{M}$ (où $M$ est un $S$-module gradué de type fini), on a $\chi_A(\mathcal{F}(n)) = \operatorname{long}(M_n)$ pour $n$ assez grand.
(iii) Le terme de plus haut degré de $\chi_A(\mathcal{F}(n))$ a un coefficient $> 0$.
(iv) Le degré de $\chi_A(\mathcal{F}(n))$ est égal à la dimension du support de $\mathcal{F}$.

Soient $s_1, \dots, s_m$ les points de $Z = \operatorname{Spec}(A)$ ; on sait que $Z$ est somme des schémas $Z_i = \operatorname{Spec}(A_i)$, où les $A_i$ sont les composantes locales de $A$, et l'espace de base de $Z_i$ est $\{s_i\}$ (I, 4, 1, prop. 4) ; si $f: X \to Z$ est le morphisme structural, $X$ est alors somme des préschémas

\newpage

%% ===== Page 371 =====

\begin{flushright}
- III-43 -
\end{flushright}

mas $X_i$ induits sur les ouverts $X_i = f^{-1}(z_i)$. Si $\mathcal{F}_i$ est le faisceau égal à $\mathcal{F} \mid X_i$ dans $X_i$, à $0$ dans les $X_j$ tels que $j \neq i$, $\mathcal{F}$ est la somme directe des $\mathcal{F}_i$, donc (prop. 13)

$$(17) \qquad \chi_A(\mathcal{F}(n)) = \sum_{i=1}^m \chi_A(\mathcal{F}_i(n))$$

Par ailleurs, il est clair que les $A_j$ d'indice $j \neq i$ annulent les modules $H^q(X, \mathcal{F}_i)$, et par suite ces derniers peuvent être considérés comme des $A_i$-modules et leur longueur est la même en tant que $A$-modules ou $A_i$-modules, d'où

$$(18) \qquad \chi_A(\mathcal{F}_i(n)) = \chi_{A_i}(\mathcal{F}_i(n)) \quad (1 \le i \le m)$$

Les relations (17) et (18) montrent que pour prouver (i), on peut supposer que $A$ est un anneau artinien local. Soit $\mathfrak{m}$ son idéal maximal ; il existe un entier $s$ tel que $\mathfrak{m}^s = 0$, et $\mathcal{F}(n)$ admet donc la filtration finie $\mathcal{F}(n) \supset \mathfrak{m}\mathcal{F}(n) \supset \dots \supset \mathfrak{m}^{s-1}\mathcal{F}(n) \supset 0$. Par récurrence, on déduit de la prop. 13 que

$$(19) \qquad \chi_A(\mathcal{F}(n)) = \sum_{k=0}^{s-1} \chi_A(\mathfrak{m}^k\mathcal{F}(n)/\mathfrak{m}^{k+1}\mathcal{F}(n))$$

Chacun des faisceaux $\mathcal{F}^k(n) = \mathfrak{m}^k\mathcal{F}(n)/\mathfrak{m}^{k+1}\mathcal{F}(n)$ (où $\mathcal{F}_k = \mathfrak{m}^k\mathcal{F}/\mathfrak{m}^{k+1}\mathcal{F}$) est annulé par $\mathfrak{m}$, donc peut être considéré comme un faisceau sur $X^0 = \operatorname{Proj}(S^0)$, où $S^0 = S/\mathfrak{m}S$ (II, 4.5, prop. 21) ; d'ailleurs, si $K = A/\mathfrak{m}$ est le corps résiduel de $A$, $S^0$ peut être considéré comme une $K$-algèbre graduée, et comme $H^q(X^0, \mathcal{F}_k(n)) = H^q(X, \mathcal{F}_k(n))$ (G, II 4.9.1), on a $\chi_K(\mathcal{F}_k(n)) = \chi_A(\mathcal{F}_k(n))$. Pour démontrer (i), on peut donc se ramener au cas où $A$ est un corps ; enfin, $X$ peut alors être considéré comme un sous-préschéma fermé de $\mathbb{P}^r_A$ pour un $r$ convenable (II, 4.6, prop. 21), et comme $H^q(X, \mathcal{F}(n)) = H^q(\mathbb{P}^r, \mathcal{F}(n))$, on voit finalement qu'on peut supposer que $S = K[T_0, \dots, T_r]$. On a alors $\mathcal{F} = \tilde{M}$, où $M$ est un $S$-module gradué de type fini (II, 3.4, cor. 1 du th. 1) ; il existe donc une résolution finie de $M$ par des $S$-modules gradués libres de type fini

$$0 \to L_q \to L_{q-1} \to \dots \to L_1 \to M \to 0$$

\newpage

%% ===== Page 372 =====

(
) et comme $\widetilde{M}$ est un foncteur additif exact (II,3,3, prop.10), on a aussi une suite exacte
$$0 \to \widetilde{L}_q \to \widetilde{L}_{q-1} \to \dots \to \widetilde{L}_1 \to \widetilde{M} \to 0$$
et par suite aussi
$$0 \to \widetilde{L}_q(n) \to \widetilde{L}_{q-1}(n) \to \dots \to \widetilde{L}_1(n) \to \widetilde{M}(n) \to 0$$
est exacte ; en appliquant la prop.13 par récurrence sur $q$, on en déduit
$$\chi_A(\widetilde{M}(n)) = \sum_{j=1}^q (-1)^{j+1} \chi_A(\widetilde{L}_j(n))$$
et pour prouver (i), on est donc ramené au cas où $M$ est libre et gradué de type fini, donc au cas où $M=S(h)$ pour un $h \in \mathbb{Z}$. Comme alors $\widetilde{M}(n) = (\widetilde{M(n)})^{\sim} = (S(n+h))^{\sim}$ (II,3,3, cor.2 de la prop.14), on est finalement ramené à prouver que $\chi_A(\mathcal{O}_X(n))$ est égal à un polynôme; mais en fait, pour $n \ge 0$ on a alors $\chi_A(\mathcal{O}_X(n)) = \operatorname{long} H^0(X, \mathcal{O}_X(n)) = \operatorname{long}_A S_n = \binom{n+r}{r}$ en vertu de la prop.5 du n°1, et on a de même $\chi_A(\mathcal{O}_X(n)) = (-1)^r \operatorname{long}_A H^r(X, \mathcal{O}_X(n))$, et si $n = -r-h$, la dimension de $H^r(X, \mathcal{O}_X(n))$ est le nombre des suites $(p_i)_{0 \le i \le r}$ d'entiers $p_i \ge 0$ tels que $\sum p_i = r+h$, ou encore le nombre des suites d'entiers $q_i \ge 0$ ($0 \le i \le r$) telles que $\sum q_i = h-1$, autrement dit $\binom{h+r-1}{r}$; comme $\binom{h+r-1}{r} = (-1)^r \binom{n+r}{r}$, on a encore $\chi_A(\mathcal{O}_X(n)) = \binom{n+r}{r}$; enfin, cette relation est aussi vérifiée pour $-r \le n < 0$, car alors $\binom{n+r}{r}=0$ et $H^i(X, \mathcal{O}_X(n)) = 0$ pour tout $i \ge 0$ (loc. cit.). La démonstration de (i) est ainsi achevée.

Comme on a $H^i(X, \mathcal{F}(n)) = 0$ pour tout $i > 0$ dès que $n$ est assez grand (n°2, cor.1 du th.1), on a bien $\chi_A(\mathcal{F}(n)) = \operatorname{long} H^0(X, \mathcal{F}(n)) = \operatorname{long} \Gamma(X, \mathcal{F}(n))$ pour ces valeurs de $n$, d'où la première assertion de (ii); la seconde résulte du n°2, cor.3 du th.2. Comme on a donc $\chi_A(\mathcal{F}(n)) \ge 0$ pour $n$ assez grand, (iii) résulte de (ii). Reste

\newpage

%% ===== Page 373 =====

- III-45 -

donc à démontrer (iv). Avec les notations du début de la démonstration, le support de $\mathcal{F}$ est réunion des supports des $\mathcal{F}_i$, donc $\dim(\operatorname{supp}(\mathcal{F})) = \sup_i(\dim(\operatorname{supp}(\mathcal{F}_i)))$ (I, 7, 1, prop. 1) et tenant compte de (17), on est ramené à prouver (iv) pour chacun des $\mathcal{F}_i$, donc on peut supposer l'anneau $A$ local. De même, le support de $\mathcal{F}$ (lorsque $A$ est local) est réunion des supports des $\mathcal{F}_i$, et en utilisant (19) et (iii), on est ramené au cas où $A$ est un corps. Enfin, comme ci-dessus, en considérant $X$ sous-préschéma fermé de $\mathbb{P}^r_A$, on se ramène au cas où $X = \mathbb{P}^r_A$. Considérons alors la catégorie $\underline{K}'$ des $\mathcal{O}_X$-modules cohérents $\mathcal{F}$ tels que $\deg(\chi_A(\mathcal{F}(n))) = \dim(\operatorname{supp}(\mathcal{F}))$; cette catégorie n'est peut-être pas exacte, mais la somme directe de deux faisceaux de $\underline{K}'$ appartient à $\underline{K}'$ (prop. 13), la relation $\mathcal{F}' \in \underline{K}'$ entraîne $\mathcal{F} \in \underline{K}'$ en vertu de la prop. 13, puisque $\operatorname{supp}(\mathcal{F}^q) = \operatorname{supp}(\mathcal{F})$; enfin, si $u : \mathcal{F} \to \mathcal{G}$ est un homomorphisme de $\mathcal{O}_X$-modules tel que $\operatorname{Ker} u$ et $\operatorname{Coker} u$ appartiennent à $\underline{K}'$ et que les supports de $\operatorname{Ker} u$ et de $\operatorname{Coker} u$ aient une dimension strictement inférieure à celle du support de $\mathcal{G}$, on en conclut que $\mathcal{F} \in \underline{K}'$. En effet, on a d'abord
$$\chi_A(\mathcal{G}(n)) = \chi_A((\operatorname{Im} u)(n)) + \chi_A((\operatorname{Coker} u)(n))$$
et comme le degré de $\chi_A((\operatorname{Coker} u)(n))$ est par hypothèse strictement inférieur au degré de $\chi_A(\mathcal{G}(n))$, les polynômes $\chi_A(\mathcal{G}(n))$ et $\chi_A((\operatorname{Im} u)(n))$ ont même degré ; d'autre part, comme $\operatorname{supp}(\mathcal{G})$ est réunion des supports de $\operatorname{Im} u$ et de $\operatorname{Coker} u$, l'hypothèse entraîne que $\operatorname{supp}(\mathcal{G})$ et $\operatorname{supp}(\operatorname{Im} u)$ ont même dimension (I, 7, 1, prop. 1), d'où $\operatorname{Im} u \in \underline{K}'$. On a ensuite
$$\chi_A(\mathcal{F}(n)) = \chi_A((\operatorname{Ker} u)(n)) + \chi_A((\operatorname{Im} u)(n))$$
et comme le degré de $\chi_A((\operatorname{Ker} u)(n))$ est strictement inférieur à celui de $\chi_A((\operatorname{Im} u)(n))$ d'après ce qu'on vient de voir, les degrés de $\chi_A(\mathcal{F}(n))$ et de $\chi_A((\operatorname{Im} u)(n))$ sont égaux, d'où on conclut comme

\newpage

%% ===== Page 374 =====

- III-46 -

précédemment que $\underline{F} \in \underline{K}'$. On prouve de même que $\underline{G} \in \underline{K}'$ lorsqu'on suppose que $\underline{F} \in \underline{K}'$ et que les supports de $\ker u$ et de $\operatorname{Coker} u$ aient une dimension strictement inférieure à celle de $\operatorname{supp}(\underline{F})$. La Remarque suivant le corollaire du lemme de dévissage (II, 1, 5) montre que ce corollaire s'applique, et pour montrer que $\underline{K}'$ est la catégorie de tous les $\mathcal{O}_X$-modules cohérents, on est ramené à prouver que pour tout sous-préschéma fermé intègre $Y$ de $X = \mathbb{P}^r_A$, le faisceau $\mathcal{O}_Y$ (identifié à un $\mathcal{O}_X$-module) appartient à $\underline{K}'$. On peut identifier le sous-préschéma $Y$ à $\operatorname{Proj}(S')$, où $S' = S/I$, $I$ étant un idéal gradué de $S_+$ (II, 3, 5, prop. 21) et il résulte de (ii) que pour $n$ grand $\chi_A(\mathcal{O}_Y(n)) = \dim_A S'_n$ ; on peut d'ailleurs remplacer $S'$ par $S'^{(d)}$, où $d > 0$, et supposer par suite que $S'_n = (S'_1)^n$ pour tout $n$ (II, 4, 1, cor. 3 de la prop. 3). D'où $\dim_A S'_n = \dim_A (S'_n / S'^{n+1}_+) - \dim_A (S'_n / S'^n_+)$.

Or, $n \to \dim_A (S'_n / S'^{n+1}_+)$ est, pour $n$ grand, un polynôme de degré égal à la dimension de l'anneau local $B = S'_{\mathfrak{m}}$. Comme $Y$ étant équidimensionnel (I, 7, 2, prop. 7), $\dim(Y')$ est égal à la dimension de chacun des anneaux locaux de ses points fermés et en particulier à $\dim(B)$, qui correspond au "sommet" de $Y$ (qui n'est autre que le cône projetant de $Y$ (II, 3, 6)) ; on voit donc que $\deg(\chi_A(\mathcal{O}_Y(n))) = \dim Y' - 1$. D'autre part, $\dim(Y') = \dim(Y' - \{c\})$, et on a un morphisme canonique surjectif de type fini $g : Y' - \{c\} \to Y$ dont toutes les fibres sont de dimension 1 (II, 3, 6, prop. 22) ; donc $\dim Y' = \dim Y + 1$ (I, 7, 3, cor. 4 du th. 1), ce qui achève de démontrer le th. 2.

Corollaire. - Soient $A$ un anneau artinien, $S$ une $A$-algèbre graduée spéciale, $X = \operatorname{Proj}(S)$. Les conditions suivantes sont équivalentes :
a) $X$ est artinien ; [unclear: a')] $X$ est fini sur $A$.
b) Il existe une $A$-algèbre $A'$, finie sur $A$, et un homomorphisme $S \to A'[T]$ de $A$-algèbres graduées, qui est un isomorphisme modulo (TN).
c) Il existe une constante telle que $\operatorname{long}(S_n) = a$ pour $n$ assez grand.

\newpage

%% ===== Page 375 =====

- III-47 -

d) Il existe une constante $a$ telle que $\operatorname{long}(S_n) \le a$ pour $n$ assez grand.

Comme un module de type fini sur un anneau artinien est artinien, il est clair que si $X$ est fini sur $A$, $X$ est artinien. Inversement, si $X$ est artinien, $A' = \Gamma(X, \mathcal{O}_X)$ est artinien et l'espace sous-jacent $X$ est fini et discret (I, 4, 1, prop. 4), donc $\mathcal{O}_X(1)$ est isomorphe à $\mathcal{O}_X$ ; par suite $\bigoplus_{n \ge 0} \Gamma(X, \mathcal{O}_X(n))$ est isomorphe à $A'[T]$. Appliquant à $S$ le cor. 3 du th. 1, $n^\circ 2$, on voit que a) implique b) ; en outre $\operatorname{Proj}(S) = \operatorname{Spec}(A')$ (II, 4, 1, cor. 3 de la prop. 2), donc $X$ est fini sur $A$. Il est clair que b) entraîne c) et c) implique trivialement d). Enfin, le th. 2 (ii) montre que d) implique que le support de $\mathcal{O}_X$ (c'est-à-dire $X$) est de dimension 0 ; les composantes irréductibles de $X$ sont donc des points, ce qui implique que les points de $X$ sont fermés, et comme $X$ est noethérien, $X$ est artinien (I, 4, 1, prop. 4), donc d) implique a).

Remarque. - La constante $a$ dans le d) du corollaire n'est autre que la longueur de $A' = \Gamma(X, \mathcal{O}_X)$ ; si $A$ est un corps, c'est le rang du schéma algébrique fini $X$. On verra que lorsque $\mathcal{F}$ est un faisceau algébrique cohérent sur $\mathbb{P}^r_k$ ($k$ étant un corps), le terme dominant $c_d n^d$ du polynôme $\chi_k(\mathcal{F}(n))$ est tel que puisse s'interpréter comme le "degré" du "cycle dominant" de $\mathcal{F}$.

\newpage

%% ===== Page 376 =====

\marginpar{Archives Grothendieck sept 59}

- III-48 -

§ 3. Etude cohomologique des morphismes propres.

1. Le théorème de finitude.

Bien que le théorème suivant soit contenu dans le théorème fondamental des morphismes propres (n° , th. ), sa démonstration directe est beaucoup plus simple, et nous la donnons à part tout d'abord.

Théorème 1 (théorème de finitude). - Soient $Y$ un schéma localement noethérien, $f: X \to Y$ un morphisme propre. Pour tout $\mathcal{O}_X$-module cohérent $\mathcal{F}$, les $\mathcal{O}_Y$-modules $R^q f_*(\mathcal{F})$ sont cohérents ($q \ge 0$).

La question étant locale sur $Y$, on peut supposer $Y$ affine, donc noethérien, et par suite $X$ noethérien. Les $\mathcal{O}_X$-modules cohérents satisfaisant à la conclusion du th. 1 forment une sous-classe exacte $\mathcal{K}'$ de la catégorie $\mathcal{K}$ des faisceaux cohérents sur $X$. En effet, soit $0 \to \mathcal{F}' \to \mathcal{F} \to \mathcal{F}'' \to 0$ une suite exacte de $\mathcal{O}_X$-modules cohérents ; supposons par exemple que $\mathcal{F}'$ et $\mathcal{F}''$ appartiennent à $\mathcal{K}'$ ; la suite exacte de cohomologie montre que la suite

$$R^{q-1} f_*(\mathcal{F}'') \to R^q f_*(\mathcal{F}') \to R^q f_*(\mathcal{F}) \to R^q f_*(\mathcal{F}'') \to R^{q+1} f_*(\mathcal{F}')$$

est exacte ; comme $R^{q-1} f_*(\mathcal{F}'')$ et $R^q f_*(\mathcal{F}')$ sont cohérents, il en est de même de l'image $I_q$ de l'homomorphisme $R^{q-1} f_*(\mathcal{F}'') \to R^q f_*(\mathcal{F}')$; de même, comme $R^q f_*(\mathcal{F}'')$ et $R^{q+1} f_*(\mathcal{F}')$ sont cohérents, il en est de même du noyau $N_{q+1}$ de l'homomorphisme $R^q f_*(\mathcal{F}'') \to R^{q+1} f_*(\mathcal{F}')$; enfin, comme la suite $0 \to I_q \to R^q f_*(\mathcal{F}) \to N_{q+1} \to 0$ est exacte et $I_q$ et $N_{q+1}$ cohérents, $R^q f_*(\mathcal{F})$ est cohérent. On montre de même que lorsque $\mathcal{F}$ et $\mathcal{F}'$ (resp. $\mathcal{F}$ et $\mathcal{F}''$) appartiennent à $\mathcal{K}'$, $\mathcal{F}''$ (resp. $\mathcal{F}'$) appartient à $\mathcal{K}'$. En outre, tout facteur direct d'un faisceau $\mathcal{F}$ appartenant à $\mathcal{K}'$ appartient aussi à $\mathcal{K}'$ : en effet, $R^q f_*(\mathcal{F}')$ est alors facteur direct de $R^q f_*(\mathcal{F})$ (G. II, 4.4.4), donc est de type fini, et par suite cohérent (FAG, I, 2.12, prop. 3). En vertu du corollaire du lemme.

\newpage

%% ===== Page 377 =====

- III-49 -

me de dévissage (II, 1, 5, cor. du th. 3), on est ramené à prouver que lorsque $X$ est irréductible, il existe un $\mathcal{O}_X$-module cohérent de support $X$ appartenant à $K'$ : en effet, si ce point est établi, on pourra l'appliquer à toute partie fermée irréductible $Y$ de $X$, car si $\mathcal{G}$ est un $\mathcal{O}_Y$-module cohérent et $j$ l'injection canonique $Y \to X$, on a $R^q(f \circ j)_*(\mathcal{G}) = R^q f_* (j_* (\mathcal{G}))$ (G, II, 4.9.1), et le corollaire du lemme de dévissage s'applique.

En vertu du lemme de Chow (II, 5, 2, cor. 1 du th. 1), il existe un préschéma irréductible $X'$ et un morphisme projectif et surjectif birationnel $g: X' \to X$ tels que $f \circ g : X' \to Y$ soit projectif. Comme $X$ est noethérien, $X'$ est isomorphe à un préschéma $\operatorname{Proj}(S)$, où $S$ est une $\mathcal{O}_Y$-algèbre quasi-cohérente graduée à degrés positifs, engendrée par $S_1$, où $S_1$ est un $\mathcal{O}_Y$-module cohérent (II, 4, 6, prop. 21 (ii)). Le th. 1 du § 2, n° 2 est applicable à $g: X' \to X$, et il existe donc un entier $n > 0$ tel que $\mathcal{F} = g_* (\mathcal{O}_{X'}(n))$ soit cohérent et $R^q g_* (\mathcal{O}_{X'}(n)) = 0$ pour tout $q > 0$ ; en outre, comme $g$ est birationnel, le support de $\mathcal{F}$ contient un ouvert partout dense de $X$, et comme ce support est fermé, il est identique à $X$. D'autre part, comme $f \circ g$ est projectif et $Y$ noethérien, le même raisonnement (II, 4, 6, prop. 21 (ii) et § 2, n° 2, th. 1) montre que les $R^p(f \circ g)_* (\mathcal{O}_{X'}(n))$ sont cohérents. Cela étant, $R^p(f \circ g)_* (\mathcal{O}_{X'}(n))$ est l'aboutissement d'une suite spectrale dont le terme $E_2$ est donné par $E_2^{pq} = R^p f_* (R^q g_* (\mathcal{O}_{X'}(n)))$ (suite spectrale de Leray ; § 0, n° ). D'après ce qui précède, on a $E_2^{pq} = 0$ pour $q > 0$ ; on sait alors que $E_2^{p0} = R^p f_* (\mathcal{F})$ est isomorphe à $R^p(f \circ g)_* (\mathcal{O}_{X'}(n))$ (G, I, 4.34.1), ce qui achève la démonstration.

Corollaire 1. - Soit $Y$ un préschéma localement noethérien. Pour tout morphisme propre $f: X \to Y$, l'image directe par $f$ de

\newpage

%% ===== Page 378 =====

- III-50 -

tout faisceau cohérent sur $X$ est un faisceau cohérent sur $Y$.

Corollaire 2. — Soient $A$ un anneau noethérien, $X$ un schéma propre sur $A$, $\mathcal{F}$ un $\mathcal{O}_X$-module cohérent. Alors les $H^p(X, \mathcal{F})$ sont des $A$-modules de type fini.

En particulier, si $X$ est un schéma algébrique propre sur un corps $k$, alors, pour tout $\mathcal{O}_X$-module cohérent $\mathcal{F}$, les $H^p(X, \mathcal{F})$ sont des espaces vectoriels de dimension finie.

Remarque. — Dans ce dernier cas, on sait que l'on a $H^p(X, \mathcal{F}) = 0$ pour $p > \dim X$; les $H^p(X, \mathcal{F})$ sont donc nuls sauf pour un nombre fini de valeurs de $p$. Nous verrons plus loin que ce résultat est encore vrai sous les conditions générales du th. 1, lorsque $Y$ est noethérien.

2. Complété d'un faisceau cohérent le long d'une partie fermée.

Soient $X$ un préschéma noethérien, $X'$ une partie fermée de l'espace de base de $X$ ; désignons par $\Phi$ l'ensemble des faisceaux d'idéaux $J$ dans $\mathcal{O}_X$ tels que le support de $\mathcal{O}_X/J$ soit $X'$. L'ensemble $\Phi$ n'est pas vide (I, 2.4, prop. 9 et I, 3.1, prop. 1) ; nous l'ordonnons par la relation $\supset$.

Lemme 1. — L'ensemble ordonné $\Phi$ est filtrant ; pour tout $J \in \Phi$, l'ensemble des puissances $J^n$ ($n > 0$) est cofinal à $\Phi$.

En effet, si $J_1$ et $J_2$ appartiennent à $\Phi$ et si on pose $J = J_1 \cap J_2$, $J$ est quasi-cohérent et on a $J(x) = J_1(x) \cap J_2(x)$ donc $J(x) = \mathcal{O}_X(x)$ pour $x \notin X'$ et $J(x) \neq \mathcal{O}_X(x)$ pour $x \in X'$, ce qui prouve que $J \in \Phi$. D'autre part, si $J$ et $J'$ sont dans $\Phi$, il existe un entier $n > 0$ tel que $J^n \subset J'$ (II, 1.4, prop. 14), ce qui signifie que $J^n \subset J'$.

Soit maintenant $\mathcal{F}$ un $\mathcal{O}_X$-module cohérent ; pour tout $J \in \Phi$, $\mathcal{F} \otimes_{\mathcal{O}_X} (\mathcal{O}_X/J)$ est un faisceau cohérent sur $X$, de support contenu dans $X'$ ; on l'identifiera le plus souvent à sa restriction à $X'$. Lorsque $J$ parcourt $\Phi$, ces faisceaux forment un système projectif de

\newpage

%% ===== Page 379 =====

- III-51 -

\textbf{Définition 1.} — On appelle complété de $\mathcal{F}$ le long de $X'$, et on désigne par $\widehat{\mathcal{F}}_{/X'}$ ou par $\widehat{\mathcal{F}}$ (lorsqu'aucune confusion n'est possible) le faisceau $\varprojlim_{\mathcal{J} \in \Phi} (\mathcal{F} \otimes_{\mathcal{O}_X} (\mathcal{O}_X/\mathcal{J}))$ sur $X$. Ses sections s'appellent les sections formelles de $\mathcal{F}$ le long de $X'$ et sa cohomologie $H^q(X', \widehat{\mathcal{F}}_{/X'})$ la cohomologie formelle de $\mathcal{F}$ le long de $X'$.

Le support de $\widehat{\mathcal{F}}_{/X'}$ est contenu dans $X'$, et on identifie d'ordinaire ce faisceau à sa restriction à $X'$; rappelons que par définition, on a, pour tout ouvert $U \subset X$ :

$$(1) \quad \Gamma(U, \widehat{\mathcal{F}}_{/X'}) = \Gamma(U \cap X', \widehat{\mathcal{F}}_{/X'}) = \varprojlim_{\mathcal{J} \in \Phi} \Gamma(U \cap X', \mathcal{F} \otimes_{\mathcal{O}_X} \mathcal{O}_X/\mathcal{J}) = \varprojlim_{n} \Gamma(U \cap X', \mathcal{F} \otimes_{\mathcal{O}_X} \mathcal{O}_X/\mathcal{J}^n)$$

pour tout $\mathcal{J} \in \Phi$ (cf. lemme 1). Il est clair que pour tout ouvert $U \subset X$ on a $(\widehat{\mathcal{F}}_{/X'})|_U = \widehat{(\mathcal{F}|_{U \cap X'})}$.

Il est immédiat, par passage à la limite projective, que $(\mathcal{O}_X/\widehat{\mathcal{O}_X})$ est un faisceau d'anneaux, et que $\widehat{\mathcal{F}}_{/X'}$ peut être considéré comme un $(\mathcal{O}_X/\widehat{\mathcal{O}_X})$-module (les sections de $(\mathcal{O}_X/\widehat{\mathcal{O}_X})$ ne sont autres que les "fonctions holomorphes" de Zariski le long de $X'$). Les homomorphismes $\mathcal{O}_X \to \mathcal{O}_X/\mathcal{J}$ ($\mathcal{J} \in \Phi$) donnent d'ailleurs par passage à la limite projective un homomorphisme canonique $\mathcal{O}_X \to \mathcal{O}_X/\widehat{\mathcal{O}_X}$, ce qui permet de considérer $\widehat{\mathcal{F}}_{/X'}$ comme un $\mathcal{O}_X$-Module, mais ce faisceau ne sera pas nécessairement quasi-cohérent. D'autre part, $\widehat{\mathcal{F}}_{/X'}$, bien qu'étant un faisceau sur l'espace topologique $X'$, ne peut être en général considéré comme un $\mathcal{O}_{X'}$-module, pour une quelconque des structures de sous-préschéma fermé de $X$ ayant pour espace de base $X'$ (car $\mathcal{O}_X/\mathcal{J}^n$ n'est pas muni d'une structure de module sur $\mathcal{O}_X/\mathcal{J}$ en général). On peut seulement dire que $X'$, muni de $(\mathcal{O}_X/\widehat{\mathcal{O}_X})$, est un espace annelé sur lequel les $\widehat{\mathcal{F}}_{/X'}$ sont des faisceaux algébriques (au § de ce chapitre, nous étudierons sous le nom de "préschémas formels" une catégorie d'espaces annelés généralisant à la fois

\newpage

%% ===== Page 380 =====

- III-52 -

les espaces annelés du type précédent et les préschémas ordinaires.
Il est clair que $\widehat{\mathcal{F}}_{/X'}$ est un foncteur additif sur $\mathcal{F}$, de la catégorie des faisceaux algébriques cohérents sur $X$, dans la catégorie des faisceaux algébriques sur l'espace annelé $(X', (\mathcal{O}_X)_{/X'})$.

Proposition 1. — Le foncteur $\widehat{\mathcal{F}}_{/X'}$ est exact.

La question étant locale, la prop. 1 est équivalente à son
Corollaire 1. — Si $0 \to \mathcal{F}' \to \mathcal{F} \to \mathcal{F}'' \to 0$ est une suite exacte de $\mathcal{O}_X$-modules cohérents, et $U$ un ouvert affine de $X$, alors la suite
$$0 \to \Gamma(U, \widehat{\mathcal{F}}'_{/X'}) \to \Gamma(U, \widehat{\mathcal{F}}_{/X'}) \to \Gamma(U, \widehat{\mathcal{F}}''_{/X'}) \to 0$$
est exacte.

On a en effet alors $U = \operatorname{Spec}(A)$, où $A$ est un anneau noethérien, et $\mathcal{F}|_U = \tilde{M}$, $\mathcal{F}'|_U = \tilde{M}'$, $\mathcal{F}''|_U = \tilde{M}''$, où $M, M', M''$ sont trois $A$-modules (tels que la suite $0 \to M' \to M \to M'' \to 0$ soit exacte (I, 1.3, cor. 3 du th. 1). Soit $I$ un idéal de $A$ tel que $V(X' \cap U) = I$; on a alors
$\Gamma(U, \mathcal{F} \otimes_{\mathcal{O}_X} \mathcal{O}_X/J) = M \otimes_A (A/I)$ en désignant par $J$ un faisceau d'idéaux tel que $\mathcal{O}_X/J$ ait pour support $X'$ et que $J|_U = I$ (I, 1.3, cor. 6 du th. 1). Il résulte donc de la formule (1) que $\Gamma(U, \widehat{\mathcal{F}}_{/X'}) = \hat{M}$, complété de $M$ pour la topologie $I$-adique, et de même $\Gamma(U, \widehat{\mathcal{F}}'_{/X'}) = \hat{M}'$, $\Gamma(U, \widehat{\mathcal{F}}''_{/X'}) = \hat{M}''$; le cor. 1 résulte alors de ce que, lorsque $A$ est noethérien, le foncteur $\hat{M}$ est exact sur la catégorie des $A$-modules de type fini (Bourbaki, Alg. comm.).

Corollaire 2. — $H^n(X', \widehat{\mathcal{F}}_{/X'})$ est un foncteur cohomologique en $\mathcal{F}$, à valeurs dans la catégorie des modules sur l'anneau $A(X/X') = \Gamma((\mathcal{O}_X)_{/X'})$ des sections formelles de $\mathcal{O}_X$ le long de $X'$.

Les homomorphismes canoniques $\mathcal{F} \otimes \mathcal{O}_X/J \to \mathcal{F} \otimes \mathcal{O}_X/J'$ pour $J, J'$ dans $\Phi$ et $J \subset J'$ définissent canoniquement des homomorphismes
$$H^n(X', \mathcal{F} \otimes \mathcal{O}_X/J) \to H^n(X', \mathcal{F} \otimes \mathcal{O}_X/J') \quad (n \ge 0)$$
qui font de $(H^n(X', \mathcal{F} \otimes \mathcal{O}_X/J))_{J \in \Phi}$ un système projectif de groupes abéliens; en vertu de l'existence des homomorphismes canoniques

\newpage

%% ===== Page 381 =====

- III-53 -
$(\mathcal{O}_X)_J \to (\mathcal{O}_X)_{J'}$ pour tout $J \in \Phi$, il s'agit même d'un système projectif
de $A(X/X')$-modules. En outre, les $A(X/X')$-homomorphismes
$\mathcal{F}_{/X'} \to \mathcal{F} \otimes \mathcal{O}_X/J$ ($J \in \Phi$) définissent un
système projectif de $A(X/X')$-homomorphismes
$$H^i(X', \mathcal{F}_{/X'}) \to H^i(X', \mathcal{F} \otimes \mathcal{O}_X/J) \quad (i \ge 0)$$
d'où, en passant à la limite projective, un homomorphisme cano-
nique de $A(X/X')$-modules
$$(2) \quad \psi_i : H^i(X', \mathcal{F}_{/X'}) \to \varprojlim H^i(X', \mathcal{F} \otimes \mathcal{O}_X/J) = \varprojlim_n H^i(X', \mathcal{F} \otimes \mathcal{O}_X/J^n) \, .$$
Nous verrons plus loin des cas divers où les $\psi_i$ sont des isomorphis-
mes.

Proposition 2. — Le système des $\hat{H}^i(X/X', \mathcal{F}) = \varprojlim H^i(X', \mathcal{F} \otimes \mathcal{O}_X/J^n)$ est
un $\partial$-foncteur en $\mathcal{F}$, à valeurs dans la catégorie des $A(X/X')$-modules
et le système des $\psi_i$ est un homomorphisme de $\partial$-foncteurs.

En effet, soit $0 \to \mathcal{F}' \to \mathcal{F} \to \mathcal{F}'' \to 0$ une suite exacte de $\mathcal{O}_X$-modules cohé-
rents ; nous devons définir un $A(X/X')$-homomorphisme
$$\partial : \hat{H}^i(X/X', \mathcal{F}'') \to \hat{H}^{i+1}(X/X', \mathcal{F}')$$
ayant les propriétés usuelles (T, II, 2.1). Or, pour tout $n$, on a
la suite exacte
$$0 \to \mathcal{F}' / (J^n\mathcal{F} \cap \mathcal{F}') = (\mathcal{F}' + J^n\mathcal{F}) / J^n\mathcal{F} \to \mathcal{F} / J^n\mathcal{F} \to (\mathcal{F}' + J^n\mathcal{F}) / J^n\mathcal{F}' \to 0$$
d'où des homomorphismes
$$H^i(X', \mathcal{F}'' / J^n\mathcal{F}'') \to H^{i+1}(X', \mathcal{F}' / (J^n\mathcal{F} \cap \mathcal{F}'))$$
qui, ainsi que leur limite projective, ont les propriétés caractéri-
sant les opérations $\partial$. Or, on a $\mathcal{F}'' / J^n\mathcal{F}'' = \mathcal{F}'' \otimes \mathcal{O}_X/J^n$, donc
$\varprojlim H^i(X', \mathcal{F}'' / J^n\mathcal{F}'') = \hat{H}^i(X/X', \mathcal{F}'')$ par définition. Pour achever la dé-
monstration, on voit que tout revient à prouver que la filtration
sur $\mathcal{F}'$ définie par les $J^n\mathcal{F} \cap \mathcal{F}'$ est cofinale à la filtration définie
par les $J^n\mathcal{F}'$. Cela résulte immédiatement du fait que l'on a
$$(3) \quad J^n\mathcal{F} \cap \mathcal{F}' \supset J^{n-h}(\mathcal{F}' \cap J^h\mathcal{F})$$
dès que $n \ge h$, pour un entier $h$ convenable. En effet, si $X$ est af...

\newpage

%% ===== Page 382 =====

- III-54 -

fine, ce résultat n'est autre que le lemme d'Artin-Rees (Bourbaki,
Alg. Comm., ) car si $X=\operatorname{Spec}(A)$, $J=\tilde{I}$, $F=\tilde{M}$, $F'=\tilde{M'}$ (I idéal de
l'anneau noethérien $A$, $M$ $A$-module de type fini, $M'$ sous-module de
$M$), on a $J^n \cap F' = (I^n M \cap M') = (I^{n-h}(I^h M \cap M')) = J^{n-h}(J^h \cap F')$ (I, 1, 3,
cor. du th. 1). Dans le cas général, il y a un recouvrement fini
de $X$ par des ouverts affines $U_i$, et pour chaque $i$, on a une relation analogue à (3) en remplaçant $J, F, F'$ par leurs restrictions à $U_i$
et $h$ par un entier $h_i$; il suffit alors de prendre pour $h$ le plus
grand des $h_i$ pour obtenir (3). Enfin, le résultat précédent montre
aussi que les $\psi_i$ constituent un homomorphisme de $\partial$-foncteurs,
car les diagrammes
$$\begin{tikzcd}
H^i(X', F''/X') \arrow[r, "\partial"] \arrow[d] & H^{i+1}(X', F'/X') \arrow[d] \\
H^i(X', F''/J^nF'') \arrow[r, "\partial"] & H^{i+1}(X', F'/(J^n \cap F'))
\end{tikzcd}$$
étant commutatifs, il en est de même, par passage à la limite projective, des diagrammes
$$\begin{tikzcd}
H^i(X', F''/X') \arrow[r, "\partial"] \arrow[d, "\psi_i"] & H^{i+1}(X', F'/X') \arrow[d, "\psi_{i+1}"] \\
H^i(X/X', F'') \arrow[r, "\partial"] & H^{i+1}(X/X', F')
\end{tikzcd}$$
Il n'est pas certain que le $\partial$-foncteur ainsi défini soit exact,
puisque une limite projective de suites exactes n'est pas nécessairement exacte; on saura seulement \emph{a posteriori} que ce $\partial$-foncteur est
exact lorsqu'on saura que les $\psi_i$ sont des isomorphismes pour tout
$\mathcal{O}_X$-module cohérent.

Proposition 3. - Si $F$ et $G$ sont deux $\mathcal{O}_X$-modules cohérents, il existe
des isomorphismes canoniques fonctoriels
$$(F \otimes G)/X \cong (F/X) \otimes (\mathcal{O}_X) (G/X)$$
et
$$(\mathcal{H}om_{\mathcal{O}_X}(F, G))/X \cong \mathcal{H}om_{(\mathcal{O}_X)/X} (F/X, G/X)$$
[unclear: et] [unclear: ...], on est ramené au cas où $X=\operatorname{Spec}(A)$,
$A$ étant noethérien, $F=\tilde{M}$, $G=\tilde{N}$, $M, N$ étant des $A$-modules de type fini.

\marginpar{Il suffit de définir ces isomorphismes sur les sections au-dessus de $V$ ouvert dans $X$}

\newpage

%% ===== Page 383 =====

- III-54 -

fine, ce résultat n'est autre que le lemme d'Artin-Rees (Bourbaki, Alg. Comm., ) car si $X=\operatorname{Spec}(A)$, $J=\tilde{I}$, $F=\tilde{M}$, $F'=\tilde{M}'$ (I idéal de l'anneau noethérien $A$, $M$ $A$-module de type fini, $M'$ sous-module de $M$), on a $J^n F \cap F' = (I^n M \cap M')^{\sim} = (I^{n-h}(I^h M \cap M'))^{\sim} = J^{n-h}(J^h F \cap F')$ (I, 1, 3, cor. du th. 1). Dans le cas général, il y a un recouvrement fini de $X$ par des ouverts affines $U_i$, et pour chaque $i$, on a une relation analogue à (3) en remplaçant $J, F, F'$ par leurs restrictions à $U_i$ et $h$ par un entier $h_i$; il suffit alors de prendre pour $h$ le plus grand des $h_i$ pour obtenir (3). Enfin, le résultat précédent montre aussi que les $\psi_i$ constituent un $\partial$-foncteur de $\partial$-foncteurs, car les diagrammes
$$\begin{tikzcd} H^i(X', F^n/J^n F') \arrow[r, "\cong"] \arrow[d] & H^{i+1}(X', F/F') \arrow[d] \\ H^i(X', F^n/J^n F') \arrow[r, "\cong"] & H^{i+1}(X', F/(J^n F \cap F')) \end{tikzcd}$$
étant commutatifs, il en est de même, par passage à la limite projective, des diagrammes
$$\begin{tikzcd} \psi_i H^i(X', F/F') \arrow[r, "\cong"] \arrow[d] & H^{i+1}(X', F/F') \psi_{i+1} \arrow[d] \\ H^i(X/X', F^n) \arrow[r, "\cong"] & H^{i+1}(X/X', F') \end{tikzcd}$$
Il n'est pas certain que le $\partial$-foncteur ainsi défini soit exact, puisqu'une limite projective de suites exactes n'est pas nécessairement exacte; on saura seulement \emph{a posteriori} que ce $\partial$-foncteur est exact lorsqu'on saura que les $\psi_i$ sont des isomorphismes pour tout $\mathcal{O}_X$-module cohérent.

Proposition 3. - Si $F$ et $G$ sont deux $\mathcal{O}_X$-modules cohérents, il existe des isomorphismes canoniques fonctoriels
$$(F \otimes_{\mathcal{O}_X} G)/X' \cong (F/X') \otimes_{(\mathcal{O}_X/X')} (G/X')$$
et
$$(\operatorname{Hom}_{\mathcal{O}_X}(F, G))/X' \cong \operatorname{Hom}_{(\mathcal{O}_X/X')} (F/X', G/X')$$
\marginpar{Il suffit de définir ces homomorphismes sur les sections au-dessus de $V$ ouvert dans $X$}
, on est ramené au cas où $X=\operatorname{Spec}(A)$, $A$ étant noethérien, $F=\tilde{M}, G=\tilde{N}$, $M, N$ étant des $A$-modules de type fini.

\newpage

%% ===== Page 384 =====

Or , on sait qu'on a alors des homomorphismes canoniques fonctoriels
$$(M \otimes_A N)^\wedge \cong \hat{M} \hat{\otimes}_{\hat{A}} \hat{N}$$
$$(\operatorname{Hom}_A(M, N))^\wedge \cong \operatorname{Hom}_{\hat{A}}(\hat{M}, \hat{N})$$
où les accents désignent les complétés des modules considérés pour une topologie $I$-adique, $I$ étant un idéal de $A$ . La conclusion en résulte aussitôt (I,1,3,cor.6 du th.1).

Soient $X,Y$ deux préschémas noethériens, $f = (\varphi, \vartheta) : X \to Y$ un morphisme, $X'$ (resp. $Y'$) une partie fermée de $X$ (resp. $Y$) telles que $\varphi(X') \subset Y'$. Nous allons définir un morphisme d'espaces annelés $\hat{f} : (X', (\mathcal{O}_X)_{/X'}) \to (Y', (\mathcal{O}_Y)_{/Y'})$ ; nous prendrons $f = \varphi|_{X'}$, où $\vartheta^\# : \varphi^{-1}((\mathcal{O}_Y)_{/Y'}) \to (\mathcal{O}_X)_{/X'}$ est un homomorphisme de faisceaux d'anneaux que nous définissons de la façon suivante . Il suffit de considérer le cas où $X=\operatorname{Spec}(A)$ et $Y=\operatorname{Spec}(B)$ sont affines, $f$ correspondant à un homomorphisme d'anneaux $\varphi : B \to A$, et on a $X'=V(\mathfrak{a})$, $Y'=V(\mathfrak{b})$, $\mathfrak{b}$ (resp. $\mathfrak{a}$) étant un idéal de $A$ (resp. $B$). L'hypothèse $\varphi(X') \subset Y'$ équivaut à $\mathfrak{b} \subset \varphi^{-1}(\mathfrak{a})$ pour tout idéal premier $\mathfrak{p} \supset \mathfrak{a}$, donc à $\mathfrak{b} \subset \varphi^{-1}(\sqrt{\mathfrak{a}})$. Comme on peut supposer $\mathfrak{a}$ égal à sa racine, on peut donc supposer $\varphi(\mathfrak{b}) \subset \mathfrak{a}$ ; alors pour tout $n$, l'homomorphisme $B \to A$ définit, par passage aux quotients, un homomorphisme $B/\mathfrak{b}^n \to A/\mathfrak{a}^n$, et en passant à la limite projective, un homomorphisme $\hat{B} \to \hat{A}$, $\hat{A}$ (resp. $\hat{B}$) étant le complété de $A$ (resp. $B$) pour la topologie $\mathfrak{a}$-adique (resp. $\mathfrak{b}$-adique). Revenant au cas général, cela définit, pour deux ouverts affines $U \subset X, V \subset Y$ tels que $\varphi(U) \subset V$, un homomorphisme $\Gamma(V, (\mathcal{O}_Y)_{/Y'}) \to \Gamma(U, (\mathcal{O}_X)_{/X'})$, et on vérifie aussitôt que ces homomorphismes sont compatibles avec les opérateurs de restriction (il suffit de le faire pour deux ouverts $U'=D(g) \subset U=\operatorname{Spec}(A)$, $V'=D(f) \subset V=\operatorname{Spec}(B)$, avec $g=\varphi(f)$, ce qui est immédiat) ; on a donc bien défini l'homomorphisme $\vartheta$ cherché .
Proposition 4 . - Pour tout $\mathcal{O}_Y$-module cohérent $\mathcal{G}$, il existe un isomorphisme canonique fonctoriel de $(\mathcal{O}_X)_{/X'}$-Modules

\newpage

%% ===== Page 385 =====

- III-57 -

au triplet $(X, X', Y, Y')$; un morphisme $(f, f') : (X, X') \to (Y, Y')$ consistant en un morphisme de préschémas $f : X \to Y$ tel que $f(X') \subset Y'$ et en un homomorphisme $f^*(G) \to F$, qui donne par complétion le long de $X'$ un homomorphisme $\widehat{f^*}(G)_{|X'} \to F_{|X'}$, et à l'aide de (4), un homomorphisme $\hat{f^*}(G_{|Y'}) \to F_{|X'} $ : qu'on vérifie aussitôt être un $\mathcal{O}_{X'}$-homomorphisme. En outre, pour chaque $n$, on a vu ci-dessus qu'on a un homomorphisme $f^*(G \otimes_{\mathcal{O}_Y} (\mathcal{O}_Y/J^n)) \to f^*(G) \otimes_{\mathcal{O}_X} (\mathcal{O}_X/I^n)$ où $I$ est un faisceau d'idéaux convenable tel que le support de $\mathcal{O}_X/I$ soit $X'$, d'où (chap. 0, § $n^o$ ) pour chaque $i$ un homomorphisme
$$H^i(Y', G \otimes_{\mathcal{O}_Y} J^n) \to H^i(X', F \otimes_{\mathcal{O}_X} I^n)$$
et par passage à la limite projective, un homomorphisme
$$H^i(Y/Y', G) \to H^i(X/X', f^*(G))$$
On a de même un homomorphisme $H^i(Y', G_{|Y'}) \to H^i(X', f^*(G_{|Y'}))$ (loc. cit.), et il résulte des définitions que pour les homomorphismes $\psi_i$ définis dans (2), les diagrammes
\begin{tikzcd}
H^i(Y', G_{|Y'}) \arrow[d] \arrow[r, "\psi_i"] & H^i(Y/Y', G) \arrow[d] \\
H^i(X', f^*(G_{|Y'})) \arrow[r, "\psi_i"] & H^i(X/X', f^*(G))
\end{tikzcd}
sont commutatifs. On peut donc dire que les $\psi_i$ sont des homomorphismes fonctoriels pour les triplets $(X, X')$.

\emph{Proposition 5}. - Pour tout $\mathcal{O}_X$-module cohérent $F$, le noyau de l'homomorphisme canonique $i : \Gamma(X, F) \to \Gamma(X', F_{|X'})$ est formé des sections nulles dans un voisinage de $X'$.

Il est évident d'après (1) que l'image canonique d'une telle section est nulle. Réciproquement, si $i(s)=0$, il suffit de voir que tout point $x \in X'$ admet un voisinage dans $X$ dans lequel $s$ est nulle, et on peut donc se ramener au cas où $X = \operatorname{Spec}(A)$ est affine, $X' = V(\mathfrak{a})$, où $\mathfrak{a}$ est un idéal de $A$, et $F = \widetilde{M}$, $M$ étant un $A$-module de type fini. Alors $\Gamma(X', F_{|X'})$ est le complété $\widehat{M}$ de $M$ pour la topologie $\mathfrak{a}$-adique, et

\newpage

%% ===== Page 386 =====

- III - 58 -

le noyau de cet homomorphisme est l'ensemble des $m$ annulés par un élément de $1+\mathfrak{a}$ (Bourbaki, Alg. comm., ). On a donc $(1+f)s=0$ avec $f \in \mathfrak{a}$ ; pour tout $x \in X'$ on en déduit $(1+f_x)s_x=0$ et comme $1+f_x$ est inversible dans $\mathcal{O}_{X,x}$, $s_x=0$, d'où la proposition.

\emph{Corollaire 1}. — Soit $u : \mathcal{F} \to \mathcal{G}$ un homomorphisme de $\mathcal{O}_X$-modules cohérents. Pour que $u_{|X'} : \mathcal{F}_{|X'} \to \mathcal{G}_{|X'}$ soit nul, il faut et il suffit que $u$ soit nul dans un voisinage de $X'$.

Tenant compte de ce que le foncteur $\mathcal{F}_{|X'}$ est exact (prop. 1), il suffit d'appliquer la prop. 5 au faisceau $\operatorname{Im} u$.

\emph{Corollaire 2}. — Soient $X, Y$ deux $S$-préschémas noethériens, $Y$ étant de type fini sur $S$. Soient $f, g : X \to Y$ deux $S$-morphismes de $X$ dans $Y$ tels que $f(X') \subset Y'$, $g(X') \subset Y'$. Pour que $f_{|X'} = g_{|X'}$, il faut et il suffit que $f$ et $g$ coïncident dans un voisinage de $X'$.

La condition est évidemment suffisante (sans condition de finitude pour $Y$). Pour voir qu'elle est nécessaire, remarquons d'abord que par définition on a $f(x)=g(x)$ pour tout $x \in X'$. D'autre part, la question étant locale, on peut supposer que $X$ et $Y$ sont des voisinages affines respectifs de $x$ et de $y=f(x)=g(x)$. Alors $f$ et $g$ correspondent à deux $\Gamma(S, \mathcal{O}_S)$-homomorphismes de $\Gamma(Y, \mathcal{O}_Y)$ dans $\Gamma(X, \mathcal{O}_X)$, soient $\varphi$ et $\sigma$, et par hypothèse, pour toute section $s \in \Gamma(Y, \mathcal{O}_Y)$, les images canoniques de $\varphi(s)$ et $\sigma(s)$ dans $\Gamma(X', (\mathcal{O}_X)_{|X'})$ sont les mêmes ; par la prop. 5 on en conclut que $\varphi(s)$ et $\sigma(s)$ coïncident dans un voisinage de $X'$, et par suite les deux homomorphismes de $\mathcal{O}_y$ dans $\mathcal{O}_x$ déduits de $f$ et de $g$ coïncident. On en conclut (I, 4, 3, prop. 15) que $f$ et $g$ coïncident dans un voisinage de $x$, d'où le corollaire.

\emph{Proposition 6}. — Soient $X$ un préschéma noethérien, $X'$ une partie fermée de $X$, $\mathcal{F}$ un $\mathcal{O}_X$-module cohérent. Alors l'homomorphisme canonique $\psi_i : H^i(X', \mathcal{F}_{|X'}) \to H^i(X/X', \mathcal{F}) = \varprojlim_n H^i(X', \mathcal{F} \otimes \mathcal{O}_X/\mathcal{I}^n)$ est sur-

\newpage

%% ===== Page 387 =====

- III-59 -

\noindent i \emph{l'homomorphisme} $I \to \mathcal{M}$. On sait que le noyau de cet homomorphisme est l'ensemble des $\mathcal{M}$ annulés par un élément de $1+\mathfrak{a}$ (Bourbaki, Alg. Comm., \dots). On a donc $(1+f)s=0$ avec $f \in \mathfrak{a}$ ; pour tout $x \in X'$ on en déduit $(1+f_x)s_x=0$ et comme $1+f_x$ est inversible dans $\mathcal{O}_{X,x}$, $s_x=0$, d'où la proposition.

\noindent Corollaire 1. — Soit $u : \mathcal{F} \to \mathcal{G}$ un homomorphisme de $\mathcal{O}_X$-modules cohérents. Pour que $\mathcal{F}|_{X'} \to \mathcal{G}|_{X'}$ soit nul, il faut et il suffit que $u$ soit nul dans un voisinage de $X'$.

\noindent Tenant compte de ce que le foncteur $\mathcal{F}|_{X'}$ est exact (prop. 1), il suffit d'appliquer la prop. 5 au faisceau $\operatorname{Im} u$.

\noindent Corollaire 2. — Soient $X, Y$ deux $S$-préschémas noethériens, $Y$ étant de type fini sur $S$. Soient $f, g$ deux $S$-morphismes de $X$ dans $Y$ tels que $f(X') \subset Y', g(X') \subset Y'$. Pour que $f|_{X'} = g|_{X'}$, il faut et il suffit que $f$ et $g$ coïncident dans un voisinage de $X'$.

\noindent La condition est évidemment suffisante (sans condition de finitude pour $Y$). Pour voir qu'elle est nécessaire, remarquons d'abord que par définition on a $f(x)=g(x)$ pour tout $x \in X'$. D'autre part, la question étant locale, on peut supposer que $X$ et $Y$ sont des voisinages affines respectifs de $x$ et de $y=f(x)=g(x)$. Alors $f$ et $g$ correspondent à deux $\Gamma(S, \mathcal{O}_S)$-homomorphismes de $\Gamma(Y, \mathcal{O}_Y)$ dans $\Gamma(X, \mathcal{O}_X)$, soient $\rho$ et $\sigma$, et par hypothèse, pour toute section $s \in \Gamma(Y, \mathcal{O}_Y)$, les images canoniques de $\rho(s)$ et $\sigma(s)$ dans $\Gamma(X', (\mathcal{O}_X)_{X'})$ sont les mêmes ; par la prop. 5 on en conclut que $\rho(s)$ et $\sigma(s)$ coïncident dans un voisinage de $X'$, et par suite les deux homomorphismes de $\mathcal{O}_{Y,y}$ dans $\mathcal{O}_{X,x}$ déduits de $f$ et de $g$ coïncident. On en conclut (I, 4, 3, prop. 15) que $f$ et $g$ coïncident dans un voisinage de $x$, d'où le corollaire.

\noindent Proposition 6. — Soient $X$ un préschéma noethérien, $X'$ une partie fermée de $X$, $\mathcal{F}$ un $\mathcal{O}_X$-module cohérent. Alors l'homomorphisme canonique $\psi_i : H^i(X', \mathcal{F}|_{X'}) \to H^i(X/X', \mathcal{F}) = \varprojlim H^i(X', \mathcal{F} \otimes \mathcal{O}_X/\mathcal{J}^n)$ est sur-

\newpage

%% ===== Page 388 =====

- III-60 -

Comme $X$ est un schéma, la cohomologie $H^q(X, \mathcal{F})$ d'un faisceau quasi-cohérent quelconque est égale à la cohomologie $H^q(\mathcal{U}, \mathcal{F})$ pour un recouvrement fini par des ouverts affines (§ 1, n$^\circ$4, prop.3), d'où résulte que les $H^q(X', \mathcal{F} \otimes \mathcal{O}_X / \mathcal{J}^n)$ sont des modules de type fini sur l'anneau $\Gamma(X', \mathcal{O}_X / \mathcal{J}^n)$. Ce dernier étant artinien, on sait (chap.0, compl.) que les systèmes projectifs $(H^q(X', \mathcal{F} \otimes \mathcal{O}_X / \mathcal{J}^n))_{n \in \mathbb{N}}$ vérifient la condition (ML) pour tout $i$, d'où le corollaire.

On notera que ce corollaire s'applique par exemple lorsque le sous-préschéma $(X', \mathcal{O}_X / \mathcal{J}^n)$ est complet pour tout $n$ (n$^\circ$2).

4. Le théorème fondamental des morphismes propres : énoncés et corollaires.

Soient $X, Y$ deux préschémas noethériens, $f: X \to Y$ un morphisme propre, $Y'$ une partie fermée de $Y$, $X'$ son image réciproque ($f^{-1}(Y')$).

Pour tout $\mathcal{O}_X$-module cohérent $\mathcal{F}$, nous allons considérer sur $Y'$, muni du faisceau d'anneaux $(\mathcal{O}_Y)_{Y'}$, les faisceaux suivants :

1$^\circ$ $R^q f_* (\mathcal{F}_{X'})$, $f$ étant le morphisme $(X', \mathcal{O}_X / \mathcal{I}) \to (Y', \mathcal{O}_Y / \mathcal{J})$ défini au n$^\circ$3, ce dernier étant un morphisme propre.

2$^\circ$ Les complétés $(R^q f_* (\mathcal{F})) \hat{}_Y$, ces derniers étant des $\mathcal{O}_Y$-modules cohérents, en vertu du n$^\circ$1, th.1.

3$^\circ$ Si $\mathcal{J}$ est un faisceau cohérent d'idéaux dans $\mathcal{O}_Y$ tel que $\operatorname{supp}(\mathcal{O}_Y / \mathcal{J}) = Y'$, on sait (II, 1, 1, cor. de la prop.2) que $\mathcal{I} = f^*(\mathcal{J})$ est un faisceau cohérent d'idéaux dans $\mathcal{O}_X$ tel que $\operatorname{supp}(\mathcal{O}_X / \mathcal{I}) = X'$ et que $\mathcal{I}^n = f^*(\mathcal{J}^n)$ ; par suite $\mathcal{F} \otimes \mathcal{O}_X / \mathcal{I}^n = \mathcal{F} \otimes \mathcal{O}_Y (\mathcal{O}_Y / \mathcal{J}^n)$. Les faisceaux $R^q f_* (\mathcal{F} \otimes \mathcal{O}_X / \mathcal{I}^n)$ ($n \in \mathbb{N}$) formant un système projectif de $\mathcal{O}_Y$-modules cohérents (n$^\circ$1, th.1). Nous noterons la limite projective de ce système $\widehat{R^q f_*} (\mathcal{F})$.

Si $V$ est une partie ouverte affine de $Y$, le module des sections de $R^q f_* (\mathcal{F})$ au-dessus de $V$ est $H^q(f^{-1}(V), \mathcal{F})$.

\newpage

%% ===== Page 389 =====

- 111 - 61 -

(§ 1, n° 4, cor. de la prop. 4). Le faisceau $(R^p f_* (\mathcal{F}))_{/Y'}$ est donc associé au préfaisceau $V \to \varprojlim_n (H^p (f^{-1} (V), \mathcal{F}) \otimes_{\Gamma (V, \mathcal{O}_Y)} \Gamma (V, \mathcal{O}_Y / \mathcal{J}^n))$, le produit tensoriel étant pris sur $\Gamma (V, \mathcal{O}_Y)$ (I, 1, 3, cor. 6 du th. 4). De même, le faisceau $R^p f_{X/X'} (\mathcal{F})$ est associé au préfaisceau $V \to \varprojlim_n H^p (f^{-1} (V), \mathcal{F} \otimes_{\mathcal{O}_X} (\mathcal{O}_Y / \mathcal{J}^n))$ (ce qui montre déjà que son support est bien contenu dans $Y'$). Enfin, on sait que $R^p f_* (\mathcal{F}_{/X'})$ est associé par définition au préfaisceau $V \to H^p (f^{-1} (V) \cap X', \mathcal{F}_{/X'})$. Cela étant, on a pour chaque $p \ge 0$, deux homomorphismes

$$(D_V) \quad \begin{tikzcd} & \varprojlim_n H^p (f^{-1} (V), \mathcal{F} \otimes_{\mathcal{O}_X} (\mathcal{O}_Y / \mathcal{J}^n)) \arrow[ld, "\varphi_p^V"'] \arrow[rd, "\psi_p^V"] & \\ \varprojlim_n (H^p (f^{-1} (V), \mathcal{F}) \otimes_{\Gamma (V, \mathcal{O}_Y)} \Gamma (V, \mathcal{O}_Y / \mathcal{J}^n)) & & H^p (f^{-1} (V) \cap X', \mathcal{F}_{/X'}) \end{tikzcd}$$

où $\psi_p^V$ a été défini au n° 3, formule (2), et $\varphi_p^V$ est l'homomorphisme obtenu par passage à la limite projective, à partir des homomorphismes $H^p (f^{-1} (V), \mathcal{F}) \to H^p (f^{-1} (V), \mathcal{F} \otimes_{\mathcal{O}_X} (\mathcal{O}_Y / \mathcal{J}^n))$, eux-mêmes déduits des homomorphismes $\mathcal{F} \to \mathcal{F} \otimes_{\mathcal{O}_X} (\mathcal{O}_Y / \mathcal{J}^n)$.

On vérifie aussitôt que les homomorphismes $\varphi_p^V$ et $\psi_p^V$ sont compatibles avec les opérations de restriction, et définissent donc deux homomorphismes de $\mathcal{O}_Y / \mathcal{J}$-modules :

$$(D) \quad \begin{tikzcd} & R^p f_{X/X'} (\mathcal{F}) \arrow[ld, "\varphi_p"'] \arrow[rd, "\psi_p"] & \\ (R^p f_* (\mathcal{F}))_{/Y'} & & R^p f_* (\mathcal{F}_{/X'}) \end{tikzcd}$$

On notera encore que le système des $(R^p f_* (\mathcal{F}))_{/Y'}$ est un foncteur cohomologique (I, II, 2. 1) en $\mathcal{F}$, en vertu de la prop. 1 du n° 3 ; il en est de même du système des $R^p f_* (\mathcal{F}_{/X'})$ (n° 3, cor. 2 de la prop. 1) ; enfin (n° 3, prop. 2) le système des $R^p f_{X/X'} (\mathcal{F})$ est un $\partial$-foncteur, qui deviendra un foncteur cohomologique lorsque nous aurons prouvé qu'il est exact. Les mêmes remarques s'appliquent aux systèmes de modules des diagrammes $(D_V)$; enfin, il est immédiat que les $\varphi_p$ et $\psi_p$ forment des systèmes d'homomorphismes compatibles a...

\newpage

%% ===== Page 390 =====

\noindent\makebox[\textwidth][c]{--- III-62 ---}

\noindent vee les opérateurs cobords, autrement dit, sont des morphismes de $\partial$-foncteurs (I, II, 2.1) ; il en est de même de $\varphi_p$ et $\psi_p$.

\noindent \textbf{Théorème 2 (théorème fondamental des morphismes propres).} --- Soient $f: X \to Y$ un morphisme propre de schémas noethériens, $Y'$ une partie fermée de $Y$, $X' = f^{-1}(Y')$. Alors, pour tout $\mathcal{O}_X$-Module cohérent $\mathcal{F}$, les $R^p f_*(\mathcal{F})$ sont des $\mathcal{O}_Y$-Modules cohérents, et les homomorphismes $\varphi_p$ et $\psi_p$ du diagramme (D) sont des isomorphismes.

\noindent Faisant en particulier $p=0$, on trouve l'énoncé suivant, généralisation naturelle du théorème fondamental des fonctions holomorphes de Zariski :

\noindent \textbf{Corollaire 1.} --- Avec les notations du th. 2, le groupe des sections formelles de $f_*(\mathcal{F})$ le long de $Y'$ est canoniquement isomorphe au groupe des sections formelles de $\mathcal{F}$ le long de $X'$.

\noindent En particulier, prenant $\mathcal{F} = \mathcal{O}_X$ (qui est le cas considéré par Zariski)

\noindent \textbf{Corollaire 2.} --- L'homomorphisme $A(Y/Y') \to A(X/X')$ déduit du morphisme $\hat{f}: (X', \mathcal{O}_{X'}/\mathcal{J}_{X'}) \to (Y', \mathcal{O}_{Y'}/\mathcal{J}_{Y'})$ est un isomorphisme lorsque $f_*(\mathcal{O}_X) = \mathcal{O}_Y$.

\noindent Nous développerons au § 4 certaines conséquences importantes de ce dernier corollaire, généralisant des résultats classiques dus à Zariski. On notera que, dans les conditions du cor. 2, l'isomorphisme du cor. 1 est un isomorphisme de modules sur $A(Y/Y') = A(X/X')$.

\noindent Nous démontrerons en fait un résultat en apparence plus fort que le th. 2, savoir :

\noindent \textbf{Corollaire 3.} --- Sous les hypothèses du th. 2, pour tout ouvert affine $V$ de $Y$, les homomorphismes $\varphi_p^V$ et $\psi_p^V$ du diagramme $(D_V)$ sont des isomorphismes.

\noindent Mais en fait, ce corollaire est équivalent au th. 2. En effet, il suffit de déduire de ce dernier le

\newpage

%% ===== Page 391 =====

\marginpar{Archives Grothendieck-sept 59}

--- III-63 ---

Corollaire 4. --- Sous les conditions du th. 2, le préfaisceau $V \cap Y' \to H^p(f^{-1}(V) \cap X', \mathcal{F}/X')$ est identifié à son faisceau associé $R^p\mathcal{F}'(\mathcal{F}/X')$.

Nous déduirons ce dernier d'un autre corollaire du th. 2 :

Corollaire 5. --- Sous les conditions du th. 2, le foncteur cohomologique $H^n(X', \mathcal{F}/X')$ est l'aboutissement d'un foncteur spectral cohomologique dont le terme $E_2$ est
$$E_2^{pq} = H^p(Y', R^qf_*(\mathcal{F})/Y').$$

En effet, la suite spectrale de Leray pour le morphisme $\hat{f} : X' \to Y'$ et le faisceau $\mathcal{F}/X'$ a un terme $H^p(Y', R^q\hat{f}_*(\mathcal{F}/X'))$ et un aboutissement $H^n(X', \mathcal{F}/X')$ (G, II, 4.17.1), et il suffit d'utiliser l'isomorphisme fonctoriel $R^q\hat{f}_*(\mathcal{F}/X') \simeq (R^qf_*(\mathcal{F}))/Y'$ défini dans le th. 2.

Pour démontrer alors le cor. 4 ; il suffit de remarquer que l'on a $H^p(V \cap Y', (R^qf_*(\mathcal{F}))/Y') = 0$ pour tout $p > 0$ et tout $q \ge 0$, en vertu du cor. 1 de la prop. 6 du n° 3, les $R^qf_*(\mathcal{F})$ étant cohérents (n° 1, th. 1). Alors on sait que $H^q(f^{-1}(V) \cap X', \mathcal{F}/X') = \Gamma(V, R^qf_*(\mathcal{F})/Y') = \Gamma(V, R^qf_*(\mathcal{F}/X'))$ (G, II, 4.4.1 appliqué avec $q(n)=n$).

Un cas particulier important du th. 2 est celui où $Y'$ est réduit à un point (fermé) ; on obtient alors l'expression des complétés, pour la topologie $m_y$-adique, des $\mathcal{O}_y$-modules de type fini $(R^pf_*(\mathcal{F}))_y$. Plus généralement, sans supposer le point $y \in Y$ fermé :

Corollaire 6. --- Soient $Y$ un schéma localement noethérien, $f : X \to Y$ un morphisme propre, $\mathcal{F}$ un $\mathcal{O}_X$-module cohérent. Alors les $\mathcal{O}_Y$-modules $R^pf_*(\mathcal{F})$ sont cohérents et pour tout $y \in Y$, on a un isomorphisme canonique de $\partial$-foncteurs
$$\varprojlim_n (R^pf_*(\mathcal{F})) \otimes (\mathcal{O}_y/m_y^n) \simeq \varprojlim_n H^p(f^{-1}(y), \mathcal{F} \otimes (\mathcal{O}_y/m_y^n))$$
où la fibre $f^{-1}(y)$ est considérée comme espace de base du préschéma $X \times_Y \operatorname{Spec}(\mathcal{O}_y/m_y^n)$ (I, 2, 7, prop. 17).

\newpage

%% ===== Page 392 =====

\noindent - II - 64 -
\emph{Y est noethérien et}
Comme on l'a déjà remarqué, lorsque $\{y\}$ est fermé, ce corollaire est un cas particulier du th. 2. Dans le cas général, posons $Y_1 = \operatorname{Spec}(\mathcal{O}_y)$, $X_1 = X \times_Y Y_1$, $\mathcal{F}_1 = \mathcal{F} \otimes_{\mathcal{O}_Y} \mathcal{O}_{Y_1}$, et soit $f_1$ le morphisme $f \times 1 : X_1 \to Y_1$; $Y_1$ est noethérien et $f_1$ est propre (II, 5, 1, prop. 2) et $\mathcal{F}_1$ est cohérent (II, 1, 1, remarque suivant la déf. 1). Soit $y_1$ l'unique point fermé de $Y_1$; le th. s'applique à $f_1, \mathcal{F}_1$ et $y_1$; on a $\mathcal{O}_{y_1} = \mathcal{O}_y$, $f_1^{-1}(y_1) = f^{-1}(y)$, les préschémas $X \times_Y \operatorname{Spec}(\mathcal{O}_y/\mathfrak{m}_y^n)$ et $X_1 \times_{Y_1} \operatorname{Spec}(\mathcal{O}_y/\mathfrak{m}_y^n)$ étant identifiés canoniquement (I, 2, 5, formule (1)), et $\mathcal{F}_1 \otimes_{\mathcal{O}_{Y_1}} (\mathcal{O}_y/\mathfrak{m}_y^n)$ s'identifie à $\mathcal{F} \otimes_{\mathcal{O}_Y} (\mathcal{O}_y/\mathfrak{m}_y^n)$ (II, 1, 1, prop. 5). Il reste donc à voir que $R^p f_{1*}(\mathcal{F}_1)$ est canoniquement isomorphe à $R^p f_*(\mathcal{F}) \otimes_{\mathcal{O}_Y} \mathcal{O}_{Y_1}$, ce qui résulte du § 1, n°4, prop. 5.

Notons aussi le corollaire suivant, qui sera obtenu au § par une méthode plus élémentaire (n'utilisant que le th. de finitude au lieu du th. 2) :

Corollaire 7. -- Soient $Y$ un préschéma localement noethérien, $f: X \to Y$ un morphisme propre, et un point de $Y$, $r$ la dimension de $f^{-1}(y)$; alors, pour tout $\mathcal{O}_X$-module cohérent $\mathcal{F}$, les faisceaux $R^p f_*(\mathcal{F})$ sont nuls au voisinage de $y$ pour tout $p > r$.

En effet, on a alors $H^p(f^{-1}(y), \mathcal{F} \otimes (\mathcal{O}_y/\mathfrak{m}_y^n)) = 0$ (G, II, 4. 15. 2) pour tout $n$, donc (cor. 6) le complété du $\mathcal{O}_y$-module $(R^p f_*(\mathcal{F}))_y$ pour la topologie $\mathfrak{m}_y$-adique est nul, et a fortiori $(R^p f_*(\mathcal{F}))_y = 0$. Le faisceau $R^p f_*(\mathcal{F})$ étant cohérent (n°1, th. 1), on en conclut le corollaire.

Corollaire 8. -- Soient $Y = \operatorname{Spec}(A)$, où $A$ est un anneau noethérien, $f: X \to Y$ un morphisme propre. Soit $\mathfrak{m}$ un idéal

\newpage

%% ===== Page 393 =====

- 111-65 -

faisceau $\varprojlim F_k$, où $F_k = F \otimes_A (A/\mathfrak{m}^k)$. Si $F$ et $G$ sont deux $\mathcal{O}_X$-modules cohérents, dénotons par $\widehat{\mathcal{H}om}_{\mathcal{O}_X}(F, G)$ le complété $\mathfrak{m}$-adique du $A$-module $\mathcal{H}om_{\mathcal{O}_X}(F, G)$ ; on a un isomorphisme canonique que

(4) $$\widehat{\mathcal{H}om}_{\mathcal{O}_X}(F, G) \xrightarrow{\sim} \varprojlim_k \mathcal{H}om_{(\mathcal{O}_X)_k}(F_k, G_k) \quad .$$

En effet, le faisceau $\mathcal{H}om_{\mathcal{O}_X}(F, G)$ est cohérent (II, 1, 1, prop. 3) et on a $\Gamma(X, L) = \mathcal{H}om_{\mathcal{O}_X}(F, G)$ (FAC, I, 2, 14, prop. 5). En vertu du cor. 3, appliqué pour $p=0$ et $Y=Y$, on a un isomorphisme canonique

$\Gamma(X, L) \xrightarrow{\sim} \varprojlim \Gamma(X, L_k) = \Gamma(X, \widehat{L})$. Or, pour tout faisceau cohérent $F$ sur $X$, $\widehat{F}$ n'est autre que $F/X'$, où $X' = f^{-1}(Y')$, comme il résulte des définitions, et on a $\widehat{F}/X' = \mathcal{H}om_{(\mathcal{O}_X)_{X'}} (F/X', G/X')$, donc $\Gamma(X, \widehat{L}) = \widehat{\mathcal{H}om}_{\mathcal{O}_X}(F, G)$. On sait d'autre part que si $B$ est une $A$-algèbre de type fini, $M, N$ des $B$-modules de type fini, $\widehat{B}, \widehat{M}, \widehat{N}$ les complétés $\mathfrak{m}$-adiques de $B, M, N$, on a $\widehat{\mathcal{H}om}_B(M, N) = \varprojlim (\mathcal{H}om_{B_k}(M_k, N_k))$, où $B_k = B \otimes_A (A/\mathfrak{m}^k)$, $M_k = M \otimes_A (A/\mathfrak{m}^k)$ et $N_k = N \otimes_A (A/\mathfrak{m}^k)$ (Bourbaki, Alg. comm.) d'où la dernière relation (4).

Corollaire 9. - Les hypothèses étant celles du cor. 6, désignons par $F_n$ le faisceau $F \otimes (\mathcal{O}_Y/\mathfrak{m}_y^n)$ sur $f^{-1}(y)$. Alors, si $\widehat{\mathcal{H}om}_{\mathcal{O}_X}(F, G)$ désigne le complété $\mathfrak{m}_y$-adique de $\mathcal{H}om_{\mathcal{O}_X}(F, G)$ : on a un isomorphisme canonique

(5) $$\widehat{\mathcal{H}om}_{\mathcal{O}_X}(F, G) \xrightarrow{\sim} \varprojlim_n \mathcal{H}om_{(\mathcal{O}_X)_n}(F_n, G_n) \quad .$$

Comme dans la démonstration du cor. 6, on se ramène au cas où le point $y$ est fermé, et le résultat est alors un cas particulier du cor. 8 (car on peut évidemment supposer $Y$ affine).

Remarques. - La partie la plus importante du th. 2 est le fait que $\varphi_p$ est un isomorphisme, ce qui peut s'exprimer sans introduire explicitement la cohomologie formelle, et est de nature locale sur $Y$. Mais pour en déduire le résultat plus fort contenu dans le cor. 5,

\newpage

%% ===== Page 394 =====

(qui est global à la fois en $X$ et $Y$), il faut avoir su interpréter les deux membres isomorphes par $\varphi_p$ comme étant $R^p f_* (F/X')$. D'ailleurs, comme on ne sait pas a priori si $R^p f_{X/X'} (F)$ est un $\delta$-foncteur exact en $F$, il semble qu'il faille commencer par démontrer que $\varphi_p$ est un isomorphisme. On notera enfin qu'on a énoncé à nouveau le th. 1 du n° 1 dans l'énoncé du th. 2, car il semble qu'il n'y ait pas d'exemple d'application du th. 2 où on ne se serve aussi du th. 1.

\subsection*{5. Démonstration du théorème fondamental : résultats préliminaires.}

Théorème 3. -- Soient $Y$ un préschéma noethérien, $S$ une $\mathcal{O}_Y$-algèbre quasi-cohérente graduée à degrés positifs, engendrée par $S_1$, ce dernier étant un $\mathcal{O}_Y$-module cohérent. Soient $f : X \to Y$ un morphisme propre, $S' = S \otimes_{\mathcal{O}_Y} f^* (\mathcal{S})$, $M = \bigoplus_k M_k$ un $S'$-module gradué, quasi-cohérent et de type fini sur $S'$. Soit enfin $Y' = \operatorname{Proj}(S)$, et soit $g : Y' \to Y$ le morphisme structural correspondant. Alors :

(i) Il existe un entier $h$ et, pour tout $p \in \mathbb{Z}$, un $\mathcal{O}_{Y'}$-module cohérent $G_p$ (unique à un isomorphisme près), tel que la relation $k \ge h$ implique :
$$R^p f_* (M_k) \cong G_p(k)$$
ces isomorphismes étant compatibles avec les structures de $S$-modules.

(ii) Supposons en outre que $X = \operatorname{Proj}(T)$, où $T$ est une $\mathcal{O}_Y$-algèbre quasi-cohérente graduée à degrés positifs, engendrée par $T_1$, ce dernier étant un $\mathcal{O}_Y$-module cohérent. Alors il existe un entier $n_0$ tel que pour tout $n \ge n_0$, on ait :
$$R^p f_* (M(n)) = 0 \text{ pour tout } p > 0 \text{ et tout } k \ge 0$$
et que l'homomorphisme canonique $f^* (f_* (M(n))) \to M(n)$ soit surjectif pour tout $k$.

Soit $X' = \operatorname{Proj}(S') = X \times_Y Y'$ (II, 4.4, prop. 14), de sorte que l'on a le diagramme commutatif :
\begin{tikzcd}
X' \arrow[r, "g'"] \arrow[d, "f'"] & X \arrow[d, "f"] \\
Y' \arrow[r, "g"] & Y
\end{tikzcd}
et $\mathcal{O}_{X'}(n) = \mathcal{O}_X(n) \otimes \mathcal{O}_{Y'}(n)$ (loc. cit.). En vertu du § 2, n° 2, cor. 3 du th. 1,

\newpage

%% ===== Page 395 =====

- III-67 -

si on pose $\mathcal{F}=\widetilde{M}$, $\mathcal{F}$ est un $\mathcal{O}_{X'}$-module cohérent et il existe un entier $h \geqslant 0$ tel que pour $k \geqslant h$, on ait
(8) $\mathcal{M}_k = g'_* (\mathcal{F}(k))$
et en outre
(9) $R^p g'_* (\mathcal{F}(k)) = 0$ pour $p > 0$
en vertu du $\S$ 2, n$^\circ$2, th. 1.

Nous utiliserons le lemme élémentaire suivant :

\emph{Lemme 2} .— Soient $Z, Z'$ deux espaces annelés, $u : Z' \to Z$ un morphisme, $\mathcal{E}$ un $\mathcal{O}_Z$-module localement libre de type fini, $\mathcal{N}$ un $\mathcal{O}_{Z'}$-module. Pour tout $p \geqslant 0$, on a un isomorphisme canonique factoriel en $\mathcal{N}$
(10) $\mathcal{E} \otimes_{\mathcal{O}_Z} R^p u_* (\mathcal{N}) \simeq R^p u_* (u^*(\mathcal{E}) \otimes_{\mathcal{O}_{Z'}} \mathcal{N})$.

On peut en effet définir un homomorphisme du premier membre de (10) dans le second sans hypothèse sur $\mathcal{E}$ ; pour tout ouvert $U \subset Z$, on définit un homomorphisme
(11) $\Gamma(U, \mathcal{E}) \otimes H^p(u^{-1}(U), \mathcal{N}) \to H^p(u^{-1}(U), u^*(\mathcal{E}) \otimes \mathcal{N})$
en considérant la résolution simpliciale canonique $\mathcal{F}^{\bullet}(\mathcal{N})$ de $\mathcal{N}$ (G, II, 6.4) ; il est clair en effet que l'on a un homomorphisme canonique
$\Gamma(U, \mathcal{E}) \otimes \Gamma(u^{-1}(U), \mathcal{F}^p(\mathcal{N})) \to \Gamma(u^{-1}(U), \mathcal{F}^p(u^*(\mathcal{E}) \otimes \mathcal{N}))$
d'où on déduit l'homomorphisme (11), et ces homomorphismes sont compatibles avec les opérations de restriction, d'où l'homomorphisme (10). Si maintenant $\mathcal{E}$ est localement libre de type fini, on peut se restreindre aux $U$ tels que $\mathcal{E}|U$ soit isomorphe à $\mathcal{O}_Z^n|U$, et comme la cohomologie commute aux sommes directes finies, on est ramené au cas où $\mathcal{E} = \mathcal{O}_Z$ ; comme alors $u^*(\mathcal{E}) = \mathcal{O}_{Z'}$, les deux membres de (10) sont identifiés canoniquement par l'homomorphisme qui vient d'être défini.

Revenant à la démonstration du th. 3, \marginpar{le lemme 2} montre que, pour tout $\mathcal{O}_{X'}$-module cohérent localement libre $\mathcal{E}$, on a
(12) $\mathcal{E} \otimes_{\mathcal{O}_{X'}} \mathcal{M}_k = g'_* (\mathcal{E} \otimes_{\mathcal{O}_{X'}} \mathcal{F}(k))$

\newpage

%% ===== Page 396 =====

(13) \quad R^p f_* (\underline{\mathcal{E}} \otimes_{\mathcal{O}_X} \mathcal{F}(k)) = 0 \quad \text{pour } p > 0,

dès que $k \ge h$. Appliquons la suite spectrale de Leray au foncteur composé $f_* g_*$ : $R^* (f \circ g)_* (\underline{\mathcal{E}} \otimes_{\mathcal{O}_X} \mathcal{F}(k))$ est l'aboutissement d'une suite spectrale de terme $E_2^{pq} = R^p f_* (R^q g_* (\underline{\mathcal{E}} \otimes_{\mathcal{O}_X} \mathcal{F}(k)))$; comme $E_2^{pq} = 0$ pour $q > 0$ d'après (13), on a donc un isomorphisme canonique (G, I, 4.4.1)

(14) \quad R^p f_* (\underline{\mathcal{E}} \otimes_{\mathcal{O}_X} \underline{\mathcal{M}}_k) \simeq R^p (f \circ g)_* (\underline{\mathcal{E}} \otimes_{\mathcal{O}_X} \mathcal{F}(k)) \quad (k \ge h).

Mais on a $f \circ g = g \circ f'$, et la suite spectrale de Leray appliquée au foncteur composé $g_* f'_*$ montre donc que, pour $k \ge h$, $R^p (f \circ g)_* (\underline{\mathcal{E}} \otimes_{\mathcal{O}_X} \mathcal{F}(k))$ est l'aboutissement d'une suite spectrale de terme $E_2^{pq} = R^p g_* (R^q f'_* (\underline{\mathcal{E}} \otimes_{\mathcal{O}_X} \mathcal{F}(k)))$. D'ailleurs $\mathcal{F}(k) = \mathcal{F} \otimes_{\mathcal{O}_X} \mathcal{O}_X(k) = (\mathcal{F} \otimes_{\mathcal{O}_X} \mathcal{O}_X') \otimes_{\mathcal{O}_{Y'}} \mathcal{O}_{Y'}(k)$, et par suite $\underline{\mathcal{E}} \otimes_{\mathcal{O}_X} \mathcal{F}(k) = (\underline{\mathcal{E}} \otimes_{\mathcal{O}_X} \mathcal{O}_X') \otimes_{\mathcal{O}_{Y'}} \mathcal{O}_{Y'}(k)$; le lemme 2 appliqué à $f'$ et au $\mathcal{O}_{Y'}$-module localement libre $\mathcal{O}_{Y'}(k)$ montre que $R^q f'_* (\underline{\mathcal{E}} \otimes_{\mathcal{O}_X} \mathcal{F}(k)) \simeq (R^q f'_* (\underline{\mathcal{E}} \otimes_{\mathcal{O}_X} \mathcal{O}_X'))(k)$. En outre, les $R^q f'_* (\underline{\mathcal{E}} \otimes_{\mathcal{O}_X} \mathcal{F})$ sont cohérents (nº 1, th. 1) puisque $f'$ est propre (II, 5.1, prop. 2).

Prenons en particulier $\underline{\mathcal{E}} = \mathcal{O}_X$, d'où $\underline{\mathcal{E}} \otimes_{\mathcal{O}_X} \mathcal{F} = \mathcal{F}$, et posons $\underline{\mathcal{G}}_q = R^q f'_* (\mathcal{F})$. On peut alors supposer $h$ assez grand pour que, dès que $k \ge h$, on ait $R^p \underline{\mathcal{G}}_q(k) = 0$ pour tout $p > 0$ (§ 2, nº 2, th. 1) ; il existe $N$ tel que $\underline{\mathcal{G}}_q = 0$ pour tout $q \ge N$ (§ 1, nº 4, prop. 4) ; si on prend $h$ plus grand des $h$ tels que $q < N$, on aura donc

$R^p g_* (R^q f'_* (\mathcal{F}(k))) = 0$ pour $p > 0, q \ge 0$ et $k \ge h$.

D'après (G, I, 4.4.1), appliqué pour $q(n) = n$, on a donc bien un isomorphisme canonique

$$ R^p f_* (\underline{\mathcal{M}}_k) \simeq g_* (\underline{\mathcal{G}}_p(k)) $$

pour tout $p \ge 0$ et tout $k \ge h$, ce qui prouve (6). L'unicité de $\underline{\mathcal{G}}$ à un isomorphisme près résulte de ce que l'on a un isomorphisme canonique $\beta^{-1} : \underline{\mathcal{G}} \to \underline{\tilde{N}}$, où $\underline{\tilde{N}} = \bigoplus_{k \in \mathbb{Z}} g_* (\underline{\mathcal{G}}(k))$ (II, 4.3, th. 1) pour tout $\mathcal{O}_Y$-module cohérent $\underline{\mathcal{G}}$ ; utilisant le fait que le foncteur $\underline{\tilde{N}}$ est exact (II, 4.2, prop. 7), on a aussi $\underline{\tilde{N}} = \underline{N}'$, où $\underline{N}' = \bigoplus_{k \ge h} g_* (\underline{\mathcal{G}}(k))$.

\newpage

%% ===== Page 397 =====

- III - 69 -

tout $h$ ; il suffit d'appliquer cette remarque aux $\underline{G}_k$, en tenant compte de (6) .

Prouvons maintenant (ii) ; la question étant locale sur $Y$, on peut supposer $Y$ affine d'anneau $A$ noethérien . Pour chaque $k$ fixé, (ii) est valable pour tout $n \geqslant n_k$ (§ 2, n°2, th.1) ; on peut se borner à ne considérer que les $k \geqslant h$ et à montrer que les $n_k$ pour $k \geqslant h$ sont bornés . Appliquant alors (14) à $\underline{E}=\mathcal{O}_X$, on est d'une part ramené à prouver que si l'on pose $\underline{F}(n,k)=\underline{F} \otimes_{\mathcal{O}_X} \mathcal{O}_X(n,k)$, où $\mathcal{O}_X(n,k) = \mathcal{O}_X(n) \otimes_{\mathcal{O}_Y} \mathcal{O}_Y(k)$, on a $R^p(f \circ g')_* (\underline{F}(n,k)) = 0$ pour tout $p > 0$, dès que $n \geqslant n_0$ et $k \geqslant h$. D'autre part, pour prouver que $f_*(\underline{M}_k(n)) \to \underline{M}_k(n)$ est surjectif, il suffit de prouver que 

$$ g^* (f_* (f_* (\underline{M}_k(n)))) \to g^* (\underline{M}_k(n)) $$

est surjectif, puisque $g^*$ est exact à droite ; comme d'après (14), et en utilisant le lemme 2, $f_* (\underline{M}_k(n)) = f_* (g' (\underline{F}(n,k)))$, on voit, en posant $u=f \circ g'$, qu'on est ramené à prouver que $u^* (u_* (\underline{F}(n,k))) \to g^* (\underline{F}(n,k))$ est surjectif pour $n \geqslant n_0$ et $k \geqslant h$ ; il suffira (chap. 0, § 2, n°) de montrer que $\underline{F}(n,k)$ est engendré par ses sections. D'ailleurs on sait que $X$ (resp. $Y'$) peut s'identifier à un sous-schéma fermé de $\mathbb{P}_A^r$ (resp. $\mathbb{P}_A^s$) (II, 4.6, prop. 21) ; $X'$ s'identifie donc à un sous-schéma fermé de $\mathbb{P}_A^r \times_A \mathbb{P}_A^s$ (I, 2.2, prop. 5) ; en outre, si $j_1$ (resp. $j_2$) est l'injection canonique de $X$ (resp. $Y'$) dans $\mathbb{P}_1 = \mathbb{P}_A^r$ (resp. $\mathbb{P}_2 = \mathbb{P}_A^s$) on a $\mathcal{O}_X(n) = j_1^* (\mathcal{O}_{\mathbb{P}_1}(n))$, $\mathcal{O}_Y(k) = j_2^* (\mathcal{O}_{\mathbb{P}_2}(k))$ (II, 4.4, cor. 1 de la prop. 15) et par suite $\mathcal{O}_X(n,k) = j^* (\mathcal{O}_{\mathbb{P}_1 \times \mathbb{P}_2}(n,k))$, où $j$ est l'injection canonique de $X'$ dans $\mathbb{P}_1 \times \mathbb{P}_2$ et $\mathcal{O}_{\mathbb{P}}(n,k) = \mathcal{O}_{\mathbb{P}_1}(n) \otimes \mathcal{O}_{\mathbb{P}_2}(k)$ ; on a donc $j^* (\underline{F}(n,k)) = j^* (\underline{F} \otimes \mathcal{O}_{\mathbb{P}}(n,k))$. Si $f$ et $g'$ sont les morphismes $\mathbb{P}_1 \times \mathbb{P}_2 \to Y$ et $\mathbb{P}_1 \times \mathbb{P}_2 \to \mathbb{P}_2$, de sorte que $f=f_0 \circ j_1$ et $j_1 \circ g' = g' \circ j$, on a $f \circ g' = (f_0 \circ g') \circ j$. Pour démontrer que $R^p(f \circ g')_* (\underline{F}(n,k)) = 0$ pour $p > 0$, il suffira donc de prouver que $R^p (f_0 \circ g')_* (j^* (\underline{F}(n,k))) = 0$ (G, II, 4.9.1), et de même

\newpage

%% ===== Page 398 =====

- III-70 -

pour montrer que $\underline{F}(n,k)$ est engendré par ses sections, il suffit de prouver que $j^*(\underline{F})(n,k)$ l'est. On est donc ramené à prouver que le Lemme 3. - Soient $A$ un anneau noethérien, $r, s$ deux entiers $>0$, $P=\mathbb{P}^r_A \times_A \mathbb{P}^s_A$. Pour tout $\mathcal{O}_P$-module cohérent $\underline{F}$, il existe deux entiers $n_0, k_0$ tels que pour $n \ge n_0, k \ge k_0$ :
(i) on ait $H^p(P, \underline{F}(n,k))=0$ pour tout $p>0$ ;
(ii) $\underline{F}(n,k)$ soit engendré par ses sections.

Prouvons d'abord (ii). Posons, comme ci-dessus, $P_1=\mathbb{P}^r_A, P_2=\mathbb{P}^s_A$; on peut identifier $P$ au fibré projectif $\mathbb{P}^r_{P_2}$ (II, 4, 5, formule (15)), et aussi au fibré $\mathbb{P}^s_{P_1}$; on peut écrire $\underline{F}(n,k)=(\underline{F}(n))(k)$, où $\underline{F}(n)$ se rapporte à la seconde identification et $(\underline{F}(n))(k)$ à la première. Il existe un entier $n_0$ tel que pour $n \ge n_0$, $\underline{F}(n)$ soit isomorphe à un quotient d'une somme directe de faisceaux isomorphes à $\mathcal{O}_{P_2}$, donc, pour $n \ge n_0$ et pour tout $k$, $\underline{F}(n,k)$ est isomorphe à un quotient d'une somme directe de faisceaux isomorphes à $\mathcal{O}_P(k)$. Mais par le même raisonnement (lorsque $P$ est identifié à $\mathbb{P}^r_{P_2}$) on conclut qu'il existe $k_0$ tel que pour $k \ge k_0$, $\mathcal{O}_P(k)$ soit isomorphe à un quotient d'une somme directe de faisceaux isomorphes à $\mathcal{O}_P$, et finalement, pour $n \ge n_0$ et $k \ge k_0$, $\underline{F}(n,k)$ est isomorphe à un quotient d'une somme directe de faisceaux isomorphes à $\mathcal{O}_P$, ce qui établit notre assertion.

Notons maintenant que l'on a $H^i(P, \underline{F})=0$ pour $i>r$ ou $i>s$ et pour tout $\mathcal{O}_P$-module cohérent $\underline{F}$ (§ 2, n°2, th. 1 et § 1, n°4, cor. de la prop. 4).
Supposons démontré le lemme 3 (i) dans le cas particulier où $\underline{F}=\mathcal{O}_P$ ; prouvons alors que $H^p(P, \underline{F}(n,k))=0$ par récurrence sur $p$. Comme $\underline{F}$ est isomorphe à un quotient d'une somme directe $\underline{G}$ d'un nombre fini de faisceaux $\mathcal{O}_P(n_j, k_j)$, on

\newpage

%% ===== Page 399 =====

- II - 71 -

a une suite exacte $0 \to R \to G \to F \to 0$, où $R$ est cohérent ($F$ et $G$ l'étant) d'où une suite exacte
$$0 \to R(n,k) \to G(n,k) \to F(n,k) \to 0$$
quels que soient $n$ et $k$ dans $\mathbb{Z}$. La suite exacte de cohomologie donne alors une suite exacte
$$H^{i-1}(P, G(n,k)) \to H^{i-1}(P, F(n,k)) \to H^i(P, R(n,k))$$
Par hypothèse de récurrence, il existe $n_i'$ et $k_i'$ tels que $H^i(P, R(n,k)) = 0$ pour $n \ge n_i'$ et $k \ge k_i'$, et puisque (i) est supposé démontré pour $F=O_P$, il existe aussi $n_i'', k_i''$ tels que $H^{i-1}(P, G(n,k)) = 0$ pour $n \ge n_i'', k \ge k_i''$; d'où $H^{i-1}(P, F(n,k)) = 0$ pour $n \ge \sup(n_i', n_i'')$ et $k \ge \sup(k_i', k_i'')$; la récurrence descendante étant finie, cela achèvera la démonstration.

Reste donc à prouver que $H^p(P, O_P(n,k)) = 0$ pour tout $p > 0$ dès que $n$ et $k$ sont assez grands. Considérons la projection $v : P \to P_A^s$; la suite spectrale de Leray montre que la cohomologie $H^*(P, O_P(n,k))$ est l'aboutissement d'une suite spectrale dont le terme $E_2^{pq} = H^p(P_A^s, R^q v_* (O_P(n,k)))$. On a d'autre part (lemme 2) $R^q v_* (O_P(n,k)) = (R^q v_* (O_P(n))) (k)$, et on est donc ramené à calculer $R^q v_* (O_P(n))$, ou, ce qui revient au même (§ 1, n°4, cor. de la prop. 4) $H^q(v^{-1}(U), O_P(n))$ pour tout ouvert affine $U$ de $P_A^s$. Mais en considérant $P$ comme fibré projectif sur $P_A^s$, on sait (§ 2, n°1, prop. 6) que ces groupes sont nuls pour $n \ge 0$ et $q > 0$, d'où $R^q v_* (O_P(n,k)) = 0$ pour $n \ge 0, k \ge 0$ et $q > 0$, et donc (G, I, 4.4.1), $H^p(P, O_P(n,k)) = H^p(P_A^s, v_* (O_P(n))(k))$. Mais, en posant $Z = P_A^s$, on a (§ 2, n°1, prop. 6) $v_* (O_P(n)) = \bigoplus_{i=0}^n (O_Z^{i-1})$, qui est évidemment de la forme $O_Z^t$, pour un entier $t$ dépendant de $n$ et $r$. Par suite $H^p(P, O_P(n,k)) = H^p(Z, (O_Z(k))^t) = (H^p(Z, O_Z(k)))^t$, qui est nul pour $p>0$ dès que $k \ge 0$ (§ 2, n°1, cor. 1 de la prop. 5). Ceci achève donc la démonstration du th. 3.

\newpage

%% ===== Page 400 =====

- III - 72 -

Corollaire 1. - Sous les conditions du th. 3 (i), les $R^p f_* (\mathcal{M})$ ($p \ge 0$) sont des $S$-modules gradués de type fini. En particulier, pour chaque $p \ge 0$, il existe $k_p$ tel que pour $k \ge k_p, r \ge 0$, on ait
$$(15) \quad R^p f_* (\mathcal{M}_{k+r}) = S_r R^p f_* (\mathcal{M}_k)$$

En effet, d'après le § 2, cor. 4 du th. 1, $f_* (\mathcal{G}) = \bigoplus_{k \ge 0} g_* (\mathcal{G}_k)$ est un $S$-module vérifiant la condition (TF') ; d'après (6) on voit donc qu'il existe un $h_p$ tel que $\bigoplus_{k \ge h_p} R^p f_* (\mathcal{M}_k)$ soit un $S$-module de type fini. Comme en outre chacun des $R^p f_* (\mathcal{M}_k)$ pour $k < h_p$ est un $\mathcal{O}_Y$-module cohérent (n°1, th. 1), donc un $\mathcal{O}_Y$-module de type fini, $R^p f_* (\mathcal{M}_k)$ est un $S$-module de type fini. En outre, il y a un recouvrement ouvert affine $(U_i)$ de $Y$ tel que les restrictions à $U_i$ des deux membres de (15) soient égaux pour $k \ge n_i$ ; il suffit de prendre pour $k_p$ le plus grand des $n_i$.

Corollaire 2. - Soient $A$ un anneau noethérien, $\mathfrak{m}$ un idéal de $A$, $X$ un $A$-schéma propre, $\mathcal{F}$ un $\mathcal{O}_X$-module cohérent. Alors, pour tout $p \ge 0$, la somme directe $\bigoplus_{k \ge 0} H^p (X, \mathfrak{m}^k \mathcal{F})$ est un module de type fini sur l'anneau $S = \bigoplus_{k \ge 0} \mathfrak{m}^k$ ; en particulier, il existe un entier $k_p$ tel que, pour $k \ge k_p$ et $r \ge 0$, on ait
$$(16) \quad H^p (X, \mathfrak{m}^{k+r} \mathcal{F}) = \mathfrak{m}^r H^p (X, \mathfrak{m}^k \mathcal{F})$$

Il suffit d'appliquer le cor. 1 avec $Y = \operatorname{Spec}(A)$, $S = \tilde{S}$, $\mathcal{M} = \mathcal{F}$, compte tenu du § 1, n°4, cor. de la prop. 4.

Corollaire 3. - Sous les hypothèses du cor. 2, on pose $\mathcal{F}_k = \mathcal{F}/\mathfrak{m}^k \mathcal{F}$ ; on munit $H^p (X, \mathcal{F})$ de la topologie $\mathfrak{m}$-adique et chacun des $A$-modules $H^p (X, \mathcal{F}_k)$ de la topologie discrète. Alors l'homomorphisme canonique
$$H^p (X, \mathcal{F}) \to \varprojlim_k H^p (X, \mathcal{F}_k)$$
déduit des homomorphismes $\mathcal{F} \to \mathcal{F}_k$, est un isomorphisme topologique de $H^p (X, \mathcal{F})$ sur un sous-groupe de $\varprojlim H^p (X, \mathcal{F}_k)$.

Par définition de la topologie d'une limite projective, il suffit de montrer que tout voisinage de 0 dans $H^p (X, \mathcal{F})$ contient $1_k$

\newpage

%% ===== Page 401 =====

- III-73 -

noyau d'un homomorphisme $H^p(X, \mathcal{F}) \to H^p(X, \mathcal{F}_k)$. Or, en vertu du cor.2, il existe $h \ge 0$ tel que pour $k \ge h$, on ait
$$H^p(X, \mathfrak{m}^k \mathcal{F}) \simeq H^p(X, \mathfrak{m}^k \mathcal{F})$$

Considérons alors la suite exacte $0 \to \mathfrak{m}^k \mathcal{F} \to \mathcal{F} \to \mathcal{F}_k \to 0$ ; il en résulte la suite exacte de cohomologie
$$H^p(X, \mathfrak{m}^k \mathcal{F}) \to H^p(X, \mathcal{F}) \to H^p(X, \mathcal{F}_k)$$
elle montre que pour $k \ge h$, le noyau de $H^p(X, \mathcal{F}) \to H^p(X, \mathcal{F}_k)$ est contenu dans l'image de $H^p(X, \mathfrak{m}^k \mathcal{F})$, et a fortiori dans $\mathfrak{m}^{k-h} H^p(X, \mathcal{F})$, ce qui prouve le corollaire.

Corollaire 4. — Sous les hypothèses du cor.2, le système projectif $(H^p(X, \mathcal{F}_k))_{k \ge 0}$ vérifie la condition (ML).

En effet, pour deux entiers $s \le r$, on a le diagramme commutatif
\begin{tikzcd}
0 \to \mathfrak{m}^r / \mathfrak{m}^s \mathcal{F} \to \mathcal{F}_r \to \mathcal{F}_s \to 0 \\
\downarrow \to \mathcal{F} \to \mathcal{F}_r \to 0
\end{tikzcd}
les lignes étant exactes ; d'où, par la suite exacte de cohomologie un diagramme commutatif de $A$-modules
\begin{tikzcd}
H^p(\mathcal{F}_r) \to H^p(\mathcal{F}_s) \xrightarrow{\partial} H^{p-1}(\mathfrak{m}^s / \mathfrak{m}^r \mathcal{F}) \\
\uparrow \\
H^{p-1}(\mathfrak{m}^r \mathcal{F})
\end{tikzcd}
où la ligne horizontale est exacte. Soit $H'$ l'image dans $H^p(\mathcal{F}_r)$ de $H^p(\mathcal{F}_r) \xrightarrow{\partial} H^{p-1}(\mathfrak{m}^s / \mathfrak{m}^r \mathcal{F})$. Comme $A$ est noethérien et que $H^{p-1}(\mathfrak{m}^r \mathcal{F})$ est un $A$-module de type fini (n°1, cor.2 du th.1), la topologie induite sur $H'$ par la topologie $\mathfrak{m}$-adique de $H^{p-1}(\mathfrak{m}^s / \mathfrak{m}^r \mathcal{F})$ est la topologie $\mathfrak{m}$-adique sur ce $A$-module ; mais comme $\mathfrak{m}^{r-s} (\mathfrak{m}^s / \mathfrak{m}^r \mathcal{F}) = 0$, $\mathcal{F}_r$ est annulé, par $\mathfrak{m}^r$, il en est de même de $H^p(\mathcal{F}_r)$, et a fortiori de $H'$; donc la topologie $\mathfrak{m}$-adique sur $H'$ est discrète. Or, le cor.2 appliqué à $\mathfrak{m}^r \mathcal{F}$ au lieu de $\mathcal{F}$, montre que la topologie de $H^{p-1}(\mathfrak{m}^r \mathcal{F})$ est induite par la topologie limite projective des $H^{p-1}(\mathfrak{m}^s / \mathfrak{m}^r \mathcal{F})$ ($s \ge r$)

\newpage

%% ===== Page 402 =====

- III-74 -

Il existe par suite un $s_0 \ge r$ tel que la restriction de l'homomorphisme $H^{p-1}(\mathfrak{m}^r\mathcal{F}) \to H^{p-1}(\mathfrak{m}^s\mathcal{F}/\mathfrak{m}^r\mathcal{F})$ soit injective pour $s \ge s_0$. Pour $s \ge s_0$, le noyau de $H^p(\mathcal{F}_s) \to H^{p-1}(\mathfrak{m}^s\mathcal{F}/\mathfrak{m}^r\mathcal{F})$ est donc égal au noyau de $H^p(\mathcal{F}_s) \to H^{p-1}(\mathfrak{m}^r\mathcal{F})$, et en particulier est indépendant de $s$, d'où la conclusion puisque ce noyau est l'image de $H^p(\mathcal{F}_s) \to H^p(\mathcal{F}_r)$.

Corollaire 5. — Sous les hypothèses du cor. 2, si $\hat{\mathcal{F}}$ désigne le complété de $\mathcal{F}$ pour la topologie $\mathfrak{m}$-adique, l'homomorphisme canonique
$$\psi_p : H^p(X, \hat{\mathcal{F}}) \to \varprojlim_k H^p(X, \mathcal{F}_k)$$
est un isomorphisme.

En effet, soit $\mathcal{J}$ le faisceau d'idéaux $\mathfrak{m}\mathcal{O}_X$, et est donc le faisceau complété $\mathcal{F}/X_k$ où $X_k$ désigne la partie fermée de $X$ définie par le faisceau $\mathcal{J}^k$. Le cor. 4, joint à la prop. 6 du n° 3, montre alors que $\psi_p$ est bijectif.

6. Fin de la démonstration du théorème fondamental.

Le th. 2 étant de caractère local sur $Y$, on peut supposer $Y = \operatorname{Spec}(A)$, où $A$ est un anneau noethérien ; $X$ est alors un $A$-schéma propre, et $Y' = V(\mathfrak{m})$, où $\mathfrak{m}$ est un idéal de $A$. Alors $X' = f^{-1}(Y')$ est défini par le faisceau d'idéaux $\mathcal{J} = \mathfrak{m}\mathcal{O}_X$, et avec les notations des cor. 2, 3, 4, 5 du th. 3, on a $\mathcal{F}/X_k = \mathcal{F}_k = \mathcal{F} \otimes \mathcal{O}_X / \mathcal{J}^k$ ; on est ramené à démontrer que dans le diagramme

$$(D')$$

$$\begin{tikzcd}
\varprojlim_k H^p(X, \mathcal{F}_k) \arrow[r, "\psi_p"] & H^p(X, \hat{\mathcal{F}}) \\
& H^p(X, \mathcal{F}) \arrow[u, "\varphi_p"]
\end{tikzcd}$$

où $H^p(X, \hat{\mathcal{F}})$ est le complété $\mathfrak{m}$-adique de $H^p(X, \mathcal{F})$. $\psi_p$ et $\varphi_p$ sont des isomorphismes. On sait déjà (cor. 5 du th. 3) qu'il en est ainsi de $\psi_p$, et il suffit donc de le démontrer pour $\varphi_p$. Le fait que $\psi_p$ soit bijectif entraîne que $\varprojlim_k H^p(X, \mathcal{F}_k)$ est un foncteur cohomologique en $\mathcal{F}$, ce qui va être utilisé dans la suite.

A) Cas où $\mathcal{F}$ est projectif. On peut alors identifier $X$ à un sous-

\newpage

%% ===== Page 403 =====

- III-75 -

préschéma fermé d'un préschéma de la forme $\mathbb{P}^r_A$ (II,4,6,prop.21) ; $X'$ s'identifie donc à un sous-préschéma fermé de $\mathbb{P}^r_A$ défini par le faisceau d'idéaux $\mathfrak{m}\mathcal{O}_{\mathbb{P}^r_A}$, et si $j$ est l'injection canonique $X \to \mathbb{P}^r_A$, $H^p(j, \mathcal{F}_k)$ et $H^p(j, \mathcal{F})$ s'identifient respectivement à $H^p(X, \mathcal{F}_k)$ et $H^p(X, \mathcal{F})$ (G,II,4,9.1) . On peut donc supposer que $X = \mathbb{P}^r_A$.

Considérons d'abord le cas particulier où $\mathcal{F} = \mathcal{O}_X(n)$ ($n \in \mathbb{Z}$) ; $\mathcal{F}_k$ s'identifie au faisceau $\mathcal{O}_X(n) \otimes (A/\mathfrak{m}^k)$ sur le sous-préschéma fermé $X_k = X \otimes (A/\mathfrak{m}^k)$ de $X$ correspondant au faisceau d'idéaux $\mathfrak{m}^k\mathcal{O}_X$ ; par ailleurs $X_k$ s'identifie à $\mathbb{P}^r_{A/\mathfrak{m}^k}$ (II,4,5,formule (15)). Or, on a alors l'expression explicite de $H^p(X, \mathcal{O}_X(n))$ et de $H^p(X_k, \mathcal{O}_{X_k}(n))$ (§ 2, n°1, prop.5) ; ce sont des modules libres sur $A$ et $A/\mathfrak{m}^k$ respectivement, dont la dimension ne dépend que de $p$ et $n$ (et non de $k$) ; le fait que $\varphi_p$ est un isomorphisme est alors trivial.

Passons au cas où $\mathcal{F}$ est un $\mathcal{O}_X$-module cohérent quelconque ; nous allons utiliser le fait que $\varprojlim_k H^p(X, \mathcal{F}_k)$ est un foncteur cohomologique et que les $(\varphi_p)$ forment un morphisme de foncteurs cohomologiques (n° 4) ; l'assertion relative aux $\varphi_p$ sera conséquence du lemme général suivant :

Lemme 4. - Soient $A$ un anneau noethérien, $X = \operatorname{Proj}(S)$, où $S$ est une $A$-algèbre graduée de type fini engendrée par $S_1$, soient $T^*$ et $T'^*$ des foncteurs cohomologiques covariants de la catégorie des $\mathcal{O}_X$-modules cohérents dans une catégorie abélienne $\mathcal{C}$, et soit $T^* \to T'^*$ un morphisme de foncteurs cohomologiques (II,2.1). On suppose que :

(i) Pour tout $n$ assez grand, $T^*( \mathcal{O}_X(-n)) \to T'^*( \mathcal{O}_X(-n))$ est un isomorphisme.

(ii) Pour tout $i$ assez grand, $T^i \to T'^i$ est un isomorphisme.

Alors $T^* \to T'^*$ est un isomorphisme.

L'hypothèse (ii) permet de raisonner par récurrence descendante sur

\newpage

%% ===== Page 404 =====

- III-76 -

le dimension $p$. Pour tout $\mathcal{F} \in \mathcal{K}$, on sait (II, 3.4, cor. 4 du th. 1) que $\mathcal{F}$ est quotient d'une somme directe finie $\mathcal{G}$ de faisceaux de la forme $\mathcal{O}_X(-n)$, où $n$ peut être pris arbitrairement grand. On peut donc, par l'hypothèse (i), supposer que $T^p(\mathcal{G}) \to T'^p(\mathcal{G})$ est un isomorphisme. La suite exacte $0 \to \mathcal{R} \to \mathcal{G} \to \mathcal{F} \to 0$ donne alors un diagramme commutatif :

$$(17) \quad \begin{tikzcd} \dots \to T^p(\mathcal{R}) \to T^p(\mathcal{G}) \to T^p(\mathcal{F}) \to T^{p+1}(\mathcal{R}) \to T^{p+1}(\mathcal{G}) \to \dots \\ \dots \to T'^p(\mathcal{R}) \to T'^p(\mathcal{G}) \to T'^p(\mathcal{F}) \to T'^{p+1}(\mathcal{R}) \to T'^{p+1}(\mathcal{G}) \to \dots \end{tikzcd}$$

où les lignes sont exactes. Dans ce diagramme : $u_2$ et $u_5$ sont des isomorphismes d'après le choix de $\mathcal{G}$, et $u_4$ d'après l'hypothèse de récurrence. D'après le lemme des cinq, $u_3$ est surjectif. Ceci s'applique à tout faisceau cohérent sur $X$, en particulier à $\mathcal{R}$, donc $u_1$ est surjectif. Le lemme des cinq prouve alors que $u_3$ est injectif, ce qui achève la démonstration du lemme 4, et par suite du th. 2 lorsque $\mathcal{F}$ est projectif.

B) Cas général. Soit $\mathcal{K}$ la catégorie des $\mathcal{O}_X$-modules cohérents. $\mathcal{K}'$ la classe des objets de $\mathcal{K}$ formés des $\mathcal{O}_X$-modules $\mathcal{F}$ tels que dans les diagrammes (D'), les $\gamma_p$ soient des isomorphismes. Comme les $\gamma_p$ forment un morphisme de $\partial$-foncteurs exacts (II, 1.5) de $\mathcal{K}$, $\mathcal{K}'$ est une sous-classe exacte de $\mathcal{K}$; comme il résulte aussitôt du lemme des cinq appliqué aux diagrammes (17). En outre, si $\mathcal{F} \in \mathcal{K}'$, tout facteur direct de $\mathcal{F}$ appartient aussi à $\mathcal{K}'$, les foncteurs $H^p(X, \mathcal{F})$ et $\varinjlim_k H^p(X, \mathcal{F}_k)$ commutant aux sommes directes finies, et $\gamma_p$ transformant une somme directe en somme directe. Par application du corollaire du lemme de dévissage (II, 1.5, cor. du th. 3), on est ramené à prouver que lorsque $X$ est irréductible, il existe un $\mathcal{O}_X$-module cohérent de support $X$ appartenant à $\mathcal{K}'$ ; en effet, si ce point est établi, on pourra l'appliquer à toute partie fermée irréductible $Z$ de $X$, car si $j$ est l'injection canonique $Z \to X$ et

\newpage

%% ===== Page 405 =====

— III-77 —

$G$ un $\mathcal{O}_Z$-module cohérent, $H^p(Z, G)$ et $H^p(Z, G)$ s'identifient respectivement à $H^p(X, j^a(G))$ et $H^p(X, j^*(G))_K$.

En vertu du lemme de Chow (II, 5, 2, cor. 1 du th. 1), il existe un préschéma irréductible $Z$ et un morphisme projectif, surjectif et birationnel $g : Z \to X$ tel que $f \circ g : Z \to Y$ soit projectif. Posons $Z' = g^{-1}(X') = (f \circ g)^{-1}(Y')$. En vertu du th. 1 du § 2, n° 2, il existe un entier $n$ tel que $F = g_*(\mathcal{O}_Z(n))$ soit cohérent et $R^qg_*(\mathcal{O}_Z(n)) = 0$ pour tout $q > 0$ ; en outre, comme $g$ est birationnel, le support de $F$ contient un ouvert partout dense de $X$, donc est identique à $X$. Le th. 2 appliqué au morphisme projectif $g$ et au faisceau cohérent $\mathcal{O}_Z(n)$ montre que $R^pg_*(\mathcal{O}_Z(n)/_{Z'}) = (R^pg_*(\mathcal{O}_Z(n)))/X'$. Cela étant, $R^p(f \circ g)_*(\mathcal{O}_Z(n)/_{Z'})$ est l'aboutissement d'une suite spectrale de terme $E_2^{pq} = R^pf_*(R^qg_*((\mathcal{O}_Z(n))/_{Z'}))$, et on en conclut des isomorphismes
$$(18) \quad R^p(f \circ g)_*(\mathcal{O}_Z(n)/_{Z'}) \cong R^pf_*(F/_{X'}).$$
Le même raisonnement appliqué à $\mathcal{O}_Z(n)$ au lieu de $(\mathcal{O}_Z(n))/_{Z'}$ donne des isomorphismes
$$(19) \quad R^p(f \circ g)_*(\mathcal{O}_Z(n)) \cong R^pf_*(F)$$
Enfin, le th. 2 appliqué au morphisme projectif $f \circ g$ et au faisceau cohérent $\mathcal{O}_Z(n)$, montre que
$$(20) \quad \gamma_p^{-1} \circ \varphi_p : (R^p(f \circ g)_*(\mathcal{O}_Z(n)))/_{Y'} \cong R^p(f \circ g)_*((\mathcal{O}_Z(n))/_{Z'})$$
est un isomorphisme. Or les homomorphismes $\varphi_p$ et $\gamma_p$ proviennent d'homomorphismes de modules, commutent avec les homomorphismes des suites spectrales, ces derniers étant fonctoriels. On conclut donc de (18), (19) et (20) que
$$\gamma_p^{-1} \circ \varphi_p : (R^pf_*(F))/_{Y'} \to R^pf_*(F/_{X'})$$
est un isomorphisme, ce qui achève la démonstration du th. 2. En outre, on a vu en cours de route (cor. 3 du th. 3) :

\emph{Corollaire.} Lorsque $Y$ est affine, les $\gamma_p$ sont des isomorphismes de groupes topologiques.

\newpage

%% ===== Page 406 =====

\marginpar{\footnotesize \textit{Archives} \\ \textit{Grothendieck-sept 59}}

\section*{-- 111-70 --}

\section*{§ 4. Application aux théorèmes fondamentaux de la géométrie algébrique qualitative.}

\noindent Théorème 1. \emph{Théorème de connexion}

Les résultats de ce n° et du suivant généralisent des théorèmes bien connus de Zariski. On notera que, bien que les énoncés de ces théorèmes ne soient pas de nature cohomologique, nous les déduisons du th. 2 du § 3, de nature cohomologique. En fait, le cor. 2 du th. 2 du § 3 est suffisant pour les applications du présent n°, et il est lui-même de nature non cohomologique, mais c'est dans sa démonstration que les méthodes cohomologiques ont joué leur rôle, en permettant une récurrence descendante sur la dimension.

Tous les résultats de ce n° sont conséquences du théorème fondamental suivant :

\noindent Théorème 1 (théorème de connexion). -- Soient $Y$ un préschéma localement noethérien, $f : X \to Y$ un morphisme propre. Alors $A(X) = f_*(\mathcal{O}_X)$ est une $\mathcal{O}_Y$-algèbre cohérente. Soit $Y'$ le préschéma au-dessus de $Y$ tel que $A(Y') = A(X)$, qui est déterminé à un $Y$-isomorphisme près (II, 2, 2, prop. 6) ; si $f' : X \to Y'$ est le $Y$-morphisme déduit de l'isomorphisme identique $e : A(Y') \to A(X)$, on a $f'_*(\mathcal{O}_X) = \mathcal{O}_{Y'}$, et les fibres $f'^{-1}(y')$ du morphisme $f'$ ( $y' \in Y'$) sont connexes et non vides.

Soit $g : Y' \to Y$ le morphisme structural. Pour prouver que l'homomorphisme $\mathcal{O}_Y \to f_*(\mathcal{O}_X)$ entrant dans la définition du morphisme $f'$ est l'identité, il suffit de prouver, puisque $Y'$ est affine sur $Y$ (II, 2, 4, prop. 8), que $g_*(\mathcal{O}_{Y'}) \to g_*(f'_*(\mathcal{O}_X)) = f_*(\mathcal{O}_X)$ est l'identité, ce qui résulte des définitions puisque $g_*(\mathcal{O}_{Y'}) = A(Y')$ et $f_*(\mathcal{O}_X) = A(X)$.
Le fait que $A(X)$ est cohérent est un cas particulier du th. de finitude (§ 3, n° 1, th. 1). Comme $f$ est propre et $g$ séparé (le texte indique : \textit{séparé (II, 5, 1, prop. 2 (v))}), $f'$ est propre (II, 5, 1, prop. 2 (v)) ; pour terminer la démonstration du th. 1, il suffit donc d'en prouver le \marginpar{\footnotesize \textit{1 et 2, 7,} \\ \textit{prop. 13}}

\newpage

%% ===== Page 407 =====

- III-79 -

Corollaire 1 .— Sous les hypothèses du th.1 , supposons de plus que
$f_*(\mathcal{O}_X) = \mathcal{O}_Y$. Alors les fibres $f^{-1}(y)$ de $f$ ($y \in Y$) sont connexes et non
vides.

L'hypothèse entraîne déjà que $f$ est dominant, et comme $f$ est une
application fermée (II, 5, 1, prop. 1), $f$ est surjectif. Appliquons le cor. 6 du th. 2 du § 3, n°4, au cas $F = \mathcal{O}_Y$ et pour $p=0$ ;
$\varprojlim ( \mathcal{O}_Y \otimes ( \mathcal{O}_Y / \mathfrak{m}_y^n ) )$ n'est autre que le complété $\hat{\mathcal{O}}_y$ pour la topologie
$\mathfrak{m}_y$-adique. Si $f^{-1}(y)$ n'est pas connexe, il est clair que chacune
des composantes connexes donne un facteur direct non nul dans
$H^0(f^{-1}(y), \mathcal{O}_{f^{-1}(y)} \otimes (\mathcal{O}_Y / \mathfrak{m}_y^n))$, qui, par passage à la limite projective,
donne un facteur direct $\neq 0$ dans cette limite ; or, comme $\mathcal{O}_y$ est un
anneau local, le $\hat{\mathcal{O}}_y$-module $\hat{\mathcal{O}}_y$ n'a pas de facteur direct $\neq 0$.

Corollaire 2 .— Sous les hypothèses du th.1, pour tout $y \in Y$,
l'ensemble des composantes connexes de la fibre $f^{-1}(y)$ est en corres-
pondance biunivoque avec l'ensemble (fini) des points de la fibre
$g^{-1}(y)$ ($g$ étant le morphisme structural $Y' \to Y$), autrement dit avec
l'ensemble des idéaux maximaux de $(A(X))_y$.

On sait en effet, puisque $Y'$ est fini sur $Y$, que $g^{-1}(y)$ est un
espace fini discret (II, 2, 7, prop. 17). Comme $f^{-1}(y) = f'^{-1}(g^{-1}(y))$,
le corollaire résulte de cette remarque et du th. 1.

On a ainsi une interprétation géométrique remarquable du $Y$-présché-
ma $Y'$. La factorisation $f = g \circ f'$ du morphisme propre $f$ est analogue
à la factorisation obtenue par K. Stein pour les applications holomor-
phes d'espaces analytiques, et nous l'appellerons par la suite la
factorisation de Stein de $f$.

Corollaire 3 .— Sous les conditions du th. 1, supposons $X$ et $Y$ inté-
gres et $f$ dominant ; soient $K(X) \supset K(Y)$ leurs corps de fractions ra-
tionnelles (I, 6, 2, prop. 3). Pour tout $y \in Y$, le nombre de composan-

\newpage

%% ===== Page 408 =====

- III-80 -

tes connexes de $f^{-1}(y)$ est au plus égal au nombre des idéaux maximaux
de la fermeture intégrale $\mathcal{O}'_y$ de $\mathcal{O}_y$ dans $R(X)$.

En effet, comme $f(X)$ est fermé, on a $f(X)=Y$, donc $g$ est aussi
surjectif. Comme $(A(X))_y$ est un $\mathcal{O}_y$-module de type fini, donc conte-
nu dans $\mathcal{O}'_y$, tout idéal maximal de $(A(X))_y$ est l'intersection de cet
anneau et d'un idéal maximal de $\mathcal{O}'_y$ (th. de Cohen-Seidenberg).

Définition 1. - On dit qu'un anneau local intègre $A$ est unibranche
si sa clôture intégrale est un anneau local. Un point $y$
d'un préschéma est unibranche si l'anneau local $\mathcal{O}_y$ est unibranche.
On notera que si le complété de $A$ est intègre (ce qu'on exprime en
disant que $A$ est analytiquement irréductible), $A$ est unibranche (
).

Corollaire 4. - Sous les hypothèses du cor. 3, supposons que la ferme-
ture algébrique de $R(Y)$ dans $R(X)$ soit radicielle sur $R(Y)$
et que $Y$ soit unibranche. Alors la fibre $f^{-1}(y)$ est connexe.

Si $\mathcal{O}''_y$ est la clôture intégrale de $\mathcal{O}_y$, la fermeture intégrale $\mathcal{O}'_y$ de
$\mathcal{O}_y$ dans $R(X)$ est aussi celle de $\mathcal{O}''_y$ ; comme le corps
des fractions de $\mathcal{O}''_y$ est radiciel sur $R(Y)$, on sait que si $\mathcal{O}''_y$ est
un anneau local, il en est de même de $\mathcal{O}'_y$ (Bourbaki, Alg. comm., )

Ce dernier corollaire est essentiellement la forme sous laquelle
Zariski énonce son "théorème de connexion" pour les schémas algébri-
ques, en supposant toutefois que $\mathcal{O}_y$ est analytiquement irréductible.

Corollaire 5. - La conclusion du cor. 4 reste valable si, au lieu de
supposer $Y$ unibranche en $y$, on suppose que $f$ est universellement ou-
vert $(I, 7, 5)$ au-dessus de $y$.

\newpage

%% ===== Page 409 =====

--- III-81 ---

Nous utiliserons la conséquence suivante du lemme de Hironaka, qui sera démontrée plus tard (§ ) : $Y$ étant noethérien et intègre, pour tout point $y \in Y$ il existe un voisinage ouvert affine $U$ de $y$ dans $Y$, un anneau local noethérien normal $A$, dont le corps des fractions est $R(Y)$, qui domine $\mathcal{O}_y$ et qui est de la forme $B_{\mathfrak{p}}$ où $B$ est une algèbre de type fini sur l'anneau $C=A(U)$, et est contenue dans $R(Y)$, et $\mathfrak{p}$ un idéal premier de $B$. Soit $Y' = \operatorname{Spec}(B)$ ; l'injection $\psi: \mathcal{O}_y \to B_{\mathfrak{p}}$ définit un morphisme $h=(Y', \theta)$ de type fini d'un ouvert $Y_1$ de $Y'$ contenant $\mathfrak{p}$, dans $U$, tel que $\psi(\mathfrak{p})=y$ et $B_{\mathfrak{p}} / \mathfrak{p} B_{\mathfrak{p}} = k(y)$ (I, 4.3, prop. 15) ; en outre ce morphisme correspond à un homomorphisme injectif de $C$ dans $A(Y')$ donc il est dominant (I, 1.2, cor. 4 de la prop. 5). Comme par ailleurs $Y'$ est intègre et $Y'$ est dense dans $Y$, on voit qu'on a défini un morphisme dominant $h: Y' \to Y$ tel que $h(\mathfrak{p})=y$. En outre, comme le corps des fractions de $B_{\mathfrak{p}}$ est $R(Y)$, $h$ est birationnel. L'hypothèse que $f$ est universellement ouvert au-dessus de $y$ entraîne alors, en posant $X'=X \times_Y Y'$ et $f'=f^{Y'}$, que $X'$ a une seule composante irréductible dont l'image par $f'$ est dense dans $Y'$, et que $f'^{-1}(\mathfrak{p})$ est contenu dans cette composante (I, 7.5, cor. 4 de la prop. 15). On peut en outre remplacer $f'$ par $f'_{\operatorname{red}}$ et supposer donc $X''$ intègre. Alors, comme l'anneau local de $\mathfrak{p}$ est intégralement clos, $\mathfrak{p}$ est unibranche ; en outre, en considérant un voisinage affine du point générique de $Y'$ dans lequel $h$ est un isomorphisme sur un voisinage du point générique de $Y$, on voit aussitôt qu'il existe une application birationnelle de $X''$ dans $X$ ; autrement dit, on a $R(Y')=R(Y)$, $R(X'')=R(X)$, et la fermeture algébrique de $R(Y')$ dans $R(X'')$ est donc radicielle sur $R(Y')$. Les conditions d'application du cor. 4 étant remplies, $f'^{-1}(\mathfrak{p})$ est connexe, et comme la projection de $f'^{-1}(\mathfrak{p})$ dans

\newpage

%% ===== Page 410 =====

$f^{-1}(y)$ est surjective (I, 2, 7, prop. 18 et I, 2, 6, prop. 10), $f^{-1}(y)$ est connexe.

\emph{Remarque}.— Le raisonnement précédent est dû en substance à Zariski (\emph{Theory and applications of Holomorphic functions}), à cela près qu'il peut prendre le normalisé de $\mathcal{O}_y$, puisque, pour un anneau local de la géométrie algébrique, on sait que son normalisé est noethérien. D'autre part, Zariski prouve que si $Y$ est la variété de Chow d'un espace projectif $\mathbb{P}^n_k$ sur un corps $k$, et si $X$ est la partie fermée de $\mathbb{P}^n_k \times Y$ qui définit la correspondance de Chow entre $\mathbb{P}^n$ et $Y$, alors la projection $X \to Y$ est un morphisme universellement ouvert (\emph{loc. cit.}, lemme de la p. 82). Il semble bien que ce soit la seule propriété formelle des "coordonnées de Chow" dont on se serve dans beaucoup d'applications ; il y a par suite intérêt, dans une telle situation, à substituer le langage : fibres d'un morphisme propre (éventuellement supposé universellement ouvert, ou soumis à d'autres restrictions de régularité locale) au langage : spécialisation de cycles de l'espace projectif.

\emph{Corollaire} 6.— Sous les conditions du cor. 3, supposons $R(Y)$ algébriquement fermé dans $R(X)$, et soit $y$ un point normal de $Y$. Alors $f^{-1}(y)$ est connexe, et il existe un voisinage ouvert $U$ de $y$ tel que $A(f^{-1}(U)) = \mathcal{O}_y|_U$. En particulier, si $Y$ est normale (et $R(Y)$ algébriquement fermé dans $R(X)$), on a $f_*(\mathcal{O}_X) = \mathcal{O}_Y$ et $f^{-1}(y)$ est connexe pour tout $y \in Y$.

La première assertion résulte du cor. 3, puisque l'hypothèse entraîne que $\mathcal{O}_y$ est intégralement fermé dans $R(X)$. D'autre part, si $f : X \to Y' \xrightarrow{f'} Y$ est la factorisation de Stein de $f$ (th. 1), on a $R(Y) \subset R(Y') \subset R(X)$, et $R(Y')$ est algébrique sur $R(Y)$, donc égal à $R(Y)$ par hypothèse ; si $y'$ est l'unique point (cor. 2) de $Y'$ au-des-

\newpage

%% ===== Page 411 =====

--- III-83 ---

Le fait que le cor.4 est établi dans le cadre des schémas permet des applications telles que la suivante :

\emph{Corollaire 7.4} Soient $A$ un anneau local noethérien unibranche, $k$ son corps résiduel, $S = \operatorname{gr}(A)$ l'anneau gradué associé à $A$, considéré comme $k$-algèbre graduée spéciale. Alors $\operatorname{Proj}(S)$ est un $k$-schéma connexe.

Soient $\mathfrak{m}$ l'idéal maximal de $A$, $Y = \operatorname{Spec}(A)$, $X = \operatorname{Proj}(S')$ le $Y$-schéma obtenu en faisant éclater l'idéal $\mathfrak{m}$ ; on a $S' = \bigoplus_{n \ge 0} \mathfrak{m}^n$ (§ 2, n° 4). On applique le cor.4 au morphisme structural $f : X \to Y$ qui est projectif, et à l'unique point fermé de $Y$ ; on a $k(f^{-1}(y)) = k$, et $f^{-1}(y)$ s'identifie à $\operatorname{Proj}(S)$ (II, 3, 5, prop. 20) ; comme $k(f^{-1}(y))$ est le corps des fractions du point générique de $f^{-1}(y)$, on voit que c'est un corps de fractions rationnelles sur $k = R(\{y\})$ ; on en conclut aussitôt que $R(X)$ ne contient pas d'élément algébrique sur $R(Y)$. Les conditions du cor.4 sont par suite remplies, d'où la conclusion.

Le résultat suivant est dû à \marginpar{due à Serre} Chevalley dans le cas des schémas algébriques ; nous en donnerons au § une démonstration plus élémentaire n'utilisant que le th. de finitude (mais utilisant le critère de Serre).

\emph{Proposition 1} Soient $Y$ un préschéma localement noethérien, $f : X \to Y$ un morphisme. Les conditions suivantes sont équivalentes :

a) $f$ est fini.

b) $f$ est affine et propre.

c) $f$ est propre et pour tout $y \in Y$, l'ensemble $f^{-1}(y)$ est fini.

\newpage

%% ===== Page 412 =====

\noindent - III-84 -

On a déjà démontré l'équivalence de a) et b) (II, 5, 1, cor. 2 de la prop. 3); il est clair que a) entraîne c), et il reste donc à prouver que c) entraîne a). En considérant la factorisation de Stein $f=g \circ f'$ de $f$, on est ramené à prouver que $g$ est fini, puisque $f'$ l'est par construction (II, 2, 7, prop. 15 (ii)) ; on peut donc supposer en plus que $f_*(\mathcal{O}_X) = \mathcal{O}_Y$. En vertu du th. de connexion, les fibres $f^{-1}(y)$ sont alors connexes ; d'autre part, $f^{-1}(y)$ étant un schéma algébrique sur $\kappa(y)$ dont l'espace de base est fini, $f^{-1}(y)$ est discret (I, 4, 2, prop. 12), donc $f^{-1}(y)$ est réduit à un point, autrement dit $f$ est injectif. Comme $f$ est une application fermée et que $f_*(\mathcal{O}_X) = \mathcal{O}_Y$, $f$ est surjective, donc un homéomorphisme de $X$ sur $Y$, et la relation $f_*(\mathcal{O}_X) = \mathcal{O}_Y$ entraîne enfin que $f$ est un isomorphisme, ce qui achève la démonstration.

Corollaire. - Soient $X, Y$ deux préschémas intègres ; on suppose en outre $Y$ normal et localement noethérien. Si $f: X \to Y$ est un morphisme propre, birationnel et tel que $f^{-1}(y)$ soit fini pour tout $y \in Y$, $f$ est un isomorphisme.

Comme les hypothèses du cor. 4 du th. 1 sont remplies pour tout $y \in Y$, on en conclut que $f^{-1}(y)$ est connexe ; donc on voit comme dans la prop. 1 que $f$ est un homéomorphisme (puisque $f$ est dominant, étant birationnel). En outre, pour tout $y \in Y$, si $x$ est l'unique point de $X$ tel que $f(x)=y$, $\mathcal{O}_x$ est entier sur $\mathcal{O}_y$ et a même corps des fractions en vertu du prop. 1, donc $\mathcal{O}_x = \mathcal{O}_y$, ce qui achève que $f$ est un isomorphisme.

\underline{Zariski} [barré : le "main theorem" de Zariski]

Proposition 2. - Soient $Y$ un préschéma localement noethérien, $f: X \to Y$ un morphisme propre. Soit $X'$ l'ensemble des points $x \in X$ qui sont isolés dans leur fibre $f^{-1}(f(x))$. Alors l'ensemble $X'$ est ouvert.

\newpage

%% ===== Page 413 =====

- III-85 -

vert dans $X$, et si $f=g f'$ est la factorisation de Stein de $f$, la restriction de $f'$ à $X'$ est un isomorphisme de $X'$ sur un sous-préschéma de $Y'$ induit sur un ouvert $U \subset Y'$, et on a $X' = f'^{-1}(U)$.

Comme $f'^{-1}(f'(x)) \subset f^{-1}(f(x))$, tout $x \in X'$ est isolé dans $f^{-1}(f(x))$ et on peut par suite se borner au cas où $f'=f$, donc $f_*(\mathcal{O}_X) = \mathcal{O}_Y$ ; on en conclut que si $x \in X'$, $f^{-1}(f(x))$, qui est connexe (n°1, cor. 1 du th. 1), est réduite au point $x$. Comme $f$ est fermée, pour tout voisinage ouvert $V$ de $x$ dans $X$, $f(X-V)$ est fermé dans $Y$ et ne contient pas $y=f(x)$, puisque $f^{-1}(y)=\{x\}$ ; si $U$ est son complémentaire dans $Y$, on a donc $f^{-1}(U) \subset V$, d'où on conclut que les images réciproques par $f$ d'un système fondamental de voisinages ouverts de $y$ forme un système fondamental de voisinages ouverts de $x$. L'hypothèse $f_*(\mathcal{O}_X) = \mathcal{O}_Y$ et la définition de l'image directe d'un faisceau entraînent alors que, si $f(\mathcal{O}_X) = \mathcal{O}_Y$, l'application $\mathcal{O}_{Y,y} \to \mathcal{O}_{X,x}$ est un isomorphisme. On en conclut qu'il existe un voisinage ouvert $V$ de $x$ et un voisinage ouvert $U$ de $y$ tel que la restriction de $f$ à $V$ soit un isomorphisme de $V$ sur $U$ (I, 4.3, prop. 16) ; en outre, d'après ce qu'on vient de voir, on peut supposer que $f^{-1}(U)=V$, d'où on conclut aussitôt, par définition, que $V \subset X'$, et cela achève la démonstration.

Théorème 2 ("Main theorem" de Zariski). — Soient $Y$ un préschéma noethérien, $f: X \to Y$ un morphisme quasi-projectif, $X'$ l'ensemble des points $x \in X$ qui sont isolés dans leur fibre $f^{-1}(f(x))$. Alors $X'$ est une partie ouverte de $X$, et le sous-préschéma induit $X'$ est isomorphe à un préschéma induit sur une partie ouverte d'un $Y$-préschéma $Y'$ fini sur $Y$.

Nous avons déjà démontré que $X'$ est ouvert, en utilisant un lemme de Nagata (I, 7.3, cor. 1 du th. 2), mais nous en obtenons ici une démonstration directe. Par hypothèse, il existe un $Y$-préschéma pro-

\newpage

%% ===== Page 414 =====

- III-86 -

jectif $Z$ tel que $X$ soit $Y$-isomorphe à un sous-préschéma induit sur un ouvert de $Z$ (II,4,7,prop.24). On est donc ramené à démontrer le théorème pour $Z$, et il résulte alors aussitôt de la prop.2.

Corollaire 1. - Soient $Y$ un schéma localement noethérien, $f: X \to Y$ un morphisme de type fini, $x$ un point de $X$ isolé dans sa fibre $f^{-1}(f(x))$. Alors il existe un voisinage ouvert de $x$ dans $X$ qui est isomorphe à une partie ouverte d'un $Y$-préschéma fini sur $Y$.

Soient en effet $y=f(x)$, $U$ un voisinage affine de $y$ dans $Y$ et $V$ un voisinage affine de $x$ dans $X$, contenu dans $f^{-1}(U)$. Comme $Y$ est séparé, $U$ est affine sur $Y$ (I,3,7,cor.3 de la prop.22), et comme $V$ est affine sur $U$, $V$ est affine sur $Y$ (II,2,3,cor.3 de la prop.7), et a fortiori quasi-projectif sur $Y$ (II,4,7,prop.27). Il suffit alors d'appliquer à $V$ le th.2.

Lorsque $Y$ est affine, le cor.1 s'énonce en termes d'algèbre commutative :

Corollaire 2. - Soient $A$ un anneau noethérien, $B$ une $A$-algèbre de type fini, $\mathfrak{q}$ un idéal premier de $B$, $\mathfrak{p}$ son image réciproque dans $A$. On suppose que $\mathfrak{q}$ soit à la fois maximal et minimal dans l'ensemble des idéaux premiers de $B$ dont l'image réciproque est égale à $\mathfrak{p}$. Alors il existe $g \in B-\mathfrak{q}$, une $A$-algèbre finie $A'$ et un élément $f' \in A'$ tels que les $A$-algèbres $B_g$ et $A'_{f'}$ soient isomorphes.

Il suffit en effet d'appliquer le cor.1 à $Y=\operatorname{Spec}(A)$ et $X=\operatorname{Spec}(B)$, l'hypothèse sur $\mathfrak{q}$ signifiant exactement qu'il est isolé dans sa fibre (I,1,2).

On en déduit le résultat suivant, moins général en apparence :

Corollaire 3. - Soient $A$ un anneau local noethérien, $B$ une $A$-algèbre de type fini, $\mathfrak{n}$ un idéal premier de $B$ dont l'image réciproque dans $A$ est l'idéal maximal $\mathfrak{m}$. On suppose que $\mathfrak{n}$ est maximal dans

\newpage

%% ===== Page 415 =====

\marginpar{III-87}

$B$ et est minimal dans l'ensemble des idéaux premiers de $B$ dont l'image réciproque est $\mathfrak{m}$ (ce qui signifie aussi que $\mathfrak{B}_{\mathfrak{n}}$ est primaire pour $\mathfrak{n}$). Alors il existe une $A$-algèbre finie $A'$ et un idéal maximal $\mathfrak{m}'$ de $A'$ (dont $\mathfrak{m}$ est l'image réciproque dans $A$) tels que $B_{\mathfrak{n}}$ soit isomorphe à la $A$-algèbre $A'_{\mathfrak{m}'}$.

C'est la forme donnée par Zariski à son "Main theorem", étendue aux anneaux locaux noethériens quelconques.

Le cas particulier suivant du cor. 3 est aussi parfois appelé "Main Theorem" :

Corollaire 4. - Sous les conditions du cor. 3, supposons en outre $A$ et $B$ intègres et ayant même corps des fractions $K$. Alors, si $A$ est normal, on a $B=A$.

En effet, la Remarque suivant le th. 2 montre que l'on peut supposer dans l'application du cor. 3 que $A'$ est intègre et a même corps des fractions $K$ que $A$ et $B$ ; l'hypothèse sur $A$ entraîne alors que $A'=A$, donc $B_{\mathfrak{n}}=A$ ; comme $B$ domine $A$ et que l'on a $A \subset B \subset B_{\mathfrak{n}}$, on en conclut bien $A=B$.

Les corollaires précédents étaient des variantes de nature locale du th. 2, qui est un résultat global. Voici une autre conséquence de nature globale :

\newpage

%% ===== Page 416 =====

III-87 -

$B$ est minimal dans l'ensemble des idéaux premiers de $B$ dont l'image réciproque est $\mathfrak{m}$ (ce qui signifie aussi que $B_{\mathfrak{m}}$ est primaire pour $\mathfrak{n}$). Alors il existe une $A$-algèbre finie $A'$ et un idéal maximal $\mathfrak{m}'$ de $A'$ (dont $\mathfrak{m}$ est l'image réciproque dans $A$) tels que $B_{\mathfrak{m}}$ soit isomorphe à la $A$-algèbre $A'_{\mathfrak{m}'}$.

C'est la forme donnée par Zariski à son "Main Theorem", étendue aux anneaux locaux noethériens quelconques.

Le cas particulier suivant du cor. 3 est aussi parfois appelé "Main Theorem" :

Corollaire 4.- Sous les conditions du cor. 3, supposons en outre $A$ et $B$ intègres et ayant même corps des fractions. Alors, si $A$ est normal, on a $B=A$.

En effet, la remarque suivant le th. 2 montre que l'on peut supposer dans l'application du cor. 3 que $A'$ est intègre et a même corps des fractions $K$ que $A$ et $B$; l'hypothèse sur $A$ entraîne alors que $A'=A$, donc $B_{\mathfrak{n}}=A$; comme $B$ domine $A$ et que l'on a $A \subset B \subset B_{\mathfrak{n}}$, on en conclut bien $A=B$.

Les corollaires précédents étaient des variantes de nature locale du th. 2, qui est un résultat global. Voici une autre conséquence de nature globale :

\newpage

%% ===== Page 417 =====

-- III-89 --

(iii) Si $X$ et $Y$ sont des $S$-préschémas, $f: X \to Y$ un $S$-morphisme quasi-fini, alors pour tout morphisme $S' \to S$, $f^{S'}: X^{S'} \to Y^{S'}$ est quasi-fini.

(iv) Si $f: X \to Y$ et $g: X' \to Y'$ sont deux $S$-morphismes quasi-finis, $f \times g: X \times_S X' \to Y \times_S Y'$ est quasi-fini.

(v) Si $f: X \to Y$ et $g: Y \to Z$ deux morphismes tels que $g \circ f$ soit quasi-fini ; si en outre $g$ est séparé, ou $X$ noethérien, ou $X \times_Z Y$ localement noethérien, $f$ est quasi-fini.

(vi) Si $f$ est quasi-fini, il en est de même de $f_{\text{red}}$.

L'assertion (i) est triviale, toute fibre étant réduite à un point, compte tenu de (I, 4, 2, prop. 7). De même, pour prouver (ii), on remarque que $h = g \circ f$ est de type fini (I, 4, 2, prop. 8) ; en outre, si $z = h(x)$ et $y = f(x)$, $y$ est isolé dans $g^{-1}(z)$, donc il existe un voisinage ouvert $V$ de $y$ ne rencontrant aucun point de $g^{-1}(z)$ distinct de $y$ ; $f^{-1}(V)$ est donc un voisinage ouvert de $x$ ne rencontrant aucun des $f^{-1}(y')$, où $y'$ est dans $g^{-1}(z)$ ; comme $x$ est isolé dans $f^{-1}(y)$, il est isolé dans $h^{-1}(z) = f^{-1}(g^{-1}(z))$. Pour prouver (iii), on remarque encore tout d'abord que $f^{S'}$ est de type fini (I, 4, 2, prop. 9) ; d'autre part, si $x' \in X^{S'}$ et si on pose $y' = f^{S'}(x')$ et $y = g(y')$, $f^{-1}(y')$ s'identifie à $f^{-1}(y) \otimes_{\kappa(y)} \kappa(y')$ (I, 2, 7, prop. 18) ; comme $f^{-1}(y)$ est un $\kappa(y)$-schéma algébrique de rang fini sur $\kappa(y)$ (I, 4, 2, prop. 12), $f^{-1}(y')$ est de rang fini sur $\kappa(y')$, et $f^{-1}(y')$ est par suite discret. Les assertions (iv), (v), (vi) se déduisent de (i), (ii), (iii) suivant le procédé habituel, sauf en ce qui concerne les hypothèses de (v) sur $g$ autres que l'hypothèse que $g$ est séparé ; il faut alors utiliser (I, 4, 2, cor. 2 de la prop. 9) et la remarque que si $x$ est isolé dans $f^{-1}(g^{-1}(g(f(x))))$, il

\newpage

%% ===== Page 418 =====

l'est a fortiori dans $f^{-1}(f(x))$.

\subsection*{3.2 Morphismes équidimensionnels et lemme de normalisation.}

Nous rappelons pour mémoire l'énoncé du lemme de normalisation de E. Noether :

\textbf{Lemme 1.} --- Soient $k$ un corps, $A$ une algèbre de type fini sur $k$, de dimension $e$. Alors il existe $e$ éléments $x_i \in A$ ($1 \le i \le e$), algébriquement indépendants et tels que $A$ soit entière sur $k[x_1, \dots, x_e]$.

L'existence des $x_i$ est démontrée par exemple dans Sém. Cartan-Chevalley 1955-56, 4-01 à 4-03. Le fait que le nombre de ces éléments soit la dimension de $A$ résulte de ce que $A' = k[x_1, \dots, x_e]$ est de dimension $e$ (I, 7.3, cor. 1 de la prop. 9) et de ce que $\dim A = \dim A'$ (II, 2.7, prop. 16 bis).

\textbf{Lemme 2.} --- Soient $A$ un anneau noethérien, $B$ une $A$-algèbre de type fini, $Y = \operatorname{Spec}(A)$, $X = \operatorname{Spec}(B)$, $f$ le morphisme structural $X \to Y$. Soit $y \in Y$, et soit $C = B \otimes_A k(y)$ l'anneau de la fibre $f^{-1}(y)$ ; soient $f_i$ ($1 \le i \le e$) des éléments de $B$ dont les images canoniques $f_i \otimes 1$ dans $C$ engendrent une sous-$k(y)$-algèbre $C'$ de $C$ telle que $C$ soit finie sur $C'$. Alors il existe un voisinage ouvert $U$ de $f^{-1}(y)$ dans $X$, quasi-fini sur $Y' = \operatorname{Spec}(A[T_1, \dots, T_e])$, et a fortiori sur $Y \otimes_{\mathbb{Z}} \mathbb{Z}[T_1, \dots, T_e]$ ($T_i$ indéterminées).

En effet, $f$ se factorise en $X \to Y' \to Y$. Par hypothèse, $f^{-1}(y)$ est fini sur $g^{-1}(y)$, dont l'anneau est $A[T_1, \dots, T_e] \otimes_A k(y)$; donc tout $x \in f^{-1}(y)$ est isolé dans sa fibre $f^{-1}(f(x))$ (II, 2.7, prop. 17). Comme $X$ est affine et $Y'$ noethérien, le lemme résulte de l'application à $f'$ du th. 2 (il est clair en effet que $Y'$ est un sous-préschéma fermé, donc fini, de $Y \otimes_{\mathbb{Z}} \mathbb{Z}[T_1, \dots, T_e]$).

\textbf{Corollaire.} --- Avec les notations du lemme 2, pour tout $z \in Y$, la dimension de $U \cap f^{-1}(z)$ est $\le e$.

\newpage

%% ===== Page 419 =====

- III-91 -

En effet, avec les mêmes notations, $U \cap f^{-1}(z)$ est fini sur $g^{-1}(z)$ donc $\dim(U \cap f^{-1}(z)) \le \dim(g^{-1}(z))$ (prop.4), et comme $g^{-1}(z)$ est un sous-préschéma fermé du spectre affine de $\kappa(z)[T_1, \dots, T_e]$, on a $\dim(g^{-1}(z)) \le e$ (I, 7, 1, prop. 1 et I, 7, 3, cor. 1 de la prop. 9).

La conjonction des lemmes 1 et 2 donne la généralisation suivante du lemme 1 (qui sera précisée ci-dessous dans le th. 4) :

Proposition 5. - Soient $A$ un anneau noethérien, $B$ une $A$-algèbre de type fini, $Y = \operatorname{Spec}(A)$, $X = \operatorname{Spec}(B)$, $f : X \to Y$ le morphisme structural ; soit $y \in Y$, et soit $e = \dim(f^{-1}(y))$. Il existe alors un voisinage ouvert $U$ de $f^{-1}(y)$ et des éléments $f_i$ ($1 \le i \le e$) de $B$ tels que $U$ soit quasi-fini sur $Y' = \operatorname{Spec}(A[f_1, \dots, f_e])$ et a fortiori sur $Y \otimes_{\mathbb{Z}} \mathbb{Z}[T_1, \dots, T_e]$.

Joint au cor. du lemme 2, cela redonne le th. 2 de I, 7, 3 :

Théorème 3. - Soient $Y$ un préschéma localement noethérien, $f : X \to Y$ un morphisme de type fini. Alors pour tout entier $e \ge 0$, l'ensemble $U_e$ des $x \in X$ tels que $\dim_x(f^{-1}(f(x))) \le e$ est ouvert.

On peut en effet supposer $Y$ affine et noethérien, et (en restreignant $X$ à un voisinage ouvert affine convenable de $x$) on peut supposer, pour $x \in U_e$, que toutes les composantes irréductibles de $f^{-1}(f(x))$ sont de dimension $\le e$ ; la prop. 5 entraîne alors aussitôt la conclusion.

On retrouve comme corollaire le lemme de Nagata qui avait servi au chap. I à démontrer le th. 2 de I, 7, 3 :

Corollaire. - Sous les hypothèses du th. 3, supposons en outre $X$ et $Y$ irréductibles et $f$ dominant. Soient $\xi$ et $\eta$ les points génériques de $X$ et $Y$, et soit $e = \dim f^{-1}(\eta) = \deg.tr._{\kappa(\eta)} \kappa(\xi)$. Alors toute composante irréductible de toute fibre $f^{-1}(y)$ a une dimension $\ge e$.

En effet, pour $n < e$, l'ensemble des $x \in X$ tels que $\dim_x(f^{-1}(f(x))) \le n$ est ouvert et ne peut contenir $\xi$, donc est nécessairement vide.

\newpage

%% ===== Page 420 =====

- III-92 -

\noindent Th\(\acute{\text{e}}\)or\(\grave{\text{e}}\)me 4. - Soient \(Y\) un pr\(\acute{\text{e}}\)sch\(\acute{\text{e}}\)ma localement noeth\(\acute{\text{e}}\)rien, \(f: X \to Y\) un morphisme de type fini. Pour tout \(x \in X\), d\(\acute{\text{e}}\)signons par \(Y_i\) (\(1 \le i \le m\)) les composantes irr\(\acute{\text{e}}\)ductibles de \(Y\) contenant \(y=f(x)\), et par \(y_i\) le point g\(\acute{\text{e}}\)n\(\acute{\text{e}}\)rique de \(Y_i\). Les conditions suivantes sont \(\acute{\text{e}}\)quivalentes :

(i) Il existe un voisinage ouvert \(U\) de \(x\) tel que toute composante irr\(\acute{\text{e}}\)ductible de \(U\) domine l'une des \(Y_i\), et que pour tout \(z \in Y\), les composantes irr\(\acute{\text{e}}\)ductibles de \(U \cap f^{-1}(z)\) aient toutes une m\(\hat{\text{e}}\)me dimension et ind\(\acute{\text{e}}\)pendante de \(z\).

(ii) Les conditions de (i) sont remplies en prenant seulement \(z=y_i\) (\(1 \le i \le m\)) et \(z=y\).

(iii) Il existe un entier \(e \ge 0\), un voisinage affine \(V\) de \(y\), un sch\(\acute{\text{e}}\)ma \(X'\) fini sur \(Y' = V \otimes_{\mathbb{Z}} \mathbb{Z}[T_1, \dots, T_e]\), dont toute composante irr\(\acute{\text{e}}\)ductible domine une composante irr\(\acute{\text{e}}\)ductible de \(Y'\), et un voisinage ouvert \(U\) de \(x\) isomorphe \(\grave{\text{a}}\) un ouvert de \(X'\).

Il est trivial que (i) entra\(\hat{\text{i}}\)ne (ii), prouvons que (ii) entra\(\hat{\text{i}}\)ne (iii). Soit \(V\) un voisinage ouvert affine de \(y\) dans \(Y\), ne rencontrant aucune des composantes irr\(\acute{\text{e}}\)ductibles de \(Y\) ne contenant pas \(y\). Les composantes irr\(\acute{\text{e}}\)ductibles de \(V\) sont donc les \(V_i = V \cap Y_i\) (\(1 \le i \le m\)), et \(y_i\) est encore le point g\(\acute{\text{e}}\)n\(\acute{\text{e}}\)rique de \(V_i\). On peut \(\acute{\text{e}}\)videmment supposer dans la condition (ii) que \(U\) est affine et contenu dans \(f^{-1}(V)\); en vertu de la prop. 5, on peut en outre supposer que \(U\) est isomorphe \(\grave{\text{a}}\) un ouvert dans un sch\(\acute{\text{e}}\)ma \(X'\) fini sur \(Y' = V \otimes_{\mathbb{Z}} \mathbb{Z}[T_1, \dots, T_e]\). En outre, en rempla\(\grave{\text{c}}\)ant au besoin \(X'\) par le plus petit sous-pr\(\acute{\text{e}}\)sch\(\acute{\text{e}}\)ma ferm\(\acute{\text{e}}\) de \(X'\) majorant \(U\) (qui est fini sur \(X'\), donc sur \(Y'\)) (II, 1, 2, prop. 8), on peut supposer \(U\) dense dans \(X'\); les composantes irr\(\acute{\text{e}}\)ductibles de \(X'\) sont donc les adh\(\acute{\text{e}}\)rences dans \(X'\) des composantes irr\(\acute{\text{e}}\)ductibles \(U_j\) de \(U\) (\(1 \le j \le n\)). Par ailleurs, les composantes irr\(\acute{\text{e}}\)ductibles de \(Y'\) sont les \(Y'_i = V_i \otimes_{\mathbb{Z}} \mathbb{Z}[T_1, \dots, T_e]\). Il faut prouver que tout

\newpage

%% ===== Page 421 =====

$U_j$ domine un $Y'_i$. Par hypothèse, il existe $i$ tel que $U_j$ domine $Y_i$, autrement dit, si $x_j$ est le point générique de $U_j$, $f(x_j)=y_i$ ; par suite (I, 7, 5, lemme 7) $x_j$ est aussi le point générique de $U_j \cap f^{-1}(y_i)$ que nous noterons pour abréger $U_j(y_i)$, et qui est égal à $U_j \times_{\operatorname{Spec}(k(y_i))} \operatorname{Spec}(k(y_i))$. De même, le point générique $y'_i$ de $Y'_i$ est aussi celui de $Y'_i(y_i) = Y'_i \times_{\operatorname{Spec}(k(y_i))} \operatorname{Spec}(k(y_i))$, qui est d'ailleurs aussi égal à $Y'_i \times_{\operatorname{Spec}(k(y_i))} \operatorname{Spec}(k(y_i))$, puisque $V$ est un voisinage de $y_i$ dans $Y$. Le morphisme quasi-fini $U \to Y'$ définit donc un morphisme quasi-fini $U_j(y_i) \to Y'_i(y_i)$ (n°2, prop. 4) ; si ce morphisme $g$ n'était pas dominant, [unclear: g_red] considéré comme un morphisme de $(U_j(y_i))_{\operatorname{red}}$ dans un sous-préschéma fermé de $Y'_i(y_i)$, dont l'espace sous-jacent est distinct de $Y'_i(y_i)$ (I, 3, 4, cor. de la prop. 9). Comme $g_{\operatorname{red}}$ est quasi-fini, $Y'_i(y_i)$ intègre et de dimension $e$ (I, 7, 3, cor. 1 de la prop. 9), on aurait $\dim(U_j(y_i)) < e$ (n°2, prop. 4), ce qui contredit l'hypothèse, puisque $U_j(y_i)$ est une composante irréductible de $U \cap f^{-1}(y_i)$. On en conclut que $g(x_j) = y'_i$, ce qui établit (iii).

Il reste à montrer que (iii) entraîne (i). On peut évidemment, en remplaçant $V$ par un voisinage plus petit, supposer que $V$ ne rencontre aucune des composantes irréductibles de $Y$ ne contenant pas $y$. Gardant les mêmes notations que ci-dessus, il résulte de (iii) que tout $U_j$ domine un $Y'_i$, puisque $U_j \to Y'_i$ est composé de $U_j \to Y'_i$, qui est dominant par hypothèse, et de $Y'_i \to Y_i$, qui est dominant par construction. Pour tout $z \in Y$, $U \cap f^{-1}(z)$ est réunion des $U_j(z) = U_j \cap f^{-1}(z)$, donc ses composantes irréductibles sont parmi celles des $U_j(z)$, et il suffit de prouver que ces dernières sont de dimension $e$. Or, d'après le cor. du th. 3, appliqué au morphisme dominant

\newpage

%% ===== Page 422 =====

- III-94 -

$U_j \to Y_i$, les composantes irréductibles de $U_j(z)$ ont toutes une dimension $\ge e$. D'autre part, comme $U \cap f^{-1}(z) = U \times_Y \operatorname{Spec}(k(z))$ (n°2, prop.3), et que ce dernier n'est autre que $\operatorname{Spec}(k(z)[T_1, \dots, T_e])$, donc est de dimension $e$, on a $\dim(U \cap f^{-1}(z)) \le e$ (n°2, prop.4), et a fortiori les composantes irréductibles de $U_j(z)$ ont une dimension $\le e$, ce qui achève la démonstration.

Définition 2. — On dit qu'un morphisme $f$ est équidimensionnel en $x$ s'il vérifie les conditions équivalentes du th.4. On dit que $f$ est équidimensionnel (resp. équidimensionnel au point $y \in Y$) s'il est équidimensionnel en tout $x \in X$ (resp. en tout $x \in f^{-1}(y)$). On dit que $f$ est strictement équidimensionnel s'il est équidimensionnel et si, pour tout $y \in Y$, $f^{-1}(y)$ a une dimension indépendante de $y$ (et par suite est non vide si $X$ est non vide).

Corollaire 1. — Soient $Y$ un préschéma noethérien irréductible, $X$ un préschéma irréductible, $f : X \to Y$ un morphisme dominant de type fini ; soit $e$ la dimension de la fibre générique de $f$. Soit $x \in X$ tel que $\mathcal{O}_{f(x)}$ soit quotient d'un anneau régulier. Les conditions suivantes sont équivalentes :

(i) $f$ est équidimensionnel en $x$.

(ii) $\dim_x(f^{-1}(f(x))) \le e$ (ce qui entraîne $\dim_x(f^{-1}(f(x))) = e$ par le cor. du th.3).

(iii) Pour toute composante irréductible $Z_i$ de $f^{-1}(f(x))$ contenant $x$, de point générique $x_i$, on a $\dim(\mathcal{O}_{x_i}) \ge \dim(\mathcal{O}_{f(x)})$ (ce qui entraîne $\dim(\mathcal{O}_{x_i}) = \dim(\mathcal{O}_{f(x)})$ par I,7,2, prop.8).

(iv) Il existe un voisinage ouvert $U$ de $x$ et un voisinage ouvert $V$ de $f(x)$ tel que $U$ soit quasi-fini au-dessus de $Y' = V \otimes_{\mathbb{Z}} [T_1, \dots, T_e]$.

L'équivalence de (i) et (iv) résulte du th.4, et il en est de même

\newpage

%% ===== Page 423 =====

- III-95 -

me de l'équivalence de (i) et (ii), car en vertu du cor. du th.3, toutes les composantes irréductibles de $f^{-1}(f(x))$ contenant $x$ ont nécessairement une dimension égale à $e$ si (ii) a lieu, et cela entraîne la condition (ii) du th.4. L'équivalence de (ii) et (iii) résulte de l'hypothèse sur $\mathcal{O}_{f(x)}$, qui entraîne l'égalité dans la "formule de dimension"
$$e + \dim(\mathcal{O}_{f(x)}) = \text{deg.tr.}_{\kappa(f(x))} \kappa(x_1) + \dim(\mathcal{O}_{x_1})$$
(I,7,2, cor.1 du th.1) (on notera que sans l'hypothèse sur $\mathcal{O}_{f(x)}$, l'inégalité de la formule de dimension prouve que (iii) entraîne (ii)).

On notera que dans le résultat de (iv), $U$ est par définition isomorphe à un ouvert d'un schéma $X'$ fini sur $Y'$; on peut naturellement supposer $X'$ irréductible, en remplaçant $X'$ par le plus petit sous-schéma fermé de $X'$ majorant $U$; si en outre $X$ est réduit, on peut aussi supposer que $X'$ est réduit (en remplaçant le morphisme $g: X' \to Y'$ par $g_{red}$). Enfin, il est clair que $g$ est toujours dominant.

Corollaire 2. - Sous les hypothèses du cor.1 pour $X, Y$ et $f$, pour que $f$ soit équidimensionnel au point $y \in Y$ (resp. équidimensionnel), il faut et il suffit que toutes les composantes irréductibles de $f^{-1}(y)$ (resp. pour tout $y \in Y$) soient de dimension $e$.

On notera que l'hypothèse que $f$ est équidimensionnel n'exclut pas le cas où certaines fibres $f^{-1}(y)$ seraient vides ; si on suppose en outre $f$ surjectif, alors $f$ est strictement équidimensionnel.

Corollaire 3. - Sous les hypothèses du cor.1 pour $X, Y$ et $f$, supposons $f$ équidimensionnel au point $y$ (resp. équidimensionnel). Alors si on pose $S = Z[T_1, \dots, T_r]$, $X^S = X \times_Z S$, $Y^S = Y \times_Z S$, le morphisme $f^S : X^S \to Y^S$ est équidimensionnel en tout point de $Y^S$ au-dessus de $y$ (resp. équidimensionnel).

Il est immédiat que $X^S$ et $Y^S$ sont localement noethériens et irréduc-

\newpage

%% ===== Page 424 =====

- III-96 -

tibles ; si $\xi, \eta, \eta'$ sont les points génériques de $X, Y,$ et $Y^S$, et $g$ le morphisme $Y^S \to Y$, on a $g(\eta') = \eta$, $f(\xi) = \eta$, donc (I, 2, 6, pro. 10), en posant $f'=f^S$, $f'^{-1}(\eta')$ n'est pas vide, et par suite $f'$ est dominant. En outre, $k(\eta') = k(\eta)(T_1, \dots, T_r)$ ; de même, si $\xi'$ est le point générique de $X^S$ on a $k(\xi') = k(\xi)(T_1, \dots, T_r)$, donc la fibre générique $f'^{-1}(\eta')$ a encore pour dimension $e$. Il suffit alors de vérifier le critère (ii) du cor. 1, savoir que pour tout $\eta'$ tel que $g(\eta') = \eta$ on a (I, 2, 7, prop. 18) $\dim(f'^{-1}(\eta')) \le e$. Or, on a $f'^{-1}(\eta') = f^{-1}(\eta) \otimes_{k(\eta)} k(\eta')$ donc $\dim f'^{-1}(\eta') = \dim(f^{-1}(\eta))$ (I, 7, 2, cor. 4 de la prop. 7), d'où la conclusion.

4. Morphismes universellement ouverts.

On a défini au n°1 la notion d'anneau local (intègre) unibranche.
Définition 3. — On dit qu'un anneau local intègre $A$ est \emph{géométriquement unibranche} si $A$ est unibranche (donc sa clôture intégrale $A'$ est un anneau local) et si de plus le corps résiduel de $A'$ est radiciel sur le corps résiduel de $A$.

Dire que $A$ est unibranche (resp. géométriquement unibranche) signifie encore que, pour le morphisme $\operatorname{Spec}(A') \to \operatorname{Spec}(A)$, la fibre du point fermé $y$ de $\operatorname{Spec}(A)$ est réduite à un point (resp. que son nombre géométrique de points (I, 4, 2)) est égal à 1).

Lemme 3. — Pour qu'un anneau local intègre $A$, de corps des fractions $K$, soit unibranche (resp. géométriquement unibranche), il faut et il suffit que tout sous-anneau $A'' \supset A$ de $K$, fini sur $A$, soit un anneau local (resp. soit un anneau local et ait un corps résiduel radiciel sur celui de $A$).

La condition est évidemment nécessaire puisque $A \subset A'' \subset A'$ (et le résultat est même valable en supposant seulement $A''$ entier (mais non

\newpage

%% ===== Page 425 =====

- III-97 -

nécessairement fini sur $A$). Inversement, si $A'$ contient deux idéaux maximaux distincts $\mathfrak{m}', \mathfrak{m}''$ tels que $\mathfrak{m}' \cap A = \mathfrak{m}'' \cap A = \mathfrak{m}$, idéal maximal de $A$, en prenant $x' \in \mathfrak{m}'$ et $x' \notin \mathfrak{m}''$ (resp. $x'' \in \mathfrak{m}''$ et $x'' \notin \mathfrak{m}'$), et en considérant le sous-anneau $A'' = A[x', x'']$, qui est fini sur $A$, $\mathfrak{m}' \cap A''$ et $\mathfrak{m}'' \cap A''$ sont des idéaux maximaux distincts de $A''$, puisque $A'$ est entier sur $A''$.

Si $A'$ est un anneau local, et si $p$ est l'exposant caractéristique de $k$, dire que le corps résiduel de $A'$ est radiciel sur $k$ signifie que pour tout $x \in A' \otimes_k k$, il existe un entier $e \geq 0$, un $y \in k$ et un entier $n > 0$ tels que $(x^{p^e} - y)^n = 0$, et $A'$ possède cette propriété lorsqu'au tout sous-anneau $A'' \supset A$, fini sur $A$, la possède.

Définition 4. — Soient $A$ un anneau local quelconque, $N$ l'idéal de ses éléments nilpotents. On dit que $A$ est unibranche (resp. géométriquement unibranche) si $A/N$ est intègre (autrement dit si $\operatorname{Spec}(A)$ est irréductible) et unibranche (resp. géométriquement unibranche). On dit qu'un point $y$ d'un préschéma $Y$ est unibranche (resp. géométriquement unibranche) si l'anneau local $\mathcal{O}_y$ est unibranche (resp. géométriquement unibranche).

Lemme 4. — Si un point $y$ d'un préschéma $Y$ est géométriquement unibranche, il en est de même de tout point du préschéma $Y' = Y \otimes_{\mathbb{Z}} \mathbb{Z}[T_1, \dots, T_r]$ au-dessus de $y$.

En effet, si $A = \mathcal{O}_y$, on peut se borner au cas où $Y = \operatorname{Spec}(A)$ (I, 2, 7, prop. 19) ; si $N$ l'idéal des éléments nilpotents de $A$, $(A/N)[T_1, \dots, T_r]$ n'a pas d'éléments nilpotents, donc on peut se borner au cas où $A$ est intègre, et il en est de même alors de $B = A[T_1, \dots, T_r]$. Il s'agit alors de montrer tout d'abord que si $\mathfrak{p}$ est un idéal premier de $B$ tel que $\mathfrak{p} \cap A$ soit l'idéal maximal $\mathfrak{m}$ de $A$, la clôture intégrale de $B_{\mathfrak{p}}$ est un anneau local. Or, si $A'$ est la clôture

\newpage

%% ===== Page 426 =====

B' = )
intégrale de $A$, on sait ( ) que $A'[T_1, \dots, T_r]$ est intégralement clos, donc est la clôture intégrale de $B$. Si $S = B \setminus \mathfrak{p}$, la clôture intégrale de $B_{\mathfrak{p}}$ est donc $S^{-1}B'$, et par suite les idéaux premiers de $S^{-1}B'$ dont l'intersection avec $B_{\mathfrak{p}}$ est $\mathfrak{p}B_{\mathfrak{p}}$ correspondent biunivoquement aux idéaux premiers $\mathfrak{p}'$ de $B'$ tels que $\mathfrak{p}' \cap B = \mathfrak{p}$. Or, pour un tel idéal, $\mathfrak{p}' \cap A'$ est nécessairement l'unique idéal maximal $\mathfrak{m}'$ de $A'$; si $k' = A'/\mathfrak{m}'$, les idéaux $\mathfrak{p}'$ considérés correspondent donc biunivoquement aux idéaux premiers $\mathfrak{q}' = \mathfrak{p}'/\mathfrak{m}B'$ de $k'[T_1, \dots, T_r]$ dont l'intersection avec $k[T_1, \dots, T_r]$ (où $k = A/\mathfrak{m}$) est l'idéal $\mathfrak{q} = \mathfrak{p}/\mathfrak{m}B$. Par hypothèse $k'$ est radiciel sur $k$, donc on sait qu'il n'y a qu'un seul idéal premier de $k'[T_1, \dots, T_r]$ au-dessus de $\mathfrak{q}/\mathfrak{m}B$; en outre, le corps des fractions de $B'/\mathfrak{p}'$, égal à celui de $k'[T_1, \dots, T_r]/\mathfrak{q}'$, est radiciel sur le corps des fractions de $k[T_1, \dots, T_r]/\mathfrak{q}$, égal au corps des fractions de $B/\mathfrak{p}$, et cela achève la démonstration.

Lemme 5 (second théorème de Cohen-Seidenberg). — Soient $A$ un anneau intègre, $K$ son corps des fractions, $B$ un anneau intègre contenant $A$, entier sur $A$ et tel que $B \cap K = A$. Alors, pour tout idéal premier $\mathfrak{p}'$ de $B$, tel que $\mathfrak{p}' \cap A = \mathfrak{p}$, et tout idéal premier $\mathfrak{q} \subset \mathfrak{p}$ de $A$, il existe un idéal premier $\mathfrak{q}' \subset \mathfrak{p}'$ de $B$, tel que $\mathfrak{q}' \cap A = \mathfrak{q}$.

L'énoncé usuel de ce résultat est restreint au cas où $A$ est intégralement clos; mais la démonstration donnée par exemple dans le Sém. Cartan-Chevalley 1955-56, 1-09 et 1-10, s'étend sans difficulté au cas qui est envisagé ici.

Proposition 6. — Soient $A$ un anneau intègre noethérien, $K$ son corps des fractions, $B$ un anneau intègre contenant $A$ et fini sur $A$; soit $X = \operatorname{Spec}(B)$, $Y = \operatorname{Spec}(A)$. Alors le morphisme $f: X \to Y$ correspondant à l'injection $A \to B$ est ouvert.

En effet, le lemme 5 montre que la condition (ii) de I, 7,5, cor. 1 de la prop. 15 est vérifiée.

\newpage

%% ===== Page 427 =====

- III-99 -

\textbf{Corollaire 1.} — Soient $A$ un anneau intègre noethérien, $K$ son corps des fractions, $B$ un anneau intègre contenant $A$ et fini sur $A$ ; soit $A' = B \cap K$. Posons $X = \operatorname{Spec}(B)$, $Y = \operatorname{Spec}(A)$, $Y' = \operatorname{Spec}(A')$, et soient $f : X \to Y$, $g : X \to Y'$, $h : Y' \to Y$ les morphismes correspondant aux injections $A \to B$, $A' \to B$, $A \to A'$. Soit $y \in Y$ tel que $h^{-1}(y)$ soit réduit à un point $y'$ (resp. ait un nombre géométrique de points égal à 1, autrement dit (I, 4, 2) soit telle que $\kappa(y')$ soit radiciel sur $\kappa(y)$). Alors $f$ est ouvert (resp. universellement ouvert) au-dessus de $y$.

La prop. 6 montre que $g$ est un morphisme ouvert ; d'autre part, on sait que $h$ est une application fermée (II, 2, 7, prop. 16). Pour tout voisinage ouvert $U$ d'un point $x \in f^{-1}(y)$, $g(U)$ est donc ouvert, et par suite $h(Y' - g(U))$ est fermé dans $Y$ ; d'ailleurs $h$ est dominant, donc surjectif, et par suite $h(g(U)) = f(U)$ contient le complémentaire de $h(Y' - g(U))$ ; d'autre part $y$ n'appartient pas à $h(Y' - g(U))$, puisque $y' \in g(U)$ et que $h^{-1}(y) = \{y'\}$ ; donc $f(U)$ est un voisinage de $y$. 

Supposons en outre que $\kappa(y')$ soit radiciel sur $\kappa(y)$, et considérons un point $y_1 = Y' \otimes_Z \mathbb{Z}[T_1, \dots, T_r]$ au-dessus de $y$ ; posons $Y_1 = Y' \otimes_Z \mathbb{Z}[T_1, \dots, T_r]$ ; l'hypothèse entraîne que l'image réciproque de $y_1$ par le morphisme $h_1 : Y_1 \to Y_1$ est encore réduite à un seul point (I, 4, 2, prop. 14). D'autre part, $K(T_1, \dots, T_r)$ est le corps des fractions de $A[T_1, \dots, T_r]$, et il est clair que $B[T_1, \dots, T_r]$ est intègre et fini sur $A[T_1, \dots, T_r]$ ; en outre $B[T_1, \dots, T_r] \cap K(T_1, \dots, T_r) = A'[T_1, \dots, T_r]$, toute système linéaire à coefficients dans $K$ qui admet des solutions dans un surcorps de $K$ admettant aussi des solutions dans $K$. Si $X_1 = X \otimes_Z \mathbb{Z}[T_1, \dots, T_r]$, on peut donc appliquer au point $y_1$ et au morphisme $f_1 : X_1 \to Y_1$ la première partie du raisonnement, et par suite $f_1$ est ouvert au-

\newpage

%% ===== Page 428 =====

- III-100 -

\noindent
\textbf{Corollaire 2.} --- Soient $Y$ un préschéma irréductible localement noethérien, $X$ un préschéma irréductible, $f : X \to Y$ un morphisme fini dominant. Soit $y$ un point unibranche (resp. géométriquement unibranche) de $Y$ ; alors $f$ est ouvert (resp. universellement ouvert) au-dessus de $y$.

Remplaçant $f$ par $f_{\text{red}}$, on peut se ramener au cas où $X$ et $Y$ sont intègres (II, 2, 6, prop. 15), et comme la question est locale sur $Y$, on peut supposer $Y$ (et par suite $X$) affine. Tout anneau intermédiaire entre $\mathcal{O}_y$ et sa clôture intégrale est alors un anneau local (resp. un anneau local) dont le corps résiduel est radiciel sur $\kappa(y)$ ; les hypothèses du cor. 1 sont par suite remplies.

\textbf{Théorème 5.} --- Soient $Y$ un préschéma localement noethérien, $f : X \to Y$ un morphisme de type fini, et un point $y$ de $Y$.

(i) Supposons $X$ et $Y$ irréductibles. Si $f$ est ouvert au-dessus de $y$, $f$ est équidimensionnel en $y$.

(ii) Supposons le point $y$ géométriquement unibranche. Si $f$ est équidimensionnel en $y$, il est ouvert au-dessus en $y$.

L'assertion (i) a déjà été démontrée dans I, 7, 5, prop. 16, compte tenu du critère (ii) du th. 4 pour que $f$ soit équidimensionnel en $y$.

Pour prouver (ii), il suffit de montrer que $f$ est ouvert au-dessus de $y$. Si ce point est établi, et en effet, si on pose $S = \mathbb{Z}[T_1, \dots, T_r]$, $f^S : X^S \to Y^S$ est équidimensionnel en tout point $y' \in Y^S$ au-dessus de $y$ (cor. 3 du th. 4), et d'autre part $y'$ est géométriquement unibranche (lemme 4), donc $f^S$ est ouvert au-dessus de $y'$, et cela montre que $f$ est universellement ouvert au-dessus de $y$ par le raisonnement de I, 7, 5, cor. 2 de la prop. 15. La question étant locale en $X$ (pour un point de $f^{-1}(y)$), on peut,

\newpage

%% ===== Page 429 =====

- III-101 -

en vertu du th.4 (iii), supposer que $X$ est fini sur $Y' = Y \otimes_{\mathbb{Z}} \mathbb{Z}[T_1, \dots, T_e]$. Si $h$ est le morphisme $Y' \to Y$, pour toute partie fermée irréductible $Z$ de $Y$, $h^{-1}(Z)$ est irréductible et son image dans $Y$ est dense dans $Z$, car on peut considérer le sous-préschéma fermé réduit de $Y$ ayant $Z$ pour espace sous-jacent (I, 3, 4, prop. 29), et il est clair alors (en se ramenant au cas où $Y$ est affine) que $Z \otimes_{\mathbb{Z}} \mathbb{Z}[T_1, \dots, T_e]$ est intègre et que $h$ transforme le point générique de ce schéma en le point générique de $Z$. On peut donc appliquer au morphisme $h$ le critère de I, 7, 5, cor. 1 de la prop. 15, et cela prouve que $h$ est ouvert. D'autre part, tout point $y'$ de $Y'$ au-dessus de $y$ est géométriquement unibranche. D'après le th. 4 (iii) et l'hypothèse sur $f$, on peut en outre supposer que le morphisme $g : X \to Y'$ est tel que pour ce morphisme, toute composante irréductible de $X$ domine une composante irréductible de $Y'$. Alors, le cor. 2 de la prop. 6 montre que la restriction de $g$ à toute composante irréductible de $X$ est ouvert au-dessus de $y'$; on en conclut aussitôt que $g$ est ouvert au-dessus de $y'$, donc que $f=h \circ g$ est ouvert au-dessus de $y$.

Corollaire 1. — Les hypothèses sur $X, Y$ et $f$ étant celles du th. 5 :

(i) Si $X$ et $Y$ sont irréductibles et si $f$ est ouvert, $f$ est équidimensionnel.

(ii) Si tout point de $Y$ est géométriquement unibranche et si $f$ est équidimensionnel, $f$ est universellement ouvert.

Corollaire 2. — Les hypothèses sur $X, Y, f$ étant celles du th. 5 :

(i) Si $X$ et $Y$ sont irréductibles et $y \in Y$ géométriquement unibranche, les conditions suivantes sont équivalentes :

a) $f$ est universellement ouvert en $y$ ; b) $f$ est ouvert en $y$ ; c) $f$ est équidimensionnel en $y$.

\newpage

%% ===== Page 430 =====

--- III-102 ---

(ii) Si $X$ et $Y$ sont irréductibles et tout point de $Y$ géométriquement unibranche, les conditions suivantes sont équivalentes :

a) $f$ est universellement ouvert ; b) $f$ est ouvert ; c) $f$ est équidimensionnel .

On notera que ces résultats s'appliquent \emph{a fortiori} lorsque $Y$ est normal .

Remarques .— 1) Le résultat (ii) du th.5 a été démontré par C. Chevalley pour les schémas algébriques, en supposant $Y$ normal .

2) L'hypothèse que $y$ est géométriquement unibranche est nécessaire pour que l'assertion (ii) du th.5 soit valable pour tout morphisme de type fini équidimensionnel en $y$, et on peut définir un préschéma $Y'$ fini sur $Y$ et tel que $f$ ne soit pas universellement ouvert au-dessus d'un point $y' \in Y'$ au-dessus de $y$. Si $y$ appartient à plus d'une composante irréductible de $Y$, il suffit de prendre pour $X$ une de ces composantes, pour $f$ le morphisme d'injection d'un sous-préschéma fermé ayant $X$ pour espace sous-jacent ; il est clair déjà ici que $f$ n'est même pas ouvert au-dessus de $y$. Si $y$ n'appartient qu'à une seule composante irréductible de $Y$, on peut, comme la question est locale sur $Y$, se borner au cas où $Y$ est irréductible, et en remplaçant $Y$ par $Y_{\text{red}}$ (ce qui laisse le point $y$ non géométriquement unibranche, par définition), on peut supposer $Y$ intègre et affine. Soit donc $Y = \operatorname{Spec}(A)$, où $A$ est noethérien et intègre, de corps des fractions $K$ ; par hypothèse, il existe un sous-anneau $A' \supset A$ (fini sur $A$ et contenu dans $K$) qui ne soit pas un anneau local, ou qui soit un anneau local dont le corps résiduel n'est pas radiciel sur le corps résiduel $k$ de $A_y$. Il est immédiat qu'il y a un système fini de générateurs de $A'$ en tant que $A$-module.

\newpage

%% ===== Page 431 =====

- III-103 -

dule, qui sont de la forme $b_i/c$, où les $b_i$ sont entiers sur $A$ et $c \in A - \mathfrak{p}_y$ ; soit $B$ le sous-anneau de $K$, fini sur $A$, engendré par les $b_i$. Nous prendrons $Y' = \operatorname{Spec}(B)$, de sorte que le morphisme $f : X \to Y$ est fini, dominant et birationnel ; en outre, l'hypothèse signifie qu'ou bien il y a deux points distincts $x_1, x_2$ de $X$ au-dessus de $y$, ou bien il n'y a qu'un seul point $x$ au-dessus de $y$, mais $\kappa(x) = L$ n'est pas radiciel sur $\kappa(y) = k$. Dans les deux cas, nous prendrons $Y' = X$, le morphisme $g : Y' \to Y$ étant égal à $f$. Notons d'abord que si $\eta$ est le point générique de $Y$, la fibre de $X \times_Y X$ au-dessus de $\eta$ est le spectre de $(B \otimes_A B) \otimes_A K( \eta) = (B \otimes_A B) \otimes_A K$ ; mais comme $B \otimes_A K$ s'identifie canoniquement à $K$, il en est de même de $(B \otimes_A B) \otimes_A K = B \otimes_A K$. On en conclut que cette fibre se réduit à un point $\xi$, qui est point générique d'une composante irréductible de $X \times_Y X$ (I, 2, 7, prop. 19). En outre comme il existe toujours au moins un point de la diagonale $Z$ de $X \times_Y X$ au-dessus de $\eta$, ce point ne peut être que $\xi$, et puisque $Z$ est fermé dans $X \times_Y X$ (le morphisme $f$ étant séparé), et isomorphe à $X$, donc irréductible, on en conclut que $Z$ est une composante irréductible de $X \times_Y X$. Cela étant, supposons d'abord qu'il y ait deux points distincts $x_1, x_2$ de $X$ au-dessus de $y$ ; alors il existe un point $t \in X \times_Y X$ dont la première projection est $x_1$ et la seconde $x_2$ (I, 2, 6, prop. 10) ; il est clair que $t \notin Z$, et par suite, si $Z'$ est une composante irréductible de $X \times_Y X$ contenant $t$, cette composante ne contient pas $\xi$, et par suite sa seconde projection sur $X$ n'est pas dense. Or cette seconde projection n'est autre que ce qui a été noté en général $f^{Y'}(t)$ et comme $f^{Y'}(t) = x_2$ est au-dessus de $y$, on voit que $f$ n'est pas universellement ouvert au-dessus de $y$, en vertu de I, 7, 5, cor. 4 de la prop. 15.

\newpage

%% ===== Page 432 =====

- III-104 -

Si maintenant $f^{-1}(y)$ est réduit à un seul point $x$, la fibre $(f^{Y'})^{-1}(x)$ est le spectre premier de $(B \otimes_A L) \otimes_A^L (L \otimes_A L)$, et par hypothèse $[L:k]_s > 1$. Comme $L \otimes_A L = L \otimes_k (L \otimes_A L) = (L \otimes_k L) \otimes_A K$, et que le spectre de $L \otimes_k L$ comporte au moins deux points distincts, il en est de même du spectre de $L \otimes_A L$, et par suite $(f^{Y'})^{-1}(x)$ a au moins 2 points distincts. L'un de ces points n'est donc pas sur la diagonale, et on peut conclure le raisonnement comme ci-dessus.

3) Il serait intéressant, pour un morphisme $f: X \to Y$ donné, d'avoir une condition nécessaire et suffisante pour que $f$ soit universellement ouverte au-dessus d'un point $y \in Y$. Si $Y = \operatorname{Spec}(A)$ est intègre noethérien, $X = \operatorname{Spec}(B)$ intègre, fini sur $Y$, on a la condition nécessaire et suffisante suivante : si $K$ est le corps des fractions de $A$, $A' = B \cap K$ et $Y' = \operatorname{Spec}(A')$, il faut et il suffit que, si on désigne par $h$ le morphisme $Y' \to Y$, $h^{-1}(y)$ ait un nombre géométrique de points égal à 1. En effet, la condition est suffisante en vertu du cor. 1 de la prop. 6, et on voit qu'elle est nécessaire comme dans la Remarque 2, en considérant $X^{Y'} = (Y' \times_Y Y') \times_{Y'} X$.

Lorsqu'on ne suppose plus $B$ entier sur $A$, on peut désigner par $A'$ la clôture intégrale de $A$ dans $B \cap K$, et se demander si la condition précédente (jointe à l'hypothèse que $f$ est équidimensionnel en $y$) est encore suffisante pour que $f$ soit universellement ouverte au-dessus de $y$ ; le même raisonnement montre toujours que la condition est nécessaire.

\newpage

%% ===== Page 433 =====

\addtocounter{page}{103}
\noindent -- III-104 --

Si maintenant $f^{-1}(y)$ est réduit à un seul point $x$, la fibre $(f^{Y^0})^{-1}(x)$ est le spectre premier de $(B \otimes_A L) \otimes_{L \otimes_A L} (L \otimes_A L)$, et par hypothèse $[L:k]_s > 1$. Comme $L \otimes_A L = (L \otimes_A L) \otimes_A k$, et que le spectre de $L \otimes_A L$ comporte au moins deux points distincts, il en est de même du spectre de $L \otimes_A L$, et par suite $(f^{Y^0})^{-1}(x)$ a au moins 2 points distincts. L'un de ces points n'est donc pas sur la diagonale, et on peut conclure le raisonnement comme ci-dessus.

3) Il serait intéressant, pour un morphisme $f: X \to Y$ donné, d'avoir une condition nécessaire et suffisante pour que $f$ soit universellement ouverte au-dessus d'un point $y \in Y$. Si $Y = \operatorname{Spec}(A)$ est intègre noethérien, $X = \operatorname{Spec}(B)$ intègre, fini sur $Y$, on a la condition nécessaire et suffisante suivante : si $K$ est le corps des fractions de $A$, $A' = B \otimes_A K$ et $Y' = \operatorname{Spec}(A')$, il faut et il suffit que, si on désigne par $h$ le morphisme $Y' \to Y$, $h^{-1}(y)$ ait un nombre géométrique de points égal à 1. En effet, la condition est suffisante en vertu du cor. 1 de la prop. 6, et on voit qu'elle est nécessaire comme dans la Remarque 2, en considérant $X^{Y^0} = (Y^0 \times_Y Y') \times_{Y'} X$.

Lorsqu'on ne suppose plus $B$ entier sur $A$, on peut désigner par $A'$ la clôture intégrale de $A$ dans $B \otimes_A K$, et se demander si la condition précédente (jointe à l'hypothèse que $f$ est équidimensionnel en $y$) est encore suffisante pour que $f$ soit universellement ouverte au-dessus de $y$ ; le même raisonnement montre toujours que la condition est nécessaire.

\newpage

%% ===== Page 434 =====

\marginpar{Archives\\Grothendieck sept 59}

\hfill --- III-1 ---

P. III-8, after line 1, add:

\emph{Remarque}. — On notera que la conclusion de la prop. 3 est encore valable lorsqu'on suppose seulement que les intersections finies des $U_\alpha$ sont des ouverts affines.

P. III-9, before n$^\circ$5, insert:

\emph{Proposition 5}. — Soient $u : X \to Y$ un morphisme séparé de type fini, $Y'$ un schéma local d'anneau $B$, $\psi : Y' \to Y$ un morphisme dont la factorisation canonique $Y' \to \operatorname{Spec}(\mathcal{O}_y) \to Y$ (I, 2, 2) est telle que $B$ soit un $\mathcal{O}_y$-module plat. Alors, pour tout faisceau cohérent $\mathcal{F}$ sur $X$, $R^q u'_* (\mathcal{F} \otimes_{\mathcal{O}_Y} \mathcal{O}_{Y'})$ est canoniquement isomorphe à $R^q u_*(\mathcal{F}) \otimes_{\mathcal{O}_Y} \mathcal{O}_{Y'}$ pour tout $q \ge 0$.

On peut évidemment supposer $Y = \operatorname{Spec}(A)$ affine et $X$ réunion finie d'ouverts affines $U_i$ ($1 \le i \le r$) ; comme $\psi$ est affine, il en est de même de $v : X' = X \times_Y Y' \to X$ (II, 2, 6, prop. 12) et les $U'_i = v^{-1}(U_i)$ forment donc un recouvrement ouvert affine $U'$ de $X'$. Conservons les notations de la démonstration de la prop. 4, et posons en outre $\mathcal{F}' = \mathcal{F} \otimes_{\mathcal{O}_Y} B$, $U'_I = U'_{i_0} \cap \dots \cap i_p = v^{-1}(U_I)$ ; compte tenu de la Remarque suivante la prop. 3 et du cor. de la prop. 4, il suffit de montrer que $H^q(U', \mathcal{F}')$ est canoniquement isomorphe à $H^q(U, \mathcal{F}) \otimes_A B$ ; or $M'_I$ s'identifie à $M_I \otimes_A B$ (II, 1, 1, prop. 7) ; donc $C^p(U', \mathcal{F}')$ s'identifie canoniquement à $C^p(U, \mathcal{F}) \otimes_A B$, et l'opérateur cobord $C^p(U', \mathcal{F}') \to C^{p+1}(U', \mathcal{F}')$ à $d \otimes 1$. Mais comme $A_I$ est un $A$-module plat, et $B$ un $A$-module plat, il résulte bien des définitions que $H^q(U', \mathcal{F}')$ s'identifie à $H^q(U, \mathcal{F}) \otimes_A B$.

Ce résultat sera généralisé au § 4 (n$^\circ$ ).

\newpage

%% ===== Page 435 =====

- III-iii -

niquement de $u$ des homomorphismes $\Gamma(U, \mathcal{F} \otimes \mathcal{O}_X/\mathcal{J}^n) \to \Gamma(U, \mathcal{G} \otimes \mathcal{O}_X/\mathcal{J}^n)$ pour tout $n$, et il est immédiat que ces homomorphismes forment un système projectif, d'où la limite un homomorphisme $\Gamma(U, \hat{\mathcal{F}}) \to \Gamma(U, \hat{\mathcal{G}})$. En outre, ces homomorphismes sont compatibles avec les opérations de restriction, et donnent donc un homomorphisme de $\mathcal{O}_X/\mathcal{I}$-modules $u_{\mathcal{I}} : \hat{\mathcal{F}} \to \hat{\mathcal{G}}$, tel que le diagramme
$$\begin{tikzcd}
\mathcal{F} \arrow[r] \arrow[d, "u"] & \hat{\mathcal{F}} \arrow[d, "u_{\mathcal{I}}"] \\
\mathcal{G} \arrow[r] & \hat{\mathcal{G}}
\end{tikzcd}$$
soit commutatif (les lignes étant les homomorphismes canoniques).
On écrira parfois $\hat{\mathcal{F}}$ au lieu de $\hat{\mathcal{F}}$ et $u_{\mathcal{I}}$ au lieu de $u_{\mathcal{I}}$.

P. III-55, before prop.4, add :

En outre il résulte de cette définition que si $\rho : \mathcal{O}_X \to (\mathcal{O}_X)_{\mathcal{I}}$ et $\sigma : \mathcal{O}_Y \to (\mathcal{O}_Y)_{\mathcal{J}}$ sont les homomorphismes canoniques, le diagramme
$$\begin{tikzcd}
\psi^*(\mathcal{O}_Y) \arrow[d, "\psi^*(\sigma)"] \arrow[r] & \mathcal{O}_X \arrow[d, "\rho"] \\
\psi^*((\mathcal{O}_Y)_{\mathcal{J}}) \arrow[r, "\omega"] & (\mathcal{O}_X)_{\mathcal{I}}
\end{tikzcd}$$
est commutatif.

\newpage

%% ===== Page 436 =====

\noindent --- III-iv ---

P.III-50 , before section 3 , insert :

2. Schémas complets .

On dit qu'un préschéma $X$ est complet s'il existe un anneau artinien $A$ et un morphisme propre $f : X \to \operatorname{Spec}(A)$ ; il est clair que $X$ est alors un schéma noethérien . Si $X$ est complet et $Z$ propre au-dessus de $X$ , il est clair que $Z$ est complet (II,5,1,prop.2 (ii)) ; en particulier tout sous-préschéma fermé d'un schéma complet est complet . Pour qu'un schéma noethérien $X$ soit complet , il faut et il suffit que ses composantes irréductibles le soient (II,5,1,cor.2 de la prop.2). Comme toute algèbre sur un anneau artinien $A$, qui est un $A$-module de type fini, est artinienne, les schémas affines complets ne sont autres que les schémas artiniens (II,5,1,cor.2 de la prop.3). Proposition 0. — Si $X$ est un schéma complet , l'anneau $\Gamma(X,\mathcal{O}_X)$ est artinien ; si $X_i$ ($1 \le i \le n$) sont les composantes connexes de $X$ , les anneaux $\Gamma(X_i, \mathcal{O}_{X_i})$ s'identifient aux composantes locaux de l'anneau $\Gamma(X, \mathcal{O}_X)$.

Comme $X$ est propre au-dessus de $Y=\operatorname{Spec}(A)$, où $A$ est un anneau artinien , $\Gamma(X, \mathcal{O}_X) = H^0(X, \mathcal{O}_X)$ est un $A$-module de type fini (n°1,cor.2 du th.1) , d'où la première assertion . Posons $B = \Gamma(X, \mathcal{O}_X)$ et $Z=\operatorname{Spec}(B)$; l'isomorphisme identique de $B$ sur lui-même définit un morphisme $g : X \to Z$ (I,2,2,prop.3) ; comme le morphisme $g : Z \to Y$ est affine , donc séparé , $g$ est propre (II,5,1,prop.2 (v)) . D'autre part , $Z$ est fini et discret (I,4,1,prop.4) ; si $z_i$ ($1 \le i \le m$) sont les points de $Z$, $X$ est somme des sous-préschémas induits sur les $X_i = g^{-1}(z_i)$ , et les morphismes $X_i \to \{z_i\}$ correspondent aux applications identiques $\mathcal{O}_{z_i} \to \Gamma(X_i, \mathcal{O}_{X_i})$. Tout revient donc à prouver que lorsque $B$ est un anneau local , $X$ est connexe , ce qui est immédiat , $B$ ne pouvant être composé direct de deux sous-anneaux non réduits à $0$.

\newpage

%% ===== Page 437 =====

\begin{center}
--- III-iv ---
\end{center}

P.III-50 , before section 3 , insert :

2. \emph{Schémas complets} .

On dit qu'un préschéma $X$ est complet s'il existe un anneau artinien $A$ et un morphisme propre $f : X \to \operatorname{Spec}(A)$ ; il est clair que $X$ est alors un schéma noethérien . Si $X$ est complet et $Z$ propre au-dessus de $X$ , il est clair que $Z$ est complet (II,5,1,prop.2 (ii)) ; en particulier tout sous-préschéma fermé d'un schéma complet est complet. Pour qu'un schéma noethérien $X$ soit complet , il faut et il suffit que ses composantes irréductibles le soient (II,5,1,cor.2 de la prop.2). Comme toute algèbre sur un anneau artinien $A$, qui est un $A$-module de type fini, est artinienne, les schémas affines complets ne sont autres que les schémas artiniens (II,5,1,cor.2 de la prop.3) .

Proposition 0 . - Si $X$ est un schéma complet , l'anneau $\Gamma(X,\mathcal{O}_X)$ est artinien ; si $X_i$ ($1 \le i \le n$) sont les composantes connexes de $X$ , les anneaux $\Gamma(X_i, \mathcal{O}_{X_i})$ s'identifient aux composantes locaux de l'anneau $\Gamma(X,\mathcal{O}_X)$ .

Nomme $X$ est propre au-dessus de $Y=\operatorname{Spec}(A)$ , où $A$ est un anneau artinien , $\Gamma(X,\mathcal{O}_X)=H^0(X,\mathcal{O}_X)$ est un $A$-module de type fini (n°1, cor.2 du th.1) , d'où la première assertion . Posons $B=\Gamma(X,\mathcal{O}_X)$ et $Z=\operatorname{Spec}(B)$ ; l'isomorphisme identique de $B$ sur lui-même définit un morphisme $g: X \to Z$ (I,2,2,prop.3) ; comme le morphisme $g: X \to Y$ est affine, donc séparé , $g$ est propre (II,5,1,prop.2 (v)) . D'autre part , $Z$ est fini et discret (I,4,1,prop.4) ; si $z_i$ ($1 \le i \le m$) sont les points de $Z$ , $X$ est somme des sous-préschémas induits sur les $X_i=g^{-1}(z_i)$ , et les morphismes $X_i \to \{z_i\}$ correspondent aux applications identiques $\mathcal{O}_{z_i} \to \Gamma(X_i, \mathcal{O}_{X_i})$ . Tout revient donc à prouver que lorsque $B$ est un anneau local , $X$ est connexe , ce qui est immédiat , $B$ ne pouvant être composé direct de deux sous-anneaux non réduits à 0 .
